% !TeX program = pdflatex
% Single-source manuscript. Compile this file twice with pdflatex.
% All proofs, references, and essential finite certificates are inline.
% Source archive: Math_Tournament_Rounds_26-50(1).zip
\documentclass[11pt,a4paper]{article}
\usepackage[T1]{fontenc}
\usepackage[utf8]{inputenc}
\usepackage{lmodern,amsmath,amssymb,amsthm}
\usepackage[a4paper,top=23mm,bottom=24mm,left=24mm,right=24mm,headheight=15pt,headsep=8mm]{geometry}
\usepackage{xcolor,microtype,booktabs,longtable,enumitem,needspace}
\usepackage{fancyhdr,titlesec,fvextra}
\usepackage[unicode=true,pdfencoding=auto,psdextra]{hyperref}
\definecolor{inkblue}{HTML}{193B58}
\definecolor{muted}{HTML}{52616B}
\definecolor{lightline}{HTML}{CED8E0}
\hypersetup{colorlinks=true,linkcolor=inkblue,citecolor=inkblue,urlcolor=inkblue,
  pdftitle={Expanded Complete Proofs from Mathematics Tournament Rounds 26--50},
  pdfauthor={Collected tournament submissions; expanded edition prepared for Keio},
  pdfsubject={Collected complete statements, proofs, and exact certificates}}
\urlstyle{same}
\setlength{\parindent}{1.15em}
\setlength{\parskip}{3pt plus 1pt minus 1pt}
\setlength{\emergencystretch}{2em}
\setlength{\jot}{5pt}
\linespread{1.015}
\allowdisplaybreaks[2]
\raggedbottom
\setcounter{tocdepth}{1}
\setcounter{secnumdepth}{2}
\setlist{topsep=4pt,itemsep=2pt,parsep=1pt}
\numberwithin{equation}{section}
\newtheorem{theorem}{Theorem}[section]
\newtheorem{thm}[theorem]{Theorem}
\newtheorem{lemma}[theorem]{Lemma}
\newtheorem{proposition}[theorem]{Proposition}
\newtheorem{prop}[theorem]{Proposition}
\newtheorem{corollary}[theorem]{Corollary}
\theoremstyle{definition}
\newtheorem{definition}[theorem]{Definition}
\theoremstyle{remark}
\newtheorem{remark}[theorem]{Remark}
\newenvironment{sectionabstract}{\par\smallskip\begingroup\small\noindent}{\par\endgroup\medskip}
\fvset{fontsize=\scriptsize,baselinestretch=1,breaklines=true,breakanywhere=true,
  breakindent=1em,breaksymbolleft={},breaksymbolright={},tabsize=4,
  xleftmargin=2mm,xrightmargin=2mm,frame=leftline,rulecolor=\color{lightline},framesep=2mm}
\titleformat{\part}[display]{\centering\color{inkblue}\sffamily\bfseries}{\large PART \thepart}{8pt}{\LARGE}
\titleformat{\section}{\Large\color{inkblue}\sffamily\bfseries}{\thesection}{0.65em}{}
\titleformat{\subsection}{\large\color{inkblue}\sffamily\bfseries}{\thesubsection}{0.65em}{}
\titleformat{\subsubsection}{\normalsize\sffamily\bfseries}{\thesubsubsection}{0.65em}{}
\titlespacing*{\section}{0pt}{20pt}{9pt}
\titlespacing*{\subsection}{0pt}{13pt}{5pt}
\titlespacing*{\paragraph}{0pt}{8pt}{0.65em}
\pagestyle{fancy}
\fancyhf{}
\fancyhead[L]{\small\sffamily\color{muted}Complete proofs | Rounds 26--50}
\fancyhead[R]{\small\sffamily\color{muted}Expanded edition}
\fancyfoot[C]{\small\sffamily\thepage}
\renewcommand{\headrulewidth}{0.3pt}
\makeatletter
\renewcommand*\l@section{\@dottedtocline{1}{0em}{2.8em}}
\renewenvironment{thebibliography}[1]
 {\subsection*{References for this section}\small
  \list{\@biblabel{\@arabic\c@enumiv}}
       {\settowidth\labelwidth{\@biblabel{#1}}\leftmargin\labelwidth
        \advance\leftmargin\labelsep\@openbib@code
        \usecounter{enumiv}\let\p@enumiv\@empty
        \renewcommand\theenumiv{\@arabic\c@enumiv}}
  \sloppy\clubpenalty4000\widowpenalty4000\sfcode`\.=1000\relax}
 {\def\@noitemerr{\@latex@warning{Empty bibliography}}\endlist}
\makeatother
\begin{document}
\hypersetup{pageanchor=false}
\begin{titlepage}
\thispagestyle{empty}
\vspace*{19mm}
{\sffamily\color{muted}\large MATHEMATICS TOURNAMENT\par}
\vspace{8mm}
{\sffamily\bfseries\color{inkblue}\fontsize{34}{39}\selectfont Complete Proofs\par}
\vspace{4mm}
{\sffamily\color{inkblue}\fontsize{25}{30}\selectfont Rounds 26--50\par}
\vspace{10mm}
{\color{lightline}\hrule height 1pt}
\vspace{8mm}
{\Large Complete statements, proofs, and exact certificates\par}
\vspace{5mm}
{\sffamily\color{inkblue}Expanded edition: 7 additional completions\par}
\vspace{15mm}
{\large Collected from the Meiji, Taisho, Showa,\par Heisei, Reiwa, and Keio submissions\par}
\vfill
{\sffamily\color{muted}Prepared for Keio\par}
\vspace{2mm}
{\sffamily 20 November 2026\par}
\end{titlepage}
\pagenumbering{roman}
\hypersetup{pageanchor=true}
\section*{Editorial scope}
\addcontentsline{toc}{section}{Editorial scope}
This expanded single-source manuscript preserves the 106 complete-proof
sections of the preceding edition and adds 7 self-contained completions
of targets left unfinished in the supplied Mathematics Tournament
Rounds 26--50 archive. The added proofs are collected in Part VI. Every included result is stated
on its exact proved domain. A separately formulated clause is identified as that
clause; its inclusion does not assert that other clauses of the same source
problem have been resolved. The selection includes complete classifications and
complete proofs of previously known results, with their prior sources retained.

Unfinished reductions, bounded numerical evidence, unproved global extensions,
and isolated counterexamples have been omitted. Repeated proofs of the same result
are consolidated where practical. When a stronger theorem needs an earlier
argument, both are retained as part of the complete proof. In particular, the
sharp refinement accompanying Round 41 is placed with its parent result.

The mathematical arguments and essential exact computational certificates are
included in full. A certificate using rational or integer interval arithmetic
establishes its assertion on every accepted box; finite point sampling is not
used as a replacement for such a proof. All essential certificate listings used
in this edition were replayed successfully during compilation. The listings are
part of this file and require no external source-file inclusion. Standard
Python is sufficient unless a listing explicitly imports SymPy.

Each main section records its round and originating seat. Notation and reference
numbers are local to that section; automatically numbered equations and theorems
use the unified section numbering. Explicit equation tags remain local.
Cited published theorems remain mathematical inputs and are
identified where used. Editorial mathematical review and certificate replay are
not claims of proof-assistant verification, external peer review, or historical
priority. The compilation does not convert earlier work into a claim of novelty.

\medskip
\noindent\textbf{Contents of this edition.} 113 main sections,
comprising the preceding 106 sections (107 source units) and 7 additional
completion manuscripts. Some sections prove several
separate statements; these counts are editorial units, not a count of newly
solved conjectures.

\medskip
\noindent\textbf{Source and reproducibility.} The source archive is
\texttt{Math\_Tournament\_Rounds\_26-50(1).zip}. The accompanying
\texttt{.tex} file is the complete source for this PDF. It is designed to compile
with a standard modern \LaTeX{} installation using \texttt{pdflatex} twice.
Source-path provenance is also recorded in comments in the \texttt{.tex} file.


\clearpage
\section*{Additional completions in this edition}
\addcontentsline{toc}{section}{Additional completions in this edition}
Every item below includes the full statement and proof in Part VI. Earlier
partial arguments are restated wherever they are needed, so the expanded
edition is usable without a separate supplement. ``Additional'' refers to
completion relative to the supplied archive, not to a claim of historical
priority. In particular, the composition result explicitly recovers an
established method and conclusion from the literature.

\medskip
\begin{longtable}{@{}p{24mm}p{130mm}@{}}
\toprule
Origin & Completed statement and proof method\\
\midrule
\endhead
Round 34 / Keio & Liu's two-product pedal-triangle inequality, for every triangle and every interior point, including equality.\par\smallskip{\small\color{muted}Six complete barycentric chambers; exact integer polynomial certificates.}\\[7pt]
Round 43 / Taisho & Equal-length composition nonalgebraicity, for all rational nonnegative-integer alphabets with at least two sizes and every tuple size at least three.\par\smallskip{\small\color{muted}Lattice local limits and the established arithmetic Puiseux obstruction; recovered known method.}\\[7pt]
Round 43 / Reiwa & Odd primitive sign density for the archived Pathak invariant: a continuous limiting law, both signs of positive density, and convergent rational bounds for the density.\par\smallskip{\small\color{muted}Mean-square approximation, prime-flip anti-concentration, and exact finite-probability bounds.}\\[7pt]
Round 47 / Taisho & Ferrers-pattern nonalgebraicity for every strict pattern with adjacent positive parts, with any number of unit gaps.\par\smallskip{\small\color{muted}An exact prefix decomposition and a dominant logarithmic term prevent cancellation.}\\[7pt]
Round 47 / Heisei & Strict decrease of Lyon's scaled gamma-median offset on the entire positive shape axis.\par\smallskip{\small\color{muted}Analytic tails and 1,344 rigorously bounded compact cells; complete directed-integer verifier.}\\[7pt]
Round 49 / Heisei & The Chernoff density's logarithmic curvature attains its unique global minimum at zero.\par\smallskip{\small\color{muted}Analytic tails, explicitly bounded Fourier errors, and a complete integer interval certificate.}\\[7pt]
Round 49 / Keio & The Liu--Chu median-to-altitude Fermat-distance inequality, for every triangle and every point of its plane, including equality.\par\smallskip{\small\color{muted}Five complete half-angle cones; exact integer polynomial certificates.}\\[7pt]

\bottomrule
\end{longtable}

\noindent\textbf{Exact scope.} The sign-density result concerns all negative
fundamental discriminants, the odd primitive population formulated in its
archived target. Its density is evaluated by converging, explicitly specified
rational bounds; an elementary closed form or an uncertified decimal is not
asserted. Other separate populations and uncompleted targets are not promoted
to solved statements.

\clearpage
\tableofcontents
\clearpage
\pagenumbering{arabic}
\part{Rounds 26--30}

% Source: /workspace/scratch/81914cec2c78/work/archive/Tournament_Round_26_Final/Tournament_Round_26/round26/meiji/paper.tex
% Edited source: /workspace/scratch/81914cec2c78/work/archive/Tournament_Round_26_Final/Tournament_Round_26/round26/meiji/paper.tex
\begingroup

\Needspace{9\baselineskip}
\section{\texorpdfstring{Two conjectures on adjacent Taylor remainders of the inverse trigonometric functions}{Two conjectures on adjacent Taylor remainders of the inverse trigonometric functions}}\label{section:r26meiji}
\noindent{\small\sffamily\color{muted}Round 26 \enspace / \enspace Meiji}\par\medskip


\begin{sectionabstract}
For each nonnegative integer $n$, the ratio of the $2n$th to the
$(2n+2)$nd normalized remainder of the inverse sine is strictly
concave on $(-1,1)$. For the inverse tangent, the corresponding
ratio is strictly increasing on $(0,\infty)$. These are the eighth and
ninth conjectures in Feng Qi's survey on normalized remainders. Both
follow from the same moment transform. Its adjacent-moment ratio
has negative first and second derivatives; the two applications use
respectively $z=x^2$ and $z=-x^2$.
\end{sectionabstract}

\subsection{Statements and normalization}
For an odd analytic function $f(x)=\sum_{j\ge0}c_jx^{2j+1}$
near zero with $c_n\ne0$, write
\[
 T_{2n}[f](x)=\frac{f(x)-\sum_{j=0}^{n-1}c_jx^{2j+1}}
                       {c_nx^{2n+1}},\qquad T_{2n}[f](0)=1.
\]
The definition by the finite remainder is valid wherever the function is
defined, even away from the interval of convergence of its series.
The empty sum for $n=0$ is zero. Qi \cite[Secs.~12.7.1.8--9]{r26meiji:Qi}
asks for the two assertions below in their non-strict forms.
The arctangent assertion also appears, with the reciprocal ratio, in
Liu and Qi \cite[Remark~5.1]{r26meiji:LQ}.

\begin{theorem}\label{r26meiji:main}
For every integer $n\ge0$:
\begin{enumerate}
\item $T_{2n}[\arcsin](x)/T_{2n+2}[\arcsin](x)$ has strictly negative
second derivative at every $x\in(-1,1)$.
\item $T_{2n}[\arctan](x)/T_{2n+2}[\arctan](x)$ has strictly positive
first derivative at every $x>0$.
\end{enumerate}
\end{theorem}
Part 1 is a statement about the reciprocal of the ratio treated in
\cite[Theorem~5.1]{r26meiji:LQ}; the absolute monotonicity proved there for the
opposite ratio does not by itself imply this concavity assertion.

\subsection{An adjacent-moment quotient}
Let $\mu$ be a finite positive measure on $[0,1]$ whose support contains
two distinct strictly positive points. For an integer $j\ge0$ put
\[
 m_j=\int s^j\,d\mu(s),\qquad
 A_j(z)=\int\frac{s^j}{1-zs}\,d\mu(s),\quad -\infty<z<1.
\]
These functions are strictly positive. On any compact subinterval of
$(-\infty,1)$ all differentiated integrands of any fixed order are
bounded uniformly in $s$; differentiation under the integral is valid.

\begin{lemma}\label{r26meiji:lem}
$B_j=A_j/A_{j+1}$ satisfies $B_j'<0$ and $B_j''<0$ on $(-\infty,1)$.
\end{lemma}
\begin{proof}
The relation $A_j=m_j+zA_{j+1}$ yields $B_j=z+m_j/A$ with
$A=A_{j+1}$. Use the measure $d\nu=s^j d\mu$ and set
$h(s)=s/(1-zs)$. Then
\[
 m_j=\int 1\,d\nu,\quad A=\int h\,d\nu,\quad
 A'=\int h^2\,d\nu,\quad A''=2\int h^3\,d\nu .
\]
Cauchy--Schwarz, applied first to $1,h$ and then to
$h^{1/2},h^{3/2}$, gives $A^2<m_jA'$ and $2(A')^2<AA''$.
Both inequalities are strict: $h$ is strictly increasing on $[0,1]$
and positive on the two distinct positive support points, which
remain in the support of $\nu$. Differentiation gives
\[
 B_j'=1-\frac{m_j A'}{A^2}<0,\qquad
 B_j''=m_j\frac{2(A')^2-AA''}{A^3}<0 .
\]
\end{proof}

\subsection{The two measures}
\paragraph{Inverse sine.}
The binomial expansion of $(1-x^2)^{-1/2}$ and integration from
zero to $x$ give $c_j=m_j=4^{-j}\binom{2j}{j}/(2j+1)$.
Take independent variables $U$ uniform on $[0,1]$ and $V$ uniform on
$[0,\pi/2]$. The law $\mu$ of $U^2\sin^2 V$ has these moments:
$\mathbb E U^{2j}=1/(2j+1)$ and
$\mathbb E\sin^{2j}V=4^{-j}\binom{2j}{j}$ by single-step integration
by parts and induction. This law has support $[0,1]$.
For $|x|<1$ the geometric series now identifies
\[
 T_{2n}[\arcsin](x)=\frac{A_n(x^2)}{m_n},\qquad
 \frac{T_{2n}[\arcsin](x)}{T_{2n+2}[\arcsin](x)}
 =\frac{m_{n+1}}{m_n}B_n(x^2).
\]
The coefficient $m_{n+1}/m_n=(2n+1)^2/\{2(n+1)(2n+3)\}$ is positive.
By Lemma~\ref{r26meiji:lem} the second derivative of $B_n(x^2)$ is
$2B_n'(x^2)+4x^2B_n''(x^2)<0$, including at zero.

\paragraph{Inverse tangent.}
Here $c_j=(-1)^j/(2j+1)$. Let $\mu$ be the law of $U^2$, with $U$
as above. Then $m_j=1/(2j+1)$. For every real $x>0$,
integrating the \emph{finite} geometric division of $1/(1+t^2)$
from $0$ to $x$ gives
\[
 \arctan x-\sum_{j=0}^{n-1}\frac{(-1)^jx^{2j+1}}{2j+1}
 =(-1)^nx^{2n+1}\int_0^1\frac{u^{2n}}{1+x^2u^2}\,du .
\]
Consequently $T_{2n}[\arctan](x)=A_n(-x^2)/m_n$, and the desired
ratio is $(2n+1)B_n(-x^2)/(2n+3)$. Its derivative is positive because
$-2xB_n'(-x^2)>0$. This proves Theorem~\ref{r26meiji:main} without appealing
to a divergent Taylor series for $x>1$.

\paragraph{Status.}
The statements are identified by the author's numbered source.
The proof is unconditional; it uses standard moment and
Cauchy--Schwarz arguments. We found no earlier explicit resolution of
these two conjectures in a bounded literature search. This is neither
a claim to priority for the underlying moment inequalities nor a
guarantee of exclusive priority for their application.

\begin{thebibliography}{9}
\bibitem{r26meiji:LQ} X.-L. Liu and F. Qi, \emph{Logarithmically absolute
monotonicity and series expansions of normalized remainders for
inverse sine and tangent functions}, author manuscript accepted by
Communications of the Korean Mathematical Society (2026),
\href{https://www.researchgate.net/publication/403516605}{author's text}.
\bibitem{r26meiji:Qi} F. Qi, \emph{Theory of normalized remainders in Taylor
series expansions}, Chapter 12, \href{https://arxiv.org/html/2512.15797v4}{arXiv:2512.15797v4}, especially
Secs.~12.7.1.8 and 12.7.1.9.
\end{thebibliography}

\par\endgroup

% Source: /workspace/scratch/81914cec2c78/work/archive/Tournament_Round_26_Final/Tournament_Round_26/round26/taisho/paper.tex
% Edited source: /workspace/scratch/81914cec2c78/work/archive/Tournament_Round_26_Final/Tournament_Round_26/round26/taisho/paper.tex
\begingroup
\newcommand{\coeff}[1]{[z^{#1}]}
\Needspace{9\baselineskip}
\section{\texorpdfstring{Exponential growth after a one-term perturbation of Euler's partition product}{Exponential growth after a one-term perturbation of Euler's partition product}}\label{section:r26taisho}
\noindent{\small\sffamily\color{muted}Round 26 \enspace / \enspace Taisho}\par\medskip


\begin{sectionabstract}
Keith, Schneider and Sills conjectured exponential coefficient growth for
the reciprocal of Euler's product after subtracting an arbitrary positive
monomial. We obtain the exponential rate and the leading constant for
every parameter in their conjecture. The underlying mechanism is a single
simple pole; consecutive block sizes rule out oscillation at the leading
rate. We also give a weighted-partition proof of an exponential lower bound.
\end{sectionabstract}
\subsection*{The conjecture and its rate}
Write $E(z)=\prod_{m\ge1}(1-z^m)$ and $P(z)=1/E(z)
=\sum_{n\ge0}p(n)z^n$, where $p(n)$ is the partition number,
including $p(0)=1$. Keith, Schneider and Sills \cite[p.~12,
Conjecture~1]{r26taisho:KSS} ask whether the coefficients of $1/(E(z)-az^j)$
grow exponentially for every real $a>0$ and positive integer $j$.
The sign is essential: their Conjecture~2 concerns a different series.
\begin{theorem}
Fix $a>0$ and $j\ge1$. Put $B(z)=az^jP(z)$ and let $r$ be the unique
real number in $(0,1)$ with $B(r)=1$. Then for some $R$ with $r<R<1$,
\[
 c_n:=\coeff{n}\frac{1}{E(z)-az^j}
 = C r^{-n}+O(R^{-n}),\qquad C=\frac{P(r)}{r B'(r)}>0 .
\]
In particular $c_n^{1/n}\to 1/r>1$ and $c_{n+1}/c_n\to 1/r$.
The error bound is for fixed $a,j$.
\end{theorem}
\begin{proof}
The Euler product converges to a nonzero analytic function on
$|z|<1$: on each closed disk of radius $s<1$ one has
$\sum_{m\ge1}s^m/(1-s^m)<\infty$, which gives locally uniform
absolute convergence of the logarithmic series. Expanding its reciprocal
factor by factor gives $p(n)\ge1$ at all degrees. Thus
$P(x)\ge 1/(1-x)$ for $0\le x<1$. On that interval $B$ increases strictly
from zero to infinity, yielding $r$, and its derivative at $r$ is positive.
In the disk $|z|<r$,
$|B(z)|\le B(|z|)<1$, so there is no competing denominator zero.

If $z=r e^{i\theta}$ and $B(z)=1$, taking real parts in
$B(r)-B(z)=0$ gives
\[
 0=\sum_{m\ge j} a p(m-j)r^m(1-\cos(m\theta)).
\]
Each summand is nonnegative. At $m=j$ and $j+1$ the weights are positive,
so both $j\theta$ and $(j+1)\theta$ are integer multiples of $2\pi$.
Subtracting shows $z=r$. This excludes every other point on the
circle, and $1-B$ has a simple zero at $r$.

There are only finitely many zeros of $1-B$ in the closed disk
of radius $(1+r)/2$: an accumulation point there would force the
nonconstant analytic function $1-B$ to vanish identically.
Consequently one can choose $R\in(r,(1+r)/2)$ with no other zero in
$|z|\le R$. Using $1/(E-az^j)=P/(1-B)$, the function
\[
 H(z)=\frac{P(z)}{1-B(z)}-\frac{C}{1-z/r}
\]
is analytic on a neighborhood of $|z|\le R$, since the residues of
the two terms at $r$ are both $-P(r)/B'(r)$. The coefficient bound
$|\coeff{n}H|\le R^{-n}\max_{|z|=R}|H(z)|$ proves the formula.
The two limits follow because $C>0$ and $r/R<1$.
\end{proof}
\subsection*{A direct combinatorial lower bound}
The formal geometric series $P\sum_{t\ge0}(az^j P)^t$ encodes a
list of $t$ partitions, each preceded by a marker of size $j$ and
weight $a$, followed by an unmarked partition. Distinct partition choices
contribute separately and all weights are positive, even for nonintegral
$a$. Choose $m$ so large that $b=a p(m)>1$; such an $m$ exists because
partitions using only ones and twos already give $\lfloor m/2\rfloor+1$
choices. Write $n=t(j+m)+s$, $0\le s<j+m$. Restrict the marked
partitions to have size $m$ and the unmarked partition to be $s$ ones.
Their total weight is $b^t$, whence
\[
  c_n\ge b^{\lfloor n/(j+m)\rfloor}
       \ge b^{-1}\bigl(b^{1/(j+m)}\bigr)^n .
\]
This gives exponential growth from below even when $a$ is arbitrarily
small. The pole argument determines the exact rate, which the restricted
model alone does not give.
\paragraph{Method and scope.}
Simple-pole transfer is standard analytic combinatorics \cite{r26taisho:FS}.
The one-pole viewpoint in VMRC manuscript~013 \cite{r26taisho:VMRC} is a
methodological parallel: our uniqueness holds only up to $R$, and we
make no arithmetic claim about $1/r$. Numerical estimates are not
used in the proof. No result for $1/(E+az^j)$ is asserted here.
\begin{thebibliography}{9}
\bibitem{r26taisho:KSS} W.~J. Keith, R.~Schneider and A.~V. Sills,
Composition-theoretic series and false theta functions,
\emph{Integers} \textbf{24A} (2024), A11, Conjecture~1.
\url{https://doi.org/10.5281/zenodo.11352821}.
\bibitem{r26taisho:FS} P.~Flajolet and R.~Sedgewick, \emph{Analytic Combinatorics},
Cambridge University Press, 2009, Ch.~IV.
\bibitem{r26taisho:VMRC} KEIO, \emph{Profinite Recognition of Mahler Equations
and D-Finite Witt Escape}, VMRC manuscript 013.
\url{https://github.com/QEEDDD/Verifiable-Mathematical-Reasoning-Corpus}.
\end{thebibliography}

\par\endgroup

% Source: /workspace/scratch/81914cec2c78/work/archive/Tournament_Round_26_Final/Tournament_Round_26/round26/showa/checkpoint.tex
% Edited source: /workspace/scratch/81914cec2c78/work/audit_26_33/cleaned_sources/r26_showa.tex
\begingroup
\newcommand{\Z}{\mathbb Z}
\Needspace{9\baselineskip}
\section{\texorpdfstring{The two arms of a truncated pentagonal series}{The two arms of a truncated pentagonal series}}\label{section:r26showa}
\noindent{\small\sffamily\color{muted}Round 26 \enspace / \enspace Showa}\par\medskip


\noindent\textbf{Attribution.} This section proves the full coefficient-positivity and exact-support statement for the symmetric truncation formulated in Merca's Theorem~16 [1]. The result was already proved by J.-C.~Liu [2]; see also [3, Remark~2]. The following is an independent pairing proof of that known statement. Its displayed summand sign is $(-1)^j$, and no claim is made about a differently printed formula.

\paragraph{Setting.}
Write $p(j)=j(3j-1)/2$, $E(q)=\prod_{n\ge1}(1-q^n)$, and $B=\{n(3+\nu_2(n)):n\ge1\}$.
Merca's partition interpretation [1, Theorem~11] gives
\[
 H(q)=\frac{\prod_{m\ge1}(1-q^{N_{2m}})}{E(q)}
     =\prod_{t\notin B,\ t\ge1}(1-q^t)^{-1},\qquad N_{2m}=m(3+\nu_2(m)).
\]
The coefficient form [1, Theorem~16] asks for the sign and exact support of
\[
 C_k(q)=(-1)^{k-1}H(q)\left(E(q)-\sum_{j=-k}^k(-1)^jq^{p(j)}\right),\qquad k\ge1.
\]
\begin{lemma} If $H(q)=V(q)/((1-q)(1-q^2))$ where $V(0)=1$ and $V$ has nonnegative coefficients, then for every integer $d\ge2$, all coefficients of $(1-q^d)H(q)$ are strictly positive.
\end{lemma}
\begin{proof}
The product $(1+q+\cdots+q^{d-1})(1-q^2)^{-1}$ covers every nonnegative degree, using an even exponent with residue $0$ or $1$. Multiplying by $V$ preserves strict positivity.
\end{proof}
\paragraph{Application.}
Both 1 and 2 are allowed, since the least member of $B$ is 3. Euler's pentagonal theorem converts the parenthesis into two alternating tails, indexed on each arm by $j\ge k+1$. Pair consecutive indices separately on the positive and negative arms. This gives the locally finite sum
\[
 C_k(q)=\sum_{\substack{j\ge k+1\\j\equiv k+1\bmod 2}}\;
 \sum_{s\in\{1,-1\}}q^{p(sj)}(1-q^{p(s(j+1))-p(sj)})H(q).
\]
The differences are $3j+1$ and $3j+2$, respectively. Thus each summand is positive at every degree starting at its own initial exponent, and zero below it. The smallest such exponent is $p(k+1)=(k+1)(3k+2)/2$. The coefficients of $C_k$ are therefore strictly positive exactly from that degree onward. Unlike an asymmetric truncation, a symmetric truncation starts both arms with the same sign, which is the only reason this argument applies.
\paragraph{Sources.}
[1] M.~Merca, \href{https://doi.org/10.1007/s00010-024-01117-6}{Overpartitions in terms of 2-adic valuation}, Aequationes Math. 99 (2025), 1153--1173.\quad
[2] J.-C.~Liu, \href{https://arxiv.org/abs/2409.19907}{A general positivity result on coefficients of certain $q$-series}, J. Algebraic Combin. 62 (2025), 46.\quad
[3] M.~P.~Saikia and A.~Sarma, \href{https://arxiv.org/abs/2609.25739}{On the positivity of truncated pentagonal number series and some conjectures of Merca}, arXiv:2609.25739.

\par\endgroup

% Source: /workspace/scratch/81914cec2c78/work/archive/Tournament_Round_26_Final/Tournament_Round_26/round26/heisei/paper.tex
% Edited source: /workspace/scratch/81914cec2c78/work/archive/Tournament_Round_26_Final/Tournament_Round_26/round26/heisei/paper.tex
\begingroup

\Needspace{9\baselineskip}
\section{\texorpdfstring{Why derivative-sign regions nest for polynomial log-functions}{Why derivative-sign regions nest for polynomial log-functions}}\label{section:r26heisei}
\noindent{\small\sffamily\color{muted}Round 26 \enspace / \enspace Heisei}\par\medskip


\textbf{Status.} This is a verification of a known result, not a new resolution.
Ma and Weigert posed the nesting property as Conjecture 4.6 in 2025.
Explicit pre-existing solutions were located during the round.

\textbf{Statement.} For a real polynomial $p$ of degree at most $n$, set
$f(x)=x^{-1}p(\log x)$ on $x>0$ and $L=-d/dx$. Let $D_k$ be the
set of such functions with $L^kf\ge0$ on $(0,\infty)$. Then
$D_{k+1}\subseteq D_k$ for every $k\ge0$, and
$\bigcap_{k\ge1}D_k$ is exactly the cone of completely monotone
functions in this family. The domain $x>0$ is the one where the defining real logarithm exists.

\textbf{Proof.} Set $p_0=p$. The chain rule gives
\[
 L^kf(x)=x^{-k-1}p_k(\log x),\qquad
 p_{k+1}(y)=(k+1)p_k(y)-p_k'(y).
\]
In particular $p_k$ is a polynomial of degree at most $n$ at every
stage. On writing $x=e^y$, the expression $e^{-(k+1)y}p_k(y)$ tends
to zero as $y\to+\infty$. For any $R>x>0$, the fundamental theorem of calculus gives
\[
 L^kf(x)-L^kf(R)=\int_x^R L^{k+1}f(t)\,dt.
\]
If $f\in D_{k+1}$, the integral is nonnegative. Taking $R\to\infty$
gives $L^kf(x)\ge0$, proving the inclusion. Membership in
all positive-order sets therefore also gives the order-zero sign.
The converse is the definition of complete monotonicity.\hfill$\square$

\textbf{Source and attribution.}
R.~Ma and J.~Weigert, \emph{Complete monotonicity of log-functions},
Anal.~Math.~Phys. 15 (2025), article 141,
\href{https://doi.org/10.1007/s13324-025-01136-9}{doi:10.1007/s13324-025-01136-9},
Conjecture 4.6. Both the
\href{https://www.trureturing.org/research/ma-weigert-2025-log-function-monotone-chain/}{trureturing proof record}
and the \href{https://github.com/pudim-project/pudim-ai-demo-zetalaw/blob/main/APPLICATIONS.md}{pudim-project applications, APP-0027},
already give the conclusion; the latter gives the integral argument.
No priority is claimed. The separation of finite checks from universal
proof follows the audit practice of the
\href{https://github.com/QEEDDD/Verifiable-Mathematical-Reasoning-Corpus}{Verifiable Mathematical Reasoning Corpus};
no theorem from that corpus is needed.

\par\endgroup

% Source: /workspace/scratch/81914cec2c78/work/archive/Tournament_Round_26_Final/Tournament_Round_26/round26/reiwa/huang_hung.tex
% Edited source: /workspace/scratch/81914cec2c78/work/archive/Tournament_Round_26_Final/Tournament_Round_26/round26/reiwa/huang_hung.tex
\begingroup

\Needspace{9\baselineskip}
\section{\texorpdfstring{Sharp weight intervals for the Huang--Hung Wilker inequalities}{Sharp weight intervals for the Huang--Hung Wilker inequalities}}\label{section:r26reiwa}
\noindent{\small\sffamily\color{muted}Round 26 \enspace / \enspace Reiwa}\par\medskip


\begin{sectionabstract}
Huang and Hung proposed four trigonometric and hyperbolic inequalities at the end of their 2023 preprint. We give a common proof of the trigonometric assertions and a direct proof of the hyperbolic assertion. The key step for the sharp upper threshold is that a rational-weight combination, seven parts of a squared cosine difference to five parts of a sinc function, is strictly increasing up to $\pi/2$. A one-variable polynomial with rational coefficients certifies this claim.
\end{sectionabstract}

Write $I=(0,\pi/2)$, and on $I$ set
\[
 A(x)=\left(\frac{1-\cos x}{x}\right)^2,\quad
 B(x)=\frac{\sin x}{x},\quad F_p(x)=pA(x)+(1-p)B(x),\qquad
 p_* =\frac{\pi(4-\pi)}{4(\pi-2)}.
\]
\begin{theorem}
The four assertions proposed in \cite[Sec.~3, (3.1)--(3.4)]{r26reiwa:HH} hold:
\begin{align}
 A(x)^{-1}+B(x)^{-1}&>A(x)+B(x)>1 &&(x\in I),\tag{3.1}\\
 p\left(\frac{\cosh x-1}{x}\right)^2+(1-p)\frac{\sinh x}{x}
 &>\tfrac12 &&(x\in\mathbb R\setminus\{0\},\ 0\le p\le\tfrac12),\tag{3.2}\\
 F_p(x)&>\tfrac12 &&(x\in I,\ 0\le p\le\tfrac12),\tag{3.3}\\
 F_p(x)&<\tfrac12 &&(x\in I,\ p_*\le p\le1).\tag{3.4}
\end{align}
Within $0\le p\le 1$, the intervals given in (3.2)--(3.4) are exact.
\end{theorem}

\paragraph{Lower bounds.}
Since $0<1-\cos x<\sin x<x$ on $I$ (use $\tan(x/2)<1$ for the middle inequality), we have $0<A<B<1$. The left inequality in (3.1) follows term by term. Put $y=x^2<5/2$, using $\pi<22/7$. The alternating Taylor bounds give
\[
 1-\cos x\ge y(1/2-y/24)>0,\qquad
 B(x)\ge 1-y/6+y^2/120-y^3/5040.
\]
Squaring the first bound and dividing by $y$ gives
\[
 (A+B-1)/y\ \ge\ 1/12-y/30+(1/576-1/5040)y^2>0.
\]
This also proves the lower bound of Huang and Hung's Theorem 2.4 independently. Since $F_p$ decreases in $p$, (3.3) follows by taking $p=1/2$. For (3.2), $\sinh x/x>1$ for every nonzero real $x$, and $1-p\ge1/2$. The limit at $x=0$ in both (3.2) and (3.3) is $1-p$, excluding $p>1/2$ from either lower inequality.

\paragraph{A uniform derivative estimate.}
The main step for the new upper bound (3.4) is the following certificate.
\begin{lemma}
For $x\in I$, both $A'(x)>0$ and $7A'(x)+5B'(x)>0$.
\end{lemma}
\begin{proof}
Writing $u=1-\cos x$ yields $x^3 A'=2u(x\sin x-u)>0$ because $x\sin x-u=\int_0^xt\cos t\,dt$. By double-angle identities,
\begin{align*}
 D(x)&=x^3(7A'+5B')=14u(x\sin x-u)+5x(x\cos x-\sin x)\\
 &=9x\sin x-7x\sin(2x)-21+28\cos x-7\cos(2x)+5x^2\cos x .
\end{align*}
For $z=x,2x$, define $C_k(z)=\sum_{j=0}^k(-1)^j z^{2j}/(2j)!$ and
$S_k(z)=\sum_{j=0}^k(-1)^j z^{2j+1}/(2j+1)!$.
The remainder after $k=5,6$ has the sign of its first term for $0<z<\pi$, by the alternating tail estimate. Thus in the last line for $D$, replace positively signed trigonometric functions by the corresponding $C_5,S_5$ and negatively signed ones by $C_6,S_6$. Expansion gives $D(x)\ge x^4 G(y)$, where
\[
 G(y)=\frac{11}{6}-y+\frac{421 y^2}{3360}-\frac{22 y^3}{2835}
      +\frac{557 y^4}{1871100}-\frac{8y^5}{868725}.
\]
With $b=5/2$, omitting negative terms after differentiation gives
\[
 G'(y)\le-1+\frac{842b}{3360}+\frac{2228b^3}{1871100}
 =-\frac{212501}{598752},\qquad G(b)=\frac{6827}{1153152}>0.
\]
This proves the lemma.
\end{proof}

\paragraph{Upper bound and optimality.}
The lemma makes $F_{7/12}$ strictly increasing, as is $F_1=A$. By affine interpolation in $p$, every $F_p$ for $7/12\le p\le1$ is strictly increasing on $I$. We have $7/12<p_*<1$: the first inequality amounts to $14+5\pi-3\pi^2>0$, since the polynomial decreases for $t\ge3$ and is $4/49$ at $22/7$; the second follows from $\pi^2>8$. Substituting the right endpoint gives
\[
 F_{p_*}(\pi/2)=p_*\frac{4}{\pi^2}+(1-p_*)\frac{2}{\pi}=\frac12.
\]
Consequently $F_p(x)\le F_{p_*}(x)<1/2$ for $p\ge p_*$. For $p<p_*$, the right endpoint value exceeds $1/2$, so (3.4) fails near that endpoint. This completes the proof.

\begin{thebibliography}{1}
\bibitem{r26reiwa:HH} Y.-C. Huang and L.-C. Hung, \emph{On Wilker-type inequalities},
viXra:2312.0112v1 (2023), Sec.~3, p.~5.
\href{https://vixra.org/pdf/2312.0112v1.pdf}{Public preprint}.
\end{thebibliography}
\noindent\small Scope: The target is the four formulas in this public preprint, not a historical open problem of Wilker himself. This is an unrefereed proof; the priority search is not exhaustive.

\par\endgroup

% Source: /workspace/scratch/81914cec2c78/work/archive/Tournament_Round_26_Final/Tournament_Round_26/round26/keio/paper.tex
% Edited source: /workspace/scratch/81914cec2c78/work/archive/Tournament_Round_26_Final/Tournament_Round_26/round26/keio/paper.tex
\begingroup

\Needspace{9\baselineskip}
\section{\texorpdfstring{The sharp exponent in a half-angle Wilker inequality}{The sharp exponent in a half-angle Wilker inequality}}\label{section:r26keio}
\noindent{\small\sffamily\color{muted}Round 26 \enspace / \enspace Keio}\par\medskip


\begin{sectionabstract}
Yang, Chu, Song and Li conjectured in 2014 that the positive exponent threshold for a mixed full-angle/half-angle Wilker inequality is $3/20$. A rational majorant for $t/\sin t$ reduces the endpoint to positivity of one integer polynomial on $(0,1)$. An exact coefficient certificate is included. This argument is not a claim of bibliographic priority.
\end{sectionabstract}
\subsection{Statement and reduction}
For $0<t<\pi/2$, write
\[
 X=\frac{t}{\sin t},\qquad Y=\left(\frac{t/2}{\tan(t/2)}\right)^2,
 \qquad F_p(t)=X^p+Y^p.
\]
\textbf{Theorem.} The inequality $F_p(t)>2$ holds for every $0<t<\pi/2$ exactly when $p<0$ or $p\ge 3/20$.

This is the conjecture following (64) in \cite{r26keio:Y}, where Corollary 15 already gives $p<0$ and $p\ge 1/5$. For necessity, the expansions at zero give
\[
 F_p(t)-2=\frac{p(20p-3)}{720}\,t^4+O(t^6).
\]
This excludes $0<p<3/20$; $p=0$ gives equality. For $p<0$, observe that $XY=X^3\cos^4(t/2)$ and
\[
 (\log(XY))'=\frac3t-\frac{2+\cos t}{\sin t}<0.
\]
Indeed $t(2+\cos t)-3\sin t$ vanishes at zero and has derivative
$2(1-\cos t)-t\sin t=\sin t\,[2\tan(t/2)-t]>0$. Thus $XY<1$ and $F_p\ge 2(XY)^{p/2}>2$.

For each fixed $t$, the function $p\mapsto X^p+Y^p$ is strictly convex, since $X>1$. Its value at zero is $2$. Thus if it exceeds $2$ at $p_*=3/20$, it exceeds $2$ for every $p>p_*$. It remains to prove the critical case.

\subsection{One-sided reduction to rational functions}
Set $z=1-\cos t\in(0,1)$ and $h=t/\sin t$. Since
\[
 (\log X)'=1/t-\cot t,\qquad (\log Y)'=2(1/t-1/\sin t),
\]
and $\sin t-t\cos t>0$ (its derivative is $t\sin t$), the inequality $F'_{p_*}(t)>0$ is equivalent to
\begin{equation}
 \left(\frac{4}{(2-z)^2h}\right)^{3/20}>
 \frac{2(h-1)}{1-(1-z)h}.\label{r26keio:eq:core}
\end{equation}
Here $X/Y=4/((2-z)^2h)$. The left side decreases with $h$; the derivative of the right side with respect to $h$ is $2z/[1-(1-z)h]^2>0$ where the denominator before squaring is positive.

The differential identity $z(2-z)h'=1-(1-z)h$, now differentiating in $z$, gives
\[
 h(z)=\sum_{n\ge0}c_n z^n,\qquad c_0=1,\quad
 c_{n+1}=\frac{n+1}{2n+3}c_n.
\]
Indeed the indicated series converges on $|z|<2$ and solves the equation by coefficient comparison. Its difference from $h$ is a multiple of $1/\sqrt{z(2-z)}$ by the homogeneous equation. Both are bounded at zero, so the multiple vanishes. As $c_{n+1}/c_n<1/2$, a strict upper bound on $(0,1)$ is
\[
 H=1+\tfrac13 z+\tfrac2{15}z^2+\tfrac2{35}z^3+
       \frac{8z^4}{315(1-z/2)}=\frac{V}{Dq},\qquad D=315,\ q=2-z.
\]
For this single majorant define
\begin{align*}
 V&=630-105z-21z^2-6z^3-2z^4,\\
 A&=(Dq-(1-z)V)/z=420-84z-15z^2-4z^3-2z^4,\\
 B&=(V-Dq)/z=210-21z-6z^2-2z^3.
\end{align*}
Each polynomial is strictly positive on $[0,1]$ by evaluation at $1$ and monotonicity. In particular $1-(1-z)H=zA/(Dq)>0$, so replacing $h$ by $H$ in \eqref{r26keio:eq:core} is sufficient. The resulting inequality is $(4D/(qV))^{3/20}>2B/A$. Raising to the twentieth power and canceling $2^6$ gives precisely
\begin{equation}
 U(z):=D^3A^{20}-2^{14}(qV)^3B^{20}>0.\label{r26keio:eq:poly}
\end{equation}
The next section certifies \eqref{r26keio:eq:poly}. Consequently $F'_{p_*}(t)>0$ on the whole interval. Its right limit at zero is $2$, which proves the required strict inequality and completes the theorem.

\subsection{Finite integer certificate}
Write $U(z)=\sum_{k=0}^{80}u_kz^k$. Under $z=s/(1+s)$ the coefficient of $s^i$ in $(1+s)^{80}U(s/(1+s))$ is
\[
 b_i=\sum_{k=0}^i u_k \binom{80-k}{i-k}.
\]
There are two initial zeros and 79 strictly positive coefficients. The minimum positive coefficient, at $i=80$, is
\[
 44308873133738660170865189097115371245117517925833764051.
\]
Thus the transformed polynomial is positive for every $s>0$, equivalent to \eqref{r26keio:eq:poly}. For reproducibility the full finite sign calculation is given below. Arrays are coefficient lists in ascending degree. All operations are on integers; neither rounding nor sampling occurs.
\begin{Verbatim}[fontsize=\small,breaklines=true,breakanywhere=true]
from math import comb
def mul(a,b):
    c=[0]*(len(a)+len(b)-1)
    for i,x in enumerate(a):
        for j,y in enumerate(b): c[i+j]+=x*y
    return c
def pw(a,n):
    c=[1]
    for _ in range(n): c=mul(c,a)
    return c
q=[2,-1]; V=[630,-105,-21,-6,-2]
A=[420,-84,-15,-4,-2]; B=[210,-21,-6,-2]
a=pw(A,20); b=mul(pw(mul(q,V),3),pw(B,20))
assert len(a)==81 and len(b)<=81
b += [0]*(81-len(b))
u=[315**3*x-2**14*y for x,y in zip(a,b)]
beta=[sum(u[k]*comb(80-k,i-k) for k in range(i+1))
      for i in range(81)]
assert beta[:2]==[0,0] and min(beta[2:])>0
\end{Verbatim}


\begin{thebibliography}{1}
\bibitem{r26keio:Y} Z.-H. Yang, Y.-M. Chu, Y.-Q. Song and Y.-M. Li,
A sharp double inequality for trigonometric functions and its applications,
\emph{Abstract and Applied Analysis} (2014), Article 592085,
\href{https://doi.org/10.1155/2014/592085}{doi:10.1155/2014/592085}.
See equations (62)--(64) and the conjecture on page 7.
\end{thebibliography}

\par\endgroup

% Source: /workspace/scratch/81914cec2c78/work/archive/Tournament_Round_27_Final/Tournament_Round_27/round27/meiji/paper.tex
% Edited source: /workspace/scratch/81914cec2c78/work/archive/Tournament_Round_27_Final/Tournament_Round_27/round27/meiji/paper.tex
\begingroup

\Needspace{9\baselineskip}
\section{\texorpdfstring{A curvature bound for power-type Toader means}{A curvature bound for power-type Toader means}}\label{section:r27meiji}
\noindent{\small\sffamily\color{muted}Round 27 \enspace / \enspace Meiji}\par\medskip


\begin{sectionabstract}
A monotonicity question posed in 2022 is a direct special case of Tobey's
1967 theorem on two-parameter homogeneous means. We make the
identification explicit and give an independent curvature proof
which yields a positive quantitative gap on compact order intervals.
This is not claimed as a new resolution of an open problem.
\end{sectionabstract}
\subsection*{The earlier theorem and the later conjecture}
For $a,b>0$, let $d\mu=(2/\pi)d\theta$ on $[0,\pi/2]$ and $w=\cos^2\theta$.
Write
\[
T_r=\left(\int (wa^r+(1-w)b^r)^{1/2}\,d\mu\right)^{2/r}
\quad (r\ne0),\qquad T_0=\sqrt{ab}.
\]
Zhao, Wang, Dai and Chu~[1, Theorem 2.5] proved strict increase in
nonzero integer orders in $[-43,43]$, for $a\ne b$. Their
Conjecture 2.6 asks for all nonzero integer orders.

With $M_s(a,b;u)=(ua^s+(1-u)b^s)^{1/s}$ and the usual continuous
conventions at zero, Tobey's mean~[2] is
\[
L(s,t)=\left(\int_0^1 M_s(a,b;u)^t P(u)\,du\right)^{1/t}.
\]
For $P(u)=1/(\pi\sqrt{u(1-u)})$, the substitution $u=\cos^2\theta$
gives $T_r=L(r,r/2)$. By Tobey's Theorem 2, $L$ increases in each
order (strictly here when $a\ne b$); the diagonal restriction is
therefore strictly increasing on the entire real line. This already
answers Conjecture 2.6. The following proof also measures the increase.

\subsection*{A quantitative derivative certificate}
Assume $a>b$ by symmetry and put $d=\log(a/b)>0$, $Z=1-w+w e^{rd}$,
and $F(r)=\log\int Z^{1/2}\,d\mu$. Then $F(0)=0$ and
$\log(T_r/b)=2F(r)/r$. Let $\mu_r$ denote the probability measure
of density $Z^{1/2}/\int Z^{1/2}d\mu$ with respect to $\mu$ and put
$q=w e^{rd}/Z$. Differentiating under the integral gives
\[
F'(r)=\tfrac d2\mathbb E_r q,\qquad
F''(r)=\tfrac{d^2}{2}\mathbb E_r[q(1-q)]
+\tfrac{d^2}{4}\operatorname{Var}_r(q).
\]
Indeed $q'=dq(1-q)$, while differentiating the density adds the
term $(d/2)(q-\mathbb E_r q)$. Compactness of the angle domain and
the positive lower bound of $Z$ on compact order intervals justify
these differentiations. In particular $F''>0$ everywhere and
$2F'(0)=d/2$, verifying the claimed continuous convention.
Since $F(r)/r=\int_0^1F'(ur)\,du$ on either side of zero, we have
\[
(\log T_r)'=2\int_0^1 uF''(ur)\,du>0
\]
also at zero. As $\mathbb E_0q(1-q)=1/8$ and
$\operatorname{Var}_0(q)=1/8$, its value there is $3d^2/32$.

The density of $\mu_r$ is at least $e^{-|r|d/2}$, by bounding $Z$
between $\min(1,e^{rd})$ and $\max(1,e^{rd})$. Likewise
$q(1-q)=w(1-w)e^{rd}/Z^2\ge w(1-w)e^{-|r|d}$. Therefore
\[
F''(r)\ge \tfrac{d^2}{16}e^{-3|r|d/2}.
\]
For all $-R\le s<t\le R$ this gives the explicit gap
\[
\log(T_t/T_s)\ge (t-s)\tfrac{d^2}{16}e^{-3Rd/2}.
\]
The elementary coefficient at zero is exact; the uniform lower
bound is not asserted sharp. For $a=b$ the means are constant.

\begin{thebibliography}{2}
\bibitem{r27meiji:z} T.-H. Zhao, M.-K. Wang, Y.-Q. Dai, and Y.-M. Chu,
On the generalized power-type Toader mean, \emph{J. Math. Inequal.}
\textbf{16} (2022), 247--264. DOI: 10.7153/jmi-2022-16-18.
\bibitem{r27meiji:t} M. D. Tobey, A two-parameter homogeneous mean value,
\emph{Proc. Amer. Math. Soc.} \textbf{18} (1967), 9--14.
DOI: 10.2307/2035216.
\end{thebibliography}

\par\endgroup

% Source: /workspace/scratch/81914cec2c78/work/archive/Tournament_Round_27_Final/Tournament_Round_27/round27/taisho/paper.tex
% Edited source: /workspace/scratch/81914cec2c78/work/archive/Tournament_Round_27_Final/Tournament_Round_27/round27/taisho/paper.tex
\begingroup
\newcommand{\E}{\mathbb E}
\newcommand{\Var}{\operatorname{Var}}
\newcommand{\Cov}{\operatorname{Cov}}
\Needspace{9\baselineskip}
\section{\texorpdfstring{A covariance proof of the Yang--Chu--Zhang digamma quotient conjecture}{A covariance proof of the Yang--Chu--Zhang digamma quotient conjecture}}\label{section:r27taisho}
\noindent{\small\sffamily\color{muted}Round 27 \enspace / \enspace Taisho}\par\medskip


\begin{sectionabstract}
The two rational constants in a conjectured monotonicity inequality for the
half-shifted digamma function are the first two moments of a single
probability law. We prove that the derivative of the quotient, in the squared
argument, is a normalized covariance determinant. Its sign follows strictly
on the whole positive axis. This gives an affirmative answer to
Yang--Chu--Zhang's Conjecture~1(i) and recovers the exact endpoint values.
The determinant argument applies to every Stieltjes probability transform
with finite second moment and at least three support points.
\end{sectionabstract}

\subsection{The exact conjecture}
Write $R(x)=\psi(x+1/2)-\log x$, where $\psi=\Gamma'/\Gamma$.
Yang's 2014 preprint \cite{r27taisho:yang} and the subsequent journal article of
Yang, Chu and Zhang \cite{r27taisho:ycz} ask whether the following function is
increasing. The source already bounds its values by its two endpoint limits;
these bounds do not compare two interior arguments.

\begin{theorem}\label{r27taisho:main}
For all real $x>0$ the denominator below is strictly positive, and
\[
 H(x)=\frac{-24(x^4-31/336)R(x)+x^2-7/40}
                 {24(x^2+7/40)R(x)-1}
\]
has strictly positive derivative. Its image is the open interval
$(155/294,11165/8284)$.
\end{theorem}
The independent complete-monotonicity assertion for a parameterized family,
printed as Conjecture~1(ii), is not a hypothesis or a conclusion here.

\subsection{A derivative identity for a positive measure}
\begin{lemma}\label{r27taisho:general}
Let $S$ be a nonnegative random variable with finite second moment and
support containing at least three distinct points. Put
$m=\E S$, $b=\E S^2$, $v=b-m^2$. For $y>0$, set $U=(y+S)^{-1}$,
$f=\E U$, $D=(y+m)f-1$, and
\[
 C(y)=\frac{y-m-(y^2-b)f}{D}.
\]
Then $D>0$ and
\[
 C'(y)=\frac{v\Var(U)-\Cov(S,U)^2}{\Cov(S,U)^2}>0.
\]
\end{lemma}
\begin{proof}
Since $SU=1-yU$, the denominator is $D=-\Cov(S,U)$.
If $S'$ is an independent copy of $S$, symmetrizing the covariance gives
\[
 D=\tfrac12\E\frac{(S-S')^2}{(y+S)(y+S')}>0 .
\]
A division of the numerator gives $C=m-y+vf/D$. Differentiation under the
expectation yields $f'=-\E U^2$; this is justified locally for $y>0$ by
the uniform bound on $(y+S)^{-2}$. Thus
\[
 f'D-fD'=-f'-f^2=\Var(U).
\]
The formula follows. Equality in Cauchy--Schwarz for the two centered
variables would imply $(y+S)^{-1}=a+cS$ almost surely. By continuity the
same identity would hold on the support. But
$(y+s)(a+cs)-1$ is a nonzero polynomial of degree at most two, and cannot
vanish at three points. This proves the strict inequality.
\end{proof}

This argument also gives an explicit positive certificate.
If $\mu$ denotes the law of $S$, then
\begin{equation}\label{r27taisho:triple}
 v\Var(U)-\Cov(S,U)^2=\frac1{6}
 \int\frac{\prod_{i<j}(s_i-s_j)^2}
 {\prod_{i=1}^3(y+s_i)^2}\,d\mu(s_1)d\mu(s_2)d\mu(s_3).
\end{equation}
Indeed, the left side is the Gram determinant of $1,s,(y+s)^{-1}$
after centering the final two functions. Expanding the squared evaluation
determinant and integrating gives $3!$ times that Gram determinant.
Multiplying each evaluation column by $y+s_i$ leaves a determinant of
polynomials $y+s,s(y+s),1$, a Vandermonde up to sign. This explains why a
two-point law gives equality and a third support point forces strictness.
The use of positive-measure Gram certificates is informed by the
moment constructions of corpus item~[001] \cite{r27taisho:corpus}; the measure and
the identification of this particular derivative are supplied here.

\subsection{Identification of the law and proof of the theorem}
For $z>0$ the differentiated Binet formula \cite{r27taisho:dlmf} reads
\[
 \psi(z)=\log z-\frac{1}{2z}
 -2\int_0^\infty\frac{t}{(z^2+t^2)(e^{2\pi t}-1)}\,dt.
\]
In $\psi(x+1/2)=2\psi(2x)-\psi(x)-2\log 2$, substitute this formula,
and replace $t$ by $2t$ in the $2x$ integral. The pole terms cancel.
Since $1/(a-1)-2/(a^2-1)=1/(a+1)$, the change $s=t^2$ gives
\[
 R(x)=2\int_0^\infty \frac{t}{(x^2+t^2)(e^{2\pi t}+1)}\,dt
     =\int_0^\infty \frac{1}{x^2+s}\frac{ds}{e^{2\pi\sqrt{s}}+1}.
\]
In particular, set $d\mu(s)=24\,ds/(e^{2\pi\sqrt{s}}+1)$ on $s>0$.
Its moments are
\[
 M_k=\frac{48(2k+1)!}{(2\pi)^{2k+2}}
       \big(1-2^{-2k-1}\big)\zeta(2k+2),\qquad k\ge0.
\]
To see this, substitute $r=2\pi\sqrt{s}$ and expand
$1/(e^r+1)$ into alternating exponentials; the integrated series is
absolutely convergent since $2k+2\ge2$. The standard even zeta values give
\[
 (M_0,M_1,M_2,M_3)=\left(1,\frac7{40},\frac{31}{336},\frac{127}{1280}\right).
\]
Consequently $\mu$ is a probability law with support $[0,\infty)$, and
$f(y)=24R(\sqrt{y})$ is precisely its transform from Lemma~\ref{r27taisho:general}.
With $y=x^2$, substituting $m=M_1$, $b=M_2$ yields $H(x)=C(x^2)$
term for term. Thus $H'(x)=2xC'(x^2)>0$.

As $y\downarrow0$, $f(y)\to\infty$ because the density is bounded below
near zero. Dividing numerator and denominator by $f$ yields
$C(0+)=b/m=155/294$. At infinity the finite geometric expansion of
$(y+s)^{-1}$, with integrated error bounded by $M_4/y^5$, gives
$f=y^{-1}-m y^{-2}+b y^{-3}-M_3 y^{-4}+O(y^{-5})$.
The first nonzero terms in the numerator and denominator of $C$ are
$(M_3-mb)y^{-2}$ and $(b-m^2)y^{-2}$, respectively. Hence
\[
 C(\infty)=\frac{M_3-mb}{b-m^2}=\frac{11165}{8284}.
\]
Continuity and strict increase complete the proof of Theorem~\ref{r27taisho:main}.

\paragraph{Independent work.}
The Heisei seat independently studied and proved the same monotonicity
assertion in this round. These are independent approaches to one conjecture,
not two separate problem closures.

\begin{thebibliography}{9}
\bibitem{r27taisho:ycz} Z.-H. Yang, Y.-M. Chu and X.-H. Zhang,
Necessary and sufficient conditions for functions involving the psi function
to be completely monotonic, \emph{J. Inequal. Appl.} (2015), article 157,
Conjecture 1(i). \url{https://doi.org/10.1186/s13660-015-0674-8}.
\bibitem{r27taisho:yang} Z.-H. Yang, Several completely monotone functions related to
DeTemple's sequence (2014), p.~14.
\url{https://arxiv.org/abs/1409.6411}.
\bibitem{r27taisho:dlmf} NIST Digital Library of Mathematical Functions, equations
\href{https://dlmf.nist.gov/5.9.E15}{5.9.15} and
\href{https://dlmf.nist.gov/5.5.E8}{5.5.8} (duplication, after logarithmic
differentiation of the gamma identity with multiplier two).
\bibitem{r27taisho:corpus} QEEDDD, \emph{Verifiable Mathematical Reasoning Corpus},
item [001], \emph{Exponential-Integral Periods across Arithmetic, Pad\'e,
Hankel, Reciprocal Geometry, and Frobenius Framing}, positive-measure
Gram determinant construction.
\href{https://github.com/QEEDDD/Verifiable-Mathematical-Reasoning-Corpus/tree/dd9c67cb72c5991cce5816de540db8c2132ca343}{Repository snapshot}.
\end{thebibliography}

\par\endgroup

% Source: /workspace/scratch/81914cec2c78/work/archive/Tournament_Round_27_Final/Tournament_Round_27/round27/showa/showa.tex
% Edited source: /workspace/scratch/81914cec2c78/work/archive/Tournament_Round_27_Final/Tournament_Round_27/round27/showa/showa.tex
\begingroup
\newcommand{\CB}{\operatorname{CB}}
\Needspace{9\baselineskip}
\section{\texorpdfstring{The mixed-parity facet bound for bicyclic graphs}{The mixed-parity facet bound for bicyclic graphs}}\label{section:r27showa}
\noindent{\small\sffamily\color{muted}Round 27 \enspace / \enspace Showa}\par\medskip


\begin{sectionabstract}
For three internally disjoint paths joining a pair of vertices, we prove the
mixed-parity facet-count bound proposed by Braun and Bruegge. Two complementary
estimates for a centered binomial sum control the cost of one contracted edge;
the cost is absorbed by a uniform additive increment of a normalized product of
odd-cycle facet counts. The argument requires only 36 small integer comparisons,
listed explicitly. Together with the known reductions this determines the
maximum for simple connected graphs with two independent cycles.
\end{sectionabstract}
\subsection{The extremal problem}
The symmetric edge polytope of a graph $G$ is the convex hull of
$\{\pm(e_i-e_j):ij\in E(G)\}$; write $N(G)$ for its number of facets.
Let $\CB(a,b,c)$ consist of three internally disjoint paths of positive lengths
$a,b,c$ between common endpoints, so $n=a+b+c-1$.
For $n\ge3$, take $t,u$ to be the two positive odd integers as close as
possible with sum $n+(n\bmod 2)$, and put $D(t)=t\binom{t-1}{(t-1)/2}$. Set
\[
 M(n)=2^{[n\ {\rm even}]}D(t)D(u).
\]
For $n\ge5$ this is the facet count of the wedge of the two odd cycles of
lengths $t,u$, with a pendant edge if $n$ is even. The displayed formula
also defines $M(3)=6$ and $M(4)=12$.
\begin{theorem}\label{r27showa:th}
If $a,b,c$ are not all of the same parity, then
$N(\CB(a,b,c))\le M(n)$.
\end{theorem}
This is Conjecture 32 of Braun and Bruegge \cite{r27showa:BB} (Conjecture 5.2 in
the preprint). The bound concerns mixed parity, not the different conjecture
identifying a maximizing theta graph. Put the two lengths of equal parity
first and arrange them so that $a\ge b$. For a centered binomial row set
\[
 R_m(j)=\binom m{(m+j)/2},\quad P_m=2^{-m}\binom m{\lfloor m/2\rfloor},
 \qquad F(r,s,v)=\sum_{j\in\mathbb Z}R_r(j)R_s(j)R_v(j),
\]
where entries with noninteger or out-of-range lower index are zero. Then
\begin{equation}\label{r27showa:count}
 N(\CB(a,b,c))=c F(a,b,c-1)+ab F(a-1,b-1,c).
\end{equation}
This is Corollary 25 of \cite{r27showa:BB}, including its length-one cases: $R_0$
vanishes away from zero, so a zero length inserts exactly the product of the
central entries of the remaining even rows. In particular $F(0,0,c)=R_c(0)$.

\subsection{An additive comparison}
For common-parity rows, symmetry and Vandermonde convolution give
$\sum_j R_a(j)R_b(j)=\binom{a+b}{(a+b)/2}$ (replace $j$ by $-j$ in
the second lower index). Bounding the remaining row by its peak proves
$F(a,b,c-1)\le 2^{a+b+c-1}P_{a+b}P_{c-1}$. Alternatively,
bound the first two rows by their peaks and sum the third. For the second
term of \eqref{r27showa:count} use this latter estimate with lengths $a-1,b-1,c$.
For $A,B$ the largest odd integers not exceeding $a,b$, respectively, and
$H(x)=xP_{x-1}$ for positive odd $x$, we obtain
\begin{equation}\label{r27showa:estimate}
 \frac{4N(\CB(a,b,c))}{2^{n+1}}
 \le H(A)H(B)+2cL,\qquad L=\min\{P_aP_b,P_{a+b}P_{c-1}\}.
\end{equation}
Indeed $aP_{a-1}=H(A)$, since for even $a$ the equality reduces to
$a\binom{a-1}{a/2-1}=2(a-1)\binom{a-2}{a/2-1}$.
The target in this normalization is $4M(n)/2^{n+1}=H(t)H(u)$.
If $a,b$ are odd then $t+u=A+B+c$; if even, then
$t+u=A+B+c+1$. Hence in either case
\begin{equation}\label{r27showa:moves}
 t+u=A+B+2h,\qquad h=\lceil c/2\rceil .
\end{equation}
\begin{lemma}\label{r27showa:increment}
Starting from odd lengths $X\ge Y\ge1$, adding two to the smaller
length increases $H(X)H(Y)$ by at least $1/2$, and by at least $9/16$ if
$Y\ge3$. At any fixed even total the product is maximal at the most
balanced odd pair.
\end{lemma}
\begin{proof}
The adjacent ratio is $H(Y+2)/H(Y)=(Y+2)/(Y+1)$. Thus the gain is
$H(X)H(Y)/(Y+1)\ge K(Y):=H(Y)^2/(Y+1)$.
Here $K(1)=1/2$, $K(3)=9/16$, and
$K(Y+2)/K(Y)=(Y+2)^2/((Y+1)(Y+3))>1$.
Finally, the adjacent ratios decrease with the length. Transferring two
from the longer length to the shorter when they differ by at least four
therefore increases the product.
\end{proof}
Starting at $(A,B)$, take $h$ steps of adding two to the smaller length,
reordering as necessary. By \eqref{r27showa:moves} and the final assertion of the
lemma, the resulting product is at most $H(t)H(u)$, even if the steps have
not brought the lengths into balance. Thus the available gain is at least
$h/2$ universally or $9h/16$ when $B\ge3$.

\subsection{Paying for the single-contraction term}
The sequence $P_m$ is nonincreasing: $P_{2j-1}=P_{2j}$ for $j\ge1$,
and $P_{2j+1}/P_{2j}=(2j+1)/(2j+2)$.
Since $2cL\le4hL$, the following three cases in \eqref{r27showa:estimate} are
covered by Lemma~\ref{r27showa:increment}.
\begin{itemize}
\item If $c\ge10$, then $P_{c-1}\le P_9=63/256<1/4$ and
$P_{a+b}\le P_2=1/2$, so $4L\le1/2$.
\item If $A\ge9$, then $P_a\le P_9<1/4$ and $P_b\le1/2$, again giving
$4L\le1/2$.
\item If $B\ge3$, then $P_a,P_b\le P_3=3/8$, so $4L\le9/16$.
\end{itemize}
What remains is $B=1$, $A\in\{1,3,5,7\}$ and $c<10$. The exact
differences $M(a+b+c-1)-N(\CB(a,b,c))$, evaluated by \eqref{r27showa:count},
are given below. For instance the entry $(a,b,c)=(3,1,2)$ is
$36-\{2F(3,1,1)+3F(2,0,2)\}=36-(12+12)=12$.
\begin{center}\small
\begin{tabular}{rr|rrrr}
$a$&$b$&$c=2$&$4$&$6$&$8$\\\hline
1&1&0&6&40&270\\
3&1&12&72&420&2100\\
5&1&80&480&2400&11900\\
7&1&480&2520&12600&58800
\end{tabular}\par\medskip
\begin{tabular}{rr|rrrrr}
$a$&$b$&$c=1$&$3$&$5$&$7$&$9$\\\hline
2&2&0&18&120&750&3864\\
4&2&12&120&800&4200&21560\\
6&2&80&750&4200&22050&105840\\
8&2&540&3864&21560&105840&508032
\end{tabular}
\end{center}
All entries are nonnegative, completing the proof of Theorem~\ref{r27showa:th}.
The length-one conventions in the source are covered even where duplicate
paths would be excluded for a simple bicyclic graph.
\paragraph{Consequence for simple graphs.}
For every $n\ge4$ the maximum facet count of a connected simple graph with
$n$ vertices and $n+1$ edges is $M(n)$. For the upper bound, Theorem 14
of \cite{r27showa:BB} handles edge-disjoint cycles, and Corollary 19 allows leaf
deletion and induction. The remaining leafless graphs are of the form
$\CB(a,b,c)$; Theorem 27 of \cite{r27showa:BB} handles equal parity and
Theorem~\ref{r27showa:th} handles mixed parity. At $n=4$ the only graph is
$K_4$ minus an edge, $\CB(2,2,1)$, which attains 12. For $n\ge5$ the
wedge construction above attains the bound. This yields the claim of
Conjecture 13 on its nonempty simple-graph domain; at $n=3$ no
four-edge simple graph exists.
\begin{thebibliography}{1}\small
\bibitem{r27showa:BB} B.~Braun and K.~Bruegge, Facets of symmetric edge polytopes
for graphs with few edges, \emph{J. Integer Seq.} \textbf{26} (2023),
23.7.2. \href{https://cs.uwaterloo.ca/journals/JIS/VOL26/Braun/braun6.pdf}{Journal version};
\href{https://arxiv.org/html/2201.13303v4}{arXiv:2201.13303v4}.
\end{thebibliography}

\par\endgroup

% Source: /workspace/scratch/81914cec2c78/work/archive/Tournament_Round_27_Final/Tournament_Round_27/round27/reiwa/paper.tex
% Edited source: /workspace/scratch/81914cec2c78/work/archive/Tournament_Round_27_Final/Tournament_Round_27/round27/reiwa/paper.tex
\begingroup

\Needspace{9\baselineskip}
\section{\texorpdfstring{All fixed logarithmic differences of five combinatorial sequences}{All fixed logarithmic differences of five combinatorial sequences}}\label{section:r27reiwa}
\noindent{\small\sffamily\color{muted}Round 27 \enspace / \enspace Reiwa}\par\medskip


\begin{sectionabstract}
For the Motzkin, Fine, central Delannoy, tree-like polyhex and Domb sequences, we obtain a strict asymptotic sign and its leading constant at every fixed order of logarithmic difference. An invariant interval for each consecutive-term ratio turns formal approximations into expansions with controlled remainders. This supplies a proof candidate for the five-family conjecture concluding the paper of Chen, Guo and Wang. The accompanying but separate Bell-number conjecture is not treated.
\end{sectionabstract}
\subsection{The target}
For an eventually positive sequence let $(Ra)_n=a_{n+1}/a_n$, and let $\Delta f(n)=f(n+1)-f(n)$. Then $\log (R^j a)_n=\Delta^j\log a_n$. Almost infinite log-monotonicity means that for each fixed order, after dropping finitely many initial terms, the successive ratio iterates are alternately log-convex and log-concave, starting with log-convexity at $j=0$. In particular, it suffices that $(-1)^k\Delta^k\log a_n>0$ eventually for every fixed $k\ge2$. To impose all orders through any fixed order one takes the maximum of finitely many thresholds.
\begin{theorem}
For each of the five sequences named above and each integer $k\ge2$,
\[
 (-1)^k\Delta^k\log a_n=\beta(k-1)!n^{-k}+O_k(n^{-k-1}).
\]
Here $\beta=1/2$ for the central Delannoy numbers and $\beta=3/2$ for the other four. This proves the first conjecture in the closing paragraphs of \cite{r27reiwa:CGW}, precisely as quantified there. No uniform starting index in $k$ is claimed.
\end{theorem}

\subsection{The five intervals}
Use the recurrences $a_n=A(n)a_{n-1}+B(n)a_{n-2}$ for $n\ge2$ recorded in Section~4 of \cite{r27reiwa:CGW}. In order the letters $M,F,C,T,D$ denote the five sequences of the theorem. In particular, $T_n$ counts the tree-like polyhexes of their recurrence with $n+1$ hexagons. All $a_0=1$.
\[
\begin{array}{c|c|c|c}
 & a_1&A(n)&B(n)\\\hline
 M&1&(2n+1)/(n+2)&3(n-1)/(n+2)\\
 F&0&(7n-5)/(2n+2)&(2n-1)/(n+1)\\
 C&3&3(2n-1)/n&-(n-1)/n\\
 T&1&3(2n-1)/(n+1)&-5(n-2)/(n+1)\\
 D&4&2(2n-1)(5n^2-5n+2)/n^3&-64(n-1)^3/n^3
\end{array}
\]
Set $r_n=a_n/a_{n-1}$ wherever defined, and $F_n(x)=A(n)+B(n)/x$. The following compact intervals $I=[h,H]$ contain $r_5$, and $F_n(I)\subset I$ for $n\ge6$. On $I$ the maps have Lipschitz constants at most $q$. The roots $L$ in the last column belong to the interiors and satisfy $L=A(\infty)+B(\infty)/L$.
\[
\begin{array}{c|c|c|c|c}
 &r_5&I&q&L\\\hline
 M&7/3 &[2,4]&3/4&3\\
 F&3 &[2,5]&1/2&4\\
 C&561/107 &[3,6]&1/9&3+2\sqrt2\\
 T&137/36 &[3,6]&5/9&5\\
 D&7876/679 &[10,20]&16/25&16
\end{array}
\]
Indeed $|F_n'(x)|\le |B(n)|/h^2\le q$ by inspection. In row order, the endpoint gaps $\min_I F_n-h$ and $H-\max_I F_n$ are
\[
\begin{array}{c|c|c}
M&3(n-5)/(4(n+2))&(n+17)/(2(n+2))\\
F&(19n-47)/(10(n+1))&(n+16)/(2(n+1))\\
C&8(n-1)/(3n)&(n+17)/(6n)\\
T&4(n-2)/(3(n+1))&(5n+44)/(6(n+1))\\
D&6(3n^3-9n^2-n+2)/(5n^3)&2(8n^3+51n^2-21n+2)/(5n^3).
\end{array}
\]
They are positive for $n\ge6$; for the last row group $3n^2(n-3)$ and $n(51n-21)$ respectively. Induction hence proves the claimed enclosure and positivity from $a_4,a_5>0$ even in the Fine case.

\subsection{Controlling the asymptotic expansion}
For each $m\ge1$ form $P_m(n)=L+\sum_{j=1}^m c_j n^{-j}$ by requiring
\[
 P_m(n)-F_n(P_m(n-1))=O(n^{-m-1}).
\]
These equations uniquely determine the $c_j$. Indeed with $t=1/n$ the shift replaces $t$ by $t/(1-t)$. All expressions are analytic near $t=0$ since $L>0$ and $A,B$ are rational and regular at infinity. Starting with $P_0=L$, changing the coefficient of $t^j$ by $d$ changes the order-$j$ defect by $(1+B(\infty)/L^2)d$, and changes no lower coefficient. This factor is nonzero because $|B(\infty)|/L^2<1$.

For all large $n$, $P_m(n-1)\in I$ since $L$ is in its interior. Thus $e_n=|r_n-P_m(n)|$ satisfies $e_n\le q e_{n-1}+K_m n^{-m-1}$. Iteration from $\lfloor n/2\rfloor$, using boundedness there, gives
\[
 e_n=O_m(q^{\lfloor n/2\rfloor})+O_m(n^{-m-1}/(1-q))=O_m(n^{-m-1}).
\]
This supplies the needed all-orders control on the actual branch, without differentiating an unspecified remainder.

Write $A(n)=A_0+A_1/n+O(n^{-2})$ and similarly for $B$. Coefficient comparison gives
\[
 c_1=(A_1+B_1/L)/(1+B_0/L^2).
\]
The pairs $(A_1,B_1)$ are respectively $(-3,-9),(-6,-3),(-3,1),(-9,15),(-30,192)$. Substitution gives $c_1/L=-\beta$ in every row. Taking $m=k$ and expanding the logarithm gives
\[
 \log r_n=\log L-\beta/n+\sum_{j=2}^k d_j n^{-j}+O_k(n^{-k-1}).
\]
Since $\Delta^k\log a_n=\Delta^{k-1}\log r_{n+1}$, only $-\beta/(n+1)$ contributes at order $n^{-k}$. Indeed, iterating the fundamental theorem of calculus gives $\Delta^p f(x)=\int_{[0,1]^p}f^{(p)}(x+\sum u_i)\,du$. Hence $\Delta^p x^{-1}=(-1)^pp! x^{-p-1}+O_p(x^{-p-2})$ and $\Delta^p x^{-j}=O_{p,j}(x^{-j-p})$. The finite difference of the remainder is still $O_k(n^{-k-1})$, without claiming any improvement. This proves the theorem.

\paragraph{Relation to earlier methods.}
Zhu \cite{r27reiwa:Zhu} studied low order inequalities for several of these sequences and all orders for derangements. Li \cite{r27reiwa:Li} already proved a general transfer from sufficiently deep quotient expansions to fixed-order log-monotonicity. Our specific contribution is to justify the expansions simultaneously in the five cases by explicit invariant intervals and contracted defects. The use of a two-solution recurrence discrepancy is informed by the discrete Abel comparison in VMRC manuscript 006 \cite{r27reiwa:VMRC}; the interval and asymptotic argument here is independent. An exact five-family prior publication was not located in the bounded search; this is not a claim of established publication priority.

\begin{thebibliography}{9}
\bibitem{r27reiwa:CGW} W. Y. C. Chen, J. J. F. Guo and L. X. W. Wang, Infinitely log-monotonic combinatorial sequences, \emph{Adv. Appl. Math.} \textbf{52} (2014), 99--120. \href{https://doi.org/10.1016/j.aam.2013.08.003}{Journal}; \href{https://arxiv.org/pdf/1304.5160}{preprint}, first closing conjecture, pp.~24--25.
\bibitem{r27reiwa:Zhu} B.-X. Zhu, Higher order log-monotonicity of combinatorial sequences, \href{https://arxiv.org/abs/1309.6025}{arXiv:1309.6025}.
\bibitem{r27reiwa:Li} G.-J. Li, $\ell$-log-monotonicity and Laguerre inequalities for P-recursive sequences, \href{https://arxiv.org/abs/2206.13922}{arXiv:2206.13922}, Thm.~2.2.
\bibitem{r27reiwa:VMRC} Keio, Minimal recurrence orders and integral solution lattices for Franel sequences, manuscript 006 in the \href{https://github.com/QEEDDD/Verifiable-Mathematical-Reasoning-Corpus}{Verifiable Mathematical Reasoning Corpus}.
\end{thebibliography}

\par\endgroup

% Source: /workspace/scratch/81914cec2c78/work/archive/Tournament_Round_27_Final/Tournament_Round_27/round27/keio/paper.tex
% Edited source: /workspace/scratch/81914cec2c78/work/archive/Tournament_Round_27_Final/Tournament_Round_27/round27/keio/paper.tex
\begingroup
\newcommand{\dd}{\,d}
\newcommand{\R}{\mathbb R}
\Needspace{9\baselineskip}
\section{\texorpdfstring{Shift quotients of the logarithmic mean and their reciprocal duals}{Shift quotients of the logarithmic mean and their reciprocal duals}}\label{section:r27keio}
\noindent{\small\sffamily\color{muted}Round 27 \enspace / \enspace Keio}\par\medskip

\begin{sectionabstract}
Feng Qi asked whether two related nontrivial logarithmic expressions are logarithmically completely monotone on the entire positive axis. We give explicit positive kernels for both negative logarithmic derivatives. In a two-exponent deformation of the pair, ordinary monotonicity already forces exactly the parameter conditions sufficient for logarithmic complete monotonicity in every order.
\end{sectionabstract}
\subsection{The paired conjecture}
All functions below have domain $(0,\infty)$. A smooth function $f$ is completely monotone (CM) if $(-1)^nf^{(n)}\ge0$ for every integer $n\ge0$. A positive smooth function is logarithmically completely monotone (LCM) if $-(\log f)'$ is CM. In Remark 4.4 of \cite{r27keio:Qi}, Qi posed both assertions for the pair
\[
 F(x)=\frac{x\log x}{(x^2-1)\log(1+x)},\qquad
 H(x)=\frac{(x^2-1)\log(1+1/x)}{x\log x}.
\]
Here $F(1)=1/(2\log2)$ and $H(1)=2\log2$. The earlier LCM results of Zhang, Ji and Qi \cite{r27keio:ZJQ} concern different functions, in particular $1/(xF(x))$, not $F$ itself. Qi's Bernstein property for $xF(x)$ \cite[Theorem 4.3]{r27keio:Qi} does imply that $F$ is CM (divide the Bernstein integral by $x$), but it does not by itself prove the all-order logarithmic assertions here. The elementary observation $H(x)=1/F(1/x)$ by itself does not transfer complete monotonicity between the functions.
\begin{theorem}\label{r27keio:main}
Both $F$ and $H$ are strictly LCM and strictly CM on $(0,\infty)$.
\end{theorem}
We prove this together with a sharp characterization. Write $L(x)=(x-1)/\log x$ with $L(1)=1$. The identity $L(x)=\int_0^1x^u\dd u$ shows that $L$ is positive and smooth, including at $1$.
\begin{theorem}\label{r27keio:general}
Fix $a>0$ and $\alpha,\beta\in\R$, and put
\[
 F_a(x)=(x+a)^{-\beta}\left(\frac{L(x+a)}{L(x)}\right)^\alpha,
 \qquad H_a(x)=1/F_a(1/x).
\]
The following are equivalent: (i) $F_a$ is nonincreasing; (ii) $H_a$ is nonincreasing; (iii) $\alpha,\beta\ge0$; (iv) both $F_a$ and $H_a$ are LCM. Under (iii) both are also CM, with strict inequalities in every positive ordinary and logarithmic derivative order unless $\alpha=\beta=0$.
\end{theorem}
\subsection{A Stieltjes kernel}
\begin{lemma}
For $q=(\log L)'$ one has
\[
 q(x)=\int_0^\infty\frac{d\mu(t)}{x+t},\qquad
 d\mu(t)=\frac{dt}{t((\log t)^2+\pi^2)} .
\]
The measure $\mu$ has total mass $1$.
\end{lemma}
\begin{proof}
For completeness apply Cauchy's formula to
\[
 q(z)=\frac{1}{z-1}-\frac{1}{z\log z}
\]
on the principal logarithmic slit plane; the singularity at $1$ is removable. On a large circle $|q(z)|=O(1/|z|)$; on a small circle of radius $r$ it is $O(1+1/(r|\log r|))$. Thus the circular parts of the keyhole formula for $q(x)$ vanish. On the upper lip $z=-t+i0$ the imaginary part is $-\pi/[t((\log t)^2+\pi^2)]$; the lower lip is its conjugate. The jump integral gives the displayed representation (first truncate the lips to $r<t<R$, then take limits). Finally the substitution $v=\log t$ gives $\int d\mu=\int_\R dv/(v^2+\pi^2)=1$. Such logarithmic Stieltjes densities also occur in \cite{r27keio:ZJQ,r27keio:Qi}; the proof here requires no conjectural input.
\end{proof}
\subsection{Proof of the characterization}
Logarithmic differentiation gives $P(x):=-(\log F_a)'(x)$ explicitly:
\begin{align}
 P(x)&=\frac{\beta}{x+a}+\alpha(q(x)-q(x+a))\notag\\
 &=\frac{\beta}{x+a}+\alpha a\int_0^\infty
 \frac{d\mu(t)}{(x+t)(x+t+a)}.\label{r27keio:kernel}
\end{align}
The reciprocal change of variable preserves precisely the product structure required here:
\begin{equation}
 -(\log H_a)'(x)=\frac{P(1/x)}{x^2}
 =\frac{\beta}{x(1+ax)}+\alpha a\int_0^\infty
 \frac{d\mu(t)}{(1+tx)(1+(t+a)x)}.\label{r27keio:dual}
\end{equation}
For $\alpha,\beta\ge0$, the elementary reciprocals on the right of both identities are CM, as are their products by the Leibniz rule. Integration preserves all signs: for $x\ge\delta>0$ the absolute value of the order $j$ derivative of each product is at most $(j+1)!\delta^{-j}$ times the product itself, which is bounded uniformly in $t$ there. This also justifies all differentiations under the finite measure. If either coefficient is positive, $(-1)^j P^{(j)}$ and the corresponding quantities for the dual are strictly positive for all $j\ge0$. If $U$ is either function and $V=-(\log U)'$, then $-U'=UV$; induction with Leibniz gives $(-1)^n U^{(n)}>0$ for $n\ge1$ in this nonconstant case. For $\alpha=\beta=0$ both functions are identically one.
    
Conversely, nonincrease of either function is equivalent to $P\ge0$ by the first equality in \eqref{r27keio:dual}. As $x\downarrow0$, $q(x)=1/(x-1)-1/(x\log x)\to+\infty$ whereas $q(x+a)$ stays bounded; thus $\alpha\ge0$. As $x\to\infty$, the integral in \eqref{r27keio:kernel} is at most $1/x^2$ since $\mu$ has mass one. Hence $\lim_{x\to\infty}xP(x)=\beta\ge0$. This proves Theorem~\ref{r27keio:general}. Taking $a=\alpha=\beta=1$ gives Theorem~\ref{r27keio:main} exactly. Neither the different conjecture in Remark 4.2 of \cite{r27keio:Qi} nor a Stieltjes assertion about $F$ or $H$ is needed or asserted.
\begin{thebibliography}{9}
\bibitem{r27keio:Qi} F. Qi, \emph{Bernstein and Stieltjes properties of some functions involving logarithmic functions}, arXiv:2610.40718v1 (2026), Remark 4.4, Eqs.~(4.3)--(4.4). \url{https://arxiv.org/html/2610.40718v1}.
\bibitem{r27keio:ZJQ} T.-Y. Zhang, A.-P. Ji and F. Qi, \emph{Two logarithmically completely monotonic functions involving logarithmic function}, Math. Inequal. Appl. \textbf{29} (2026), 313--318. \url{https://doi.org/10.7153/mia-2026-29-18}.
\end{thebibliography}

\par\endgroup

% Source: /workspace/scratch/81914cec2c78/work/archive/Tournament_Round_28_Final/Tournament_Round_28/round28/meiji/paper.tex
% Edited source: /workspace/scratch/81914cec2c78/work/archive/Tournament_Round_28_Final/Tournament_Round_28/round28/meiji/paper.tex
\begingroup
\newcommand{\cyc}{\mathrm{cyc}}
\Needspace{9\baselineskip}
\section{\texorpdfstring{A coefficient certificate for Liu's parameterized triangle inequality}{A coefficient certificate for Liu's parameterized triangle inequality}}\label{section:r28meiji}
\noindent{\small\sffamily\color{muted}Round 28 \enspace / \enspace Meiji}\par\medskip


\begin{sectionabstract}
Liu's 2014 conjecture on distances to the vertices and opposite sidelines
was proved by Yang, Chen, Huang, Wang and Lin in 2017. This note gives
an independent verification of the same result, not a new open-problem
resolution. Positivity on the line at infinity reduces the proof to a
single determinant. An ordered triangle parametrization turns its sign
into a check of 988 positive integer coefficients. The full
exact-arithmetic verifier is included.
\end{sectionabstract}
\subsection*{Statement and status}
Let $a,b,c$ be the respective lengths of $BC,CA,AB$ in a nondegenerate
Euclidean triangle. For any point $P$ in its plane, write $R_i$ for
its distance to vertex $i$ and $r_i$ for its distance to the
\emph{line} containing the opposite side. If $0<\lambda<1$, then
\begin{equation}
 \sum_{\cyc}\frac{R_1^2-r_1^2}{b^2+c^2+\lambda a^2}
 \ \geq\ \frac{3}{4(\lambda+2)}. \label{r28meiji:goal}
\end{equation}
This is precisely Conjecture~1, equation (4.1), in Liu~[1]. The exact
statement is Proposition~2 in Yang et al.~[2], who give a proof.
Thus the present calculation verifies a known theorem.
The bound is sharp: in an equilateral triangle at its center each
summand is $1/[4(\lambda+2)]$.

\subsection*{The reduction to one determinant}
Relabel so $a\ge b\ge c$ and scale to make $a=1$. Put
$A_1=1$, $A_2=b^2$, $A_3=c^2$, $K=A_1A_2A_3$, and set
\[
 D_i=A_j+A_k+\lambda A_i,\quad Q=D_1D_2D_3,\qquad
 H=2\!\sum_{i<j}A_i A_j-\sum_i A_i^2=16[ABC]^2.
\]
Indices $(i,j,k)$ run in cyclic order.
For barycentric coordinates $p=(p_1,p_2,p_3)$ with $\sum p_i=1$,
the cosine law and the altitude formula give
\[
 R_i^2=A_k p_j^2+A_j p_k^2+(A_j+A_k-A_i)p_jp_k,
 \qquad r_i^2=H p_i^2/(4A_i).
\]
No sign conditions on $p_i$ are imposed. Write $h=\sum_i p_i$.
Define $M$ as the symmetric matrix of the degree-two form obtained
by multiplying the gap in \eqref{r28meiji:goal} by $4KQ(2+\lambda)$, using
the two homogeneous formulas above and multiplying the constant term
by $h^2$. The scale factor is strictly positive.

On $h=0$, identify $p$ with the plane vector $v=\sum_i p_i V_i$,
where $V_i$ are the vertex position vectors. This identification is
injective since the triangle is nondegenerate. The restriction of the
unscaled quadratic form is
\[
 \sum_i \big(\lVert v\rVert^2-(n_i\cdot v)^2\big)/D_i,
\]
where $n_i$ is a unit normal to sideline $i$: these are the
degree-two parts in the displacement of the squared distances.
They are squared tangential projections, and no nonzero vector is
perpendicular to all the sides. Hence the restriction is positive
definite. In a basis beginning with two vectors spanning $h=0$, take
its positive definite $2\times2$ block $E$. The remaining scalar Schur
complement has the sign of $\det M$, since congruence only multiplies
the determinant by a positive square. Thus $\det M\geq0$ implies
$M\succeq0$ and proves \eqref{r28meiji:goal} for every $P$.

\subsection*{The positive coefficient certificate}
For the computation alone write $B=A_2$, $C=A_3$.
Exact division yields
\[
 \det M=H^2 B^2 C^2(1-\lambda)(2+\lambda)^2 F(B,C,\lambda),
\]
where $F$ has 234 nonzero terms and degrees at most 8 in $B,C$
combined and 6 in $\lambda$. To cover all admissible triangles,
divide the original ordered side lengths by $b+c-a>0$ and put
$x=(a-b)/(b+c-a)$ and $y=(b-c)/(b+c-a)$. The rescaled sides are
$2x+y+1,x+y+1,x+1$, in that order. Their squares $U,V,W$ satisfy
$B=V/U$, $C=W/U$. Writing $z=\lambda/(1-\lambda)$, define
\[
 G(x,y,z)=U^8(1+z)^6 F\!\left(V/U,W/U,z/(1+z)\right).
\]
The table gives the count of nonzero coefficients and the smallest
such coefficient in each $z$-degree. In particular $G\geq0$ for
$x,y,z\geq0$, so $F\geq0$ and $\det M\geq0$. This proves the claim.
\[
\begin{array}{c|rrrrrrr}
z\text{-degree}&0&1&2&3&4&5&6\\\hline
\text{count}&88&150&150&150&150&150&150\\
\text{minimum}&40&4&40&160&320&320&128
\end{array}
\]

\subsection*{References}
[1] J. Liu, A new inequality for a point in the plane of a triangle,
\emph{J. Math. Inequal.} \textbf{8} (2014), 597--611.
DOI: 10.7153/jmi-08-44.\\[3pt]
[2] Y. Yang, S. Chen, D. Huang, X. Wang and X. Lin,
Proof of an inequality conjecture for a point in the plane of a
triangle, \emph{J. Math. Inequal.} \textbf{11} (2017), 399--411.
DOI: 10.7153/jmi-11-34.


\subsection*{Reproducible calculation}
This Python/SymPy listing works throughout in integer polynomial
rings. In the first ring, the variables named \texttt{b,c} denote
squared normalized side lengths, not ordinary side lengths.
\begingroup\scriptsize
\begin{Verbatim}
# Exact coefficient certificate. Run with Python 3 and SymPy.
from sympy import ZZ
from sympy.polys.rings import ring
S,b,c,l = ring('b,c,l',ZZ)
A = [S.one,b,c]
D = [1+b+c-(1-l)*a for a in A]
H = 2*(b+c+b*c)-1-b*b-c*c
K = b*c
Q = D[0]*D[1]*D[2]
M = [[-3*K*Q for j in range(3)] for i in range(3)]
for i in range(3):
    j,k = (i+1)%3,(i+2)%3
    w = D[j]*D[k]*(2+l)
    M[i][i] -= A[j]*A[k]*H*w
    M[j][j] += 4*K*A[k]*w
    M[k][k] += 4*K*A[j]*w
    M[j][k] += 2*K*(A[j]+A[k]-A[i])*w
    M[k][j] = M[j][k]
u,t,s = M[0]
_,v,r = M[1]
_,_,w = M[2]
det = u*(v*w-r*r)-t*(t*w-r*s)+s*(t*r-v*s)
f,rem = det.div(H**2*b**2*c**2*(1-l)*(2+l)**2)
assert not rem and len(f)==234
d = max(i+j for i,j,k in f)
m = max(k for i,j,k in f)
assert (d,m)==(8,6)
T,x,y,z = ring('x,y,z',ZZ)
U,V,W = (2*x+y+1)**2,(x+y+1)**2,(x+1)**2
def pows(h,n): return [h**i for i in range(n+1)]
U,V,W = [pows(h,d) for h in (U,V,W)]
Z,B = pows(z,m),pows(1+z,m)
g = T.zero
for (i,j,k),h in f.items():
    g += h*U[d-i-j]*V[i]*W[j]*Z[k]*B[m-k]
assert len(g)==988 and min(g.values())>0
print("Exact quotient: 234 terms; degrees:",d,m)
for k in range(m+1):
    a = [h for (i,j,n),h in g.items() if n==k]
    print("z-degree:",k,"count:",len(a),"minimum:",min(a))
print("PASS:",len(g),"strictly positive coefficients.")
\end{Verbatim}

\endgroup

\par\endgroup

% Source: /workspace/scratch/81914cec2c78/work/archive/Tournament_Round_28_Final/Tournament_Round_28/round28/taisho/paper.tex
% Edited source: /workspace/scratch/81914cec2c78/work/archive/Tournament_Round_28_Final/Tournament_Round_28/round28/taisho/paper.tex
\begingroup
\newcommand{\slm}{\operatorname{arcsl}}
\newcommand{\slh}{\operatorname{arcslh}}
\newcommand{\tlm}{\operatorname{arctl}}
\newcommand{\tlh}{\operatorname{arctlh}}
\newcommand{\E}{\mathbb E}
\Needspace{9\baselineskip}
\section{\texorpdfstring{Coefficient reflection for the diagonal Pad\'e approximants of four lemniscate functions}{Coefficient reflection for the diagonal Pad\'e approximants of four lemniscate functions}}\label{section:r28taisho}
\noindent{\small\sffamily\color{muted}Round 28 \enspace / \enspace Taisho}\par\medskip

\begin{sectionabstract}
Liu and Chen conjectured two coefficientwise reflection identities for the normalized diagonal Pad\'e approximants of the ordinary and hyperbolic lemniscate functions. Both identities hold in every order. Reflection of the power series accounts for the signs; an explicit product measure proves that the normalized numerator and denominator are individually unique, so the identities hold before taking their quotient.
\end{sectionabstract}
\subsection{The two conjectures}
Put $z=x^4$. For an analytic germ $G$ with $G(0)=1$, call $(P,Q)$ a normalized diagonal pair of order $n$ if $\deg P,\deg Q\le n$, $P(0)=Q(0)=1$, and
\[
Q(z)G(z)-P(z)=O(z^{2n+1}).
\]
Equivalently, $G(x^4)-P(x^4)/Q(x^4)=O(x^{8n+4})$. For the four germs obtained by dividing the arc functions by $x$, use the following names for the coefficients of $z^j$ in $(P,Q)$, suppressing $n$:
\[
\begin{array}{c|cccc}
 &\slm x/x &\slh x/x &\tlm x/x &\tlh x/x\\\hline
 ([z^j]P,[z^j]Q)&(a_j,b_j)&(\alpha_j,\beta_j)&(p_j,q_j)&(r_j,s_j)
\end{array}
\]
\begin{theorem}[Liu--Chen Conjectures 2.1 and 2.2]
For all integers $n\ge1$ the four normalized pairs exist uniquely. For $1\le j\le n$,
\[
(a_j,b_j)=(-1)^j(\alpha_j,\beta_j),
\qquad (p_j,q_j)=(-1)^j(r_j,s_j).
\]
\end{theorem}
The assertions are (2.23)--(2.24) and (2.27)--(2.28) in \cite{r28taisho:LC}. The Sun--Chen inequality labelled Conjecture 1.1 in the same article was proved there and is not at issue here.

\subsection{A common pair of reflected germs}
Write $F_a(z)=\int_0^1(1-zu^4)^{-a}\,du$ for $|z|<1$. The arc sines $\slm x$ and $\slh x$ are the integrals from $0$ to $x$ of $(1-t^4)^{-1/2}$ and $(1+t^4)^{-1/2}$, respectively. For the arc tangents use
\[
\tlm x=\slm\big(x(1+x^4)^{-1/4}\big),\qquad
\tlh x=\slh\big(x(1-x^4)^{-1/4}\big).
\]
Differentiating the inner maps gives $(1\pm x^4)^{-5/4}$; the derivative of the corresponding outer map at the inner value is $(1\pm x^4)^{1/2}$. Integrating the products and scaling $t=xu$ gives all four germs:
\begin{align*}
(\slm x/x,\slh x/x)&=(F_{1/2}(z),F_{1/2}(-z)),\\
(\tlm x/x,\tlh x/x)&=(F_{3/4}(-z),F_{3/4}(z)).
\end{align*}

\subsection{Nondegeneracy and coefficient reflection}
Why not immediately substitute $-z$? That gives one matching pair, but equality of rational functions need not give equality of normalized polynomial pairs: both members of a lower-degree pair could share a new factor with constant term one. The following argument removes that ambiguity in every order.
\begin{lemma}
For $0<a<1$ and each $n\ge1$, the normalized diagonal pair for $F_a$ exists uniquely.
\end{lemma}
\begin{proof}
Termwise integration of the binomial series gives
\[
m_k=[z^k]F_a(z)=\frac{(a)_k}{k!(4k+1)}\qquad(k\ge0),
\]
where $(a)_k$ is the rising factorial. Let $U$ be uniform on $(0,1)$, and independently let $V$ have the density $v^{a-1}(1-v)^{-a}/\mathrm B(a,1-a)$ on $(0,1)$. The beta integral yields $\E V^k=(a)_k/k!$; also $\E U^{4k}=1/(4k+1)$. Thus $m_k=\E T^k$ for $T=U^4V$. The support of $T$ contains $(0,1)$: every open subinterval there contains products $u^4v$ from a nonempty open set in $(0,1)^2$ of positive probability.

Write $Q=1+\sum_{j=1}^n c_jz^j$. The $n$ coefficients of $QF_a$ above degree $n$ that must vanish give, for $0\le i<n$,
\[
\sum_{h=0}^{n-1} m_{1+i+h} c_{n-h}=-m_{n+1+i}.
\]
This real symmetric matrix is positive definite. Indeed, for any nonzero real polynomial $D$ of degree $<n$ its quadratic form on the coefficient vector of $D$ is $\E[T D(T)^2]>0$. Invertibility uniquely determines $Q$. The truncation of $QF_a$ through degree $n$ uniquely supplies $P$.
\end{proof}

\begin{proof}[Proof of the theorem]
If $(P,Q)$ matches $F_a(z)$, then $(P(-z),Q(-z))$ matches $F_a(-z)$ to the same order. Applying this involution to any other matching pair for the reflected germ would give another pair for $F_a$, contradicting the lemma. Consequently both polynomial coefficients, not just the quotient, acquire the factor $(-1)^j$. The two germ identifications give the stated identities with the indicated signs.
\end{proof}

\paragraph{An exact check.} For $a=1/2$, order $n=1$ gives $m_1=1/10$ and $m_2=1/24$, hence $Q=1-(5/12)z$ and $P=1-(19/60)z$. Both signs reverse for the hyperbolic sine. The accompanying implementation solves the original, unreversed systems separately for all four functions through order 16. These finite checks are not used in the proof.

\paragraph{Scope and attribution.}
Heisei independently found the same moment argument in this round. The proof is self-contained and claims no general priority for the Pad\'e or moment methods.
\begin{thebibliography}{1}
\bibitem{r28taisho:LC} J. Liu and C.-P. Chen, Pad\'e approximant related to inequalities for Gauss lemniscate functions, \emph{J. Inequal. Appl.} (2016), article 320. \url{https://doi.org/10.1186/s13660-016-1262-2}.
\end{thebibliography}

\par\endgroup

% Source: /workspace/scratch/81914cec2c78/work/archive/Tournament_Round_28_Final/Tournament_Round_28/round28/showa/showa.tex
% Edited source: /workspace/scratch/81914cec2c78/work/archive/Tournament_Round_28_Final/Tournament_Round_28/round28/showa/showa.tex
\begingroup
\newcommand{\one}{{\bf 1}}
\newcommand{\Var}{\operatorname{Var}}
\newcommand{\E}{\operatorname{E}}
\Needspace{9\baselineskip}
\section{\texorpdfstring{Quadratic variance convergence without a power remainder}{Quadratic variance convergence without a power remainder}}\label{section:r28showa}
\noindent{\small\sffamily\color{muted}Round 28 \enspace / \enspace Showa}\par\medskip

\begin{sectionabstract}
Pasteczka asked whether his exact variance ratio for iterated mean-type maps survives when their components are merely twice continuously differentiable near the diagonal. We prove it does: the variance itself controls the Taylor remainder once the expansion is made at the current arithmetic center. We also characterize residuality for quasiarithmetic means and construct a nonresidual mean with a twice continuously differentiable generator, as requested in his Remark~1. These are proof candidates for review, not a claim of priority.
\end{sectionabstract}
\subsection{The variance ratio}
Let $I$ be an open real interval. A mean on $I^p$ takes each input to a number between its smallest and largest entries. Write $\E z=p^{-1}\sum z_i$,
$Pz=z-(\E z)\one$ and $\Var(z)=p^{-1}\|Pz\|_2^2$. A mean $K$ is invariant for $F=(M_1,\ldots,M_p)$ when $K\circ F=K$. Local $C^2$ here means $C^2$ in an open neighborhood of the whole diagonal, not necessarily on all of $I^p$.
\begin{theorem}\label{r28showa:t:main}
Suppose the $M_i$ are symmetric, locally $C^2$ means on $I^p$ with a unique invariant mean $K$ for $F$. For each $z\in I^p$, either an iterate of $z$ is constant, or, putting $\mu=K(z)$,
\[
\lim_{n\to\infty}
\frac{\Var(F^{n+1}(z))}{\Var(F^n(z))^2}
=\tfrac14\Var\big(\xi_{M_1}(\mu),\ldots,\xi_{M_p}(\mu)\big).
\]
Here, for $p\ge2$, $\xi_M(x)=-p^2\partial_1\partial_2 M(x\one)$; the case $p=1$ always falls in the constant alternative.
\end{theorem}
The question is item~3 of Section~5 of \cite{r28showa:p}; that article proved this exact formula under a power bound on the remainders. Quadratic convergence as an order estimate is older. Our proof uses no additional smoothness assumption on $K$ and allows the limit above to be zero.

\begin{proof}
Write $z^{(n)}=F^n(z)$. Its coordinate minimum is nondecreasing and its coordinate maximum nonincreasing; both sequences stay in the closed interval spanned by $z$. Their limits, viewed as functions of $z$, are themselves invariant means (shifting an orbit by one does not change its limits). By uniqueness they both equal $K(z)$. Thus $z^{(n)}\to\mu\one$ without requiring continuity away from the diagonal.

For any component $M$, symmetry makes all first derivatives equal at $x\one$; their sum equals $1$ by $M(t\one)=t$. Write $d,b$ for the common diagonal and off-diagonal entries of its Hessian at $x\one$. Differentiating along the diagonal twice gives $d+(p-1)b=0$. Hence the Hessian annihilates $\one$, and its quadratic form on $h$ equals
\[
 (d-b)\|Ph\|_2^2 = -p^2 b\,\Var(h)=\xi_M(x)\Var(h).
\]
Set $m_n=\E z^{(n)}$ and $h_n=Pz^{(n)}$. The second-order Taylor formula with integral remainder at $m_n\one$ gives
\[
 M(z^{(n)})=m_n+\tfrac12\xi_M(m_n)\Var(z^{(n)})+r_{M,n}.
\]
Indeed $|r_{M,n}|$ is at most $\|h_n\|_2^2/2$ times the maximum difference, in operator norm, between the Hessians along the segment and at its starting point. The entire segment approaches $\mu\one$, where the Hessian is continuous, so this factor tends to zero. Since $\|h_n\|_2^2=p\Var(z^{(n)})$, we have $r_{M,n}=o(\Var(z^{(n)}))$ in the nonconstant alternative. In vector form, writing $\xi=(\xi_{M_i})_{i=1}^p$, this yields
\[
 \frac{Pz^{(n+1)}}{\Var(z^{(n)})}=\tfrac12 P\xi(m_n)+o(1)
 \longrightarrow\tfrac12 P\xi(\mu).
\]
Squaring and dividing by $p$ finishes the proof.
\end{proof}
\subsection{Testing the power remainder at two basepoints}
Fix $p\ge2$. For a symmetric $C^2$ mean define its remainder by subtracting
\[
 x+\E h+\tfrac12\xi_M(x)\Var(h)\quad\text{from }M(x\one+h).
\]
Say the remainder has order $2+\alpha$ locally if near each diagonal point its absolute value is at most $C\|h\|_\infty^{2+\alpha}$ uniformly for neighboring basepoints $x$ and sufficiently small $h$. Residuality in \cite{r28showa:p} is the existence of a positive exponent locally. Compactness of the unit sphere gives uniformity in directions; a multiplicative constant is absorbed by decreasing the exponent if necessary. For a fixed exponent we always use the displayed big-$O$ convention.

\begin{theorem}\label{r28showa:t:q}
For $0<\alpha\le1$ let $f\in C^2(I)$ with $f'$ nowhere zero. The quasiarithmetic mean $A_f(z)=f^{-1}(\E f(z_i))$ has local remainder of order $2+\alpha$ if and only if $f''$ is locally $\alpha$-H\"older. Thus $A_f$ is residual if and only if around each point $f''$ is H\"older of some positive order.
\end{theorem}
\begin{proof}
The inverse function theorem shows that $A_f$ is $C^2$. Differentiating $f(A_f(z))=p^{-1}\sum_i f(z_i)$ twice in two distinct coordinates at $x\one$ gives $\partial_1\partial_2 A_f(x\one)=-f''(x)/(p^2 f'(x))$, hence $\xi_{A_f}=f''/f'$.

For necessity there is a useful general observation. Given nearby distinct $x,y$, take the same tuple $z=(x,y,\ldots,y)$ and expand a mean at each of the basepoints $x\one,y\one$. The affine term at either is $\E z$; the variance in either expansion is $(p-1)|x-y|^2/p^2$. Subtract and divide by half this variance. The two remainders prove that $|\xi_M(x)-\xi_M(y)|\le 4Cp^2|x-y|^\alpha/(p-1)$. Therefore $f''=f'\xi_{A_f}$ is locally $\alpha$-H\"older, since $f'$ is locally Lipschitz and both factors are locally bounded.

Conversely, if $f''$ is locally $\alpha$-H\"older then $f$ and its inverse are locally $C^{2,\alpha}$: indeed $(f^{-1})''=-(f''/(f')^3)\circ f^{-1}$ and $f^{-1}$ is locally Lipschitz. Composing these maps in $A_f$ gives a locally $\alpha$-H\"older Hessian, so the Taylor integral remainder gives the asserted order, uniformly locally.
\end{proof}
\begin{corollary}
On $I=(-1/2,1/2)$ put $g(0)=0$, $g(x)=1/\log(e/|x|)$ otherwise, and take
$f(x)=x+\int_0^x(x-s)g(s)\,ds$. Then $A_f$ is locally $C^2$ but not residual.
\end{corollary}
\begin{proof}
We have $f''=g$, $f'(x)=1+\int_0^x g(s)\,ds>1/2$, and $g$ is continuous. Along $x=e^{-k}$, $g(x)/|x|^\alpha=e^{\alpha k}/(1+k)\to\infty$ for every $\alpha>0$. Apply Theorem~\ref{r28showa:t:q}. Moreover $g$ is monotone on each side of zero, hence locally of bounded variation. Thus bounded variation alone does not suffice, contrary to the separate assertion in \cite{r28showa:p}, Section~5, item~2.
\end{proof}
\noindent This answers the existential request in Remark~1, not a universal assertion that non-Lipschitz second derivatives force nonresiduality. We do not assert that the invariant mean itself is residual.
\begin{thebibliography}{9}
\bibitem{r28showa:p} P. Pasteczka, On the new smoothness class of means and its impact on mean-type mappings, \emph{Aequationes Math.} 100 (2026), 95. \href{https://doi.org/10.1007/s00010-026-01338-x}{DOI: 10.1007/s00010-026-01338-x}. Earlier version: \href{https://arxiv.org/abs/2406.12491}{arXiv:2406.12491}.
\end{thebibliography}

\par\endgroup

% Source: /workspace/scratch/81914cec2c78/work/archive/Tournament_Round_28_Final/Tournament_Round_28/round28/keio/paper.tex
% Edited source: /workspace/scratch/81914cec2c78/work/archive/Tournament_Round_28_Final/Tournament_Round_28/round28/keio/paper.tex
\begingroup
\newcommand{\bin}[2]{\binom{#1}{#2}}
\newcommand{\dd}{\displaystyle}
\Needspace{9\baselineskip}
\section{\texorpdfstring{Scalar log-convexity of rising-binomial power sums}{Scalar log-convexity of rising-binomial power sums}}\label{section:r28keio}
\noindent{\small\sffamily\color{muted}Round 28 \enspace / \enspace Keio}\par\medskip


\begin{sectionabstract}
The coefficientwise log-convexity of certain rising-binomial polynomials fails for the fifth power. After this counterexample, Z.-W. Sun conjectured separately that their values at one form a log-convex sequence for every integer power at least five. We prove this scalar assertion for every real power at least two, simultaneously at every real argument at least one. The proof compares the discrete curvature of the largest term with the change of the entire normalized tail. No assumption on the signs of individual coefficients of a Tur\'an difference is required.
\end{sectionabstract}
\subsection{A scalar conjecture}
Set
\[
 S_n^{(m)}(x)=\sum_{k=0}^n \bin{n+k}{k}^{m}x^k
 \qquad(n\ge0).
\]
Sun proposed coefficientwise log-convexity for integral $m\ge2$ in \cite{r28keio:mo,r28keio:sun}. Zhou exhibited a negative coefficient of $x^2$ at $m=5,n=7$ in \cite{r28keio:mo}. In a subsequent comment to Zhou's answer, Sun posed a separate conjecture: for integral $m\ge5$, the sequence $S_n^{(m)}(1)$ is log-convex. It is this separately stated scalar conjecture, not the coefficientwise conjecture, that the following theorem resolves.
\begin{theorem}
For any real $m\ge2$, $x\ge1$, and integer $n\ge1$,
\[
 S_{n-1}^{(m)}(x)S_{n+1}^{(m)}(x)> S_n^{(m)}(x)^2 .
\]
\end{theorem}
\subsection{Geometric control from the last summand}
Write $C_n=\bin{2n}n$ and $U_n=S_n^{(m)}(x)/(x^n C_n^m)$. Reversing the summation index gives
\[
 U_n=\sum_{j=0}^n x^{-j} R_{n,j}^m,\qquad
 R_{n,j}=\frac{\bin{2n-j}{n}}{C_n}=\prod_{\ell=0}^{j-1}f_\ell(n),
 \qquad f_\ell(n)=\frac{n-\ell}{2n-\ell}.
\]
In particular $U_n\ge1$. The initial factor is $1/2$ for positive $n$, and each factor in its indicated range lies in $[0,1/2]$.
\begin{lemma}
For $n\ge2$ and $t=2^{-m}/x$,
\[
 0\le U_n-U_{n-1}\le
 \frac{m}{(n-1)(n+1)}\frac{2t^2}{(1-t)^3}.
\]
\end{lemma}
\begin{proof}
For $0\le\ell\le n-1$ we have
\[
 f_\ell(n)-f_\ell(n-1)=
 \frac{\ell}{(2n-\ell)(2n-2-\ell)}
 \le\frac{\ell}{(n+1)(n-1)} .
\]
When $j=n$, the product at $n-1$ is interpreted as zero: its last factor $f_{n-1}(n-1)$ vanishes. Hence the comparison covers the new term as well. By telescoping a difference of products, every other factor being at most $1/2$,
\[
 R_{n,j}-R_{n-1,j}\le
 2^{1-j}\frac{j(j-1)}{2(n+1)(n-1)},\qquad j\ge1.
\]
Both products belong to $[0,2^{-j}]$. The derivative bound for $u\mapsto u^m$ on that interval, followed by summation, gives
\[
 U_n-U_{n-1}\le
 \frac{m}{(n+1)(n-1)}\sum_{j\ge1}j(j-1)t^j .
\]
The geometric-series identity for this sum proves the upper bound. Nonnegativity follows from the nonnegative differences of the factors.
\end{proof}
\begin{proof}[Proof of the theorem]
Assume $n\ge2$. Since $0<t\le1/4$, the geometric factor in the lemma is at most $8/27<1/3\le(n-1)/(2n-1)$. Since $U_{n-1}\ge1$, this gives
\[
 \frac{U_n}{U_{n-1}}<1+\frac{m}{(n+1)(2n-1)}
 \le \left(1+\frac{1}{(n+1)(2n-1)}\right)^m
 =\left(\frac{C_{n+1}C_{n-1}}{C_n^2}\right)^m .
\]
The middle inequality is Bernoulli's inequality; the last equality follows from $C_n/C_{n-1}=2(2n-1)/n$. Multiplying by $U_{n+1}/U_n\ge1$ in the comparison
$(C_{n+1}C_{n-1}/C_n^2)^m(U_{n-1}/U_n)>1$ yields the theorem.
For $n=1$, the difference is
\[
 (3^m-2\cdot2^m)x+(6^m-4^m)x^2>0
\]
because $(3/2)^m\ge9/4>2$.
\end{proof}
\paragraph{Scope.} Evaluation at each real $x\ge1$ is distinct from coefficientwise positivity. Nothing here asserts a new solution of the coefficientwise question. The local comparison with $C_n$ works even when $m$ is not an integer, whereas a counting interpretation for $m$ identical choices would not extend directly.
\begin{thebibliography}{9}
\bibitem{r28keio:mo} Z.-W. Sun, \emph{On the polynomials $\sum_{k=0}^n\binom{n+k}k^m q^k$}, MathOverflow question 479301 (2024), specifically the author's July 25, 2026 comment beneath Q. Zhou's answer. \url{https://mathoverflow.net/questions/479301/}.
\bibitem{r28keio:sun} Z.-W. Sun, A family of polynomials and related congruences and series, \emph{Frontiers in Combinatorics and Number Theory} \textbf{2} (2026), 1--22. \url{https://doi.org/10.3934/fcnt.2026009}.
\end{thebibliography}

\par\endgroup

% Source: /workspace/scratch/81914cec2c78/work/archive/Tournament_Round_29_Final/Tournament_Round_29/round29/meiji/paper.tex
% Edited source: /workspace/scratch/81914cec2c78/work/archive/Tournament_Round_29_Final/Tournament_Round_29/round29/meiji/paper.tex
\begingroup
\newcommand{\sech}{\operatorname{sech}}
\newcommand{\arcosh}{\operatorname{arcosh}}
\newcommand{\R}{\mathbb R}
\Needspace{9\baselineskip}
\section{\texorpdfstring{When do the two Li--Qi harmonic-mean bounds coexist?}{When do the two Li--Qi harmonic-mean bounds coexist?}}\label{section:r29meiji}
\noindent{\small\sffamily\color{muted}Round 29 \enspace / \enspace Meiji}\par\medskip


\begin{sectionabstract}
Li and Qi asked for cutoffs giving two strict lower bounds and a common upper bound for the harmonic mean of generalized cosines at reciprocal arguments. For a fixed base $a>1$ the requested configuration exists exactly when $a<4.32030\ldots$, where the threshold is specified by a scalar equation. We give cutoffs and an upper bound in the entire admissible range. A quartic obstruction excludes the critical base itself.
\end{sectionabstract}
\subsection{Statement and classification}
In Section 4 of \cite{r29meiji:LQ}, the same three constants $z_1,z_2,\alpha$ are sought for two pairs of strict inequalities. Write $t=\log a>0$, $C=\cosh t=(a^2+1)/(2a)$ and set
\[
 H_t(z)=\frac{2\cosh(tz)\cosh(t/z)}{\cosh(tz)+\cosh(t/z)},\qquad z>0.
\]
Since the source defines $\cosh_a z=(a^z+a^{-z})/2$, this is exactly its harmonic mean. The question, for each fixed $a>1$ and nonempty inner intervals $0<z_1<1<z_2$, is to arrange
\begin{align*}
 C < H_t(z)<\alpha &\quad (z_1<z<1\text{ or }1<z<z_2),\\
 2 < H_t(z)<\alpha &\quad (0<z<z_1\text{ or }z_2<z<\infty).
\end{align*}
The source does not request a least possible $\alpha$ or demand that the split points be critical points. Nantomah \cite[Thm.~2.6]{r29meiji:N} treated the base $a=e$.
\paragraph{Theorem.} Let $T$ be the unique positive zero of
\[
 L(t)=1/t+\coth t-2\tanh t .
\]
The requested constants exist if and only if $0<t<T$. In every admissible case take $\alpha=2C$ and $z_1=1/z_2$. Valid values of $z_2$ are:
\begin{itemize}
\item if $C<2$, the unique $z>1$ such that $2\sinh(tz/2)\sinh(t/(2z))=1$;
\item if $C\ge2$ and $t<T$, take $z_2=T/t$.
\end{itemize}
In particular, the parameter restriction cannot be omitted.

\subsection{Proof of sufficiency}
By reciprocal symmetry it is enough to work on $z>1$. Since the harmonic mean is less than twice its smaller argument, $H_t(z)<2\cosh(t/z)\le2C$.

\paragraph{The level $2$.}
Set $u=t(z+z^{-1})/2$, $v=t(z-z^{-1})/2$, $A=\cosh u$, $B=\cosh v$ and $\Delta=A-B$. Addition formulas give the exact identities
\begin{equation*}
 H_t(z)-2=\frac{\Delta^2-1}{AB},\qquad
 \Delta=2\sinh(tz/2)\sinh(t/(2z)).
\end{equation*}
Indeed, $\cosh(u+v)+\cosh(u-v)=2AB$ and the corresponding product is $A^2+B^2-1$. For $z>1$, $AB$ is strictly increasing, as are both arguments $u,v$. The logarithmic derivative of $\Delta$ with respect to $\log z$ is $g(tz/2)-g(t/(2z))>0$, where $g(y)=y\coth y$. Here $g'(y)>0$ follows from $\sinh(2y)>2y$ for $y>0$. We have $\Delta(1)=C-1$ and $\Delta(z)\to\infty$ because $\sinh(t/(2z))\ge t/(2z)$.

If $C<2$, $\Delta$ crosses $1$ exactly once. Until that crossing, $(1-\Delta^2)/(AB)$ strictly decreases, so $H_t(z)>H_t(1)=C$ for $1<z<z_2$. Beyond the crossing $H_t(z)>2$. If $C\ge2$, $H_t(z)>2$ for every $z>1$.

\paragraph{The central bound.}
Let $q(x)=x\sech x\tanh x$. Its logarithmic derivative is $L(x)$, which strictly decreases from $+\infty$ to $-1$, since
\[
 L'(x)=-x^{-2}-\sinh^{-2}x-2\sech^2x<0.
\]
Thus $q$ strictly increases up to $T$ and decreases thereafter. In the coordinate $s=\log z$, write
\[
 S_t(s)=\sech(te^s)+\sech(te^{-s})=2/H_t(e^s),
 \qquad S'_t(s)=q(te^{-s})-q(te^s).
\]
If $t<T$ and $0<s<\log(T/t)$, both arguments are below $T$ and hence $S_t'(s)<0$. This gives $H_t(z)>C$ throughout the claimed central interval in the second case. For consistency the two ranges above do not overlap any excluded parameter: at $t_0=\arcosh 2$ one has $L(t_0)=1/t_0-1/\sqrt 3>0$, since $\cosh\sqrt3>1+3/2>2$. In particular $t_0<T$.

\subsection{Necessity and the endpoint}
For $t>T$, $S_t''(0)=-2tq'(t)>0$, so $H_t(z)<C$ on a deleted neighborhood of $1$. No central cutoff can satisfy the required inequality. At $T$ this quadratic test vanishes. We compute the next even derivative.

Put $D=x\,d/dx$, $w=x\tanh x$, $b=x^2$ and $P_j=(D^j\sech x)/\sech x$. Starting with $P_0=1$, these polynomials are obtained by
\[
 P_{j+1}=(w+b-w^2)\partial_w P_j+2b\partial_bP_j-wP_j .
\]
In particular $P_2=2w^2-w-b$. Since $S_T''(0)=0$, at $x=T$ we have $b=2w^2-w$. For reproducibility, the fourth polynomial is
\begin{align*}
 P_4={}&5b^2-28bw^2+30bw-7b\\
       &+24w^4-36w^3+14w^2-w .
\end{align*}
Substitution gives
\[
 P_4\big|_{b=2w^2-w}=-w(2w-3)(2w-1)(3w-2)>0 .
\]
Indeed $b>w^2$ implies $w>1$. Also $w<T<3/2$: the sum of the first seven terms of $e^3$ exceeds $19$, so $y=\tanh(3/2)>9/10$. Hence $L(3/2)=2/3+1/y-2y<2/3+10/9-9/5<0$. Evenness and analyticity now give
\[
 S_T(s)-S_T(0)=\tfrac1{24}(2\sech T)P_4(T)s^4+O(s^6)>0
\]
for sufficiently small $s\ne0$. Thus the endpoint fails as well. This proves necessity without choosing the cutoffs symmetrically.

\paragraph{Numerical orientation.} $T\simeq1.46332548$ and $e^T\simeq4.32030274$. The split for $C<2$ is a level crossing, not a purported maximizer. Neither decimals nor numerical searches enter the proof.
\begin{thebibliography}{9}
\bibitem{r29meiji:LQ} W.-H. Li and F. Qi, \emph{Harmonic mean inequalities for generalized hyperbolic functions}, Montes Taurus J. Pure Appl. Math. 6(3) (2024), 199--207, Definition 2.1 and Section 4.
\href{https://mtjpamjournal.com/papers/article_id_mtjpam-d-23-00053/}{Publisher record}.
\bibitem{r29meiji:N} K. Nantomah, \emph{Harmonic mean inequalities for hyperbolic functions}, Earthline J. Math. Sci. 6(1) (2021), 117--129. \href{https://doi.org/10.34198/ejms.6121.117129}{doi:10.34198/ejms.6121.117129}.
\end{thebibliography}

\par\endgroup

% Source: /workspace/scratch/81914cec2c78/work/archive/Tournament_Round_29_Final/Tournament_Round_29/round29/taisho/paper.tex
% Edited source: /workspace/scratch/81914cec2c78/work/audit_26_33/cleaned_sources/r29_taisho.tex
\begingroup
\newcommand{\R}{\mathbb R}
\Needspace{9\baselineskip}
\section{\texorpdfstring{Strict parameter log-concavity of the generalized inverse hyperbolic sine}{Strict parameter log-concavity of the generalized inverse hyperbolic sine}}\label{section:r29taisho}
\noindent{\small\sffamily\color{muted}Round 29 \enspace / \enspace Taisho}\par\medskip

\begin{sectionabstract}
We prove the strict three-term product inequality in Conjecture 2.4 of Baricz, Bhayo and Vuorinen on its entire parameter range. More generally, the generalized inverse hyperbolic sine is strictly log-concave as a function of its positive parameter for every fixed positive argument. The proof uses logarithmic integration coordinates, a positive Hessian, and an explicit one-dimensional transport argument.
\end{sectionabstract}
\subsection{The product inequality}
For $x>0$ and $p>0$ define
\[
 F_x(p)=\operatorname{arcsinh}_p(x)=\int_0^x(1+t^p)^{-1/p}\,dt.
\]
The product clause in \cite[Conjecture~2.4]{r29taisho:bbv} asks whether, for each
$0<x<1$ and every $p>1$, one has $F_x(p)^2>F_x(p-1)F_x(p+1)$.
The following theorem proves this whole clause and a wider statement.


\begin{thm} For every $x>0$, $p\mapsto F_x(p)$ is strictly log-concave on $(0,\infty)$. In particular, for every $p>h>0$,
\[
 F_x(p)^2>F_x(p-h)F_x(p+h).
\]
\end{thm}
\begin{proof}
Put $b=\log x$. Substituting $t=e^s$ gives $F_x(p)=\int_{-\infty}^b k(p,s)\,ds$, where
\[
 H(p,s)=\frac{\log(1+e^{ps})}{p},\qquad k(p,s)=e^{s-H(p,s)} .
\]
For $u=e^{ps}/(1+e^{ps})$ put $V=u(1-u)>0$ and $E=-u\log u-(1-u)\log(1-u)>0$. We have $H_s=u$ and $H_p=-E/p^2$; hence
\[
 H_{ss}=pV,\qquad H_{ps}=sV,\qquad H_{pp}=s^2V/p+2E/p^3.
\]
Thus $\det\nabla^2H=2EV/p^2>0$, so $H$ is strictly convex on the convex domain $(0,\infty)\times\R$. Consequently $k$ is strictly jointly log-concave.

Here is a direct proof that strictness persists on integrating. Fix $p_0\ne p_1$, $0<\theta<1$, and write $k_i(s)=k(p_i,s)$, $M_i=F_x(p_i)$, $p_\theta=(1-\theta)p_0+\theta p_1$. Let $T$ be the increasing map of $(-\infty,b)$ onto itself that matches the cumulative masses of $k_0/M_0$ and $k_1/M_1$. Both densities are positive and continuous, so $T$ is continuously differentiable and
\[
 T'(s)=\frac{k_0(s)M_1}{k_1(T(s))M_0}>0.
\]
The function $z(s)=(1-\theta)s+\theta T(s)$ is also increasing onto $(-\infty,b)$, and $z'(s)\ge T'(s)^\theta$ by weighted AM--GM. Strict joint log-concavity at $(p_0,s)$ and $(p_1,T(s))$ gives
\[
 \begin{aligned}
 k(p_\theta,z(s))z'(s)
 &>k_0(s)^{1-\theta}k_1(T(s))^\theta T'(s)^\theta\\
 &=k_0(s)(M_1/M_0)^\theta .
 \end{aligned}
\]
Both sides are integrable, and the difference is continuous and positive. Integration proves $F_x(p_\theta)>M_0^{1-\theta}M_1^\theta$. Take $p_0=p-h$, $p_1=p+h$ and $\theta=1/2$.
\end{proof}
Taking $h=1$ proves the product inequality on the entire original range. The substitution $t=e^s$ is essential to this argument: pointwise log-concavity only in $p$ would not by itself imply log-concavity of the integral.

\paragraph{Scope.} The theorem proves the product clause; no ordinary-concavity conclusion is asserted.

\begin{thebibliography}{9}
\bibitem{r29taisho:bbv} \'A.~Baricz, B.~A.~Bhayo and M.~Vuorinen, Tur\'an type inequalities for generalized inverse trigonometric functions, \emph{Filomat} 29 (2015), 303--313, Conjecture 2.4. DOI: \url{10.2298/FIL1502303B}.
\end{thebibliography}

\par\endgroup

% Source: /workspace/scratch/81914cec2c78/work/archive/Tournament_Round_29_Final/Tournament_Round_29/round29/showa/paper.tex
% Edited source: /workspace/scratch/81914cec2c78/work/audit_26_33/cleaned_sources/r29_showa.tex
\begingroup

\Needspace{9\baselineskip}
\section{\texorpdfstring{Sparse Complements and Well-Total-Domination Conditions}{Sparse Complements and Well-Total-Domination Conditions}}\label{section:r29showa}
\noindent{\small\sffamily\color{muted}Round 29 \enspace / \enspace Showa}\par\medskip


\begin{sectionabstract}
We give elementary arguments for three well-total-domination implications from the Graffiti.pc list. The even-distance implication follows from a structural lemma: deleting one vertex to leave at most one missing edge forces every minimal total dominating set to have size two. The other two implications are independently reproved with their prior attribution retained.
\end{sectionabstract}

All graphs below are finite, simple and connected with at least two
vertices. A total dominating set $D$ supplies a neighbor in $D$ to
every vertex, even those in $D$. Write WTD if every inclusion-minimal
such set has minimum size. The three affirmative targets were listed on 4 March 2007
in E.~DeLaVi\~na's \emph{Written on the Wall II} [1].

\subsection*{The even-distance condition (318)}
Let $m$ be the number of nonedges in $G$. Entry 318 proposes that
$G$ is WTD if some vertex $v$ has at least $m$ vertices at even
distance from it. All such vertices except possibly $v$ are
nonneighbors of $v$. Consequently at most one nonedge avoids $v$,
regardless of whether the distance-zero vertex is counted. The
following structural lemma therefore proves the published implication.

\paragraph{Lemma.} If $R=V(G)\setminus\{v\}$ spans a clique minus at
most one edge, any total dominating set $D$ contains a total
dominating pair.
\paragraph{Proof.}
Take $p\in D$ adjacent to $v$. If another member $q\in D\cap R$
is adjacent to $p$, the pair $p,q$ totally dominates: they cover one
another, $p$ covers $v$, and no third vertex in $R$ misses both.
Otherwise domination of $p$ forces $v\in D$. If $D$ contains only
$p,v$, we are done. Every additional member of $D$ lies in $R$ and
misses $p$, so there is at most one, say $q$. It must be adjacent to
$v$ because it is not adjacent to $p$. Now $v,p$ cover one another,
$v$ covers $q$, and all remaining vertices are adjacent to $p$.
This proves the lemma.\hfill$\square$

\subsection*{The induced-tree condition (317; previously resolved)}
Write $t$ for the maximum order of an induced tree. If $m\le t$, an
induced $t$-vertex tree has $(t-1)(t-2)/2$ nonedges, so $t\le4$ and
$m\le4$. To prove WTD consider a minimal total dominating set $D$
of size $k$. For each $u\in D$, minimality supplies a distinct witness
$w$ with $N(w)\cap D=\{u\}$ (remove $u$ to see this).
If $k=3$, the induced subgraph on $D$ is a triangle or a three-vertex
path, as it has no isolates. In a triangle all witnesses lie outside,
giving at least six distinct nonedges. In a path the two ends need
distinct outside witnesses, giving four nonedges in addition to the
nonedge between the ends.

If $k\ge4$, each witness has at least $k-2$ actual nonedges into $D$:
a self-pair is not counted. Each missing edge occurs in at most two of
these counts. Hence $m\ge k(k-2)/2$. Only $k=4$ and $m=4$ are possible.
Equality requires all witnesses to lie in $D$ (an outside witness
has $k-1$ nonedges). Thus $G[D]$ is two disjoint edges, its four
nonedges exhausting the complement. No four vertices induce a tree:
on $D$ the graph is disconnected; any other set of four contains
a universal vertex and can induce a tree only if the remaining three
are independent, which is impossible. Thus $t<4=m$, a contradiction.
We conclude $k=2$ for every minimal $D$, proving 317.
This is a rederivation, not a new resolution [2].

\subsection*{The matching condition (326; previously resolved)}
Write $\mu$ for matching number and $S$ for the degree-two vertices.
Entry 326 asserts WTD if $3\mu(G)\le |E(G[S])|$. Each component
of $G[S]$ is a path or cycle. For a nontrivial path on $a$ vertices,
$a-1<3\lfloor a/2\rfloor$; for a cycle,
$a\le3\lfloor a/2\rfloor$ with equality only at $a=3$. Summing,
$|E(G[S])|\le3\mu(G[S])\le3\mu(G)$. Under the hypothesis equality
holds, and there must be a nontrivial component since $\mu(G)>0$.
It is a triangle. Each of its vertices already has degree two
inside this triangle; connectedness gives $G=K_3$, which is WTD.

\subsection*{Source and status notes}
\small
[1] E.~DeLaVi\~na, \href{https://www.uhd.edu/academics/sciences/mathematics-statistics/wowii/files/wowiiaii.pdf}{\emph{Written on the Wall II}}, entries 317, 318, 326. The November 2026 snapshot still marks them open; this is not evidence of priority.\par
[2] \href{https://www.researchgate.net/publication/411012117_Three_Well-Total-Domination_Conjectures_from_the_Written_on_the_Wall_II_List_Structural_resolutions_of_Conjectures_315_317_and_322}{\emph{Three Well-Total-Domination Conjectures from the Written on the Wall II List}}, Theorem 4.5.\par
[3] AI Village, \href{https://grok-ai-village-news-496089.gitlab.io/articles/opus5-wow2-conj327-false-one-hundred-fifty-seven}{August 2026 record of the previous resolution of 326}, citing commit b386990.\par
No exact earlier solution of 318 was located in the pre-stop search.
A similarly numbered entry in the original \emph{Written on the Wall}
(not volume II) concerns a different conjecture. No assertion of first
discovery is made here.

\par\endgroup

% Source: /workspace/scratch/81914cec2c78/work/archive/Tournament_Round_29_Final/Tournament_Round_29/round29/reiwa/paper.tex
% Edited source: /workspace/scratch/81914cec2c78/work/archive/Tournament_Round_29_Final/Tournament_Round_29/round29/reiwa/paper.tex
\begingroup

\Needspace{9\baselineskip}
\section{\texorpdfstring{A sharp lower bound for a weighted sum of triangle medians}{A sharp lower bound for a weighted sum of triangle medians}}\label{section:r29reiwa}
\noindent{\small\sffamily\color{muted}Round 29 \enspace / \enspace Reiwa}\par\medskip


\begin{sectionabstract}
For every nonobtuse Euclidean triangle the sum of $a^2m_a$, divided by
$abc$, is at least $1+\sqrt{10}/2$. Equality characterizes the right
isosceles triangle. We prove the bound proposed by Jian Liu by averaging the
squares of the two shorter sides while holding the longest side fixed. The
resulting one-variable profile has an interior maximum, not an interior
minimum.
\end{sectionabstract}

\subsection*{The inequality}
Let $a,b,c$ be the side lengths of a nondegenerate Euclidean triangle,
and let $m_a,m_b,m_c$ be the lengths of the medians to the corresponding
sides. Write $\sum_{\rm cyc}$ for the three-term cyclic sum.
\begin{theorem}
If the triangle is nonobtuse, then
\[
 \sum_{\rm cyc}a^2m_a\ \geq\ \frac{2+\sqrt{10}}2\,abc .
\]
Equality holds exactly when the triangle is right isosceles. On the
strictly acute domain the same constant is the unattained infimum of
$\sum_{\rm cyc}a^2m_a/(abc)$.
\end{theorem}
This gives the inequality in \cite[Conjecture 4, (4.4)]{r29reiwa:liu}. The source
states it for acute triangles while naming a right-isosceles equality
case; working on the nonobtuse domain makes the distinction explicit.

\subsection*{Averaging the shorter squared sides}
The ratio is symmetric and invariant under scaling. Set $c=1$ and
$0<b\le a\le1$, and write
\[
 a^2=u(1+t),\qquad b^2=u(1-t).
\]
Then $1/2\le u\le 1$, by nonobtuseness and $a,b\le1$, and $0\le t<1$.
Define
\[
 A=\sqrt{2+u},\qquad B=\frac{\sqrt{4u-1}}{u},\qquad
 k=\frac{3u}{2+u},\qquad q=\sqrt{1-t^2}.
\]
\begin{lemma}
At fixed $u$, the doubled ratio $D=2\sum_{\rm cyc}a^2m_a/(abc)$
satisfies $D\ge 2A+B$, with equality exactly for $t=0$.
\end{lemma}
\begin{proof}
Using $2m_a=\sqrt{2b^2+2c^2-a^2}$ and its cyclic forms gives
\[
 D=\frac A q\big[(1+t)\sqrt{1-kt}+(1-t)\sqrt{1+kt}\big]+\frac B q.
\]
For $t>0$ the expression in brackets decreases in $k\in[0,1]$:
its derivative is
\[
 \frac t2\left(\frac{1-t}{\sqrt{1+kt}}-
                      \frac{1+t}{\sqrt{1-kt}}\right)<0.
\]
As $k\le1$, lowering the bracket to its value at $k=1$ and dividing
by $q$ gives $\sqrt{1+t}+\sqrt{1-t}=\sqrt{2+2q}$.
Also $B\ge A$, since
\[
 u^2(B^2-A^2)=(1-u)(u^2+3u-1)\ge0.
\]
Thus
\[
 D-(2A+B)\ \ge\ A\left(\sqrt{2+2q}+q^{-1}-3\right).
\]
The function in parentheses is zero at $1$ and strictly decreasing
for $0<q\le1$, since its derivative is
$1/\sqrt{2+2q}-1/q^2<0$. Hence the bound is strict for $t>0$.
Direct substitution gives equality at $t=0$.
\end{proof}

\subsection*{The lower envelope}
It remains to minimize $h(u)=2\sqrt{2+u}+\sqrt{4u-1}/u$ on
$[1/2,1]$. We have
\[
 h'(u)=\frac1{\sqrt{2+u}}-\frac{2u-1}{u^2\sqrt{4u-1}} .
\]
Both competing terms are nonnegative. By clearing the positive
denominators, and then taking the difference of their squares, the
sign of $h'$ is the sign of
\begin{align*}
 u^4(4u-1)-(2u-1)^2(2+u)&=(u-1)Q(u),\\
 Q(u)&=4u^4+3u^3-u^2-5u+2.
\end{align*}
Now $Q(1/2)=-1/8$, $Q(1)=3$ and $Q''(u)=48u^2+18u-2>0$ throughout
this interval. Thus $Q$ has exactly one zero in $(1/2,1)$: its
nonpositive sublevel set is an interval, and strict convexity gives
strict negativity between its left endpoint $1/2$ and the zero.
Consequently $h$ rises and then falls. Its endpoint values are
\[
 h(1/2)=2+\sqrt{10}<3\sqrt3=h(1),
\]
where the strict comparison follows by squaring and using $13^2>160$.
The unique minimum of $h$ is therefore at $u=1/2$. Combining this
with the lemma proves the theorem, including the equality condition
$u=1/2$, $t=0$, or $(a^2,b^2,c^2)=(1/2,1/2,1)$ after normalization.
Finally $t=0$, $u\downarrow1/2$ with $u>1/2$ gives acute triangles
approaching equality, so the constant cannot be improved even there.

\begin{thebibliography}{1}
\bibitem{r29reiwa:liu}
J. Liu, Two inequalities involving circumradius, inradius and medians
of an acute triangle, \textit{Advances in Inequalities and Applications}
\textbf{2023} (2023), Article 1, p.~20.
\href{https://doi.org/10.28919/aia/8232}{doi:10.28919/aia/8232}.
\end{thebibliography}

\par\endgroup

% Source: /workspace/scratch/81914cec2c78/work/archive/Tournament_Round_29_Final/Tournament_Round_29/round29/keio/paper.tex
% Edited source: /workspace/scratch/81914cec2c78/work/audit_26_33/cleaned_sources/r29_keio.tex
\begingroup
\newcommand{\Nar}{\operatorname{Nar}}
\Needspace{9\baselineskip}
\section{\texorpdfstring{Two diagonals in the symmetric decomposition of refined Narayana polynomials}{Two diagonals in the symmetric decomposition of refined Narayana polynomials}}\label{section:r29keio}
\noindent{\small\sffamily\color{muted}Round 29 \enspace / \enspace Keio}\par\medskip


\begin{sectionabstract}
A coefficient reflection converts the symmetric decomposition of a Dyck-path enumerator into a first-difference problem. On the central diagonal and its upper neighbor, three components are evaluated by telescoping binomial products. This proves both identities of Conjecture~37(a) of B\'ona et al. and the first identity of Conjecture~37(b), with all original quantifiers.
\end{sectionabstract}
\subsection{Three affirmative evaluations}
For Dyck paths of semilength $n$, let $w_i$ count paths having $h$ occurrences of $UD$ and $i$ occurrences of $UUD$. Write $W(t)=W_{n,h}(t)=\sum_i w_it^i$. For $1\le h<n$, the published enumeration \cite[Theorem~2]{r29keio:bona} is
\begin{equation}
 w_i=\frac1h\binom{n}{h-1}\binom{n-h-1}{i-1}\binom{h}{i}.\label{r29keio:eq:counts}
\end{equation}
We extend binomial coefficients with nonnegative upper index by zero outside their usual range. Put $C_k=\binom{2k}{k}/(k+1)$ and define
\[
 N_m(t)=\frac1m\sum_i\binom mi\binom{m}{i-1}t^i\quad(m\ge1),\qquad N_0(t)=1,
\]
Thus $N_m$ is the ordinary Narayana polynomial. Let $A,D$ be the unique polynomials such that
\[
 W=A+tD,\qquad A(t)=t^hA(1/t),\qquad D(t)=t^{h-1}D(1/t).
\]
Following the definitions after Eq.~(11) in \cite{r29keio:bona}, $W^+_{n,h}$ and $W^-_{n,h}$ have coefficients respectively $|[t^i]A|$ and $|[t^i]D|$.
\begin{theorem}\label{r29keio:th:main}
For every integer $k\ge1$,
\[
 W^+_{2k,k}=(k-1)C_kN_{k-1}(t),\qquad
 W^-_{2k,k}=C_k\sum_{i=0}^{k-1}\binom{k-1}{i}^{\!2}t^i.
\]
For every integer $k\ge2$,
\[
 W^-_{2k,k+1}=\binom{2k}{k}N_{k-1}(t).
\]
\end{theorem}
The separate remaining identity of Conjecture~37(b) is not asserted here.
\subsection{Proof using coefficient reflection}
Let $I_h W=t^hW(1/t)$. Applying $I_h$ to $W=A+tD$ yields $I_hW=A+D$. Hence
\[
 A=\frac{W-t I_hW}{1-t},\qquad D=\frac{I_hW-W}{1-t}.
\]
These formulas also give existence (both numerators vanish at $1$), symmetry by reversal, and uniqueness. For $a_i=[t^i]A$, $e_i=[t^i]D$, we have
\begin{equation}
 a_i-a_{i-1}=w_i-w_{h+1-i},\qquad e_i-e_{i-1}=w_{h-i}-w_i.\label{r29keio:eq:diff}
\end{equation}
The initial coefficients $a_{-1}$ and $e_{-1}$ vanish. In both cases put $b_i=\binom{k-1}{i}$.
\paragraph{Central diagonal $h=k$.}
Equation~\eqref{r29keio:eq:counts} and Pascal's identity give $w_i=C_k b_{i-1}(b_i+b_{i-1})$. The sequences $C_k b_{i-1}b_i$ and $C_k b_i^2$ are symmetric under $i\mapsto k-i$ and $i\mapsto k-1-i$, respectively. Their corresponding polynomials therefore give $A$ and $D$ by uniqueness, since $a_i+e_{i-1}=w_i$. This is exactly the first pair of expressions (including $a_i=0$ for $k=1$).
\paragraph{Upper neighbor $h=k+1$, $k\ge2$.}
Write $d_i=\binom{k-2}{i}$, $v_i=\binom{k+1}{i}$ and $c=\binom{2k}{k}/(k+1)$. Equation~\eqref{r29keio:eq:counts} gives $w_i=c\,d_{i-1}v_i$, whence $w_{h-i}-w_i=c(d_{i-2}-d_{i-1})v_i$. We use
\[
 d_{i-1}-d_{i-2}=\frac{k+1-2i}{k-1}b_{i-1},\qquad
 b_i-b_{i-2}=\frac{k+1-2i}{k+1}v_i.
\]
The first follows by expressing each summand as a multiple of $b_{i-1}$; the second follows by writing the two multiples of $v_i$ with common denominator $k(k+1)$, with numerators $(k+1-i)(k-i)$ and $i(i-1)$. These identities give
\[
 w_{h-i}-w_i=-\frac{c(k+1)}{k-1}(b_i b_{i-1}-b_{i-1}b_{i-2}).
\]
Equation~\eqref{r29keio:eq:diff} now telescopes. Taking absolute values proves the third expression.
\begin{thebibliography}{2}
\bibitem{r29keio:bona} M. B\'ona, S. Dimitrov, G. Labelle, Y. Li, J. Pappe, A. R. Vindas-Mel\'endez and Y. Zhuang, A combinatorial proof of a symmetry for a refinement of the Narayana numbers, \emph{Electron. J. Combin.} 32 (2025), P2.46. \url{https://doi.org/10.37236/12681}.
\end{thebibliography}

\par\endgroup

% Source: /workspace/scratch/81914cec2c78/work/archive/Tournament_Round_30_Final/Tournament_Round_30/round30/meiji/paper.tex
% Edited source: /workspace/scratch/81914cec2c78/work/audit_26_33/cleaned_sources/r30_meiji.tex
\begingroup

\Needspace{9\baselineskip}
\section{\texorpdfstring{Monotonicity under every positive Hurwitz shift}{Monotonicity under every positive Hurwitz shift}}\label{section:r30meiji}
\noindent{\small\sffamily\color{muted}Round 30 \enspace / \enspace Meiji}\par\medskip


\begin{sectionabstract}
For Gao's difference $F_q(s)=\zeta(s,q)-\psi(1-s)$, we prove the monotonicity clause of Conjecture~2 for every $q>0$ and every $0<s<1$. A uniform positive derivative bound follows from interpolation between adjacent shifts and the established unshifted base case.
\end{sectionabstract}
\subsection{The exact question}
Let $\zeta(s,q)$ be the Hurwitz zeta function, and $\psi$ the digamma function. For $0<s<1$ and $q>0$, put
\[
 F_q(s)=\zeta(s,q)-\psi(1-s).
\]
The monotonicity clause of Gao's Conjecture~2 \cite{r30meiji:G} asks whether
$F_q$ is increasing throughout this interval for every positive shift.
The underlying theorem for $q=1$ is already given in \cite{r30meiji:G}. In particular it asserts $F_1''>0$ and $F'_1(0)=\zeta(2)-\tfrac12\log(2\pi)$. We use this established base case in the positive result below.
\paragraph{Theorem.} For every $q>0$ and $0<s<1$,
\[
 F'_q(s)\ge\tfrac58\zeta(2)-\tfrac12\log(2\pi)>0.
\]

\subsection{Interpolation between neighboring shifts}
The Hurwitz shift $F_{q+1}(s)=F_q(s)-q^{-s}$ holds by analytic continuation. Fix $s\in(0,1)$ and write $g(q)=F_q'(s)$. Differentiating with respect to $s$ gives
\[
 g(q+1)=g(q)+q^{-s}\log q,\qquad g(2)=g(1).
\]
For $1\le q\le2$, two shift derivatives of the Hurwitz function satisfy $\partial_q^2\zeta(s,q)=s(s+1)\zeta(s+2,q)$. It follows that
\[
 g''(q)=(2s+1)\zeta(s+2,q)+s(s+1)\partial_s\zeta(s+2,q)
 \le 3\zeta(2).
\]
In this display $\partial_s$ differentiates the first argument. The Dirichlet series for $\zeta(s+2,q)$ has summands at most $(n+1)^{-2}$ and its derivative has nonpositive summands. Subtracting $\tfrac32\zeta(2)(q-1)(q-2)$ now gives a concave function with both endpoint values $g(1)$. Consequently
\[
 g(q)\ge g(1)-\tfrac32\zeta(2)(q-1)(2-q)
       \ge g(1)-\tfrac38\zeta(2).
\]
Gao's base-case result quoted above bounds $g(1)$ below by $F_1'(0)$. Since $\pi>3.1$ and $2\pi<e^2$, the constant in the theorem is positive. For $q>2$, successive positive shifts from $[1,2]$ can only increase $g$. For $0<q<1$, the backward step $g(q)=g(q+1)-q^{-s}\log q$ increases it as well. This proves the claimed universal derivative bound.
\begin{thebibliography}{9}
\bibitem{r30meiji:G} L. Gao, \emph{Inequalities for $\zeta(s)-\psi(1-s)$ Related to a Conjecture of Henry}, Axioms 15 (2026), 577. \href{https://doi.org/10.3390/axioms15080577}{doi:10.3390/axioms15080577}.
\bibitem{r30meiji:D} NIST Digital Library of Mathematical Functions, \href{https://dlmf.nist.gov/25.11}{Section 25.11} (recurrence, continuation, special values and derivatives) and \href{https://dlmf.nist.gov/5.7}{Section 5.7} (digamma expansions).
\end{thebibliography}

\par\endgroup

% Source: /workspace/scratch/81914cec2c78/work/archive/Tournament_Round_30_Final/Tournament_Round_30/round30/taisho/paper.tex
% Edited source: /workspace/scratch/81914cec2c78/work/archive/Tournament_Round_30_Final/Tournament_Round_30/round30/taisho/paper.tex
\begingroup

\Needspace{9\baselineskip}
\section{\texorpdfstring{A short coefficient-tail proof for a known elliptic remainder}{A short coefficient-tail proof for a known elliptic remainder}}\label{section:r30taisho}
\noindent{\small\sffamily\color{muted}Round 30 \enspace / \enspace Taisho}\par\medskip

\begin{sectionabstract}
We give an alternative derivation for Zhao and Yang's question on absolute monotonicity of a normalized elliptic remainder. The 2025 theorem of Wang, Yang and Zhao had already answered the question: this note is not a new resolution. The argument combines the 2026 coefficient theorem of Deng, Kou and Zhong with endpoint normalization and two rational estimates.
\end{sectionabstract}
For $0<x<1$, put $p=\log 4$, and let $K(r)=\int_0^{\pi/2}(1-r^2\sin^2t)^{-1/2}\,dt$. Zhao and Yang \cite[Remark 8(iii)]{r30taisho:zy} asked whether
\[
    G(x)=\frac{F(x)-1}{1-x},\qquad
    F(x)=\frac{K(\sqrt{x})}{p-\frac12\log(1-x)}
\]
has nonnegative derivatives of every order there, including order zero.

\paragraph{An alternative affirmative proof.}
Write $F(x)=\sum_{n\ge0}A_nx^n$ for $|x|<1$. Deng, Kou and Zhong \cite[Theorem 1.1]{r30taisho:dkz} prove $A_n<0$ for $n\ge3$ (their normalization multiplies coefficients by $2/\pi$). The classical elliptic endpoint formula $K(\sqrt{x})=\log(4/\sqrt{1-x})+O((1-x)\log(1/(1-x)))$ gives $F(1^-)=1$. Monotone convergence applied to $\sum_{k\ge3}(-A_k)x^k$ as $x\uparrow1$ now yields $\sum_{k\ge0}A_k=1$, with absolute convergence. Consequently, for $n\ge2$,
\[
    [x^n]G=\sum_{k=0}^n A_k-1=-\sum_{k>n}A_k>0.
\]
The first two coefficients follow by expansion of $K(\sqrt x)=(\pi/2)(1+x/4+O(x^2))$:
\[
  [x^0]G=\frac{\pi}{2p}-1,\qquad
  [x^1]G=\frac{\pi(5p-2)}{8p^2}-1.
\]
Both are strictly positive: $1.386<p<1.387$, $\pi>3.14$, and $3.14(5\cdot1.386-2)>8(1.387)^2$. For instance, these bounds on $p$ follow from the first five terms of $4\sum_{j\ge0} 3^{-2j-1}/(2j+1)$ and a geometric tail estimate. Termwise differentiation gives $G^{(k)}(x)>0$ for every $k\ge0$ and $0<x<1$.

\paragraph{Prior-result check.}
Theorem 1 of \cite{r30taisho:wyz} establishes absolute monotonicity of $[(B\,{}_2F_1(a,b;a+b;x))/(P-\log(1-x))-1]/(1-x)$ for $0<a,b<1$, $a+b\le1$ and $2\le P\le R(a,b)$. Here $B=B(a,b)$ and $R=-\psi(a)-\psi(b)-2\gamma$. At $a=b=1/2$, $B=\pi$, $R=4\log 2=2p$; taking $P=R$ gives exactly $G$. Thus the original question was not open when this argument was obtained.
\begin{thebibliography}{9}
\bibitem{r30taisho:zy} T.~Zhao and Z.-H.~Yang, \emph{Absolutely monotonic functions related to the asymptotic formula for the complete elliptic integral of the first kind}, arXiv:2405.19651 (2024).
\bibitem{r30taisho:wyz} M.-K.~Wang, Z.-H.~Yang and T.-H.~Zhao, \emph{An affirmative answer to an open problem on Ramanujan's asymptotic formula of zero-balanced hypergeometric function}, arXiv:2504.04521 (2025).
\bibitem{r30taisho:dkz} Q.~Deng, Y.~Kou and F.~Zhong, \emph{Endpoint absolute monotonicity for the complete elliptic integral of the first kind}, arXiv:2609.32366 (2026).
\end{thebibliography}

\par\endgroup

% Source: /workspace/scratch/81914cec2c78/work/archive/Tournament_Round_30_Final/Tournament_Round_30/round30/showa/paper.tex
% Edited source: /workspace/scratch/81914cec2c78/work/archive/Tournament_Round_30_Final/Tournament_Round_30/round30/showa/paper.tex
\begingroup
\newcommand{\gt}{\gamma_{\mathrm t}}
\Needspace{9\baselineskip}
\section{\texorpdfstring{Total Domination in a Graph and Its Degree Complement}{Total Domination in a Graph and Its Degree Complement}}\label{section:r30showa}
\noindent{\small\sffamily\color{muted}Round 30 \enspace / \enspace Showa}\par\medskip


\begin{sectionabstract}
We prove Pinhe Chen's Conjecture 4 for total domination in a graph
and its degree complement. In fact, if the graph has order $n$, has
$k$ degree classes, and $s_1$ of those classes are singletons, the sum of
the total domination numbers is at most $n+2k-s_1$. The essential choice
is to cover a class from the ambient graph instead of seeking a total
dominating set inside a class that contains an isolate. The proof
applies to any pair of graphs that induce ordinary complements on each
block of a partition; no agreement between blocks is necessary.
\end{sectionabstract}

\subsection*{The published question}
For an isolate-free finite simple graph $G$, a set $D$ is totally
dominating if every vertex, including every vertex of $D$, has a
neighbor in $D$. The minimum size is $\gt(G)$. Its degree complement
$G_\delta$ toggles adjacency between different vertices of equal
degree in $G$; adjacency of unequal-degree pairs is preserved.
Chen [1, Conjecture~4] proposed that whenever both ambient graphs are
isolate-free,
\[
                 \gt(G)+\gt(G_\delta)\le n+2k,
\]
where $k$ is the number of distinct degrees in $G$. Theorem~6 in [1]
proves this when neither graph has an isolate inside any degree class.
That assumption is not automatic: the leaves of a star form an
independent degree class. We do not assume it.

\subsection*{The partition theorem}
\textbf{Theorem.}
Let $A,B$ be isolate-free finite simple graphs on the same set of $n$
vertices. Partition the set into $k$ nonempty blocks, $s_1$ of them
singletons. Suppose $B[C]=\overline{A[C]}$ on each block $C$.
Then
\[
                    \gt(A)+\gt(B)\le n+2k-s_1 .
\]

\textit{Proof.}
To cover $C$ in $A$ means to choose a set $U_C$ of ambient vertices
such that every vertex of $C$ has a neighbor in $U_C$. Vertices of
$U_C$ outside $C$ need not be covered at this step. Define $V_C$
likewise for $B$. We shall pay at most $|C|+2$ for the two covers
together, or two if $C$ is a singleton.

For $C=\{c\}$, choose one neighbor of $c$ in each ambient graph.
Now suppose $|C|\ge2$. If both induced graphs on $C$ are isolate-free,
the classical total-domination inequality of Cockayne, Dawes and
Hedetniemi [2], as stated in [1, Theorem~5], gives
\[
         \gt(A[C])+\gt(B[C])\le |C|+2 .
\]
Minimum total dominating sets inside the block serve as covers.

Otherwise exchange $A,B$ if needed and let $I\ne\varnothing$ be the
set of isolated vertices of $A[C]$. All the vertices of $C\setminus I$
together cover $C\setminus I$ in $A$: a neighbor inside $C$ is
necessarily a nonisolate there. To cover each $x\in I$, choose an
ambient $A$-neighbor $a_x$, possible since $A$ has no isolates.
Thus $U_C=(C\setminus I)\cup\{a_x:x\in I\}$ covers $C$ in $A$ at cost
at most $|C|$. In $B[C]$, any $x\in I$ is universal. For any other
$y\in C$, the adjacent pair $x,y$ covers $C$ in $B$, at cost two.
This includes the case $I=C$.

Finally $\bigcup_C U_C$ totally dominates $A$: every vertex belongs
to a covered block. In particular this covers the outside witnesses
used by other blocks. The analogous union totally dominates $B$.
Summing the block costs proves the result, as a singleton costs two
instead of three.\hfill$\square$

\subsection*{Consequences and scope}
Take $A=G$ and $B=G_\delta$ and partition by the original degrees in $G$.
The induced graphs are complementary by definition. The theorem proves
the full conjecture, with the singleton correction. No connectedness is
needed. The proof neither counts a vertex as its own neighbor nor
invokes the ordinary-complement theorem on a graph with isolated vertices.

For comparison, Chen's Example~3 [1] is $G=2K_2\sqcup C_4$. The two
classes have orders four, with respective degrees one and two; they
exchange induced subgraphs under degree complementation. In both
graphs the two isolated edges force all their endpoints, and a pair
of adjacent cycle vertices is necessary and sufficient on the cycle.
Hence both total domination numbers are six, and their sum is
$12=8+2\cdot2$. This existing sharpness example is not claimed as new.
We do not classify equality or claim a historical priority certificate.

\begin{thebibliography}{9}
\bibitem{r30showa:Chen} P. Chen, Refined Nordhaus--Gaddum-type bounds for Roman
and total domination on $\delta$-complement graphs, \emph{Mathematics}
\textbf{14} (2026), 1797.
\href{https://doi.org/10.3390/math14111797}{doi:10.3390/math14111797}.
\bibitem{r30showa:CDH} E. J. Cockayne, R. M. Dawes and S. T. Hedetniemi,
Total domination in graphs, \emph{Networks} \textbf{10} (1980), 211--219.
\end{thebibliography}

\par\endgroup

% Source: /workspace/scratch/81914cec2c78/work/archive/Tournament_Round_30_Final/Tournament_Round_30/round30/heisei/paper.tex
% Edited source: /workspace/scratch/81914cec2c78/work/audit_26_33/cleaned_sources/r30_heisei.tex
\begingroup
\newcommand{\dd}{\,\mathrm d}
\Needspace{9\baselineskip}
\section{\texorpdfstring{The sharp common interval of logarithmic concavity for normalized arctangent tails}{The sharp common interval of logarithmic concavity for normalized arctangent tails}}\label{section:r30heisei}
\noindent{\small\sffamily\color{muted}Round 30 \enspace / \enspace Heisei}\par\medskip

\begin{sectionabstract}
All normalized even-order Taylor remainders of the arctangent have strictly
negative logarithmic curvature on $[-1,1]$, proving the Seventh Conjecture of
Feng Qi's survey. For each fixed $x>1$ the curvature is positive for all
sufficiently high orders. We give two proofs of the strict sign: a short
Gauss-quotient argument and a direct computation of an increasing phase
on the branch cut of a beta integral.
\end{sectionabstract}
\subsection{The normalized remainder}
For $n=0,1,\ldots$, denote the normalized arctangent remainder of Liu and
Qi \cite{r30heisei:LQ} by $T_n(x)=T_{2n}[\arctan x]$. Thus $T_n(0)=1$, and otherwise
\[
 T_n(x)=\frac{(-1)^n(2n+1)}{x^{2n+1}}
 \left(\arctan x-\sum_{k=0}^{n-1}\frac{(-1)^kx^{2k+1}}{2k+1}\right).
\]
The empty sum is zero. Qi's Seventh Conjecture \cite[\S12.7.1.7]{r30heisei:Q} asks
whether $T_n$ is log-concave on $(0,1)$ for every $n$.
\begin{theorem}\label{r30heisei:main}
For every integer $n\ge0$ and $|x|\le1$, $(\log T_n(x))''<0$.
For each $|x|>1$, $(\log T_n(x))''>0$ for all sufficiently large $n$.
Thus $[-1,1]$ is the exact common set of negative logarithmic curvature.
\end{theorem}
Integrating finite geometric division of $1/(1+u^2)$ from $0$ to $x$ yields
$T_n(x)=(2n+1)\int_0^1 v^{2n}/(1+x^2v^2)\dd v$. With
$b=n+1/2$ and $t=v^2$, this becomes $T_n(x)=F_b(-x^2)$, where
\[
 F_b(z)= b\int_0^1\frac{t^{b-1}}{1-zt}\dd t
       ={}_2F_1(1,b;b+1;z),\quad z\notin[1,\infty),\qquad b>0 .
\]
The power series has coefficients $b/(b+j)$ at degree $j$, giving the
hypergeometric identity near zero and hence throughout the slit plane.
In particular $T_n(x)>0$ for all real $x$.
\subsection{A proof using Gauss quotients}
Set $c=b+1$. By K\"ustner's Theorem 1.5 \cite[p.~606]{r30heisei:K}, there exists
a probability measure $\nu_b$ on $[0,1]$ for which
\[
 \frac{F_b'(z)}{F_b(z)}
 =\frac{b}{c}\frac{{}_2F_1(2,b+1;c+1;z)}{{}_2F_1(1,b;c;z)}
 =\frac{b}{c}\int_0^1\frac{\dd\nu_b(t)}{1-tz}.
\]
The first equality is Gauss's differentiation formula; the measure theorem
requires $-1\le 1\le c$ and $0<b\le c$, both satisfied.
Differentiation at $z=-x^2$ gives
\begin{equation}\label{r30heisei:curv}
 (\log T_n(x))''=-\frac{2b}{c}\int_0^1
                   \frac{1-tx^2}{(1+tx^2)^2}\dd\nu_b(t).
\end{equation}
Differentiation under this finite measure is justified locally by bounded
continuous derivatives. The integrand is strictly positive for $|x|<1$.
Even at $|x|=1$ the integral is strictly positive: $\nu_b$ is not the unit
atom at $1$. Indeed its first moment, obtained by differentiating the
normalized quotient at $0$, is
\[
 2(b+1)/(b+2)-b/(b+1)=1-\frac{b}{(b+1)(b+2)}<1.
\]
\subsection{Direct proof from the jump of a beta integral}
For completeness we derive the sign without the hypergeometric quotient
theorem. For $z$ in the upper half-plane $F_b(z)$ has strictly positive
imaginary part, by its beta integral. Hence the principal logarithm of
$F_b$ is analytic there and its imaginary part lies between $0$ and $\pi$.
Its boundary argument vanishes on $(-\infty,1)$. For $s>1$ set
\[
 A_b(s)= -\log(s-1)+
                \int_0^s\frac{u^{b-1}-1}{1-u}\dd u,\qquad
 \theta_b(s)=\arg(A_b(s)+i\pi).
\]
By substituting $u=st$ in the upper boundary limit of the beta integral,
and splitting off the value $u^{b-1}=1$ at its pole, we obtain
$F_b(s+i0)=b s^{-b}\{A_b(s)+i\pi\}$. More explicitly, the constant
piece integrates to $-\log(1-s-i0)=-\log(s-1)+i\pi$; the remaining
integrand has a removable singularity, and converges by domination
(local Lipschitz continuity at $u=1$ suffices).
The omitted factor $1/(1+i\epsilon/s)$ in the constant-piece integral
tends to $1$. Thus $\theta_b$ is the upper boundary argument.
Since $A_b(1+)=+\infty$ and $A_b'(s)=s^{b-1}/(1-s)$, we have
\begin{equation}\label{r30heisei:jump}
 \theta_b(1+)=0,\qquad
 \theta_b'(s)=\frac{\pi s^{b-1}}{(s-1)(A_b(s)^2+\pi^2)}>0 .
\end{equation}
The bounded harmonic function $\operatorname{Im}\log F_b$ is the Poisson
integral of its boundary values. Comparing imaginary parts, this gives
\[
 \log F_b(z)=C+\frac1\pi\int_1^\infty\theta_b(s)
       \left(\frac{1}{s-z}-\frac{s}{1+s^2}\right)\dd s
\]
in the upper half-plane, for a real constant $C$. Indeed the imaginary part
of the first Cauchy kernel is the Poisson kernel times $\pi$ with its usual
normalization, and the subtracted term ensures convergence at infinity.
This formula continues to real $z<1$: the integral extends analytically
there, and $F_b>0$ there. Differentiate and integrate by parts using
\eqref{r30heisei:jump} and $0<\theta_b<\pi$:
\[
 \frac{F_b'(z)}{F_b(z)}=\frac1\pi\int_1^\infty
                      \frac{\theta_b'(s)}{s-z}\dd s .
\]
In particular $\theta_b'$ is integrable since its primitive is bounded
and increasing. Two differentiations of the composite logarithm now yield
\[
 (\log F_b(-x^2))''=-\frac2\pi\int_1^\infty
         \frac{s-x^2}{(s+x^2)^2}\theta_b'(s)\dd s.
\]
This is strictly negative on $|x|\le1$, including both endpoints.
\subsection{Sharpness over the order}
The measure $bt^{b-1}\dd t$ converges weakly to the unit mass at $1$ as
$b\to\infty$: it gives $[0,r]$ mass $r^b$ for $0\le r<1$.
The functions $(1+x^2t)^{-1}$ and its first two $x$ derivatives
are continuous bounded functions of $t\in[0,1]$ for every fixed $x$.
Consequently $T_n$ and its first two derivatives converge pointwise to
those of $(1+x^2)^{-1}$. Since the limit is positive,
\[
 \lim_{n\to\infty}(\log T_n(x))''
     =\frac{2(x^2-1)}{(1+x^2)^2}.
\]
This proves the second assertion of Theorem~\ref{r30heisei:main}. It does not assert
that $|x|=1$ is the inflection point of any fixed remainder.
\begin{thebibliography}{9}
\bibitem{r30heisei:Q} F. Qi, Theory of normalized remainders in Taylor series expansions,
\S12.7.1.7, p.~385. \href{https://arxiv.org/abs/2512.15797}{arXiv:2512.15797}.
\bibitem{r30heisei:LQ} X.-L. Liu and F. Qi, Logarithmically absolute monotonicity and
series expansions of normalized remainders for inverse sine and tangent
functions, \emph{Commun. Korean Math. Soc.} (2026).
\bibitem{r30heisei:K} R. K\"ustner, Mapping properties of hypergeometric functions and
convolutions of starlike or convex functions of order $\alpha$, \emph{Comput.
Methods Funct. Theory} \textbf{2} (2002), 597--610.
\href{https://doi.org/10.1007/BF03321867}{doi:10.1007/BF03321867}.
\end{thebibliography}

\par\endgroup

% Source: /workspace/scratch/81914cec2c78/work/archive/Tournament_Round_30_Final/Tournament_Round_30/round30/keio/paper.tex
% Edited source: /workspace/scratch/81914cec2c78/work/archive/Tournament_Round_30_Final/Tournament_Round_30/round30/keio/paper.tex
\begingroup
\newcommand{\ii}{\mathrm i}
\Needspace{9\baselineskip}
\section{\texorpdfstring{The zeros of the row polynomials of the Delannoy matrix square}{The zeros of the row polynomials of the Delannoy matrix square}}\label{section:r30keio}
\noindent{\small\sffamily\color{muted}Round 30 \enspace / \enspace Keio}\par\medskip


\begin{sectionabstract}
We prove the Delannoy assertion in Mao and Wang's Conjecture~4.1: in every positive degree the row polynomial of the ordinary matrix square of the Delannoy triangle has only negative real zeros. In fact they are simple. A previously documented divided-difference and circle argument counts the zeros in a central interval. We finish the count by comparing two Chebyshev amplitudes between consecutive zeros of an ordinary Delannoy polynomial. The Eulerian matrix-square assertion is not addressed.
\end{sectionabstract}

\subsection{The matrix square as a divided difference}
Let $D(n,j)$ count paths from $(0,0)$ to $(n-j,j)$ with steps $(1,0),(0,1),(1,1)$. With $D$ the lower triangular matrix of these counts, write
\[
 d_n(x)=\sum_{j=0}^n D(n,j)x^j,\qquad
 G_n(x)=\sum_{k=0}^n (D^2)(n,k)x^k
       =\sum_{j=0}^n D(n,j)d_j(x).
\]
All coefficients of $G_n$ are positive; its degree is $n$, and it is monic because $D(n,n)=1$.
\begin{theorem}
For every $n\ge1$, $G_n$ has $n$ simple, strictly negative real zeros. For $n=0$, $G_0=1$.
\end{theorem}
This proves the full $D^2$ clause of \cite[Conjecture~4.1]{r30keio:mw}, not the distinct $A^2$ clause for Eulerian numbers. The partial proof in~\cite{r30keio:partial} develops the quadratic divided difference and the central-circle count, and proposes complementary real samples. We give a self-contained account and supply a strict amplitude comparison for those samples.

Splitting paths by their final step gives $d_0=1$, $d_1=1+x$, and $d_j=(1+x)d_{j-1}+xd_{j-2}$ for $j\ge2$. If $z,s$ are distinct roots of $u^2-(1+x)u-x=0$, this yields $d_j(x)=(z^{j+1}-s^{j+1})/(z-s)$. Set $H(z)=zd_n(z)$. Then
\begin{equation}
 G_n(x)=\frac{H(z)-H(s)}{z-s},\qquad
 (z+1)(s+1)=2,\quad x=z+s-1.\label{r30keio:dd}
\end{equation}
Conversely, given $z\ne-1$ and $s=2/(z+1)-1$, these variables satisfy the required quadratic.

\subsection{Two disjoint root counts}
Put $L=-3-2\sqrt2$, $R=-3+2\sqrt2$ and $a=-1-\sqrt2$. For $z<0$ let $v(z)=(1+z)/(2\sqrt{-z})$. It is strictly increasing and maps $[L,R]$ onto $[-1,1]$. By the path recurrence,
\begin{equation}
 d_n(z)=(-z)^{n/2} U_n(v(z)),\qquad
 U_n(\cos\theta)=\frac{\sin((n+1)\theta)}{\sin\theta}\quad(0<\theta<\pi),\label{r30keio:cheb}
\end{equation}
where $U_n$ is the Chebyshev polynomial of the second kind; the identity follows also by its recurrence $U_n(y)=2yU_{n-1}(y)-U_{n-2}(y)$ with $U_0=1,U_1=2y$. In particular all $n$ zeros of $d_n$ are simple and belong to $(L,R)$. None equals $a$: since $d_n$ has rational coefficients the positive conjugate $-1+\sqrt2$ would also be a zero. Denote by $m$ the number of its zeros less than $a$.

\paragraph{The central interval.}
On the upper half of the circle $z=-1+\sqrt2 e^{\ii\phi}$, $0<\phi<\pi$, the partner $s$ equals $\overline z$. In~\eqref{r30keio:dd}, $x=-3+2\sqrt2\cos\phi$ decreases strictly through $(L,R)$ and $G_n(x)=\operatorname{Im}H(z)/\operatorname{Im}z$. The argument of $H$ on this arc, continuously chosen to start at $0$, ends at $N\pi$, where $N=n+1-m$. Indeed all $n+1$ factors $z-r$ in the monic polynomial $H$ start positive at $z=-1+\sqrt2$. For $0<\phi<\pi$ their principal arguments lie in $(0,\pi)$. At $z=a$ their limiting arguments are $\pi$ for a zero $r>a$ (including $r=0$), and $0$ otherwise. Their sum gives the claimed argument. Thus for $j=0,\ldots,N-1$ we may select points in increasing order with arguments $(j+1/2)\pi$: iterated application of continuity suffices, without assuming monotonicity of the phase. At consecutive points the imaginary parts alternate sign. This gives at least $N-1=n-m$ distinct zeros of odd multiplicity for $G_n$ in $(L,R)$.

\paragraph{The exterior interval.}
Suppose $m\ge1$. Take $L\le z<a$, $s=2/(z+1)-1$. Then $L\le z<a<s<-1$, so
\[
 |z|>|s|,\qquad -1\le v(z)<v(s)<0.
\]
Moreover $x=z+s-1$ is strictly increasing in $z$ since its derivative is $1-2/(z+1)^2>0$. As $z$ increases to $a$, $x$ increases to $2a-1=L$. Thus all the samples below give $x<L$.

List in increasing order the $m$ zeros of $d_n$ in $(L,a)$. Use $L$ as the first sample; between each consecutive pair take the point at which $\theta(z)=\arccos v(z)$ equals the arithmetic mean of the two root angles. This gives $m$ samples. For any interior sample the root angles are consecutive multiples of $\pi/(n+1)$, so $|\sin((n+1)\theta(z))|=1$. Since $-1<v(z)<v(s)<0$, we have
$0<\sin\theta(z)<\sin\theta(s)$. It follows that
\begin{equation}
 |H(z)|=\frac{|z|^{1+n/2}}{\sin\theta(z)}
       >\frac{|s|^{1+n/2}}{\sin\theta(s)}\ge |H(s)|.\label{r30keio:key}
\end{equation}
At the endpoint $z=L$ the same strict domination follows from $|U_n(-1)|=n+1$, $|U_n(v(s))|\le n+1$ and $|z|>|s|$. The latter polynomial bound follows, for instance, by expanding $\sin((n+1)\theta)/\sin\theta$ as a sum of $n+1$ unit complex numbers.

Exactly one simple zero of $H$ separates successive samples, giving alternating signs of $H(z)$. By~\eqref{r30keio:key} these are the signs of $H(z)-H(s)$, and $z-s<0$ throughout. Thus $G_n$ changes sign between successive $x$-samples and has at least $m-1$ distinct zeros of odd multiplicity in $(-\infty,L)$.

\paragraph{Degree closure.}
If $m=0$, the central count supplies $n$ distinct negative zeros. Otherwise the disjoint counts give at least $n-1$ distinct negative zeros. After removing their linear factors, the real residual polynomial has degree at most one; it has no nonnegative root since $G_n$ has positive coefficients and nonzero constant term. This alone proves the conjectured real-rootedness without any assumption on multiplicities. In fact the roots selected above have odd multiplicity. None can have multiplicity three or more since that would exceed the total degree; nor can the residual linear factor duplicate a selected root, since then its multiplicity would be even. All roots are therefore simple.\hfill$\square$

\begin{thebibliography}{9}
\bibitem{r30keio:mw} J. Mao and L. Wang, \emph{The Narayana transformation}, arXiv:2607.01572v1 (2026), Section~4, Conjecture~4.1. \url{https://arxiv.org/abs/2607.01572v1}.
\bibitem{r30keio:partial} The Omega Institute, trureturing, \emph{Negative roots of the original Delannoy matrix square}, public partial-proof record, October 2026. \url{https://the-omega-institute.github.io/trureturing-mdbook/Problems/mao-wang-2026-delannoy-square-negative-roots.html}.
\end{thebibliography}

\par\endgroup
\clearpage
\part{Rounds 31--35}

% Source: /workspace/scratch/81914cec2c78/work/archive/Tournament_Round_31_Final/Tournament_Round_31/round31/meiji/paper.tex
% Edited source: /workspace/scratch/81914cec2c78/work/archive/Tournament_Round_31_Final/Tournament_Round_31/round31/meiji/paper.tex
\begingroup

\Needspace{9\baselineskip}
\section{\texorpdfstring{Sharp increments of a reciprocal Hurwitz power}{Sharp increments of a reciprocal Hurwitz power}}\label{section:r31meiji}
\noindent{\small\sffamily\color{muted}Round 31 \enspace / \enspace Meiji}\par\medskip


\begin{sectionabstract}
We derive the four range extensions proposed in Bat\i r's Remark 4.3. The key monotonicity and convexity already follow from Alzer's 2005 theorem. A direct integral argument explains the convexity as positive covariance, and arbitrary positive increments have sharp bounds. Thus the affirmative conclusions are recorded as consequences of prior work, not as a new tournament closure.
\end{sectionabstract}
\subsection{Order and midpoint comparison}
For $p,x>0$ put $V_p(x)=[p\zeta(p+1,x)]^{1/p}$. Alzer \cite{r31meiji:A} proves strict decrease in $p$. The Laurent expansion at $p=0$ and the dominant summand as $p\to\infty$ give
\[
 V_{0+}(x)=e^{-\psi(x)},\qquad V_\infty(x)=1/x .
\]
For the latter limit one may bound $\zeta(p+1,x)$ between $x^{-p-1}$ and $x^{-p-1}+x^{-p}/p$ using the integral test. By strict convexity of $t\mapsto(x+t)^{-p-1}$, its integral over each unit interval strictly exceeds its midpoint value. Summation yields $V_p(x+1/2)<1/x$. Passing to $p\downarrow0$ yields $e^{-\psi(x+1/2)}\le 1/x$. Hence for all $p>0$,
\[
 e^{-\psi(x+1/2)}<V_p(x)<e^{-\psi(x)},
\]
and for $0<p<q$,
\[
 V_p(x+1/2)<1/x<V_q(x)<V_p(x).
\]
These are the proposed extensions of Theorems 3.1 and 3.2 in \cite{r31meiji:B}; the first uses strictness from $1/x<V_p(x)$.
\paragraph{Optimal universal shift.}
A uniform replacement of $1/2$ by $c<1/2$ in the first lower bound would force $\psi(x+c)\ge\log x$ on letting $p\to\infty$. This is impossible since $\psi(x+c)=\log x+(c-1/2)/x+O(x^{-2})$. For the second lower bound let $q\to\infty$ with any fixed $p>1$; the Euler--Maclaurin expansion gives $V_p(x+c)=x^{-1}[1+(1/2-c)/x+O(x^{-2})]$, again ruling out $c<1/2$. This addresses the parameter-uniform reading of Remark 4.4, not the optimal shift for particular exponents.

\subsection{Convexity and sharp increments}
Let $U(x)=\zeta(p+1,x)^{-1/p}$, and $s=p+1$. Since
\[
 U'(x)=\frac{s}{p}\frac{\zeta(s+1,x)}{\zeta(s,x)^{s/p}},
\]
strict convexity follows directly from part (ii) of Alzer's monotonicity theorem \cite{r31meiji:A}: his parameters are $(\alpha,\beta;a,b)=(s+1,s;1,s/p)$ and $\alpha a-\beta b=a-b=-1/p$. Here is an alternative integral proof.
Set $h(t)=t e^{-xt}/(1-e^{-t})$, $L=-\log h$, and
\[
 I_k=\int_0^\infty t^{k-1}h(t)\,dt=\Gamma(k+1)\zeta(k+1,x).
\]
We have $L''(t)=t^{-2}-[4\sinh^2(t/2)]^{-1}>0$. Under the probability density $t^ph(t)/I_{p+1}$, the functions $t$ and $L'(t)$ have strictly positive covariance, since their increments on distinct ordered points have the same strict sign. Integration by parts, with vanishing boundary terms since $p>0$, gives
\[
 E[L']=p I_p/I_{p+1},\qquad E[tL']=p+1.
\]
It follows that $(p+1)I_{p+1}^2>pI_p I_{p+2}$. Substituting the gamma factors gives
\[
 (p+1)^2\zeta(p+2,x)^2>p(p+2)\zeta(p+1,x)\zeta(p+3,x),
\]
which gives $U''>0$ by the displayed derivative formula.
\paragraph{Sharp inequalities.} Since $U(0+)=0$ and $U(x)/x\to p^{1/p}$ by the integral test, for any fixed $d>0$ strict convexity gives
\[
 U(d)< U(x+d)-U(x)<d\,p^{1/p}.
\]
Indeed increments increase strictly with $x$; their limit divided by $d$ equals the asymptotic slope by telescoping at $x=nd$ and taking arithmetic means. Both constants are sharp by the endpoint limits. Setting $d=1$ proves the proposed extension of Theorem 3.3 to all $s>1$.
\paragraph{Inverse form.} Let $u=\zeta(s,y)$, $a=p^{1/p}$ and $b=\zeta(s)^{-1/p}$. Since $\zeta(s,y+1)=u-y^{-s}$, solving the increment inequality for $y$ gives
\[
 [u-(u^{-1/p}+a)^{-p}]^{-1/s}<y<
 [u-(u^{-1/p}+b)^{-p}]^{-1/s}.
\]
All $u>0$ occur because $y\mapsto\zeta(s,y)$ decreases continuously from infinity to zero. The invertible steps preserve sharpness of $a,b$, giving Theorem 3.4 for all $s>1$.
\begin{thebibliography}{9}
\bibitem{r31meiji:B} N. Bat\i r, \emph{New inequalities for the Hurwitz zeta function}, Proc. Indian Acad. Sci. Math. Sci. 118 (2008), 495--503. \href{https://rgmia.org/papers/v12n1/Hurwitz-Piac.pdf}{Author report, Remark 4.3}.
\bibitem{r31meiji:A} H. Alzer, \emph{Monotonicity properties of the Hurwitz zeta function}, Canad. Math. Bull. 48 (2005), 333--339. \href{https://doi.org/10.4153/CMB-2005-031-3}{doi:10.4153/CMB-2005-031-3}.
\end{thebibliography}

\par\endgroup

% Source: /workspace/scratch/81914cec2c78/work/archive/Tournament_Round_31_Final/Tournament_Round_31/round31/showa/paper.tex
% Edited source: /workspace/scratch/81914cec2c78/work/archive/Tournament_Round_31_Final/Tournament_Round_31/round31/showa/paper.tex
\begingroup
\newcommand{\eone}{e_1}
\newcommand{\etwo}{e_2}
\Needspace{9\baselineskip}
\section{\texorpdfstring{Two Sharp Median-to-Side Ratios in Nonobtuse Triangles}{Two Sharp Median-to-Side Ratios in Nonobtuse Triangles}}\label{section:r31showa}
\noindent{\small\sffamily\color{muted}Round 31 \enspace / \enspace Showa}\par\medskip


\begin{sectionabstract}
Among nondegenerate nonobtuse triangles, the ratio of the sum of
medians to the perimeter is minimized precisely by the right
isosceles triangle. The same extremizer minimizes the ratio of the
sums of pairwise products. We prove two conjectures of Jian Liu by
showing that, at fixed largest normalized side square, replacing the
other two side squares by their mean can only decrease either ratio.
The comparison uses positive even-power coefficients to control the
loss of each numerator by that of its denominator.
\end{sectionabstract}

\subsection*{Statements}
Write $a,b,c$ for the sides and $m_a,m_b,m_c$ for the respective
medians. For a triple $v$, let $e_1(v)=v_1+v_2+v_3$ and
$e_2(v)=v_1v_2+v_2v_3+v_3v_1$. Put $s=(a,b,c)$ and
$m=(m_a,m_b,m_c)$. For every nondegenerate nonobtuse triangle,
\begin{align}
 \frac{e_1(m)}{e_1(s)}&\geq \frac{\sqrt2+2\sqrt5}{4+2\sqrt2},
 \label{r31showa:one}\\
 \frac{e_2(m)}{e_2(s)}&\geq \frac{5+2\sqrt{10}}{4+8\sqrt2}.
 \label{r31showa:two}
\end{align}
Equality in either bound holds precisely at the right isosceles shape.
Thus both inequalities are strict in the open acute domain, with
unchangeable constants by approximation. These are Liu's Conjectures
2 and 3, respectively, equations (4.2) and (4.3) in [1, p.~20].
The author's Conjecture 4 involves a different ratio. An earlier
tournament argument on that ratio also smoothed the two shorter side
squares. The numerator and denominator comparisons below are specific
to the first and second elementary symmetric functions.

\subsection*{A common smoothing step}
Scale so $a^2+b^2+c^2=1$, and label the normalized squared sides
$x\ge y$ and $z\ge x$. Nonobtuseness implies $1/3\le z\le1/2$.
Write
\[
 u=\tfrac12(1-z),\quad x=u(1+t),\quad y=u(1-t),\quad
 A=2-3u,\quad B=2-3z,\quad b_0=3u/A .
\]
Here $1/4\le u\le1/3$, $0\le t<1$, and $0<b_0\le1$. The triangle
with side squares $(u,u,z)$ is nondegenerate and nonobtuse, since
$z\le2u$. Apollonius' median formula gives $(2m_a)^2=2-3a^2$
and its cyclic analogues. Set $L(t)=\sqrt{1+t}+\sqrt{1-t}$ and
$E(t)=\sqrt{1-t^2}$.

For $0\le d\le1$ and $0\le t<1$ the binomial expansions give
\begin{equation}
 2-L(dt)\le d^2[2-L(t)],\qquad
 1-E(dt)\le d^2[1-E(t)]. \label{r31showa:loss}
\end{equation}
Indeed each expression on the left is a sum of nonnegative multiples
of $(dt)^{2j}$, $j\ge1$, and $d^{2j}\le d^2$. All series converge
in the range of use. Write $P=A/u$ and $Q=B/z$.
Then $P\ge Q>0$, since $(2-3w)/w$ decreases with $w$ and $u\le z$.
The key comparison is $P^3Q\ge81$, with strict inequality if $u<1/3$.
To see it, multiply by $u^3z$ and factor:
\[
 B A^3-81u^3 z=2(3u-1)(45u^2-30u+4).
\]
The second factor is negative: it is decreasing on $[1/4,1/3]$
and equals $-11/16$ at $u=1/4$.

For $i=1,2$, let $N_i(t)=e_i(2m)$ and $D_i(t)=e_i(s)$ at the shape
$(x,y,z)$. Direct expansion now gives
\begin{align*}
 N_1(t)&=\sqrt A L(b_0t)+\sqrt B,
 &D_1(t)&=\sqrt u L(t)+\sqrt z,\\
 N_2(t)&=A E(b_0 t)+\sqrt{AB}L(b_0t),
 &D_2(t)&=uE(t)+\sqrt{uz}L(t).
\end{align*}
For the first ratio, \eqref{r31showa:loss} yields
\[
 N_1(0)-N_1(t)\le M_1[D_1(0)-D_1(t)],\qquad
 M_1=\sqrt P b_0^2=9/P^{3/2}\le\sqrt Q.
\]
The last inequality is exactly $P^3Q\ge81$. Moreover,
$N_1(0)/D_1(0)$ is a weighted average of $\sqrt P$ and $\sqrt Q$,
so it is at least $M_1$.

For the second ratio apply both estimates in \eqref{r31showa:loss}. In the
$E$ and $L$ terms the coefficients of denominator loss are bounded
respectively by $Pb_0^2$ and $\sqrt{PQ}b_0^2$, both at most
$M_2=Pb_0^2=9/P$. Since $P^3Q\ge81$, this is at most $\sqrt{PQ}$.
In turn $N_2(0)/D_2(0)$ is a weighted average of $P$ and $\sqrt{PQ}$.
Consequently, for either $i$,
\[
 N_i(0)-N_i(t)\le \frac{N_i(0)}{D_i(0)}
                     \bigl(D_i(0)-D_i(t)\bigr),
 \qquad \frac{N_i(t)}{D_i(t)}\ge \frac{N_i(0)}{D_i(0)}.
\]
If $t>0$, then $u<1/3$ (otherwise all side squares would equal $1/3$),
the comparisons using $P^3Q$ are strict, and each denominator loss is
positive. Thus equality in smoothing requires $t=0$.

\subsection*{The isosceles interval}
The remaining sides may be scaled to $(1,1,r)$, $1\le r\le\sqrt2$.
The first scaled ratio is
\[
 h(r)=\frac{2\sqrt{2r^2+1}+\sqrt{4-r^2}}{2+r}.
\]
The sign of its derivative is the sign of
$(8r-2)/\sqrt{2r^2+1}-(2r+4)/\sqrt{4-r^2}$. Both subtracted
quantities are positive, and their squared difference has numerator
\[
 (8r-2)^2(4-r^2)-(2r+4)^2(2r^2+1)
       =-72r(r-1)^2(r+2).
\]
Thus $h$ decreases strictly from $1$ to $\sqrt2$. For the second
scaled ratio put $W=(2r^2+1)(4-r^2)$ and
\[
 H(r)=\frac{2r^2+1+2\sqrt W}{1+2r}.
\]
The sign of $H'(r)$ is the sign of $X\sqrt W+Y$, where
\[
 X=4r^2+4r-2,\qquad Y=(1+2r)W'-4W=-2(4r^4+4r^3-7r+8).
\]
Here $X>0$ and $Y<0$ since $4r^4+4r^3\ge8r$.
Direct multiplication yields
\begin{align*}
 X^2W-Y^2=-24(r-1)^2
 (4r^6+16r^5+26r^4+16r^3+5r^2+4r+10).
\end{align*}
Hence $H$ also decreases strictly on the interval.
Substituting $r=\sqrt2$ gives $h(\sqrt2)/2$ and $H(\sqrt2)/4$
equal to the right sides of \eqref{r31showa:one} and \eqref{r31showa:two}.
The smoothing equality and this strict decrease prove the claimed
uniqueness. Taking $r\uparrow\sqrt2$ gives sharpness even for acute
triangles.\hfill$\square$

\begin{thebibliography}{9}
\bibitem{r31showa:liu} J. Liu, Two inequalities involving circumradius, inradius
and medians of an acute triangle, \emph{Advances in Inequalities and
Applications} (2023), article 1, p.~20.
\href{https://doi.org/10.28919/aia/8232}{doi:10.28919/aia/8232}.
\end{thebibliography}

\par\endgroup

% Source: /workspace/scratch/81914cec2c78/work/archive/Tournament_Round_31_Final/Tournament_Round_31/round31/heisei/paper.tex
% Edited source: /workspace/scratch/81914cec2c78/work/audit_26_33/cleaned_sources/r31_heisei.tex
\begingroup
\newcommand{\RR}{\mathbb R}
\newcommand{\sgn}{\operatorname{sgn}}
\Needspace{9\baselineskip}
\section{\texorpdfstring{The extremizing phase in the Kalaj--Elkies higher-derivative integral}{The extremizing phase in the Kalaj--Elkies higher-derivative integral}}\label{section:r31heisei}
\noindent{\small\sffamily\color{muted}Round 31 \enspace / \enspace Heisei}\par\medskip


\begin{sectionabstract}
Folding an angular integral onto a circle converts the phase dependence in a sharp real-part estimate into a convolution. Its weight is a Chebyshev-root sum whose derivative has the sign of a single falling factorial. We prove the endpoint-maximum conjecture of Kalaj and Elkies, including the strict decrease conjectured for every finite exponent above two. We give the exact parity law below two as part of the same all-parameter classification. A separate appendix gives a signed-kernel classification for the Arora--Sahoo gamma quotient, including the order-dependent zero hierarchy outside its region of logarithmic complete monotonicity.
\end{sectionabstract}
\subsection{The angular maximization}
Use the angle and normalization of (1.7) of Kalaj--Elkies \cite{r31heisei:EK}. For an integer $n\ge1$ and $1\le q<\infty$, let
\[
 F_q(\beta)=\int_0^\pi \sin^{(n+1)q-2}v\,
 \left|\cos\!\left((n+1)v+\beta-\tfrac{n-1}{2}\pi\right)\right|^q\,dv.
\]
In Remark 1.3 the authors ask whether the maximum on $[0,\pi/2]$ always occurs at one of its ends. They discuss $q=1$ separately and record known results for $n=1$ and for the Hilbert exponent $q=2$. Conjecture 4.1 of the same preprint predicts strict decrease for $q>2$ and a parity-dependent direction for lower exponents. We use preprint numbering throughout. Known results for a related upper-half-plane optimization include special exponent families \cite{r31heisei:K}.
\begin{theorem}\label{r31heisei:main}
Set $N=n+1$, $s=Nq/2-1$, and $A_N(s)=\prod_{k=0}^{N-1}(s-k)$. The sign of $-F_q'(\beta)$ for $0<\beta<\pi/2$ equals the sign of $A_N(s)$. In particular $F_q$ is strictly decreasing, strictly increasing, or constant on $[0,\pi/2]$ according as $A_N(s)$ is positive, negative, or zero.
\end{theorem}
Consequently the endpoint-maximum conjecture holds, and $\beta=0$ is the unique maximizer in this interval whenever $q>2$. For $1\le q\le2$ a zero factor occurs precisely when $Nq/2$ is integral. Otherwise the sign of $A_N(s)$ is $(-1)^{N-\lfloor Nq/2\rfloor}$. This parity formula is the exact direction supplied by the theorem.


At $q=\infty$ the original remark lists the endpoint optimization as known; equation (1.7) is an ordinary integral only for finite $q$. If one instead defines a phase functional by the natural limiting norm, the endpoint result holds directly: it is
$\max_v \sin^N(v) |\cos(Nv+\beta-(n-1)\pi/2)|$, bounded by one, and is one at $v=\pi/2$, $\beta=0$. This follows also from uniform convergence of $\sin^{N-2/q}(v)$ to $\sin^N(v)$ and the $L^q$-to-uniform norm limit; no identification with a differently scaled endpoint constant is needed.

\subsection{The folding and sign calculation}
\begin{proof}[Proof of Theorem~\ref{r31heisei:main}]
Write $v=w+\pi/2$, so that the angular cosine becomes $-\cos(Nw+\beta)$. The variable $w$ now runs through a period of length $\pi$ for the product $|\cos w|^{2s}|\cos(Nw+\beta)|^q$. For exponent zero interpret the first factor as the constant one. On $\mathbb T=\RR/\pi\mathbb Z$ let
\[
 K(u)=|\cos u|^q,\qquad W(u)=\sum_{k=0}^{N-1}
 \left|\cos\frac{u+k\pi}{N}\right|^{2s}.
\]
Substitute $u=Nw$ and split the resulting $N$ periods of the kernel. Since $W$ is even (permutation of the indices) and $\pi$-periodic, this gives
\[
 F_q(\beta)=\frac1N\int_{\mathbb T}K(\beta-u)W(u)\,du.
\]
Here $du$ is ordinary arc length on one period, not probability measure.

For $0<u<\pi/2$ set $y=\cos(2u)$ and let $T_N$ denote the Chebyshev polynomial defined by $T_N(\cos v)=\cos(Nv)$. The numbers
\[
 z_k=\cos\frac{2u+2k\pi}{N},\qquad 0\le k<N,
\]
are the $N$ distinct roots of $T_N(z)=y$ in $(-1,1)$: equality of two such cosines forces equality of indices modulo $N$ or $2u\in\pi\mathbb Z$. Set $f(z)=((1+z)/2)^s$. Then $W(u)=\sum_k f(z_k)$, and $dz_k/dy=1/T_N'(z_k)$. The leading coefficient of $T_N$ is $2^{N-1}$, whence
\[
 \frac{dW}{dy}=2^{1-N}\sum_{k=0}^{N-1}
 \frac{f'(z_k)}{\prod_{\ell\ne k}(z_k-z_\ell)}
 = \frac{2^{1-N}}{(N-1)!} f^{(N)}(\xi)
\]
for some $\xi$ between the least and greatest nodes. Indeed the sum is the leading coefficient of the Lagrange interpolant to $f'$ at the nodes; repeated Rolle's theorem applied to the interpolation error proves the last equality. Only interior smoothness is needed, even for nonintegral $s$. Since
\[
 f^{(N)}(z)=2^{-s}A_N(s)(1+z)^{s-N},
\]
$W$ strictly decreases (respectively increases) with $u\in(0,\pi/2)$ if $A_N(s)>0$ (respectively $<0$), and is constant when $A_N(s)=0$. Continuity extends these conclusions to the endpoints. This weight is, up to rotation and scale, a regular-polygon powered-distance sum. Indeed $2|\cos x|=|e^{2ix}+1|$. Extremal results for such sums predate the conjecture; see Nikolov and Rafailov \cite{r31heisei:NR}. We include the derivative computation so that we need no separate polygon theorem.

Let $|u|_{\mathbb T}$ be the distance to $\pi\mathbb Z$. For $0<r<\pi/2$, $K'(r)=-q\cos^{q-1}r\sin r<0$, and $K(\pi/2)=0$; hence $K(u)$ is the integral of $-K'(r)$ over $|u|_{\mathbb T}<r<\pi/2$. For fixed $0<r<\pi/2$, the sliding interval average of $W$ (without normalization) has derivative
\[
 \frac{d}{d\beta}\int_{\beta-r}^{\beta+r}W(u)\,du
 =W(\beta+r)-W(\beta-r).
\]
For $0<\beta<\pi/2$ we have $|\beta+r|_{\mathbb T}>|\beta-r|$:
both $\beta+r$ and $\pi-\beta-r$ exceed $|\beta-r|$. Thus this derivative has the strict sign of $W$'s direction in every nonconstant case. Integrating it against $-K'(r)\,dr/N$ preserves that sign. Interchange is valid by the boundedness of $W$ and integrability of $-K'$. This proves the theorem.
\end{proof}
The substantive passage is from the original phase integral to a monotone convolution on the folded circle. It is not a claim of priority for divided differences, polygon extremizers or convolution monotonicity by themselves.

\subsection*{Appendix: an independent gamma-quotient problem}
Let $x=c-4$ and define
\[
 G_a(x)=\frac{\Gamma(x)\Gamma(x+4)}
 {\Gamma(x+4-a)\Gamma(x+3+a)},\qquad
 D_j(x)=(-1)^{j+1}(\log G_a)^{(j+1)}(x)\quad(j\ge0).
\]
We consider only the domain on which all four individual gamma arguments are positive, namely $x>d_0=\max(0,a-4,-3-a)$. Arora--Sahoo \cite{r31heisei:AS}, Conjecture 3.14, asks for strict decrease of $G_a$ for $0\le a\le1$, $x>0$. The positive-kernel region below is already implied by Alzer's general weak-majorization criterion \cite{r31heisei:Alzer}: the numerator shifts have minimum $0$ and sum $4$, whereas the denominator shifts then have nonnegative minimum and sum $7$. The threshold analysis in the complementary region is a separate extension, not part of their conjecture.
\begin{theorem}\label{r31heisei:gamma}
For $-3\le a\le4$, all $D_j$ are strictly positive on $x>0$. If $a\notin[-3,4]$, put $d=d_0>0$. There exists $\tau>0$ such that for every $j\ge0$, $D_j$ has exactly one zero $r_j$ on $(d,\infty)$, is negative below it and positive above it, and
\[
 d<r_0<r_1<\cdots,\qquad \lim_{j\to\infty}\frac{r_j}{j}=\frac1\tau.
\]
\end{theorem}
\begin{proof}
The integral for differences of digamma functions and differentiation under the integral give
\[
 D_j(x)=\int_0^\infty t^j e^{-xt}h(t)\,dt,\qquad
 h(t)=\frac{1+e^{-4t}-e^{-(4-a)t}-e^{-(3+a)t}}{1-e^{-t}}.
\]
Here $h(0)=3$ by continuity; exponential decay of the integrand is uniform on compact subintervals of $x>d_0$. On $[-3,4]$ the sum of the last two exponentials is convex as a function of $a$, with the same value $1+e^{-7t}$ at both endpoints. Thus $h(t)\ge(e^{-4t}-e^{-7t})/(1-e^{-t})>0$ for every $t>0$. Taking $j=0$ gives the original conjecture.

Outside this interval the denominator shifts are $-d,7+d$ in some order. Then the numerator of $h(t)$ equals
\[
 e^{-4t}(1-e^{-(3+d)t})-(e^{dt}-1).
\]
The ratio of the subtracted to the first quantity is
\[
 e^{4t}\,\frac{\int_0^d e^{st}\,ds}{\int_0^{3+d}e^{-st}\,ds}.
\]
It increases strictly from $d/(3+d)<1$ to $\infty$. Its unique crossing of one defines $\tau$, and $h$ changes sign from positive to negative there. Consequently
\[
 \frac{d}{dx}(e^{\tau x}D_j(x))
 =e^{\tau x}\int_0^\infty (\tau-t)t^je^{-xt}h(t)\,dt>0.
\]
As $x\downarrow d$, the negative tail diverges, since $e^{-dt}h(t)\to-1$ and the positive part on $[0,\tau]$ stays bounded. As $x\to\infty$, $x^{j+1}D_j(x)\to 3j!$ by $t=u/x$, the bound $|h(t)|\le C(1+e^{dt})$, and dominated convergence. This gives exactly one sign change at $r_j$. Further,
\[
 D_{j+1}(r_j)=\int_0^\infty(t-\tau)t^je^{-r_jt}h(t)\,dt<0,
\]
which proves the strict ordering.

Fix $v>0$. For $j\ge1$ let $T_j$ have gamma density with shape $j+1$ and rate $jv$. Then $D_j(jv)$ has the sign of $\mathbb E h(T_j)$ for large $j$. Its mean and variance are $(j+1)/(jv)$ and $(j+1)/(jv)^2$, respectively, so $T_j\to 1/v$ in probability. Moreover $\mathbb E e^{2dT_j}=(1-2d/(jv))^{-(j+1)}$ stays bounded for all sufficiently large $j$. The bound on $h$ gives uniform integrability, thus $\mathbb E h(T_j)\to h(1/v)$. For $v<1/\tau$ this limit is negative, giving $jv<r_j$ eventually, and for $v>1/\tau$ it is positive, giving $r_j<jv$ eventually. Squeezing from both sides proves the limit.
\end{proof}

\begin{thebibliography}{9}
\bibitem{r31heisei:EK} N. D. Elkies and D. Kalaj, A real part theorem for the higher derivatives of analytic functions in the unit disk, \emph{Comput. Methods Funct. Theory} \textbf{13} (2013), 189--203, \href{https://doi.org/10.1007/s40315-013-0016-0}{doi:10.1007/s40315-013-0016-0}. Numbering and statements used here: \href{https://arxiv.org/pdf/1202.2520}{arXiv:1202.2520v5}, (1.7), Remark 1.3 and Conjecture 4.1.
\bibitem{r31heisei:NR} N. Nikolov and R. Rafailov, On the sum of powered distances to certain sets of points on the circle, \emph{Pacific J. Math.} \textbf{253} (2011), 157--168, \href{https://doi.org/10.2140/pjm.2011.253.157}{doi:10.2140/pjm.2011.253.157}.
\bibitem{r31heisei:K} G. Kresin, Explicit real-part estimates for high order derivatives of analytic functions, \href{https://arxiv.org/abs/1509.01176}{arXiv:1509.01176} (2015).
\bibitem{r31heisei:AS} H. M. Arora and S. K. Sahoo, Interpolation on Gauss hypergeometric functions with an application, \emph{Involve} \textbf{11} (2018), 625--641, \href{https://doi.org/10.2140/involve.2018.11.625}{doi:10.2140/involve.2018.11.625}, Conjecture 3.14.
\bibitem{r31heisei:Alzer} H. Alzer, On some inequalities for the gamma and psi functions, \emph{Math. Comp.} \textbf{66} (1997), 373--389, Theorem 10, \href{https://doi.org/10.1090/S0025-5718-97-00807-7}{doi:10.1090/S0025-5718-97-00807-7}.
\end{thebibliography}

\par\endgroup

% Source: /workspace/scratch/81914cec2c78/work/archive/Tournament_Round_31_Final/Tournament_Round_31/round31/reiwa/paper.tex
% Edited source: /workspace/scratch/81914cec2c78/work/archive/Tournament_Round_31_Final/Tournament_Round_31/round31/reiwa/paper.tex
\begingroup

\Needspace{9\baselineskip}
\section{\texorpdfstring{Which rational degree roots occur at a fixed graph order?}{Which rational degree roots occur at a fixed graph order?}}\label{section:r31reiwa}
\noindent{\small\sffamily\color{muted}Round 31 \enspace / \enspace Reiwa}\par\medskip


\begin{sectionabstract}
We determine the rational numbers that are roots of degree polynomials of
graphs on exactly $n$ vertices. For odd $n$, a reduced negative fraction
$-p/q$ occurs precisely when $p+q$ divides $n$; for even $n$, it occurs
precisely when $p+q$ divides $n/2$. The classification is the same for
loopless multigraphs. The parity obstruction comes from differentiating
the primitive linear factor at $1$ and using the handshaking identity.
Matchings and paths suffice to realize every permitted negative root.
\end{sectionabstract}

\subsection*{The fixed-order question}
For a nonempty finite undirected graph $G$, its degree polynomial is
$D_G(x)=\sum_{v\in V(G)}x^{\deg v}$. A degree root is a zero of this
polynomial. George \cite[Sec.~4.1, pp.~96--97]{r31reiwa:george} asks:
\textit{for graphs of some order $n$ what rational degree roots may we see?}
The question is not about the union over all orders, which already
contains every negative rational. Brown and George \cite[Sec.~4, p.~181]{r31reiwa:bg}
also identify the rational-root question for graphs and multigraphs; it is
separate from their Conjecture~4.1 on the region containing complex roots.

\begin{theorem}
Fix an integer $n\ge1$ and put $N=n/\gcd(n,2)$. For relatively prime
positive integers $p,q$, the following are equivalent:
\begin{enumerate}
\item a simple graph on $n$ vertices has degree root $-p/q$;
\item a loopless multigraph on $n$ vertices has degree root $-p/q$;
\item $p+q$ divides $N$.
\end{enumerate}
No positive rational occurs. Zero occurs if and only if $n\ge2$.
\end{theorem}
Parallel edges are allowed in a multigraph and counted with
multiplicity in the degrees. Graphs may be disconnected, as in the source.

\begin{proof}
A simple graph is a loopless multigraph. For necessity let $G$ be such a
multigraph and suppose $D_G(-p/q)=0$. Write $s=p+q$. By the factor theorem
over $\mathbb Q$ and Gauss's lemma for the primitive polynomial $p+qx$,
there is an $H\in\mathbb Z[x]$ with $D_G(x)=(p+qx)H(x)$.
Evaluation at $1$ gives $n=sH(1)$, so $m=n/s$ is an integer. Further,
\[
 2|E(G)|=D_G'(1)=qm+sH'(1).
\]
If $s$ is even then $q$ is odd, so $m$ must be even. When $n$ is even
and $s$ is odd, $m$ is even automatically. Hence $s\mid N$ for even
$n$; for odd $n$ we already have $s\mid n=N$.

Conversely, assume $s\mid N$, and let $m=n/s$. We seek a simple graph
with $mp$ vertices of degree $d$ and $mq$ of degree $d+1$. If $mq$ is
even, take a matching on $mq$ vertices and add $mp$ isolated vertices;
this works with $d=0$. If $mq$ is odd, then $m$ and $q$ are odd.
The hypothesis $s\mid N$ excludes even $n$ in this case, so $s$ is odd
and $p$ even. In particular $mp\ge2$ is even. Take one path with $mq$
internal vertices and two end vertices, and add $(mp-2)/2$ disjoint
edges. Exactly $mq$ vertices have degree two and $mp$ have degree one.
In either case $D_G(x)=m x^d(p+qx)$ has the required root.

All coefficients are nonnegative and $D_G$ is not zero, excluding positive
roots. For $n\ge2$ the path on $n$ vertices has no isolates, so $D_G(0)=0$.
At order one the only loopless graph has polynomial $1$.
\end{proof}

\paragraph{A way to list the roots.}
The nonzero rational roots at order $n$ are precisely
$-(N-j)/j$ for $1\le j<N$. Indeed, if $p+q\mid N$ take
$j=Nq/(p+q)$; conversely, reducing such a ratio yields a numerator
and denominator whose sum divides $N$. For example at $n=12$, $N=6$
and the list is $-5,-2,-1,-1/2,-1/5$. This illustrates why an equality
$p+q=n/2$ for reduced fractions, suggested by early even-order data, is too
restrictive: the relevant relation is divisibility. George already proved
the odd-order realizations in his thesis; the fixed-order necessity and
the even-order criterion above supply the classification.

\begin{thebibliography}{9}
\bibitem{r31reiwa:george} I. George, \emph{On the Degree Polynomial of Graphs},
M.Sc. thesis, Dalhousie University, 2023, Sec.~4.1, pp.~96--97.
\href{https://dalspaceb.library.dal.ca/server/api/core/bitstreams/b38bd694-b1b7-4993-a49f-b0d2dc04cc0b/content}{University archive}.
\bibitem{r31reiwa:bg} J. I. Brown and I. C. George, On the roots of degree
polynomials, \emph{Australas. J. Combin.} \textbf{92} (2025), 164--183.
\href{https://ajc.maths.uq.edu.au/pdf/92/ajc_v92_p164.pdf}{Published article}.
\end{thebibliography}

\par\endgroup

% Source: /workspace/scratch/81914cec2c78/work/archive/Tournament_Round_31_Final/Tournament_Round_31/round31/keio/paper.tex
% Edited source: /workspace/scratch/81914cec2c78/work/archive/Tournament_Round_31_Final/Tournament_Round_31/round31/keio/paper.tex
\begingroup

\Needspace{9\baselineskip}
\section{\texorpdfstring{A three-coordinate variation for power sums on a moment sphere}{A three-coordinate variation for power sums on a moment sphere}}\label{section:r31keio}
\noindent{\small\sffamily\color{muted}Round 31 \enspace / \enspace Keio}\par\medskip


\begin{sectionabstract}
Fix the first and second power sums of $n$ real numbers. We give a direct proof that a maximum of any integer power sum of degree at least three is attained with at least $n-1$ coordinates equal. The argument uses complete homogeneous polynomials on a varying triple, followed by a second-order test at repeated coordinates. We distinguish this alternative proof of C\^irtoaje's stated conjecture from a claim of historical priority: sharp standardized sample-moment bounds have a much older literature.
\end{sectionabstract}
\subsection{Eliminating triples of distinct coordinates}
\begin{theorem}
Let $n\ge3$ and $m\ge3$ be integers and suppose $\sum_i x_i=S$ and $\sum_i x_i^2=Q$ are feasible over the reals. The maximum of $\sum_i x_i^m$ is attained at a point with at least $n-1$ equal coordinates.
\end{theorem}
This is the assertion posed as Conjecture~2.1 in \cite{r31keio:c}. We provide a proof, not a claim that no earlier result implies it. The feasible set is compact. If $nQ=S^2$ it is a singleton, so assume a nonconstant configuration from now on.

For a triple $a,b,c$ of distinct coordinates, fix $p=a+b+c$ and $q=ab+bc+ca$ while varying $r=abc$. The roots of $P_r(z)=z^3-pz^2+qz-r$ remain real and distinct for all $r$ in an open neighborhood of their present value, since they are simple. This is an admissible variation of the original constraints. Write $s_m=a^m+b^m+c^m$ and let $h_d(u_1,\ldots,u_l)$ denote the coefficient of $t^d$ in $\prod_i(1-u_i t)^{-1}$, with $h_d=0$ if $d<0$. The identity
\[
 \sum_{m\ge1}\frac{s_m}{m}t^m=-\log(1-pt+qt^2-rt^3)
\]
implies
\begin{equation}
 s_m'=m h_{m-3}(a,b,c),\qquad
 s_m''=m h_{m-6}(a,a,b,b,c,c),\label{r31keio:triple}
\end{equation}
where primes are derivatives in $r$. If $d$ is even and the real tuple $u$ is nonzero, then $h_d(u)>0$: for independent unit exponential variables $E_i$ the coefficient is $\mathbb E[(\sum_i u_i E_i)^d]/d!$ by the moment generating function. For $d=0$ it is one, and otherwise the linear form is nonzero with positive probability.

Thus for odd $m$ the first derivative in~\eqref{r31keio:triple} is positive. For even $m\ge6$ the second derivative is positive. Neither case permits a maximum at a triple of three distinct values. For $m=4$, choose among the maximizers a point maximizing $\sum_i x_i^3$ as a tiebreaker. On a triple, $s_4'=4p$; if this is nonzero it excludes a maximum, and otherwise $s_4$ is constant while $s_3'=3$ contradicts the tiebreaker. Consequently in all cases there exists a maximizing point with just two distinct coordinate values $a<b$.

\subsection{Ruling out two repeated clusters}
Write $f(x)=x^m$. At a nonconstant feasible point the gradients of the constraints are independent. The Lagrange multiplier equations at an extremum are $f'(x_i)=\lambda+\eta x_i$, where
\[
 \eta=\frac{f'(b)-f'(a)}{b-a}=\frac{1}{b-a}\int_a^b f''(u)\,du.
\]
If $a$ occurs at least twice, the vector with entries $1,-1$ on a pair of equal $a$ coordinates and zero elsewhere is tangent to both constraints. The second-order necessary condition gives $2(f''(a)-\eta)\le0$. The same applies to $b$ if it is repeated. But $f''$ is strictly increasing for odd $m$ and strictly convex for even $m\ge4$. In either case its integral average on $[a,b]$ is strictly less than the larger endpoint value. Thus the two endpoint tests cannot both hold. At least one cluster is a singleton, proving the theorem.\hfill$\square$

For example, with $\mu=S/n$ and $\sigma^2=Q/n-\mu^2$, only the values of $\sum_i x_i^m$ at two explicit vectors need be compared:
\[
 \left(\mu\pm\sigma\sqrt{n-1},\quad
       \underbrace{\mu\mp\sigma/\sqrt{n-1},\ldots,\mu\mp\sigma/\sqrt{n-1}}_{n-1}\right).
\]

\paragraph{Scope and prior results.}
Dal\'en \cite{r31keio:d} reported sharp standardized sample-moment bounds of arbitrary order long before \cite{r31keio:c}. The connection is important: for an even raw power, reflection allows $\mu\ge0$, and the binomial expansion in centered powers then has nonnegative coefficients. Whenever the sharp centered upper bounds for all orders are realized by the same positive-outlier sample, their sum gives the claimed even raw bound. For odd powers the strictly convex derivative already fits the equal-variables principle. We do not count the theorem as a demonstrated new resolution; the proof here does not need any centered-moment estimate.

\paragraph{A separate sufficient condition.}
The variation also proves the same maximizer statement for $f\in C^6(\mathbb R)$ with $f^{(6)}>0$ and $f''$ strictly convex. Indeed along a distinct-root triple, $a'=1/P_r'(a)$ and similarly for $b,c$. Thus the first derivative of $\sum f(a)$ (sum over the triple) is $\sum f'(a)/P_r'(a)$. Its derivative at a fixed $r$ depends only on $f'$ and $f''$ at the three roots. Replace $f'$ by its degree-at-most-five Hermite interpolant $B$ matching its values and first derivatives there. The expression being differentiated now equals the coefficient of $z^{-1}$ at infinity in $B(z)/P_r(z)$ by partial fractions (use the polynomial part plus the simple-pole terms). Its derivative in $r$ is the $z^{-1}$ coefficient of $B(z)/P_r(z)^2$, namely the $z^5$ coefficient of $B$. By repeated Rolle applied to $f'-B$ with three double zeros, this coefficient is $f^{(6)}(\xi)/5!>0$ for some $\xi$ between the extreme roots. The triple objective is strictly convex in $r$ and the preceding repeated-cluster test finishes the proof. This criterion is not the statement of Conjecture~2.1.

\begin{thebibliography}{9}
\bibitem{r31keio:c} V. C\^irtoaje, On the equal variables method applied to real variables, \emph{Creative Math. Inform.} 24 (2015), 145--154, Conjecture~2.1.
\bibitem{r31keio:d} J. Dal\'en, Algebraic bounds on standardized sample moments, \emph{Statist. Probab. Lett.} 5 (1987), 329--331. \url{https://doi.org/10.1016/0167-7152(87)90005-8}.
\end{thebibliography}

\par\endgroup

% Source: /workspace/scratch/81914cec2c78/work/archive/Tournament_Round_32_Final/Tournament_Round_32/round32/taisho/paper.tex
% Edited source: /workspace/scratch/81914cec2c78/work/audit_26_33/cleaned_sources/r32_taisho.tex
\begingroup
\newcommand{\cyc}{\mathop{\sum}_{\rm cyc}}
\Needspace{9\baselineskip}
\section{\texorpdfstring{A two-cone proof of Liu's cubic-denominator median inequality}{A two-cone proof of Liu's cubic-denominator median inequality}}\label{section:r32taisho}
\noindent{\small\sffamily\color{muted}Round 32 \enspace / \enspace Taisho}\par\medskip

\textbf{Abstract.} We prove Conjecture 4.3 of Jian Liu (2021). In median coordinates, the acute-triangle condition has a quadratic boundary, whereas the proposed inequality has no radicals. Replacing the curved boundary by two rational linear constraints gives two cones, on each of which a single degree-nine polynomial has nonnegative coefficients. The proof also includes non-obtuse triangles and determines the equality case.

\subsection*{Statement and reduction}
Let $a,b,c$ be the sides of a non-obtuse triangle and $m_a,m_b,m_c$ the corresponding median lengths. Then
\begin{equation}
 \frac{a^2}{(m_b+m_c)^3}+\frac{b^2}{(m_c+m_a)^3}
 +\frac{c^2}{(m_a+m_b)^3}
 \ \le\ \frac3{2(m_a+m_b+m_c)}, \label{r32taisho:main}
\end{equation}
with equality exactly for equilateral triangles. For acute triangles this is precisely \cite[Conjecture 4.3, (66)]{r32taisho:liu}. It strengthens the original paper's Theorem 1.1: multiplying the right side of \eqref{r32taisho:main} by $\sum_{\rm cyc}(m_b+m_c)$ and using Cauchy--Schwarz recovers $\sum_{\rm cyc}a/(m_b+m_c)\le\sqrt3$.

Relabel so that $u=m_a\ge v=m_b\ge w=m_c$. Apollonius's formula, applied to the three sides, gives
\[
 a^2=\tfrac49(2v^2+2w^2-u^2),\qquad
 b^2=\tfrac49(2w^2+2u^2-v^2),\qquad
 c^2=\tfrac49(2u^2+2v^2-w^2).
\]
Thus $c$ is the largest side, and $a^2+b^2-c^2=\tfrac49(5w^2-u^2-v^2)\ge0$.
In particular $u\le2w$ since $v\ge w$, while Cauchy--Schwarz gives $u+v\le\sqrt{10}\,w<16w/5$. If $h=u-v$ and $k=v-w$, the two useful consequences are
\begin{equation}
 w\ge h+k,\qquad w\ge\tfrac56(h+2k). \label{r32taisho:cone}
\end{equation}
The two linear bounds agree at $h=4k$.

Let $S=u+v+w$ and $D=(u+v)(v+w)(w+u)$. Multiplication of the right side minus the left side in \eqref{r32taisho:main} by the positive quantity $18SD^3$ gives the polynomial
\begin{equation}
 P(u,v,w)=27 D^3-8S\cyc(2v^2+2w^2-u^2)(u+v)^3(u+w)^3 .\label{r32taisho:poly}
\end{equation}
It remains to prove its nonnegativity on the enlarged region \eqref{r32taisho:cone}.


\subsection*{Positive coefficient certificate}
The following two triples of linear forms cover that region as $X,Y,Z$ range over the nonnegative reals:
\[
\begin{array}{c|ccc}
 &u&v&w\\ \hline
 {\rm I}&10X+2Y+Z&6X+Y+Z&5X+Y+Z\\
 {\rm II}&10X+8Y+Z&6X+8Y+Z&5X+5Y+Z
\end{array}
\]
Indeed, for $h\ge4k$ take $(X,Y,Z)=(k,h-4k,w-h-k)$. For $h\le4k$ take $(X,Y,Z)=(h/4,(k-h/4)/3,w-5(h+2k)/6)$. Nonnegativity follows from \eqref{r32taisho:cone}, and direct substitution gives the stated medians in each case.

For each chart write $P=\sum_{j=0}^9\sum_{i=0}^{9-j} c_{ij}X^iY^{9-j-i}Z^j$. The following table gives $\min_{0\le i\le9-j} c_{ij}$:
\begin{center}
\begin{tabular}{r|r|r}
$j$ & Chart I & Chart II\\ \hline
0&33\,048&1\,793\,332\,008\\
1&155\,304&3\,020\,596\,488\\
2&309\,960&1\,841\,309\,496\\
3&339\,984&573\,493\,392\\
4&220\,832&101\,918\,304\\
5&84\,672&10\,507\,968\\
6&17\,664&587\,520\\
7&1\,536&13\,824\\
8&0&0\\
9&0&0
\end{tabular}
\end{center}
For an explicit reproducible expansion rule, attach to each factor $L=aX+bY+cZ$ its coefficient vector $(a,b,c)$. If $A_{ijk}$ are coefficients before multiplying by $L$, the new coefficient at $(i,j,k)$ is
\[
 aA_{i-1,j,k}+bA_{i,j-1,k}+cA_{i,j,k-1},
\]
with out-of-range entries zero and initial array $A_{000}=1$. Apply this rule nine times to each product of linear forms in \eqref{r32taisho:poly}, expanding $2v^2+2w^2-u^2$ into its three terms, and add with the displayed integer multipliers. This yields the table; the accompanying standard-library checker independently implements ordinary multivariate convolution.

All coefficients are nonnegative. The $j=0$ row has strictly positive coefficients, so $P>0$ if $X+Y>0$. The $Z^9$ coefficient vanishes in both charts, and $X=Y=0$ corresponds exactly to $u=v=w$, equivalently $a=b=c$. This proves \eqref{r32taisho:main} and its equality case.\hfill$\square$

\subsection*{Complete integer certificate}
The following self-contained Python listing expands both homogeneous
polynomials using integer arithmetic and verifies every coefficient sign
and every minimum in the table. No numerical approximation is used.
\begingroup\scriptsize
\begin{Verbatim}
from collections import defaultdict

def add(a, b):
    out = dict(a)
    for k, v in b.items():
        out[k] = out.get(k, 0) + v
    return out

def scale(a, s):
    return {k: s*v for k, v in a.items()}

def mul(*polys):
    out = {(0, 0, 0): 1}
    for p in polys:
        new = defaultdict(int)
        for (i, j, k), v in out.items():
            for (a, b, c), w in p.items():
                new[i+a, j+b, k+c] += v*w
        out = new
    return out

def L(a, b, c):
    return {(1, 0, 0): a, (0, 1, 0): b, (0, 0, 1): c}

def cube(x):
    return mul(x, x, x)

def certificate(u, v, w):
    D = mul(add(u, v), add(v, w), add(w, u))
    P = scale(cube(D), 27)
    S = add(add(u, v), w)
    for i, j, k in [(u, v, w), (v, w, u), (w, u, v)]:
        base = add(scale(add(mul(j, j), mul(k, k)), 2),
                   scale(mul(i, i), -1))
        term = mul(S, base, cube(add(i, j)), cube(add(i, k)))
        P = add(P, scale(term, -8))
    return P

charts = [(L(10,2,1), L(6,1,1), L(5,1,1)),
          (L(10,8,1), L(6,8,1), L(5,5,1))]
expected = [
    [33048,155304,309960,339984,220832,84672,17664,1536,0,0],
    [1793332008,3020596488,1841309496,573493392,101918304,
     10507968,587520,13824,0,0]]
for chart, minima in zip(charts, expected):
    P = certificate(*chart)
    assert all(sum(k) == 9 for k in P)
    rows = [[P.get((i,9-j-i,j),0) for i in range(10-j)]
            for j in range(10)]
    assert [min(row) for row in rows] == minima
    assert all(min(row) > 0 for row in rows[:8])
    assert all(not any(row) for row in rows[8:])
print("Both full coefficient certificates verified exactly.")
\end{Verbatim}

\endgroup

\begin{thebibliography}{1}
\bibitem{r32taisho:liu} J. Liu, An inequality involving medians and sides of an acute triangle, \emph{Int. J. Open Problems Comput. Math.} \textbf{14} (2021), no.~3, 49--63; \href{https://www.ijopcm.org/Vol/2021/3.5.pdf}{Conjecture 4.3, p.~62}.
\end{thebibliography}

\par\endgroup

% Source: /workspace/scratch/81914cec2c78/work/archive/Tournament_Round_32_Final/Tournament_Round_32/round32/heisei/paper.tex
% Edited source: /workspace/scratch/81914cec2c78/work/archive/Tournament_Round_32_Final/Tournament_Round_32/round32/heisei/paper.tex
\begingroup
\newcommand{\D}{\mathbb D}
\newcommand{\T}{\mathbb T}
\newcommand{\R}{\mathbb R}
\Needspace{9\baselineskip}
\section{\texorpdfstring{Sharp odd derivatives for holomorphic functions with bounded real part}{Sharp odd derivatives for holomorphic functions with bounded real part}}\label{section:r32heisei}
\noindent{\small\sffamily\color{muted}Round 32 \enspace / \enspace Heisei}\par\medskip


\begin{sectionabstract}
We prove Kresin's conjectured sharp estimate for every odd-order derivative at every point of the unit disk. The normalized phase of the conjugate Schwarz derivative kernel is a finite Blaschke product. It sends arclength weighted by the modulus of that kernel to a constant multiple of uniform arclength, so every direction has exactly the same norm. This gives the sharp coefficient and an explicit analytic extremizer for every direction.
\end{sectionabstract}

\subsection{The sharp pointwise norm}
For a holomorphic function $f$ on $\D=\{z\in\mathbb C:|z|<1\}$, write $M(f)=\sup_{w\in\D}|\Re f(w)|$. The imaginary part of $f$ need not be bounded. For $n\ge1$, let $K_n(z)=\sup_{M(f)\le1}|f^{(n)}(z)|$.

\begin{theorem}\label{r32heisei:main}
For every integer $m\ge0$ and all $z\in\D$, with $r=|z|$,
\[
 K_{2m+1}(z)=\frac{4}{\pi}\,\frac{(2m+1)!}{(1-r^2)^{2m+1}}\,
                  \sum_{j=0}^m\binom{m}{j}^{\!2}r^{2j}.
\]
In particular $|f^{(2m+1)}(z)|\le K_{2m+1}(z)M(f)$, and for every $z$ the factor cannot be reduced. For each fixed $z$, the possible values of $f^{(2m+1)}(z)$ under $M(f)\le1$ form the full closed disk of the displayed radius.
\end{theorem}

For $m\ge1$ this is the sharp estimate (25) of Kresin, Conjecture 2 \cite{r32heisei:K}; $m=0$ recovers the classical first-derivative constant. Sz\'asz's odd-order theorem for the different hypothesis $\sup|f|\le1$ yields the same radial polynomial but lacks the factor $4/\pi$. The change from modulus-bounded to real-part-bounded analytic functions is essential. Earlier work also determined odd-order coefficients in the upper half-plane and boundary asymptotics for the disk \cite{r32heisei:KM}; neither is a pointwise evaluation at all disk radii.

\subsection{A weighted uniform phase}
Write $\zeta=e^{it}$, $0\le t<2\pi$. Rotating the input reduces the problem to $z=r$ real. Put
\[
 H(\zeta)=\frac{\zeta}{(\zeta-r)^{2m+2}},\quad
 A(\zeta)=|\zeta-r|^{-2m-2},\quad
 B(w)=w^m\left(\frac{w-r}{1-rw}\right)^{m+1}.
\]
The function $B$ is holomorphic near the closed disk, maps $\D$ into $\D$, and has unit modulus on $\T=\partial\D$. The identity $\overline{\zeta-r}=(1-r\zeta)/\zeta$ on $\T$ gives
\begin{equation}\label{r32heisei:angle}
 H(\zeta)=A(\zeta)\overline{B(\zeta)}.
\end{equation}

\begin{lemma}\label{r32heisei:push}
Let $J=\int_0^{2\pi} A(e^{it})\,dt$. For every continuous $h:\T\to\mathbb C$,
\[
 \int_0^{2\pi} h(B(e^{it})) A(e^{it})\,dt
       =\frac{J}{2\pi}\int_0^{2\pi}h(e^{is})\,ds.
\]
\end{lemma}
\begin{proof}
For each integer $k\ge1$, the boundary function $AB^k$ agrees with
\[
 \frac{w^{mk+m+1}(w-r)^{(k-1)(m+1)}}
 {(1-rw)^{(k+1)(m+1)}}.
\]
Indeed $A=\zeta^{m+1}/[(\zeta-r)^{m+1}(1-r\zeta)^{m+1}]$ on the circle. The displayed rational function is analytic near $\overline\D$ and vanishes at zero. Its circular mean vanishes by the mean value identity. All negative-index moments vanish by conjugation, and the constant moment is $J$. The asserted equality follows first for trigonometric polynomials, then for all continuous $h$ by uniform approximation.
\end{proof}

It follows from \eqref{r32heisei:angle} and Lemma~\ref{r32heisei:push} that for every real $\alpha$,
\begin{equation}\label{r32heisei:average}
 \int_0^{2\pi}\bigl|\Re(e^{i\alpha}H(e^{it}))\bigr|\,dt
   =J\,\frac{1}{2\pi}\int_0^{2\pi}|\cos(s-\alpha)|\,ds
   =\frac{2J}{\pi}.
\end{equation}
No optimization over a distinguished direction is required.

\subsection{The norm, its value and equality}
A bounded harmonic function on $\D$ is the Poisson integral of real boundary data $u\in L^\infty(\T)$ with the same norm (here the function is $\Re f$). Conversely every such real datum gives the holomorphic Schwarz integral
\[
 S[u](w)=\frac1{2\pi}\int_0^{2\pi}
    \frac{e^{it}+w}{e^{it}-w}\,u(e^{it})\,dt,
\]
whose real part is the Poisson integral. Thus $f=S[u]+iC$ for some real $C$. Termwise differentiation, justified on each interior compact set, gives, with $n=2m+1$,
\[
 f^{(n)}(r)=\frac{n!}{\pi}\int_0^{2\pi} H(e^{it})u(e^{it})\,dt.
\]
For each $\alpha$ and $\|u\|_\infty\le1$, the maximum of $\Re(e^{i\alpha}f^{(n)}(r))$ is $n!/\pi$ times the left side of \eqref{r32heisei:average}. It is attained by
\begin{equation}\label{r32heisei:sgn}
 u(e^{it})=\operatorname{sgn}\Re(e^{i\alpha}H(e^{it})).
\end{equation}
Its values at zeros do not affect the integral. Taking the supremum over $\alpha$ therefore gives $K_n(r)=2n!J/\pi^2$. Also every direction attains its supporting line. Since the whole range lies in the disk of this radius, a point on that line must equal the disk's point of tangency. Multiplication of its boundary datum by any scalar in $[0,1]$ then fills the corresponding radius.

To compute $J$, put $\zeta=(w+r)/(1+rw)$ on the circle. Then
\[
 |\zeta-r|=\frac{1-r^2}{|1+rw|},\qquad
 |d\zeta|=\frac{1-r^2}{|1+rw|^2}|dw|.
\]
Consequently
\[
 J=(1-r^2)^{-2m-1}\int_0^{2\pi}|1+re^{is}|^{2m}\,ds
  =\frac{2\pi}{(1-r^2)^{2m+1}}\sum_{j=0}^m\binom{m}{j}^{\!2}r^{2j},
\]
the final identity being termwise Fourier orthogonality of the binomial expansion. This proves Theorem~\ref{r32heisei:main}.

\paragraph{An explicit extremizer.}
One can avoid the boundary integral in writing an equality function. For $r\in[0,1)$ and prescribed real $\alpha$, take
\[
 f_{r,\alpha}(w)=\frac4\pi\arctan\!\left(e^{-i\alpha}B(w)\right),
\]
where the branch of $\arctan$ is real on $(-1,1)$ and zero at $0$. The identity
$\arctan v=(2i)^{-1}\log[(1+iv)/(1-iv)]$ shows that $|\Re\arctan v|<\pi/4$ in $\D$ and that on $\T$ the boundary real part equals $(\pi/4)\operatorname{sgn}(\Re v)$ except at $v=\pm i$. The boundary values therefore match \eqref{r32heisei:sgn} almost everywhere, by \eqref{r32heisei:angle}. Conversely, equality for a prescribed direction forces those normalized boundary values almost everywhere: $A>0$, and $\Re(e^{i\alpha}\overline B)$ vanishes at only finitely many points, since $B$ is nonconstant rational. Hence, for $M(f)\le1$, these are all supporting-direction extremizers up to an imaginary additive constant, after rotating the input as needed.

\begin{thebibliography}{9}
\bibitem{r32heisei:K} G. Kresin, On sharp pointwise estimates for derivatives of analytic functions with bounded real part, \emph{Functional Differential Equations} \textbf{26} (2020), 193--201, Conjecture~2, p.~200, \href{https://doi.org/10.26351/FDE/26/3-4/4}{doi:10.26351/FDE/26/3-4/4}.
\bibitem{r32heisei:KM} G. Kresin and V. Maz'ya, Sharp real-part theorems for high order derivatives, \emph{Journal of Mathematical Sciences} \textbf{181} (2012), 107--125, \href{https://users.mai.liu.se/vlama82/pdf/fulltext2.pdf}{authors' text}.
\end{thebibliography}

\par\endgroup

% Source: /workspace/scratch/81914cec2c78/work/archive/Tournament_Round_32_Final/Tournament_Round_32/round32/reiwa/paper.tex
% Edited source: /workspace/scratch/81914cec2c78/work/archive/Tournament_Round_32_Final/Tournament_Round_32/round32/reiwa/paper.tex
\begingroup
\newcommand{\mr}{\operatorname{mr}}
\Needspace{9\baselineskip}
\section{\texorpdfstring{The second exponential term for minimum-first reverse shuffles}{The second exponential term for minimum-first reverse shuffles}}\label{section:r32reiwa}
\noindent{\small\sffamily\color{muted}Round 32 \enspace / \enspace Reiwa}\par\medskip


\begin{sectionabstract}
For each fixed positive number $k$ of MIN-REV passes, the second exponential
term in the number of permutations sorted in exactly $k$ passes is
$-((k-1)n+1)k^{n-2}$. This answers Open Question 5.2 of Pudwell and Smith.
Passing to the exponential generating function transforms their binomial
convolution into a product, and summing adjacent Eulerian columns cancels
one constant contribution at the second rate.
\end{sectionabstract}

\subsection*{Exact exponential expansion}
In a MIN-REV pass, take the longest increasing prefix, reverse the remaining
suffix, and merge these two lists by taking the smaller of their two heads
until one is exhausted; then append the other. Let $\mr(n,k)$ count permutations
of $[n]$ needing \emph{exactly} $k$ passes to reach the identity, with the identity
requiring zero passes. The question concerns the second exponential base in
the expansion at fixed $k$, not the number sorted in at most $k$ passes.

\begin{theorem}
For $k\ge2$ there are rational polynomials $P_{k,j}$, independent of $n$,
such that for every $n\ge1$,
\[
 \mr(n,k)=(k+1)^{n-1}-((k-1)n+1)k^{n-2}
       +\sum_{j=1}^{k-1} P_{k,j}(n)j^{n-1}.
\]
For $k=1$, $\mr(n,1)=2^{n-1}-1$, giving the same second-term rule.
\end{theorem}

\paragraph{Step 1: the product.} Let $A(i,d)$ be the Eulerian number of
permutations with $d$ descents on $i$ letters, taking $A(0,d)=[d=0]$.
Use the series
\[
 E_d(t)=\sum_{i\ge0} A(i,d)\frac{t^i}{i!},\qquad
 S_h(t)=\sum_{d=0}^hE_d(t),\qquad
 F_k(t)=\sum_{n\ge1}\mr(n,k)\frac{t^{n-1}}{(n-1)!}.
\]
By Theorem 3.16 of \cite{r32reiwa:ps}, for $m\ge1$,
\[
 F_{2m}=E_mS_{m-1},\qquad F_{2m+1}=E_mS_m .
\]
Indeed, their formula for $k=2m$ partitions the $n-1$ labels other than
the maximum into an ordered pair of lists. The first list has exactly $m$
descents, the second at most $m-1$ ascents. For $k=2m+1$, the first list
has at most $m$ descents and the second exactly $m$ ascents. Ascents and
descents are equidistributed. Choosing which $i$ labels are in the
first list produces the binomial factor in exponential series multiplication.

\paragraph{Step 2: Eulerian columns.} Use the classical finite formula
\[
 A(i,d)=\sum_{\ell=0}^d(-1)^\ell
   \binom{i+1}{\ell}(d+1-\ell)^i .
\]
This also respects our convention at $i=0$. Set $B_0(u)=1$ and for
$\ell\ge1$ let $B_\ell(u)=u^\ell/\ell!+u^{\ell-1}/(\ell-1)!$.
Pascal's identity gives
$\sum_{i\ge0}\binom{i+1}{\ell}u^i/i!=e^u B_\ell(u)$.
Consequently the exact expansion of one column is
\[
 E_d(t)=\sum_{\ell=0}^d(-1)^\ell e^{(d+1-\ell)t}
              B_\ell\big((d+1-\ell)t\big).
\]
Write $L_b$ for an unspecified finite sum of terms $P(t)e^{jt}$
with $P\in\mathbb Q[t]$ and integer $1\le j\le b$. An empty
sum is zero. For $d\ge1$ the first two terms give
\[
 E_d=e^{(d+1)t}-(1+dt)e^{dt}+L_{d-1}.
\]
For $h\ge1$, adding $E_{h-1}$ and the preceding columns contributes
$+e^{ht}$ at the second rate, hence
\[
 S_h=e^{(h+1)t}-ht e^{ht}+L_{h-1}.
\]
Also $S_0=E_0=e^t$, so the formula with vanishing second coefficient
applies in the even case $m=1$ as well.

\paragraph{Step 3: the two leading bases.} In $E_mS_h$, for $h=m-1$ or
$h=m$, the largest rate is $m+h+2=k+1$. The next rate is
$k=m+h+1$, obtained by lowering exactly one factor by one unit. Thus
\[
 F_k=e^{(k+1)t}-[1+(m+h)t]e^{kt}+L_{k-1}
     =e^{(k+1)t}-[1+(k-1)t]e^{kt}+L_{k-1}.
\]
The $(n-1)$st derivative of the second term at zero equals
$-k^{n-1}-(k-1)(n-1)k^{n-2}$. Combining the summands gives exactly
the coefficient claimed. Each term in $L_{k-1}$ yields a polynomial
in $n$ times $j^{n-1}$, $1\le j<k$: for example the $N$th
derivative of $t^r e^{jt}$ at zero is $(N)_r j^{N-r}$, where
$(N)_r$ is the falling factorial. Thus no part of the discarded
finite sum competes with the second base as $n\to\infty$ for fixed $k$.
The case $k=1$ is the one-pass count of \cite[Theorem~3.16]{r32reiwa:ps}.

\begin{thebibliography}{1}
\bibitem{r32reiwa:ps} L. Pudwell and R. Smith, Sorting via shuffles with a cut after
the longest increasing prefix, \emph{Theoretical Computer Science}
\textbf{1008} (2024), 114679. Algorithm 3.3, Theorem 3.16, Open Question 5.2.
\href{https://doi.org/10.1016/j.tcs.2024.114679}{doi:10.1016/j.tcs.2024.114679}.
\end{thebibliography}

\par\endgroup

% Source: /workspace/scratch/81914cec2c78/work/archive/Tournament_Round_32_Final/Tournament_Round_32/round32/keio/paper.tex
% Edited source: /workspace/scratch/81914cec2c78/work/archive/Tournament_Round_32_Final/Tournament_Round_32/round32/keio/paper.tex
\begingroup
\newcommand{\A}{\widetilde A}
\newcommand{\E}{\widetilde E}
\Needspace{9\baselineskip}
\section{\texorpdfstring{Strict interlacing across ranks for binomial Eulerian polynomials}{Strict interlacing across ranks for binomial Eulerian polynomials}}\label{section:r32keio}
\noindent{\small\sffamily\color{muted}Round 32 \enspace / \enspace Keio}\par\medskip


\begin{sectionabstract}
The Haglund--Zhang row interlacing theorem can be promoted to strict interlacing between consecutive ranks. We give an explicit recovery of an old row from selected evaluations of the new row. It prevents common roots for every strictly increasing inversion alphabet starting at two or more. Ordinary binomial Eulerian polynomials form a special case. We present this as a deduction from the literature, without a claim of priority for the classical statement.
\end{sectionabstract}
\subsection{The rank question}
For the symmetric group $S_n$ let
\[
 \A_n(x)=\sum_{\pi\in S_n}(1+x)^{\#\{i:\pi(i)=i\}}x^{\#\{i:\pi(i)>i\}}.
\]
Consecutive interlacing of these polynomials is listed as a conjecture in SymCat \cite{r32keio:symcat}. We establish the strict version: if $u_1<\cdots<u_n$ and $v_1<\cdots<v_{n+1}$ are the respective zeros of $\A_n$ and $\A_{n+1}$, then
\[
 v_1<u_1<v_2<\cdots<u_n<v_{n+1}<0.
\]
The proof applies to a wider class. For positive integer alphabets $s_1,\ldots,s_m$ consider vectors $0\le e_j<s_j$ with integer entries. Set $w_j=e_j/s_j$ and append $w_0=w_{m+1}=0$. Give each comparison $w_j$ to $w_{j+1}$ weight $x$, $1+x$, or $1$, according as it increases, ties, or decreases. Let $\E_m$ be the sum of the products of weights, with $\E_0=1+x$ (one tie). Haglund and Zhang \cite[Theorem 3.1 and its proof]{r32keio:hz} identify $\E_{n-1}=\A_n$ when $s_j=j+1$.

\begin{theorem}\label{r32keio:main}
If $2\le s_1<s_2<\cdots$, then for every $m\ge0$, $\E_m$ has $m+1$ simple negative zeros and strictly interlaces $\E_{m+1}$.
\end{theorem}
For clarity, the row interlacing used below is an established result, not the new step of this argument. In their terminology $f\ll g$ for equal-degree polynomials with positive leading coefficients means that the sorted zeros alternate, beginning with a zero of $f$ (weak inequalities allowed).
\subsection{Recovering an old row}
Let $f_{m,j}$ count vectors ending at $e_m=j$ with the \emph{last} comparison $w_m$ to $0$ omitted, keeping all other weights. Their degrees are exactly $m$: the all-zero prefix followed by $j$ contributes a degree-$m$ term, and all coefficients are nonnegative. The established last-entry refinement \cite[Theorem 1.1]{r32keio:hz} gives $f_{m,0}\ll f_{m,j}$ for every $j$. The only vector with no ascents is all zero, so $f_{m,0}(0)=1$ and $f_{m,j}(0)=0$ for $j>0$. At every rank,
\begin{equation}
 \E_m=(1+x)f_{m,0}+\sum_{j>0}f_{m,j}=f_{m+1,0}.\label{r32keio:eq:first}
\end{equation}
Here the last identity means that extending by the entry zero restores exactly the discarded comparison.

\begin{lemma}
For increasing alphabets as in Theorem \ref{r32keio:main}, the evaluations $f_{m,j}(r)$ do not all vanish at any real $r<0$.
\end{lemma}
\begin{proof}
At rank one the row is $(1+x,x,\ldots,x)$ and has at least two entries. Write $p=s_m$, $q=s_{m+1}$, $f_i=f_{m,i}$ and $g_k=f_{m+1,k}$. For each $0\le j<p$, choose $k_j$ with $j/p<k_j/q<(j+1)/p$. This is possible because the rescaled open interval has length $q/p>1$. With $h_j=g_{k_j}$, appending $k_j$ gives
\[
 h_j=x\sum_{i=0}^j f_i+\sum_{i=j+1}^{p-1}f_i.
\]
Consequently $g_0-h_0=f_0$ and $h_j-h_{j-1}=(x-1)f_j$ for $1\le j<p$. If all entries of the new row vanished at $r<0$, every old entry would vanish there, which induction excludes.
\end{proof}
\begin{proof}[Proof of Theorem \ref{r32keio:main}]
The first polynomial $f_{1,0}=1+x=\E_0$ has one simple negative zero. Suppose $f_{m,0}$ has simple negative zeros $r_1<\cdots<r_m$. Since $f_{m,0}\ll f_{m,j}$ and the degrees are $m$, each value $f_{m,j}(r_i)$ has weak sign $(-1)^{m-i+1}$: unless it is zero, precisely $m-i+1$ of its zeros exceed $r_i$. The lemma shows that at least one of these values is nonzero. Equation \eqref{r32keio:eq:first} thus gives that same \emph{strict} sign for $f_{m+1,0}(r_i)$. Its constant term is one, its leading coefficient is positive, and its degree is $m+1$. Its signs alternate at $-\infty,r_1,\ldots,r_m,0$, giving one zero in each of the $m+1$ disjoint open intervals. This exhausts the degree, so each zero is simple. Applying \eqref{r32keio:eq:first} again finishes the induction.
\end{proof}
\paragraph{Provenance.}
The weak conclusion follows readily from the interlacing cone in \cite{r32keio:hz}; the recovery argument supplies strictness uniformly for growing alphabets. Liu and Yan \cite{r32keio:ly} give generalized Sturm criteria for binomial-transform diagonals. We have not excluded coverage of the classical specialization by their criteria. Neither the online open-problem listing nor this proof establishes historical novelty.
\begin{thebibliography}{9}
\bibitem{r32keio:symcat} P. Alexandersson, \emph{Real-rooted polynomials: open problems}, SymCat, section on interlacing of Eulerian-type polynomials, \url{https://www.symmetricfunctions.com/realRootedOpenProblems.htm}.
\bibitem{r32keio:hz} J. Haglund and P. B. Zhang, \emph{Real-rootedness of variations of Eulerian polynomials}, Adv. Appl. Math. 109 (2019), 38--54. \url{https://zhangbiaomath.github.io/papers/2019-aam-bino_eulerian.pdf}.
\bibitem{r32keio:ly} L. L. Liu and X. Yan, \emph{Some interlacing properties related to the Eulerian and derangement polynomials}, Adv. Appl. Math. 162 (2025), 102776. \url{https://doi.org/10.1016/j.aam.2024.102776}.
\end{thebibliography}

\par\endgroup

% Source: /workspace/scratch/81914cec2c78/work/archive/Tournament_Round_33_Final/Tournament_Round_33/round33/showa/paper.tex
% Edited source: /workspace/scratch/81914cec2c78/work/audit_26_33/cleaned_sources/r33_showa.tex
\begingroup
\newcommand{\sinc}{\operatorname{sinc}}
\Needspace{9\baselineskip}
\section{\texorpdfstring{Sharp bounds for the real-parameter sine quotient}{Sharp bounds for the real-parameter sine quotient}}\label{section:r33showa}
\noindent{\small\sffamily\color{muted}Round 33 \enspace / \enspace Showa}\par\medskip

\begin{sectionabstract}
We give the optimal upper and lower constants for the D'Aurizio--S\'andor sine quotient at every real $p>1$. For $p\ge2$, including nonintegers, the quotient decreases strictly on $(0,\pi/2)$. A fractional-frequency remainder is bounded on one initial band, and a two-step difference identity covers all larger parameters. For $1<p<2$, both optimal constants are characterized by endpoint tests and at most one interior critical point. The previously established initial bands are credited and their proofs are included.
\end{sectionabstract}
\subsection{The sine quotient}
Male\v sevi\'c et al.\ \cite{r33showa:mm} ask in their concluding section for a real-parameter extension, with best constants, of their sine inequality (15). Hung and Li \cite{r33showa:hl} proved the bounds for integers $p\ge2$. Write
\[
F_p(x)=\frac{p-\sin(x)/\sin(x/p)}{x^2},\quad
T=\frac{\pi}{2p},\quad
R_p(t)=\frac{\sin(pt)}{\sin t},\quad H_p(u)=\frac{p-R_p(\sqrt{u})}{u}.
\]
Then $F_p(x)=p^{-2}H_p(x^2/p^2)$. At $u=0,T^2$, respectively, the continuous endpoint values of $H_p$ are
\[
L=\frac{p(p^2-1)}6,\qquad U=\frac{p-\csc T}{T^2}.
\]
\begin{theorem}
For all real $p\ge2$ and $0<x<\pi/2$, the function $F_p$ decreases strictly and its sharp bounds are
\[
\frac4{\pi^2}\left(p-\csc\frac\pi{2p}\right)< F_p(x)<\frac{p^2-1}{6p}.
\]
For $1<p<2$ the supremum of $F_p$ is $p^{-2}\max(L,U)$. Put
\[
 D_0=-\frac{p(p^2-1)(3p^2-7)}{360},\qquad
 D_T=\frac{T\csc T\cot T-2(p-\csc T)}{2T^4}.
\]
If $D_0\ge0$ then $\inf F_p=L/p^2$; if $D_T\le0$ then $\inf F_p=U/p^2$.
Otherwise there is exactly one $u_*\in(0,T^2)$ with $H_p'(u_*)=0$ and $\min F_p=H_p(u_*)/p^2$. Endpoint extrema are not attained on the open interval. For negative parameters of absolute value greater than one, $F_{-p}=-F_p$, so the two bounds exchange sign and order.
\end{theorem}

\begin{proof}
Put $K_a(u)=(1-\cos(a\sqrt u))/u$ and $J_q(u)=(R_q(\sqrt u)-q)/u$ for $q\in[0,1]$. Euler's product gives
\[
 R_q(\sqrt u)=q\prod_{n=1}^{\infty}
 \left(1+\frac{(1-q^2)u}{\pi^2 n^2-u}\right).
\]
Local uniform convergence for $|u|<\pi^2$ and coefficientwise nonnegativity of the finite products show that all Taylor coefficients $b_j$ of $R_q(\sqrt u)$ are nonnegative. In particular $J_q',J_q''\ge0$ for $0<u<\pi^2$; $J_0=J_1=0$. Writing $\sinc z=\sin(z)/z$ with value $1$ at $0$ yields
\[
K_a(u)=a^2\int_0^1(1-v)\cos(av\sqrt u)\,dv,\qquad
-2K_a'(u)=a^4\int_0^1(1-v)v^2\sinc(av\sqrt u)\,dv.
\]
On $0<a\sqrt u<\pi/2$, $K_a'<0$ and $K_a''>0$; the latter follows because $\sin z-z\cos z$ is positive there (its derivative is $z\sin z$), so $\sinc z$ is strictly decreasing.

The addition formula $R_p(t)=2\cos((p-1)t)+R_{p-2}(t)$ implies the two-step identity
\[
 H_p=2K_{p-1}+H_{p-2}.
\]
For $2\le p\le3$, $H_p=2K_{p-1}-J_{p-2}$ is strictly decreasing. For $1<p<2$, $H_p=2K_{p-1}+J_{2-p}$ is strictly convex. Its endpoint derivatives are $D_0,D_T$, obtained from the Taylor expansion at zero and $R_p'(T)=-\csc T\cot T$, respectively. This proves the asserted classification on $1<p\le3$ from \cite{r33showa:taisho}.

It remains to extend strict decrease to all larger real parameters. First take $3\le p\le4$. In this band
\[
 H_p=2K_{p-1}+2K_{p-3}+J_{4-p}.
\]
For $q=4-p$, the nonnegative coefficients just established obey $b_j4^j\le R_q(2)=\sin(2q)/\sin 2<3$, since $\sin 2\ge 2-2^3/6>1/3$. Thus with $0<u<T^2\le\pi^2/36<1/3$,
\[
 J_q'(u)=\sum_{j\ge2}(j-1)b_ju^{j-2}
 \le \frac{3}{16(1-u/4)^2}<\frac14.
\]
For $a=p-1\ge2$, use $\sinc z\ge2/\pi$ for $0\le z\le\pi/2$ to get
\[
 -2K_{p-1}'(u)\ge \frac{a^4}{6\pi}>\frac14.
\]
The other cosine term has nonpositive derivative, so $H_p'<0$. For $p>4$ iterate the two-step identity until the remaining parameter is in $[2,4]$. Every intermediate parameter $b$ has $T=\pi/(2p)\le\pi/(2b)$; every added $K_{b-1}$ is strictly decreasing on the required interval. This finishes the proof, including sharpness by the endpoint limits.
\end{proof}

\begin{thebibliography}{9}
\bibitem{r33showa:mm} B. Male\v sevi\'c, M. Mi\'covi\'c, B. Mihailovi\'c and T. Lutovac, A new approach to selecting constants for some analytic inequalities, \emph{J. Math. Inequal.} 19 (2025), 223--245. \url{https://doi.org/10.7153/jmi-2025-19-15}.
\bibitem{r33showa:hl} L.-C. Hung and P.-Y. Li, On generalization of D'Aurizio--S\'andor inequalities involving a parameter, \emph{J. Math. Inequal.} 12 (2018), 853--860. \url{https://doi.org/10.7153/jmi-2018-12-63}.
\bibitem{r33showa:taisho} Taisho, \emph{Endpoint tests for a real-parameter sine quotient}, internal Round 25 archive, 19 Nov. 2026. Source of the range $1<p\le3$ and the two-step identity; not a journal publication.
\end{thebibliography}

\par\endgroup

% Source: /workspace/scratch/81914cec2c78/work/archive/Tournament_Round_33_Final/Tournament_Round_33/round33/heisei/paper.tex
% Edited source: /workspace/scratch/81914cec2c78/work/archive/Tournament_Round_33_Final/Tournament_Round_33/round33/heisei/paper.tex
\begingroup
\newcommand{\R}{\mathbb R}
\newcommand{\Z}{\mathbb Z}
\Needspace{9\baselineskip}
\section{\texorpdfstring{Reciprocal-power curvature and the Guo--Qi square-root polygamma conjecture}{Reciprocal-power curvature and the Guo--Qi square-root polygamma conjecture}}\label{section:r33heisei}
\noindent{\small\sffamily\color{muted}Round 33 \enspace / \enspace Heisei}\par\medskip


\begin{sectionabstract}
For every positive integral order we prove the monotonicity and both endpoint limits conjectured by Guo and Qi for a difference of polygamma functions at $1+x$ and $\sqrt{1+x}$. The derivative is a tangent-line remainder for the Hurwitz zeta function in a reciprocal-square coordinate. More generally, we classify all reciprocal powers giving global convexity or concavity and show uniqueness of the inflection in the intermediate regime. A separate appendix gives a finite-difference proof of Kresin's sign law for an alternating trigonometric sum.
\end{sectionabstract}
\subsection{The square-root conjecture}
Put $\psi=\Gamma'/\Gamma$. In Remark 1.11, (1.22)--(1.23), of \cite{r33heisei:GQ}, Guo and Qi conjectured a monotonicity statement and endpoint limits for
\[
 g_i(x)=\psi^{(i)}(x+1)-x\psi^{(i+1)}(\sqrt{x+1}),\qquad x>-1.
\]
The linear-argument function of their preceding Theorem 1.10 is different. We establish the stated conjecture without additional restrictions.
\begin{theorem}\label{r33heisei:gq}
For each positive odd integer $i$, $g_i$ is strictly decreasing on $(-1,\infty)$; for each positive even integer $i$, it is strictly increasing. For every such $i$,
\[
 \lim_{x\downarrow-1}g_i(x)=(-1)^{i+1}\infty,\qquad
 \lim_{x\to\infty}g_i(x)=
 \begin{cases}1,&i=1,\\0,&i\ge2.\end{cases}
\]
Its derivative vanishes at $x=0$ and nowhere else in the domain.
\end{theorem}
For $s>1$ and $t>0$, let $Z_s(t)=\sum_{k=0}^\infty(t+k)^{-s}$ be the Hurwitz zeta function. We first give the short curvature estimate actually required by the conjecture.

\begin{lemma}\label{r33heisei:curve}
For $s\ge3$ the function $U(v)=Z_s(v^{-1/2})$ satisfies $U''(v)>0$ on $(0,\infty)$.
\end{lemma}
\begin{proof}
Writing $t=v^{-1/2}$ and differentiating twice yields
\[
 U''(v)=\frac{t^5}{4}\left[t Z_s''(t)+3Z_s'(t)\right]
 =\frac{s t^5}{4}\big[(s+1)tZ_{s+2}(t)-3Z_{s+1}(t)\big].
\]
For $s>1$, use the defining series and the gamma integral for each summand to get
\begin{equation}\label{r33heisei:int}
 t^sZ_{s+1}(t)=\frac1{\Gamma(s+1)}
 \int_0^\infty\frac{u^s e^{-u}}{t(1-e^{-u/t})}\,du.
\end{equation}
For each $u>0$, the denominator has derivative
$1-(1+u/t)e^{-u/t}>0$, so the integral has strictly negative derivative. Differentiation is valid on compact $t$-intervals by an exponential majorant at infinity and local integrability at zero (the integrand is $O(u^{s-1})$ there, and its $t$-derivative is $O(u^{s})$). It follows that $(s+1)tZ_{s+2}(t)>sZ_{s+1}(t)$, which proves the lemma including $s=3$.
\end{proof}

\begin{proof}[Proof of Theorem~\ref{r33heisei:gq}]
Let $t=\sqrt{1+x}$ and $p=i+1$. The usual polygamma series gives
\[
 G(x):=\frac{(-1)^{i+1}}{i!}g_i(x)
 =Z_p(t^2)+p(t^2-1)Z_{p+1}(t).
\]
All differentiations of $Z_s$ follow locally uniformly from the series; in particular $Z_s'=-sZ_{s+1}$. Setting $s=p+1=i+2$, we obtain
\begin{equation}\label{r33heisei:rem}
 -\frac1pG'(x)=Z_s(t^2)-Z_s(t)-\frac{t^2-1}{2t}Z_s'(t).
\end{equation}
In the variable $v=t^{-2}$ the right-hand side is precisely
$U(v^2)-U(v)-(v^2-v)U'(v)$: the product of $v^2-v$ and $dt/dv=-t^3/2$ equals $(t^2-1)/(2t)$. By Lemma~\ref{r33heisei:curve} this tangent remainder is strictly positive unless $v^2=v$, which here means $x=0$. Thus $G'(x)\le0$ everywhere and is strictly negative off a single point. Integration on each nondegenerate subinterval gives strict decrease of $G$ and the specified parity rule for $g_i$.

For $a>1$, the defining series gives $Z_a(u)=u^{-a}+O(1)$ at zero. At infinity the elementary inequalities
\[
 \frac{u^{1-a}}{a-1}\le Z_a(u)\le u^{-a}+\frac{u^{1-a}}{a-1}
\]
follow by integral comparison. Consequently, at $t\downarrow0$ the $t^{-2p}$ leading term of $Z_p(t^2)$ dominates the $O(t^{-p-1})$ second term of $G$, because $p\ge2$. At $t\to\infty$, the first term tends to zero, and the second is asymptotic to $t^{2-p}$. For $p=2$ the normalizing sign and factorial are both one; this gives the exceptional limit. All parts of the theorem follow.
\end{proof}

\subsection{The threshold for a general reciprocal power}
The exponent two in Lemma~\ref{r33heisei:curve} is sharp when $s=3$. The following classification identifies precisely where the tangent argument is available.
\begin{theorem}\label{r33heisei:threshold}
For real $s>1$, $h>0$, set $U_{s,h}(v)=Z_s(v^{-1/h})$ on $v>0$. It has everywhere positive second derivative exactly when $h\le s-1$, and everywhere negative second derivative exactly when $h\ge s$. For $s-1<h<s$, its second derivative changes sign exactly once, from negative to positive as $v$ increases.
\end{theorem}
\begin{proof}
Putting $t=v^{-1/h}$ gives
\[
 U_{s,h}''(v)=\frac{s t^{2h+1}Z_{s+1}(t)}{h^2}\big(R_s(t)-h-1\big),
 \qquad R_s(t)=\frac{(s+1)t Z_{s+2}(t)}{Z_{s+1}(t)}.
\]
We show that $R_s$ decreases strictly from $s+1$ to $s$ as $t$ goes from zero to infinity. The limits follow from the same endpoint estimates as in Theorem~\ref{r33heisei:gq}. For the strict monotonicity, put $b=s+1$. Up to $1/\Gamma(b)$, $Z_b(e^y)$ is the integral in $u>0$ of
\[
 u^{b-1}e^{-u}\,\frac{e^{-by}}{1-\exp(-ue^{-y})}.
\]
For fixed $u>0$, differentiation in $y$ of the log of the second factor gives $-b+z/(e^z-1)$, with $z=ue^{-y}$. The function $z/(e^z-1)$ decreases strictly for $z>0$ (differentiate, using $1-z<e^{-z}$), so this log derivative increases strictly in $y$. Thus the integrand is strictly log-convex in $y$ pointwise. For any two distinct $y$'s, multiply their integrands in geometric proportions; strict pointwise log-convexity followed by H\"older's inequality proves strict log-convexity of the integral. A differentiable strictly convex function has a strictly increasing derivative, so $R_s(e^y)=-(\log Z_b(e^y))'$ decreases strictly. Since $v\mapsto t$ reverses the order, the assertions now follow.
\end{proof}

\subsection*{Appendix: an independent finite-difference proof of Kresin's sign law}
Conjecture 1 in \cite{r33heisei:Kr}, separate from that paper's odd-derivative disk Conjecture 2, concerns the literal finite sum
\[
 \Lambda_m(\varphi)=\sum_{\ell=1}^m(-1)^\ell(2\ell-1)
 \binom{2m-1}{m-\ell}
 \frac{\sin((2\ell-1)\varphi/(2m+1))}
 {\sin((2\ell-1)\pi/[2(2m+1)])}.
\]
We prove $(-1)^m\Lambda_m(\varphi)>0$ for every integer $m\ge1$ and $0<\varphi<\pi/2$. The even-order upper-half-plane phase-direction consequence also follows from earlier folded-weight arguments for Kalaj--Elkies; this appendix provides a different proof, not a second independent open-problem credit.

Put $d=2m-1$, $N=2m+1$, $a=\varphi/N$, $b=\pi/(2N)$ and $f(u)=u\sin(au)/\sin(bu)$, continued at zero. Let $\Delta_2f(u)=f(u+2)-f(u)$. Since $f$ and $d$ are odd, pairing opposite nodes gives
\[
 \Delta_2^d f(-d)=\sum_{j=0}^d(-1)^{d-j}\binom dj f(-d+2j)
                 =2(-1)^m\Lambda_m(\varphi).
\]
Indeed the node $2\ell-1$ corresponds to $j=m+\ell-1$. On $|u|<\pi/b$, Euler's sine product gives
\[
 \frac{\sin(au)}{\sin(bu)}
 =\frac ab\prod_{k\ge1}\frac{1-a^2u^2/(k\pi)^2}{1-b^2u^2/(k\pi)^2}.
\]
Since $0<a<b$, each individual factor has constant coefficient one and positive coefficients in every positive even degree; all odd coefficients vanish. Compact convergence of the products on this disk gives convergence of the Taylor coefficients, and the first factor already supplies strict positivity in each even degree. Thus $f(u)=\sum_{k\ge0}c_k u^{2k+1}$ with $c_k>0$. Every surviving monomial of $f^{(d)}$ has nonnegative even power on the real interval $[-d,d]$ contained in the disk; the constant term is strictly positive. Finally,
\[
 \Delta_2^d f(-d)=\int_{[0,2]^d} f^{(d)}(-d+t_1+\cdots+t_d)\,dt_1\cdots dt_d>0
\]
by the fundamental theorem of calculus, yielding the asserted sign. In particular \cite[(18)--(20)]{r33heisei:Kr} gives the maximizing phase $\pi/2$ for even $m$ and $0$ for odd $m$, relative to $[0,\pi/2]$.

\begin{thebibliography}{9}
\bibitem{r33heisei:GQ} B.-N. Guo and F. Qi, Some properties of the psi and polygamma functions, \emph{Hacet. J. Math. Stat.} \textbf{39} (2010), 219--231, \href{https://dergipark.org.tr/en/download/article-file/86737}{published text}, Remark 1.11 on p.~222.
\bibitem{r33heisei:Kr} G. Kresin, On sharp pointwise estimates for derivatives of analytic functions with bounded real part, \emph{Funct. Differ. Equ.} \textbf{26} (2020), 193--201, \href{https://doi.org/10.26351/FDE/26/3-4/4}{doi:10.26351/FDE/26/3-4/4}, (18)--(20) and Conjecture~1.
\end{thebibliography}

\par\endgroup

% Source: /workspace/scratch/81914cec2c78/work/archive/Tournament_Round_33_Final/Tournament_Round_33/round33/reiwa/paper.tex
% Edited source: /workspace/scratch/81914cec2c78/work/archive/Tournament_Round_33_Final/Tournament_Round_33/round33/reiwa/paper.tex
\begingroup
\newcommand{\pr}{\mathcal P}
\newcommand{\av}{\operatorname{Av}}
\Needspace{9\baselineskip}
\section{\texorpdfstring{An irredundant classical basis for iterated PRE-REV sorting}{An irredundant classical basis for iterated PRE-REV sorting}}\label{section:r33reiwa}
\noindent{\small\sffamily\color{muted}Round 33 \enspace / \enspace Reiwa}\par\medskip


\begin{sectionabstract}
For all $k>2$, we give the explicit classical patterns avoided by permutations sortable in at most $k$ iterations of PRE-REV, as asked by Pudwell and Smith. Every minimal forbidden permutation is obtained by joining two descent-minimal permutations at a common valley, with the second read backwards. Their sizes lie between $k+2$ and $2k+1$, and both bounds are attained. The result is a basis formula, not just a finite list requiring a subsequent minimality test.
\end{sectionabstract}
\subsection{Question and answer}
Let $\pr_k$ be the permutations sorted in at most $k$ applications of the deterministic shuffle PRE-REV from \cite{r33reiwa:PS}. A classical pattern occurs if some subsequence has its relative order, and $\av(B)$ denotes avoidance of every pattern in $B$. The basis of a class consists of the minimal forbidden permutations. Open Question 5.5 of \cite{r33reiwa:PS} asks for the explicit permutation classes when $k>2$. The original paper treats the at-most-$k$ classes; identity permutations stay sorted. If instead one wants permutations needing exactly $k$ passes, the answer is $\pr_k\setminus\pr_{k-1}$, not generally a classical avoidance class.

Write $M_d$ for permutations having exactly $d$ adjacent descents such that every single-entry deletion decreases the number of descents. Deletion cannot increase this number, so these are precisely the permutations minimal for having at least $d$ descents. These ingredients are explicit: by Bouvel and Pergola \cite[Thm.~3]{r33reiwa:BP}, a permutation with $d$ descents is in $M_d$ if and only if every ascent at $i$ lies strictly inside the word and its consecutive four-letter window at indices $i-1,i,i+1,i+2$ has pattern $2143$ or $3142$.

\begin{theorem}\label{r33reiwa:T}
Let $k\ge1$, $a=\lceil(k+1)/2\rceil$, $b=\lfloor(k+1)/2\rfloor$. Form $B_{a,b}$ as follows. Choose $u\in M_a$ and $v\in M_b$. Glue their final entries together and read first $u$, then the reversal of $v$, counting the common point only once. Take every permutation labeling of this shape that preserves the orders of the two pieces. Then $B_{a,b}$ is the unique irredundant basis of $\pr_k$. Its minimum and maximum lengths are $k+2$ and $2k+1$.
\end{theorem}
The two sides may interleave in value; the common entry is a local minimum and is not required to be the global minimum. For instance $k=1$ yields $213,312$. For $k=2$ the rule gives the eight patterns in \cite[Prop.~4.1]{r33reiwa:PS}.

\subsection{The common-valley decomposition}
Use $D,A$ for the adjacent comparison signs down and up; an occurrence of a sign word below is as a subsequence of this list of signs, not necessarily consecutively.

\begin{lemma}\label{r33reiwa:L}
A permutation fails to be a juxtaposition of $a$ increasing blocks followed by $b$ decreasing blocks if and only if its sign word contains $D^a A^b$. Empty blocks are allowed.
\end{lemma}
\begin{proof}
Fill the first block greedily with the maximal initial monotone segment in its required direction. If a legal segmentation exists, this choice preserves feasibility: if it extends past the original first cut, remove the additionally consumed letters from the initial parts of the later blocks. Repeat in the remaining suffix. A cut is made at the first adjacent comparison violating the current block's direction; the second endpoint of that comparison begins the next block. Thus successive cuts are at the earliest possible occurrences of $D,\ldots,D,A,\ldots,A$, with each comparison strictly later than the preceding one. All blocks are exhausted while leaving a nonempty remainder exactly if all $a+b$ letters of the test word are found.
\end{proof}
This obstruction is an instance of the increase--decrease pattern formalism \cite{r33reiwa:HR}. By the juxtaposition characterization in the proof of \cite[Thm.~4.5]{r33reiwa:PS}, the lemma's class with the displayed $a,b$ is exactly $\pr_k$. We analyze the minimal obstructions specific to the sign change.

\begin{proof}[Proof of Theorem~\ref{r33reiwa:T}]
Consider any permutation minimal outside the class. Let $e$ be its $a$th down-edge starting from the left, and select the first up-edge after $e$. It exists by Lemma~\ref{r33reiwa:L}, and there are at least $b$ up-edges from it onward. The two selected edges must be adjacent: otherwise $e$ is immediately followed by another down-edge. Deleting the middle point of those two down-edges leaves one down-edge in their place, still completing $a$ down-edges before all $b$ up-edges. This contradicts minimality. Denote their shared valley by $z$ and let $L,R$ be the prefix ending and the suffix starting at $z$. Then $L$ has exactly $a$ descents and $R$ has exactly $b$ ascents: if $R$ had more, deleting the last point would retain the witness. Deletion of any non-valley point of $L$ must reduce its descent count, because otherwise the entire witness persists. Deletion of the final point of $L$ also reduces its count because $L$ ends in a descent. Thus $L$ has pattern in $M_a$. Equally, any deletion in $R$ must cost an ascent (including deletion of its first point since it starts with an ascent). Its reversal therefore has pattern in $M_b$. In particular, the pivot is determined by the $a$th descent, independent of any choice of subsequence.

Conversely take any gluing in the statement. Each descent-minimal half ends in a descent (else its last point could be deleted). Thus $L$ ends down and $R$ starts up, and together they supply $D^aA^b$. If a point other than $z$ is removed from $L$, at most $a-1$ descents remain in $L$. Any $a$th descent now lies in $R$, strictly after its initial ascent, leaving fewer than $b$ following ascents. If a non-valley point is removed from $R$, at most $b-1$ ascents remain there; the earliest $a$th descent is still the last edge of $L$. Finally, deleting $z$ replaces its two incident edges with one new comparison. Before that edge there are only $a-1$ descents; after it there are only $b-1$ ascents. A copy of the witness would have to use this edge both as a descent and as an ascent, which is impossible. Every single deletion thus belongs to the class. Since deleting points preserves membership of a juxtaposition, these are exactly the basis elements.

By \cite[Prop.~2]{r33reiwa:BP}, each member of $M_d$ has length between $d+1$ and $2d$. The decreasing permutation of length $d+1$ and the word $21\,43\,65\cdots(2d)(2d-1)$ realize the endpoints, by the four-letter test above. Any chosen pair of total orders with just one identified point admits a consistent labeling (merge their segments below and above the common point separately). The two extreme combined lengths, with the pivot counted once, are therefore $a+b+1$ and $2a+2b-1$.
\end{proof}

\paragraph{Scope and checks.}
Finite basing of juxtaposition is classical \cite{r33reiwa:At}; the assertion here is the irredundant basis with no filtering, not merely finite basing. No universal bibliographic priority claim is made. The companion script compares the two-half local criterion with deletion-minimality computed independently from the greedy prefix--suffix pass count \cite[Thm.~3.6]{r33reiwa:PS}, for all permutations through size eight and $1\le k\le6$.
\begin{thebibliography}{9}
\bibitem{r33reiwa:PS} L.~Pudwell and R.~Smith, Sorting via shuffles with a cut after the longest increasing prefix, \emph{Theoret. Comput. Sci.} 1008 (2024), 114679. \href{https://doi.org/10.1016/j.tcs.2024.114679}{doi:10.1016/j.tcs.2024.114679}.
\bibitem{r33reiwa:BP} M.~Bouvel and E.~Pergola, Posets and permutations in the duplication-loss model: minimal permutations with $d$ descents, \href{https://arxiv.org/abs/0806.1494}{arXiv:0806.1494}.
\bibitem{r33reiwa:At} M.~D.~Atkinson, Restricted permutations, \emph{Discrete Math.} 195 (1999), 27--38.
\bibitem{r33reiwa:HR} S.~Huczynska and N.~Ru\v{s}kuc, Pattern classes of permutations via bijections between linearly ordered sets, \emph{European J. Combin.} 29 (2008), 118--139, Sec.~4.
\end{thebibliography}

\par\endgroup

% Source: /workspace/scratch/81914cec2c78/work/archive/Tournament_Round_33_Final/Tournament_Round_33/round33/keio/paper.tex
% Edited source: /workspace/scratch/81914cec2c78/work/audit_26_33/cleaned_sources/r33_keio.tex
\begingroup
\newcommand{\R}{\mathbb R}
\newcommand{\norm}[1]{\lVert #1\rVert}
\Needspace{9\baselineskip}
\section{\texorpdfstring{A single fourth-power potential characterizes the regular simplex}{A single fourth-power potential characterizes the regular simplex}}\label{section:r33keio}
\noindent{\small\sffamily\color{muted}Round 33 \enspace / \enspace Keio}\par\medskip


\begin{sectionabstract}
If the sum of fourth powers of distances from $d+1$ sites in $\mathbb R^d$ is constant on one sphere of positive radius, the sites either are all its center or form a regular simplex centered there. No common radius for the sites is assumed. This gives the full distinct-sites conclusion of the simplex conjecture of Nikolov and Rafailov without its quadratic-potential hypothesis.
\end{sectionabstract}
\subsection{A quartic inverse problem}
Let $d\ge1$, $O\in\R^d$, $R>0$, and let $a_1,\ldots,a_{d+1}$ denote site vectors relative to $O$. On $\norm{x}=R$ define $U_\lambda(x)=\sum_i\norm{x-a_i}^\lambda$. Nikolov and Rafailov ask whether constancy of both $U_2$ and $U_4$ on a sphere forces the sites to be the vertices of a regular simplex; see \cite[arXiv version, Conjecture 2, p.~11]{r33keio:nr}, which gives the low-dimensional cases. The case of repeated sites is handled explicitly below.
\begin{theorem}\label{r33keio:simplex}
The function $U_4$ is constant on $\norm{x}=R$ if and only if all sites are $O$, or they are the vertices of a regular $d$-simplex with center $O$.
\end{theorem}
\begin{proof}
Let $A$ be the $d$ by $d+1$ matrix with columns $a_i$ and put $w_i=R^2+\norm{a_i}^2>0$. On the sphere,
\[
 U_4(Ru)=\sum_i w_i^2-4R u^T A w+4R^2u^TAA^Tu,\qquad \norm u=1.
\]
Comparison at $u$ and $-u$ yields $Aw=0$. The even part now shows $AA^T=cI_d$ for some $c\ge0$, since the Rayleigh quotient is constant on all unit vectors. If $c=0$, every column vanishes. Otherwise $A$ has rank $d$, so its nullspace is exactly $\R w$. The matrix $c^{-1}A^TA$ is the orthogonal projection onto $w^\perp$: it is symmetric and idempotent and its nullspace is $\R w$. Consequently,
\[
 A^TA=c\left(I_{d+1}-\frac{ww^T}{\norm w^2}\right).
\]
The diagonal entries give, for each $i$,
\[
 w_i+\frac{c}{\norm w^2}w_i^2=R^2+c.
\]
The left-hand side is strictly increasing in $w_i>0$. Thus all $w_i$ are equal, $\sum_i a_i=0$, and the Gram matrix is $c(I_{d+1}-\mathbf1\mathbf1^T/(d+1))$. It follows that all squared edge lengths equal $2c>0$. The columns are affinely independent since their sole linear dependence is a multiple of $\mathbf1$; their barycenter and circumcenter are $O$.

Conversely, a regular simplex centered at $O$ has common column norm and Gram matrix of the exhibited form (constant diagonal and off-diagonal, with zero row sum). It has $A\mathbf1=0$ and $AA^T=cI_d$. The displayed expansion then gives a constant. The all-center case is immediate.
\end{proof}
In particular $U_2$ is automatically constant for every quartic solution. If the $d+1$ sites are understood as distinct points, the all-center exception is impossible. If a list with repetitions is allowed, omitting this exception would be incorrect. Setting $d=2$ treats the already known three-site planar case; no assertion about an arbitrary number of planar sites is made.

\begin{thebibliography}{9}
\bibitem{r33keio:nr} N. Nikolov and R. Rafailov, \emph{On extremums of sums of powered distances to a finite set of points}, Geom. Dedicata \textbf{167} (2013), 69--89. \url{https://doi.org/10.1007/s10711-012-9804-3}. Conjecture numbers in this note refer to the accessible author version, \url{https://arxiv.org/pdf/1211.2975v1}.
\end{thebibliography}

\par\endgroup

% Source: /workspace/scratch/81914cec2c78/work/archive/Tournament_Round_34_Final/Tournament_Round_34/round34/meiji/paper.tex
% Edited source: /workspace/scratch/81914cec2c78/work/audit_34_41/cleaned/r34_meiji.tex
\begingroup

\Needspace{9\baselineskip}
\section{\texorpdfstring{Two fourth-degree triangle inequalities for medians and angle bisectors}{Two fourth-degree triangle inequalities for medians and angle bisectors}}\label{section:r34meiji}
\noindent{\small\sffamily\color{muted}Round 34 \enspace / \enspace Meiji}\par\medskip


\begin{sectionabstract}
Let $m_a$ and $w_a$ denote the median and internal bisector at vertex $A$, and let $s$ be the semiperimeter. We prove two conjectures of Jian Liu. The product $(\sum w_a)(\sum m_a^3)$ is at least $s^4$ in every triangle by a known pointwise comparison. For the other product, $(\sum m_a)(\sum w_a^3)$, we prove the same lower bound on the entire nonobtuse domain. The latter argument combines a local Hessian bound at the equilateral triangle with a complete rational interval certificate on the remaining two-parameter domain.
\end{sectionabstract}

\begin{theorem}
For any nondegenerate triangle,
\begin{equation}\label{r34meiji:eq1}
 \left(\sum w_a\right)\left(\sum m_a^3\right)\ge s^4 .
\end{equation}
For any nondegenerate nonobtuse triangle,
\begin{equation}\label{r34meiji:eq2}
 \left(\sum m_a\right)\left(\sum w_a^3\right)\ge s^4 .
\end{equation}
In each case equality holds exactly for the equilateral triangle. The sums extend over the three vertices.
\end{theorem}
These are \cite[p.~78, Conjectures 5.1 and 5.2]{r34meiji:liu}, with the second extended to right triangles. The source proves a different inequality where altitudes occur in place of bisectors in the linear factor. Its Section~5 also proposes a third conjecture; we do not address it here.

\paragraph{Proof of \eqref{r34meiji:eq1}.}
Put $X=s(s-a)$, $D=(b-c)^2$. The standard median and bisector formulas give
\[
 m_a^2=X+D/4,\qquad w_a^2=X(1-D/(b+c)^2).
\]
In particular $m_aw_a\ge X$, the known pointwise lemma of Liu \cite[Lemma 6, Eq.~(27)]{r34meiji:gaz}:
\[
 (m_aw_a)^2-X^2
   =\frac{XD}{4(b+c)^2}\bigl(a^2-D\bigr)\ge0 .
\]
By Cauchy--Schwarz,
\[
 \left(\sum w_a\right)\left(\sum m_a^3\right)
 \ge\left(\sum m_a\sqrt{w_am_a}\right)^2
 \ge\left(\sum s(s-a)\right)^2=s^4 .
\]
Any nonzero $D$ makes the corresponding $m_a^2>X$ strictly, and equality in the equilateral case is immediate. Thus this first conjecture is a short consequence of the pre-existing local bound.

\paragraph{Proof of \eqref{r34meiji:eq2}: reduction and local bound.}
By symmetry and homogeneity take $a\le b\le c$, $b=1$, and put
\[
 a=1-x,\qquad c=1+z,\qquad
 F(x,z)=\left(\sum m_a\right)\left(\sum w_a^3\right)-s^4 .
\]
The nonobtuse domain lies in
\[
 0\le x\le1,\quad 0\le z\le1/2,\quad 2z+z^2\le(1-x)^2
\]
because $c^2\le a^2+1\le 2$. We use the following expressions for every cyclic triple $(i,j,k)$ of side lengths:
\[
 m_i=\tfrac12\sqrt{2j^2+2k^2-i^2},\qquad
 w_i^3=\left(\frac{jk(j+k-i)(j+k+i)}{(j+k)^2}\right)^{3/2}.
\]
Outside the triangle domain, whenever necessary, interpret the last operation as $v\mapsto(\max(0,v))^{3/2}$. It is continuously differentiable on $\mathbb R$ with derivative $3\sqrt{\max(0,v)}/2$. No box with a nonpositive lower bound on a median radicand is accepted by derivative evaluation.

Put $K=[0,2^{-10}]^2$. In a neighborhood of $K$ the ingredients are smooth. The interval calculation below gives
\[
 1024F_{xx}\ge582,\qquad 1024F_{xz}\ge192,\qquad
 1024F_{zz}\ge574\qquad ((x,z)\in K).
\]
At the origin, $F(0,0)=0$ and its gradient vanishes: the original three-variable difference is symmetric and homogeneous of degree four, so at $(1,1,1)$ Euler's identity makes its three equal partial derivatives zero. Along every nonconstant segment $t\mapsto (tx,tz)$ in $K$ the second derivative is strictly positive, because $x,z\ge0$. Integration twice proves strict positivity there except at $(0,0)$.

\paragraph{Proof of \eqref{r34meiji:eq2}: the finite certificate.}
Represent an interval as integer endpoints divided by $S=2^{64}$ and round each arithmetic operation outward. For a positive interval $[L,U]/S$, square root is enclosed by
$[\lfloor\sqrt{LS}\rfloor,\lfloor\sqrt{US}\rfloor+1]/S$. Starting with the displayed rectangle, discard a box as exterior only if the least possible $2z+z^2$ exceeds the largest possible $(1-x)^2$. Accept boxes within $K$ by the local proof. Otherwise try the natural value enclosure, and next the enclosure
\[
 F(c_1,c_2)+[F_x]([x]-c_1)+[F_z]([z]-c_2)
\]
at the box midpoint $c$. The latter follows by integrating along the segment from the midpoint; brackets bound the indicated derivatives throughout the whole box. A box is accepted by either enclosure only if its lower endpoint is strictly positive. If neither test succeeds, bisect the longer side and test \emph{both} children. This proves positivity on a continuum, not just at the midpoints.

The complete verifier is printed below. The same integer code bounds the three Hessian entries on $K$ by two-direction automatic differentiation; the mixed component of the product uses the ordinary product rule twice. For the scalar composition $g(v)$ it uses $g'(v)v_{ij}+g''(v)v_i v_j$. In the global test only first derivatives are needed. Python's \texttt{//} is floor division, so negating a floor of the negative gives ceiling division. All certificate arithmetic is integral.
\[
\begin{array}{lr}
\text{Boxes accepted by the centered bound}&2241\\
\text{Exterior boxes}&36\\
\text{Local boxes}&1\\
\text{Greatest subdivision depth}&27
\end{array}
\]
These exhaust the initial rectangle. Together with the argument at the origin this proves \eqref{r34meiji:eq2} and its equality statement. \hfill$\square$

\begin{thebibliography}{9}
\bibitem{r34meiji:liu} J. Liu, \emph{Improvements of two inequalities in an acute triangle}, International Journal of Geometry 13 (2024), no.~4, 63--79. \href{https://ijgeometry.com/wp-content/uploads/2024/09/6.-63-79.pdf}{Publisher PDF}.
\bibitem{r34meiji:gaz} J. Liu, \emph{A class of inequalities involving length elements of a triangle}, Gazeta Matematic\u{a}, Seria A, XLI (2023), nos.~1--2, 1--17. \href{https://ssmr.ro/gazeta/gma/2023/gma1-2-2023-continut.pdf}{Publisher PDF}.
\end{thebibliography}
\subsection*{Appendix: complete verifier (Python 3)}
\begingroup\scriptsize
\begin{Verbatim}
from math import isqrt
from collections import Counter
S=1<<64
class I:
 __slots__='l','h'
 def __init__(self,l,h=None):
  if h is None:self.l=self.h=l*S
  else:self.l=l;self.h=h
 def __add__(x,y):
  if not isinstance(y,I):y=I(y)
  return I(x.l+y.l,x.h+y.h)
 __radd__=__add__
 def __neg__(x):return I(-x.h,-x.l)
 def __sub__(x,y):return x+-y
 def __rsub__(x,y):return -x+y
 def __mul__(x,y):
  if not isinstance(y,I):return I(min(x.l*y,x.h*y),max(x.l*y,x.h*y))
  if x.l>=0 and y.l>=0:return I(x.l*y.l//S,-(-x.h*y.h//S))
  p=(x.l*y.l,x.l*y.h,x.h*y.l,x.h*y.h);return I(min(p)//S,-(-max(p)//S))
 __rmul__=__mul__
 def __truediv__(x,y):
  if not isinstance(y,I):
   assert y>0;return I(x.l//y,-(-x.h//y))
  assert y.l>0, 'unsafe denominator'
  return x*I(S*S//y.h,-(-S*S//y.l))
 def rt(x):
  if x.l<=0:raise ValueError
  return I(isqrt(x.l*S),isqrt(x.h*S)+1)
 def pow32(x):
  y=I(max(x.l,0),max(x.h,0))
  return y*I(isqrt(y.l*S),isqrt(y.h*S)+1)
class J:
 __slots__='v','g'
 def __init__(self,v,g):
  self.v=v if isinstance(v,I) else I(v);self.g=[d if isinstance(d,I) else I(d) for d in g]
 def __add__(x,y):
  if not isinstance(y,J):return J(x.v+y,x.g)
  return J(x.v+y.v,[p+q for p,q in zip(x.g,y.g)])
 __radd__=__add__
 def __neg__(x):return J(-x.v,[-p for p in x.g])
 def __sub__(x,y):return x+-y
 def __rsub__(x,y):return -x+y
 def __mul__(x,y):
  if not isinstance(y,J):return J(x.v*y,[p*y for p in x.g])
  return J(x.v*y.v,[p*y.v+q*x.v for p,q in zip(x.g,y.g)])
 __rmul__=__mul__
 def __truediv__(x,y):
  if not isinstance(y,J):return J(x.v/y,[p/y for p in x.g])
  v=x.v/y.v;return J(v,[(p-v*q)/y.v for p,q in zip(x.g,y.g)])
 def rt(x):
  y=x.v.rt();return J(y,[p/(2*y) for p in x.g])
 def pow32(x):
  y=I(max(x.v.l,0),max(x.v.h,0));r=I(isqrt(y.l*S),isqrt(y.h*S)+1)
  return J(y*r,[3*r*p/2 for p in x.g])
def f(x,z):
 a,b,c=1-x,1,1+z; s=(a+b+c)/2
 A=[]
 for i,j,k in ((a,b,c),(b,c,a),(c,a,b)):
  A.append(((2*j*j+2*k*k-i*i).rt()/2,((j*k)*((j+k-i)*(j+k+i))/((j+k)*(j+k))).pow32()))
 return sum(p for p,q in A)*sum(q for p,q in A)-s*s*s*s
def prove(eps=1024):
 counts=Counter();stack=[(0,S,0,S//2,0)]
 while stack:
  l,h,u,v,d=stack.pop();counts['max depth']=max(d,counts['max depth'])
  if 2*u*S+u*u>(S-l)**2:counts['exterior']+=1;continue
  if max(h,v)<=S//eps:counts['local']+=1;continue
  left=I(l,h);right=I(u,v)
  try:
   if f(left,right).l>0:counts['value']+=1;continue
   midx=(l+h)//2;midz=(u+v)//2
   val=f(I(midx,midx),I(midz,midz))
   jet=f(J(left,[1,0]),J(right,[0,1]))
   if (val+jet.g[0]*I(l-midx,h-midx)+jet.g[1]*I(u-midz,v-midz)).l>0:counts['mean']+=1;continue
  except ValueError:pass
  assert d<50,(d,l/S,h/S,u/S,v/S,counts)
  if h-l>=v-u:
   w=(l+h)//2;stack.extend([(l,w,u,v,d+1),(w,h,u,v,d+1)])
  else:
   w=(u+v)//2;stack.extend([(l,h,u,w,d+1),(l,h,w,v,d+1)])
 print(counts)

class B:
 def __init__(self,v,g=0,h=0,j=0):
  self.v,self.g,self.h,self.j=[i if isinstance(i,I) else I(i) for i in (v,g,h,j)]
 def __add__(x,y):
  if not isinstance(y,B):y=B(y)
  return B(x.v+y.v,x.g+y.g,x.h+y.h,x.j+y.j)
 __radd__=__add__
 def __neg__(x):return B(-x.v,-x.g,-x.h,-x.j)
 def __sub__(x,y):return x+-y
 def __rsub__(x,y):return -x+y
 def __mul__(x,y):
  if not isinstance(y,B):y=B(y)
  return B(x.v*y.v,x.g*y.v+x.v*y.g,x.h*y.v+x.v*y.h,x.j*y.v+x.g*y.h+x.h*y.g+x.v*y.j)
 __rmul__=__mul__
 def chain(x,y,a,b):return B(y,a*x.g,a*x.h,a*x.j+b*x.g*x.h)
 def __truediv__(x,y):
  if not isinstance(y,B):return B(x.v/y,x.g/y,x.h/y,x.j/y)
  r=I(1)/y.v;return x*y.chain(r,-r*r,2*r*r*r)
 def rt(x):
  r=x.v.rt();d=I(1)/(2*r);return x.chain(r,d,-d/(2*x.v))
 def pow32(x):
  r=x.v.rt();return x.chain(x.v*r,3*r/2,I(3)/(4*r))
def local(eps):
 box=I(0,S//eps)
 for p,q in [(0,0),(0,1),(1,1)]:
  j=f(B(box,p==0,q==0),B(box,p==1,q==1)).j
  print('local Hessian lower > numerator/1024:',p,q,j.l*1024//S)
  assert j.l>0
if __name__=='__main__':
 local(1024)
 prove()
\end{Verbatim}

\endgroup

\par\endgroup

% Source: /workspace/scratch/81914cec2c78/work/archive/Tournament_Round_34_Final/Tournament_Round_34/round34/taisho/paper.tex
% Edited source: /workspace/scratch/81914cec2c78/work/audit_34_41/cleaned/r34_taisho.tex
\begingroup

\Needspace{9\baselineskip}
\section{\texorpdfstring{A hyperbolic arctangent bound}{A hyperbolic arctangent bound}}\label{section:r34taisho}
\noindent{\small\sffamily\color{muted}Round 34 \enspace / \enspace Taisho}\par\medskip


\subsection*{A hyperbolic counterpart of the arctangent bound}
For \(0<x<\sqrt{105/109}=\sqrt{11445}/109\), one has
\[
 \operatorname{arctanh}x < U(x):=\frac{53x}{18+35(1-109x^2/105)^{53/109}} .
\]
This is the hyperbolic-tangent inequality of Xu, Sun and Chen,
\href{https://doi.org/10.2298/AADM250324010X}{AADM 20 (2026), 81--102},
Conjecture 3, Eq.~(35). The same paper poses a different inequality for
the ordinary arctangent in its Conjecture 2. Crucially, the polynomial
used below is \emph{identical} to Meiji's earlier Round 21 polynomial for that
ordinary arctangent. The complete root count and arithmetic certificate needed here
are included below.

\textit{Proof.} Put \(t=(1-109x^2/105)^{1/109}\in(0,1)\),
\(D=18+35t^{53}\), and \(E=U-\operatorname{arctanh}\).
Since \(1-x^2=(4+105t^{109})/109\) and
\(xD'=3710(t^{53}-t^{-56})/109\), differentiation gives
\[
 E'=\frac{P(t)}{109t^{56}(4+105t^{109})D^2},\qquad
 P=53(1962t^{56}+105t^{109}+3710)(4+105t^{109})-109^2t^{56}D^2 .
\]
Writing \(P=\sum_j a_jt^j\), its nonzero terms are:
\[
\begin{array}{c|rrrrrr}
j&0&56&109&162&165&218\\\hline
a_j &786520&-3433500&5698350&-14554225&10918530&584325 .
\end{array}
\]
The same Descartes argument is included here: there are four sign changes,
and direct differentiation of the six displayed terms gives
\[
 P^{(j)}(1)=0\quad(0\le j\le3),\qquad
 P^{(4)}(1)=158952396481200.
\]
Thus the root at \(t=1\) has multiplicity four, and no other positive
root exists. Since \(P(0)>0\), this proves \(E'>0\) on the desired interval.
For an alternative exact arithmetic certificate set \(t=1/(1+y)\). In
\((1+y)^{218}P(1/(1+y))\) the coefficient of \(y^k\) is
\(c_k=\sum_j a_j\binom{218-j}{k}\), with out-of-range binomials zero.
The sums give \(c_0=\cdots=c_3=0,\ c_4=6623016520050\), and
\(\min_{4\le k\le218}c_k=786520\) (all 215 checks are reproduced by the
inline certificate below). The transformed polynomial is positive for \(y>0\),
and also certifies the asserted multiplicity at \(t=1\).
\begin{Verbatim}
"""Reconstruct the derivative numerator and verify its positivity certificate."""
from math import comb
def mul(p,q):
    z={}
    for a,b in p.items():
        for c,d in q.items(): z[a+c]=z.get(a+c,0)+b*d
    return z
d=mul({56:1962,109:105,0:3710},{109:105,0:4})
e=mul({0:18,53:35},{0:18,53:35})
p={j:53*n for j,n in d.items()}
for j,n in e.items(): p[j+56]=p.get(j+56,0)-109**2*n
print('P(t) ascending:',sorted(p.items()))
C=[sum(v*comb(218-j,k) for j,v in p.items() if k<=218-j) for k in range(219)]
assert C[:4]==[0]*4 and C[4]==6623016520050
assert min(C[4:])==786520
print('shifted coefficients c0..c4:',C[:5])
print('smallest of c4..c218:',min(C[4:]),'; checked',len(C[4:]),'strict inequalities')
\end{Verbatim}


Continuity and \(E(0)=0\) complete the proof. \hfill\(\square\)

\textbf{Status.} The R21 manuscript explicitly stated only the ordinary
arctangent conclusion, but proved that no positive root other than 1
exists. Thus the present result is a direct companion consequence, not
an independent new tournament closure.


\par\endgroup

% Source: /workspace/scratch/81914cec2c78/work/archive/Tournament_Round_34_Final/Tournament_Round_34/round34/showa/paper.tex
% Edited source: /workspace/scratch/81914cec2c78/work/audit_34_41/cleaned/r34_showa.tex
\begingroup

\Needspace{9\baselineskip}
\section{\texorpdfstring{The extremal root in every Gaussian Mills-ratio Tur\'an bound}{The extremal root in every Gaussian Mills-ratio Tur\'an bound}}\label{section:r34showa}
\noindent{\small\sffamily\color{muted}Round 34 \enspace / \enspace Showa}\par\medskip

\begin{sectionabstract}
In Remark 3.5 of his 2020 paper, Steven G. From asked which positive root of the uncorrected derivative Tur\'an quadratic bounds the Gaussian Mills ratio at each derivative order. We prove the predicted rule for every order and every nonnegative argument, including the points where the quadratic becomes linear. The positive moment-ratio branch gives the bound. To order it relative to the second branch, we use a Gaussian measure for one coefficient sequence and a signed, separately normalized two-half-line Gaussian measure for the other. The signed measure becomes positive after multiplication by an odd power. This argument is about root selection, not the first construction of alternating bounds for the Mills ratio.
\end{sectionabstract}

\subsection{The precise root question}
Let $r(t)=(1-\Phi(t))/\phi(t)$, where $\Phi,\phi$ are the standard normal distribution function and density. For $j\ge0$ define polynomials $P_j,Q_j$ by $P_0=1$, $P_1=t$, $Q_0=0$, $Q_1=1$, and, in both cases,
\[
 U_{j+2}=tU_{j+1}+(j+1)U_j .
\]
For a fixed $t\ge0$ put $S_j(s)=(-1)^j(P_js-Q_j)$. The uncorrected polynomial in From's Remark 3.5 \cite{r34showa:from}, equations (3.1)--(3.2) with $R_L^*=0$, is
\[
 F_L(s)=S_L(s)S_{L+2}(s)-S_{L+1}(s)^2 .
\]
Its value at $r(t)$ is $r^{(L)}(t)r^{(L+2)}(t)-(r^{(L+1)}(t))^2$. From gave root-selection proofs for the first seven orders. A root formula need not by itself identify the \emph{largest or smallest positive} root: at some arguments there are two positive roots, at others only one.
\begin{theorem}
For every integer $L\ge0$ and every $t\ge0$, put
\[
 z_+=\frac{-t+\sqrt{t^2+4(L+1)}}2,\qquad
 b_L(t)=\frac{Q_{L+1}+z_+Q_L}{P_{L+1}+z_+P_L}.
\]
If $L$ is even, $b_L(t)$ is the largest positive real zero of $F_L$ and
$b_L(t)<r(t)$. If $L$ is odd, $b_L(t)$ is the smallest positive real zero
and $r(t)<b_L(t)$. If $F_L$ is linear, the assertion concerns its unique finite root.
\end{theorem}

\subsection{The branch giving the bound}
Define $M_j(t)=\int_0^\infty u^j e^{-tu-u^2/2}\,du$. Then $M_0=r$, $M_1=1-tr$ and
$M_{j+2}=(j+1)M_j-tM_{j+1}$ by integration by parts. Consequently $M_j=S_j(r)=(-1)^jr^{(j)}(t)$.
Cauchy--Schwarz for $u^L e^{-tu-u^2/2}du$ is strict, so
\[
y=M_{L+1}/M_L>0,\qquad y^2+ty<L+1.
\]
The quadratic in $y$ has roots $z_+$ and $z_-=-t-z_+$. Solving $y=S_{L+1}(s)/S_L(s)$ for $s$ gives
\[
 m(y)=\frac{Q_{L+1}+yQ_L}{P_{L+1}+yP_L},\qquad r=m(y).
\]
The continuant determinant satisfies
\[
 W_L=Q_LP_{L+1}-Q_{L+1}P_L=(-1)^{L+1}L!,
\]
by $W_0=-1$ and $W_{L+1}=-(L+1)W_L$. The polynomials have nonnegative coefficients, and successive $P$'s (or $Q$'s) cannot both vanish at a fixed $t$ by $W_L\ne0$. Thus $m(v)>0$ for $v>0$, also at $t=0$. Its derivative has sign $(-1)^{L+1}$, since $m'(v)=W_L/(P_{L+1}+vP_L)^2$. As $0<y<z_+$, the strict bound follows.

For completeness, if $S_L(s)=0$, then $S_{L+1}(s)\ne0$ by $W_L\ne0$, so $F_L(s)<0$ and $s$ is not a root. At all other $s$,
\[
 \frac{F_L(s)}{S_L(s)^2}
 =(L+1)-t\frac{S_{L+1}(s)}{S_L(s)}
 -\left(\frac{S_{L+1}(s)}{S_L(s)}\right)^2.
\]
Thus the only projective roots are $m(z_+)$ and $m(z_-)$. They are distinct by $W_L\ne0$. The latter root is at infinity if and only if $P_{L+1}+z_-P_L=0$. In that case the numerator is nonzero, and, on clearing denominators in the last display, one of the two affine linear factors is a nonzero constant. The remaining factor is nonconstant because $P_{L+1}+z_+P_L>0$. Thus $F_L$ is linear there.

\subsection{The sign of the other branch}
Let $\mu=N(t,1)$, considered as a measure on the full line. Gaussian integration by parts and the first two moments give
$P_j=\int u^j\,d\mu$. There is a second representation
\[
 Q_j=\int u^j\,d\nu,\quad
 d\nu=\big(c_+\mathbf 1_{u>0}-c_-\mathbf 1_{u<0}\big)d\mu,\qquad c_+,c_->0.
\]
Indeed choose $c_+,c_-$ so that the two weighted half-line masses agree, then rescale them together to give first moment one. This is possible because $u\,d\nu$ is positive and not zero. Integration by parts on each half-line gives
$Q_{j+2}=tQ_{j+1}+(j+1)Q_j$ for $j\ge0$: the boundary term at zero includes $u^{j+1}$ and hence vanishes.

When $L$ is even, $u^L\,d\mu$ is a positive measure with interval support; when $L$ is odd, the same is true of $u^L\,d\nu$. Therefore with $U=P$ for even $L$ and $U=Q$ for odd $L$,
\[
 U_L>0,\qquad U_{L+1}^2<U_LU_{L+2}.
\]
Writing $w=U_{L+1}/U_L$ and using the recurrence yields $w^2-tw<L+1$. Its larger boundary root is $t+z_+$, and thus $U_{L+1}+z_-U_L<0$.

Write $D_\pm=P_{L+1}+z_\pm P_L$ and $N_\pm=Q_{L+1}+z_\pm Q_L$. We already have $D_+,N_+>0$. If $D_-\ne0$, then
\[
 m(z_+)-m(z_-)=\frac{(z_+-z_-)W_L}{D_+D_-}.
\]
For even $L$, $D_-<0$, and this difference is positive, so $m(z_+)$ is the larger root. For odd $L$, $N_-<0$. If $D_-<0$, both roots are positive and the difference is negative. If $D_->0$, the other root is negative; if $D_-=0$, it is infinite. In each case $m(z_+)$ is the smallest positive finite root. All arguments include $t=0$ without division by $t$. This completes the proof.

\paragraph{Scope and priority.}
The moment ratio and the alternating nature of related Gaussian continued-fraction bounds are not claimed as original. In particular the public Pudim Project record APP-0032 \cite{r34showa:pudim} describes an all-order alternating-bound family for From's Remark 6.5, not a proof of the ordering of all positive roots of the polynomial in Remark 3.5. The second-branch classification above is the point of the present result. No conclusion is drawn here for From's separate inequality with a nonzero exponential correction following Theorem 4.8, or for Remark 6.5 itself. A bounded literature check is not a formal priority certificate.
\begin{thebibliography}{9}
\bibitem{r34showa:from} S. G. From. Some new upper and lower bounds for the Mills ratio. \emph{J. Math. Anal. Appl.} 486 (2020), 123872. \url{https://doi.org/10.1016/j.jmaa.2020.123872}.
\bibitem{r34showa:pudim} Pudim Project. Application APP-0032: From's all-$L$ Mills-ratio bound family is explicit. Public application inventory, retrieved 19 Nov. 2026. \url{https://github.com/pudim-project/pudim-ai-demo-zetalaw/blob/main/APPLICATIONS.md}.
\end{thebibliography}

\par\endgroup

% Source: /workspace/scratch/81914cec2c78/work/archive/Tournament_Round_34_Final/Tournament_Round_34/round34/heisei/paper.tex
% Edited source: /workspace/scratch/81914cec2c78/work/audit_34_41/cleaned/r34_heisei.tex
\begingroup

\Needspace{9\baselineskip}
\section{\texorpdfstring{A Pad\'e-point sign certificate for From's fourth-difference bounds on the Gaussian Mills ratio}{A Pad\'e-point sign certificate for From's fourth-difference bounds on the Gaussian Mills ratio}}\label{section:r34heisei}
\noindent{\small\sffamily\color{muted}Round 34 \enspace / \enspace Heisei}\par\medskip


\begin{sectionabstract}
For the Gaussian Mills ratio, a fourth-difference inequality in the
derivatives gives a polynomial of degree at most two in the ratio itself.
We prove that for every shift the polynomial changes sign between the
Mills ratio and the corresponding Pad\'e approximant. This gives the
alternating root bounds proposed by From in Remark 6.5 of his 2020 paper.
An explicit root also covers the endpoint at which an odd Pad\'e approximant
becomes infinite. An earlier public solution is recorded under APP-0032;
the present proof does not establish historical priority.
\end{sectionabstract}
\subsection{The conjecture and the sign certificate}
For $t\ge0$, let
\[
 r(t)=e^{t^2/2}\int_t^\infty e^{-s^2/2}\,ds
     =\int_0^\infty e^{-tu-u^2/2}\,du.
\]
Fix $t$. For a variable $y$, define $m_0(y)=y$, $m_1(y)=1-ty$ and
\begin{equation}
m_{k+2}(y)=(k+1)m_k(y)-t m_{k+1}(y)\qquad(k\ge0).
\label{r34heisei:rec}
\end{equation}
Put
\[
Q_L(y)=m_L(y)m_{L+4}(y)-4m_{L+1}(y)m_{L+3}(y)+3m_{L+2}(y)^2 .
\]
At $y=r(t)$, integration by parts and differentiation under the integral
sign give $m_k(r(t))=(-1)^kr^{(k)}(t)
=\int_0^\infty u^ke^{-tu-u^2/2}\,du>0$.
Gaussian decay dominates all these differentiated integrands locally in $t$.
Thus $Q_L$ is exactly the polynomial of From's Remark 6.5.
\begin{theorem}
For every integer $L\ge0$ and every $t\ge0$, the equation $Q_L(y)=0$
has a positive root $y_*$ such that $y_*<r(t)$ if $L$ is even and
$y_*>r(t)$ if $L$ is odd.
\end{theorem}
\begin{proof}
Write $m_k(y)=(-1)^k(p_k y-s_k)$ with $p_0=1$, $p_1=t$, $s_0=0$,
$s_1=1$. Both coefficient sequences satisfy
$a_{k+2}=t a_{k+1}+(k+1)a_k$ and are nonnegative. At $y=r(t)$, expansion
of $(u-v)^4$ and symmetry give
\[
 Q_L(r(t))=\tfrac12\iint_{(0,\infty)^2}
 (uv)^L(u-v)^4e^{-t(u+v)-(u^2+v^2)/2}\,du\,dv>0 .
\]
For $p_L>0$, let $h=s_L/p_L$. Since $m_L(r(t))>0$, we have
$0\le h<r(t)$ for even $L$, and $h>r(t)>0$ for odd $L$.
At $h$, $m_L=0$ but $m_{L+1}\ne0$: the determinant
$p_Ls_{L+1}-p_{L+1}s_L=(-1)^L L!$ follows from the common recurrence.
Taking $V=m_{L+1}(h)$ in \eqref{r34heisei:rec} gives
$m_{L+2}(h)=-tV$ and $m_{L+3}(h)=(L+2+t^2)V$. Consequently
\[
 Q_L(h)=-(4L+8+t^2)V^2<0 .
\]
A root lies strictly between $h$ and $r(t)$, proving the claim here.

The only exceptional case is $t=0$, $L$ odd: indeed for $t>0$ every
$p_k$ is positive, and for $t=0$ exactly the odd $p_k$ vanish. In this
case $U=m_L(y)$ is a positive constant, and $V=m_{L+1}(y)$ has nonzero
slope. From \eqref{r34heisei:rec},
\[
Q_L(y)=(L+1)(4L+6)U^2-(4L+8)V^2.
\]
It tends to $-\infty$ as $y\to+\infty$, whereas $Q_L(r(0))>0$.
There is therefore a root larger than $r(0)$.
\end{proof}
\subsection{An explicit positive root}
One can identify the root without comparing the sizes of the two roots
of $Q_L$. Set $U=m_L(y)$ and $V=m_{L+1}(y)$. Three applications of
\eqref{r34heisei:rec}, including $m_{L+4}=(L+3)m_{L+2}-t m_{L+3}$, yield
\[
 Q_L(y)=A U^2-BUV-CV^2,\quad
 \begin{cases}
 A=(L+1)(4L+6+t^2),\\
 B=t(4L+7+t^2),\\
 C=4L+8+t^2.
 \end{cases}
\]
Let $b=2A/(B+\sqrt{B^2+4AC})$, the positive zero of $A-Bz-Cz^2$.
At $y=r(t)$ we have $U,V>0$ and $Q_L>0$, hence $V/U<b$. The equation
$V=bU$ defines the positive root
\[
 y_* = \frac{s_{L+1}+b s_L}{p_{L+1}+b p_L}.
\]
Both sums are strictly positive since their entries are nonnegative and
the adjacent determinant never vanishes. The sign of
$V-bU=(-1)^{L+1}(p_{L+1}+b p_L)(r(t)-y_*)$ at $r(t)$ yields the
claimed strict inequalities directly. This formula selects the bound
even when $Q_L$ itself degenerates to a linear polynomial.
\begin{thebibliography}{9}
\bibitem{r34heisei:F} S.~G. From, Some new upper and lower bounds for the Mills ratio,
\emph{J. Math. Anal. Appl.} \textbf{486} (2020), 123872. Eq.~(6.8) and
Remark~6.5. DOI: 10.1016/j.jmaa.2020.123872.
\bibitem{r34heisei:P} Pudim Project, \emph{Public research application register},
APP-0032, June 3, 2026, on the all-$L$ Mills-ratio bound family.
\url{https://github.com/pudim-project/pudim-ai-demo-zetalaw}
(see APPLICATIONS.md). The register is public prior art, not evidence of
peer review.
\end{thebibliography}

\par\endgroup

% Source: /workspace/scratch/81914cec2c78/work/archive/Tournament_Round_35_Final/Tournament_Round_35/round35/meiji/paper.tex
% Edited source: /workspace/scratch/81914cec2c78/work/audit_34_41/cleaned/r35_meiji.tex
\begingroup

\Needspace{9\baselineskip}
\section{\texorpdfstring{A fifth-power inequality for medians, altitudes and exradii}{A fifth-power inequality for medians, altitudes and exradii}}\label{section:r35meiji}
\noindent{\small\sffamily\color{muted}Round 35 \enspace / \enspace Meiji}\par\medskip


\begin{sectionabstract}
For a nonobtuse triangle, the sum of its medians bounds the fifth power mean of its altitude sum and exradius sum. This proves Conjecture 5.3 of Jian Liu, including an extension to right triangles. A change of variables cancels the apparent singularity in the exradius sum as one side tends to zero. This gives a compact parameter domain on which the proof reduces to a local second-derivative bound at the unique equality case and an exact subdivision certificate elsewhere.
\end{sectionabstract}

\begin{theorem}
Let $m_a$, $h_a$ and $r_a$ be the median, altitude and exradius associated with side $a$, and define the corresponding quantities cyclically. On every nondegenerate nonobtuse triangle,
\[
 2\left(\sum m_a\right)^5\ge
 \left(\sum h_a\right)^5+\left(\sum r_a\right)^5 .
\]
Equality holds if and only if the triangle is equilateral.
\end{theorem}
The original open statement is Liu \cite[p.~78, Conjecture 5.3, Eq.~(5.6)]{r35meiji:liu}, for acute triangles. Eq.~(5.5) in the same paper treats only the arithmetic mean. Here the fifth powers are taken \emph{after} the respective sums.

\paragraph{Eliminating the small-side singularity.}
Scale and order the sides as $0<a\le b=1\le c$. Then $0\le c-1\le a^2/(c+1)$ by nonobtuseness. This suggests the bounded variables
\[
 x=1-a,\qquad y=(c-1)/a^2,\qquad
 0\le x\le1,\quad 0\le y\le\tfrac12,\quad
 2y+(1-x)^2y^2\le1 .
\]
Every desired triangle maps into this set; the limiting edge $x=1$ is harmless. Write $s=(a+1+c)/2$ and set
\[
 U=\sqrt{((1+c)^2-a^2)(1-a^2y^2)} .
\]
By the symmetric area identity $16\Delta^2=((1+c)^2-a^2)(a^2-(c-1)^2)$, the area is $\Delta=aU/4$ and $U$ stays positive up to $a=0$. The altitude sum $H$ and exradius sum $T$ therefore admit the formulas
\[
 H=\frac{U}{2}(1+a+a/c),\qquad
 T=\frac{4c}{U}+\frac{aU}{4s}.
\]
In fact $H=2\Delta(1/a+1+1/c)$ and $T=4R+r$, using circumradius $R=ac/(4\Delta)$ and inradius $r=\Delta/s$. For the latter identity one may put $p=s-a,q=s-b,t=s-c$. Heron's identity gives
$\Delta^2=spqt$ and
\[
 4R+r=\frac{(p+q)(q+t)(t+p)+pqt}{\Delta}
      =\frac{s(pq+qt+tp)}{\Delta}
      =\Delta(1/p+1/q+1/t)=\sum r_a .
\]
Obtain the median sum $M$ from $2m_i=\sqrt{2j^2+2k^2-i^2}$ for each cyclic triple $(i,j,k)$ of sides. The difference to be bounded is the regular expression
\[
 F(x,y)=2M^5-H^5-T^5,\qquad a=1-x,\quad c=1+a^2y .
\]

\paragraph{The equilateral neighborhood.}
At $(x,y)=(0,0)$, $F=0$ and its first derivatives vanish. Indeed the original three-side expression is symmetric, smooth near the equilateral triangle and homogeneous of degree five; at $(1,1,1)$ its three equal partial derivatives sum to $5F=0$, and the pullback therefore also has zero gradient. On the box $K=[0,2^{-13}]^2$ the integer verifier in the appendix proves
\[
 1024F_{xx}\ge266207,\qquad
 1024F_{xy}\ge129594,\qquad
 1024F_{yy}\ge264019 .
\]
Along every nonconstant ray in the nonnegative quadrant the second directional derivative is positive throughout this box. Integration twice from the origin proves strict positivity on $K\setminus\{(0,0)\}$.

\paragraph{The remaining continuum of triangles.}
The calculation uses closed intervals with integer endpoints in units of $1/S$, $S=2^{64}$. Every primitive interval operation rounds outwards. For example, for positive $[L,U]/S$ the respective square-root endpoints in these units are $\lfloor\sqrt{LS}\rfloor$ and $1+\lfloor\sqrt{US}\rfloor$. All divisions have certified positive denominators. The local second derivatives are evaluated by mixed-direction automatic differentiation using the ordinary product and chain rules; only gradients are required outside $K$.

Start with the full rectangle $[0,1]\times[0,1/2]$. On a box with lower $y$-endpoint $u$ and upper $x$-endpoint $h$, the triangle restriction is impossible if $2u+(1-h)^2u^2>1$. Discard precisely such boxes as exterior. Boxes within $K$ are already proved. For another box $B$, accept a strictly positive lower bound from direct evaluation or from the centered enclosure
\[
 F(m)+[F_x](B_x-m_x)+[F_y](B_y-m_y),
\]
where $m$ is the box midpoint and each derivative is enclosed over all of $B$. Integrating the gradient along the segment from $m$ proves this bound. An inconclusive enclosure, including an unverified positive radicand, leads to bisection along the longer side. Both children are checked; there is no sampling step.

The complete certificate replay yields $1452$ boxes accepted by the centered bound, $3$ by the direct bound, $221$ excluded as exterior and $1$ in $K$, with maximum depth $31$. These cover the original rectangle, and prove the required sign on its feasible part. As equality outside the origin was excluded, this completes the theorem. \hfill$\square$

\begin{thebibliography}{1}
\bibitem{r35meiji:liu} J. Liu, \emph{Improvements of two inequalities in an acute triangle}, International Journal of Geometry 13 (2024), no.~4, 63--79, \href{https://ijgeometry.com/wp-content/uploads/2024/09/6.-63-79.pdf}{publisher PDF}.
\end{thebibliography}

\subsection*{Appendix: full integer verifier (Python 3)}
\begingroup\scriptsize
\begin{Verbatim}
from math import isqrt
from collections import Counter
S=1<<64
class I:
 __slots__='l','h'
 def __init__(self,l,h=None):
  if h is None:self.l=self.h=l*S
  else:self.l=l;self.h=h
 def __add__(x,y):
  if not isinstance(y,I):y=I(y)
  return I(x.l+y.l,x.h+y.h)
 __radd__=__add__
 def __neg__(x):return I(-x.h,-x.l)
 def __sub__(x,y):return x+-y
 def __rsub__(x,y):return -x+y
 def __mul__(x,y):
  if not isinstance(y,I):return I(min(x.l*y,x.h*y),max(x.l*y,x.h*y))
  if x.l>=0 and y.l>=0:return I(x.l*y.l//S,-(-x.h*y.h//S))
  p=(x.l*y.l,x.l*y.h,x.h*y.l,x.h*y.h);return I(min(p)//S,-(-max(p)//S))
 __rmul__=__mul__
 def __truediv__(x,y):
  if not isinstance(y,I):
   assert y>0;return I(x.l//y,-(-x.h//y))
  assert y.l>0, 'unsafe denominator'
  return x*I(S*S//y.h,-(-S*S//y.l))
 def rt(x):
  if x.l<=0:raise ValueError
  return I(isqrt(x.l*S),isqrt(x.h*S)+1)
 def pow32(x):
  y=I(max(x.l,0),max(x.h,0))
  return y*I(isqrt(y.l*S),isqrt(y.h*S)+1)
class J:
 __slots__='v','g'
 def __init__(self,v,g):
  self.v=v if isinstance(v,I) else I(v);self.g=[d if isinstance(d,I) else I(d) for d in g]
 def __add__(x,y):
  if not isinstance(y,J):return J(x.v+y,x.g)
  return J(x.v+y.v,[p+q for p,q in zip(x.g,y.g)])
 __radd__=__add__
 def __neg__(x):return J(-x.v,[-p for p in x.g])
 def __sub__(x,y):return x+-y
 def __rsub__(x,y):return -x+y
 def __mul__(x,y):
  if not isinstance(y,J):return J(x.v*y,[p*y for p in x.g])
  return J(x.v*y.v,[p*y.v+q*x.v for p,q in zip(x.g,y.g)])
 __rmul__=__mul__
 def __truediv__(x,y):
  if not isinstance(y,J):return J(x.v/y,[p/y for p in x.g])
  v=x.v/y.v;return J(v,[(p-v*q)/y.v for p,q in zip(x.g,y.g)])
 def rt(x):
  y=x.v.rt();return J(y,[p/(2*y) for p in x.g])
 def pow32(x):
  y=I(max(x.v.l,0),max(x.v.h,0));r=I(isqrt(y.l*S),isqrt(y.h*S)+1)
  return J(y*r,[3*r*p/2 for p in x.g])
def f(x,y):
 a=1-x;b=1;c=1+a*a*y;s=(a+b+c)/2
 U=(((1+c)*(1+c)-a*a)*(1-a*a*y*y)).rt()
 H=(U/2)*(1+a+a/c)
 R=(4*c)/U + a*U/(4*s)
 M=sum((2*j*j+2*k*k-i*i).rt()/2 for i,j,k in ((a,b,c),(b,c,a),(c,a,b)))
 def p5(v):return v*v*v*v*v
 return 2*p5(M)-p5(H)-p5(R)
def prove(eps=8192):
 counts=Counter();stack=[(0,S,0,S//2,0)]
 while stack:
  l,h,u,v,d=stack.pop();counts['max depth']=max(d,counts['max depth'])
  if 2*u*S**3+(S-h)**2*u*u>S**4:counts['exterior']+=1;continue
  if max(h,v)<=S//eps:counts['local']+=1;continue
  left=I(l,h);right=I(u,v)
  try:
   if f(left,right).l>0:counts['value']+=1;continue
   midx=(l+h)//2;midz=(u+v)//2
   val=f(I(midx,midx),I(midz,midz))
   jet=f(J(left,[1,0]),J(right,[0,1]))
   if (val+jet.g[0]*I(l-midx,h-midx)+jet.g[1]*I(u-midz,v-midz)).l>0:counts['mean']+=1;continue
  except ValueError:pass
  assert d<50,(d,l/S,h/S,u/S,v/S,counts)
  if h-l>=v-u:
   w=(l+h)//2;stack.extend([(l,w,u,v,d+1),(w,h,u,v,d+1)])
  else:
   w=(u+v)//2;stack.extend([(l,h,u,w,d+1),(l,h,w,v,d+1)])
 print(counts)

class B:
 def __init__(self,v,g=0,h=0,j=0):
  self.v,self.g,self.h,self.j=[i if isinstance(i,I) else I(i) for i in (v,g,h,j)]
 def __add__(x,y):
  if not isinstance(y,B):y=B(y)
  return B(x.v+y.v,x.g+y.g,x.h+y.h,x.j+y.j)
 __radd__=__add__
 def __neg__(x):return B(-x.v,-x.g,-x.h,-x.j)
 def __sub__(x,y):return x+-y
 def __rsub__(x,y):return -x+y
 def __mul__(x,y):
  if not isinstance(y,B):y=B(y)
  return B(x.v*y.v,x.g*y.v+x.v*y.g,x.h*y.v+x.v*y.h,x.j*y.v+x.g*y.h+x.h*y.g+x.v*y.j)
 __rmul__=__mul__
 def chain(x,y,a,b):return B(y,a*x.g,a*x.h,a*x.j+b*x.g*x.h)
 def __truediv__(x,y):
  if not isinstance(y,B):return B(x.v/y,x.g/y,x.h/y,x.j/y)
  r=I(1)/y.v;return x*y.chain(r,-r*r,2*r*r*r)
 def rt(x):
  r=x.v.rt();d=I(1)/(2*r);return x.chain(r,d,-d/(2*x.v))
 def pow32(x):
  r=x.v.rt();return x.chain(x.v*r,3*r/2,I(3)/(4*r))
def local(eps):
 box=I(0,S//eps)
 for p,q in [(0,0),(0,1),(1,1)]:
  j=f(B(box,p==0,q==0),B(box,p==1,q==1)).j
  print('local Hessian lower > numerator/1024:',p,q,j.l*1024//S)
  assert j.l>0
if __name__=='__main__':
 local(8192)
 prove()
\end{Verbatim}

\endgroup

\par\endgroup

% Source: /workspace/scratch/81914cec2c78/work/archive/Tournament_Round_35_Final/Tournament_Round_35/round35/taisho/paper.tex
% Edited source: /workspace/scratch/81914cec2c78/work/audit_34_41/cleaned/r35_taisho.tex
\begingroup

\Needspace{9\baselineskip}
\section{\texorpdfstring{A shifted digamma remainder in Chen's Vernescu problem}{A shifted digamma remainder in Chen's Vernescu problem}}\label{section:r35taisho}
\noindent{\small\sffamily\color{muted}Round 35 \enspace / \enspace Taisho}\par\medskip


\textbf{Theorem.} Let \(\psi=\Gamma'/\Gamma\) be the digamma function and put
\[
 C(x)=\log x-\psi(x+1)+\frac{1}{2x},\qquad H_s(x)=(x+s)^2C(x),\qquad x>0.
\]
For every \(s\ge\sqrt{6}/(2\pi)\), the function \(H_s\) is strictly completely
monotone: \((-1)^mH_s^{(m)}(x)>0\) for every integer \(m\ge0\) and every
\(x>0\). The case \(s=1\) answers Conjecture 2 of Chen [1] affirmatively.
We do not claim the stated shift threshold is optimal. Nor do we assert
priority for the conjecture: a published theorem for another remainder
already entails it by the duplication argument on the following page.

\textit{Proof.} Subtract the standard integral formulas for \(\log x\)
and \(\psi(x+1)\) to obtain
\[
 C(x)=\int_0^\infty e^{-xt}K(t)\,dt,\quad
 K(t)=\frac{1}{e^t-1}-\frac1t+\frac12
     =\frac12\coth(t/2)-\frac1t .
\]
The partial-fraction expansion of \(\coth\), equivalently the logarithmic
derivative of Euler's product for \(\sinh(t/2)/(t/2)\), gives
\[
 K(t)=2\sum_{n=1}^{\infty} f_{2\pi n}(t),\qquad f_a(t)=\frac{t}{t^2+a^2}.
\]
On compact \(t\)-intervals the rational series and its first two derivatives
converge uniformly, by bounds \(O(n^{-2})\). The following exact identity
controls the second derivative for all \(a,t>0\):
\[
 \frac{f_a''}{f_a}+\frac{6}{a^2}
 =\frac{2(t^2-3a^2)}{(t^2+a^2)^2}+\frac6{a^2}
 =\frac{2t^2(7a^2+3t^2)}{a^2(t^2+a^2)^2}>0.
\]
All poles have \(a\ge2\pi\) and \(f_a>0\), so
\(K''+6K/(2\pi)^2>0\). Also \(K>0\) and
\[
 K'(t)=\frac{1}{t^2}-\frac{1}{4\sinh^2(t/2)}>0
\]
because \(\sinh(t/2)>t/2\). For every shift in the theorem we conclude
\[
 W_s(t):=K''(t)+2sK'(t)+s^2K(t)>0\qquad (t>0).
\]
Expansion at the origin gives \(K(0)=0\), \(K'(0)=1/12\).
At infinity \(K,K',K''\) are bounded by their explicit hyperbolic formulas.
Integrating by parts twice thus yields
\[
 H_s(x)=\frac1{12}+\int_0^\infty e^{-xt}W_s(t)\,dt .
\]
The integral and all its \(x\)-derivatives converge locally uniformly for
\(x>0\). The positive density proves strictness for every order. \hfill\(\square\)


\subsection*{What follows from prior work?}
Chen's conjecture is for the Vernescu-associated remainder \(C\), not the
DeTemple remainder \(R(x)=\psi(x+1/2)-\log x\). Nevertheless, Theorem 1
of Yang, Chu and Zhang [2] implies the case \(s=1\). More generally, they
prove \((u+a)^2R(u)\) is strictly completely monotone for
\(a\ge a_0=0.4847\ldots\). To see the indirect implication,
use \(\psi(x+1)=\psi(x)+1/x\) and digamma duplication to get
\[
 C(x)-2C(2x)=\psi(x+1/2)-\log x=R(x).
\]
The expansion \(C(u)=1/(12u^2)+O(u^{-4})\) at infinity follows from
the kernel at zero or the standard expansion of \(\psi\).
Iteration and passage to the limit therefore give
\[
 H_s(x)=\sum_{j=0}^\infty 2^{-j}
     (2^j x+2^j s)^2 R(2^j x).
\]
For \(s\ge a_0\) each term, as a function of \(x\), is completely
monotone by [2]. Finite nonnegative sums and finite-valued pointwise
limits preserve complete monotonicity (equivalently by Bernstein's
representation theorem). The first summand is itself strictly completely
monotone, so the limiting sum is strict.

Thus the exact question in [1] can already be answered using a published
result. The partial-fraction proof is separate and yields a sufficient
interval below the one available from that dyadic argument, namely
\[
 \frac{\sqrt6}{2\pi}=0.3898484\ldots\ \le s< a_0.
\]
This comparison concerns the two displayed sufficient criteria, not
a claim of a best-known threshold or a new result about the arithmetic
nature of Euler's constant.

\begin{thebibliography}{9}
\bibitem{r35taisho:chen} C.-P. Chen, \emph{The best bounds in Vernescu's inequalities
for the Euler's constant}, RGMIA Res.\ Rep.\ Collection \textbf{12}
(2009), no.\ 3, article 11, Eq.~(10) and Conjecture~2.
\href{https://rgmia.org/papers/v12n3/Euler-inequality.pdf}{Author's report}.
\bibitem{r35taisho:yang} Z.-H. Yang, Y.-M. Chu, and X.-H. Zhang,
\emph{Necessary and sufficient conditions for functions involving the
psi function to be completely monotonic}, J.\ Inequal.\ Appl.\
\textbf{2015}, article 157, Theorem 1.
\href{https://doi.org/10.1186/s13660-015-0674-8}{doi:10.1186/s13660-015-0674-8}.
\end{thebibliography}

\par\endgroup

% Source: /workspace/scratch/81914cec2c78/work/archive/Tournament_Round_35_Final/Tournament_Round_35/round35/showa/paper.tex
% Edited source: /workspace/scratch/81914cec2c78/work/audit_34_41/cleaned/r35_showa.tex
\begingroup

\Needspace{9\baselineskip}
\section{\texorpdfstring{Extremal roots of the corrected Gaussian Mills-ratio quadratics}{Extremal roots of the corrected Gaussian Mills-ratio quadratics}}\label{section:r35showa}
\noindent{\small\sffamily\color{muted}Round 35 \enspace / \enspace Showa}\par\medskip

\begin{sectionabstract}
Steven G. From conjectured that, at every derivative order, the largest or smallest positive root of his exponentially corrected quadratics bounds the Gaussian Mills ratio in the appropriate direction. We prove this for all nonnegative arguments and all nonnegative integer orders. More generally, the same positive-root selection rule holds at every level between zero and the actual derivative Tur\'an determinant. For odd orders the key fact is that the value of the quadratic at zero strictly exceeds that determinant. We prove this by comparing moment matrices for the two reflected Gaussian half-lines and their convex mixture. The corrected equality holds at the endpoint, and the bounds are strict at every positive argument.
\end{sectionabstract}

\subsection{Statement and the correction}
Let $r(t)=(1-\Phi(t))/\phi(t)$ for the standard normal distribution function and density. For $t\ge0$, set
\[
 M_j(t)=\int_0^\infty u^j e^{-tu-u^2/2}\,du=(-1)^j r^{(j)}(t),\qquad
 D_L=M_L M_{L+2}-M_{L+1}^2 .
\]
In particular $M_0=r$. For $P_0=1$, $P_1=t$, $Q_0=0$, $Q_1=1$, use the recurrence
$U_{j+2}=tU_{j+1}+(j+1)U_j$ for both sequences. Writing $S_j(s)=(-1)^j(P_js-Q_j)$, define the polynomial in a trial value $s$,
\[
 F_L(s)=S_L(s)S_{L+2}(s)-S_{L+1}(s)^2 .
\]
Integration by parts on the integral defining $M_j$ gives $M_1=1-tr$ and
$M_{j+2}=(j+1)M_j-tM_{j+1}$. Thus $M_j=S_j(r)$ and $F_L(r)=D_L>0$, the latter inequality being strict Cauchy--Schwarz.

The conjecture following Theorem 4.8 of From \cite{r35showa:from} concerns the positive roots of
$F_L(s)=R^*_L(t)$ for the exact correction from equation (4.10):
\[
 R^*_L(t)=2^L(L+1)!V_L e^{-\sqrt2 U_L t},\quad
 U_L=\sqrt2\,\frac{\Gamma(L+5/2)}{\Gamma(L+2)},\quad
 V_L=\int_0^{\pi/2}(\cos\theta\sin\theta)^L
 (\cos\theta-\sin\theta)^2\,d\theta .
\]
This source conjecture is \emph{not} its Remark 3.5, where the right side is identically zero. Moving a parabola vertically can move a previously negative root through zero, changing which root is smallest positive. Our theorem rules that out in the relevant case.
\begin{theorem}
Fix $L\ge0$ an integer, $t\ge0$ and $0\le K\le D_L(t)$. For even $L$ the largest positive real root $s_K$ of $F_L(s)=K$ exists and satisfies $s_K\le r(t)$. For odd $L$ the smallest positive real root $s_K$ exists and satisfies $r(t)\le s_K$. In both cases equality holds exactly for $K=D_L(t)$. A linear polynomial is allowed.
\end{theorem}
\begin{corollary}
For From's correction $K=R^*_L(t)$, the parity-directed bounds hold for every order and are strict for $t>0$; equality holds at $t=0$.
\end{corollary}
Indeed, by symmetrizing the double integral for $D_L$, its weight is
$\tfrac12 (uv)^L(u-v)^2e^{-(u^2+v^2)/2}$ and the remaining factor is $e^{-t(u+v)}$ on the positive quadrant. In polar coordinates, the radial density after normalization is proportional to $z^{2L+3}e^{-z^2/2}$, with mean $U_L$; the total mass is $2^L(L+1)!V_L=D_L(0)$. As $u+v\le\sqrt2 z$, Jensen's inequality gives
$D_L(t)\ge D_L(0)\exp(-\sqrt2 U_Lt)=R^*_L(t)$. Strictness holds for $t>0$ by strict convexity and the nondegenerate radial law, and equality is immediate at $t=0$. This recovers From's needed Lemma 4.7(a).

\subsection{The odd-order intercept}
\begin{lemma}
If $L$ is odd, then $F_L(0)>D_L(t)$ for every $t\ge0$.
\end{lemma}
\begin{proof}
On the real line let $d\omega(x)=e^{tx-x^2/2}\,dx$, $a=\omega((-\infty,0))$, $b=\omega((0,\infty))$, and $\lambda=a/(a+b)$. Define
\[
 d\nu=\left(\lambda\mathbf 1_{x>0}-(1-\lambda)\mathbf 1_{x<0}\right)d\omega.
\]
This signed measure has total mass zero. Since $x\,d\omega=t\,d\omega-d(e^{tx-x^2/2})$, its first moment is one (the density jump factor at the origin is one). Integration by parts on each half-line shows
$\int x^j\,d\nu=Q_j$: for the recursion, the boundary term at the origin contains $x^{j+1}$ and vanishes. This is the signed moment representation used for the uncorrected problem in \cite{r35showa:r34}.

For $h(x)=(1,x)^{\mathsf T}$ let $H_+$ and $H_-$ be the positive definite matrices
$\int |x|^L h(x)h(x)^{\mathsf T}\,d\omega(x)$ over the positive and negative half-lines, respectively. Odd parity allows us to identify
\[
 \begin{pmatrix}Q_L&Q_{L+1}\\Q_{L+1}&Q_{L+2}\end{pmatrix}
 =\lambda H_+ +(1-\lambda)H_- .
\]
The determinant of this mixture equals $F_L(0)$. Reflecting the negative half-line gives $\det H_-=D_L(t)$, since the off-diagonal changes sign but the determinant does not. Moreover
\[
 \det H_{\pm}=\frac12\iint_{u,v>0}
 (uv)^L(u-v)^2 e^{\pm t(u+v)-(u^2+v^2)/2}\,du\,dv .
\]
Thus $\det H_+>\det H_-$ when $t>0$. For two positive definite $2\times2$ matrices,
\[
 \sqrt{\det(\lambda A+(1-\lambda)B)}\ge
 \lambda\sqrt{\det A}+(1-\lambda)\sqrt{\det B}.
\]
To see this, conjugate by $B^{-1/2}$ and apply the arithmetic-geometric mean inequality to the two positive eigenvalues of $B^{-1/2}AB^{-1/2}$ after squaring. Hence at $t>0$ the determinant of the mixture strictly exceeds $\det H_-$. At $t=0$, $\lambda=1/2$ and $H_+,H_-$ have equal diagonal entries and opposite nonzero off-diagonal entries. The mixture is diagonal and its determinant is again strictly larger. This proves the lemma.
\end{proof}

\subsection{Selection of the positive root}
The identity
\[
 Q_LP_{L+1}-Q_{L+1}P_L=(-1)^{L+1}L!
\]
follows from its value at $L=0$ and the common recurrence. When $P_L>0$, put $c=Q_L/P_L\ge0$. Then $S_L(c)=0$ but $S_{L+1}(c)\ne0$, so $F_L(c)<0$. The polynomials $P_j$ are the moments of the normal distribution $N(t,1)$ by Gaussian integration by parts.

For even $L$, $P_L>0$, and strict Cauchy--Schwarz for the positive measure $x^L dN(t,1)$ shows that $P_L P_{L+2}-P_{L+1}^2>0$, the leading coefficient of $F_L$. Also $r>c$ because $M_L=P_Lr-Q_L>0$. Since $F_L(c)<K$, the equation $F_L=K$ has a unique root to the right of $c$. It is positive, the larger of the two, and lies at or to the left of $r$ because $F_L(r)=D_L\ge K$.

For odd $L$ and $t>0$, $P_L>0$ (the recurrence has nonnegative coefficients and the leading monomial is $t^L$), and $0<r<c$ because $M_L=Q_L-P_Lr>0$. By the lemma $F_L(0)>D_L\ge K$ while $F_L(c)<K$. A polynomial of degree at most two with opposite strict signs at the endpoints has exactly one root in $(0,c)$: it must cross there, and two simple crossings would preserve the endpoint sign. This root is the smallest positive one. As $F_L(r)\ge K$, we have $r\le s_K$. The same argument applies when the leading coefficient vanishes.

Finally if $t=0$ and $L$ is odd, then $P_L=P_{L+2}=Q_{L+1}=0$ by parity, whereas $P_{L+1}>0$. Hence $F_L(s)=F_L(0)-P_{L+1}^2s^2$. The intercept lemma again shows that there is a unique positive root and that $r\le s_K$. In every case equality holds if and only if $F_L(r)=K$, i.e. $K=D_L(t)$. This proves the theorem.

\paragraph{Scope.} From treated the first seven orders in the corrected setting. His nesting observation explains how an existing branch moves; for arbitrary odd orders, the positive-root designation also requires control of the intercept. Round 34's uncorrected proof gave $F_L(0)>0$, not the all-level intercept comparison above. The public Pudim APP-0032 \cite{r35showa:pudim} targets the separate Remark 6.5. We do not claim novelty for the determinant representation of $D_L$, the correction, Jensen's bound, or all-order uncorrected continued-fraction estimates.
\begin{thebibliography}{9}
\bibitem{r35showa:from} S. G. From, Some new upper and lower bounds for the Mills ratio, \emph{J. Math. Anal. Appl.} 486 (2020), 123872. \url{https://doi.org/10.1016/j.jmaa.2020.123872}.
\bibitem{r35showa:r34} Showa, The extremal root in every Gaussian Mills-ratio Tur\'an bound, internal Round 34 manuscript (2026); identically zero correction.
\bibitem{r35showa:pudim} Pudim Project, APP-0032, public list of applications. \url{https://github.com/pudim-project/pudim-ai-demo-zetalaw/blob/main/APPLICATIONS.md}.
\end{thebibliography}

\par\endgroup

% Source: /workspace/scratch/81914cec2c78/work/archive/Tournament_Round_35_Final/Tournament_Round_35/round35/heisei/paper.tex
% Edited source: /workspace/scratch/81914cec2c78/work/audit_34_41/cleaned/r35_heisei.tex
\begingroup

\Needspace{9\baselineskip}
\section{\texorpdfstring{Divisibility by \(z^k+1\) in the sparse binomial polynomials of Dilcher and Ulas}{Divisibility by z^k+1 in the sparse binomial polynomials of Dilcher and Ulas}}\label{section:r35heisei}
\noindent{\small\sffamily\color{muted}Round 35 \enspace / \enspace Heisei}\par\medskip


\begin{sectionabstract}
For \(m=2^\mu\ge2\), let \(F_{m,n}(z)=\sum_j\binom nj z^{\binom jm}\).
For every \(k\ge1\) we give an arithmetic progression of indices \(n\)
for which \(z^k+1\) divides \(F_{m,n}(z)\) in \(\mathbb Z[z]\).
This proves Conjecture 1 of Dilcher and Ulas. A translation shifts every
exponent by \(k\) modulo \(2k\); a reflection then matches terms having
the same binomial coefficient. The progression is explicit but not asserted
to be optimal.
\end{sectionabstract}

\subsection{The divisibility conjecture}
For \(m,n\ge0\), put \(F_{m,n}(z)=\sum_{j=0}^n\binom nj z^{\binom jm}\),
where \(\binom jm=0\) when \(0\le j<m\).
Dilcher and Ulas~\cite{r35heisei:DU} established infinitely many occurrences of
\(z^k+1\mid F_{2^\mu,n}(z)\) when \(k\) has no odd prime divisor less
than \(2^\mu\). Conjecture~1 of their published paper
(Conjecture~3.6 in the preprint) asks for infinitely many such indices
for every \(k\ge1\). The restriction on the small prime factors of \(k\)
is removed here.
\begin{theorem}\label{r35heisei:main}
Let \(\mu,k\ge1\) be integers, \(m=2^\mu\), and
\(D=\operatorname{lcm}(1,2,\ldots,m)\). For every integer \(h\ge0\),
\[
 n=m-1+(2h+1)kD \quad\Longrightarrow\quad z^k+1\mid F_{m,n}(z).
\]
\end{theorem}
The restriction to a power-of-two lower index is useful for two separate
reasons: the polynomial binomial coefficient is invariant under a
reflection for even \(m\), and \(m\) is the unique integer in
\([1,m]\) having the largest 2-adic valuation.

For integer \(x\) and \(r\ge0\) we use the polynomial definition
\(\binom xr=x(x-1)\cdots(x-r+1)/r!\); it takes integral values,
also at negative integers. Write \(v_2\) for the exponent of 2 in a
nonzero integer.
\begin{lemma}\label{r35heisei:shift}
With \(T=(2h+1)kD\) as above, every integer \(x\) satisfies
\[
 \binom{x-T}{m}\equiv\binom xm+k\pmod {2k}.
\]
\end{lemma}
\begin{proof}
For \(1\le r\le m\), let \(c_r=\binom{-T}{r}\).
The identity \(c_r=(-T/r)\binom{-T-1}{r-1}\) implies \(k\mid c_r\),
since \(r\mid D\). For \(r<m\), we have \(v_2(r)<\mu=v_2(D)\);
therefore \(T/r\) is even after division by \(k\), and \(2k\mid c_r\).
For \(r=m\), \(T/(mk)\) is odd. The other factor
\[
 \binom{-T-1}{m-1}=(-1)^{m-1}\binom{T+m-1}{m-1}
\]
is odd as well: over \(\mathbb F_2\), \(m\mid T\) and
\((1+X)^T=(1+X^m)^{T/m}\), so the coefficient of \(X^{m-1}\) in
\((1+X)^{T+m-1}\) is one. Thus \(c_m\equiv k\bmod 2k\).
Vandermonde's identity
\[
 \binom{x-T}{m}=\sum_{r=0}^m c_r\binom{x}{m-r}
\]
now gives the congruence. For integer arguments of either sign the
identity follows by taking the coefficient of \(X^m\) in the formal power
series identity \((1+X)^{x-T}=(1+X)^{-T}(1+X)^x\).
\end{proof}

\begin{proof}[Proof of Theorem~\ref{r35heisei:main}]
Because \(m\) is even,
\[
 \binom{n-j}{m}=\binom{T+m-1-j}{m}
               =(-1)^m\binom{j-T}{m}
               \equiv \binom jm+k\pmod{2k}
\]
by Lemma~\ref{r35heisei:shift}. Pair the indices \(j\) and \(n-j\); none is
fixed because \(T\) is even and \(n\) is odd. In the quotient ring
\(\mathbb Z[z]/(z^k+1)\), \(z^{2k}=1\) and \(z^k=-1\).
The two nonnegative exponents in each pair therefore contribute opposite
monomials, with equal coefficients \(\binom nj=\binom n{n-j}\).
All pairs cancel. As \(h\) varies, the indices are distinct.
\end{proof}
For example, \(m=4,k=3,n=39\) illustrates the first prime excluded
by the source's general theorem. This proof uses the source's paired-index
idea, not a new general theory of periods of binomial coefficients.
The specific ingredient is a shift that is constant but nonzero modulo
\(2k\), rather than an ordinary period; compare the classical work of
Kwong~\cite{r35heisei:KW} on ordinary periods.

\subsection*{Separate observation on all forward differences}
This observation concerns the other Dilcher--Ulas conjecture (Conjecture~2),
whose sign and first-difference statements were already proved in Round~17.
It plays no role in the divisibility proof. Let
\(g_n=F_{m,n}(-1)\), \(\Delta g_n=g_{n+1}-g_n\), and
\(G(x)=\operatorname{sgn}\sin(\pi(x+1)/(2m))\) for \(x\in\mathbb Z\),
with zero sign at the zeros. Put \(A H(x)=H(x-1)+H(x+1)\).
Reduction of \((1+X)^j\) modulo 2 shows that
\(\binom jm\bmod2\) equals the binary digit of \(j\) of weight \(m\);
hence \(G(2j)=(-1)^{\binom jm}\).
Expansion gives \(g_n=(A^nG)(n)\). Applying the definition of \(A\)
to each forward difference proves inductively that
\[
 \Delta^d g_n=(A^nG)(n+2d)\qquad(d,n\ge0).
\]
Odd reflection about \(-1\) and sign reversal under translation by \(2m\)
are preserved under \(A\). Thus \(A^nG\) vanishes at the stripe endpoints
\(-1\) and \(2m-1\). It is strictly positive on \(0,\ldots,2m-2\) by
induction: each vertex in that interval has at least one neighbor there
and both neighbors lie there or on a zero endpoint. Translation by \(2m\)
now proves \(\operatorname{sgn}\Delta^d g_n=G(n+2d)\) for every order.
Only this all-orders formulation extends beyond the first-difference
scope addressed in the earlier tournament proof.
\begin{thebibliography}{9}
\bibitem{r35heisei:DU} K. Dilcher and M. Ulas,
Divisibility and arithmetic properties of a class of sparse polynomials,
\emph{Integers} \textbf{22} (2022), A8.
\url{https://math.colgate.edu/~integers/w8/w8.pdf}.
\bibitem{r35heisei:KW} Y. H. H. Kwong, Minimum periods of binomial coefficients
modulo \(M\), \emph{Fibonacci Quarterly} \textbf{27} (1989).
DOI: 10.1080/00150517.1989.12429549.
\end{thebibliography}

\par\endgroup

% Source: /workspace/scratch/81914cec2c78/work/archive/Tournament_Round_35_Final/Tournament_Round_35/round35/keio/paper.tex
% Edited source: /workspace/scratch/81914cec2c78/work/audit_34_41/cleaned/r35_keio.tex
\begingroup
\newcommand{\cyc}{\sum_{\rm cyc}}
\Needspace{9\baselineskip}
\section{\texorpdfstring{A median--exradius inequality beyond acute triangles}{A median--exradius inequality beyond acute triangles}}\label{section:r35keio}
\noindent{\small\sffamily\color{muted}Round 35 \enspace / \enspace Keio}\par\medskip

\begin{sectionabstract}
We extend a reciprocal inequality of Jian Liu to right and obtuse triangles, answering his Conjecture 6.3. A rational half-angle chart makes the exradii rational and isolates a full boundary segment on which the cleared gap vanishes. Exact dyadic enclosures prove positivity, using transverse derivatives along that segment. The complete verifier is included.
\end{sectionabstract}
\begin{theorem}
In every nondegenerate Euclidean triangle, with medians $m_a,m_b,m_c$ and exradii $r_a,r_b,r_c$ opposite the respective vertices,
\[
   (m_a+m_b+m_c)\cyc\frac{1}{r_a+2m_a}\leq 3.
\]
Equality holds precisely for equilateral triangles.
\end{theorem}
For acute triangles the assertion, with equality, is Liu's Theorem~5.1 [1]. We prove strictness for the rest. By symmetry let $A\geq\pi/2$, $B\geq C$. Write $t=\tan(B/2)$ and $v=\tan(C/2)/t$. Then $0<t,v\leq1$ and, from $B+C\leq\pi/2$,
\[
 t(1+v)+vt^2\leq1 .
\]
Set $T=t^2$, $f=1-vT$, $k=1+v$. Homogeneity and the law of sines let us take
\[
 (a,b,c)=(kf,\ 1+Tv^2,\ v(1+T)).
\]
Indeed, the common multiplier of $(\sin A,\sin B,\sin C)$ is $(1+T)(1+Tv^2)/(2t)$. The semiperimeter is $k$, and Heron's formula gives the area $vtkf$. Dividing it by $k-a,k-b,k-c$ gives
\[
 (r_a,r_b,r_c)=(f/t,\ tk,\ vtk).
\]
Put $n_a=\sqrt{2(b^2+c^2)-a^2}=2m_a$ and likewise cyclically. With $(d_a,d_b,d_c)=(f+tn_a,tk+n_b,vtk+n_c)$, our claim is that
\[
 D=6-NH>0,\qquad N=n_a+n_b+n_c,\quad H=t/d_a+1/d_b+1/d_c.
\]

\subsection*{A finite positivity certificate}
For $v=0$ we have $(a,b,c)=(1,1,0)$, $(n_a,n_b,n_c)=(1,1,2)$ and $H=3/2$, so that $D(t,0)=0$ on the entire segment. This allows a certificate near the collapsed side: if $D_v>0$ throughout $I\times[0,h]$, then $D>0$ on that rectangle at every $v>0$. To specify this derivative test, use primes for $\partial_v$. Then
\[
 (a',b',c')=(f-kT,2Tv,1+T),\qquad
 n'_a=\frac{2(bb'+cc')-aa'}{n_a}
\]
and cyclically for $n'_b,n'_c$ wherever the denominator is positive. Thus $d'_a=-T+tn'_a$, $d'_b=t+n'_b$, and $d'_c=t(1+2v)+n'_c$. The tested expression is
\[
 D_v=N\left(t\frac{d'_a}{d_a^2}+\frac{d'_b}{d_b^2}+\frac{d'_c}{d_c^2}\right)
       -(n'_a+n'_b+n'_c)H.
\]

\paragraph{Cover.}
Starting with $[0,1]^2$, discard a box if the lower bound for $tk+vT$ exceeds one. Otherwise accept it if the lower bound for $D$ is strictly positive. Failing that, when its upper $v$ endpoint $h\leq1/4$, use the derivative test on the \emph{enlarged} box $I\times[0,h]$. If neither succeeds, bisect, alternating the $t$ and $v$ directions and retaining both closed halves. A failed reciprocal test never accepts a box. The exact computation below gives:
\[
\begin{array}{lr}
\text{positive gap} &8504\\
\text{outside domain}&303\\
\text{positive transverse derivative}&259\\ \hline
\text{maximum bisection depth} & 20
\end{array}
\]
Endpoints have denominator $2^{42}$; each integer operation encloses the exact answer by rounding outward. For a square root we may intersect the radicand enclosure with the nonnegative axis: the median radicands are nonnegative for the sides in this chart. In fact $k-a=kvt^2$, $k-b=vf$, $k-c=f$ are nonnegative on the entire unit square, so degenerate limits and enlarged boxes cause no problem. Divisions require strictly positive lower bounds. Consequently each retained domain point with $v>0$ has strictly positive gap, which proves the theorem.

\paragraph{Reference.} [1] J. Liu, A new geometric inequality in acute triangles and its applications, \emph{International Journal of Geometry} \textbf{12} (2023), no.~4, 47--63, Theorem 5.1 and Conjecture 6.3.
\url{https://ijgeometry.com/wp-content/uploads/2023/09/5.-47-63.pdf}.
\subsection*{Exact verifier (Python 3)}
\noindent The two listings are consecutive parts of one file. Integer endpoints represent the corresponding endpoints divided by $P$.
\small
\begin{Verbatim}
from math import isqrt
from collections import Counter
P=2**42
class I:
 def __init__(s,l,h):s.l=l;s.h=h
 def __add__(s,o):
  o=iv(o);return I(s.l+o.l,s.h+o.h)
 __radd__=__add__
 def __neg__(s):return I(-s.h,-s.l)
 def __sub__(s,o):return s+-iv(o)
 def __rsub__(s,o):return iv(o)+-s
 def __mul__(s,o):
  o=iv(o);v=(s.l*o.l,s.l*o.h,s.h*o.l,s.h*o.h);return I(min(v)//P,-(-max(v)//P))
 __rmul__=__mul__
 def inv(s):
  if s.l<=0:raise ValueError
  return I(P*P//s.h,-(-P*P//s.l))
 def __truediv__(s,o):return s*iv(o).inv()
 def __rtruediv__(s,o):return iv(o)*s.inv()
 def root(s):
  assert s.h>=0
  l=isqrt(max(0,s.l)*P);n=s.h*P;h=isqrt(n);return I(l,h+(h*h<n))
def iv(x):return x if isinstance(x,I) else I(x*P,x*P)
\end{Verbatim}

\subsection*{Verifier, continued}
\footnotesize\begin{Verbatim}
def evalbox(t,v,der):
 T=t*t; f=1-v*T; k=1+v;a=k*f;b=1+T*v*v;c=v*(1+T)
 S=[a,b,c];m=[]
 for i in range(3):
  x,y,z=S[i:]+S[:i];m.append((2*(y*y+z*z)-x*x).root())
 den=[f+t*m[0],t*k+m[1],v*t*k+m[2]]
 inv=[1/x for x in den]
 lhs=sum(m);rhs=t*inv[0]+inv[1]+inv[2]
 if not der:return 6-lhs*rhs
 ds=[f-k*T,2*T*v,1+T]
 dm=[]
 for i in range(3):
  j=(i+1)%3;k0=(i+2)%3
  dm.append((2*(S[j]*ds[j]+S[k0]*ds[k0])-S[i]*ds[i])/m[i])
 dd=[-T+t*dm[0],t+dm[1],t*(1+2*v)+dm[2]]
 return lhs*(t*dd[0]*inv[0]*inv[0]+dd[1]*inv[1]*inv[1]+dd[2]*inv[2]*inv[2])-sum(dm)*rhs
def solve():
 box=[(I(0,P),I(0,P),0)];ct=Counter();md=0
 while box:
  t,v,d=box.pop();md=max(md,d)
  if (t*(1+v)+v*t*t).l>P:ct['out']+=1;continue
  try:
   if evalbox(t,v,False).l>0:ct['value']+=1;continue
  except ValueError:pass
  if v.h<=P//4:
   try:
    if evalbox(t,I(0,v.h),True).l>0:ct['derivative']+=1;continue
   except ValueError:pass
  if d>35:raise RuntimeError(t.l/P,t.h/P,v.l/P,v.h/P)
  if d%2==0:
   m=(t.l+t.h)//2;box.extend([(I(t.l,m),v,d+1),(I(m,t.h),v,d+1)])
  else:
   m=(v.l+v.h)//2;box.extend([(t,I(v.l,m),d+1),(t,I(m,v.h),d+1)])
 print(ct,md)
if __name__=='__main__':solve()
\end{Verbatim}


\par\endgroup
\clearpage
\part{Rounds 36--40}

% Source: /workspace/scratch/81914cec2c78/work/archive/Tournament_Round_36_Final/Tournament_Round_36/round36/meiji/paper.tex
% Edited source: /workspace/scratch/81914cec2c78/work/audit_34_41/cleaned/r36_meiji.tex
\begingroup
\newcommand{\cyc}{\mathop{\sum_{\rm cyc}}}
\Needspace{9\baselineskip}
\section{\texorpdfstring{Reflected mixtures of the three medians and exradii}{Reflected mixtures of the three medians and exradii}}\label{section:r36meiji}
\noindent{\small\sffamily\color{muted}Round 36 \enspace / \enspace Meiji}\par\medskip


\begin{sectionabstract}
For a triangle, let $f(t)$ be the sum of the reciprocals of the three linear mixtures $(1-t)m_i+t r_i$ of corresponding medians and exradii. We show that $f(1-t)>f(t)$ for $0\le t<1/2$ unless the triangle is equilateral. Convexity then proves both parameter regimes in Jian Liu's Conjecture 6.5. The proof removes the unbounded parameters: after reflection the relevant numerator is quadratic, and the signs of its three Bernstein coefficients follow from an exact two-variable certificate.
\end{sectionabstract}

\begin{theorem}
Let $m_i,r_i$ be the medians and corresponding exradii of any nondegenerate triangle. If $p\ge q\ge1$, then
\[
 (p+1)\sum_i\frac1{p r_i+m_i}
 \ \ge\ (q+1)\sum_i\frac1{q m_i+r_i}.
\]
If $0<q\le p\le1$, the inequality is reversed. Equality in either regime holds precisely when $p=q=1$ or the triangle is equilateral.
\end{theorem}
This is Liu's Conjecture 6.5 \cite[p.~62, Eq.~(6.6)]{r36meiji:liu}, not the distinct reciprocal-sum extension posed as Conjecture 6.3 in the same paper.

\paragraph{Removing the two parameters.}
Let $f(t)=\sum_i((1-t)m_i+tr_i)^{-1}$ on $[0,1]$. Each summand is convex. For $0\le t<1/2$, set $d=t(1-t)$; direct subtraction gives
\[
 \frac{f(1-t)-f(t)}{1-2t}
   =\sum_i\frac{m_i-r_i}{m_ir_i+d(m_i-r_i)^2}.
 \tag{1}
\]
We prove that the right side is positive for all $0\le d\le1/4$ when the triangle is not equilateral. Its value at $d=1/4$ is $f'(1/2)$, so convexity also makes $f$ strictly increasing on $[1/2,1]$. Write $a_p=p/(p+1)$ and $b_q=1/(q+1)$. If $p\ge q\ge1$, then $a_p\ge1-b_q\ge1/2$ and $b_q\le1/2$, giving
$f(a_p)\ge f(1-b_q)\ge f(b_q)$. If $0<q\le p\le1$, then $b_q\ge1-a_p\ge1/2$ and $a_p\le1/2$, giving $f(a_p)\le f(1-a_p)\le f(b_q)$. These are the claimed inequalities including their strictness, apart from $p=q=1$ or a constant $f$ in the equilateral case.

\paragraph{A square chart and three coefficients.}
Let $s$ be the semiperimeter. Sort the tangent lengths (semiperimeter minus side), and rescale them to
\[
 (t_0,t_1,t_2)=(u^2v^2,v^2,1),\qquad 0<u,v\le1.
\]
Then $s=1+v^2(1+u^2)$; set $g=\sqrt{s}$. Heron's formula and the median identity give
\[
 (r_0,r_1,r_2)=(g/u,ug,uv^2g),\qquad
 m_i^2=st_i+(t_j-t_k)^2/4 .
\]
Put $e=(u,1,1)$, $N_i=e_i^2(m_i-r_i)$, and let
\[
 A_i=e_i^2 m_ir_i/u,\qquad C_i=e_i^2(m_i+r_i)^2 .
\]
For $\lambda=4d$, four times $e_i^2$ times the denominator of the $i$-th term in (1) equals $G_i=4(1-\lambda)u A_i+\lambda C_i$. Hence the sign of (1) is the sign of $\sum_i N_iG_jG_k$. Write $L(A,B)=\cyc N_i A_j B_k$, with cyclic order $0,1,2$, and define
\[
 K_0=L(A,A),\quad K_1=L(A,C)+L(C,A),\quad K_2=L(C,C).
\]
The numerator becomes
\[
 16u^2(1-\lambda)^2 K_0+
 4u\lambda(1-\lambda)K_1+\lambda^2K_2.
 \tag{2}
\]
Thus it suffices to prove strict positivity of these three functions of the shape alone, except at the equilateral point. No parameter was fixed or sampled.

\paragraph{The equilateral corner.}
Use $x=1-u,\ z=1-v$. At $x=z=0$ the $N_i$ vanish and both $A$ and $C$ have constant entries. Thus $K_i=0$ there, and their gradients are constant multiples of $\nabla\sum_i(m_i-r_i)=0$. For this last assertion the sum is symmetric, smooth and homogeneous of degree one in the original sides. Its value at the equilateral triangle is zero, so Euler's identity makes its three equal partial derivatives zero. The factors $e_i^2$ contribute no additional first derivatives since $m_i-r_i=0$ there.

On $[0,2^{-14}]^2$ the verified lower bounds for $1024$ times the corresponding Hessian entries are:
\[
\begin{array}{c|rrr}
 &\partial_{xx}&\partial_{xz}&\partial_{zz}\\ \hline
 K_0&42268&20785&41812\\
 K_1&337225&165491&334178\\
 K_2&672257&329349&667710
\end{array}
\]
Each $K_i$ therefore has strictly positive second derivative along any nonconstant segment from the origin in the nonnegative quadrant. Integrating twice proves the required local sign.

\paragraph{The degenerate corner and the rest of the square.}
As $(u,v)\to(0,1)$, $m_0$ can vanish. Here use $m_0\ge uvg$, $N_0\ge-ug$, $C_0\ge g^2$ and $A_0=g m_0\ge uv g^2$. The special test is used only when $N_1,N_2\ge0$ are enclosed. A lower bound for $K_0/u$ is then
\[
 -gA_1A_2+vg^2(N_1A_2+N_2A_1).
\]
For $K_1,K_2$ substitute $-ug,0,g^2$ respectively for $N_0,A_0,C_0$: in each term where $A_0$ or $C_0$ occurs its coefficient is nonnegative. All other terms retain their original values. These tests handle genuine triangles arbitrarily close to the degenerate corner without dividing by a small parameter.

Every remaining box in the unit square is certified by an outward interval bound on $K_i$ or its centered form
\[
 K_i(c)+[\partial_x K_i](B_x-c_x)+[\partial_z K_i](B_z-c_z)
\]
on the box $B$ with midpoint $c$. This enclosure follows by integration along the full segment from $c$, not by a sample at $c$. Failing boxes are bisected along their longer side, with both children retained. The complete verifier is printed below. It uses fixed-point integer intervals with endpoints in units of $2^{-64}$, exact outward rounding, and integer square roots. A derivative test requiring a radicand not proved positive is not accepted. Natural interval differentiation uses the chain and product rules, including the mixed product rule for the local Hessian. The exact partition counts are:
\[
\begin{array}{c|rrrr|r}
 &\text{centered}&\text{direct}&\text{special}&\text{local}&\text{max depth}\\ \hline
 K_0&1015&26&14&1&34\\
 K_1&1074&35&2&1&34\\
 K_2&1148&36&3&1&34
\end{array}
\]
They give strictly positive $K_0,K_1,K_2$ on the entire required domain except at the equilateral point. Substitution in (2) and (1) finishes the proof. \hfill$\square$

\begin{thebibliography}{1}
\bibitem{r36meiji:liu} J. Liu, \emph{A new geometric inequality in acute triangles and its applications}, Int. J. Geometry 12 (2023), no.~4, 47--63, \href{https://ijgeometry.com/wp-content/uploads/2023/09/5.-47-63.pdf}{publisher PDF}.
\end{thebibliography}
\subsection*{Appendix: standalone verifier (Python 3)}
\begingroup\scriptsize
\begin{Verbatim}
from math import isqrt
from collections import Counter
S=1<<64
class I:
 __slots__='l','h'
 def __init__(self,l,h=None):
  if h is None:self.l=self.h=l*S
  else:self.l=l;self.h=h
 def __add__(x,y):
  if not isinstance(y,I):y=I(y)
  return I(x.l+y.l,x.h+y.h)
 __radd__=__add__
 def __neg__(x):return I(-x.h,-x.l)
 def __sub__(x,y):return x+-y
 def __rsub__(x,y):return -x+y
 def __mul__(x,y):
  if not isinstance(y,I):return I(min(x.l*y,x.h*y),max(x.l*y,x.h*y))
  if x.l>=0 and y.l>=0:return I(x.l*y.l//S,-(-x.h*y.h//S))
  p=(x.l*y.l,x.l*y.h,x.h*y.l,x.h*y.h);return I(min(p)//S,-(-max(p)//S))
 __rmul__=__mul__
 def __truediv__(x,y):
  if not isinstance(y,I):
   assert y>0;return I(x.l//y,-(-x.h//y))
  assert y.l>0, 'unsafe denominator'
  return x*I(S*S//y.h,-(-S*S//y.l))
 def rt(x):
  if x.l<=0:raise ValueError
  return I(isqrt(x.l*S),isqrt(x.h*S)+1)
 def pow32(x):
  y=I(max(x.l,0),max(x.h,0))
  return y*I(isqrt(y.l*S),isqrt(y.h*S)+1)
class J:
 __slots__='v','g'
 def __init__(self,v,g):
  self.v=v if isinstance(v,I) else I(v);self.g=[d if isinstance(d,I) else I(d) for d in g]
 def __add__(x,y):
  if not isinstance(y,J):return J(x.v+y,x.g)
  return J(x.v+y.v,[p+q for p,q in zip(x.g,y.g)])
 __radd__=__add__
 def __neg__(x):return J(-x.v,[-p for p in x.g])
 def __sub__(x,y):return x+-y
 def __rsub__(x,y):return -x+y
 def __mul__(x,y):
  if not isinstance(y,J):return J(x.v*y,[p*y for p in x.g])
  return J(x.v*y.v,[p*y.v+q*x.v for p,q in zip(x.g,y.g)])
 __rmul__=__mul__
 def __truediv__(x,y):
  if not isinstance(y,J):return J(x.v/y,[p/y for p in x.g])
  v=x.v/y.v;return J(v,[(p-v*q)/y.v for p,q in zip(x.g,y.g)])
 def rt(x):
  y=x.v.rt();return J(y,[p/(2*y) for p in x.g])
 def pow32(x):
  y=I(max(x.v.l,0),max(x.v.h,0));r=I(isqrt(y.l*S),isqrt(y.h*S)+1)
  return J(y*r,[3*r*p/2 for p in x.g])
def parts(x,z,corner=False):
 u,v=1-x,1-z
 v2=v*v;u2=u*u;s=1+v2*(1+u2);r=s.rt()
 X=(1-v2)/2;Y=(1-u2*v2)/2;Z=v2*(1-u2)/2
 if corner:
  m0=u*v*r
 else: m0=(s*u2*v2+X*X).rt()
 m1=(s*v2+Y*Y).rt();m2=(s+Z*Z).rt()
 T=[u*m0-r,m1-u*r,m2-u*v2*r]
 N=[u*T[0],T[1],T[2]]
 A=[r*m0,r*m1,v2*r*m2]
 L=[u*m0+r,m1+u*r,m2+u*v2*r]
 C=[e*e for e in L]
 return N,A,C
def cyc(N,A,B):
 return sum(N[i]*A[(i+1)%3]*B[(i+2)%3] for i in range(3))
def f(x,z):
 N,A,C=parts(x,z)
 if test==0:return cyc(N,A,A)
 if test==1:return cyc(N,A,C)+cyc(N,C,A)
 return cyc(N,C,C)
def special(l,h,u,v):
 # x near 1, z near 0; check N1,N2>=0 and use m0>=u0*v0*r.
 if l<3*S//4 or v>S//4:return False
 x,z=I(l,h),I(u,v); uu,vv=1-x,1-z
 N,A,C=parts(x,z,True)
 if min(N[1].l,N[2].l)<0:return False
 r=(1+vv*vv*(1+uu*uu)).rt()
 if test==0:
  # replace N0>=-uu*r and divide out uu>0
  return (-r*A[1]*A[2]+vv*r*r*(N[1]*A[2]+N[2]*A[1])).l>0
 # C0>=r^2 and A0>=0; signed N0 bounded below independently.
 N[0]=-uu*r;A[0]=I(0);C[0]=r*r
 if test==1:return (cyc(N,A,C)+cyc(N,C,A)).l>0
 return cyc(N,C,C).l>0
def prove(eps=16384):
 counts=Counter();stack=[(0,S,0,S,0)]
 while stack:
  l,h,u,v,d=stack.pop();counts['max depth']=max(d,counts['max depth'])
  if special(l,h,u,v):counts['corner']+=1;continue
  if max(h,v)<=S//eps:counts['local']+=1;continue
  left=I(l,h);right=I(u,v)
  try:
   if f(left,right).l>0:counts['value']+=1;continue
   midx=(l+h)//2;midz=(u+v)//2
   val=f(I(midx,midx),I(midz,midz))
   jet=f(J(left,[1,0]),J(right,[0,1]))
   if (val+jet.g[0]*I(l-midx,h-midx)+jet.g[1]*I(u-midz,v-midz)).l>0:counts['mean']+=1;continue
  except ValueError:pass
  assert d<50,(d,l/S,h/S,u/S,v/S,counts)
  if h-l>=v-u:
   w=(l+h)//2;stack.extend([(l,w,u,v,d+1),(w,h,u,v,d+1)])
  else:
   w=(u+v)//2;stack.extend([(l,h,u,w,d+1),(l,h,w,v,d+1)])
 print(counts)

class B:
 def __init__(self,v,g=0,h=0,j=0):
  self.v,self.g,self.h,self.j=[i if isinstance(i,I) else I(i) for i in (v,g,h,j)]
 def __add__(x,y):
  if not isinstance(y,B):y=B(y)
  return B(x.v+y.v,x.g+y.g,x.h+y.h,x.j+y.j)
 __radd__=__add__
 def __neg__(x):return B(-x.v,-x.g,-x.h,-x.j)
 def __sub__(x,y):return x+-y
 def __rsub__(x,y):return -x+y
 def __mul__(x,y):
  if not isinstance(y,B):y=B(y)
  return B(x.v*y.v,x.g*y.v+x.v*y.g,x.h*y.v+x.v*y.h,x.j*y.v+x.g*y.h+x.h*y.g+x.v*y.j)
 __rmul__=__mul__
 def chain(x,y,a,b):return B(y,a*x.g,a*x.h,a*x.j+b*x.g*x.h)
 def __truediv__(x,y):
  if not isinstance(y,B):return B(x.v/y,x.g/y,x.h/y,x.j/y)
  r=I(1)/y.v;return x*y.chain(r,-r*r,2*r*r*r)
 def rt(x):
  r=x.v.rt();d=I(1)/(2*r);return x.chain(r,d,-d/(2*x.v))
 def pow32(x):
  r=x.v.rt();return x.chain(x.v*r,3*r/2,I(3)/(4*r))
def local(eps):
 box=I(0,S//eps)
 for p,q in [(0,0),(0,1),(1,1)]:
  j=f(B(box,p==0,q==0),B(box,p==1,q==1)).j
  print('local Hessian lower > numerator/1024:',p,q,j.l*1024//S)
  assert j.l>0
if __name__=='__main__':
 for test in range(3):
  print('coefficient',test,flush=True);local(16384);prove()
\end{Verbatim}

\endgroup

\par\endgroup

% Source: /workspace/scratch/81914cec2c78/work/archive/Tournament_Round_36_Final/Tournament_Round_36/round36/showa/paper.tex
% Edited source: /workspace/scratch/81914cec2c78/work/audit_34_41/cleaned/r36_showa.tex
\begingroup

\Needspace{9\baselineskip}
\section{\texorpdfstring{Three Bessel-zero constants in Gauss--Legendre quadrature: phase convexity and the heat kernel of the disk}{Three Bessel-zero constants in Gauss--Legendre quadrature: phase convexity and the heat kernel of the disk}}\label{section:r36showa}
\noindent{\small\sffamily\color{muted}Round 36 \enspace / \enspace Showa}\par\medskip

\begin{sectionabstract}
We prove Conjectures 2, 5 and 8 of L\'opez Pouso and Segura on three sequences built from zeros of $J_0$ and the values of $J_1$ at those zeros. The first and third sequences increase strictly to $\pi^2/12$, and the second decreases strictly to $-\pi^2/6$. Convexity of a continuous inverse-phase derivative turns two differences into trapezoidal and midpoint errors of opposite signs. The constant left undetermined by local Bessel asymptotics is recovered from the short-time Dirichlet heat kernel at the center of the unit disk. These theorems concern the three coefficient sequences, not the paper's separate conjectures on finite Gauss--Legendre nodes and weights.
\end{sectionabstract}

\subsection{The three assertions}
Let $j_k$ be the $k$th positive zero of $J_0$ and $w_k=J_1(j_k)^{-2}$, for $k\ge1$. With $W_i=\sum_{k=1}^i w_k$ and $W_0=0$, the coefficients in equations (16), (18) and (19) of \cite{r36showa:lps} are
\[
\begin{aligned}
 C_{k-1}&=2W_{k-1}+w_k-j_k^2/2 &&(k\ge1),\\
 E_i&=2W_i-(j_i^2+j_{i+1}^2)/4 &&(i\ge1),\\
 K_i&=W_i-\frac{\sum_{k=1}^i w_k j_k^2}{2W_i} &&(i\ge1).
\end{aligned}
\]
The original paper's Appendix D establishes finiteness of the limit of $C_i$. We determine its exact value.
\begin{theorem}
The sequence $(C_i)_{i\ge0}$ is positive and strictly increasing to $\pi^2/12$. The sequence $(E_i)_{i\ge1}$ is negative and strictly decreasing to $-\pi^2/6$. The sequence $(K_i)_{i\ge1}$ is positive and strictly increasing to $\pi^2/12$.
\end{theorem}

\subsection{A convex integrand in zero-index coordinates}
Choose the continuous phase $\theta(x)=\arg(J_0(x)+iY_0(x))$ for $x>0$. Set $f(x)=x(J_0(x)^2+Y_0(x)^2)$. The Bessel Wronskian gives $\theta'=2/(\pi f)>0$. Consecutive zeros of $J_0$ occur at consecutive half-integer multiples of $\pi$ in this phase. On $s\ge1$ define $x(s)$ by $\theta(x(s))=\theta(j_1)+\pi(s-1)$. Then
\[
 x(k)=j_k,\qquad x'=c f(x),\qquad c=\pi^2/2.
\]
At a zero $j$, $J_1(j)Y_0(j)=2/(\pi j)$. It follows that $h(s):=x(s)x'(s)$ satisfies $h(k)=2w_k$.

Nicholson's integral for $J_0^2+Y_0^2$ \cite{r36showa:dlmf}, with $v=2x\sinh t$, yields
\[
 f(x)=\frac4{\pi^2}\int_0^\infty K_0(v)
           \frac{x}{\sqrt{x^2+v^2/4}}\,dv .
\]
For $g(x)=x/\sqrt{x^2+a^2}$, $a>0$, we have $g'(x)=a^2/(x^2+a^2)^{3/2}>0$ and
$3g'(x)+xg''(x)=3a^4/(x^2+a^2)^{5/2}>0$.
Integration against the strictly positive kernel shows $f'>0$ and $3f'+xf''>0$; differentiation is justified on compact $x$-subintervals by domination, using integrability of $K_0$ at zero and exponential decay at infinity. In particular $x''=c^2ff'>0$ and
\[
 h''(s)=c^3 f\big(f(3f'+xf'')+x(f')^2\big)>0.
\]
Thus
\[
 C_k-C_{k-1}=\frac{h(k)+h(k+1)}2-\int_k^{k+1} h(s)\,ds>0
\]
by strict convexity and $\int_a^b h=(x(b)^2-x(a)^2)/2$. The sign reverses for the midpoint:
\[
 E_{i+1}-E_i=h(i+1)-\tfrac12\int_i^{i+2}h(s)\,ds<0 .
\]

Here are explicit checks of the initial signs. The alternating Bessel power series give $J_0(2.40)>0$, $J_0(2.41)<0$ (lower and upper sums to degrees $10$ and $8$, respectively), and $J_1(x)\ge(x/2)(1-x^2/8)>0$ for $0<x\le2.41$. Thus $2.40<j_1<2.41$, since $J_0'=-J_1$. Furthermore
\[
 J_1(j_1)\le 2.41/2-2.40^3/16+2.41^5/384<0.56 .
\]
Since $(2.41\cdot0.56)^2<2$, we have $C_0>0$. For the other initial sign, strict convexity of $x$ yields $j_2>j_1+x'(1)=j_1+2w_1/j_1$. Now
\[
 j_1^2+(j_1+2w_1/j_1)^2-8w_1
 =j_1^2+ (2w_1/j_1-j_1)^2>0 ,
\]
so $E_1<0$.

\subsection{The heat-kernel sum rule}
Set $a_k=k-1/4$. McMahon's zero expansion \cite{r36showa:dlmf} gives
\[
 j_k=\pi a_k+\frac1{8\pi a_k}+O(a_k^{-3}),\qquad
 j_k^2=\pi^2 a_k^2+\tfrac14+O(k^{-2}) .
\]
In Nicholson's integral expand $(1+y)^{-1/2}=1-y/2+O(y^2)$ uniformly for $y\ge0$, since its second derivative is bounded there. The classical Mellin integral of $K_0$ gives $\int K_0(v)\,dv=\int v^2 K_0(v)\,dv=\pi/2$ over the positive axis, and its fourth moment is finite. Consequently
\[
 f(x)=\frac2\pi\left(1-\frac1{8x^2}+O(x^{-4})\right),\qquad
 w_k=\frac{\pi^2}4 j_k f(j_k)=\frac{\pi^2}2 a_k+e_k,\quad
 e_k=O(k^{-3}).
\]
The terms of order $a_k^{-1}$ cancel. Determining the sum of these summable remainders requires global information.

In the unit disk, let $p_D(t;0,0)$ be the center diagonal of the Dirichlet heat kernel with generator the Laplacian. Separation of variables leaves only angular order zero at the center; the squared norm there of the normalized $k$th radial eigenfunction is $w_k/\pi$ because $2\pi\int_0^1 J_0(j_kr)^2r\,dr=\pi J_1(j_k)^2$ and $J_0(0)=1$. The center is unit distance from the boundary. The interior short-time heat-kernel estimate \cite{r36showa:heat} gives $p_D(t;0,0)=(4\pi t)^{-1}+O(e^{-c_0/t})$ for some $c_0>0$, where powers of $1/t$ may be absorbed by decreasing $c_0$. Hence
\[
 \sum_{k=1}^{\infty} w_k e^{-t j_k^2}=\frac1{4t}+o(1)\qquad(t\downarrow0).
\]
The summable corrections contribute $\sum e_k+o(1)$ by dominated convergence. In the remaining sum replacing $j_k^2$ by $\pi^2a_k^2+1/4$ costs at most $O(t\sum_{k\ge1}k^{-1}e^{-c_1t k^2})=o(1)$ by the mean-value theorem; $j_k^2$ and the model are both bounded below by $c_1 k^2$ for some $c_1>0$.

Use the shifted Euler--Maclaurin formula with $\alpha=3/4$, $g(z)=ze^{-z^2}$, and $h_0=\pi\sqrt t$. Explicitly
\[
 \sum_{m=0}^\infty (m+\alpha)e^{-h_0^2(m+\alpha)^2}
 =\frac1{2h_0^2}-\frac{B_2(\alpha)}{2}+O(h_0).
\]
Indeed the integral term is $h_0^{-2}\int_0^\infty g$, the next term vanishes as $g(0)=0$, and $g'(0)=1$. The third-order remainder is $O(h_0)$ after division by $h_0$, since $g^{(3)}$ is integrable. As $B_2(3/4)=-1/48$, we find
\[
 e^{-t/4}\sum_{k\ge1}\frac{\pi^2}2a_k e^{-t\pi^2a_k^2}
 =\frac1{4t}+\frac{\pi^2}{192}-\frac1{16}+o(1),\qquad
 \sum_{k\ge1} e_k=\frac1{16}-\frac{\pi^2}{192}.
\]
The polynomial terms in $C_{k-1}$ simplify using
$\sum_{l<k} a_l+a_k/2-a_k^2/2=3/32$. Taking the limit now gives
\[
 \lim C_{k-1}=3\pi^2/32-1/8+2\sum e_k=\pi^2/12.
\]

\subsection{The companion limits}
Since $E_i=C_{i-1}+w_i-(j_{i+1}^2-j_i^2)/4$, the preceding asymptotics yield $E_i-C_{i-1}\to-\pi^2/4$. This proves the claimed limit of $E_i$. Finally
\[
 \sum_{k=1}^i w_k C_{k-1}=W_i^2-\tfrac12\sum_{k=1}^i w_k j_k^2 ,
\]
so $K_i$ is precisely the weighted average of $C_0,\ldots,C_{i-1}$ with weights $w_1,\ldots,w_i$. This average is positive and strictly increasing because the sequence averaged is. As $W_i\to\infty$, it has the same limit $\pi^2/12$.

\begin{thebibliography}{9}
\bibitem{r36showa:lps} O. L\'opez Pouso and J. Segura, Uniform relations between the Gauss--Legendre nodes and weights, \emph{J. Inequal. Appl.} 2025, article 40. \url{https://doi.org/10.1186/s13660-025-03283-w}.
\bibitem{r36showa:dlmf} NIST Digital Library of Mathematical Functions, Chapter 10: Nicholson's integral (10.9.30), McMahon expansions (10.21.19), Bessel orthogonality (10.22.37), and the Mellin transform (10.43.19). \url{https://dlmf.nist.gov/10}.
\bibitem{r36showa:heat} L. Li and A. Strohmaier, Heat kernel estimates for general boundary problems, \emph{J. Spectr. Theory} 6 (2016), 903--919. \url{https://arxiv.org/abs/1604.00784}.
\end{thebibliography}

\par\endgroup

% Source: /workspace/scratch/81914cec2c78/work/archive/Tournament_Round_36_Final/Tournament_Round_36/round36/heisei/paper.tex
% Edited source: /workspace/scratch/81914cec2c78/work/audit_34_41/cleaned/r36_heisei.tex
\begingroup
\newcommand{\Prb}{\mathbb P}
\newcommand{\E}{\mathbb E}
\newcommand{\Var}{\operatorname{Var}}
\newcommand{\Cov}{\operatorname{Cov}}
\Needspace{9\baselineskip}
\section{\texorpdfstring{The sharp order threshold for log-concavity of the lower Marcum function in the radius}{The sharp order threshold for log-concavity of the lower Marcum function in the radius}}\label{section:r36heisei}
\noindent{\small\sffamily\color{muted}Round 36 \enspace / \enspace Heisei}\par\medskip


\begin{sectionabstract}
For order \(\nu\ge1/2\), the distribution function \(1-Q_\nu(a,b)\) of the
noncentral chi radius is jointly log-concave in the displacement \(a\)
and radius \(b\). For every \(0\le\nu<1/2\) and all sufficiently large
\(a\), it fails to be log-concave in \(b\). The positive result settles the
radius range left open in Sun, Baricz and Zhou's Conjecture 2. The
squared-radius assertion in that conjecture is a theorem of Finner and
Roters; a separate covariance proof is included. No log-concavity of an
upper tail in its order is asserted.
\end{sectionabstract}

\subsection{Statement and earlier results}
For \(a\ge0\), let \(X\) have a noncentral chi-square distribution with
\(2\nu\) degrees of freedom and noncentrality \(a^2\). Write
\(P_\nu(a,b)=1-Q_\nu(a,b)=\Prb(X\le b^2)\), \(b>0\).
This includes \(\nu=0\) with the usual zero-degree Poisson mixture and
its atom at zero. Conjecture 2 of Sun, Baricz and Zhou~\cite{r36heisei:SBZ}
asserts that, for every \(a>0\), \(P_\nu(a,\sqrt b)\) is log-concave
in \(b>0\) for \(0\le\nu<1\), and \(P_\nu(a,b)\) is log-concave in
\(b>0\) for \(1/2\le\nu<3/2\). The authors' Theorem 2 gives the
respective ranges \(\nu\ge1\) and \(\nu\ge3/2\). In fact the whole
squared-radius result had already been proved by Finner and Roters
\cite[Theorem 3.4]{r36heisei:FR}, including the zero-degree case.
Yu~\cite[Corollary 1]{r36heisei:Yu} proved the radius assertion for all \(a\ge0\)
if \(\nu\ge\nu_0\simeq0.78449776\), and for \(\nu\ge1/2\) if
\(a\le1\). He explicitly identified the residual region.

\begin{theorem}\label{r36heisei:thm}
For every \(\nu\ge1/2\), the map \((a,b)\mapsto P_\nu(|a|,b)\)
is jointly log-concave on \(\mathbb R\times(0,\infty)\). For every
\(0\le\nu<1/2\) and all sufficiently large \(a>0\), the function
\(b\mapsto P_\nu(a,b)\) is not log-concave.
\end{theorem}
The failure assertion is about the range uniform in \(a\); it does
not say that all positive displacements fail below the threshold.

\subsection{Fractional degrees via a one-dimensional slice}
For \(d>0\), let \(H_d(s)=\Prb(\sqrt{\chi^2_d}\le s)\), \(s>0\).
This function is positive, nondecreasing and log-concave. Indeed its
density is a positive constant times \(s^{d-1}e^{-s^2/2}\). If
\(0<d<1\), the density decreases; its integral from zero is a
positive concave function, hence log-concave. If \(d\ge1\), the density
is log-concave, as is its integral from zero by the marginal theorem
of Pr\'ekopa~\cite{r36heisei:Pr}: multiply the density at \(u\) by the indicator
of the convex set \(0<u<s\), then integrate in \(u\).

Let \(\delta=2\nu\ge1\). Take independent \(Z\sim N(a,1)\) and
\(Y\sim\chi^2_{\delta-1}\), using \(Y=0\) if \(\delta=1\).
Then \(Z^2+Y\) has the law of \(X\). For instance its Laplace transform
at \(t\ge0\) is
\[
(1+2t)^{-\delta/2}\exp[-a^2t/(1+2t)],
\]
the noncentral chi-square transform, also obtained from its
Poisson mixture. For \(\delta>1\) this gives
\begin{equation}\label{r36heisei:slice}
P_\nu(|a|,b)=\int_{-b}^b \phi(z-a)\,
H_{\delta-1}(\sqrt{b^2-z^2})\,dz,\qquad
\phi(v)=(2\pi)^{-1/2}e^{-v^2/2}.
\end{equation}
On the convex cone \(b>|z|\), the function \(s=\sqrt{b^2-z^2}\)
is concave: its second differential is
\(- (z\,db-b\,dz)^2/s^3\).
Composing it with the concave nondecreasing \(\log H_{\delta-1}\)
gives a concave function of \((b,z)\).
Multiplication by \(\phi(z-a)\) therefore leaves a jointly log-concave
integrand in \((a,b,z)\); set it to zero off the cone.
Marginalization in \(z\) proves joint log-concavity.
For \(\delta=1\) the same proof simply integrates
\(\phi(z-a)\mathbf1_{\{|z|<b\}}\), without introducing \(H_0\).
This proves the first part of Theorem~\ref{r36heisei:thm}.

\subsection{Why the threshold cannot be lowered}
Fix \(s>0\) and let \(a\to\infty\). For \(\nu>0\) the density in the
radius variable \(b>0\) is
\(a^{1-\nu} b^\nu e^{-(a^2+b^2)/2}I_{\nu-1}(ab)\), in the notation
of~\cite{r36heisei:SBZ}. With \(u=ab\), this gives
\[
e^{a^2/2}a^{2\nu}P_\nu(a,s/a)
 =\int_0^s e^{-u^2/(2a^2)}u^\nu I_{\nu-1}(u)\,du
 \longrightarrow s^\nu I_\nu(s).
\]
Indeed \((u^\nu I_\nu(u))'=u^\nu I_{\nu-1}(u)\), and its primitive
tends to zero at \(0\) for \(\nu>0\). These identities also follow from
the power series for \(I_\nu\); the integrand is integrable at zero and
convergence follows by dominated convergence. At \(\nu=0\) the atom of
mass \(e^{-a^2/2}\) contributes \(1\) on the scaled left. The integral is
then \(\int_0^s e^{-u^2/(2a^2)}I_1(u)\,du\), and the limit is
\(1+I_0(s)-I_0(0)=I_0(s)\).

Set \(H(s)=s^\nu I_\nu(s)>0\). The standard asymptotic
\(I_\nu(s)\sim e^s/\sqrt{2\pi s}\) for fixed order and positive
\(s\to\infty\)~\cite[10.40(i)]{r36heisei:DLMF} implies
\[
 \lim_{s\to\infty}\frac{H(2s)^2}{H(s)H(3s)}
   = (4/3)^{\nu-1/2}.
\]
For \(\nu<1/2\), choose a fixed \(s\) with this quotient below one.
The same strict inequality holds with \(H(js)\) replaced by
\(P_\nu(a,js/a)\), \(j=1,2,3\), for all sufficiently large \(a\).
This contradicts midpoint log-concavity and proves the second assertion.

\subsection{A covariance proof for the known square-threshold result}
Write \(x=b/2,\lambda=a^2/2>0\), so
\(F(x)=P_\nu(a,\sqrt{2x})\). A Poisson random variable of mean
\(\lambda\) mixes gamma laws of shapes \(\nu+s\), \(s=0,1,\ldots\),
and scale one for \(X/2\). For \(t>0\) the shape-\(t\) gamma CDF is
\(e^{-x}\sum_{\ell\ge0} x^{t+\ell}/\Gamma(t+\ell+1)\), as one
checks by differentiating and taking the limit at zero. This formula
is \(1\) when \(t=0\), as required for the atom. By nonnegativity,
\[
 F(x)=e^{-\lambda-x}\sum_{j\ge0}c_j\frac{x^{\nu+j}}{\Gamma(\nu+j+1)},
 \qquad c_j=\sum_{s=0}^j\frac{\lambda^s}{s!}.
\]
Normalize the summands after the exponential to a probability
distribution \(p_j\) on \(j\ge0\), and set \(M=\nu+j\) and
\(R_j=x c_{j+1}/c_j\). The \(R_j\) decrease, since
\[
 \frac{c_j}{\lambda^{j+1}/(j+1)!}
 =\sum_{\ell=1}^{j+1}\frac{(j+1)j\cdots(j+2-\ell)}{\lambda^\ell}
\]
increases in \(j\).
The identity \(p_j(\nu+j)=p_{j-1}R_{j-1}\) for \(j\ge1\) gives
\(\E M=\E R_j+\nu p_0\) and
\(\E[M(M-1)]=\E[R_j M]+\nu(\nu-1)p_0\).
Termwise differentiation is valid locally for \(x>0\) by
\(c_j\le e^\lambda\) and factorial decay. Consequently
\[
 x^2(\log F)''=\Var(M)-\E M
 =\Cov(R_j,M)+\nu p_0(\nu-1-\E M)\le0.
\]
The covariance is nonpositive because \(M\) increases and \(R_j\)
decreases in \(j\) (compare two independent copies); also \(\E M\ge\nu\).
This alternate argument gives the known square-threshold assertion
for all \(\nu\ge0\), independently of the radius proof.
\begin{thebibliography}{9}
\bibitem{r36heisei:SBZ} Y. Sun, \'A. Baricz, S. Zhou, On the monotonicity,
log-concavity, and tight bounds of the generalized Marcum and Nuttall
\(Q\)-functions, \emph{IEEE Trans. Inform. Theory} \textbf{56}
(2010), 1166--1186. \url{https://arxiv.org/abs/1003.2760}.
\bibitem{r36heisei:FR} H. Finner, M. Roters, Log-concavity and inequalities for
chi-square, F and beta distributions with applications in multiple
comparisons, \emph{Statistica Sinica} \textbf{7} (1997), 771--787.
\bibitem{r36heisei:Yu} Y. Yu, On log-concavity of the generalized Marcum
\(Q\) function, 2011, \url{https://arxiv.org/abs/1105.5762}.
\bibitem{r36heisei:Pr} A. Pr\'ekopa, On logarithmic concave measures and
functions, \emph{Acta Sci. Math. (Szeged)} \textbf{34} (1973), 335--343.
\bibitem{r36heisei:DLMF} NIST Digital Library of Mathematical Functions,
\url{https://dlmf.nist.gov/10.40}.
\end{thebibliography}

\par\endgroup

% Source: /workspace/scratch/81914cec2c78/work/archive/Tournament_Round_36_Final/Tournament_Round_36/round36/reiwa/paper.tex
% Edited source: /workspace/scratch/81914cec2c78/work/audit_34_41/cleaned/r36_reiwa.tex
\begingroup

\Needspace{9\baselineskip}
\section{\texorpdfstring{The reciprocal cubes of pairwise sums of triangle medians}{The reciprocal cubes of pairwise sums of triangle medians}}\label{section:r36reiwa}
\noindent{\small\sffamily\color{muted}Round 36 \enspace / \enspace Reiwa}\par\medskip


\begin{sectionabstract}
We prove Liu's 2024 inequality for the sum of reciprocal cubes of pairwise median sums. Refining a known geometric lower bound once gives a rational lower bound sharp enough for a global polynomial certificate. All calculations are over integers, without sampling triangles.
\end{sectionabstract}
\begin{theorem}
For every nondegenerate triangle of circumradius $R$, inradius $r$, and medians $m_a,m_b,m_c$,
\[
\sum_{\mathrm{cyc}}\frac1{(m_b+m_c)^3}
\le \frac1{72}\left(\frac1R+\frac1{2r}\right)^3 .
\]
Equality holds exactly for equilateral triangles.
\end{theorem}
\noindent This proves Conjecture 1 in \cite[Eq.~(5.1)]{r36reiwa:liu}; the author's inverse-square bound (1.6) is weaker.

\begin{proof}
Use tangent lengths $x=s-a$, $y=s-b$, $z=s-c$, where $s=x+y+z$ is the semiperimeter. Then $a=y+z$, $b=z+x$, $c=x+y$. Put $D=sxyz=\Delta^2$ for area $\Delta$, $J=abc$, and $A_a=4m_a^2=4sx+(y-z)^2$, with analogous cyclic definitions. Define
\[
 L_a=9a^4+16D,\qquad K_a=2(A_b+A_c)L_a-a^2(A_b-A_c)^2 .
\]

We first claim $(m_b+m_c)^2\ge L_a/(4a^2)$; this is \cite[Lemma 2.3]{r36reiwa:liu}. For completeness place $B=(0,0)$, $C=(a,0)$, $A=(w,h)$, where $h=2\Delta/a$. The vectors $(a+w,h)/2$ and $(2a-w,h)/2$ have lengths $m_b$ and $m_c$ respectively (the second is a horizontal-coordinate reflection of the median vector). Their sum has squared length $9a^2/4+h^2$. The triangle inequality proves the claim for all projections $w$, including obtuse triangles.

The useful refinement is obtained from
\[
 (m_b+m_c)^2=2(m_b^2+m_c^2)-\frac{(m_b^2-m_c^2)^2}{(m_b+m_c)^2}.
\]
Using the first lower bound in the denominator of the subtracted term now yields
\[
 (m_b+m_c)^2\ \ge\ \frac{K_a}{4L_a}.
\]
We will verify $K_a,K_b,K_c>0$ in the coefficient calculation below. Put
\[
 C=\frac1R+\frac1{2r}=\frac{8D+sJ}{2\sqrt D J},\qquad
 U=(8D+sJ)^2,\qquad V=576DJ^2,\qquad
 t_a=\frac{C^2K_a}{144L_a}=\frac{UK_a}{VL_a}.
\]
It follows that
\[
 \frac{216}{C^3(m_b+m_c)^3}\le t_a^{-3/2}
 \le \frac{t_a^{-1}+t_a^{-2}}{2},
\]
the last inequality being arithmetic--geometric mean. Thus it suffices to prove $\sum_{\rm cyc}(t_a^{-1}+t_a^{-2})\le6$. Multiplication by $U^2(K_aK_bK_c)^2$ reduces this to
\begin{equation}\label{r36reiwa:poly}
 G=6U^2(K_aK_bK_c)^2-
 \sum_{\rm cyc}VL_a(UK_a+VL_a)(K_bK_c)^2\ge0.
\end{equation}

By permutation symmetry and homogeneity of the target we may take $x=1$, $y=1+u$, $z=1+u+v$ with $u,v\ge0$. In these variables each $K$ has only nonnegative coefficients and constant term $9216$; $G$ has degree 52 with 1380 nonzero coefficients, all positive. Its constant term is zero and its pure quadratic coefficients are positive. The following exact polynomial-arithmetic calculation establishes these assertions and is included so that no separate data file is needed.
\begin{small}
\begin{Verbatim}
from sympy import ZZ
from sympy.polys.rings import ring
B,u,v = ring('u,v',ZZ)
x,y,z = B.one,1+u,1+u+v
s=x+y+z; sides=y+z,z+x,x+y
J=sides[0]*sides[1]*sides[2]; D=s*x*y*z
U=(8*D+s*J)**2; V=576*D*J**2
A=[4*s*X+(Y-Z)**2 for X,Y,Z in [(x,y,z),(y,z,x),(z,x,y)]]
K=[]; L=[]
for i,a in enumerate(sides):
    h=9*a**4+16*D
    L.append(h)
    K.append(2*(A[i-1]+A[i-2])*h-a**2*(A[i-1]-A[i-2])**2)
for h in K:
    assert min(h.values())>0 and h[0,0]==9216
G=6*U**2*(K[0]*K[1]*K[2])**2
for i in range(3):
    G-=V*L[i]*(U*K[i]+V*L[i])*(K[i-1]*K[i-2])**2
assert len(G)==1380 and max(i+j for i,j in G)==52
assert min(G.values())==1057916215296
assert G.get((0,0),0)==0 and G[2,0]>0 and G[0,2]>0
\end{Verbatim}

\end{small}
In particular \eqref{r36reiwa:poly} is strict unless $u=v=0$. An equilateral triangle has $C(m_b+m_c)=6$, giving equality in the original statement.
\end{proof}
\begin{thebibliography}{1}
\bibitem{r36reiwa:liu} J.~Liu, Two inequalities for the medians of a triangle, \emph{Turkish J. Inequal.} \textbf{8} (2024), no.~1, 35--47. \href{https://tjinequality.com/articles/08-01-002.pdf}{Publisher PDF}, Conjecture 1, p.~46.
\end{thebibliography}

\par\endgroup

% Source: /workspace/scratch/81914cec2c78/work/archive/Tournament_Round_36_Final/Tournament_Round_36/round36/keio/paper.tex
% Edited source: /workspace/scratch/81914cec2c78/work/audit_34_41/cleaned/r36_keio.tex
\begingroup
\newcommand{\Cyc}{\sum_{\rm cyc}}
\Needspace{9\baselineskip}
\section{\texorpdfstring{Two-term domination in three weighted Walker inequalities}{Two-term domination in three weighted Walker inequalities}}\label{section:r36keio}
\noindent{\small\sffamily\color{muted}Round 36 \enspace / \enspace Keio}\par\medskip


\begin{sectionabstract}
Jian Liu conjectured three weighted cyclic inequalities for acute triangles, with exponents on three unbounded real intervals. We prove each by showing that the term indexed by the shortest side already covers the entire negative term. The deciding comparison reduces to a right triangle; in the last two families it becomes a monotone rational function in a half-angle parameter. The proofs require no numerical estimates.
\end{sectionabstract}

Write $L_a=(a-b)(a-c)(b^2+c^2-a^2)$, with $L_b,L_c$ defined cyclically. The three conjectures appear in Section~7 of [1], p.~6671. The original includes special instances of the first family: $k=0$ in (5.12), $k=1$ in (5.9), and $k=2$ in (7.2), the last from Walker's inequality.

\begin{theorem}[Liu's Conjecture 1, equation (7.3)]
For any acute triangle with side lengths $a,b,c$ and every real $k\le4$,
\[
                     \Cyc L_a a^k\ge0.
\]
\end{theorem}
\begin{proof}
By symmetry and homogeneity, order the sides $a\ge b\ge c$ and set $b=1$, $c=z$. Then $z\le1$ and $a\le A:=\sqrt{1+z^2}$ by acuteness. First suppose $a>1>z$. All cosine factors are positive, so $L_a,L_c>0$ and $L_b<0$. Cancelling the common gap $1-z$ gives
\[
 P(a):=\frac{L_c}{-L_b}=
 \frac{a-z}{a-1}\,\frac{a^2+1-z^2}{a^2+z^2-1}.
\]
The factors equal $1+(1-z)/(a-1)$ and $1+2(1-z^2)/(a^2+z^2-1)$; both decrease with $a$. Thus $P(a)\ge P(A)=(A-z)/((A-1)z^2)$. Since $z^k\ge z^4$, it suffices that $z^2(A-z)\ge A-1$, or $A\le(1+z+z^2)/(1+z)$. Both sides are positive, and
\[
       (1+z+z^2)^2-(1+z)^2(1+z^2)=z^2>0.
\]
So $L_cz^k>-L_b$. If two sides tie, the corresponding two terms vanish and the third is nonnegative; it is positive unless all three sides are equal. This proves the claim, including equality.
\end{proof}


\begin{theorem}[Liu's Conjecture 2, equation (7.4)]
For every acute triangle and every real $k\ge-10$,
\[
                        \Cyc L_a(b+c)^k\ge0 .
\]
\end{theorem}
\begin{theorem}[Liu's Conjecture 3, equation (7.5)]
For every acute triangle and every real $k\ge-3$,
\[
                        \Cyc L_a(b+c-a)^k\ge0 .
\]
\end{theorem}
We prove both statements by the same reduction. Use the normalization in the proof of the first Walker inequality above. The tied-side cases are identical since the triangle inequalities keep all weights positive. For scalene triangles define
\[
 (q,h)=\left(\frac{a+z}{a+1},10\right)
    \quad\text{or}\quad
 (q,h)=\left(\frac{a+z-1}{a+1-z},3\right),
\]
respectively. In either case $0<q\le1$. The shortest-to-middle weight ratio equals $q^{-k}\ge q^h$ on $k\ge-h$. We need only prove $P(a)q(a)^h\ge1$. As the second factor of $P$ already decreases in $a$, it suffices to show that $(a-z)q^h/(a-1)$ also decreases, and to compare at $a=A$.

\paragraph{The monotonicity step.}
Logarithmic differentiation reduces the two derivative signs to the respective inequalities
\begin{align*}
 (a+z)(a+1)&\ge10(a-z)(a-1),\\
 (a+z-1)(a+1-z)&\ge6(a-z)(a-1).
\end{align*}
Indeed, the derivative of $\log((a-z)/(a-1))$ is $-(1-z)/((a-z)(a-1))$, while that of $h\log q$ is $10(1-z)/((a+z)(a+1))$ or $6(1-z)/((a+z-1)(a+1-z))$, respectively.

Put $x=a-1\le z^2/(a+1)\le z^2/2$. In the first inequality the right side is at most $5z^2(1-z/2)$, since $a-z\le1-z+z^2/2\le1-z/2$. This is at most $5z/2$ because $z(1-z/2)\le1/2$ for $0\le z\le1$. The left side is at least $2(1+z)>5z/2$. For the second inequality, the left side minus the right is
\[
              2z-z^2-4x+6xz-5x^2.
\]
This is concave in $x$. At the left endpoint $x=0$ it is positive; at the right endpoint $x=z^2/2$ it equals
\[
 z(2-3z+3z^2-5z^3/4)\ \ge\ z(2-3z+7z^2/4)>0,
\]
where the last quadratic has negative discriminant. The two derivative signs follow.


\paragraph{The endpoint step.}
At $a=A$ write $p=A+z$, which lies in $(1,1+\sqrt2]$. Since $A-z=1/p$, we have
\[
 A=\tfrac12(p+1/p),\quad z=\tfrac12(p-1/p),\qquad
 \frac1{P(A)}=\frac{(A-1)z^2}{A-z}
 =\frac{(p-1)^4(p+1)^2}{8p^2}.
\]
For the second and third Walker inequalities the corresponding values of $q(A)$ are $2p^2/(p+1)^2$ and $p(p-1)/(p+1)$. Therefore the quantity $1/(P(A)q(A)^h)$ to be bounded above by one is respectively
\[
    H_2(p)=\frac{(p-1)^4(1+1/p)^{22}}{8192},
    \qquad
    H_3(p)=\frac{(p-1)(1+1/p)^5}{8}.
\]
These functions increase on this interval: with common positive denominator $(p-1)p(p+1)$, the logarithmic derivatives have numerators
\[
         4p^2-18p+22 \qquad\text{and}\qquad p^2-4p+5.
\]
Both quadratics have positive leading coefficient and negative discriminant. At $p=1+\sqrt2$ one has $z=1$, so both ratios $q(A)$ and $P(A)$ in their limiting formulas are exactly one. It follows that $H_i(p)\le1$ throughout. Since $P(a)q(a)^h\ge P(A)q(A)^h$, the shortest-side term dominates the negative middle-side term in each family. The positive longest-side term gives strict inequality in every scalene acute triangle. The tied-side calculation completes both proofs. \hfill$\square$

\paragraph{Scope.}
All three conclusions concern ordinary real powers of positive lengths. They require no integrality, parity or avoidance of an endpoint exponent. The proof compares with a right triangle without including degenerate triangles or asserting optimality of the exponents.

\begin{thebibliography}{1}
\bibitem{r36keio:liu} J. Liu, Further generalization of Walker's inequality in acute triangles and its applications, \emph{AIMS Mathematics} \textbf{5} (2020), 6657--6672. \url{https://doi.org/10.3934/math.2020428}.
\end{thebibliography}

\par\endgroup

% Source: /workspace/scratch/81914cec2c78/work/archive/Tournament_Round_37_Final/Tournament_Round_37/round37/meiji/paper.tex
% Edited source: /workspace/scratch/81914cec2c78/work/audit_34_41/cleaned/r37_meiji.tex
\begingroup
\newcommand{\cyc}{\mathop{\sum_{\rm cyc}}}
\Needspace{9\baselineskip}
\section{\texorpdfstring{The sharp Euler-deficit coefficient for reciprocal sums of pairs of medians}{The sharp Euler-deficit coefficient for reciprocal sums of pairs of medians}}\label{section:r37meiji}
\noindent{\small\sffamily\color{muted}Round 37 \enspace / \enspace Meiji}\par\medskip


\begin{sectionabstract}
For a nonobtuse triangle, how much of the Euler deficit $R-2r$ can be added to $2r$ while leaving its reciprocal at least as large as the sum of the three reciprocals of pairs of medians? Liu proved that the coefficient $1/2$ works and conjectured the exact maximum. We prove his proposed coefficient, distinguish the two equality shapes, and identify the strictly acute sharpness limit. The proof uses a shape square on which the circumradius and inradius stay smooth at the degenerate edge, followed by an exact interval certificate.
\end{sectionabstract}
\begin{theorem}
Let $m_a,m_b,m_c$ be the medians and $R,r$ the circumradius and inradius. The largest $k>0$ for which
\[
 \cyc\frac1{m_b+m_c}\le \frac1{2r+k(R-2r)}
 \tag{1}
\]
holds for every acute triangle is
\[
 K=\frac{23\sqrt{10}-42}{41}+\frac{92\sqrt5-166\sqrt2}{123}.
 \tag{2}
\]
It holds also on the nondegenerate nonobtuse closure. At $k=K$, equality holds there precisely for the equilateral and right isosceles shapes. For a nonequilateral triangle the shape-dependent bound on $k$ attains its minimum on that closure only at the right isosceles shape.
\end{theorem}

This is Conjecture 1 in Liu \cite[p.~20, (4.1)]{r37meiji:liu}. The source's Theorem 1 treats $k=1/2$ and does not give (2). A right angle is not strictly acute: the bound on the coefficient on the strictly acute domain is a limit. The equilateral triangle, in contrast, satisfies (1) identically in $k$ and places no restriction on the coefficient.

\paragraph{Shape coordinates.}
Relabel and scale so that $a\le b=1\le c$. Put $x=1-a$ and $y=(c^2-1)/a^2$. Nonobtuseness is exactly $0\le y\le1$. Conversely all $0<a\le1$ and $0\le y\le1$ give genuine triangles since $\sqrt{1+a^2}<1+a$. We work on the full closed square $0\le x,y\le1$. Write
\[
 a=1-x,\qquad c=\sqrt{1+a^2y},\qquad
 U=\sqrt{4-a^2(1-y)^2}.
\]
The cosine law at the angle opposite $c$ gives $4\Delta=aU$ for area $\Delta$. Cancellation gives the following smooth continuation to $x=1$:
\[
 R=c/U,\qquad r=\frac{aU}{2(1+a+c)}.
\]
Apollonius's median identities give
\[
 2m_a=\sqrt{4+a^2(2y-1)},\quad
 2m_b=\sqrt{1+a^2(2+2y)},\quad
 2m_c=\sqrt{1+a^2(2-y)}.
\]
Let $A=m_a+m_b$, $B=m_b+m_c$, $C=m_c+m_a$. Since $R\ge 2r$ and $K>0$, clearing denominators in (1) gives exactly
\[
 F(x,y)=ABC-\{2r+K(R-2r)\}(AB+BC+CA)\ \ge 0. \tag{3}
\]

\paragraph{Two different zeros.}
At the equilateral point $(0,0)$, (1) is an identity: its two sides are $\sqrt{3}$ for side length $1$. Consequently $F=0$ there. As a function of the original three sides $F$ is symmetric and homogeneous of degree three. Its partial derivatives at the equilateral shape are equal, hence zero by Euler's identity. Thus both chart derivatives also vanish.

At the right isosceles point $(0,1)$ the sides are $1,1,\sqrt2$. Here
\[
 R=\sqrt2/2,\quad r=1-\sqrt2/2,\quad
 Q:=\cyc\frac1{m_b+m_c}=\frac1{\sqrt5}+\frac4{\sqrt5+\sqrt2}.
\]
Rationalizing $(1/Q-2r)/(R-2r)$ gives exactly (2), hence $F(0,1)=0$. On the right edge let $H(a)=F(1-a,1)$. Extending the edge formula to both sides of $a=1$, interchanging the two legs and rescaling gives $H(a)=a^3 H(1/a)$. Differentiation at $a=1$ gives $2H'(1)=3H(1)=0$. Thus $F_x(0,1)=0$ but we do not need $F_y$ to vanish.

\paragraph{Positivity certificate.}
Put $e=1/8192$. The integer verification below proves the following rational lower bounds, in the indicated coordinates:
\[
\begin{array}{c|c|r}
\text{region}&\text{quantity}&\text{lower bound after multiplication by }1024\\ \hline
{[0,e]}^2&F_{xx},F_{xy},F_{yy}&1021,\ 244,\ 252\\
{[0,e]}\times\{1\}&F_{xx}&1984\\
{[0,e]}\times[1-e,1]&-F_y&137
\end{array}
\]
On the first patch, for a nonzero vector $(x,y)$ pointing into the patch, these three Hessian bounds make the second derivative along $t\mapsto(tx,ty)$ strictly positive. Integration from the equilateral zero proves strict positivity off the corner. On the second patch, the edge second derivative and the vanishing tangential derivative give $F(x,1)>0$ for $x>0$. The sign $F_y<0$ gives $F(x,y)>F(x,1)$ for $y<1$ in the patch.

On the remaining boxes of a dyadic subdivision, an interval for $F$ or its centered enclosure is strictly positive:
\[
 F(c_0)+[\partial_x F](B_0)([x]-c_{0x})
       +[\partial_y F](B_0)([y]-c_{0y}),\tag{4}
\]
where $B_0=[x]\times[y]$ and $c_0$ is its center. This follows by integration along the line segment from the center; there is no approximation of the domain by a grid of points. Each failed test bisects the longer coordinate. Starting with the entire square yields 4 directly positive boxes, 3,016 positive by (4), and 2 local boxes. The maximum depth is 35. These tests prove (3) throughout the closed square, with just the two described zeros.

For $0<k\le K$, (1) follows by $R\ge 2r$. Conversely $R>2r$ at the right isosceles triangle. Thus for every $k>K$, (1) fails strictly there and, by continuity, also at nearby strictly acute isosceles triangles. For a nonequilateral triangle Euler's inequality is strict, so the individual upper bound $(1/Q-2r)/(R-2r)$ equals $K$ only at the right isosceles shape by (3). This proves the theorem.

\begin{thebibliography}{9}
\bibitem{r37meiji:liu} J. Liu, Two inequalities involving circumradius, inradius and medians of an acute triangle, \emph{Adv. Inequal. Appl.} (2023), Article 1.
\href{https://doi.org/10.28919/aia/8232}{doi:10.28919/aia/8232}.
\end{thebibliography}

\paragraph{Complete verifier (Python 3).}
Intervals encode $[l/2^{64},h/2^{64}]$. Integer floor and ceiling implement directed rounding, and integer square root supplies root enclosures. The first-order class stores partials; the second-order class stores two directional derivatives and their mixed derivative. The local tests print the integer floors of the certified lower bounds after scaling by $1024$. An exception for nonpositive square-root input triggers subdivision; unsafe division or exhausted depth instead stops the computation. Run with assertions enabled.
\begingroup\fontsize{7}{8.3}\selectfont
\begin{Verbatim}
from math import isqrt
from collections import Counter
S=1<<64
class I:
 __slots__='l','h'
 def __init__(self,l,h=None):
  if h is None:self.l=self.h=l*S
  else:self.l=l;self.h=h
 def __add__(x,y):
  if not isinstance(y,I):y=I(y)
  return I(x.l+y.l,x.h+y.h)
 __radd__=__add__
 def __neg__(x):return I(-x.h,-x.l)
 def __sub__(x,y):return x+-y
 def __rsub__(x,y):return -x+y
 def __mul__(x,y):
  if not isinstance(y,I):return I(min(x.l*y,x.h*y),max(x.l*y,x.h*y))
  if x.l>=0 and y.l>=0:return I(x.l*y.l//S,-(-x.h*y.h//S))
  p=(x.l*y.l,x.l*y.h,x.h*y.l,x.h*y.h);return I(min(p)//S,-(-max(p)//S))
 __rmul__=__mul__
 def __truediv__(x,y):
  if not isinstance(y,I):
   assert y>0;return I(x.l//y,-(-x.h//y))
  assert y.l>0, 'unsafe denominator'
  return x*I(S*S//y.h,-(-S*S//y.l))
 def rt(x):
  if x.l<=0:raise ValueError
  return I(isqrt(x.l*S),isqrt(x.h*S)+1)
class J:
 __slots__='v','g'
 def __init__(self,v,g):
  self.v=v if isinstance(v,I) else I(v);self.g=[d if isinstance(d,I) else I(d) for d in g]
 def __add__(x,y):
  if not isinstance(y,J):return J(x.v+y,x.g)
  return J(x.v+y.v,[p+q for p,q in zip(x.g,y.g)])
 __radd__=__add__
 def __neg__(x):return J(-x.v,[-p for p in x.g])
 def __sub__(x,y):return x+-y
 def __rsub__(x,y):return -x+y
 def __mul__(x,y):
  if not isinstance(y,J):return J(x.v*y,[p*y for p in x.g])
  return J(x.v*y.v,[p*y.v+q*x.v for p,q in zip(x.g,y.g)])
 __rmul__=__mul__
 def __truediv__(x,y):
  if not isinstance(y,J):return J(x.v/y,[p/y for p in x.g])
  v=x.v/y.v;return J(v,[(p-v*q)/y.v for p,q in zip(x.g,y.g)])
 def rt(x):
  y=x.v.rt();return J(y,[p/(2*y) for p in x.g])
class B:
 def __init__(self,v,g=0,h=0,j=0):
  self.v,self.g,self.h,self.j=[i if isinstance(i,I) else I(i) for i in (v,g,h,j)]
 def __add__(x,y):
  if not isinstance(y,B):y=B(y)
  return B(x.v+y.v,x.g+y.g,x.h+y.h,x.j+y.j)
 __radd__=__add__
 def __neg__(x):return B(-x.v,-x.g,-x.h,-x.j)
 def __sub__(x,y):return x+-y
 def __rsub__(x,y):return -x+y
 def __mul__(x,y):
  if not isinstance(y,B):y=B(y)
  return B(x.v*y.v,x.g*y.v+x.v*y.g,x.h*y.v+x.v*y.h,x.j*y.v+x.g*y.h+x.h*y.g+x.v*y.j)
 __rmul__=__mul__
 def chain(x,y,a,b):return B(y,a*x.g,a*x.h,a*x.j+b*x.g*x.h)
 def __truediv__(x,y):
  if not isinstance(y,B):return B(x.v/y,x.g/y,x.h/y,x.j/y)
  r=I(1)/y.v;return x*y.chain(r,-r*r,2*r*r*r)
 def rt(x):
  r=x.v.rt();d=I(1)/(2*r);return x.chain(r,d,-d/(2*x.v))

K=(23*I(10).rt()-42)/41+(92*I(5).rt()-166*I(2).rt())/123
def f(x,y):
 a=1-x; c=(1+a*a*y).rt()
 # 4*area/a from Heron with c^2=1+a^2*y
 U=(4-a*a*(1-y)*(1-y)).rt()
 R=c/U; r=a*U/(2*(1+a+c))
 ma=(4+a*a*(2*y-1)).rt()/2
 mb=(1+a*a*(2+2*y)).rt()/2
 mc=(1+a*a*(2-y)).rt()/2
 A=ma+mb;Bv=mb+mc;C=mc+ma
 return A*Bv*C-(2*r+(R-2*r)*K)*(A*Bv+Bv*C+C*A)
def prove(eps=8192):
 counts=Counter();stack=[(0,S,0,S,0)]
 while stack:
  l,h,u,v,d=stack.pop();counts['max depth']=max(d,counts['max depth'])
  if h<=S//eps and (v<=S//eps or u>=S-S//eps):
   counts['local']+=1;continue
  left=I(l,h);right=I(u,v)
  try:
   if f(left,right).l>0:counts['value']+=1;continue
   midx=(l+h)//2;midz=(u+v)//2
   val=f(I(midx,midx),I(midz,midz))
   jet=f(J(left,[1,0]),J(right,[0,1]))
   if (val+jet.g[0]*I(l-midx,h-midx)+jet.g[1]*I(u-midz,v-midz)).l>0:counts['mean']+=1;continue
  except ValueError:pass
  assert d<60,(d,l/S,h/S,u/S,v/S,counts)
  if h-l>=v-u:
   w=(l+h)//2;stack.extend([(l,w,u,v,d+1),(w,h,u,v,d+1)])
  else:
   w=(u+v)//2;stack.extend([(l,h,u,w,d+1),(l,h,w,v,d+1)])
 print(counts,flush=True)
def local(eps):
 box=I(0,S//eps)
 for p,q in [(0,0),(0,1),(1,1)]:
  j=f(B(box,p==0,q==0),B(box,p==1,q==1)).j
  print('equilateral Hessian lower / 1024:',p,q,j.l*1024//S,flush=True)
  assert j.l>0
 # symmetry of right triangle function at equal legs gives F_x(0,1)=0
 j=f(B(box,1,1),B(I(1),0,0)).j
 print('right edge Fxx lower /1024:',j.l*1024//S,flush=True)
 assert j.l>0
 j=f(J(box,[0,0]),J(I(S-S//eps,S),[0,1])).g[1]
 print('right patch -Fy lower /1024:',-j.h*1024//S,flush=True)
 assert j.h<0
if __name__=='__main__':
 local(8192);prove(8192)
\end{Verbatim}
\endgroup

\par\endgroup

% Source: /workspace/scratch/81914cec2c78/work/archive/Tournament_Round_37_Final/Tournament_Round_37/round37/taisho/paper.tex
% Edited source: /workspace/scratch/81914cec2c78/work/audit_34_41/cleaned/r37_taisho.tex
\begingroup
\newcommand{\at}{\operatorname{arctan}}
\newcommand{\ath}{\operatorname{arctanh}}
\Needspace{9\baselineskip}
\section{\texorpdfstring{Two poles and the asymptotic formula for OEIS A012200}{Two poles and the asymptotic formula for OEIS A012200}}\label{section:r37taisho}
\noindent{\small\sffamily\color{muted}Round 37 \enspace / \enspace Taisho}\par\medskip

V\'aclav Kote\v{s}ovec proposed in 2015 [1] an asymptotic for the coefficients in
\[
 F(z)=\ath(\at(\arcsin z))
     =\sum_{n\ge0}\frac{a(n)}{(2n+1)!}\,z^{2n+1}.
\]
The dominant singularities have modulus $r=\sin(\tan 1)=0.9999103740\ldots$.
Their proximity to the inverse sine's branch points makes a short coefficient table a poor guide. We prove the conjecture and quantify the effect of those branch points.
\subsection*{The conjecture and its proof}
\textbf{Theorem 1.} For every fixed $R$ satisfying $r<R<1$,
\[
 \frac{a(n)}{(2n)!}=r^{-2n-1}+O_R(R^{-2n}).
 \qquad\text{In particular,}\qquad
 a(n)\sim\frac{(2n)!}{(\sin(\tan 1))^{2n+1}} .
\]
\textit{Proof.} Let $h(z)=\arcsin z$ be the branch vanishing at $0$ in $|z|<1$.
Since $\sin h(z)=z$, an equality $h(z)=\pm i$ would imply $z=\pm i\sinh 1$, outside this disk. Therefore the germ $b(z)=\at(h(z))$ extends analytically to the disk: take the primitive of $h'/(1+h^2)$ with $b(0)=0$. The identity $\tan b=h$ persists by continuation as an identity of meromorphic functions.

The only possible points at which $b(z)=\pm1$ are $z=\pm r$, by applying $\tan$ and then $\sin$. Conversely these points do attain those values, since $0<\tan 1<\pi/2$. For example, the alternating Taylor bounds give
$\sin 1<101/120$, $\cos1>389/720$, and $606/389<1.56<\pi/2$.
In particular $0<r<1$. At each point $b'=h'/(1+h^2)\ne0$, since
$h'=1/\sqrt{1-z^2}$. The meromorphic continuation
\[
 F'(z)=\frac{b'(z)}{1-b(z)^2}
\]
thus has precisely two poles in $|z|<1$: both are simple, with residues $-1/2$ at $r$ and $+1/2$ at $-r$. The rational function $r/(r^2-z^2)$ has exactly these principal parts. Cauchy's estimate on $|z|=R$ gives
\[
 [z^{2n}]\!\left(F'(z)-\frac{r}{r^2-z^2}\right)=O_R(R^{-2n}).
\]
Now $[z^{2n}]F'=a(n)/(2n)!$ and $[z^{2n}]r/(r^2-z^2)=r^{-2n-1}$. The relative error tends to zero exponentially for each such $R$. \hfill$\square$


\subsection*{The first correction}
Put
\[
 g(s)=\frac{1}{(1+s^2)(1-(\at s)^2)},\qquad
 \kappa=g(\pi/2)=-37.04691538\ldots.
\]
\textbf{Theorem 2.} The normalized coefficients satisfy
\[
 \frac{a(n)}{(2n)!}
   =r^{-2n-1}+\kappa\,4^{-n}\binom{2n}{n}+O(n^{-3/2}).
\]
Thus the first correction is negative and of absolute order $1/\sqrt{n}$.

\textit{Proof.} Choose $1<\rho<\sinh1$ and slit $|z|<\rho$ along $[1,\rho)$ and its negative. This is a simply connected domain on which $h$ continues as the primitive of $(1-z^2)^{-1/2}$ with value $0$ at the origin. Again $\sin h=z$ rules out $h=\pm i$, and the primitive defining $b$ continues. The preimage test in Theorem 1 once more rules out additional poles. Because $F'$ is even as a meromorphic continuation, the function
\[
 P(t)=F'(\sqrt t)-\frac{r}{r^2-t}
\]
is single-valued and holomorphic in $|t|<\rho^2$ slit from $1$ to the right. The removable singularity at $t=r^2$ was handled in Theorem 1. In a slit neighborhood of $t=1$, take the square root positive just to the left of $1$. The real inverse sine identity followed by analytic continuation gives
\[
 \arcsin\sqrt{t}=\frac{\pi}{2}-\arcsin\sqrt{1-t}.
\]
The function $g$ is analytic near $\pi/2$ since $\at(\pi/2)>1$. Expanding $g$ locally and using the chain rule gives
\[
 P(t)=\frac{g(\arcsin\sqrt t)}{\sqrt{1-t}}-\frac{r}{r^2-t}
      =\kappa(1-t)^{-1/2}+C+O((1-t)^{1/2}).
\]
The slit continuation permits the standard big-$O$ transfer theorem [2], giving coefficients $O(n^{-3/2})$ for the last term. The constant does not contribute for $n>0$, and $[t^n](1-t)^{-1/2}=4^{-n}\binom{2n}{n}$. This proves the formula. \hfill$\square$

\textbf{Verification and scope.} Exact rational Taylor arithmetic reproduces the first 14 tabulated integers; these finite values are not used in the proofs. The conjecture appears in a public sequence entry rather than a journal. Searches by sequence number and exact expression found the conjectural entry but not a proof; no exhaustive priority claim is made.

{\small\raggedright
[1] V. Kote\v{s}ovec, February 7, 2015, conjectural formula in P. Demichel's
\href{https://oeis.org/A012200/internal}{OEIS A012200}, \emph{The On-Line Encyclopedia of Integer Sequences}.

[2] P. Flajolet and A. Odlyzko, \emph{Singularity analysis of generating functions}, SIAM J. Discrete Math. \textbf{3} (1990), 216--240.
\href{https://doi.org/10.1137/0403019}{DOI: 10.1137/0403019}.\par}

\par\endgroup

% Source: /workspace/scratch/81914cec2c78/work/archive/Tournament_Round_37_Final/Tournament_Round_37/round37/showa/paper.tex
% Edited source: /workspace/scratch/81914cec2c78/work/audit_34_41/cleaned/r37_showa.tex
\begingroup
\newcommand{\E}{\mathbb E}
\newcommand{\Var}{\operatorname{Var}}
\Needspace{9\baselineskip}
\section{\texorpdfstring{Endpoint matching, curvature, and coefficient tails Three previously covered questions on means and elliptic integrals}{Endpoint matching, curvature, and coefficient tails Three previously covered questions on means and elliptic integrals}}\label{section:r37showa}
\noindent{\small\sffamily\color{muted}Round 37 \enspace / \enspace Showa}\par\medskip


\begin{sectionabstract}
We record a hyperbolic-coordinate comparison for a sharp power mean bound, a short curvature argument for Toader means, and a tail-coefficient argument for an elliptic quotient. Each target was already covered by public prior work. The arguments are recorded for reuse, not as new resolutions of open problems.
\end{sectionabstract}

\subsection{A sharp comparison for symmetric beta averages}
Barnard, Richards, and Tiedeman [1, Conjecture II] ask for the sharp lower
power-mean bound on a Gauss hypergeometric function. For $a>1$, $b>0$ put
\[
 m=\frac{\Gamma(2b)\Gamma(b+a)}{\Gamma(b)\Gamma(2b+a)},\qquad
 p_E=-\frac{a\log 2}{\log m}.
\]
Then, uniformly for $0<r<1$, the inequality
\[
 M_p(1,1-r)\le {}_2F_1(-a,b;2b;r)^{1/a}
\]
holds precisely when $p\le p_E$. Here $M_p(x,y)=((x^p+y^p)/2)^{1/p}$ with the geometric convention at $p=0$. This theorem and the nonsymmetric companion are covered by Deng, Peng, and Zou [2] using a weighted Wronskian. Here is a derivation reducing the symmetric residual to a single ratio of hyperbolic sines.

For $X\sim\operatorname{Beta}(b,b)$, $U=2X-1$, write $r=1-e^{-2t}$. By
the beta integral and symmetry the hypergeometric function is
$e^{-at}f(t)$, where $f(t)=\E h^a$, $h=\cosh t+U\sinh t$. The $a$th power of the comparison mean at $p=p_E$ is $e^{-at}g(t)$, where
$g(t)=\cosh(p_Et)^{a/p_E}$. Define
\[
 Ly=y''+2b\coth(t)y'-a(a+2b)y .
\]
Indeed $Lf=0$: use $h^2-(\partial_t h)^2=1-U^2$ and
$\coth(t)\partial_t h-h=U/\sinh t$, followed by the identity
\[
 (a-1)\sinh(t)\E[(1-U^2)h^{a-2}]=2b\E[Uh^{a-1}].
\]
It follows by integrating the derivative of $(1-u^2)^b h^{a-1}$ on $(-1,1)$; boundary values vanish since $b>0$. Differentiation at fixed positive $t$ is justified by $h\ge e^{-t}$.

Now $m=\E X^a>2^{-a}$ by strict Jensen, so $p=p_E>1$.
Both $g(0)=f(0)=1$ and $g/f\to1$ at infinity: $e^{-at}f\to m$ by bounded convergence, and $e^{-at}g\to2^{-a/p}=m$. Direct differentiation gives
\[
 \frac{Lg}{a g\,\mathrm{sech}^2(pt)}
 =p-a+b\left(\frac{\sinh((2p-1)t)}{\sinh t}-1\right).
\]
The ratio inside parentheses is strictly increasing since $z\coth z$ is
strictly increasing for $z>0$, and $2p-1>1$. Hence the residual can
only change sign from minus to plus. The weighted Wronskian
$W=(\sinh t)^{2b}(fg'-gf')$ has $W(0+)=0$ and
$W'=(\sinh t)^{2b} f Lg$. Thus $W$ can also change sign only once,
from negative to positive. Since $(g/f)'$ has its sign and $g/f$ has identical endpoint limits, both signs must occur (the residual precludes constancy). Hence $g/f<1$ throughout. The strict increase of the power means settles $p<p_E$. For $p>p_E>1$, the endpoint limit at $r=1$ violates the inequality.

\subsection{A real interpolation for the power-type Toader sequence}
Zhao, Wang, Dai, and Chu [3, Conjecture 2.6] ask for strict increase in
nonzero integer order; their direct comparison shows it for
$-43\le n\le43$, $n\ne0$. For fixed distinct $a,b>0$ define $w=\cos^2\theta$,
where $\theta$ is uniform on $[0,\pi/2]$. For all real $s$ let
\[
 Z(s)=\E[(w a^s+(1-w)b^s)^{1/2}],\quad H(s)=\log Z(s),
 \qquad T_s=e^{2H(s)/s}\quad(s\ne0).
\]
With $\ell(s)=\tfrac12\log(wa^s+(1-w)b^s)$,
\[
 \ell''=\frac{w(1-w)(ab)^s(\log a-\log b)^2}
 {2(wa^s+(1-w)b^s)^2}>0\quad\text{almost surely}.
\]
For the probability tilted by $e^\ell/Z$, differentiation under the integral
gives $H''=\E_s\ell''+\Var_s(\ell')>0$; the derivatives are bounded on
each compact parameter interval. Since $H(0)=0$, $H(s)/s$ is the
strictly increasing secant slope through zero, continuously extended by
$H'(0)$ there. As $\E w=1/2$, $T_0=\sqrt{ab}$. This also covers the
jump from negative to positive integer orders. The result is already a
special case of Tobey [4, Theorem 2]: with arcsine-distributed $w$, it is
the diagonal $L(s,s/2)$ in his coordinatewise increasing two-order mean.

\subsection{An elliptic quotient via coefficient tails}
In [5, Remark 8(iii)], Zhao and Yang ask about absolute monotonicity on
$(0,1)$ of
\[
 G(x)=\frac{K(\sqrt{x})/L(x)-1}{1-x},\quad
 L(x)=\log 4-\tfrac12\log(1-x).
\]
Here $K(k)=\int_0^{\pi/2}(1-k^2\sin^2 t)^{-1/2}\,dt$.
The general zero-balanced hypergeometric problem is treated in [6].
The following elementary argument gives the strict coefficient signs directly.

Write $W_n=4^{-n}\binom{2n}{n}$,
$\delta_n=1/(2n)-(\pi/2)W_n^2$ for $n\ge1$, and
$d_n=\sum_{k>n}\delta_k$ for $n\ge0$.
We have $\pi(n+1/4)W_n^2<1$: the left side increases to 1 by the
Wallis asymptotic, with the sign of its consecutive difference given by
$(n+5/4)(n+1/2)^2-(n+1/4)(n+1)^2=1/16$.
In particular $\delta_n>0$. Using $K(\sqrt{x})=(\pi/2)\sum W_n^2 x^n$
and the classical asymptotic $K(\sqrt{x})-L(x)\to0$ as $x\uparrow1$, monotone convergence yields $d_0=\pi/2-\log 4$ and
\[
 (K(\sqrt{x})-L(x))/(1-x)=\sum_{n\ge0}d_nx^n .
\]
For $k\ge1$ set $C_k=2k(k-1)\delta_k$. Its increment is positive:
\[
 C_{k+1}-C_k
 =1-\pi W_k^2\frac{k(4k+5)}{4(k+1)}
 >1-\frac{k(4k+5)}{(4k+1)(k+1)}>0.
\]
Thus for $n\ge1$ telescoping $1/(k(k-1))$ gives
$d_{n+1}>C_{n+1}/(2(n+1))=n\delta_{n+1}$. In particular
$(n+1)d_{n+1}>nd_n$.

Let $l_n=[x^n]L(x)$, so $l_0=\log4$, $l_n=1/(2n)$ for $n\ge1$.
Then $d_n/l_n$ strictly increases: the remaining first step
$2d_1>d_0/\log4$ follows upon inserting $d_1=5\pi/8-\log4-1/2$
(e.g. $3.14159<\pi<3.14160$, $1.38629<\log4<1.38630$).
The $l_n$ are log-convex, including $l_1^2<l_0 l_2$.
Write $(\sum d_nx^n)/L(x)=\sum c_nx^n$. For $n\ge1$ put
$q_n=l_n/l_{n-1}$. Subtracting $q_n$ times the previous convolution
from the current one gives
\[
 l_0c_n=d_n-q_nd_{n-1}
 +\sum_{j=0}^{n-1}c_j(q_n l_{n-1-j}-l_{n-j}).
\]
Since $q_n$ is nondecreasing, induction from $c_0>0$ proves $c_n>0$
for every $n$. The logarithm defining $L$ has positive real part after adding the indicated constant, in the sense that $\Re L(z)=\log4-\tfrac12\log|1-z|>0$ for $|z|<1$.
Consequently the quotient series converges on the disk; termwise
differentiation proves strict positivity of every derivative on $(0,1)$.

\begin{thebibliography}{8}
\bibitem{r37showa:B} R. W. Barnard, K. C. Richards, and H. C. Tiedeman,
A survey of some bounds for Gauss' hypergeometric function and related bivariate means, J. Math. Inequal. 4 (2010), 45--52.
\bibitem{r37showa:D} Q. Deng, H. Peng, and J. Zou, Sharp power-mean comparisons for Gauss hypergeometric functions, \href{https://arxiv.org/abs/2609.28944}{arXiv:2609.28944} (2026).
\bibitem{r37showa:Z} T.-H. Zhao, M.-K. Wang, Y.-Q. Dai, and Y.-M. Chu,
On the generalized power-type Toader mean, J. Math. Inequal. 16 (2022), 247--264.
\bibitem{r37showa:T} M. D. Tobey, A two-parameter homogeneous mean value, Proc. Amer. Math. Soc. 18 (1967), 9--14.
\bibitem{r37showa:E} T. Zhao and Z.-H. Yang, Absolutely monotonic functions related to the asymptotic formula for the complete elliptic integral of the first kind, \href{https://arxiv.org/abs/2405.19651}{arXiv:2405.19651} (2024).
\bibitem{r37showa:W} M.-K. Wang, Z.-H. Yang, and T.-H. Zhao, An affirmative answer to an open problem on Ramanujan's asymptotic formula of zero-balanced hypergeometric function, \href{https://arxiv.org/abs/2504.04521}{arXiv:2504.04521} (2025).
\end{thebibliography}

\par\endgroup

% Source: /workspace/scratch/81914cec2c78/work/archive/Tournament_Round_37_Final/Tournament_Round_37/round37/heisei/paper.tex
% Edited source: /workspace/scratch/81914cec2c78/work/audit_34_41/cleaned/r37_heisei.tex
\begingroup
\newcommand{\R}{\mathbb R}
\newcommand{\vol}{\operatorname{vol}}
\newcommand{\pr}{\operatorname{pr}}
\Needspace{9\baselineskip}
\section{\texorpdfstring{No endpoint jumps in the intrinsic-volume growth of super-convolutive convex bodies}{No endpoint jumps in the intrinsic-volume growth of super-convolutive convex bodies}}\label{section:r37heisei}
\noindent{\small\sffamily\color{muted}Round 37 \enspace / \enspace Heisei}\par\medskip


\begin{sectionabstract}
For a sequence of convex bodies that contains every Cartesian product of its earlier members, we give finite-dimensional upper bounds for the two ends of its intrinsic-volume profile. At the left endpoint one uses many copies of a fixed body with nearly optimal normalized log-volume. At the right endpoint one uses a section of an inscribed product cube. These estimates settle both endpoint conditions in the jump problem posed by Jog. This is a research proof for review, not a claim of historical priority.
\end{sectionabstract}

\subsection{The question}
Let $K_n\subset\R^n$ be compact convex sets satisfying
\begin{equation}
 K_m\times K_n\subseteq K_{m+n}\qquad(m,n\ge1).
 \label{r37heisei:product}
\end{equation}
Write $u_{n,i}=V_i(K_n)$ for the intrinsic volumes with $V_0=1$ for nonempty convex sets and $V_n=\vol_n$ in dimension $n$. In Jog's notation the four assumptions are: (A) $v_n=n^{-1}\log u_{n,n}$ has a finite real limit $v$; (B) $n^{-1}\log u_{n,0}$ has a finite real limit; (C) $u_{n,0}$ and $u_{n,n}$ are positive at every $n$; and (D) $d_n=n^{-1}\log\sum_i u_{n,i}$ has a finite real limit. In particular, $K_1$ is a nondegenerate interval of length $\ell>0$, (B) has limit zero, and $v=\sup_s v_s$ by the volume consequence of \eqref{r37heisei:product}. The conjecture is about sets satisfying \eqref{r37heisei:product}, not just arrays of numbers satisfying a convolution inequality.

\begin{theorem}\label{r37heisei:main}
Under \eqref{r37heisei:product} and (A)--(D), the limits
\[
 G(\theta)=\lim_{n\to\infty}n^{-1}\log u_{n,\lfloor n\theta\rfloor},
 \qquad 0\le\theta\le1,
\]
form a continuous function, with $G(0)=0$ and $G(1)=v$.
\end{theorem}

This is Conjecture 1 in Section 5.1 of \cite{r37heisei:jog}. We recall the established interior argument to make clear exactly what the endpoint estimates must accomplish. Join the values $n^{-1}\log u_{n,i}$ linearly at the grid points $i/n$ to obtain $f_n$. Log-concavity of intrinsic volumes (Alexandrov--Fenchel) makes $f_n$ concave. By its endpoint chord, $f_n(x)\ge x v_n$; also $f_n\le d_n$. Thus on every closed interval strictly inside $(0,1)$ these functions have a common Lipschitz bound, by the secant inequalities for concave functions.

The super-convolution inequality required by Jog's Theorems 2.2.2 and 2.2.4 follows from monotonicity of intrinsic volumes and their convolution formula for Cartesian products. Those theorems give the closed-interval upper and open-interval lower large-deviation estimates for the grid measures with masses $u_{n,i}$ at $i/n$, with exponent $-\Lambda^*$, finite and concave in the interior. Explicitly $\Lambda(t)=\lim_n n^{-1}\log\sum_k u_{n,k}e^{kt}$, and $\Lambda^*(x)=\sup_t(xt-\Lambda(t))$. For an interior $x$, take any small interval about $x$. The logarithm of its mass divided by $n$ differs from its maximal grid value by at most $\log(n+1)/n$; the latter differs from $f_n(x)$ by at most the common local Lipschitz constant times the interval radius. Apply both deviation bounds and let the radius go to zero. Continuity of $-\Lambda^*$ inside $(0,1)$ gives $f_n(x)\to G(x)=-\Lambda^*(x)$, and the floor convention has the same limit. In particular $G(x)\ge xv$. What remains is to show
\[
 \limsup_{x\downarrow0}G(x)\le0,\qquad
 \limsup_{x\uparrow1}G(x)\le v.
\]

\subsection{A large projection and a thick fiber}
For a plane $E$ write $P_E$ for orthogonal projection. Vertical bars denote Lebesgue volume in the indicated space or affine flat. Write $\kappa_j$ for the volume of the Euclidean unit $j$-ball.

\begin{lemma}\label{r37heisei:geometry}
Let $K$ be a convex body in $\R^N$ and $1\le j<N$. There is a $j$-plane $E$ with $|P_EK|\ge V_j(K)/\binom Nj$. For any $j$-plane $E$, convex body $B\subseteq K$, $0<\lambda<1$, and $y\in E$,
\[
 |K|\ge \lambda^j(1-\lambda)^{N-j}|P_EK|\,
 |B\cap(y+E^\perp)|.
\]
\end{lemma}
\begin{proof}
Kubota's formula says $V_j(K)$ is the mean $j$-projection volume times $\binom Nj\kappa_N/(\kappa_j\kappa_{N-j})$. A unit $N$-ball lies in the Cartesian product of the two unit balls, so the prefactor is at most $\binom Nj$. For the second assertion, take any $x\in P_EK$ and choose $z\in K$ with $P_Ez=x$. Convexity puts
$\lambda z+(1-\lambda)(B\cap(y+E^\perp))$ in the fiber of $K$ over $\lambda x+(1-\lambda)y$. Integrate the resulting lower bound on fiber volumes over $\lambda P_EK+(1-\lambda)y$.
\end{proof}

Write $H(t)=-t\log t-(1-t)\log(1-t)$ for binary entropy. We use $\binom Nj\le\exp(NH(j/N))$.

\subsection{Ruling out the left jump}
Fix $s\ge1$. For $N=st$ the body $K_N$ contains the product $B=(K_s)^t$. Its log-volume is $Nv_s$. After a translation, $B$ lies in a ball of radius $R_s\sqrt N$, where $R_s=\operatorname{diam}(K_s)/\sqrt{s}$. Every $j$-projection of $B$ consequently satisfies
\[
 |P_EB|\le\kappa_j(R_s\sqrt N)^j\le (C_s\sqrt{N/j})^j,
 \qquad C_s=\sqrt{2\pi e}\,R_s .
\]
Indeed, integrating $e^{-j|x|^2/2}$ shows $\kappa_j\le(2\pi e/j)^{j/2}$. Fubini gives a supremum fiber volume of $B$ at least $e^{Nv_s}/|P_EB|$. Choose $E$ as in the first assertion of Lemma~\ref{r37heisei:geometry} for $K_N$, use its second assertion with $\lambda=j/N$, and take the supremum over fibers. With $z=j/N$ this yields
\begin{equation}
 N^{-1}\log V_j(K_N)
 \le v_N-v_s+2H(z)+z\log(C_s/\sqrt z).
 \label{r37heisei:left}
\end{equation}
One entropy term bounds the binomial coefficient; the other equals the cost of the convex mixing. Fix $x\in(0,1)$ and let $N$ tend to infinity through multiples of $s$ with $j=\lfloor Nx\rfloor$. By interior convergence the left side tends to $G(x)$. Next let $x\downarrow0$ with $s$ fixed: $\limsup_{x\downarrow0}G(x)\le v-v_s$. Only now send $s\to\infty$, using (A). This proves the first endpoint bound, with no uniform diameter bound on the bodies $K_s$.

\subsection{Ruling out the right jump}
By \eqref{r37heisei:product}, $K_N$ contains $K_1^N$, a translate of an $N$-cube of side $\ell$. Set $i=N-j$ with $1\le j<N$. For every $i$-plane $E$, Vaaler's cube section theorem \cite{r37heisei:vaaler} gives a fiber of this cube perpendicular to $E$ with $j$-volume at least $\ell^j$: take the fiber through the cube center. Choose $E$ with $|P_EK_N|\ge V_i(K_N)/\binom Ni$, and use Lemma~\ref{r37heisei:geometry} with projection dimension $i$ and $\lambda=i/N$. With $z=j/N$,
\begin{equation}
 N^{-1}\log V_{N-j}(K_N)\le v_N+2H(z)-z\log\ell.
 \label{r37heisei:right}
\end{equation}
For fixed $x\in(0,1)$ set $i=\lfloor Nx\rfloor$ and take the limit. Then $z\to1-x$; letting $x\uparrow1$ gives the second endpoint bound. The chord lower bound from the initial reduction completes Theorem~\ref{r37heisei:main}.

\paragraph{Scope.} For the geometric problem, no symmetry, common outradius, typical-set representation, or reverse product inclusion has been used. For abstract super-convolution arrays alone, both the projection and fiber steps are unavailable; Jog's Remark 2.2.5 gives a boundary discrepancy in that setting. Earlier typical-set continuity results of Jog and Anantharam \cite{r37heisei:ja} are not statements about all bodies obeying \eqref{r37heisei:product}.

\begin{thebibliography}{9}
\bibitem{r37heisei:jog} V. Jog, \emph{Convex geometric tools in information theory}, Ph.D. dissertation, University of California, Berkeley, EECS-2015-192 (2015). Equations (1.1), (1.7), Section 2.2, and Section 5.1, Conjecture 1.
\href{https://www2.eecs.berkeley.edu/Pubs/TechRpts/2015/Archive/EECS-2015-192.pdf}{Berkeley repository}.
\bibitem{r37heisei:ja} V. Jog and V. Anantharam, \emph{Intrinsic entropies of log-concave distributions}, \href{https://arxiv.org/abs/1702.01203}{arXiv:1702.01203}.
\bibitem{r37heisei:vaaler} J. D. Vaaler, A geometric inequality with applications to linear forms, \emph{Pacific J. Math.} \textbf{83} (1979), 543--553. \href{https://doi.org/10.2140/pjm.1979.83.543}{doi:10.2140/pjm.1979.83.543}.
\end{thebibliography}

\par\endgroup

% Source: /workspace/scratch/81914cec2c78/work/archive/Tournament_Round_37_Final/Tournament_Round_37/round37/reiwa/paper.tex
% Edited source: /workspace/scratch/81914cec2c78/work/audit_34_41/cleaned/r37_reiwa.tex
\begingroup
\newcommand{\dd}{\Delta}
\newcommand{\SSS}{S_j}
\Needspace{9\baselineskip}
\section{\texorpdfstring{Two partial-sum orders for triangle medians and altitudes}{Two partial-sum orders for triangle medians and altitudes}}\label{section:r37reiwa}
\noindent{\small\sffamily\color{muted}Round 37 \enspace / \enspace Reiwa}\par\medskip


\begin{sectionabstract}
We prove Jian Liu's 2012 Conjecture 3 for every triangle and every real exponent in both prescribed ranges. Positive powers follow from an additive comparison of two four-entry vectors; negative powers follow from a separate multiplicative comparison of their smallest entries. The two comparisons allow arbitrary increasing convex test functions on their respective scales.
\end{sectionabstract}
\begin{theorem}\label{r37reiwa:pow}
For a nondegenerate triangle with medians $m_a,m_b,m_c$, altitudes $h_a,h_b,h_c$, circumradius $R$, and inradius $r$, one has
\[
  \sum_{\rm cyc}m_a^k-\sum_{\rm cyc}h_a^k
      \le (2R)^k-(4r)^k\qquad (k>1\ \text{or}\ k<0).
\]
Equality holds if and only if the triangle is equilateral.
\end{theorem}
\noindent This is the complete exponent range of \cite[Conj.~3, Eq.~(4.3)]{r37reiwa:Liu}; Liu noted that $k=-1$ follows from an earlier bisector inequality. His Theorem~A, Eq.~(3.1), gives the result at $k=1$.

\paragraph{Geometry.}
Write $s$ for the semiperimeter and $\dd$ for the area. Number the sides in increasing order and denote their corresponding medians and altitudes by $m_1\ge m_2\ge m_3$ and $h_1\ge h_2\ge h_3$. The order follows from $4m_a^2=2(a^2+b^2+c^2)-3a^2$ and $h_a=2\dd/a$. Put $U=2R$ and $L=4r$. We use the following facts:
\begin{align}
&h_i\le m_i\le U,\quad h_2,h_3\le L\le U,\label{r37reiwa:e:basic}\\
&d=\sum_{i=1}^3(m_i-h_i)\le U-L,\label{r37reiwa:e:add}\\
&U\prod_{i=1}^3 h_i\le L\prod_{i=1}^3m_i.\label{r37reiwa:e:prod}
\end{align}
The first pointwise bound uses distance to the opposite line rather than distance to a particular point on the line. The second uses that a vertex and the opposite midpoint lie in the closed circumdisk. At most one side is less than $s/2$: two such sides would sum to less than $s$, so the third would exceed $s$. Hence the two smallest altitudes are at most $4\dd/s=4r$; Euler's inequality gives $4r\le 2R$. Inequality \eqref{r37reiwa:e:add} is precisely Liu's proved Theorem A; it is strict except in the equilateral case.

For \eqref{r37reiwa:e:prod}, Apollonius's identity
\[
 4m_a^2=4s(s-a)+(b-c)^2
\]
and Heron's formula imply $\prod m_i\ge s\dd$. Also $abc=4R\dd$ gives $\prod h_i=2\dd^2/R$, so $U\prod h_i=4\dd^2=L s\dd$. This proves \eqref{r37reiwa:e:prod}, strictly unless the triangle is equilateral.

\begin{theorem}\label{r37reiwa:vec}
Let $A=(U,h_1,h_2,h_3)$ and $B=(L,m_1,m_2,m_3)$. For each $j=1,2,3,4$, the sum of the largest $j$ entries of $A$ is at least the corresponding sum for $B$. The product of the smallest $j$ entries of $A$ is at most the corresponding product for $B$.
\end{theorem}
\begin{proof}
The entries of $A$ are displayed in decreasing order. If the largest $j$ entries of $B$ include $L$, their sum is $L+\sum_{i<j}m_i$. The difference from $U+\sum_{i<j}h_i$ is in the claimed direction by \eqref{r37reiwa:e:add}, since $m_i-h_i\ge0$. Otherwise $j\le3$. For $j=1$ use $m_1\le U$; for $j=2,3$ use
\[
 \sum_{i\le j}m_i\le U-L+\sum_{i\le j}h_i\le U+\sum_{i<j}h_i ,
\]
as $h_j\le L$.

For the product comparison at $j=1,2,3$, if the smallest $j$ entries of $B$ omit $L$, compare each median with the corresponding altitude. If they include $L$, at $j=1$ use $h_3\le L$; at $j=2$ use $h_2h_3\le Lm_3$. At $j=3$, \eqref{r37reiwa:e:prod} and $m_1\le U$ give $h_1h_2h_3\le Lm_2m_3$. The assertion at $j=4$ is precisely \eqref{r37reiwa:e:prod}.
\end{proof}

\paragraph{From ordered sums to powers.}
Recall the elementary convex-test principle: if $S_j(X)\ge S_j(Y)$ for all $j$, where $S_j$ sums the largest $j$ entries, then $\sum f(X_i)\ge\sum f(Y_i)$ for any increasing convex $f$ on an interval containing the entries. Indeed,
\[
 \sum_i(X_i-t)_+=\max_{0\le j\le n}\{S_j(X)-jt\},\qquad S_0=0.
\]
The increasing convex piecewise-linear interpolation of $f$ at the finite set of entries is a constant plus a nonnegative linear term plus a nonnegative combination of $(x-t)_+$. The principle follows term by term. For strictly increasing $f$ a strict total-sum comparison gives a strict test inequality, since the linear slope on the first interpolation interval is positive.

Apply the sum part of Theorem~\ref{r37reiwa:vec} with $f(x)=x^k$ on the positive reals for $k>1$. For $k<0$ apply the product part first: it says precisely that for every $j$ the sum of the largest $j$ entries of $(-\log A_i)_i$ dominates that of $(-\log B_i)_i$. The increasing convex function $f(t)=\exp((-k)t)$ gives $\sum A_i^k\ge\sum B_i^k$. Strictness follows respectively from \eqref{r37reiwa:e:add} and \eqref{r37reiwa:e:prod} at $j=4$ in a nonequilateral triangle. For an equilateral triangle the two vectors agree as multisets. This proves Theorem~\ref{r37reiwa:pow}.

\begin{thebibliography}{1}
\bibitem{r37reiwa:Liu} J.~Liu, On an inequality for the medians of a triangle, \emph{Journal of Science and Arts} 12 (2012), no.~2(19), 127--136. \href{https://josa.ro/docs/josa_2012_2/a_03_Liu_J.pdf}{Publisher PDF}.
\end{thebibliography}

\par\endgroup

% Source: /workspace/scratch/81914cec2c78/work/archive/Tournament_Round_37_Final/Tournament_Round_37/round37/keio/paper.tex
% Edited source: /workspace/scratch/81914cec2c78/work/audit_34_41/cleaned/r37_keio.tex
\begingroup

\Needspace{9\baselineskip}
\section{\texorpdfstring{Which powers preserve a triangle inequality chain?}{Which powers preserve a triangle inequality chain?}}\label{section:r37keio}
\noindent{\small\sffamily\color{muted}Round 37 \enspace / \enspace Keio}\par\medskip


\begin{sectionabstract}
We verify Conjectures 3.3 and 3.4 of Jian Liu by reducing all four families of cevians to a two-point weighted mean. The affirmative conclusions follow directly from the author's base inequality and Jensen's inequality.
\end{sectionabstract}

Let $ABC$ be a nondegenerate triangle with opposite sides $a,b,c$, exradii $r_a,r_b,r_c$, heights $h_a,h_b,h_c$ and medians $m_a,m_b,m_c$. As in [1], let $\ell_b,\ell_c$ be the corresponding lengths in any one of the four families of altitudes, internal angle bisectors, medians or symmedians. Define
\[
 J(p,q)=\frac{\ell_b^p/r_c^q+\ell_c^p/r_b^q}{\ell_b^p+\ell_c^p}.
\]
All bases are positive and all powers here are real.
\begin{theorem}[Liu, Conjecture 3.3]\label{r37keio:t:33}
If $p>0$ and $0<q\le1$, then $J(p,q)\le h_a^{-q}$. If $p>0>q$, or $p<0$ and $q\ge1$, then $J(p,q)\ge h_a^{-q}$.
\end{theorem}
For the first assertion the source uses two ranges: $0<p\le1$, $0<q\le1$, or $p\ge q>0$, $q\le1$. Their union is precisely the range displayed above.
\begin{theorem}[Liu, Conjecture 3.4]\label{r37keio:t:34}
If $q\ge1\ge p$, then $J(p,q)\ge m_a^{-q}$.
\end{theorem}
\paragraph{Reduction.}
Interchange $b,c$ if necessary so that $b\ge c$, and set
\[
 u=1/r_c,\qquad v=1/r_b,\qquad t=\ell_c/\ell_b .
\]
Writing $K$ for the area and $s$ for the semiperimeter, $1/r_i=(s-i)/K$ for $i=b,c$. Thus $u\ge v$ and $(u+v)/2=a/(2K)=1/h_a$. The first part of [1, Theorem 1.1, (1.7)] says that
\[
       \ell_b v+\ell_c u\ \ge\ (\ell_b+\ell_c)(u+v)/2 .
\]
Equivalently, $(\ell_c-\ell_b)(u-v)\ge0$. In particular $t\ge1$ for $b>c$. When $b=c$ symmetry gives $t=1$, and $m_a=h_a=r_b=r_c$, settling this case. Below we assume $b>c$.


\paragraph{Proofs.}
The exact representation
\[
 J(p,q)=\frac{u^q+t^p v^q}{1+t^p}
\]
makes $J(p,1)$ nonincreasing in $p$: increasing the weight on the smaller number $v$ cannot raise the mean. In particular $p>0$ implies $J(p,1)\le J(0,1)=1/h_a$, with the opposite relation when $p<0$.

If $p>0$ and $0<q\le1$, concavity and monotonic increase of $x^q$ give $J(p,q)\le J(p,1)^q\le h_a^{-q}$. For $p>0>q$, convexity and monotonic decrease give $J(p,q)\ge J(p,1)^q\ge h_a^{-q}$. Finally, for $p<0$ and $q\ge1$, convexity and monotonic increase give $J(p,q)\ge J(p,1)^q\ge h_a^{-q}$. These are all parts of Theorem~\ref{r37keio:t:33}.

The last part of [1, Theorem 1.1, (1.7)] states $J(1,1)\ge1/m_a$. For all $p\le1$ monotonicity gives $J(p,1)\ge J(1,1)$. Raising the mean to $q\ge1$ and applying Jensen once more yields
\[
 J(p,q)\ge J(p,1)^q\ge J(1,1)^q\ge m_a^{-q}.
\]
This proves Theorem~\ref{r37keio:t:34}.\hfill$\square$

\begin{thebibliography}{1}
\bibitem{r37keio:liu} J. Liu, A beautiful inequality chain for a triangle, \emph{Int. J. Open Probl. Comput. Sci. Math.} \textbf{15} (2022), 54--69. \url{https://www.ijopcm.org/Vol/2022/1.5.pdf}.
\end{thebibliography}

\par\endgroup

% Source: /workspace/scratch/81914cec2c78/work/archive/Tournament_Round_38_Final/Tournament_Round_38/round38/taisho/paper.tex
% Edited source: /workspace/scratch/81914cec2c78/work/audit_34_41/cleaned/r38_taisho.tex
\begingroup
\newcommand{\E}{\mathbb E}
\newcommand{\p}{\mathbb P}
\Needspace{9\baselineskip}
\section{\texorpdfstring{A signed Lambert series and the asymptotic formula for A330388}{A signed Lambert series and the asymptotic formula for A330388}}\label{section:r38taisho}
\noindent{\small\sffamily\color{muted}Round 38 \enspace / \enspace Taisho}\par\medskip

In 2019 V\'aclav Kote\v{s}ovec conjectured an asymptotic for the integers $a(n)$ defined in Ilya Gutkovskiy's OEIS entry A330388 [1] by
\[
 F(x)=\sum_{n\ge1}a(n)\frac{x^n}{n!}
 =\sum_{m\ge1}\frac{(-1)^{m+1}}{m}
       \frac{\log(1+x)^m}{1-\log(1+x)^m}.
\]
The outer Lambert series has arithmetic coefficients that oscillate in sign after reflection. We average them against a row of unsigned Stirling numbers, obtaining an error estimate for the individual coefficients.

\textbf{Theorem.} As $n\to\infty$,
\[
 a(n)=(-1)^{n+1}n!\,
 \frac{\pi^2 e^n}{24(e-1)^{n+1}}
 \left(1+O\!\left(\frac{\log n}{\sqrt n}\right)\right).
\]
In particular, the conjectured asymptotic equivalence holds.

\subsection*{The signed divisor weights}
Put $L(z)=-\log(1-z)$, $r=1-e^{-1}$, so $L(r)=1$. The coefficient of $w^k$ in $\sum_{m\ge1}(-1)^{m+1}w^m/[m(1-w^m)]$ is $\sum_{m\mid k}(-1)^{m+1}/m=\sigma_{\rm odd}(k)/k$, where $\sigma_{\rm odd}$ sums the odd divisors. Indeed the equivalent product identity is
\[
 \prod_{j\ge1}(1+w^j)=\prod_{d\ {\rm odd}\ge1}(1-w^d)^{-1};
\]
taking logarithms gives exactly those coefficient formulas. Thus
\[
 -F(-z)=\sum_{k\ge1} b_k L(z)^k,\qquad
 b_k=(-1)^{k+1}\frac{\sigma_{\rm odd}(k)}k .
\]
Writing $k=dm$ with $d$ odd, so that $k$ and $m$ have the same parity, gives
\[
 \sum_{k\le N} b_k
 =\sum_{m\le N}\frac{(-1)^{m+1}}m
          \#\{d\le N/m:\ d\text{ odd}\}
 =\frac{N}{2}\sum_{m\ge1}\frac{(-1)^{m+1}}{m^2}+O(\log(2N)).
\]
Here the count differs from $N/(2m)$ by at most one, and extending the square-summable series costs $O(1)$ after multiplication by $N$. Its alternating sum is $\pi^2/12$. Hence, with $B=\pi^2/24$, the partial sums of $b_k-B$ are $O(\log(2N))$.


\subsection*{A Stirling row smooths the fluctuations}
Set $W_{n,k}=r^n[z^n]L(z)^k$ for $1\le k\le n$, zero otherwise. For each $n$ this is a nonnegative unimodal row. Indeed, let $c(n,k)=[u^k]u(u+1)\cdots(u+n-1)$. Comparing powers of $u$ in $e^{uL(z)}=(1-z)^{-u}$ gives
\[
 W_{n,k}=\frac{r^n}{n!}\,k!\,c(n,k).
\]
By Newton's inequalities, $c(n,k)/\binom{n}{k}$ is log-concave in $k$, since the rising factorial has real nonpositive zeros. The second factor $k!\binom{n}{k}$ has decreasing consecutive ratios $n-k$. Their product is log-concave with contiguous support, so it is unimodal.

We next show $\max_k W_{n,k}=O(n^{-1/2})$. Consider an integer-valued random variable $X$ with $\p(X=j)=p_j=r^j/j$ for $j\ge1$ (these sum to $L(r)=1$). For independent copies $X_i$ one has $W_{n,k}=\p(\sum_{i=1}^k X_i=n)$. If $\phi(\theta)=\E e^{i\theta X}$, then, for $|\theta|\le\pi$,
\[
 1-|\phi(\theta)|^2
 =\sum_{i,j\ge1}p_i p_j[1-\cos((i-j)\theta)]
 \ge 2p_1p_2(1-\cos\theta)\ge 4p_1p_2\theta^2/\pi^2.
\]
Consequently Fourier inversion bounds $W_{n,k}$ by $(2\pi)^{-1}\int_{-\pi}^{\pi}|\phi(\theta)|^k\,d\theta=O(k^{-1/2})$. For $k\ge c n$ this suffices. To cover $k< c n$, fix $s\in(r,1)$ and use positivity:
\[
 W_{n,k}\le(r/s)^n L(s)^k\le[(r/s)L(s)^c]^n .
\]
Choose $c>0$ sufficiently small that the bracket is less than one. This proves the claimed uniform bound.

Abel summation, the divisor-weight discrepancy proved above, and unimodality now give
\[
 \left|\sum_{k=1}^n(b_k-B)W_{n,k}\right|
 \le O(\log(2n))\sum_{k=1}^n|W_{n,k}-W_{n,k+1}|
 =O(n^{-1/2}\log(2n)),
\]
as this variation is at most twice the largest entry. It remains to find the total mass:
\[
 \sum_{k=1}^n W_{n,k}
  =r^n[z^n]\frac{L(z)}{1-L(z)}
  =\frac{1}{rL'(r)}+O_s((r/s)^n)=\frac{1}{e-1}+O_s((r/s)^n).
\]
Indeed, $L(z)=1$ in $|z|<1$ forces $z=r$ by exponentiation. The pole is simple, and subtracting $(rL'(r))^{-1}/(1-z/r)$ leaves a holomorphic remainder. Cauchy's estimate on $|z|=s$ applies. Since $(-1)^{n+1}r^na(n)/n!=\sum b_kW_{n,k}$, the two estimates prove the theorem. $\square$

{\small\raggedright [1] I. Gutkovskiy, \href{https://oeis.org/A330388}{OEIS A330388} (December 12, 2019); asymptotic conjecture by V. Kote\v{s}ovec (December 16, 2019). A named public sequence-entry conjecture, not a journal problem; no exhaustive priority claim is made.\par}

\par\endgroup

% Source: /workspace/scratch/81914cec2c78/work/archive/Tournament_Round_38_Final/Tournament_Round_38/round38/heisei/paper.tex
% Edited source: /workspace/scratch/81914cec2c78/work/audit_34_41/cleaned/r38_heisei.tex
\begingroup

\Needspace{9\baselineskip}
\section{\texorpdfstring{Monotonicity of all positive zeros of Bessel derivatives}{Monotonicity of all positive zeros of Bessel derivatives}}\label{section:r38heisei}
\noindent{\small\sffamily\color{muted}Round 38 \enspace / \enspace Heisei}\par\medskip

\begin{sectionabstract}
The positive zeros of every derivative of the Bessel function increase strictly with the order on their full original parameter range. This is a complete known theorem of Dimitrov and Lun; the archived alternative proof is included without a novelty claim.
\end{sectionabstract}
\begin{theorem}
For every integer $n\ge0$ and every index $k\ge1$, the $k$th positive zero of $J_\nu^{(n)}$ is strictly increasing as a function of $\nu$ on $\nu>n-1$.
\end{theorem}
This is the complete positive-zero order-monotonicity assertion of Baricz, Kokologiannaki and Pog\'any \cite{r38heisei:bkp}, preprint Conjecture 3 and published Open Problem 6. Its prior resolution is due to Dimitrov and Lun \cite{r38heisei:dl}.

The normalized product for $J_\nu^{(m)}$, valid when $\nu>m-1$, has real simple nonzero zeros $\pm b_k$ and supplies the logarithmic-derivative equation for zeros of $J_\nu^{(m+1)}$ away from those zeros:
\begin{equation}\label{r38heisei:r38heisei:poles}
 \sum_{k\ge1}\frac{2w}{b_k^2-w}=\nu-m,\qquad w=z^2.
\end{equation}
This follows by differentiating the product
$Cz^{\nu-m}\prod_{k\ge1}(1-z^2/b_k^2)$, where $C\ne0$; the series is locally uniformly convergent away from its poles. The product and simplicity are the established results in \cite[Eq.~(2.1)]{r38heisei:bkp}.
\begin{proof}
Induct from the classical strict increase of the positive zeros of $J_\nu$ for $\nu>-1$. If $m\ge0$ and $\nu>m$, apply \eqref{r38heisei:r38heisei:poles}. Its right side is positive. For $w\ge0$, the left side is strictly increasing between poles, with one solution before the first pole and one in every later gap: it starts at zero at $w=0$ and goes from negative to positive infinity on each interior gap. As $\nu$ increases, the poles all move to the right by induction. In corresponding gaps, if the old root is at or left of the new left endpoint, it already lies below the new root. Otherwise it lies inside both gaps, and each term in \eqref{r38heisei:r38heisei:poles} at the old root strictly decreases, since no pole passes it and the two denominators have the same sign. Meanwhile the right side strictly increases. Hence every new root is strictly larger. This proves the rank-by-rank assertion in precisely $\nu>n-1$ for the $n$th derivative.

\end{proof}
\begin{thebibliography}{9}
\bibitem{r38heisei:bkp} \'A. Baricz, C. G. Kokologiannaki and T. K. Pog\'any, Zeros of Bessel function derivatives, \emph{Proc. Amer. Math. Soc.} \textbf{146} (2018), 209--222. \href{https://doi.org/10.1090/proc/13725}{doi:10.1090/proc/13725}; \href{https://arxiv.org/abs/1602.04295}{arXiv:1602.04295}.
\bibitem{r38heisei:kms} C. G. Kokologiannaki, M. E. Muldoon and P. D. Siafarikas, A unimodal property of purely imaginary zeros of Bessel and related functions, \emph{Canad. Math. Bull.} \textbf{37} (1994), 365--373. \href{https://doi.org/10.4153/CMB-1994-054-3}{doi:10.4153/CMB-1994-054-3}.
\bibitem{r38heisei:dl} D. K. Dimitrov and Y. C. Lun, Monotonicity of zeros of derivatives of Bessel functions, \emph{J. Approx. Theory} \textbf{305} (2025), 106102. \href{https://doi.org/10.1016/j.jat.2024.106102}{doi:10.1016/j.jat.2024.106102}.
\end{thebibliography}

\par\endgroup

% Source: /workspace/scratch/81914cec2c78/work/archive/Tournament_Round_38_Final/Tournament_Round_38/round38/reiwa/paper.tex
% Edited source: /workspace/scratch/81914cec2c78/work/audit_34_41/cleaned/r38_reiwa.tex
\begingroup

\Needspace{9\baselineskip}
\section{\texorpdfstring{A two-sided bound for the sum of triangle symmedians}{A two-sided bound for the sum of triangle symmedians}}\label{section:r38reiwa}
\noindent{\small\sffamily\color{muted}Round 38 \enspace / \enspace Reiwa}\par\medskip


\begin{sectionabstract}
We prove both inequalities in Jian Liu's Conjecture 5 (2012). The lower bound uses a rational approximation only after subtracting a common leading term; the upper bound uses a weighted Cauchy inequality. Two integer-polynomial positivity certificates cover the entire triangle domain.
\end{sectionabstract}
\begin{theorem}
Let $k_a,k_b,k_c$ be the internal symmedian lengths of a nondegenerate triangle, $m_a,m_b,m_c$ its medians, $r_a,r_b,r_c$ its exradii, $R$ its circumradius and $r$ its inradius. Then
\[
 \frac{\sum k_a}{\sum m_a}\ \ge\ \frac12+\frac rR
 \ \ge\ \frac{\sum k_a}{\sum r_a}.
\]
Either equality occurs exactly for equilateral triangles.
\end{theorem}
This is the double inequality in \cite[Conj.~5, Eq.~(4.8)]{r38reiwa:Liu}; all sums run over the three corners.

\begin{proof}
Write $x=s-a$, $y=s-b$, $z=s-c$ for the positive tangent lengths and $s=x+y+z$ for the semiperimeter; thus $a=y+z$, $b=z+x$, $c=x+y$. Put
\[
 P=xyz,\qquad q=xy+yz+zx,\qquad J=abc,\qquad C_0=\frac{J+8P}{2J}.
\]
Heron's formula gives area $\Delta=\sqrt{sP}$. From $r=\Delta/s$ and $R=J/(4\Delta)$, it follows that $C_0=1/2+r/R\le1$. Further,
$\sum r_a=\Delta(1/x+1/y+1/z)=q\sqrt{s/P}$.

For each corner form the cyclic quantities
\[
 \begin{aligned}
 A&=4sx+(y-z)^2=4m_a^2,&\quad H&=b+c,&\quad N&=b^2+c^2,\\
 Q&=A/H^2,& d&=(b-c)^2/N,& T&=2bcH/N.
\end{aligned}
\]
Then $2m_a=H\sqrt Q$ and $2k_a=T\sqrt Q=(1-d)H\sqrt Q$. Indeed the internal symmedian divides the opposite side in the ratio $c^2:b^2$. Writing $\mathbf u,\mathbf v$ for unit vectors from its vertex toward the two other vertices, its length is $bc\lVert b\mathbf u+c\mathbf v\rVert/(b^2+c^2)$. The norm is unchanged on swapping $b,c$, giving $2bc\,m_a/(b^2+c^2)$.

For the lower bound, subtract the common leading term first. It suffices to show
\begin{equation}\label{r38reiwa:l}
 (1-C_0)\sum H\ell(Q)\ \ge\ \sum dH g(Q),
 \qquad
 \ell(t)=\frac{4t(1+t)}{1+6t+t^2},\quad
 g(t)=\frac{1+6t+t^2}{4(1+t)}.
\end{equation}
Indeed $\ell(t)\le\sqrt t\le g(t)$ for $t>0$: writing $h=\sqrt t$, the numerator of $g(t)-h$ is $(h-1)^4$, and $\ell(t)g(t)=t$. Both coefficients $1-C_0,d$ are nonnegative. Thus \eqref{r38reiwa:l} implies $(1-C_0)\sum H\sqrt Q\ge\sum dH\sqrt Q$, precisely the left inequality.

For the upper bound, $4(\sum k_a)^2=(\sum T\sqrt Q)^2\le(\sum T)(\sum TQ)$ by Cauchy--Schwarz. It suffices to show
\begin{equation}\label{r38reiwa:r}
 (J+8P)^2 s q^2\ \ge\ PJ^2(\sum T)(\sum TQ).
\end{equation}
Both \eqref{r38reiwa:l} and \eqref{r38reiwa:r} are rational inequalities with positive denominators. By symmetry and homogeneity it suffices to fix $x=1$, $y=1+u$, $z=1+u+v$ where $u,v\ge0$. The following complete calculation in $\mathbb Z[u,v]$ clears the denominators and constructs the two residual polynomials $F,G$, respectively.
\begin{small}
\begin{Verbatim}
from sympy import ZZ
from sympy.polys.rings import ring
B,u,v=ring('u,v',ZZ)
x,y,z=B.one,1+u,1+u+v
s=x+y+z; P=x*y*z; q=x*y+y*z+z*x; a,b,c=y+z,z+x,x+y
J=a*b*c; ll=[]; rr=[]; t=[]; tq=[]
for X,Y,Z,B,C in [(x,y,z,b,c),(y,z,x,c,a),(z,x,y,a,b)]:
    A=4*s*X+(Y-Z)**2; H=B+C; N=B*B+C*C
    K=A*A+6*A*H**2+H**4; W=A+H*H
    ll.append((4*H*A*W,K))
    rr.append(((B-C)**2*K,4*N*H*W))
    t.append((2*B*C*H,N))
    tq.append((2*B*C*A,N*H))
def add(rows):
    (i,a),(j,b),(k,c)=rows
    return i*b*c+a*j*c+a*b*k,a*b*c
n,d=add(ll); m,e=add(rr); F=(J-8*P)*n*e-2*J*m*d
n,d=add(t); m,e=add(tq); G=(J+8*P)**2*s*q*q*d*e-P*J**2*n*m
for p,nn,degree,minimum in [(F,520,31,536870912),(G,363,26,8)]:
    assert len(p)==nn and max(i+j for i,j in p)==degree
    assert min(p.values())==minimum
    assert p.get((0,0),0)==0 and p[2,0]>0 and p[0,2]>0
\end{Verbatim}

\end{small}
Each residual has exclusively positive nonzero coefficients, including those of $u^2$ and $v^2$. Both inequalities are therefore strict unless $u=v=0$. In that case the triangle is equilateral, and the medians, exradii and symmedians all have the same length, with $C_0=1$.
\end{proof}
\begin{thebibliography}{1}
\bibitem{r38reiwa:Liu} J.~Liu, On an inequality for the medians of a triangle, \emph{Journal of Science and Arts} 12 (2012), no.~2(19), 127--136. \href{https://www.josa.ro/docs/josa_2012_2/a_03_Liu_J.pdf}{Publisher PDF}, p.~135.
\end{thebibliography}

\par\endgroup

% Source: /workspace/scratch/81914cec2c78/work/archive/Tournament_Round_38_Final/Tournament_Round_38/round38/keio/paper.tex
% Edited source: /workspace/scratch/81914cec2c78/work/audit_34_41/cleaned/r38_keio.tex
\begingroup
\newcommand{\cyc}{\sum_{\rm cyc}}
\Needspace{9\baselineskip}
\section{\texorpdfstring{Convex interpolation in a median-ratio inequality}{Convex interpolation in a median-ratio inequality}}\label{section:r38keio}
\noindent{\small\sffamily\color{muted}Round 38 \enspace / \enspace Keio}\par\medskip


\begin{sectionabstract}
A change of parameter reduces a continuum of inequalities for the medians of acute triangles to two endpoint comparisons. We prove each endpoint by a finite cover using exact integer intervals, with an analytic treatment of the equilateral equality case. This establishes Jian Liu's Conjecture~4.2 over its full stated interval; no optimality of the endpoints is claimed.
\end{sectionabstract}
\begin{theorem}
Let $a,b,c$ be the sides of an acute triangle, and $m_a,m_b,m_c$ their corresponding medians. For every real $-21/100\le k\le 55/100$,
\[
        \cyc\frac{a}{m_b+m_c+km_a}\ \le\ \frac{2\sqrt3}{k+2}.
\]
Equality holds precisely for equilateral triangles.
\end{theorem}
This is Conjecture~4.2, equation~(65) of [1]. Theorem~1.1 of [1] proves the instance $k=0$. Neither of the neighboring conjectures is needed.

\paragraph{Why two parameter values suffice.}
The median lengths satisfy the triangle inequalities: the three vectors from the centroid to the vertices sum to zero and have lengths $2m_i/3$. In particular $m_b+m_c+km_a>0$ on the given interval. Set $\lambda=k/(k+2)$, an increasing function of $k$. Then
\[
 \frac{m_b+m_c+km_a}{k+2}
  =(1-\lambda)\frac{m_b+m_c}{2}+\lambda m_a.
\]
Each such positive affine function of $\lambda$ has convex reciprocal, even though the left endpoint of the $\lambda$-interval is negative. Thus the scaled cyclic sum is convex, and its maximum on the interval is at an endpoint. Strictness at both endpoints for a fixed nonequilateral triangle implies strictness throughout.

\paragraph{A square chart.}
Order and rescale the sides so that $a\le b=1\le c$. Acuteness gives $c^2\le1+a^2$. Write $a=1-x$, $c=\sqrt{1+a^2 y}$, so that $(x,y)\in[0,1]^2$. The doubled medians $u,v,w$ opposite $a,1,c$ satisfy, by Apollonius' formula,
\[
 u^2=4+a^2(2y-1),\quad v^2=1+a^2(2+2y),\quad w^2=1+a^2(2-y).
\]
It is harmless to certify the closed square, which includes limiting triangles.


\paragraph{Endpoint certificate.}
For $K=-21,55$ put $U=100(v+w)+Ku$, $V=100(u+w)+Kv$, $W=100(u+v)+Kw$. At $k=K/100$ the assertion is exactly
\[
 F_K(x,y)=\sqrt3-(200+K)\left(\frac aU+\frac1V+\frac cW\right)\ge0.
\]
The origin is equilateral, giving $F_K(0,0)=0$. Its gradient there vanishes: the residual as a function of three independent side lengths is smooth, symmetric and homogeneous of degree zero near $(1,1,1)$. Symmetry equates its three first derivatives there, and differentiating along common scaling shows that their sum is zero. Composition with the chart preserves the zero gradient.

On $[0,1/256]^2$, the following are rigorous integer lower bounds for $10^6$ times the indicated second partial derivatives:
\[
 \begin{array}{r|rrr}
 K&F_{xx}&F_{xy}&F_{yy}\\\hline
 -21&131246&19222&29359\\
 55&85609&8082&18666
 \end{array}
\]
For any nonzero $(x,y)$ in that square, $x,y\ge0$, and the second derivative along $t\mapsto(tx,ty)$ is strictly positive. Together with the zero value and gradient at the origin this proves the desired strict inequality locally. Entrywise positivity is used only on nonnegative directions.

The remaining boxes are checked by one of two enclosures: direct evaluation of $F_K$, or the centered enclosure $F_K(z)+\nabla F_K(B)\cdot(B-z)$ for a box $B$ with midpoint $z$. A box is accepted only when its lower bound is strictly positive or it belongs to the verified local square. Otherwise both closed halves are processed. The complete cover has the following statistics.
\[
\begin{array}{r|rrr|r}
 K& \text{direct}&\text{centered}&\text{local}&\text{maximum depth}\\\hline
 -21&6&643&1&23\\
 55&5&859&1&23
\end{array}
\]
All interval endpoints are integers divided by $2^{64}$. Multiplication and inversion round outwards, and square roots are bounded using integer square roots. The derivative rules are the ordinary product, quotient and chain rules applied to intervals. Failed domain checks trigger subdivision, never acceptance. The following complete Python~3 verifier derives the tables. The normalized chart and generic interval differentiation routines are adapted from Meiji's Round~37 work on a different inequality; the convex reduction and the two residual certificates here address~(65). This completes the proof.

\begin{thebibliography}{1}
\bibitem{r38keio:liu} J. Liu, An inequality involving medians and sides of an acute triangle, \emph{Int. J. Open Probl. Comput. Math.} \textbf{14}(3) (2021), 49--63.
\url{https://www.ijopcm.org/Vol/2021/3.5.pdf}.
\end{thebibliography}
\subsection*{Verifier: Intervals and first derivatives}
\scriptsize\begin{Verbatim}
from math import isqrt
from collections import Counter
S=1<<64
class I:
 __slots__='l','h'
 def __init__(self,l,h=None):
  if h is None:self.l=self.h=l*S
  else:self.l=l;self.h=h
 def __add__(x,y):
  if not isinstance(y,I):y=I(y)
  return I(x.l+y.l,x.h+y.h)
 __radd__=__add__
 def __neg__(x):return I(-x.h,-x.l)
 def __sub__(x,y):return x+-y
 def __rsub__(x,y):return -x+y
 def __mul__(x,y):
  if not isinstance(y,I):return I(min(x.l*y,x.h*y),max(x.l*y,x.h*y))
  if x.l>=0 and y.l>=0:return I(x.l*y.l//S,-(-x.h*y.h//S))
  p=(x.l*y.l,x.l*y.h,x.h*y.l,x.h*y.h);return I(min(p)//S,-(-max(p)//S))
 __rmul__=__mul__
 def __truediv__(x,y):
  if not isinstance(y,I):
   assert y>0;return I(x.l//y,-(-x.h//y))
  if y.l<=0:raise ValueError
  return x*I(S*S//y.h,-(-S*S//y.l))
 def rt(x):
  if x.l<=0:raise ValueError
  return I(isqrt(x.l*S),isqrt(x.h*S)+1)
class J:
 __slots__='v','g'
 def __init__(self,v,g):
  self.v=v if isinstance(v,I) else I(v);self.g=[d if isinstance(d,I) else I(d) for d in g]
 def __add__(x,y):
  if not isinstance(y,J):return J(x.v+y,x.g)
  return J(x.v+y.v,[p+q for p,q in zip(x.g,y.g)])
 __radd__=__add__
 def __neg__(x):return J(-x.v,[-p for p in x.g])
 def __sub__(x,y):return x+-y
 def __rsub__(x,y):return -x+y
 def __mul__(x,y):
  if not isinstance(y,J):return J(x.v*y,[p*y for p in x.g])
  return J(x.v*y.v,[p*y.v+q*x.v for p,q in zip(x.g,y.g)])
 __rmul__=__mul__
 def __truediv__(x,y):
  if not isinstance(y,J):return J(x.v/y,[p/y for p in x.g])
  v=x.v/y.v;return J(v,[(p-v*q)/y.v for p,q in zip(x.g,y.g)])
 def rt(x):
  y=x.v.rt();return J(y,[p/(2*y) for p in x.g])
\end{Verbatim}

\subsection*{Verifier: Mixed second derivatives}
\scriptsize\begin{Verbatim}
class B:
 def __init__(self,v,g=0,h=0,j=0):
  self.v,self.g,self.h,self.j=[i if isinstance(i,I) else I(i) for i in (v,g,h,j)]
 def __add__(x,y):
  if not isinstance(y,B):y=B(y)
  return B(x.v+y.v,x.g+y.g,x.h+y.h,x.j+y.j)
 __radd__=__add__
 def __neg__(x):return B(-x.v,-x.g,-x.h,-x.j)
 def __sub__(x,y):return x+-y
 def __rsub__(x,y):return -x+y
 def __mul__(x,y):
  if not isinstance(y,B):y=B(y)
  return B(x.v*y.v,x.g*y.v+x.v*y.g,x.h*y.v+x.v*y.h,x.j*y.v+x.g*y.h+x.h*y.g+x.v*y.j)
 __rmul__=__mul__
 def chain(x,y,a,b):return B(y,a*x.g,a*x.h,a*x.j+b*x.g*x.h)
 def __truediv__(x,y):
  if not isinstance(y,B):return B(x.v/y,x.g/y,x.h/y,x.j/y)
  r=I(1)/y.v;return x*y.chain(r,-r*r,2*r*r*r)
 def rt(x):
  r=x.v.rt();d=I(1)/(2*r);return x.chain(r,d,-d/(2*x.v))


\end{Verbatim}

\subsection*{Verifier: Residual and complete cover}
\scriptsize\begin{Verbatim}
def f(x,y,k):
 a=1-x; aa=a*a; c=(1+aa*y).rt()
 u=(4+aa*(2*y-1)).rt()
 v=(1+aa*(2+2*y)).rt()
 w=(1+aa*(2-y)).rt()
 # u,v,w are twice the medians opposite a,1,c.
 U=100*(v+w)+k*u;V=100*(u+w)+k*v;W=100*(u+v)+k*w
 return -(a/U+(v*0+1)/V+c/W)*(200+k)+I(3).rt()
def local(k,eps):
 z=I(0,S//eps)
 bounds=[]
 for p,q in [(0,0),(0,1),(1,1)]:
  H=f(B(z,p==0,q==0),B(z,p==1,q==1),k).j
  bounds.append(H.l*10**6//S)
 if min(bounds)>0:
  print('positive Hessian bounds /1e6:',k,eps,bounds,flush=True)
  return True
 return False
def prove(k):
 eps=8
 while True:
  try:
   if local(k,eps):break
  except ValueError:pass
  eps*=2;assert eps<2**20
 ct=Counter();stack=[(0,S,0,S,0)]
 while stack:
  l,h,u,v,d=stack.pop();ct['depth']=max(ct['depth'],d)
  if h<=S//eps and v<=S//eps:ct['Taylor']+=1;continue
  xx=I(l,h);yy=I(u,v)
  try:
   if f(xx,yy,k).l>0:ct['direct']+=1;continue
   m=(l+h)//2;n=(u+v)//2
   val=f(I(m,m),I(n,n),k)
   jet=f(J(xx,[1,0]),J(yy,[0,1]),k)
   if (val+jet.g[0]*I(l-m,h-m)+jet.g[1]*I(u-n,v-n)).l>0:ct['mean']+=1;continue
  except ValueError:pass
  assert d<50,(k,d,l/S,h/S,u/S,v/S,ct)
  if h-l>=v-u:
   m=(l+h)//2;stack.extend([(l,m,u,v,d+1),(m,h,u,v,d+1)])
  else:
   n=(u+v)//2;stack.extend([(l,h,u,n,d+1),(l,h,n,v,d+1)])
 print('endpoint',k,'/100',dict(ct),flush=True)
for k in (-21,55):prove(k)
\end{Verbatim}


\par\endgroup

% Source: /workspace/scratch/81914cec2c78/work/archive/Tournament_Round_39_Final/Tournament_Round_39/round39/meiji/paper.tex
% Edited source: /workspace/scratch/81914cec2c78/work/audit_34_41/cleaned/r39_meiji.tex
\begingroup

\Needspace{9\baselineskip}
\section{\texorpdfstring{A determinant inequality for all powers of three exradius--median ratios above $-12$}{A determinant inequality for all powers of three exradius--median ratios above -12}}\label{section:r39meiji}
\noindent{\small\sffamily\color{muted}Round 39 \enspace / \enspace Meiji}\par\medskip


\begin{sectionabstract}
Liu asked whether a cubic determinant formed from three corresponding exradius--median ratios is nonpositive for every real exponent $k\ge-12$. We prove that it is, with equality for nonzero $k$ exactly at isosceles triangles. The only power at which a geometric certificate is needed is $k=-12$: Descartes' sign rule and a forced double zero at the origin propagate the endpoint sign to the whole half-line. The endpoint certificate is an exact positive-coefficient expansion in two ordered tangency-length differences.
\end{sectionabstract}

\begin{theorem}
For a nondegenerate Euclidean triangle let $m_i$ be its medians and $r_i$ its exradii. Put $U_i=(r_j+r_\ell)/(2m_i)$, with $\{i,j,\ell\}=\{1,2,3\}$. For every real $k\ge-12$,
\[
 D(k):=1+2(U_1U_2U_3)^k-\sum_{i=1}^3U_i^{2k}\le0 .
 \tag{1}
\]
If $k\ne0$, equality holds exactly when the triangle is isosceles (including equilateral). At $k=0$ equality is automatic.
\end{theorem}
This is Conjecture 2, Eq.~(4.4), of Liu \cite[p.~73]{r39meiji:liu}. The distinct Conjecture 1 uses altitude ratios and is not used. No optimality claim about the cutoff $-12$ is made here.

\paragraph{Geometry at a single exponent.}
Let $x_i=s-a_i$, where $a_i$ are the sides and $s$ is the semiperimeter, so that $s=x_1+x_2+x_3$ and $a_i=x_j+x_\ell$. The area $\Delta$ satisfies $\Delta^2=sx_1x_2x_3$ and $r_j=\Delta/x_j$. Apollonius's identity gives $4m_i^2=4sx_i+(x_j-x_\ell)^2$. Consequently
\[
 U_i^2=\frac{N_i}{E_i},\qquad
 N_i=sx_i(x_j+x_\ell)^2,\qquad
 E_i=x_jx_\ell\{4sx_i+(x_j-x_\ell)^2\}. \tag{2}
\]
Both are strictly positive. Moreover
\[
 N_i-E_i=(x_j-x_\ell)^2(sx_i-x_jx_\ell). \tag{3}
\]
All quantities in (1) are scale invariant and symmetric under relabelling; hence we can use
\[
 (x_1,x_2,x_3)=(1,1+u,1+u+v),\qquad u,v\ge0.
 \tag{4}
\]
The scalene locus is precisely $u,v>0$. Put $n=N_1N_2N_3$ and $e=E_1E_2E_3$. At $k=-12$ the polynomial
\[
 P=-n^{12}D(-12)
   =\sum_i(E_i N_jN_\ell)^{12}-n^{12}-2(ne)^6
 \tag{5}
\]
detects the sign of the determinant without radicals or division.

\begin{lemma}
With the substitution (4), $P(u,v)$ has exactly $8{,}788$ nonzero coefficients. Each is a positive integer, and each corresponding monomial is divisible by $u^2v^2$. In particular $D(-12)<0$ for every scalene triangle.
\end{lemma}
\begin{proof}
Here and at the end of the article we give the complete finite expansion certificate. Let $P_+$ be the sum in (5) and $P_-=n^{12}+2(ne)^6$. Each of these has nonnegative coefficients: in (2) every squared difference in (4) is $u^2,v^2$ or $(u+v)^2$. The total degrees are at most $144$. Choose $B=2^{512}$. Direct integer evaluation of $P_\pm(1,1)$ gives two integers of bit length $267$. Therefore every coefficient of either polynomial is less than $B$. For each polynomial evaluated at $(B,B^{145})$, the base-$B$ digit numbered $i+145j$ is exactly the coefficient of $u^iv^j$. No carries occur and different exponent pairs cannot share a digit. Subtracting corresponding digits of \emph{separately} expanded $P_+$ and $P_-$ gives the asserted count and signs, with minimum nonzero coefficient $1$. The code below performs every operation including both degree and nonnegativity checks. This proves an identity of polynomials and its consequences on the whole quadrant; it is not a sign test at finitely many shapes.
\end{proof}

\paragraph{Propagation to arbitrary powers.}
We recall the sign rule for finite exponential sums: if $\theta_1<\cdots<\theta_h$ and $c_j$ are nonzero, then $\sum_j c_j e^{\theta_j k}$ has no more real zeros, counting multiplicity, than sign changes in $(c_1,\ldots,c_h)$; see \cite{r39meiji:jameson}. Here is a direct argument. Substitute $t=e^k$ and divide by $t^{\theta_1}$. The expression now has constant term $c_1$. Its derivative has successive coefficient signs $\operatorname{sgn}c_2,\dots,\operatorname{sgn}c_h$ and strictly increasing real exponents, so induction applies. If there are $q$ positive roots counted with multiplicity, Rolle's theorem gives at least $q-1$ derivative roots. If $c_1,c_2$ have the same sign and $q>0$, there is an extra derivative root before the first positive root: near $0$ the function starts with sign $c_1$ and initially moves further in that direction. Thus the number of roots can increase from the derivative's bound only when the extra first coefficient adds a sign change. The single-term base case is immediate.

Now assume the triangle scalene. For the two greatest $x_i$, (3) implies $U_i>1$: if $x_i$ is not the smallest, $x_i\max(x_j,x_\ell)\ge x_jx_\ell$ and $s>\max(x_j,x_\ell)$. If the remaining $U_i=1$, (1) becomes $-(U_j^k-U_\ell^k)^2$. The two bases must differ by the strict endpoint sign, so equality occurs only at zero. If the remaining $U_i<1$, choose an index $j$ with $U_j>1$; writing $z_i=U_i^k$ gives
\[
 D(k)=(1-z_i^2)(1-z_j^2)-(z_\ell-z_iz_j)^2<0\qquad(k\ne0).
\]
This leaves only the case where every $U_i>1$. Set $L_i=\log U_i>0$. Then
\[
 D(k)=1-\sum_i e^{2L_i k}+2e^{(L_1+L_2+L_3)k}.
\]
The term at exponent $0$ is positive. The remaining terms all have negative coefficients except for one positive coefficient, so there are at most three sign changes in exponent order, also after combining ties and removing zeros. By the sign rule, $D$ has at most three real zeros with multiplicity. As $D(-\infty)=1$ in the limiting sense and $D(-12)<0$, one zero occurs strictly before $-12$. Also $D(0)=D'(0)=0$. Thus zero itself has exactly multiplicity two and there are no other zeros. Its sign does not change at zero; hence $D(k)<0$ for every $k\ge-12$ other than zero.

For an isosceles triangle two of the $x_i$ are equal. Formula (3) makes the remaining $U_i$ equal to one, while the two other $U$'s are equal by symmetry. Thus $D(k)=0$ identically. This concludes the proof.

\begin{thebibliography}{9}
\bibitem{r39meiji:liu} J. Liu, Inequalities for the product of medians in a triangle, \emph{International Journal of Geometry} \textbf{14} (2025), no.~3, 57--74. \href{https://ijgeometry.com/wp-content/uploads/2025/06/5.-57-74.pdf}{Publisher PDF}.
\bibitem{r39meiji:jameson} G.~J.~O. Jameson, Counting zeros of generalised polynomials: Descartes' rule of signs and Laguerre's extensions, \emph{Mathematical Gazette} \textbf{90} (2006). \href{https://doi.org/10.1017/S0025557200179628}{doi:10.1017/S0025557200179628}.
\end{thebibliography}

\paragraph{Exact coefficient verifier (Python 3).} Only the standard interpreter is needed; keep assertions enabled.
\begingroup\fontsize{8}{9.7}\selectfont
\begin{Verbatim}
"""Pure Python 3, exact carry-free coefficient verification."""
def eval_parts(u,v):
    x=[1,1+u,1+u+v]
    s=sum(x); N=[]; D=[]
    for i,a in enumerate(x):
        b,c=x[(i+1)%3],x[(i+2)%3]
        N.append(s*a*(b+c)**2)
        D.append(b*c*(4*s*a+(b-c)**2))
    n=N[0]*N[1]*N[2]
    d=D[0]*D[1]*D[2]
    plus=sum((D[i]*N[(i+1)%3]*N[(i+2)%3])**12 for i in range(3))
    return plus,n**12+2*(n*d)**6

B=1<<512
stride=145
assert all(0<m<B for m in eval_parts(1,1))
p,q=eval_parts(B,B**stride)
length=64*stride**2
# Every coefficient is at most its whole polynomial's value at (1,1).
# The exponent address is i+145*j for u**i*v**j, with i,j<=144.
p=p.to_bytes(length,'little'); q=q.to_bytes(length,'little')
def digit(a,h):
    return int.from_bytes(a[64*h:64*(h+1)],'little')
cs=[]
for j in range(stride):
    for i in range(stride):
        h=i+stride*j
        c=digit(p,h)-digit(q,h)
        if c: cs.append((i,j,c))
assert len(cs)==8788 and min(c for _,_,c in cs)==1
assert all(i+j<=144 and i>=2 and j>=2 for i,j,c in cs)
print('verified 8788 positive coefficients, minimum 1; all other coefficients 0')
\end{Verbatim}
\endgroup

\par\endgroup

% Source: /workspace/scratch/81914cec2c78/work/archive/Tournament_Round_39_Final/Tournament_Round_39/round39/taisho/paper.tex
% Edited source: /workspace/scratch/81914cec2c78/work/audit_34_41/cleaned/r39_taisho.tex
\begingroup
\newcommand{\ee}{\mathrm e}
\Needspace{9\baselineskip}
\section{\texorpdfstring{A growing-dimensional first-zero path count}{A growing-dimensional first-zero path count}}\label{section:r39taisho}
\noindent{\small\sffamily\color{muted}Round 39 \enspace / \enspace Taisho}\par\medskip

Start from $(n,\ldots,n)$ in $n$ labelled coordinates. A step decrements any nonempty subset of coordinates by one. Stop immediately when a coordinate first reaches zero; simultaneous first zeros terminate a single path. Write $a(n)$ for the number of terminal paths. Kote\v{s}ovec's November 13, 2026 formula for OEIS A263160 specifies an asymptotic equivalent with a numerical multiplier [1,2]. We give its closed form.

\textbf{Theorem.} Put $L=\log 2$. As $n\to\infty$,
\[
 a(n)\sim
 \frac{2\sqrt{2\pi}}{L}\exp\!\left(\frac5{12}-\frac{L^2}{24}\right)
 \frac{n^{\,n(n-3/2)+1}}{\pi^{n/2}L^{\,n(n-1)}} .
\]
The prefactor is $10.753662835613\ldots$. The separate, weaker statement $a(n)^{1/n^2}/n\to1/L$ already has a public proof [2].

\subsection*{Counting before the first zero}
Let $m=n-1$, $M=nm$, and let $I$ be the number of internal paths, counting the empty prefix. Each internal node of the path tree has $2^n-1$ children, and internal paths have at most $M$ steps. The edge count of a finite tree gives
\[
 a(n)=1+(2^n-2)I.
\]
An internal path uses each coordinate at most $m$ times. Marking coordinate uses by $x_i$, the generating function for all words of nonempty subsets is
\[
 \frac{1}{2-\prod_{i=1}^n(1+x_i)}
 =\frac12\sum_{k\ge0}2^{-k}\prod_{i=1}^n(1+x_i)^k .
\]
This expansion converges in a neighborhood of the origin; the identity also follows by summing nonnegative coefficients. Summing coefficients over the box of all exponents at most $m$ yields the known binomial-moment representation
\[
 I=\frac12\sum_{k\ge0}e^{-Lk}B(k,m)^n,\qquad
 B(k,m)=\sum_{j=0}^m\binom{k}{j},
\]
where $\binom{k}{j}=0$ for $j>k$. It remains to estimate the sum when $m=n-1$, not in fixed dimension.


\subsection*{The displaced saddle}
Write $k_0=M/L$, $\sigma=\sqrt M/L$ and $y=(k-k_0)/\sigma$. The base weight $e^{-Lk}k^M$ is maximized at $k_0$. We first restrict to the middle interval $k_0/2\le k\le2k_0$. There $k\asymp n^2$. Relative to the last binomial $\binom{k}{m}$, the penultimate one has size $m/(k-m+1)=m/k+O(n^{-2})$. Each consecutive backward ratio is no greater, so the sum of all earlier terms relative to the last is $O(n^{-2})$. Thus, with $D(k)=n\log(B(k,m)m!/k^m)$,
\[
 D(k)= -\frac{nm(m-1)}{2k}
 -\frac{nm(m-1)(2m-1)}{12k^2}+\frac{nm}{k}+O(n^{-1}). \tag{1}
\]
Indeed, apply $\log(1-u)=-u-u^2/2+O(u^3)$ to $u=j/k$, $0\le j<m$, in the falling factorial. The cubic remainders total $O(nm^4/k^3)=O(n^{-1})$. All errors are uniform on the middle interval.

For bounded $y$, substituting $k=k_0(1+y/\sqrt M)$ into (1) gives
\[
 D(k)+\frac{Ln}{2}=2L-\frac{L^2}{6}+\frac L2y+o(1). \tag{2}
\]
For clarity, the first term of (1) at $k_0$ is $-L(m-1)/2=-Ln/2+L$, while the third is $L$. The normalized base weight is
\[
 \frac{e^{-Lk}k^M}{e^{-M}k_0^M}
 =\exp\!\left(M[\log(1+y/\sqrt M)-y/\sqrt M]\right).
\]
It tends to $e^{-y^2/2}$ uniformly on compact sets and is bounded on the entire middle interval by $e^{-c y^2}$ for some $c>0$. Equation (1) likewise gives $D(k)+Ln/2\le C(1+|y|)$ there: the change of its first term from $k_0$ is $O(n|k-k_0|/k_0)=O(|y|)$ and the other terms are bounded. The mesh in $y$ is $1/\sigma\to0$. The integrable Gaussian bound $e^{-cy^2+C(1+|y|)}$ controls the lattice tails uniformly; hence the normalized middle sum is a Riemann sum with limit
\[
 \int_{\mathbb R}e^{-y^2/2+Ly/2+2L-L^2/6}\,dy
 =\sqrt{2\pi}\,\exp(2L-L^2/24). \tag{3}
\]
Here the normalized sum is the sum of $e^{-Lk}B(k,m)^n$ divided by
$\sigma e^{-M}k_0^M e^{-Ln/2}/(m!)^n$.

\textbf{Outer tails.} By Vandermonde, $B(k,m)\le\binom{k+m}{m}\le(k+m)^m/m!$. With $h=k+m$, the normalized outer tails are bounded by $e^{Lm+Ln/2}/\sigma$ times the corresponding sums of $e^{-Lh}h^M/(e^{-M}k_0^M)$. For $k<k_0/2$, eventually $h/k_0\le 3/4$, and there are $O(M)$ terms, each $e^{-\Omega(M)}$ by $\log u-u+1<0$. For $k>2k_0$, $h/k_0>2$ and the base is again $e^{-\Omega(M)}$ at the first term. The log weight has derivative $M/h-L\le -L/2$, giving a geometric tail. These estimates survive the $e^{O(n)}$ factor.


\subsection*{Evaluation of the multiplier}
Combining (3) and the tail estimates, and using $\sqrt{2\pi}\sigma e^{-M}k_0^M\sim M!/L^{M+1}$, we get
\[
 a(n)\sim 2^{n-1}\frac{M!}{L^{M+1}(m!)^n}
                  \exp(-Ln/2+2L-L^2/24).
\]
The additive $1$ in the tree identity is negligible (for example $I\ge1$ and $2^n-2\to\infty$). The uniform Stirling calculation here needs the first correction for the \emph{denominator} factorial: an error of order $1/m$ accumulates $n$ times. Since $M=nm$ and $m=n-1$,
\[
 \frac{M!}{(m!)^n}
 \sim\sqrt{2\pi M}\,(2\pi m)^{-n/2}n^{M}
       \exp\!\left(-\frac{n}{12m}\right).
\]
In fact $\log(j!)=(j+\tfrac12)\log j-j+\tfrac12\log(2\pi)+1/(12j)+O(j^{-3})$; here $n/m^3\to0$. Now
\[
 m^{-n/2}\sim n^{-n/2}e^{1/2},\qquad
 \sqrt M\sim n,\qquad e^{-n/(12m)}\to e^{-1/12}.
\]
Finally $2^{n-1}\exp(-Ln/2+2L)=2^{1+n/2}$. Cancelling against $(2\pi)^{-n/2}$ gives the formula in the theorem. $\square$

\textbf{Source status.} The November 9 entry called the root limit a conjecture. The November 13 entry \emph{asserts} the stronger equivalent and displays a decimal for the multiplier; it does not label the stronger assertion a conjecture. We do not claim that the author had no proof, or claim priority for the formula. The cited public proof of the root limit explicitly does not cover the equivalent. This note supplies an independent proof with a closed-form constant. Exact finite differences reproduce small entries of A263160; numerical saddle windows are illustrative only.

\subsection*{References}
\small\raggedright
[1] A.~P. Heinz, \href{https://oeis.org/A263159}{OEIS A263159} and
\href{https://oeis.org/A263160}{A263160}, \emph{The On-Line Encyclopedia of Integer Sequences}; V. Kote\v{s}ovec, Nov. 13, 2026, asymptotic formula,
\href{https://oeis.org/A263160/history}{revision history}.

[2] Trureturing, \href{https://www.trureturing.org/research/oeis-a263160-first-zero-diagonal-limit/}{First-zero lattice paths and the A263160 diagonal limit}; see the
\href{https://the-omega-institute.github.io/trureturing-mdbook/Library/FiniteGeometry/kotesovec2026a263160.html}{dated source comparison} for the two revisions, full stronger formula, and scope of the historical limit proof.

[3] S. Eger, \href{https://arxiv.org/abs/1511.00622}{On the Number of Many-to-Many Alignments of Multiple Sequences}, section 5.5, for the binomial moment for endpoint counts. Summing those endpoint coefficients over a box is not claimed as new.

\par\endgroup

% Source: /workspace/scratch/81914cec2c78/work/archive/Tournament_Round_39_Final/Tournament_Round_39/round39/showa/paper.tex
% Edited source: /workspace/scratch/81914cec2c78/work/audit_34_41/cleaned/r39_showa.tex
\begingroup

\Needspace{9\baselineskip}
\section{\texorpdfstring{Complementary points in three triangle inequalities}{Complementary points in three triangle inequalities}}\label{section:r39showa}
\noindent{\small\sffamily\color{muted}Round 39 \enspace / \enspace Showa}\par\medskip


\noindent\textbf{Status.} The three inequalities from Chu--Zhang (2005), Section 4, are proved below on their full domains as consequences of known inequalities.

\subsection{The complementary point}
Write $A_i$ for the three vertices, $a_i$ for the opposite sides, and $R_i,r_i,h_i$ for the distances from an interior $P$ to $A_i$, from $P$ to the opposite line, and from $A_i$ to that line, respectively. Let $K_i$ be the distance between the midpoints of $PA_i$ and the opposite side. The complementary point $Q=(A_1+A_2+A_3-P)/2$ has $QA_i=K_i$. If $P$ has normalized barycentrics $u_i$, those of $Q$ are $(1-u_i)/2$, all strictly positive. Thus twice the distance from $Q$ to line $i$ is $h_i-r_i$.

With complex coordinates $z_i,p,q$ for these points, put $d_i=\prod_{j\ne i}(z_i-z_j)$, so $|d_i|=\prod a_j/a_i$. Lagrange interpolation of the monic quadratic $(z-p)(z-q)$ gives, by comparing leading coefficients,
\[
 1=\sum_i\frac{(z_i-p)(z_i-q)}{d_i}.
\]
The triangle inequality proves $\sum a_iR_iK_i\ge\prod a_i$, the proposed (4.1). Interpolating the constant 1 at $q$ instead yields
\[
 1=\sum_i\frac{\prod_{j\ne i}(q-z_j)}{d_i},
 \qquad \sum_i a_i\prod_{j\ne i}K_j\ge\prod_i a_i,
\]
which is (4.2). Finally Erd\H{o}s--Mordell at $Q$ gives $\sum K_i\ge\sum(h_i-r_i)$, proving (4.3). The underlying first two inequalities occur in complex geometry references, and the third is a direct application of the classical theorem; none is claimed as a new open-problem solution.

\paragraph{Reference.} Chu--Zhang, \emph{An Klamkin Type Inequality on the Triangle} (2005), Section 4, equations (4.1)--(4.3). Primary source: \texttt{https://rgmia.org/papers/v8n1/cxgzzh01.pdf}.

\par\endgroup

% Source: /workspace/scratch/81914cec2c78/work/archive/Tournament_Round_39_Final/Tournament_Round_39/round39/heisei/paper.tex
% Edited source: /workspace/scratch/81914cec2c78/work/audit_34_41/cleaned/r39_heisei.tex
\begingroup

\Needspace{9\baselineskip}
\section{\texorpdfstring{A uniform disk bound for the diagonal Pad\'e approximants to the wave kernel}{A uniform disk bound for the diagonal Pad\'e approximants to the wave kernel}}\label{section:r39heisei}
\noindent{\small\sffamily\color{muted}Round 39 \enspace / \enspace Heisei}\par\medskip

\begin{sectionabstract}
For every $m\ge0$, the diagonal Pad\'e approximant $r_m$ to $\cosh\sqrt z$ maps the closed disk of radius three into the closed disk of radius two centered at one. Equality is attained only at $m=1$, $z=3$. This proves Conjecture 4.6 of Nadukandi and Higham. The proof expresses the denominator coefficients as normalized Schur functions and bounds them by deleting a row of a tableau. In this way only the tail of the reciprocal zeros enters the denominator estimate.
\end{sectionabstract}

\subsection{A tail estimate for Pad\'e denominators}
Let $x_1\ge x_2\ge\cdots>0$ be summable and define
\[
 F(z)=\prod_{j\ge1}(1+x_jz)=\sum_{i\ge0}e_i(x)z^i,\qquad T_m=\sum_{j>m}x_j.
\]
Here $e_i$ is the elementary symmetric function, $e_0=1$. An $[m/m]$ Pad\'e approximant at zero is a quotient $P_m/Q_m$ of polynomials of degree at most $m$, with $Q_m(0)=1$ and $Q_mF-P_m=O(z^{2m+1})$.

\begin{theorem}\label{r39heisei:tail}
Such a normalized pair exists uniquely. If $Q_m=\sum_{k=0}^m q_kz^k$, then for $0\le k\le m$
\[
 0<(-1)^kq_k\le h_k(x_{m+1},x_{m+2},\ldots)\le T_m^k ,
\]
where $h_k$ is the complete homogeneous symmetric function. For $R>0$ with $RT_m<1/2$ the denominator has no zero in $|z|\le R$ and
\begin{equation}
 \max_{|z|\le R}\left|F(z)-\frac{P_m(z)}{Q_m(z)}\right|
 \le\frac{1}{1-2RT_m}\sum_{i\ge m+1} e_i(x)R^i.
 \label{r39heisei:error}
\end{equation}
\end{theorem}
\begin{proof}
The case $m=0$ is immediate (the tail bound on the error is deliberately not optimized). Assume $m\ge1$. At degrees $m+1,\ldots,2m$ the Pad\'e equations are
\[
 \sum_{c=0}^m e_{m+j-c}(x)q_c=0,\qquad j=1,\ldots,m.
\]
Write $D_k$ for the determinant obtained by deleting column $k$ in this $m$ by $m+1$ system. Let $c_1<\cdots<c_m$ be the surviving column numbers. The transpose of that minor has entries $e_{m+j-c_i}$, and $c_i=i-1$ for $i\le k$, $c_i=i$ for $i>k$. By the dual Jacobi--Trudi identity
\[
 D_k=s_{(m^m,k)}(x).
\]
Indeed, the conjugate of the displayed partition has $k$ parts of size $m+1$ and $m-k$ of size $m$, so that the determinant entries are $e_{\lambda'_i-i+j}$. The notation $(m^m,k)$ means $m$ rows of length $m$ followed by a row of length $k$, with a zero last part omitted. The tableau formula makes $D_0=s_{(m^m)}>0$. The signed-minor null vector therefore gives $q_k=(-1)^kD_k/D_0$, uniquely after normalization. The coefficients of $P_m$ are then fixed by truncating $Q_mF$ through degree $m$.

Delete the bottom row of length $k$ from a semistandard tableau of shape $(m^m,k)$. What remains is a tableau of the $m$ by $m$ square. The removed row is weakly increasing, and all its entries are at least $m+1$ because there were $m$ strictly smaller entries above each one. The pair comprised of the remaining tableau and the row recovers the original tableau uniquely. Taking products of entry weights and summing yields
\[
 D_k\le D_0\sum_{m<i_1\le\cdots\le i_k}x_{i_1}\cdots x_{i_k}
 =D_0 h_k(x_{m+1},x_{m+2},\ldots)\le D_0T_m^k.
\]
For $k=0$ the empty row has weight one. For all $k$ positivity follows by using row-constant tableaux with row entries equal to the row number. The identities can first be read in finitely many variables: as the alphabet increases the elementary symmetric functions converge by summability, the determinants converge entrywise, and the nonnegative tableau sums converge monotonically to the same limits.

Set $\eta=RT_m<1/2$. On the closed disk $|z|\le R$,
\[
 |Q_m(z)|\ge1-\frac{\eta}{1-\eta}=\frac{1-2\eta}{1-\eta},
 \qquad \sum_{k=0}^m|q_k|R^k\le\frac1{1-\eta}.
\]
Only degrees at least $2m+1$ remain in $Q_mF-P_m$. Such a coefficient from $q_kz^kF$ uses an index $i\ge2m+1-k\ge m+1$ of $F$. Thus the absolute numerator error is at most $(1-\eta)^{-1}\sum_{i\ge m+1}e_iR^i$. Dividing by the lower bound for $|Q_m|$ proves \eqref{r39heisei:error}.
\end{proof}

\subsection{The wave kernel}
\begin{corollary}[Nadukandi--Higham, Conjecture 4.6]
For every $m\ge0$, the $[m/m]$ Pad\'e approximant $r_m$ to the entire function $f(z)=\cosh\sqrt z$ satisfies $|r_m(z)-1|\le2$ on $|z|\le3$. The inequality is strict except when $m=1$ and $z=3$.
\end{corollary}
\begin{proof}
Euler's cosine product gives $f(z)=\prod_{j\ge1}(1+x_jz)$ with $x_j=4/[\pi^2(2j-1)^2]$. Its coefficients are $e_i(x)=1/(2i)!$, so $\sum_jx_j=1/2$. For $m\ge2$ one has
\[
  T_m\le T_2=\frac12-\frac{40}{9\pi^2}<\frac1{18},
\]
using $\pi^2<10$. At $R=3$ it follows that $\eta<1/6$ and $1/(1-2\eta)<3/2$. To bound the coefficient tail put
\[
 U=\sum_{i\ge3}\frac{3^i}{(2i)!}
 \le \frac{27}{720}\frac1{1-3/56}<\frac1{25};
\]
each ratio of successive terms in this tail is at most $3/56$. By \eqref{r39heisei:error}, for $m\ge2$,
\[
 |r_m(z)-1|
 \le |f(z)-1|+|r_m(z)-f(z)|
 \le\frac32+\frac9{24}+U+\frac32 U
 <\frac{79}{40}<2 .
\]
For $m=0$, $r_0=1$. Direct matching gives
\[
 r_1(z)-1=\frac{z}{2(1-z/12)}.
\]
For $|z|\le3$ this has modulus at most $3/[2(1-3/12)]=2$, with equality precisely at $z=3$.
\end{proof}

\paragraph{Relation to earlier work.} Nadukandi and Higham verified the range $m\le20$ before stating Conjecture 4.6 \cite{r39heisei:nh}. The present proof does not need that finite verification. Dzyadyk's earlier work \cite{r39heisei:dz} treats asymptotics of approximants to trigonometric and hyperbolic functions. An asymptotic statement without explicit bounds down to a controlled finite degree is distinct from the disk assertion proved here. No priority for Schur determinant representations or Pad\'e convergence in general is suggested.

\begin{thebibliography}{9}
\bibitem{r39heisei:nh} P. Nadukandi and N. J. Higham, Computing the wave-kernel matrix functions, \emph{SIAM J. Sci. Comput.} \textbf{40} (2018), Conjecture 4.6 after Lemma 4.5. \href{https://doi.org/10.1137/18M1170352}{doi:10.1137/18M1170352}.
\bibitem{r39heisei:dz} V. K. Dzyadyk, Asymptotics of diagonal Pad\'e approximants of the functions $\sin z$, $\cos z$, $\sinh z$ and $\cosh z$, \emph{Math. USSR-Sb.} \textbf{36} (1980), 231--249. \href{https://doi.org/10.1070/SM1980v036n02ABEH001802}{doi:10.1070/SM1980v036n02ABEH001802}.
\bibitem{r39heisei:mac} I. G. Macdonald, \emph{Symmetric Functions and Hall Polynomials}, second ed., Oxford University Press, 1995, Chapter I, Sections 3 and 5 (dual Jacobi--Trudi identity and tableau expansion).
\end{thebibliography}

\par\endgroup

% Source: /workspace/scratch/81914cec2c78/work/archive/Tournament_Round_39_Final/Tournament_Round_39/round39/reiwa/paper.tex
% Edited source: /workspace/scratch/81914cec2c78/work/audit_34_41/cleaned/r39_reiwa.tex
\begingroup

\Needspace{9\baselineskip}
\section{\texorpdfstring{A lower bound for the difference between the sums of medians and altitudes}{A lower bound for the difference between the sums of medians and altitudes}}\label{section:r39reiwa}
\noindent{\small\sffamily\color{muted}Round 39 \enspace / \enspace Reiwa}\par\medskip


\textbf{Theorem.} In every nondegenerate Euclidean triangle with semiperimeter
$s$, inradius $r$, medians $m_a,m_b,m_c$ and corresponding altitudes
$h_a,h_b,h_c$, we have
\[
 \sum_{\rm cyc}(m_a-h_a)\ \geq\ s-3\sqrt3\,r .
\]
This proves the inequality posed by J.~Liu in Conjecture 4, (4.7), of [1].
An equilateral triangle attains equality.

\textbf{Proof.} Write $t=\sqrt3$. The main device is a rational lower
bound for each median. Put
\[
 H=b+c,\quad P=2b^2+2c^2-a^2=4m_a^2,\quad
 q=P/H^2,\quad w=(2q+t)/(2+t).
\]
The identity $q+w^2-2w\sqrt q=(w-\sqrt q)^2$ gives
\[
 2m_a=H\sqrt q\ \geq\ H\frac{2qw}{q+w^2}.
\]
The choice of $w$ makes this bound exact both at $q=3/4$ (the equilateral
value) and at $q=1$ (a limiting value). For use without nested fractions put
\[
 N=2P+tH^2,\qquad T_a=2PH(2+t)N,\qquad
 B_a=P H^2(2+t)^2+N^2.
\]
The rational lower bound above equals $T_a/B_a$, with $B_a>0$.
Define $d=B_aB_bB_c$ and $n=\sum_{\rm cyc}T_aB_bB_c$, so
$2\sum m_a\ge n/d$.

In tangent lengths $x=s-a$, $y=s-b$, $z=s-c$ the area is $\sqrt D$,
where $D=sxyz$ by Heron's formula; $s=x+y+z$ and $a=y+z$ cyclically.
Write $J=abc$, $Q=ab+bc+ca$. The desired inequality follows from
\[
 (n-2sd)sJ\ \geq\ 2\sqrt D\,d(2sQ-3tJ),
\]
using $h_a=2\sqrt D/a$ and $r=\sqrt D/s$. It suffices to prove
\[
 K=n-2sd\ge0,\qquad
 F=(sJ K)^2-4Dd^2(2sQ-3tJ)^2\ge0.
\]
Indeed the second inequality compares the squares, and the first guarantees
the required sign for the left side. These are homogeneous symmetric
inequalities: after relabeling and scaling, write
$x=1$, $y=1+u$, $z=1+u+v$, $u,v\ge0$. The following exact coefficient
verification completes the proof. \hfill$\square$


\textbf{Coefficient verification.} In the above substitution, each nonzero
coefficient is positive in the real embedding of $\mathbb Z[t]$:
\[
\begin{array}{lrrr}
 &\text{maximum total degree}&\text{positive coefficients}&\text{negative coefficients}\\
 K&13&105&0\\ F&34&622&0
\end{array}
\]
Here is complete Python 3 code (with SymPy) for the finite arithmetic in the table.
Each coefficient $A+Bt$ is stored as two integers; the comparison uses only
integer operations. No sampled evaluations are used to infer positivity.

\begin{Verbatim}
"""Exact coefficient certificate for Liu (2012), Conjecture 4. Python 3 + SymPy."""
from sympy import ZZ
from sympy.polys.rings import ring
B, u, v = ring('u,v', ZZ)
# Pair (a,b) denotes a+b sqrt(3). Coefficients are integer polynomials.
def add(A,B): return A[0]+B[0], A[1]+B[1]
def mul(A,B): return A[0]*B[0]+3*A[1]*B[1], A[0]*B[1]+A[1]*B[0]
def sc(A,k): return A[0]*k,A[1]*k
def sq(A): return mul(A,A)
def certificates(x,y,z):
    s=x+y+z; D=s*x*y*z
    a,b,c=y+z,z+x,x+y; J=a*b*c; Q=a*b+b*c+c*a
    n,d=(0,0),(1,0)
    for A,B,C in [(a,b,c),(b,c,a),(c,a,b)]:
        H=B+C; P=2*B**2+2*C**2-A**2; N=(2*P,H**2)
        top=sc(mul(N,(2,1)),2*P*H)
        bot=add(sc(sq((2,1)),P*H**2),sq(N))
        n,d=add(mul(n,bot),mul(top,d)),mul(d,bot)
    K=add(n,sc(d,-2*s))
    F=add(sc(sq(K),(s*J)**2),sc(mul(sq(d),sq((2*s*Q,-3*J))),-4*D))
    return K,F
def sign(a,b):
    if not a and not b: return 0
    if a>=0 and b>=0: return 1
    if a<=0 and b<=0: return -1
    k=a*a-3*b*b
    return (1 if k>0 else -1)*(1 if a>0 else -1)
allcert=certificates(B.one,1+u,1+u+v)
for name,(p,q),expect in zip(('K','F'),allcert,[105,622]):
    support=set(p)|set(q)
    signs=[sign(p.get(e,0),q.get(e,0)) for e in support]
    assert len(support)==expect and min(signs)==1
    print(name, 'positive coefficients', len(support), 'negative coefficients', 0,
          'max total degree', max(i+j for i,j in support))
# Check exact polynomial evaluation at independent encoded integers against
# the same construction done directly over integer pairs.
bits=max(abs(int(c)).bit_length() for C in allcert for p in C for c in p.values())+3
stride=1+max(i for C in allcert for p in C for i,j in p)
U,V=1<<bits, 1<<(stride*bits)
for C,E in zip(allcert, certificates(1,1+U,1+U+V)):
    for p,q in zip(C,E):
        assert sum(int(c)<<(bits*(i+stride*j)) for (i,j),c in p.items())==q
print('Integer-pair Kronecker checks passed; stride',stride,'bits',bits)
\end{Verbatim}


\textbf{Reference.} [1] Jian Liu, On an inequality for the medians of a
triangle, \emph{Journal of Science and Arts} 12, no.~2(19) (2012), 127--136.
See Conjecture 4, p.~135, Eq.~(4.7).
\url{https://www.josa.ro/docs/josa_2012_2/a_03_Liu_J.pdf}.

\par\endgroup

% Source: /workspace/scratch/81914cec2c78/work/archive/Tournament_Round_39_Final/Tournament_Round_39/round39/keio/paper.tex
% Edited source: /workspace/scratch/81914cec2c78/work/audit_34_41/cleaned/r39_keio.tex
\begingroup
\newcommand{\cyc}{\sum_{\mathrm{cyc}}}
\Needspace{9\baselineskip}
\section{\texorpdfstring{A weighted triangle graph bounds the altitudes of a cevian triangle}{A weighted triangle graph bounds the altitudes of a cevian triangle}}\label{section:r39keio}
\noindent{\small\sffamily\color{muted}Round 39 \enspace / \enspace Keio}\par\medskip


\begin{sectionabstract}
We prove Jian Liu's conjectured quadratic inequality comparing three exradii with the altitudes of a cevian triangle. Each normalized altitude is bounded by the corresponding entry of the normalized adjacency matrix of a triangle whose edge weights are the barycentric coordinates of the point of concurrence. A sum of three squares proves the inequality for all real weights and identifies all cases of equality.
\end{sectionabstract}
Let $ABC$ be a nondegenerate Euclidean triangle of area $S$ and side lengths $a=BC$, $b=CA$, $c=AB$. For a point $P$ in its interior, put $L=AP\cap BC$, $M=BP\cap CA$, $N=CP\cap AB$. Write $h_l,h_m,h_n$ for the altitudes of $LMN$ opposite $L,M,N$. In particular $h_l$ is the distance from $L$ to the \emph{line} $MN$. Write $r_a,r_b,r_c$ for the exradii of $ABC$.

\begin{theorem}[Liu's Conjecture 5]
For all $x,y,z\in\mathbb R$ one has
\[
     r_ax^2+r_by^2+r_cz^2\ \ge\ 2(h_l yz+h_m zx+h_n xy).
\]
Apart from $x=y=z=0$, equality occurs exactly when $ABC$ is equilateral, $P$ is its center, and $x=y=z$.
\end{theorem}
The inequality was proposed in [1, p.~118, (3.10)]. The adjacent classical weighted Erd\H os--Mordell inequality involves distances from $P$ and does not directly give this claim.

\begin{lemma}
Let $(u:v:w)$ be the positive barycentric coordinates of $P$. Then
\[
 \frac{h_l}{\sqrt{r_br_c}}\le
 \frac{2\sqrt{vw}}{v+w}\,\frac{u}{\sqrt{(u+v)(u+w)}}.
\]
The other two bounds follow by cyclic permutation.
\end{lemma}

\begin{proof}
The barycentric coordinate rows of $L,M,N$, normalized to sum to one, are
\[
 \left(0,\frac{v}{v+w},\frac{w}{v+w}\right),\quad
 \left(\frac{u}{u+w},0,\frac{w}{u+w}\right),\quad
 \left(\frac{u}{u+v},\frac{v}{u+v},0\right).
\]
Their determinant gives the ratio of oriented areas; direct expansion gives
\[
         [LMN]=\frac{2Suvw}{(u+v)(u+w)(v+w)}.
\]
\noindent
Let $\alpha=w/(u+w)$ and $\beta=v/(u+v)$. Then $M-A=\alpha(C-A)$ and $N-A=\beta(B-A)$. Consequently,
\[
  |MN|^2=(\alpha b-\beta c)^2+4\alpha\beta bc\sin^2(A/2).
\]
If $s=(a+b+c)/2$, the half-angle identity $bc\sin^2(A/2)=(s-b)(s-c)$ and $r_b=S/(s-b)$, $r_c=S/(s-c)$ imply
\[
   \sqrt{r_br_c}=\frac{S}{\sqrt{bc}\sin(A/2)}.
\]
Inserting $|MN|\ge2\sqrt{\alpha\beta bc}\sin(A/2)$ in $h_l=2[LMN]/|MN|$ and dividing by this last expression proves the lemma.
\end{proof}

\begin{proof}[Proof of the theorem]
Write $X=|x|\sqrt{r_a}$, $Y=|y|\sqrt{r_b}$, $Z=|z|\sqrt{r_c}$ and set
\[
           d_a=v+w,\qquad d_b=w+u,\qquad d_c=u+v.
\]
They are the vertex degrees of a triangle graph whose edges $BC,CA,AB$ carry weights $u,v,w$, respectively. By the lemma and $2\sqrt{vw}\le v+w$,
\[
 \frac{h_l}{\sqrt{r_br_c}}\le\frac{u}{\sqrt{d_bd_c}}.
\]
This is an entrywise bound, not a claim that arbitrary entrywise comparisons preserve positive semidefiniteness. The signs have already been removed, so that the actual quadratic deficit has the following explicit lower bound:
\begin{align*}
 \cyc r_a x^2-2\cyc h_l yz
 &\ge X^2+Y^2+Z^2-2\cyc\frac{uYZ}{\sqrt{d_bd_c}}\\
 &=\cyc u\left(\frac{Y}{\sqrt{d_b}}-\frac{Z}{\sqrt{d_c}}\right)^2\ge0.
\end{align*}
The middle equality follows by collecting the three square terms and using the definitions of the degrees.

If equality holds and the weights are not all zero, the three ratios $X/\sqrt{d_a},Y/\sqrt{d_b},Z/\sqrt{d_c}$ are equal and strictly positive. The three sign replacements therefore require $x,y,z$ to have a common sign. Moreover each altitude bound must be tight, so the three arithmetic--geometric mean comparisons give $u=v=w$. Thus $P$ is the centroid. Now $\alpha=\beta=1/2$, and equality in each side-norm bound forces $b=c$ and its cyclic analogues. Hence $ABC$ is equilateral, and all exradii and all three degrees are equal, giving $x=y=z$. Conversely these data give equality throughout.
\end{proof}
\begin{thebibliography}{1}
\bibitem{r39keio:liu} J. Liu, A new proof of the Erd\H os--Mordell inequality, \emph{International Electronic Journal of Geometry} \textbf{4} (2011), no.~2, 114--119. \url{https://dergipark.org.tr/en/download/article-file/775559}.
\end{thebibliography}

\par\endgroup

% Source: /workspace/scratch/81914cec2c78/work/archive/Tournament_Round_40_Final/Tournament_Round_40/round40/meiji/paper.tex
% Edited source: /workspace/scratch/81914cec2c78/work/audit_34_41/cleaned/r40_meiji.tex
\begingroup

\Needspace{9\baselineskip}
\section{\texorpdfstring{A sharp refinement of the Lessels--Pelling inequality}{A sharp refinement of the Lessels--Pelling inequality}}\label{section:r40meiji}
\noindent{\small\sffamily\color{muted}Round 40 \enspace / \enspace Meiji}\par\medskip


\begin{sectionabstract}
We prove Liu's conjectured bound for one median and the two remaining angle bisectors of an arbitrary triangle. In place of the classical upper bound $\sqrt3$ times the semiperimeter, one may take $\tfrac32\sqrt{ab+bc+ca}$, strictly less for every nonequilateral shape. The proof identifies the desired radical sum as the largest root of a quartic and certifies that every shifted coefficient is nonnegative in four ordered shape charts.
\end{sectionabstract}

\begin{theorem}
In any nondegenerate triangle $ABC$ with side lengths $a,b,c$, median $m_a$ from $A$, and internal angle bisectors of lengths $w_b,w_c$ from $B,C$, one has
\[
 m_a+w_b+w_c\le\frac32\sqrt{ab+bc+ca}.
 \tag{1}
\]
Equality holds exactly for equilateral triangles, so the factor $3/2$ is best possible.
\end{theorem}
This is Conjecture 1, Eq.~(71), of Liu \cite[p.~1373]{r40meiji:liu}. Let $s=(a+b+c)/2$ and $p=ab+bc+ca$. The Lessels--Pelling inequality \cite{r40meiji:lp} gives $m_a+w_b+w_c\le\sqrt3s$; see also the later weighted refinements \cite{r40meiji:bd}. Its upper bound is larger than the one in (1) except at the equilateral triangle, since $4s^2-3p=\tfrac12\sum_{\rm cyc}(a-b)^2$.

\paragraph{A polynomial for the desired branch.}
Let $x=s-a$, $y=s-b$, $z=s-c$, so $a=y+z$, $b=z+x$, $c=x+y$ and $s=x+y+z$. Set $d=(a+b)^2(a+c)^2$. Apollonius's median identity and the bisector identity $w_b^2=4ac\,s(s-b)/(a+c)^2$ give
\[
\begin{aligned}
 X&=4d\,m_a^2=\{4sx+(y-z)^2\}d,\\
 Y&=4d\,w_b^2=16acsy(a+b)^2,\qquad
 Z=4d\,w_c^2=16absz(a+c)^2,\\
 W&=9pd .
\end{aligned}
\tag{2}
\]
These are positive. Inequality (1) is equivalent to $\sqrt X+\sqrt Y+\sqrt Z\le\sqrt W$. Put $L=X+Y+Z$ and $S=XY+YZ+ZX$. The polynomial
\[
 H(t)=\bigl((t-L)^2-4S\bigr)^2-64XYZt \tag{3}
\]
has as its four roots
\[
 (\sqrt X+\epsilon\sqrt Y+\delta\sqrt Z)^2,\qquad \epsilon,\delta\in\{-1,1\} .
\]
Indeed the monic product of these four factors $t-(\sqrt X+\epsilon\sqrt Y+\delta\sqrt Z)^2$ expands to (3). The largest root is $(\sqrt X+\sqrt Y+\sqrt Z)^2$. Testing only $H(W)$ would not locate $W$ relative to all four roots. Instead we exclude positive roots of $H(W+T)$.

For $V=W-L$ and $F=V^2-4S$, the coefficients of this shifted polynomial, in ascending powers of $T$, are
\[
 \underbrace{F^2-64XYZW}_{C_0},\quad
 \underbrace{4VF-64XYZ}_{C_1},\quad
 \underbrace{2F+4V^2}_{C_2},\quad
 \underbrace{4V}_{C_3},\quad 1. \tag{4}
\]

\paragraph{Covering every shape.}
Since (1) is homogeneous, rescale the shortest tangency length to one. Write the sorted lengths as $1,1+g,1+g+h$, $g,h\ge0$. The length $x$ is distinguished, but $y,z$ can be interchanged. Thus we can assign $(x,y,z)$ by rotating the three sorted lengths until $x$ has the required rank. For smallest and largest $x$, take $(g,h)=(u,v)$. For middle $x$, subdivide into $(g,h)=(u,u+v)$ and $(u+v,v)$. In every chart $u,v\ge0$ without restriction.

The following table gives counts of nonzero coefficients of each $C_j$ in $\mathbb{Z}[u,v]$. All these nonzero coefficients are \emph{positive}. In each chart $C_0$ has constant coefficient zero and contains both $u^2$ and $v^2$ with positive coefficients. The exact expansion code below proves these statements.
\[
\begin{array}{c|c|rrrr}
 x\text{ rank}&(g,h)&C_0&C_1&C_2&C_3\\\hline
\text{smallest}&(u,v)&319&189&91&28\\
\text{middle}&(u,u+v)&319&189&91&28\\
\text{middle}&(u+v,v)&322&190&91&28\\
\text{largest}&(u,v)&322&190&91&28
\end{array}
\]
It follows that $H(W+T)>0$ for all $T>0$. When $(u,v)\ne(0,0)$, the same holds at $T=0$ by the two pure squares in $C_0$. Thus the largest root of $H$ is strictly less than $W$ for nonequilateral triangles. For an equilateral triangle of side $a$, all three lengths in the left side of (1) equal $\sqrt3 a/2$, giving equality. This proves the theorem.

\paragraph{Consequences listed by Liu.}
Write $M=m_a+w_b+w_c$. For an interior point $P$, let $R_a,R_b,R_c$ denote the circumradii of $BPC,CPA,APB$. Liu's (6) gives $\sum R_i\ge\sqrt p$, hence $\sum R_i\ge 2M/3$. Next (1) gives $M\le\sqrt3s$ as above and $M\le3(R+r)$ by $p\le4(R+r)^2$ \cite[Eq.~(62)]{r40meiji:liu}, where $R,r$ are circumradius and inradius. Finally
\[
  2M\le 3\sqrt p\le w_a+w_b+w_c+r_a+r_b+r_c
\]
by Liu's (9), where $r_a,r_b,r_c$ are the exradii. Subtraction gives $2m_a+w_b+w_c\le w_a+r_a+r_b+r_c$. These recover (72)--(75) of the source as consequences; the semiperimeter bound (73) is classical, and they are not being counted as four independent discoveries.

\begin{thebibliography}{9}
\bibitem{r40meiji:liu} J. Liu, Proof of a double inequality in triangles, \emph{Journal of Mathematical Inequalities} \textbf{15} (2021), 1361--1374. \href{https://doi.org/10.7153/jmi-2021-15-92}{doi:10.7153/jmi-2021-15-92}.
\bibitem{r40meiji:lp} G.~S. Lessels and M.~J. Pelling, An inequality for the sum of two angle bisectors and a median, \emph{Univ. Beograd Publ. Elektrotehn. Fak. Ser. Mat. Fiz.} 577--598 (1977), 59--62.
\bibitem{r40meiji:bd} M. Bencze and M. Dr\u{a}gan, Two refinements of power Lessels--Pelling inequality, \emph{Arhimede Mathematical Journal} \textbf{5} (2018), 110--117.
\end{thebibliography}

\paragraph{Complete certificate (Python 3 with SymPy).}
The operations are additions and multiplications in the integer polynomial ring. No numerical sign tests or external files are needed. Output fields are the rank (starting at zero), chart number, index $j$, monomial count and least nonzero coefficient.
\begingroup\fontsize{8}{9.7}\selectfont
\begin{Verbatim}
"""Exact expansion over ZZ[u,v] (requires SymPy)."""
from sympy.polys.rings import ring
from sympy.polys.domains import ZZ
Q,u,v=ring("u,v",ZZ)
for rank,case in [(0,0),(1,1),(1,2),(2,0)]:
    g,h=(u,v) if case==0 else ((u,u+v) if case==1 else (u+v,v))
    tang=[Q.one,1+g,1+g+h]
    x,y,z=tang[rank:]+tang[:rank]
    a,b,c=y+z,z+x,x+y
    s=x+y+z
    d=(a+b)**2*(a+c)**2
    X=(4*s*x+(y-z)**2)*d
    Y=16*a*c*s*y*(a+b)**2
    Z=16*a*b*s*z*(a+c)**2
    W=9*(a*b+a*c+b*c)*d
    V=W-X-Y-Z; F=V**2-4*(X*Y+Y*Z+Z*X)
    polys=[F**2-64*X*Y*Z*W,4*V*F-64*X*Y*Z,2*F+4*V**2,4*V]
    for j,p in enumerate(polys):
        assert min(p.values())>0
        if j==0:
            assert p.get((0,0),0)==0
            assert p.get((2,0),0)>0 and p.get((0,2),0)>0
        print(rank,case,j,len(p),min(p.values()))
\end{Verbatim}
\endgroup

\par\endgroup

% Source: /workspace/scratch/81914cec2c78/work/archive/Tournament_Round_40_Final/Tournament_Round_40/round40/taisho/paper.tex
% Edited source: /workspace/scratch/81914cec2c78/work/audit_34_41/cleaned/r40_taisho.tex
\begingroup
\newcommand{\E}{\mathbb E}
\newcommand{\PP}{\mathbb P}
\Needspace{9\baselineskip}
\section{\texorpdfstring{Local asymptotics for reciprocal partition weights}{Local asymptotics for reciprocal partition weights}}\label{section:r40taisho}
\noindent{\small\sffamily\color{muted}Round 40 \enspace / \enspace Taisho}\par\medskip

The first result proves Kote\v{s}ovec's two 2020 conjectures on tangent compositions of partition products, and identifies the signed logarithmic correction in the unrestricted case. The second result proves Lagarias and Sun's three conjectures for partition supernorms at \emph{exact} size or perimeter. Their known cumulative size asymptotic is an input, not a substitute for the local estimate.

\subsection*{Two tangent compositions}
Set $r=\pi/4$, write $\gamma$ for Euler's constant, and define
\[
 F_-(w)=\prod_{j\ge1}(1-w^j/j)^{-1},\qquad
 F_+(w)=\prod_{j\ge1}(1+w^j/j),\qquad
 a_\pm(n)=n![x^n]F_\pm(\tan x).
\]
\textbf{Theorem 1.} The unrestricted and distinct cases, respectively, satisfy
\[
 \frac{r^na_-(n)}{n!}=\frac{e^{-\gamma}}{(2r)^2}
               (n-2r\log n+o(\log n)),\qquad
 \frac{r^na_+(n)}{n!}\longrightarrow\frac{e^{-\gamma}}{2r}.
\]
The resulting leading terms are exactly the separately posted conjectures for OEIS A335636 and A335638 [1].

\textit{Factorization.} Let $Q_\pm(w)=\sum_{i,j\ge2}\epsilon_{\pm,i}w^{ij}/(i j^i)$, where $\epsilon_{-,i}=1$ and $\epsilon_{+,i}=(-1)^{i+1}$. Expanding the logarithms and separating $i=1$ and $j=1$ gives
\[
 \begin{aligned}
 F_-&=(1-w)^{-2}H_-,&\quad H_-&=\exp(-w+Q_-),\\
 F_+&=(1-w)^{-1}H_+,& H_+&=(1+w)\exp(-w+Q_+).
\end{aligned}
\]
Both $H_\pm$ have absolutely summable coefficients: $\sum_{j\ge2,i\ge2}j^{-i}<\infty$, and the coefficient $\ell^1$ norm is submultiplicative. At $w=1$, the unrestricted product for $H_-(1)$ is $e^{-1}\lim_{J\to\infty}\prod_{j=2}^J(1-1/j)^{-1}e^{-1/j}$; the distinct product is $2e^{-1}\lim\prod_{j=2}^J(1+1/j)e^{-1/j}$. Telescoping gives respectively $\lim J e^{-\sum_{j=1}^J1/j}$ and $\lim (J+1)e^{-\sum_{j=1}^J1/j}$, both $e^{-\gamma}$.


\subsection*{Individual coefficients and the logarithmic term}
Put $T(z)=\tan(rz)$. Its coefficients are nonnegative (by $T'=r(1+T^2)$) and sum to $1$. Hence $H_\pm(T(z))$ still have absolutely summable coefficients. For $d=1,2$ let $p_{d,n}=[z^n](1-T(z))^{-d}$. The nearest pole is $z=1$ and is isolated: $\tan(rz)=1$ precisely when $rz=\pi/4+\pi j$, $j\in\mathbb Z$, and $\tan(rz)$ is holomorphic for $|z|<2$. Since $T'(1)=2r$, subtraction of the principal part gives
\[
 p_{1,n}=(2r)^{-1}+o(1),\qquad
 p_{2,n}=(2r)^{-2}(n+1)+O(1).
\]
Convolution and dominated convergence now prove both leading equivalents. This does not assume that either original infinite product admits an isolated dominant singularity.

For the finer estimate let $G(z)=\exp(Q_-(T(z)))=\sum_{k\ge0}g_kz^k$ with $g_k\ge0$, $E(z)=\exp(-T(z))$, and write $H_-(T(z))=E(z)G(z)=\sum h_kz^k$. As $w\uparrow1$, the $i=2$ term of $Q_-'(w)$ is
$\sum_{j\ge2}w^{2j-1}/j=(-\log(1-w^2)-w^2)/w\sim\log(1/(1-w))$.
All $i\ge3$ derivative terms have bounded total, since $\sum_{j,i\ge2,\, i\ge3}j^{1-i}$ converges. Integration and $1-T(x)\sim2r(1-x)$ yield
\[
 G(1)-G(x)\sim 2rG(1)(1-x)\log(1/(1-x))
 \qquad(x\uparrow1).
\]
The nonnegative increasing sequence $S_n=\sum_{k\ge0}\min(n,k)g_k$ therefore satisfies
\[
 \sum_{n\ge0}S_nx^n=\frac{x(G(1)-G(x))}{(1-x)^2}
 \sim\frac{2rG(1)}{1-x}\log\frac1{1-x}.
\]
By the power-series Hardy--Littlewood--Karamata theorem, $\sum_{j\le n}S_j\sim2rG(1)n\log n$. Monotonicity gives $S_n\sim2rG(1)\log n$ by differencing at $n/2,n$ and at $n,2n$ and dividing by the interval lengths. The coefficients $e_j$ of $E$ satisfy $\sum j|e_j|<\infty$ by analyticity beyond the closed unit disk. Since $\min(n,k)$ is $1$-Lipschitz in $k$,
\[
 \left|\sum_{k\ge0}\min(n,k)h_k-E(1)S_n\right|
 \le G(1)\sum_j j|e_j|.
\]
Finally, convolving $p_{2,n}$ against $h_k$ gives
\[
 \frac{r^na_-(n)}{n!}
 =(2r)^{-2}\left((n+1)H_-(1)-\sum_{k\ge0}\min(n+1,k)h_k\right)+O(1).
\]
This proves the negative logarithmic correction. The same argument without composition ($T(z)=z$, $T'(1)=1$) gives the uncomposed unrestricted coefficient $e^{-\gamma}(n-\log n+o(\log n))$. Lehmer's uncomposed leading terms are prior work [2]. The conjectured composed leading terms can alternatively be transferred from them using positivity: the contribution of each fixed power $T^j$ tends to zero exponentially, while the sums against $1$ and $j$ have a simple and a double pole respectively.


\subsection*{The individual perimeter conjecture}
For a nonempty integer partition $\lambda=(\lambda_1,\ldots,\lambda_\ell)$ in decreasing order, let $|\lambda|=\sum\lambda_i$, $\operatorname{per}(\lambda)=\lambda_1+\ell-1$, $N(\lambda)=\prod\lambda_i$ and $\widehat N(\lambda)=\prod p_{\lambda_i}$, where $p_j$ is the $j$th prime. Write $W_{\rm size}(n),W_{\rm per}(n)$ for the sums of $1/N$ over the indicated exact ensembles; a hat denotes $1/\widehat N$ instead. For large $n$ conventions for the empty partition do not matter.

\textbf{Theorem 2 (Lagarias--Sun, Conjectures 3.6--3.8 [3]).} As $n\to\infty$,
\[
 \widehat W_{\rm per}(n)\sim e^\gamma/n,\qquad
 \widehat W_{\rm size}(n)\sim e^\gamma/n,\qquad
 W_{\rm size}(n)\widehat W_{\rm size}(n)\longrightarrow1 .
\]

\textit{Perimeter proof.} A partition of perimeter $n$ with largest part $k$ leaves, after removing one $k$, exactly $n-k$ parts of sizes at most $k$. Take independent random variables with $\PP(X_j=h)=(1-p_j^{-1})p_j^{-h}$, $h\ge0$, and put $S_k'=\sum_{j=1}^k X_j$. The geometric generating functions give the exact identity, for $n\ge2$,
\[
 \widehat W_{\rm per}(n)=\sum_{k=1}^n w_k\,\PP(S_k'=n-k),\qquad
 w_k=\frac1{p_k}\prod_{j=1}^k(1-p_j^{-1})^{-1}\sim\frac{e^\gamma}{k}.
\]
The last equivalence follows from Mertens' product theorem and $p_k\sim k\log k$, the prime number theorem. In particular the $w_k$ are bounded.

Take $t=3/2$, $D=\lceil10\log n\rceil$. Uniformly for $k\le n$, $\E t^{S_k'}=\prod_{j\le k}(1-1/p_j)/(1-t/p_j)\le O((\log(2n))^C)$ for some constant $C$: the log of each factor is $O(1/p_j)$ and $\sum_{j\le n}1/p_j=O(\log\log(3n))$. Consequently all $n-k>D$ together contribute $O((\log(2n))^C t^{-D})=o(n^{-1})$ by a geometric sum.

For $0\le n-k\le D$, couple $S_k'$ and $S_n'$ with the same $X_j$. Their probability of being unequal is at most $\sum_{k<j\le n}1/p_j=O(D/(n\log n))$. Summing over $D+1$ possible $k$ gives an $o(1)$ error. Hence
\[
 \sum_{n-D\le k\le n}\PP(S_k'=n-k)
 =\PP(S_n'\le D)+o(1)=1+o(1)
\]
by the same exponential moment bound. In this range $w_k\sim e^\gamma/n$ uniformly. This proves the perimeter assertion, without inferring local coefficients by differencing asymptotic equivalents.


\subsection*{Exact size and the ordinary perimeter question}
\textit{Size proof.} Let $s_n=\widehat W_{\rm size}(n)$, $s_0=1$, and let $A(N)=\sum_{j=0}^N s_j$. Lagarias and Sun [3, Thm.~3.4] establish $A(N)\sim e^\gamma\log N$. Logarithmic differentiation of the formal product $\sum s_n z^n=\prod_{j\ge1}(1-z^j/p_j)^{-1}$ gives
\[
 ns_n=\sum_{k=1}^n b_k s_{n-k},\qquad b_k=\sum_{d\mid k}\frac{d}{p_d^{\,k/d}}.
\]
The term $d=k$ is $k/p_k\sim1/\log k$. For $e=k/d$ between $2$ and $\sqrt k$, $d/p_d^e\le (1/\sqrt k)^{e-1}$ as $p_d\ge d$. The terms $e>\sqrt k$ total at most $k^2 2^{-\sqrt k}$. Thus $b_k\sim1/\log k$, so $b_k$ is bounded and the recurrence itself gives $s_n=O(\log(2n)/n)$.

For $k\le\sqrt n$ the convolution contributes $O(\log n/\sqrt n)$. For $\sqrt n<k<n/2$ it contributes at most
$O(1/\log n)\left(A(n)-A(\lfloor n/2\rfloor)\right)=o(1)$ by the cumulative equivalent. In the last range $n/2\le k\le n$ we have $b_k\sim1/\log n$ uniformly and the $s_{n-k}$ sum to $A(\lfloor n/2\rfloor)$. Consequently $n s_n\to e^\gamma$. Lehmer's theorem $W_{\rm size}(n)\sim e^{-\gamma}n$ gives the product assertion, as the source itself observes. $\square$

\textbf{Ordinary perimeter.} An additional question in the appendix of [3] asks for the leading growth of $W_{\rm per}(n)$. It is $W_{\rm per}(n)\sim n$. In the same largest-part decomposition the ordinary part $1$ has weight $1$ and can fill the unused length. For independent geometrical variables $Y_j$ with $\PP(Y_j=h)=(1-1/j)j^{-h}$ for $j\ge2$, this gives
\[
 W_{\rm per}(n)=\sum_{k=1}^n\PP\left(\sum_{j=2}^{k}Y_j\le n-k\right)\quad(n\ge2),
\]
because the normalization is $k^{-1}\prod_{j=2}^k(1-1/j)^{-1}=1$. The mean of the inner sum is $\sum_{j=2}^k1/(j-1)=O(\log(2n))$. For $k\le n-\sqrt n$ Markov's inequality makes the displayed probabilities $1-o(1)$ uniformly. The other $O(\sqrt n)$ terms are bounded by one, proving the result.

\textbf{Scope.} No arithmetic property of $\gamma$ is claimed. The conjectures are settled here by unconditional arguments; focused searches do not certify worldwide priority.

{\small\raggedright
[1] V. Kote\v{s}ovec, Oct.\ 4, 2020 conjectures, in S. Manyama's
\href{https://oeis.org/A335636}{OEIS A335636} and \href{https://oeis.org/A335638}{A335638}; \href{https://oeis.org/A080130}{A080130} is $e^{-\gamma}$.

[2] D.~H. Lehmer, \emph{On reciprocally weighted partitions}, Acta Arith. \textbf{21} (1972), 379--388. For the uncomposed distinct case see also \href{https://arxiv.org/abs/math/0606370}{Flajolet, Fusy, Gourdon, Panario and Pouyanne}.

[3] J.~C. Lagarias and C. Sun, \href{https://doi.org/10.1007/s11139-024-00893-8}{\emph{Asymptotics of reciprocal supernorm partition statistics}}, Ramanujan J. \textbf{65} (2024), 487--520. See Section 3 and Appendix A.\par}

\par\endgroup

% Source: /workspace/scratch/81914cec2c78/work/archive/Tournament_Round_40_Final/Tournament_Round_40/round40/showa/paper.tex
% Edited source: /workspace/scratch/81914cec2c78/work/audit_34_41/cleaned/r40_showa.tex
\begingroup

\Needspace{9\baselineskip}
\section{\texorpdfstring{The sharp common exponent in two of Liu's ternary quadratic inequalities}{The sharp common exponent in two of Liu's ternary quadratic inequalities}}\label{section:r40showa}
\noindent{\small\sffamily\color{muted}Round 40 \enspace / \enspace Showa}\par\medskip


\begin{sectionabstract}
The best exponent in Problems 1 and 3 of Jian Liu's 2008 paper is
$\log 2/\log(2/\sqrt3)$. The two apparently different quadratic forms become identical after diagonal scaling. At the limiting exponent their positivity reduces to a sum of three single-angle functions. Concavity proves the claim near the equilateral triangle, a quadratic deficit controls small angles, and an explicitly reproducible finite rational certificate treats the remaining compact set.
\end{sectionabstract}

\subsection{The two problems and their common upper bound}
Let $A,B,C$ be the angles of an arbitrary nondegenerate triangle, $a,b,c$ its opposite side lengths, and $h_a,h_b,h_c$ its altitudes. Liu [1, Sec.~4, Problems 1 and 3, (4.5) and (4.9)] asks for the largest real $k$ such that, for every such triangle and all real $x,y,z$, one has, respectively,
\begin{align}
 \sum_{\rm cyc}\frac{x^2}{\sin^k A}&\ge
 (2/\sqrt3)^k(yz+zx+xy),\label{r40showa:e1}\\
 \sum_{\rm cyc}\frac{x^2}{h_a^k}&\ge
 (2/\sqrt3)^k\sum_{\rm cyc}\frac{yz}{a^k}.\label{r40showa:e2}
\end{align}
He gives $k=4$ and a numerical indication of approximately $4.82$ for the maximum in Problem 1. We prove that the maximum in both problems is
\[
 K=\frac{\log 2}{\log(2/\sqrt3)}=4.8188416793\ldots .
\]

In \eqref{r40showa:e1} put $x=(\sin A)^{k/2}X$ cyclically, and in
\eqref{r40showa:e2} put $x=h_a^{k/2}X$ cyclically. Both acquire diagonal terms $X^2+Y^2+Z^2$, and corresponding $YZ$ terms agree because
\[
 \frac{\sqrt{h_bh_c}}{a}
 =\frac{2S}{a\sqrt{bc}}=\frac{\sqrt{bc}}{2R}
 =\sqrt{\sin B\sin C},
\]
using the area $S$, circumradius $R$, and sine rule. Thus there is equivalence for every $k$ without any sign restriction on $x,y,z$. In \eqref{r40showa:e1} take $x=0$, $y=z=1$, and let $B=C$ tend to $\pi/2$ from below. The bound $2\ge(2/\sqrt3)^k$ is necessary. It remains to prove \eqref{r40showa:e1} at $k=K$.

\subsection{The critical scalar inequality}
At $k=K$ rewrite \eqref{r40showa:e1} as $\sum_{\rm cyc}(1+\sin^{-K}A)x^2\ge (x+y+z)^2$. By weighted Cauchy--Schwarz, with necessity obtained by choosing $x=(1+\sin^{-K}A)^{-1}$ cyclically, this is equivalent to
\[
 \sum_{\rm cyc} f(A)\le 1,\qquad f(t)=\frac{\sin^Kt}{1+\sin^Kt}.
\]
Assume $A\le B\le C$. We use $53/11<K<5$: the two inequalities are respectively the integer comparisons $2^{84}<3^{53}$ and $2^8>3^5$.

\paragraph{The small-angle region: $A\le\pi/16$.}
Here $A<1/4$, so $f(A)\le A^4$. For any angle $t$ one has
\[
 \tfrac12-f(t)=\frac{1-\sin^Kt}{2(1+\sin^Kt)}
 \ge \tfrac14\cos^2t.
\]
Also $\cos^2B+\cos^2C=1-\cos A\cos(B-C)\ge1-\cos A\ge A^2/3$; the last bound follows from the cosine Taylor bound for $A\le1/4$. Thus the two deficits are at least $A^2/12>A^4$ in sum.

\paragraph{The near-equilateral region: $A\ge3\pi/10$.}
Every angle lies in $[3\pi/10,2\pi/5]$. Throughout this interval $\sin t\ge (1+\sqrt5)/4>4/5$. The sign of $f''(t)$ is that of
\[
 K\cos^2t(1-u)-(1+u),\qquad u=\sin^K t .
\]
Indeed this is its factor after removing the positive factor $K\sin^{K-2}t/(1+u)^3$. Since $u\ge(4/5)^5$, $\cos^2t\le9/25$, and
$5(9/25)(1-(4/5)^5)<1+(4/5)^5$, the function is strictly concave. Hence $\sum f(A)\le3f(\pi/3)=1$.

\paragraph{The remaining region: a rational certificate.}
Put $\alpha=A/\pi$, $\beta=B/\pi$. It remains to check
\[
 (\alpha,\beta)\in[1/16,3/10]\times[1/16,1/2],\qquad
 \alpha\le\beta,\quad \alpha+2\beta\le1.
\]
The third sine is $\sin(\pi(\alpha+\beta))$. For a rational interval $I=[l,h]\subset[0,1]$, put $t=1/2$ if $1/2\in I$ and $t=\min(h,1-l)$ otherwise. The sine on $I$ is bounded above by $\sin(\pi t)$ by monotonicity and reflection. Write $D=10^9$, $v=355t/113$, and compute
\[
 P_9(v)=v-v^3/6+v^5/120-v^7/5040+v^9/362880,\quad
 n=\min\{D,\lceil D(P_9(v)+10^{-6})\rceil\}.
\]
The sine Taylor series alternates with decreasing terms at $v\le 355/226$, and $|\pi-355/113|<10^{-6}$, so $n/D$ is a rigorous upper bound. Let $h_0$ be the least nonnegative integer with
$h_0^{11}D^{42}\ge n^{53}$. Since $K>53/11$ and $n/D\in[0,1]$, an upper bound for $f$ on this interval of normalized angle is $h_0/(D+h_0)$.

The algorithm below checks the sum of the three upper bounds on each rectangle, discards only rectangles wholly outside the two order constraints, and otherwise quarters the rectangle. It completes at depth eight: 275 accepted leaves and 47 exterior leaves, so the closed rectangles cover the entire domain, including boundaries. All comparisons in this certificate are exact rational comparisons; no floating point is used to accept any result. This proves the endpoint and both maxima.

\paragraph{Reference.} [1] J. Liu, \emph{The Composite Theorem of Ternary Quadratic Inequalities and its Applications}, RGMIA Research Report Collection 11 (2008), no. 4, art. 13, Sec.~4, pp.~9--10.


\subsection*{Certificate listing (Python 3, standard library)}
\begingroup\scriptsize\begin{Verbatim}
from fractions import Fraction as Q
from math import factorial
from functools import lru_cache
from collections import Counter
B=10**9
assert 2**84<3**53 and 2**8>3**5 # 53/11 < k < 5
def ceildiv(n,d):return -(-n//d)
@lru_cache(None)
def value(t):
    # 0<=t<=1/2, majorize sin(pi*t) using monotonicity/Lipschitz,
    # or alternating Taylor upper at pi*t followed by pi upper (see paper).
    assert 0<=t<=Q(1,2)
    x=t*Q(355,113)
    # use pi approximation error plus sin(x) <= its Taylor upper.
    s=sum((-1)**j*x**(2*j+1)/factorial(2*j+1) for j in range(5))+Q(1,1000000)
    n=min(B,ceildiv(B*s.numerator,s.denominator))
    # least integral h for (h/B)^11 >= (n/B)^53
    lo,hi=0,B
    while hi-lo>1:
        m=(lo+hi)//2
        if m**11*B**42>=n**53:hi=m
        else:lo=m
    return Q(hi,B+hi)
def maj(lo,hi):
    t=Q(1,2) if lo<=Q(1,2)<=hi else min(hi,1-lo)
    return value(t)
counts=Counter()
def test(a,A,b,B0,d):
    if a>B0 or a+2*b>1:
        counts['outside']+=1;return
    upper=maj(a,A)+maj(b,B0)+maj(a+b,A+B0)
    if upper<=1:
        counts['pass']+=1
        counts['max-depth']=max(counts['max-depth'],d)
        return
    if d>=24:raise AssertionError((a,A,b,B0,float(upper)))
    u=(a+A)/2;v=(b+B0)/2
    for e,E in [(a,u),(u,A)]:
        for f,F in [(b,v),(v,B0)]:test(e,E,f,F,d+1)
test(Q(1,16),Q(3,10),Q(1,16),Q(1,2),0)
print('certified cells, discarded cells, maximum depth',counts['pass'],counts['outside'],counts['max-depth'])
print('cached trigonometric bounds',value.cache_info().currsize)
\end{Verbatim}
\endgroup

\par\endgroup

% Source: /workspace/scratch/81914cec2c78/work/archive/Tournament_Round_40_Final/Tournament_Round_40/round40/heisei/paper.tex
% Edited source: /workspace/scratch/81914cec2c78/work/audit_34_41/cleaned/r40_heisei.tex
\begingroup
\newcommand{\Lap}{\mathcal L}
\Needspace{9\baselineskip}
\section{\texorpdfstring{A positive Laplace density for the Guo--Zhang--Qi power of the digamma function}{A positive Laplace density for the Guo--Zhang--Qi power of the digamma function}}\label{section:r40heisei}
\noindent{\small\sffamily\color{muted}Round 40 \enspace / \enspace Heisei}\par\medskip


\begin{sectionabstract}
For $q(x)=x^{x[\psi(x)-\log x]-\gamma}$, we construct a strictly positive inverse Laplace density for $-(\log q)'$. This proves the global logarithmic complete monotonicity conjectured in 2008, beyond the previously justified unit interval. The key step bounds a logarithmic convolution by the monotonicity of the derivative of the digamma kernel. A separate appendix records a known-theorem implication for a Fuss--Dyck path polynomial.
\end{sectionabstract}
\subsection{The whole positive half-line}
Here $\psi=\Gamma'/\Gamma$ and $\gamma$ is Euler's constant. Guo, Zhang and Qi \cite[Remark 8, eq.~(21)]{r40heisei:gzq} asked whether $q$ is logarithmically completely monotone on $(0,\infty)$. Krasniqi and Shabani \cite{r40heisei:ks} claimed this, but Guo and Qi \cite[Theorem 1.1 and preceding discussion]{r40heisei:gq} identified the faulty global assertion that $-\log x$ is completely monotone, proved the result for $0<x<1$ and proved that $q$ is not a Stieltjes function. Our proof does not require a sign for $-\log x$.
\begin{theorem}
For every integer $k\ge1$ and every real $x>0$,
\[
(-1)^k\frac{d^k}{dx^k}\log q(x)>0.
\]
In particular $q$ is strictly decreasing from infinity to zero, with $q(1)=1$.
\end{theorem}
\begin{proof}
Write $\Lap v(x)=\int_0^\infty e^{-xt}v(t)\,dt$. The classical digamma representation gives $\log x-\psi(x)=\Lap h(x)$ for
\[
 h(t)=\frac{1}{1-e^{-t}}-\frac1t,\qquad
 d(t)=h'(t)=\frac1{t^2}-\frac1{4\sinh^2(t/2)}.
\]
For example, this follows by subtracting the integral for $\psi(x)$ from Frullani's integral for $\log x$, with common cutoffs. At zero $h(0)=1/2$ and $d(0)=1/12$ by Taylor expansion. At infinity $h(t)\to1$ and $d(t)=O(t^{-2})$. Put $g(t)=td(t)$, $c=\tfrac12+\gamma$, and $f(x)=x(\log x-\psi(x))$. Integration by parts gives $f=1/2+\Lap d$, whence $f'=-\Lap g$. Thus
\[
 H(x):=-(\log q)'(x)=\frac{\gamma+f(x)}x+f'(x)\log x.
\]
The transform of $\log t+\gamma$ is $-(\log x)/x$ by differentiating Euler's gamma integral. Since $g(0)=0$, $\Lap g'=x\Lap g$. The convolution theorem gives $H=\Lap K$, where
\begin{align}
 K(t)&=c+\int_0^t d(s)\,ds+\int_0^t g'(s)\,[\log(t-s)+\gamma] \,ds\notag\\
 &=c+\int_0^t d(s)\,ds+g(t)(\log t+\gamma)
       +\int_0^t\frac{g(s)-g(t)}{t-s}\,ds.\label{r40heisei:eq:key}
\end{align}
For the second equality, integrate by parts after subtracting $g(t)$ from $g(s)$; the upper boundary is zero since $(g(s)-g(t))\log(t-s)\to0$. All transforms and the convolution are absolutely integrable with each weight $e^{-xt}$, $x>0$: $g'$ is bounded near zero and $O(t^{-2})$ at infinity, while $\log t+\gamma$ is locally integrable and grows logarithmically.

We claim $0<d(t)\le1/12$ and $d$ decreases. The first strict sign follows from $\sinh u>u$ for $u=t/2>0$. Direct differentiation shows $d'(t)<0$ is equivalent to $\sinh^3 u>u^3\cosh u$. To verify this, the derivative of $3\log(\sinh u/u)-\log\cosh u$ has sign equal to that of
\[
 A(u)=u(3+2\sinh^2 u)-3\sinh u\cosh u.
\]
Here $A(0)=0$ and $A'(u)=4\sinh u(u\cosh u-\sinh u)>0$, since the last parenthesis vanishes at zero and has derivative $u\sinh u>0$. The desired log difference starts at zero as well.

For $0<s<t$, monotonicity of $d$ gives $g(s)-g(t)\ge-(t-s)d(t)$, and $\int_0^t d(s)\,ds\ge td(t)$. In~\eqref{r40heisei:eq:key} these two lower bounds cancel, leaving
\[
 K(t)\ge c+t d(t)(\log t+\gamma).
\]
For $t\ge1$ the right side is at least $c$; for $0<t<1$ it is at least $c-1/(12e)>0$ by $t\log t\ge-1/e$. Thus the inverse density is strictly positive everywhere. Absolute exponential integrability at $x/2$ justifies differentiating every order:
\[
 (-1)^k(\log q)^{(k)}(x)
 =\int_0^\infty t^{k-1}e^{-xt}K(t)\,dt>0.
\]
Finally $d$ is integrable and $\int_0^\infty d=h(\infty)-h(0)=1/2$. Consequently $f(0+)=1$, $f(\infty)=1/2$ by dominated convergence. Since $\log q=-(\gamma+f)\log x$, the stated limits and normalization follow.
\end{proof}


\subsection{A known-theorem implication for the Fuss--Dyck question}
This appendix is independent of the preceding theorem. Alexandersson's list \cite{r40heisei:alex} asks for real-rootedness of the polynomial $W^{(m)}_{n,k}(z)$ counting $m$-Dyck paths with $k$ occurrences of $NE$ by occurrences of $NNE$. Take $m,n\ge1$ and $1\le k\le n$; the paths have $n$ north steps and $mn$ east steps with $m\#N-\#E\ge0$ in every prefix. An $NE$ ends a north run, and an $NNE$ ends one of length at least two.

Append one east step and read the resulting word cyclically, with increments $+m$ and $-1$. Every cyclic word with these step counts has a unique cut with all proper partial sums nonnegative. In fact cut just after the first occurrence of the minimum among the partial sums at times $0,\ldots,L-1$, where $L=(m+1)n+1$. The earlier partial sums are strictly larger than the minimum, hence are still at least the minimum after adding the full sum $-1$ on wrapping. Conversely these weak and strict comparisons force exactly that first minimum. The last step of the cut word is east; its removal recovers the path. Total sum $-1$ also precludes a nontrivial period.

Mark one of the $k$ north runs. Reading from the marked run gives independently a composition of $n$ and a composition of $mn+1$, both into $k$ positive parts, alternating cyclically. Every path yields $k$ distinct marked words because there is no nontrivial period. To have $\ell$ long north runs, choose which $\ell$ parts exceed one, then distribute $n-k$ extra units positively amongst them. Thus for $k<n$,
\[
 W^{(m)}_{n,k}(z)=\frac1k\binom{mn}{k-1}
 \sum_{\ell\ge1}\binom{k}{\ell}\binom{n-k-1}{\ell-1}z^{\ell}.
\]
For $k=n$ it is the constant $\binom{mn}{n-1}/n$. In each case the polynomial is a positive scalar multiple of its $m=1$ counterpart. Real-rootedness follows from the ordinary Dyck result \cite{r40heisei:wz}.

\paragraph{Prior scope.} This count is already implied by the rational Kreweras enumeration of vertical run multisets \cite{r40heisei:br}: appending one east step identifies the given paths with rational $(n,mn+1)$-Dyck paths. Indeed at a height $y<n$ the rational bound on the abscissa is $\lfloor(mn+1)y/n\rfloor=my$. Group the known vertical-run count by number of runs of length at least two. Thus neither the enumeration nor its combination with the ordinary case should be advertised as an independent primary discovery.

\begin{thebibliography}{9}\small
\bibitem{r40heisei:gzq} B.-N. Guo, Y.-J. Zhang and F. Qi, Refinements and sharpenings of some double inequalities for bounding the gamma function, \emph{J. Inequal. Pure Appl. Math.} 9 (2008), no. 1, Art. 17, Remark 8.
\bibitem{r40heisei:ks} V. Krasniqi and A. Sh. Shabani, On a conjecture of a logarithmically completely monotonic function, \emph{Aust. J. Math. Anal. Appl.} 11 (2014), no. 1, Art. 5.
\bibitem{r40heisei:gq} B.-N. Guo and F. Qi, Logarithmically complete monotonicity of a power-exponential function involving the logarithmic and psi functions, \emph{Global J. Math. Anal.} 3 (2015), 77--80. \href{https://doi.org/10.14419/gjma.v3i2.4605}{doi:10.14419/gjma.v3i2.4605}.
\bibitem{r40heisei:alex} P. Alexandersson, \href{https://www.symmetricfunctions.com/realRootedOpenProblems.htm}{Real-rooted polynomials: open problems}, SymCat, entry on generalized Dyck paths.
\bibitem{r40heisei:br} M. Bodnar and B. Rhoades, Cyclic sieving and rational Catalan theory, \emph{Electron. J. Combin.} 23 (2016), no. 2, P2.4, eq.~(2).
\bibitem{r40heisei:wz} B. Wang and C. X. Zhang, Interlacing property of a family of generating polynomials over Dyck paths, \href{https://arxiv.org/abs/2309.05903}{arXiv:2309.05903} (2023).
\end{thebibliography}

\par\endgroup

% Source: /workspace/scratch/81914cec2c78/work/archive/Tournament_Round_40_Final/Tournament_Round_40/round40/keio/paper.tex
% Edited source: /workspace/scratch/81914cec2c78/work/audit_34_41/cleaned/r40_keio.tex
\begingroup
\newcommand{\cyc}{\sum_{\rm cyc}}
\Needspace{9\baselineskip}
\section{\texorpdfstring{A change of direction for a reciprocal exradius--median sum}{A change of direction for a reciprocal exradius--median sum}}\label{section:r40keio}
\noindent{\small\sffamily\color{muted}Round 40 \enspace / \enspace Keio}\par\medskip


\begin{sectionabstract}
We establish both directions of Jian Liu's Conjecture 6.4 on its stated parameter ranges for arbitrary triangles. Convexity reduces the forward range to two endpoint inequalities and propagates the reversed inequality from a third endpoint. Exact two-dimensional certificates verify the three shape inequalities; a degenerate limit and the equilateral equality point are treated separately.
\end{sectionabstract}

\begin{theorem}[Liu 2023, Conjecture 6.4]
In every nondegenerate triangle with sides $a,b,c$, medians $m_a,m_b,m_c$ and exradii $r_a,r_b,r_c$, put $Q=ab+bc+ca$. Then
\[
          (1+k)\cyc\frac{1}{r_a+km_a}\le\frac{6}{\sqrt Q}
          \qquad(1\le k\le 6.12).
\]
The inequality reverses for $0<k\le0.37$. In either range equality holds precisely for equilateral triangles.
\end{theorem}
This is equation (6.5) of [1]. The preceding equation gives the $k=2$ bound for acute triangles. The full conjecture, unlike many inequalities in that article, allows obtuse triangles in both ranges.

\paragraph{Interpolating and extrapolating.}
Set $T=1/(1+k)$ and define a convex function on $[0,1]$ by
\[
       H(T)=\sqrt Q\cyc\frac{1}{(1-T)m_a+Tr_a}.
\]
Indeed, each reciprocal has a positive affine denominator. We will certify $H\le6$ at the two points for $k=1$ and $6.12$ and $H\ge6$ at the point for $k=0.37$. The forward assertion then follows by convex interpolation. For the reverse assertion take $T_0=1/2$ and $T_1=1/1.37>T_0$. Since $H(T_1)\ge6\ge H(T_0)$, the slope of this secant is nonnegative. Convexity implies $H(T)\ge H(T_1)$ for $T_1<T<1$. Strict endpoint inequalities for a nonequilateral triangle give strictness throughout; an equilateral triangle gives equality by substitution.


\subsection*{Three exact endpoint certificates}
Write $s=(a+b+c)/2$. Permute and rescale the tangent lengths $(s-a,s-b,s-c)$ to $(u^2t,t,1)$ with $0<u,t\le1$. Then
\[
   a=1+t,\qquad b=1+u^2t,\qquad c=(1+u^2)t,\qquad g=\sqrt{s}=\sqrt{1+c}.
\]
Heron's formula gives the exradii $(g/u,ug,utg)$. For tangent lengths $t_i$ the usual median identity takes the form $m_i^2=st_i+(t_j-t_l)^2/4$. Thus in the indicated order
\[
 m_0^2=su^2t+\tfrac14(1-t)^2,\quad
 m_1^2=st+\tfrac14(1-u^2t)^2,\quad m_2^2=s+\tfrac14t^2(1-u^2)^2.
\]
With $k=K/100$ set
\[
 D_K=6-(100+K)\sqrt Q\left(
 \frac{u}{100g+Kum_0}+\frac1{100ug+Km_1}
                  +\frac1{100utg+Km_2}\right).
\]
The first term clears the singularity of $g/u$ on the compactification. Put $F_K=D_K$ for $K=100,612$ and $F_{37}=-D_{37}$. We certify the positivity of these three functions on the square $x=1-u$, $z=1-t$, except at specified boundary zeros.

\paragraph{Equilateral corner.}
At $x=z=0$, $F=0$ and $\nabla F=0$. The gradient vanishes because in the original side variables the function is smooth, symmetric and homogeneous of degree zero: at the diagonal the three derivatives coincide and their sum vanishes by common scaling. The chart preserves this zero. On $[0,1/e]^2$ the verifier proves the following lower bounds (entries are multiplied by $10^6$ and rounded down):
\[
\begin{array}{r|r|rrr}
K&e&F_{xx}&F_{xz}&F_{zz}\\\hline
100&512&894191&36492&155852\\
612&8192&102773&11987&23761\\
37&512&819451&24052&112477
\end{array}
\]
As both coordinates are nonnegative, these entrywise signs give strictly positive second derivative on every nonconstant ray from $(0,0)$, so integration twice proves strict positivity away from that origin.

\paragraph{Degenerate edge for $K=100$.}
On $u=0$ we have
\[
 6-D_{100}=\frac{12(1+t)\sqrt{1+3t+t^2}}{(1+2t)(2+t)}.
\]
After squaring the claimed comparison with $6$ and multiplying by the squared denominator, the difference is exactly $36t^2$. Thus $D_{100}(1,z)\ge0$ in $(x,z)$ coordinates. If an interval proof certifies $F_x<0$ on the rectangle extending from the left edge of the tested box all the way to $x=1$, then every actual triangle in that box has $F(x,z)>F(1,z)\ge0$. The test uses this entire rectangle, not just the box itself.


\paragraph{Exhaustion of the square.}
On the remaining boxes the verifier certifies either $F(B)>0$ by direct interval evaluation, or a positive lower bound for
\[
               F(c_B)+\nabla F(B)\mathbin{\cdot}(B-c_B)
\]
with $c_B$ the midpoint. This enclosure follows by integrating the gradient on each segment from $c_B$. A rejected or inconclusive box is bisected, with both closed halves retained. All three resulting covers terminate:
\[
\begin{array}{r|rrrr|r}
 K&\text{direct}&\text{centered}&\text{local}&\text{edge}&\text{maximum depth}\\\hline
100&4&1170&1&3&26\\
612&10&18094&1&0&37\\
37&10&1395&1&0&26
\end{array}
\]
The rest of the paper includes all the code needed to regenerate these tables. Intervals use a common denominator $2^{64}$ with outward rounding. Square roots use integer square roots; a known nonnegative radicand may be intersected with $[0,\infty)$. All radicals here are either positive by definition or are lengths computed as sums of nonnegative terms. A derivative requiring division by a box containing zero is never accepted. The first and second derivatives are enclosed using the product and chain rules. Positivity is certified for boxes, not sampled points. The sorted tangent-length coordinates and generic interval primitives were adapted from earlier tournament checks on different inequalities (Meiji, Rounds 36--37; Keio, Round 38). No claim of optimal parameter endpoints is made.\hfill$\square$

\begin{thebibliography}{1}
\bibitem{r40keio:liu} J. Liu, A new geometric inequality in acute triangles and its applications, \emph{International Journal of Geometry} \textbf{12} (2023), no.~4, 47--63. \url{https://ijgeometry.com/wp-content/uploads/2023/09/5.-47-63.pdf}.
\end{thebibliography}
\subsection*{Verifier: Intervals and gradients}
\scriptsize\begin{Verbatim}
from math import isqrt
from collections import Counter
S=1<<64
class I:
 __slots__='l','h'
 def __init__(self,l,h=None):
  if h is None:self.l=self.h=l*S
  else:self.l=l;self.h=h
 def __add__(x,y):
  if not isinstance(y,I):y=I(y)
  return I(x.l+y.l,x.h+y.h)
 __radd__=__add__
 def __neg__(x):return I(-x.h,-x.l)
 def __sub__(x,y):return x+-y
 def __rsub__(x,y):return -x+y
 def __mul__(x,y):
  if not isinstance(y,I):return I(min(x.l*y,x.h*y),max(x.l*y,x.h*y))
  if x.l>=0 and y.l>=0:return I(x.l*y.l//S,-(-x.h*y.h//S))
  p=(x.l*y.l,x.l*y.h,x.h*y.l,x.h*y.h);return I(min(p)//S,-(-max(p)//S))
 __rmul__=__mul__
 def __truediv__(x,y):
  if not isinstance(y,I):
   assert y>0;return I(x.l//y,-(-x.h//y))
  if y.l<=0:raise ValueError
  return x*I(S*S//y.h,-(-S*S//y.l))
 def rt(x):
  if x.h<0:raise ValueError
  return I(isqrt(max(0,x.l)*S),isqrt(x.h*S)+1)
class J:
 __slots__='v','g'
 def __init__(self,v,g):
  self.v=v if isinstance(v,I) else I(v);self.g=[d if isinstance(d,I) else I(d) for d in g]
 def __add__(x,y):
  if not isinstance(y,J):return J(x.v+y,x.g)
  return J(x.v+y.v,[p+q for p,q in zip(x.g,y.g)])
 __radd__=__add__
 def __neg__(x):return J(-x.v,[-p for p in x.g])
 def __sub__(x,y):return x+-y
 def __rsub__(x,y):return -x+y
 def __mul__(x,y):
  if not isinstance(y,J):return J(x.v*y,[p*y for p in x.g])
  return J(x.v*y.v,[p*y.v+q*x.v for p,q in zip(x.g,y.g)])
 __rmul__=__mul__
 def __truediv__(x,y):
  if not isinstance(y,J):return J(x.v/y,[p/y for p in x.g])
  v=x.v/y.v;return J(v,[(p-v*q)/y.v for p,q in zip(x.g,y.g)])
 def rt(x):
  y=x.v.rt();return J(y,[p/(2*y) for p in x.g])
\end{Verbatim}

\subsection*{Verifier: Mixed derivatives}
\scriptsize\begin{Verbatim}
class B:
 def __init__(self,v,g=0,h=0,j=0):
  self.v,self.g,self.h,self.j=[i if isinstance(i,I) else I(i) for i in (v,g,h,j)]
 def __add__(x,y):
  if not isinstance(y,B):y=B(y)
  return B(x.v+y.v,x.g+y.g,x.h+y.h,x.j+y.j)
 __radd__=__add__
 def __neg__(x):return B(-x.v,-x.g,-x.h,-x.j)
 def __sub__(x,y):return x+-y
 def __rsub__(x,y):return -x+y
 def __mul__(x,y):
  if not isinstance(y,B):y=B(y)
  return B(x.v*y.v,x.g*y.v+x.v*y.g,x.h*y.v+x.v*y.h,x.j*y.v+x.g*y.h+x.h*y.g+x.v*y.j)
 __rmul__=__mul__
 def chain(x,y,a,b):return B(y,a*x.g,a*x.h,a*x.j+b*x.g*x.h)
 def __truediv__(x,y):
  if not isinstance(y,B):return B(x.v/y,x.g/y,x.h/y,x.j/y)
  r=I(1)/y.v;return x*y.chain(r,-r*r,2*r*r*r)
 def rt(x):
  r=x.v.rt();d=I(1)/(2*r);return x.chain(r,d,-d/(2*x.v))



\end{Verbatim}

\subsection*{Verifier: Endpoint and local test}
\scriptsize\begin{Verbatim}
def f(x,z,k):
 u,t=1-x,1-z
 a=1+t;b=1+u*u*t;c=t*(1+u*u)
 ss=1+c;g=ss.rt()
 d0=(1-t)/2;d1=(1-u*u*t)/2;d2=t*(1-u*u)/2
 m0=(ss*u*u*t+d0*d0).rt()
 m1=(ss*t+d1*d1).rt()
 m2=(ss+d2*d2).rt()
 terms=u/(100*g+k*u*m0)+(m1*0+1)/(100*u*g+k*m1)+(m2*0+1)/(100*u*t*g+k*m2)
 F=6-(a*b+b*c+c*a).rt()*(100+k)*terms
 return -F if k==37 else F
def local(k,eps):
 z=I(0,S//eps)
 bounds=[]
 for p,q in [(0,0),(0,1),(1,1)]:
  H=f(B(z,p==0,q==0),B(z,p==1,q==1),k).j
  bounds.append(H.l*10**6//S)
 if min(bounds)>0:
  print('positive Hessian bounds /1e6:',k,eps,bounds,flush=True)
  return True
 return False
\end{Verbatim}

\subsection*{Verifier: Adaptive cover}
\scriptsize\begin{Verbatim}
def prove(k):
 eps=8
 while True:
  try:
   if local(k,eps):break
  except ValueError:pass
  eps*=2;assert eps<2**20
 ct=Counter();stack=[(0,S,0,S,0)]
 while stack:
  l,h,u,v,d=stack.pop();ct['depth']=max(ct['depth'],d)
  if h<=S//eps and v<=S//eps:ct['Taylor']+=1;continue
  if k==100 and l>=3*S//4 and u>=3*S//4:
   try:
    edge=f(J(I(l,S),[1]),J(I(u,v),[0]),k)
    if edge.g[0].h<0:ct['corner']+=1;continue
   except ValueError:pass
  xx=I(l,h);yy=I(u,v)
  try:
   if f(xx,yy,k).l>0:ct['direct']+=1;continue
   m=(l+h)//2;n=(u+v)//2
   val=f(I(m,m),I(n,n),k)
   jet=f(J(xx,[1,0]),J(yy,[0,1]),k)
   if (val+jet.g[0]*I(l-m,h-m)+jet.g[1]*I(u-n,v-n)).l>0:ct['mean']+=1;continue
  except ValueError:pass
  assert d<50,(k,d,l/S,h/S,u/S,v/S,ct)
  if h-l>=v-u:
   m=(l+h)//2;stack.extend([(l,m,u,v,d+1),(m,h,u,v,d+1)])
  else:
   n=(u+v)//2;stack.extend([(l,h,u,n,d+1),(l,h,n,v,d+1)])
 print('endpoint',k,'/100',dict(ct),flush=True)
for k in (100,612,37):prove(k)
\end{Verbatim}


\par\endgroup
\clearpage
\part{Rounds 41--45}

% Source: /workspace/scratch/81914cec2c78/work/archive/Tournament_Round_41_Final/Tournament_Round_41/meiji/paper.tex
% Edited source: /workspace/scratch/81914cec2c78/work/audit_34_41/cleaned/r41_meiji.tex
\begingroup

\Needspace{9\baselineskip}
\section{\texorpdfstring{A coefficient certificate for Liu's weighted side inequality}{A coefficient certificate for Liu's weighted side inequality}}\label{section:r41meiji}
\noindent{\small\sffamily\color{muted}Round 41 \enspace / \enspace Meiji}\par\medskip

\begin{sectionabstract}
We prove Conjecture 4.1 of Jian Liu (2010) for every nondegenerate Euclidean triangle and every triple of positive weights. Six simultaneous order charts turn the cleared inequality into polynomials with positive nonzero coefficients. In fact the published inequality is strict on its domain.
\end{sectionabstract}
\paragraph{Theorem.}
Let $a,b,c$ be the sides of a nondegenerate triangle, $s=(a+b+c)/2$, and $x,y,z>0$. Then
\[
\frac{(s-b)(s-c)}{y+z}+\frac{(s-c)(s-a)}{z+x}
+\frac{(s-a)(s-b)}{x+y}
<\frac{3(xa^2+yb^2+zc^2)}{4(xy+yz+zx)}.
\]
The weak form is precisely \cite[Conj.~4.1, (4.30)]{r41meiji:liu}. We address the weighted algebraic conjecture, not the separate angle-bisector Conjecture~4.2.

\paragraph{Proof.} Set $(u,v,w)=(s-a,s-b,s-c)$, so that all three are positive and $(a,b,c)=(v+w,w+u,u+v)$. Write $H=xy+yz+zx$, $D=(x+y)(y+z)(z+x)$. Multiplication of the proposed right minus left by $4HD>0$ gives exactly
\begin{align*}
 F={}&3D\{x(v+w)^2+y(w+u)^2+z(u+v)^2\}\\
  &-4H\{vw(z+x)(x+y)+wu(x+y)(y+z)+uv(y+z)(z+x)\}.
\end{align*}
Simultaneous permutation of the triples of variables preserves $F$. Thus we may put
\[
 (x,y,z)=(p+q+r,p+q,p),\qquad p>0,\quad q,r\ge0 .
\]
The order of the three side parameters remains free: write $(u,v,w)=(T_i,T_j,T_k)$, where $(T_0,T_1,T_2)=(t+m+n,t+m,t)$, $t>0$, $m,n\ge0$, and $(i,j,k)$ ranges over all permutations of $(0,1,2)$. Ties may occur in multiple charts; no case is lost.

The integer coefficients obtained by distributivity have the following certificate. Here $N$ is the number of nonzero coefficients, $L$ is their minimum, $S$ is their sum and $B=[p^4t^2]F$ (the remaining exponents are zero).
\begin{center}
\begin{tabular}{c|rrrr}
$(i,j,k)$&$N$&$L$&$S$&$B$\\\hline
012&78&2&8212&144\\021&78&3&9436&144\\
102&84&3&9252&144\\120&84&3&11604&144\\
201&84&3&11436&144\\210&84&3&12564&144
\end{tabular}
\end{center}
For completeness, the following exact polynomial-arithmetic verification regenerates the table (Python with SymPy). No point sampling or root tolerance is involved.
\begin{Verbatim}
from itertools import permutations
from sympy import ZZ
from sympy.polys.rings import ring
R,p,q,r,t,m,n = ring('p,q,r,t,m,n', ZZ)
x,y,z = p+q+r,p+q,p
D = (x+y)*(y+z)*(z+x); H=x*y+y*z+z*x
for order in permutations(range(3)):
    T=t+m+n,t+m,t
    u,v,w=[T[i] for i in order]
    G=x*(v+w)**2+y*(w+u)**2+z*(u+v)**2
    J=(v*w*(z+x)*(x+y)+w*u*(x+y)*(y+z)
       +u*v*(y+z)*(z+x))
    F=3*D*G-4*H*J
    print(order,len(F),min(F.values()),sum(F.values()),
          F[(4,0,0,2,0,0)])
\end{Verbatim}

Positive $L$ and nonnegative variables give $F\ge144p^4t^2>0$. This proves the theorem.\hfill$\square$

\begin{thebibliography}{1}
\bibitem{r41meiji:liu} J. Liu, A weighted inequality involving the sides of a triangle,
\textit{Creative Math. Inform.} \textbf{19} (2010), no.~2, 160--168, Sec.~4, Conjecture 4.1.
\end{thebibliography}
\par\endgroup

% Source: /workspace/scratch/81914cec2c78/work/archive/Tournament_Round_41_Final/Tournament_Round_41/meiji/supplement/paper.tex
% Edited source: /workspace/scratch/81914cec2c78/work/audit_34_41/cleaned/r41_meiji_supplement.tex
\begingroup

\subsection*{Sharp refinement: The sharp constant in Liu's weighted side inequality}

\paragraph{Theorem.} In Conjecture 4.1 of Jian Liu, \textit{A weighted inequality involving the sides of a triangle}, Creative Math. Inform. 19 (2010), Eq.~(4.30), the coefficient $3/4$ can be replaced by the best possible coefficient $3/8$. Namely, for $a,b,c$ the sides of a nondegenerate triangle, $s=(a+b+c)/2$ and $x,y,z>0$,
\[
 \sum_{\rm cyc}\frac{(s-b)(s-c)}{y+z}
 \le\frac38\,\frac{xa^2+yb^2+zc^2}{xy+yz+zx}.
\]
Equality holds exactly at equilateral triangles with equal weights.

\paragraph{Proof.} Set $(u,v,w)=(s-a,s-b,s-c)$, $W=(x,y,z)$, $H=xy+yz+zx$, $S=x+y+z$, $D=(x+y)(y+z)(z+x)$. Multiplying the difference of sides (right minus left) by $8H$ gives
\[
 Q=3\sum_{\rm cyc}x(v+w)^2-8H\sum_{\rm cyc}\frac{vw}{y+z}.
\]
In the coordinates $(u,v,w)$ its symmetric matrix $A$ has entries
\[
 A_{ii}=3(S-W_i),\qquad
 A_{ij}=3W_k-4H/(S-W_k)\quad(\{i,j,k\}=\{1,2,3\}).
\]
Its first leading principal minor is positive. Its second is
\[
 \frac{H}{(x+y)^2}\big(9x^2+2xy+8xz+9y^2+8yz\big)>0.
\]
For the third, simultaneously permute coordinates to put $(x,y,z)=(p+q+r,p+q,p)$ with $p>0$, $q,r\ge0$. By expansion $\det A=4P/D$, with
\[
 P=\sum_{h=0}^4p^h\sum_{j=0}^{6-h} C_{h,j} q^{6-h-j}r^j.
\]
The rows below list $C_{h,j}$ in increasing order of $j$. They follow by substituting the entries into $A_{11}A_{22}A_{33}+2A_{12}A_{23}A_{31}-A_{11}A_{23}^2-A_{22}A_{31}^2-A_{33}A_{12}^2$ and collecting terms.
\begin{center}\begin{tabular}{c|rrrrrrr}
$h$ &$j=0$&1&2&3&4&5&6\\\hline
4&72&72&72\\3&202&303&273&86\\2&198&396&375&177&24\\
1&78&195&204&111&24&0\\0&10&30&36&22&6&0&0
\end{tabular}\end{center}
Thus $P\ge72p^4(q^2+r^2)$. Unless $q=r=0$, Sylvester's criterion proves $A$ positive definite. If $q=r=0$, then $A=3p(3I-\mathbf1\mathbf1^{\mathsf T})$ is positive semidefinite, with kernel the span of $(1,1,1)$. As $(u,v,w)$ is nonzero, these cases give both the inequality and its exact equality statement. An equilateral triangle and equal weights attain equality with a strictly positive numerator on the right, so no smaller constant can work.\hfill$\square$

\par\endgroup

% Source: /workspace/scratch/81914cec2c78/work/archive/Tournament_Round_41_Final/Tournament_Round_41/taisho/paper.tex
% Edited source: /workspace/scratch/81914cec2c78/work/audit_34_41/cleaned/r41_taisho.tex
\begingroup
\newcommand{\D}{\mathbb D}
\newcommand{\N}{\mathbb N}
\Needspace{9\baselineskip}
\section{\texorpdfstring{Infinitely many poles for the generating function of $(4321,4231)$-avoiding involutions}{Infinitely many poles for the generating function of (4321,4231)-avoiding involutions}}\label{section:r41taisho}
\noindent{\small\sffamily\color{muted}Round 41 \enspace / \enspace Taisho}\par\medskip


\begin{sectionabstract}
Barnabei, Castronuovo and Silimbani asked whether the length generating
function of involutions avoiding both $4321$ and $4231$ is non-D-finite.
We prove that it is: a meromorphic continuation has infinitely many real
poles approaching $1/2$ from below. We use their continued fraction as
input. The central step is to count negative directions of an associated
tridiagonal matrix and to show that every real singularity so detected
survives in its boundary resolvent. Specialization settles the question
for the descent-refined series as well.
\end{sectionabstract}

Let $a_{n,k}$ count involutions on $n$ letters avoiding both indicated
classical patterns and having $k$ descents. Set
$F(x,y)=\sum_{n,k}a_{n,k}x^ny^k$ and $f(x)=F(x,1)$.
Our target is Question 8.17 of \cite{r41taisho:BCS}, also recorded in OEIS
A399659 \cite{r41taisho:OEIS}. All coefficients in $x$ are polynomials in $y$.
\begin{theorem}
The germ $G(s)=f(s/(1+s^2))$ continues meromorphically to $|s|<1$
and has infinitely many distinct positive poles accumulating at $s=1$.
Thus a continuation of $f$ has infinitely many real poles below $1/2$
accumulating there. Neither $f$ nor $F$ is D-finite.
\end{theorem}

\subsection*{A matrix for the known continued fraction}
Write $q=s^2$, $C=(1-q)^2$ and $x=s/(1+s^2)$; this inverts the
branch $s=(1-\sqrt{1-4x^2})/(2x)$ regular at zero.
Define a symmetric tridiagonal matrix $J(s)$ with indices $n\ge0$ by
\[
 d_0=1-q-s,\qquad d_n=C-sq^n(1+q)\quad(n\ge1),\qquad
 b_n=sq^{n+1}\quad(n\ge0).
\]
Here $d_n$ is the diagonal entry and $b_n$ joins indices $n,n+1$.
These formulas give a norm-holomorphic bounded operator on
$\ell^2(\N_0)$ throughout $\D=\{|s|<1\}$, since its three diagonals
after subtracting $C I$ have locally uniformly vanishing tails.
At $s=0$ it equals the identity. The continued fraction of Corollary 8.19 of \cite{r41taisho:BCS}, also reproduced in the author's OEIS entry \cite{r41taisho:OEIS}, reads
\[
 G(s)=(1-q^2)\left(d_0-\frac{b_0^2}{
 d_1-\frac{b_1^2}{d_2-\cdots}}\right)^{-1}.
\]
Indeed $b_0^2=s^2q^2$ and $b_n^2=s^2q^{2n+2}$ for $n\ge1$.
Finite Schur elimination identifies the truncated fraction on levels
$0,\ldots,M$ with $(1-q^2)(J_{[0,M]}^{-1})_{00}$.
Extend the finite matrix by $CI$ on the omitted tail. The extensions
converge to $J$ in norm locally uniformly in $\D$; their inverses thus
converge on a sufficiently small closed disk about zero. The formal
convergents stabilize coefficientwise because each new partial
numerator has increasing positive order and each denominator constant
term one. Uniform analytic convergence and Cauchy's formula identify
the Taylor series of the limit with the formal fraction in \cite{r41taisho:BCS}.
This establishes the germ identity
\[
 G(s)=(1-q^2)(J(s)^{-1})_{00}.
 \tag{1}
\]

\subsection*{Tail elimination and uncancelled poles}
Fix $0<r<1$. For $|s|\le r$, $|C|\ge (1-r^2)^2$;
the compression $T$ of $J$ to indices $n\ge M$ satisfies
$\|T-C I\|=O_r(r^{2M})$. Choose $M\ge2$ large enough that
this is smaller than $(1-r^2)^2/2$. A Neumann series shows that
$T^{-1}$ is holomorphic on a neighborhood of the closed disk; for real
$0\le s\le r$, $T$ is positive definite. Write, with transposition
rather than conjugation for complex $s$,
\[
 J=\begin{pmatrix}A&E\\E^{\mathsf T}&T\end{pmatrix},
 \qquad H=A-E T^{-1}E^{\mathsf T}.
\]
Here $A$ is $M$ by $M$ and $E$ has one possibly nonzero entry.
Thus the holomorphic, finite tridiagonal matrix $H$ differs from $A$
only at its last diagonal entry. Block elimination gives
$(J^{-1})_{00}=(H^{-1})_{00}$ whenever $H$ is invertible;
$J$ is invertible if and only if $H$ is. Because $H(0)=I$,
$\det H$ has isolated zeros. Equation (1) therefore defines a single
meromorphic continuation on $\D$ (the diskwise formulas agree near
zero and hence agree on overlaps).

At any real $s\in(0,r]$ the off-diagonals $b_0,\ldots,b_{M-2}$
are nonzero. Let $D_j$ be the determinant of the trailing principal
block of $H$ beginning at $j$, setting $D_M=1$. For $j\le M-2$,
\[
 D_j=H_{jj}D_{j+1}-b_j^2D_{j+2}.
\]
If $D_0=D_1=0$ the recurrence forces $D_M=0$, a contradiction.
By Cramer's rule $(H^{-1})_{00}=D_1/D_0$, so every positive
real zero of $\det H$ is a genuine pole of $G$: the prefactor
$1-q^2$ is not zero. Moreover the nullity of $H$ is at most one:
the tridiagonal row recurrence determines every coordinate of a null
vector from its first coordinate.

\subsection*{Why the poles cannot terminate}
Fix an integer $m\ge1$. As $s\uparrow1$, each of
$d_2,d_4,\ldots,d_{2m}$ tends to $-2$. Choose $r<1$ so these
entries are strictly negative at $s=r$. In the preceding construction
choose $M>2m+1$; increasing $M$ only improves the tail estimate.
The compression of $H(r)$ to the even coordinate vectors
$e_2,e_4,\ldots,e_{2m}$ is diagonal with these entries, and is
negative definite. Hence $H(r)$ has at least $m$ strictly negative
eigenvalues, whereas $H(0)$ has none. The ordered eigenvalues of this
finite real symmetric matrix vary continuously on $[0,r]$. Each of
its $m$ lowest eigenvalues therefore vanishes somewhere in $(0,r)$.
The zero locations for these ordered eigenvalues are distinct, since
the nullity at every point of $(0,r)$ is at most one. By the tail-elimination argument above
each such location is a pole of $G$.

As $m$ was arbitrary, there are infinitely many poles, and on every
compact subdisk there are only finitely many. Hence there is a sequence
of real poles tending to $1$. The map $s\mapsto s/(1+s^2)$ is
one-to-one on $\D$ (two distinct inverse images would have product
one), is holomorphic there and has nonzero derivative. It transports
the continuation to its open image; on $(0,1)$ it increases onto
$(0,1/2)$. Thus $f$ has the asserted infinitely many poles.

A univariate D-finite function satisfies a nonzero linear ordinary
differential equation with polynomial coefficients. At any point where
the leading coefficient is nonzero, a meromorphic solution cannot
have a pole: the highest derivative would create a uniquely highest
pole order. Applied on the connected image domain, this contradicts
our pole count. Finally, D-finiteness of $F$ would give an equation
$\sum_{j=0}^h P_j(x,y)\partial_x^jF=0$ with polynomial coefficients
not all zero. Divide their largest common power of $y-1$ out of the
identity and substitute $y=1$, coefficientwise legitimately in
$\mathbb C[y][[x]]$. This gives a nonzero differential equation for
$f$, again impossible. This proves both assertions of Question 8.17.

\paragraph{Checks and scope.}
The accompanying exact check compares 33 coefficients of the boundary
resolvent with direct Motzkin-path dynamic programming and checks
the first ten term counts against direct involution enumeration. At
$s=99/100$ the test subspace gives at least 211 distinct poles before
$x=9900/19801$. This last bound also follows without numerical
diagonalization by checking
$sq^{422}(1+q)>(1-q)^2$ at that rational point.
The proof does not identify the individual pole orders and does not
require experimental estimates of their locations. The published
continued fraction, not the pole-count argument, is attributed to
the original authors.

\begin{thebibliography}{9}
\bibitem{r41taisho:BCS} M. Barnabei, N. Castronuovo and M. Silimbani,
\emph{Involutions avoiding 4321 and another pattern of length four},
arXiv:2609.39952v1 (30 September 2026), Question 8.17 and
Corollary 8.19.
\bibitem{r41taisho:OEIS} N. Castronuovo, OEIS A399659, 5 October 2026,
\texttt{https://oeis.org/A399659/internal}.
\end{thebibliography}

\par\endgroup

% Source: /workspace/scratch/81914cec2c78/work/archive/Tournament_Round_41_Final/Tournament_Round_41/showa/paper.tex
% Edited source: /workspace/scratch/81914cec2c78/work/audit_34_41/cleaned/r41_showa.tex
\begingroup

\Needspace{9\baselineskip}
\section{\texorpdfstring{A sharp exponent for Liu's quadratic angle-bisector inequality}{A sharp exponent for Liu's quadratic angle-bisector inequality}}\label{section:r41showa}
\noindent{\small\sffamily\color{muted}Round 41 \enspace / \enspace Showa}\par\medskip


\begin{sectionabstract}
We give a computer-assisted proof that the best exponent in Jian Liu's 2008 Problem 2 is $\log 2/\log(32/27)$. Its limiting obstruction is a flattened isosceles triangle, not the configuration determining the different constant in the altitude and sine problems of the same paper. The proof uses an integer-arithmetic cover for a determinant on the closed triangle of shapes, with separate derivative checks at both vanishing shapes.
\end{sectionabstract}
\paragraph{Problem and theorem.}
Write $w_a,w_b,w_c$ for the internal angle-bisector lengths of a triangle of sides $a,b,c$. Problem 2, equation (4.7), in J. Liu [1] asks for the largest real $k$ for which
\[
 a^kx^2+b^ky^2+c^kz^2\;\ge\;
 (2/\sqrt3)^k\,(w_a^kyz+w_b^kzx+w_c^kxy)
\]
holds for all nondegenerate triangles and all real $x,y,z$. The source supplies the case $k=4$. We prove that the maximum is $K=\log 2/\log(32/27)=4.0797593302\ldots$.

\paragraph{Reduction.}
Set $u_i=a_i/(a+b+c)$, $a_1=a,a_2=b,a_3=c$. Normalize the diagonal by setting $X=a^{k/2}x$, and likewise for $Y,Z$. The elementary identity $w_a^2=bc(1-a^2/(b+c)^2)$ gives off-diagonal matrix entries $-p_i$ opposite diagonal $i$, where
\[
 g(u)=\frac43\left(1-\frac{u^2}{(1-u)^2}\right),
 \qquad p_i=\tfrac12g(u_i)^{k/2}.
\]
The determinant of the unit-diagonal matrix is
$D=1-\sum p_i^2-2p_1p_2p_3$. For $k>0$ on the closed triangle of perimeter fractions, all $p_i$ are nonnegative, so $D\ge0$ forces $p_i^2\le1$ individually and hence is equivalent to positive semidefiniteness by all principal minors. At the limit $u=(1/4,1/4,1/2)$ the determinant is $1-(32/27)^k/2$. Every $k>K$ therefore fails by continuity along nondegenerate isosceles triangles. Only the endpoint proof remains.

\paragraph{Domain and special shapes.}
Order $u_1\le u_2\le u_3$ and put
\[
 (u_1,u_2,u_3)=(1/4+t/2-v,\;1/4+t/2+v,\;1/2-t).
\]
The closed ordered domain is $0\le t\le E=1/6$, $0\le v\le1/4-3t/2$. For $k=K$ both $D(0,0)$ and $D(E,0)$ vanish exactly; at the latter all $p_i=1/2$. At every $t$ the function is even in $v$; at $(E,0)$ its first derivatives in both coordinates vanish by permutation symmetry and the fixed sum. At $u=1/2$ the scalar function $p(u)$ has two continuous one-sided derivatives since $h=K/2>2$.

We certify $D>0$ on closed rectangles not containing a special shape whenever its direct lower estimate is positive. For rectangles near $(0,0)$ we may instead verify $D_t>0$ and $D_{vv}>0$ on the larger rectangle spanned by their upper endpoints and the origin. Integrating the latter derivative along $t=0$ using $D_v(0,0)=0$, then integrating $D_t$ at fixed $v$, proves nonnegativity. Near $(E,0)$ we instead check $D_{tt},D_{vv}>0$ and $D_{tt}D_{vv}-D_{tv}^2>0$ on the hull with $(E,0)$, so the second derivative on every segment from $(E,0)$ is positive. These local hulls stay in the region with all three fractions between $0$ and $1/2$ (also when their ordering changes).

\paragraph{Rigorous bounds.}
Writing $g=g(u)$ and $r=g^{h-2}$, we use
\[
 p=rg^2/2,\qquad p'=hrgg'/2,\qquad
 p''=(h/2)r\big((h-1)(g')^2+gg''\big).
\]
Here $g'=-8u/[3(1-u)^3]$, $g''=-8(1+2u)/[3(1-u)^4]$. For distinct $i,j,l$ the derivatives with respect to independent fractions are
\[
 D_i=-2(p_i+p_jp_l)p_i',\quad
 D_{ii}=-2(p_i+p_jp_l)p_i''-2(p_i')^2,\quad
 D_{ij}=-2p_i'p_j'p_l.
\]
Their contraction with $(1/2,1/2,-1)$ and $(-1,1,0)$ yields the derivatives in the rectangle coordinates. Exact integer comparisons give $1/26<h-2<1/25$. On an interval of nonnegative bases $g$, the minimum and maximum of $r$ thus follow from monotonicity in the base and checking which side of $1$ each endpoint lies on. The entire enclosed program uses integer endpoints of scale $10^{10}$ with outward rounding; $n$th roots are enclosed by bisection using integer powers. Rational arguments of $g$ at the enclosed interval endpoints suffice since $g$ is decreasing. Intersection with $[0,1/2]$ is valid on the true domain; on the local derivative hulls the intersection does not remove any point in those hulls. Each remaining box is quartered or discarded if its least $v$ exceeds the order boundary even at its least $t$. The complete tree uses 14,825 strictly positive leaves, 5 boundary leaves, 144 equilateral leaves and 309 wholly exterior leaves, to maximum depth 12. This proves nonnegativity at the endpoint for every triangle. This finite certificate bounds continuous sets, not points selected therefrom.

\paragraph{Scope.}
This concerns only the distinct bisector problem, not Problems 1 or 3. The closure was not adjudicated as a race-winning proof. No external refereeing or claim of absolute priority.
\paragraph{Reference.}
[1] J. Liu, \emph{The Composite Theorem of Ternary Quadratic Inequalities and its Applications}, RGMIA Research Report Collection 11 (2008), no. 4, art. 13, Sec.~4, eqs.~(4.6)--(4.7).

\subsection*{Reproducible certificate (Python 3)}
\begingroup\scriptsize\begin{Verbatim}
"""Exact outward-rounded integer interval certificate, Python 3 standard library."""
from fractions import Fraction as Q
from functools import lru_cache
from collections import Counter
S=10**10
def ceil(n,d):return -((-n)//d)
class I:
    def __init__(self,l,h): self.l,self.h=l,h
    def __add__(a,b):
        b=iv(b);return I(a.l+b.l,a.h+b.h)
    __radd__=__add__
    def __neg__(a):return I(-a.h,-a.l)
    def __sub__(a,b):return a+-iv(b)
    def __rsub__(a,b):return iv(b)+-a
    def __mul__(a,b):
        b=iv(b);v=[x*y for x in (a.l,a.h) for y in (b.l,b.h)]
        return I(min(v)//S,ceil(max(v),S))
    __rmul__=__mul__
    def __truediv__(a,b):
        b=iv(b);assert b.l>0;return a*I(S*S//b.h,ceil(S*S,b.l))
def iv(x,y=None):
    if isinstance(x,I):return x
    x=Q(x);y=x if y is None else Q(y)
    return I((S*x.numerator)//x.denominator,ceil(S*y.numerator,y.denominator))
assert 32**106 < 2**26*27**106 and 32**102 > 2**25*27**102
H=iv(Q(53,26),Q(51,25))
@lru_cache(None)
def floorroot(a,n):
    # floor of S*(a/S)^(1/n)
    k=a*S**(n-1);l=0;h=2*S
    assert h**n>k
    while h-l>1:
        m=(h+l)//2
        if m**n<=k:l=m
        else:h=m
    return l
def rho(g):
    assert g.l>=0
    return I(floorroot(g.l,25 if g.l<=S else 26),
             floorroot(g.h,26 if g.h<=S else 25)+1)
@lru_cache(None)
def at(l,u,der):
    x=I(l,u);a=1-x
    # monotone rational map g(x) evaluated at endpoints to avoid dependency
    def gp(j):
        j=Q(j,S);v=Q(4,3)*(1-(j/(1-j))**2);return iv(v)
    g=I(gp(u).l,gp(l).h);r=rho(g)
    P=r*g*g/2
    if not der:return P
    gp1=(-8*x)/(3*a*a*a)
    gp2=(-8*(1+2*x))/(3*a*a*a*a)
    p1=H*r*g*gp1/2
    p2=(H/2)*r*((H-1)*gp1*gp1+g*gp2)
    return P,p1,p2
def bounds(t,T,v,V,der=False):
    xu=[(Q(1,4)+t/2-V,Q(1,4)+T/2-v),
        (Q(1,4)+t/2+v,Q(1,4)+T/2+V),(Q(1,2)-T,Q(1,2)-t)]
    out=[]
    for l,u in xu:
        x=iv(max(0,l),min(Q(1,2),u))
        assert 0<=x.l<=x.h<=S//2
        out.append(at(x.l,x.h,der))
    if not der:
        a,b,c=out
        return (1-a*a-b*b-c*c-2*a*b*c).l>0
    p,d,e=zip(*out);grad=[]
    h=[[0]*3 for _ in range(3)]
    for i in range(3):
        j=(i+1)%3;k=(i+2)%3
        w=-2*(p[i]+p[j]*p[k]);grad.append(w*d[i])
        h[i][i]=w*e[i]-2*d[i]*d[i]
        h[i][j]=h[j][i]=-2*d[i]*d[j]*p[k]
    ts=[Q(1,2),Q(1,2),-1];vs=[-1,1,0]
    def hess(a,b):
        return sum(a[i]*b[j]*h[i][j] for i in range(3) for j in range(3))
    Dt=sum(a*b for a,b in zip(grad,ts));Dvv=hess(vs,vs)
    Dtt=hess(ts,ts);Dtv=hess(ts,vs)
    return Dt,Dvv,Dtt,Dtv
C=Counter()
E=Q(1,6)
def rect(t,T,v,V,lev):
    C['depth']=max(C['depth'],lev)
    if v>Q(1,4)-3*t/2:
        C['exterior']+=1;return
    if bounds(t,T,v,V):
        C['positive']+=1;return
    if T<=Q(1,32) and V<=Q(1,16):
        Dt,Dvv,_,_=bounds(Q(0),T,Q(0),V,True)
        if Dt.l>0 and Dvv.l>0:
            C['boundary']+=1;return
    if t>=E-Q(1,32) and V<=Q(1,32):
        _,Dvv,Dtt,Dtv=bounds(t,E,Q(0),V,True)
        if Dvv.l>0 and Dtt.l>0 and (Dtt*Dvv-Dtv*Dtv).l>0:
            C['equilateral']+=1;return
    assert lev<22,(t,T,v,V)
    m=(t+T)/2;n=(v+V)/2
    for a,b in [(t,m),(m,T)]:
        for c,d in [(v,n),(n,V)]:rect(a,b,c,d,lev+1)
if __name__=='__main__':
    rect(Q(0),E,Q(0),Q(1,4),0)
    print('verified rectangles:',dict(sorted(C.items())))
\end{Verbatim}
\endgroup

\par\endgroup

% Source: /workspace/scratch/81914cec2c78/work/archive/Tournament_Round_41_Final/Tournament_Round_41/keio/paper.tex
% Edited source: /workspace/scratch/81914cec2c78/work/audit_34_41/cleaned/r41_keio.tex
\begingroup

\Needspace{9\baselineskip}
\section{\texorpdfstring{The sharp upper bound for the sum of cubes of the medians in a nonobtuse triangle}{The sharp upper bound for the sum of cubes of the medians in a nonobtuse triangle}}\label{section:r41keio}
\noindent{\small\sffamily\color{muted}Round 41 \enspace / \enspace Keio}\par\medskip


\begin{sectionabstract}
We prove Conjecture 5 of Jian Liu's 2023 paper on acute triangles. At fixed longest normalized squared side, averaging the two shorter squared sides increases a suitable difference of power sums. The resulting one-variable problem has a derivative controlled by a perfect square. The optimal coefficient is approached but not attained by strictly acute triangles.
\end{sectionabstract}

\textbf{Theorem.} Let $a,b,c$ be the sides of a nondegenerate nonobtuse triangle and $m_a,m_b,m_c$ its medians from the opposite vertices. Then
\[
 m_a^3+m_b^3+m_c^3\le
 \frac{\sqrt2+5\sqrt5}{8(1+\sqrt2)}(a^3+b^3+c^3).
\]
Equality holds exactly at right isosceles triangles. The same coefficient is sharp if only strictly acute triangles are admitted. This is the statement in \cite[Conjecture 5, (4.5)]{r41keio:liu}, with its boundary interpretation made explicit.

\textbf{Proof.} Write $C=(\sqrt2+5\sqrt5)/(1+\sqrt2)$. Scale to $a^2+b^2+c^2=1$ and, after reordering, put $x=a^2\ge y=b^2\ge z=c^2$. Nonobtuseness gives $1/3\le x\le 1/2$. Apollonius's median formula gives $4m_a^2=2-3x$ and its cyclic analogues. Thus the desired inequality is
\[
 f(x)+f(y)+f(z)\le0,\qquad
 f(s)=(2-3s)^{3/2}-Cs^{3/2}.
\]

\emph{The isosceles profile.} Suppose the two shorter squared sides agree. In unnormalized units let their common value be $1$ and the remaining squared side $t$. Then $1\le t\le2$. The ratio of eight times the median cube sum to the side cube sum is
\[
 F(t)=\frac{(4-t)^{3/2}+2(2t+1)^{3/2}}{t^{3/2}+2}.
\]
Writing $u=\sqrt t$, direct differentiation gives
\[
 \frac{(t^{3/2}+2)^2}{3}F'(t)
  =(4-u)\sqrt{2t+1}-(2u+1)\sqrt{4-t}.
\]
Both terms on the right are positive before subtraction. The difference of their squares is
\[
 (4-u)^2(2u^2+1)-(2u+1)^2(4-u^2)
 =6(u-1)^2(u^2+2).
\]
Therefore $F$ is strictly increasing on $[1,2]$, so $F(t)\le F(2)=C$, and also $C>F(1)=3\sqrt3$.

\emph{Averaging the shorter squared sides.} Fix $x$, put $s=1-x$, and write $y=s/2+h$, $z=s/2-h$. Along the segment from $h=0$ to the specified $h\ge0$, the ordering and the nonobtuse bound are preserved. Here $y+z=s\le2/3$, so $2-3y\ge3z$ and $2-3z\ge3y$. Rationalization yields
\[
 \frac{d}{dh}[f(y)+f(z)]=\frac32(y-z)
 \left(\frac{9}{\sqrt{2-3y}+\sqrt{2-3z}}
              -\frac C{\sqrt y+\sqrt z}\right).
\]
Since the first denominator is at least $\sqrt3(\sqrt y+\sqrt z)$, the parenthesis is at most $(3\sqrt3-C)/(\sqrt y+\sqrt z)<0$. Hence averaging increases the sum strictly unless $y=z$. It follows that
\[
 f(x)+f(y)+f(z)\le f(x)+2f((1-x)/2)\le0;
\]
the second inequality is the isosceles result with $t=2x/(1-x)$. Equality requires both $y=z$ and $t=2$, giving squared side ratio $2:1:1$, and this ratio indeed gives equality. Finally $t\uparrow2$ through $1<t<2$ supplies strictly acute triangles with $F(t)\to C$.\hfill$\square$

\begin{thebibliography}{1}
\bibitem{r41keio:liu} J. Liu, Two inequalities involving circumradius, inradius and medians of an acute triangle, \emph{Advances in Inequalities and Applications} (2023), article 1, Conjecture 5, p.~21.
\href{https://doi.org/10.28919/aia/8232}{doi:10.28919/aia/8232}.
\end{thebibliography}

\par\endgroup

% Source: /workspace/scratch/81914cec2c78/work/archive/Tournament_Round_42_Final/Tournament_Round_42/taisho/paper.tex
% Edited source: /workspace/scratch/81914cec2c78/work/audit_42_50/cleaned/r42_taisho.tex
\begingroup

\Needspace{9\baselineskip}
\section{\texorpdfstring{The anisotropic generating function of pseudo-square parallelogram polyominoes is not D-finite}{The anisotropic generating function of pseudo-square parallelogram polyominoes is not D-finite}}\label{section:r42taisho}
\noindent{\small\sffamily\color{muted}Round 42 \enspace / \enspace Taisho}\par\medskip


\begin{sectionabstract}
We prove Conjecture 1 of Brlek, Frosini, Rinaldi and Vuillon (2006).
Encoding path separation by the uniqueness of a maximal cyclic window
turns a fixed-width section into a rational polytope enumeration.
The leading coefficient forces a prime to divide every quasiperiod
when the width equals that prime. The resulting large cyclotomic
factors cannot be introduced by a recurrence of bounded degree.
\end{sectionabstract}

Let $q_{m,n}$ count pseudo-square parallelogram polyominoes with
$m$ rows and $n$ columns, up to translation. Put
\[
 H_n(x)=\sum_{m\ge1}q_{m,n}x^m,\qquad
 Q(x,y)=\sum_{n\ge1}H_n(x)y^n.
\]
A parallelogram polyomino has northeast upper and lower boundary
paths meeting only at their endpoints. It is pseudo-square when its
boundary factors into two nonempty paths followed by their opposite
reversals. We use the structural result of Brlek et al.\
\cite[Prop.~3.3 and the preceding characterization]{r42taisho:BFRV}: in the
northeast coding the upper and lower words of such a polyomino have
the form $U=XY$, $V=YX$, respectively, and the pair $(X,Y)$
is unique. Conversely, if these two northeast paths from the origin
are strictly separated until their common end, their boundary has the
required factorization. We count the full class, not only shapes with
flat bottom.
\begin{theorem}
$Q$ is not D-finite in $y$ over $\mathbb C(x)$.
In particular it is not bivariate D-finite.
\end{theorem}

\subsection*{The circular-window bijection}
Use $E,N$ for east and north. In a cyclic word $w$ with $m$
letters $N$ and $n$ letters $E$, let $L=m+n$. For a fixed length
$1\le k<L$, let $h_t$ count letters $N$ in the $k$ successive
positions starting at site $t$ (cyclic indices).
Cut at $0$ and $k$, giving $U=XY$ with $|X|=k$. The difference of
the north-prefix counts at time $t$ for $U$ and $V=YX$ is $h_0-h_t$:
the two quantities agree at zero and both increase by
$\mathbf1_{w_t=N}-\mathbf1_{w_{t+k}=N}$ at the next step. Since the
paths advance from each diagonal $x+y=t$ to the next, they have
the required strict separation exactly when
\[
 h_0>h_t\quad(1\le t<L). \tag{1}
\]

For an oriented circular arrangement list the lengths $g_1,\ldots,g_n$
of runs of $N$ following successive $E$'s; zeros are allowed.
With indices modulo $n$, set $S_{i,l}=\sum_{a=0}^{l-1} g_{i+a}$.
The full arc of $l$ gaps starting after the $i$th $E$ has length
$S_{i,l}+l-1$ and contains $l-1$ east steps: start after an exterior
$E$ and stop just before the $l$th subsequent $E$.
For $1\le l\le n$ put $K_l=l-1+\max_i S_{i,l}$. This is the maximum
length of an arc with at most $l-1$ east steps. Indeed, every such arc
can be extended to a full arc of $l$ gaps by including the two end runs
and, if necessary, further east steps and runs.
In particular $K_l>K_{l-1}$ for $l\ge2$, by nonnegativity.
For $l=n$ the $n$ candidates all have length $L-1$.
\begin{lemma}
For $m\ge1$ and $n\ge2$,
\[
 q_{m,n}=\frac1n
 \sum_{\substack{g_1+\cdots+g_n=m\\ g_i\ge0}}
 \;\sum_{l=1}^{n-1}
 {\bf1}\{S_{1,l},\ldots,S_{n,l}\hbox{ have a unique maximum}\}.
 \tag{2}
\]
\end{lemma}
\begin{proof}
A unique winning $k$-window in (1) must begin and end with $N$,
with an $E$ immediately before and immediately after it: each of the
two shifts by one site decreases the count. If the window contains
$r$ east steps it is thus a full arc with $l=r+1$ gaps. A longer
arc containing at most $r$ east steps would contain at least two
length-$k$ windows with at most $r$ east steps, contradicting
uniqueness, since $r$ is the minimum among all such windows.
Consequently $k=K_l$, and the winning full arc is unique.
In particular $l<n$ since $n\ge2$ and all $n$ full arcs for $l=n$
tie, starting at distinct sites.

Conversely suppose the maximum of $S_{i,l}$ is unique and $l<n$.
Its full arc has length $K_l$, which is positive: $m>0$ and
every gap occurs in some $l$-block. No window of that length has
fewer than $l-1$ east steps, since $K_{l-1}<K_l$ for $l>1$.
Any other length-$K_l$ window with exactly $l-1$ east steps extends
to a full $l$-gap arc; as its length already equals the largest
full-arc length, it must itself be full and equal to the winner.
It follows that (1) holds when the word is cut at the winner's
start. We have $K_l<L$ since $l<n$.

Label the $E$ immediately preceding the winner as the one followed by
$g_1$. Hence shapes are in bijection with $(g,l)$ for which
$S_{1,l}$ is the strict maximum. For reconstruction start the upper
word just after that labeled $E$ and take its length-$K_l$ prefix as
$X$. Conversely the unique factorization of a shape supplies the cut
and the labeling. Thus no marking is overcounted. By cyclic symmetry,
among ordered vectors each winning label occurs equally often,
giving (2).
\end{proof}

\subsection*{Period divisibility from the leading coefficient}
For fixed $n,l$ and winning index $i$, the contribution counts integer
vectors in $m$ times the region
\[
 z_j\ge0,\quad \sum_j z_j=1,\quad
 S_{i,l}(z)>S_{a,l}(z)\quad(a\ne i).
\]
This is a rational polytope with bounding faces removed. By
inclusion-exclusion on the removed equalities and Ehrhart's theorem
for closed rational polytopes \cite{r42taisho:BR}, its counting function is a
quasipolynomial in $m\ge1$. Thus $H_n\in\mathbb Q(x)$ has only
root-of-unity poles: on each residue class of a common period $T$
one sums a polynomial in $m$ times $x^m$, with denominator a power
of $1-x^T$.

For each $1\le l<n$ the $n$ cyclic blocks of length $l$ are distinct.
A tie is contained in a finite union of hyperplane sections
$S_{i,l}=S_{j,l}$ of the simplex. For $i\ne j$
the coefficient difference is nonzero and sums to zero, so this
equality is independent of $\sum g_a=m$. There are
$O_n((m+1)^{n-2})$ solutions to each equality: assign all but two
coordinates and solve two independent linear equations for the
remaining two, choosing an invertible pair of columns. Therefore
(2) gives, without restricting the residue class,
\[
 q_{m,n}=\frac{n-1}{n}\binom{m+n-1}{n-1}+O_n(m^{n-2})
 =\frac{n-1}{n(n-1)!}\,m^{n-1}+O_n(m^{n-2}).
 \tag{3}
\]
Take $n=p$ prime and let $T$ be any eventual quasiperiod. On every
class modulo $T$ the corresponding polynomial has degree $p-1$
and the leading coefficient in (3). Its $(p-1)$st forward difference
with step $T$ is $(p-1)T^{p-1}/p$. Since the entries $q_{m,p}$ are
integers, so is this difference. Hence $p\mid T$.

In particular some pole of $H_p$ has multiplicative order $d$
divisible by $p$. Indeed, if $T$ is the least common multiple
of the orders of all poles, then by partial fractions the coefficients
are eventually polynomial on residue classes modulo $T$, regardless
of pole multiplicities. Since $H_p$ is rational over $\mathbb Q$,
all $\varphi(d)$ primitive $d$th roots are poles.
The orders produced for increasing $p$ have unbounded cyclotomic
degrees $\varphi(d)\ge p-1$.

\subsection*{A differential equation cannot create these pole orbits}
Suppose $Q$ satisfied a nontrivial polynomial-coefficient linear
differential equation in $y$. Coefficient extraction (the usual
D-finite to P-recursive implication) supplies a recurrence
\[
 \sum_{j=0}^r P_j(h,x)H_{h+j}(x)=0
 \quad(h\hbox{ sufficiently large}),\qquad
 P_j\in\mathbb C[h,x],\quad P_r\ne0.
\]
The $x$-degrees are bounded independently of $h$. The polynomial
$P_r(h,\cdot)$ is identically zero at only finitely many integers $h$.
For each order $d$ above let $j_d$ be the first index with a primitive
$d$th-root pole. All primitive $d$th roots occur at this same first
index. A finite set of sections has finitely many poles, so we may
choose $d$ with $j_d$ large enough for the recurrence to apply at
$h=j_d-r$, $P_r(h,\cdot)\ne0$, and
$\varphi(d)>\deg_x P_r$. Earlier sections are regular at these
$\varphi(d)$ points. The recurrence forces $P_r(h,\cdot)$ to vanish
at every one, a contradiction. This proves the theorem.

\paragraph{Scope and check.}
This proves the authors' Conjecture 1 without assuming Conjecture 2.
We force an order divisible by $p$, not necessarily order $p$
itself. No conclusion about $Q(x,x)$ is implied. Independently of
the proof, direct boundary-pair enumeration agrees with (2) for all
105 positive dimension pairs with $m+n\le15$; the first nine
semiperimeter totals are $1,2,3,6,11,22,45,90,184$. For $n=1$,
omitted from (2), every shape is a rectangle and $H_1=x/(1-x)$.
\begin{thebibliography}{9}
\bibitem{r42taisho:BFRV} S. Brlek, A. Frosini, S. Rinaldi and L. Vuillon,
Tilings by translation: enumeration by a rational language approach,
\emph{Electronic J. Combin.} \textbf{13} (2006), R15.
\bibitem{r42taisho:BR} M. Beck and S. Robins, \emph{Computing the Continuous
Discretely: Integer-Point Enumeration in Polyhedra}, 2nd ed.,
Springer, 2015, Ch.~3.
\end{thebibliography}

\par\endgroup

% Source: /workspace/scratch/81914cec2c78/work/archive/Tournament_Round_42_Final/Tournament_Round_42/showa/paper.tex
% Edited source: /workspace/scratch/81914cec2c78/work/audit_42_50/cleaned/r42_showa.tex
\begingroup

\Needspace{9\baselineskip}
\section{\texorpdfstring{A signed polygonal duality for Liu's circumradius power conjecture}{A signed polygonal duality for Liu's circumradius power conjecture}}\label{section:r42showa}
\noindent{\small\sffamily\color{muted}Round 42 \enspace / \enspace Showa}\par\medskip


\begin{sectionabstract}
We prove both ranges of Jian Liu's polygonal circumradius conjecture (2011, Conjecture 3.6). An antiperiodic spectral inequality controls powers of the cosine of half an angle even when successive angles have different signs. Doubling the circumcenters produces a polygonal loop whose supporting-line distances are the original vertex distances. This loop retains the original winding number but need not be convex. Reciprocal powers admit a direct proof on the original polygon.
\end{sectionabstract}
\subsection{Statement}
Let $A_1,\ldots,A_n$ be the successive vertices of a simple polygon, $n\ge3$, and $P$ interior. Write $R_i=|PA_i|$, $\rho_i$ for the circumradius of $PA_iA_{i+1}$, and $c=\cos(\pi/n)$. Indices are cyclic. The source says polygon, without explicitly imposing convexity. We do not impose convexity. Expressions involving ordinary finite circumradii require noncollinear incident triples, automatically true for convex polygons with $P$ strictly interior. In a nonconvex polygon an interior point can lie on the extension of an edge, and that member's circumradius is then undefined. The source calls the $\rho_i$ circumcenters after (3.15); the triangle analogue (3.14), scalar inequality and dimensional consistency identify the intended quantities as circumradii.
\begin{theorem}
For every $0<p\le1$,
\[
 \sum_i(2\rho_i)^p\ge c^{-p}\sum_i R_i^p,
 \qquad
 \sum_i(2\rho_i)^{-p}\le c^p\sum_i R_i^{-p}.
\]
Each constant is attained at a centered regular polygon, and hence is sharp.
\end{theorem}
This is equation (3.17) of Liu [1], including its reversed inequality for negative exponents. The endpoint $p=1$ on the positive side is also (3.16).
\subsection{A power bound on a signed angle cycle}
\begin{lemma}
If $|\theta_j|<\pi$, $\sum_j\theta_j=2\pi$, and $s_j\ge0$, then
\[
 \sum_j s_j s_{j+1}\cos^p(\theta_j/2)\le c^p\sum_j s_j^2
 \qquad(0<p\le1).
\]
\end{lemma}
\begin{proof}
Put $z_j=s_j\exp(\frac{i}{2}\sum_{l<j}\theta_l)$, where $i$ is the imaginary unit. With $z_{n+1}=-z_1$, the real part of $\sum_{j=1}^n\overline z_j z_{j+1}$ is the desired left-hand side at $p=1$. The unitary cyclic shift with negative wrap has Fourier eigenvectors with step multipliers $\exp((2m+1)\pi i/n)$, $0\le m<n$. The Hermitian part of the shift has largest eigenvalue $c$, giving the result. For $p<1$, H\"older bounds the left-hand side by
\[
 \Big(\sum_j s_j s_{j+1}\cos(\theta_j/2)\Big)^p
 \Big(\sum_j s_j s_{j+1}\Big)^{1-p}.
\]
Apply the first case and $2s_j s_{j+1}\le s_j^2+s_{j+1}^2$.
\end{proof}
Suppose two nonzero vectors of lengths $a,b$ have signed principal angular difference $\theta$, $|\theta|<\pi$. Their distance $d$ satisfies
\[
 d^2=(a-b)^2+4ab\sin^2(\theta/2).
\]
Thus $|\sin\theta|/d\le\cos(\theta/2)/\sqrt{ab}$ for distinct vectors, with $\theta=0$ immediate. The distance from the origin to their joining line, $ab|\sin\theta|/d$, is at most $\sqrt{ab}\cos(\theta/2)$. For coincident vectors the same line-distance bound holds for any line through their common endpoint.
\subsection{Proof of both signs}
Take $P$ as origin, orient the polygon counterclockwise, and put $e_i=A_i/R_i$. The signed principal direction changes $\theta_i$ from $e_i$ to $e_{i+1}$ sum to $2\pi$, by winding number one. Indeed along each edge avoiding the origin the direction varies through the signed principal angle. Collinear triangles have already been excluded. The extended sine rule gives $1/(2\rho_i)=|\sin\theta_i|/|A_i-A_{i+1}|$. The local estimate and the lemma with $s_i=R_i^{-p/2}$ prove the reciprocal-power bound.

Let $Q_i$ be twice the circumcenter vector of $0,A_i,A_{i+1}$. By equal distances to the origin and the endpoints,
\[
 |Q_i|=2\rho_i,\qquad e_i\cdot Q_i=R_i,\qquad e_{i+1}\cdot Q_i=R_{i+1}.
\]
In particular $Q_{i-1}$ and $Q_i$ lie on the same line $e_i\cdot X=R_i$. Let $\psi_i$ be the signed principal change from $Q_{i-1}$ to $Q_i$. We verify that these changes sum to $2\pi$, without assuming that the $Q$-polygon is convex or simple. Insert the directions of $Q_i$ between $e_i$ and $e_{i+1}$ in the cyclic direction list. Each new increment has absolute value strictly smaller than $\pi/2$, by the positive dot products above. The two increments from $e_i$ via $Q_i$ to $e_{i+1}$ therefore sum to the principal increment $\theta_i$, not to one differing by a full turn. Regrouping the same increments from $Q_{i-1}$ via $e_i$ to $Q_i$ gives $\psi_i$ for the same reason. Thus $\sum_i\psi_i=2\pi$ and $|\psi_i|<\pi$.

The local line-distance bound yields $R_i\le\sqrt{|Q_{i-1}|\,|Q_i|}\cos(\psi_i/2)$, also when $Q_{i-1}=Q_i$. Take $p$th powers and sum using the lemma on the $Q$-cycle with weights $|Q_i|^{p/2}$. This proves the positive-power bound. For a centered regular polygon of radial distance $R$, all $2\rho_i=R/c$, proving sharpness.
\paragraph{Relation to earlier work.} Classical polygonal Erd\H{o}s--Mordell and Lenhard inequalities bound side distances or angle-bisector lengths in terms of vertex distances [2]. The unit-power spectral estimate is a weighted cyclic cosine inequality, not a novelty claim here. The application to [1] uses the doubled circumcenter directions, signed winding, and both exponent branches. No external priority guarantee or journal refereeing is claimed.
\begin{thebibliography}{9}
\bibitem{r42showa:liu} J. Liu, A new sharpening of the Erd\H{o}s--Mordell inequality and related inequalities, \emph{RGMIA Research Report Collection} 14 (2011), Article 85, esp. Conjecture 3.6, eq. (3.17).
\bibitem{r42showa:len} H.-C. Lenhard, Verallgemeinerung und Versch\"arfung der Erd\H{o}s--Mordellschen Ungleichung f\"ur Polygone, \emph{Arch. Math.} 12 (1961), 311--314. DOI: 10.1007/BF01650566.
\end{thebibliography}

\par\endgroup

% Source: /workspace/scratch/81914cec2c78/work/archive/Tournament_Round_42_Final/Tournament_Round_42/heisei/paper.tex
% Edited source: /workspace/scratch/81914cec2c78/work/audit_42_50/cleaned/r42_heisei.tex
\begingroup
\newcommand{\sgn}{\operatorname{sgn}}
\Needspace{9\baselineskip}
\section{\texorpdfstring{A sharp curvature transition for Richter's generalized sine and its inverse}{A sharp curvature transition for Richter's generalized sine and its inverse}}\label{section:r42heisei}
\noindent{\small\sffamily\color{muted}Round 42 \enspace / \enspace Heisei}\par\medskip


\begin{sectionabstract}
An explicit critical exponent separates the one-, two- and no-inflection regimes of the angular parametrization of the planar $\ell_p$ unit curve. Ordinary curvature and complete monotonicity are different here: even when an inverse is strictly convex in its argument, its first derivative already rules out complete monotonicity. A flat-endpoint obstruction settles the one parameter-dependent inverse that survives the first-derivative test. These results answer items (i) and (ii) of Yin, Huang and Lin's Open Problem~5.1, not items (iii)--(v).
\end{sectionabstract}

\subsection{The four functions}
The problem is about Richter's Euclidean-angle parametrization \cite{r42heisei:richt,r42heisei:yhl}, not the generalized sine arising from the $p$-Laplacian. At $p>0$ set
\[
 N_p(\theta)=(|\sin\theta|^p+|\cos\theta|^p)^{1/p},\quad
 S_p(\theta)=\frac{\sin\theta}{N_p(\theta)},\quad
 C_p(\theta)=\frac{\cos\theta}{N_p(\theta)}.
\]
When $0<\theta<\pi/2$, use $t=\tan\theta$. Then
\[
 S_p(\theta)=t(1+t^p)^{-1/p},\qquad
 C_p(\theta)=(1+t^p)^{-1/p}=S_p(\pi/2-\theta).
\]
The sine is strictly increasing and the cosine strictly decreasing here. Denote their inverses on this quadrant by $A_p,D_p:(0,1)\to(0,\pi/2)$, respectively. Solving for $t$ gives
\begin{equation}\label{r42heisei:inv}
 A_p(y)=\frac\pi2-\arctan\frac{(1-y^p)^{1/p}}y,\qquad
 D_p(y)=\frac\pi2-A_p(y).
\end{equation}
A smooth real function on an open interval is called completely monotone (CM) there if $(-1)^kf^{(k)}\ge0$ for all integers $k\ge0$. Our negative CM statements concern entire natural intervals. They are not assertions about every proper subinterval.

\subsection{Classification of angular curvature}
\begin{theorem}\label{r42heisei:th}
Write
\[
 H_p(t)=2-(p+1)t^{p-2}-(p-1)t^p.
\]
The signs of $S_p''(\theta)$ as $\theta$ runs from 0 to $\pi/2$ are as follows. Each sign change occurs at $\arctan t$ for a simple positive root of $H_p$:
\[
\begin{array}{c|c|c}
 p & \text{successive signs} & \text{number of inflections}\\ \hline
 0<p\le1 & -,+ & 1\\
 1<p<p_* & -,+,- & 2\\
 p=p_* & -\text{ except at one double zero} & 0\\
 p_*<p\le2 & - & 0\\
 p>2 & +,- & 1
\end{array}
\]
Here $p_*$ is the unique root in $(1,(1+\sqrt5)/2)$ of $H_p(t_p)=0$ where
\[
 t_p=\sqrt{\frac{(p+1)(2-p)}{p(p-1)}}.
\]
Numerically $p_*\simeq1.2676354$. In particular, $S_p$ and $C_p$ are strictly concave on the whole quadrant exactly for $p_*\le p\le2$, and never convex there.
\end{theorem}

\begin{proof}
Differentiating $S_p$ gives $S'_p=(1+t^2)(1+t^p)^{-1-1/p}$. Since $dt/d\theta>0$, its logarithmic derivative with respect to $t$ has the sign of $S''_p$. Multiplying this logarithmic derivative by $(1+t^2)(1+t^p)/t$ gives exactly $H_p(t)$. Further,
\[
 H_p'(t)=t^{p-3}\big((p+1)(2-p)-p(p-1)t^2\big).
\]
For $0<p\le1$ this is positive; $H_p$ starts at $-\infty$ and ends at a positive limit (2 when $p=1$). For $p>2$ the derivative is negative, and $H_p$ decreases from 2 to $-\infty$. At $p=2$, $H_p=-1-t^2<0$.

For $1<p<2$ the limits at both ends are $-\infty$. The unique maximum is at $t_p$; it is nondegenerate because the bracket in the formula for $H_p'$ has a simple zero. On $0<t\le1$ we have $(p+1)t^{p-2}\ge p+1>2$, hence $H_p<0$. On $t\ge1$ both subtracted terms strictly increase with $p$; the derivative of the last term with respect to $p$ is at least 1. Consequently $M(p)=\max_{t\ge1}H_p(t)$ is strictly decreasing: $M(q)\le M(p)-(q-p)$ for $1<p<q<2$. It is continuous because $M(p)=H_p(\max\{1,t_p\})$. As $p\downarrow1$, $H_p(2)\to1$, so $M(p)>0$ near 1. At the golden ratio $p=(1+\sqrt5)/2$, $t_p=1$ and $M(p)=2-2p<0$. This proves existence and uniqueness of the transition, and strict monotonicity on either side of the maximum gives the table. A single zero of the second derivative in an otherwise strictly negative interval does not spoil strict concavity.
\end{proof}

\paragraph{Cosine and inverses.} Reflect the table by $\theta\mapsto\pi/2-\theta$ for $C_p''$. At $y=S_p(\theta)$,
\[
 A_p''(y)=-\frac{S_p''(\theta)}{S_p'(\theta)^3},\qquad
 D_p''(y)=+\frac{S_p''(\theta)}{S_p'(\theta)^3}.
\]
These formulas give the convexity and concavity intervals, with all endpoints, for all four functions. In particular $A_p$ is strictly convex and $D_p$ strictly concave on the entire interval $(0,1)$ exactly for $p_*\le p\le2$. Neither has the opposite global curvature for any $p$. No angular or argument function is CM on its first-quadrant domain: $S_p$ and $A_p$ increase; $C_p$ and $D_p$ have strictly negative second derivative somewhere, by the table.

\paragraph{Quadrants, principal branches and their axes.}
For each open quadrant fold $\theta$ to $v=\arctan|\tan\theta|\in(0,\pi/2)$. Multiply $S_p(v)$ or $C_p(v)$ by the sign of the original numerator. Since $v$ is affine with slope $\pm1$, their second derivatives are transformed by exactly the same signs. This specifies all curvature intervals on every quadrant. No nonconstant periodic coordinate can be CM over a full period. The usual principal inverses at negative arguments are $A_p(-y)=-A_p(y)$ and $D_p(-y)=\pi-D_p(y)$; their curvatures change sign under this reflection. Neither principal branch is CM on $(-1,1)$ because it already fails on $(0,1)$.

For clarity where quadrant tables meet, $S_p(0)=0$, $C_p(0)=1$, the other axes following by symmetry. Near zero $S_p(\theta)=\tan\theta(1+|\tan\theta|^p)^{-1/p}$ has slope 1 from both sides. Its derivative formula gives $S_p'(\theta)-1=O(\theta^2)+[-1-1/p+o(1)]|\theta|^p$; thus $S_p''(0)=0$ for $p>1$ and no finite second derivative exists there for $p\le1$. Likewise $C_p(\theta)=1-(1/p)|\theta|^p(1+o(1))$, with one-sided slopes from its explicit derivative. It is not differentiable at zero for $p\le1$, has a first but no finite second derivative for $1<p<2$, and $C_p''(0)=-1$ for $p=2$, zero for $p>2$. Locally across zero, $S_p$ is neither convex nor concave (opposite curvatures on the two sides). The peak of $C_p$ is locally strictly concave for $p>1$, and neither convex nor concave for $p\le1$ (a strict peak with positive one-sided curvature). Signs and shifts give the same description at all other axes. The principal inverses at $y=0$ have the same differentiability order as $S_p$ by inversion at slope 1 and complementation.

\subsection{Complete monotonicity in the parameter}
\begin{theorem}
For every $0<\theta<\pi/2$ and $0<y<1$, none of the four functions $p\mapsto S_p(\theta), C_p(\theta), A_p(y), D_p(y)$ is CM on $(0,\infty)$.
\end{theorem}
\begin{proof}
For fixed $a,b>0$ let $N=(a^p+b^p)^{1/p}$. If $q>p$, $a/N,b/N$ lie strictly between 0 and 1 and their $p$-powers sum to 1; their $q$-powers sum to less than 1. Thus $(a^q+b^q)^{1/q}<N$. Both positive angular coordinates strictly increase with $p$. For fixed $y$ the inverse $A_p(y)$ strictly decreases with $p$ because $S_p$ is strictly increasing both in $\theta$ and in $p$. Its complement $D_p(y)$ strictly increases. The three increasing functions are not CM.

For the remaining inverse set $c=-\log y>0$. Since $1-y^p=1-e^{-cp}\le cp$, formula \eqref{r42heisei:inv} gives
\[
 0<\frac\pi2-A_p(y)\le y^{-1}(cp)^{1/p}=o(p)\qquad(p\downarrow0).
\]
Thus $A_p(y)$ extends continuously to $p=0$ with value $\pi/2$ and right secant slope 0. A convex function would have nondecreasing secant slopes from this endpoint, as follows by the chord inequality and continuity; yet every such slope is negative. Hence $A_p(y)$ is not even convex, and so is not CM.
\end{proof}

For arbitrary fixed angles off the axes, positive coordinates still strictly increase and negative ones already fail nonnegativity. At an axis each coordinate is constant in $p$, CM precisely if that constant is 0 or 1. For $0<y<1$, $A_p(-y)<0$, whereas $D_p(-y)=\pi/2+A_p(y)$ is not convex in $p$. At $y=0,\pm1$ the principal inverses are constants in $p$, CM precisely when nonnegative. These exhaust the boundary and sign exceptions.

\begin{thebibliography}{9}
\bibitem{r42heisei:yhl} L. Yin, L. Huang and X. Lin, Inequalities and bounds for the $p$-generalized trigonometric functions, \emph{AIMS Mathematics} 6 (2021), 11097--11108, Definition~1.1 and Open Problem~5.1(i),(ii), p.~11105. \href{https://doi.org/10.3934/math.2021644}{doi:10.3934/math.2021644}.
\bibitem{r42heisei:richt} W.-D. Richter, Generalized spherical and simplicial coordinates, \emph{J. Math. Anal. Appl.} 336 (2007), 1187--1202. \href{https://doi.org/10.1016/j.jmaa.2007.03.047}{doi:10.1016/j.jmaa.2007.03.047}.
\end{thebibliography}

\par\endgroup

% Source: /workspace/scratch/81914cec2c78/work/archive/Tournament_Round_42_Final/Tournament_Round_42/reiwa/paper.tex
% Edited source: /workspace/scratch/81914cec2c78/work/audit_42_50/cleaned/r42_reiwa.tex
\begingroup
\newcommand{\C}{\mathbb C}
\newcommand{\Z}{\mathbb Z}
\Needspace{9\baselineskip}
\section{\texorpdfstring{Limits of independent-domination roots for cozero-divisor graphs}{Limits of independent-domination roots for cozero-divisor graphs}}\label{section:r42reiwa}
\noindent{\small\sffamily\color{muted}Round 42 \enspace / \enspace Reiwa}\par\medskip


\begin{sectionabstract}
Fix primes $p<q$. We determine the finite limits of nonzero independent-domination roots for the cozero-divisor graphs of $\Z/(p^a q)$ as $a\to\infty$. Outside the exceptional pair $(2,3)$, the limit set consists of the whole unit circle together with the zeros, in its exterior, of one explicit convergent lacunary series. There are no other limits. The exterior component is discrete and admits effective tail bounds; it is nonempty, for example, for every $p=2$, $q\ge5$. For $(2,3)$ the sole nonzero root is $-a$, so no finite limit exists. The connection is an exact identification of reciprocal graph roots with the zeros of Taylor sections.
\end{sectionabstract}

\subsection{The limit set}
For a finite commutative ring $R$, its cozero-divisor graph has vertex set the nonzero nonunits of $R$; distinct vertices are adjacent if neither principal ideal contains the other. Write $P_a(z)=\sum_S z^{|S|}$ for the independent-domination polynomial of this graph for $R=\Z/(p^a q)$. Here $S$ ranges over maximal independent sets: failure to dominate is exactly the possibility of adding a further vertex while preserving independence. Rather \cite{r42reiwa:SMJ} asks for the limit set of nonzero zeros in this family indexed by $a\ge1$. Precisely, let $L_{p,q}$ consist of $\lim_{i\to\infty} z_i\in\C$ where $a_i\to\infty$, $z_i\ne0$ and $P_{a_i}(z_i)=0$.

\begin{theorem}\label{r42reiwa:main}
If $(p,q)=(2,3)$, then $L_{p,q}=\varnothing$. For all other prime pairs $p<q$, put $c=q-p-1$. Then
\[
 L_{p,q}=\{z:|z|=1\}\ \cup\
 \left\{z:|z|>1,\quad 1+\sum_{j=0}^\infty z^{-(1+c p^j)}=0\right\}.
\]
The exterior component lies in $1<|z|<2$, is discrete there, and is finite on each closed subannulus separated from the unit circle. For $p=2$, $q\ge5$, it contains exactly one negative real point.
\end{theorem}

This describes membership without referring to the graphs, their polynomials, or a sequence of root choices. On $|z|\ge R>1$ the absolute error after retaining indices $0,\ldots,m-1$ is at most
$R^{-(1+c p^m)}/(1-R^{-1})$ because all remaining exponents are distinct increasing integers. In particular the exterior zeros can be isolated on zero-free contours by Rouch\'e's theorem with a certified tail bound.

\subsection{Maximal chains and reversal}
For each divisor $d\mid N=p^a q$ with $1<d<N$, vertices of gcd $d$ with $N$ form an independent class of size $\varphi(N/d)$. Each such class generates one ideal, and incomparability among ideals is equivalent to incomparability of gcds by divisibility. A maximal independent set therefore consists of every vertex of each class on one maximal chain of the proper-divisor poset. This recovers the family enumeration of \cite[Thm.~2.2]{r42reiwa:DML}, not claimed here as new.

Adjoin the bottom and top divisors $1,N$. A maximal chain is now a saturated path in $[0,a]\times[0,1]$ from $(0,0)$ to $(a,1)$, where the coordinates denote the powers of $p$ and $q$. Exactly one step increases the second coordinate. Let $E_k$ be the size of the set from the path taking that step at first coordinate $k$ ($0\le k\le a$). Counting each chain gives $P_a(z)=\sum_{k=0}^a z^{E_k}$, including terms whose exponents coincide. The first path uses all gcds $p^i q$ for $0\le i<a$, so
\[
 E_0=\sum_{i=0}^{a-1}\varphi(p^{a-i})=p^a-1.
\]
Moving the step from $k$ to $k+1$ replaces the class at $(k,1)$ by the class at $(k+1,0)$. Thus
\[
 E_{k+1}-E_k=
 \begin{cases}
 c(p-1)p^{a-k-2}, &0\le k\le a-2,\\
 c+1, & k=a-1 .
 \end{cases}
\]
Indeed the difference in the first line is $(q-1)\varphi(p^{a-k-1})-\varphi(p^{a-k})$; the last line is $(q-1)-(p-1)$. Telescoping, $E_a-E_k=1+c p^{a-k-1}$ for $0\le k<a$. It follows that
\[
 w^{E_a}P_a(w^{-1})=1+\sum_{j=0}^{a-1}w^{1+c p^j}=:F_a(w).
\]
Always $c\ge0$. The only pair with $c=0$ is $(2,3)$, for which $F_a(w)=1+aw$. Hence the only nonzero graph root is $-a$ (including the case $a=1$). This also shows that an annulus uniform in $a$ cannot contain the nonzero roots for the exceptional pair; compare the explicit $a=4$ example in \cite[Ex.~2.2]{r42reiwa:DML}. Henceforth $c>0$.

\subsection{The three regions}
Set $d_j=1+cp^j$ and $F(w)=1+\sum_{j\ge0}w^{d_j}$. Its power-series radius is exactly $1$; coefficients equal one at arbitrarily large indices and vanish elsewhere. The nonconstant distinct ordinary Taylor sections of $F$ are exactly $F_a$.
\begin{enumerate}
\item On compact subsets of $|w|<1$, the sections converge uniformly to $F$. Any limit there of roots is therefore a zero of $F$. Conversely take a sufficiently small closed disk around a zero with no other zeros and a zero-free boundary. By Rouch\'e's theorem all sufficiently late sections have precisely as many zeros in this disk, with multiplicity, as $F$ does. Shrinking the disk gives a sequence tending to that zero. For $|w|\le1/2$ we have $|F_a(w)-1|,|F(w)-1|\le\sum_{k=2}^\infty (1/2)^k=1/2$; thus reciprocal inversion loses no point at the origin.
\item Jentzsch's theorem \cite{r42reiwa:J} states that every point of the circle of convergence of a power series is an accumulation point of zeros of its Taylor sections. It applies here to give the entire unit circle, not merely some points on it. The distinct sections exhaust our sequence; each fixed section has only finitely many roots.
\item Fix $r>1$. For $a\ge2$ and $|w|\ge r$, the sum of all terms before the leading one, divided by its modulus, is bounded by
\[
 a r^{-(d_{a-1}-d_{a-2})}
   =a r^{-c(p-1)p^{a-2}}\longrightarrow0.
\]
Consequently $F_a$ eventually has no zeros anywhere in $|w|\ge r$. This rules out finite exterior accumulation in the $w$-plane as well as escape to infinity.
\end{enumerate}
Inversion $z=1/w$ now proves the set formula, including the exclusion of $z=0$. Discreteness follows because $F$ is analytic in the disk and $F(0)=1$. Finally for $p=2$, $q\ge5$, $c=q-3$ is even; all $d_j$ are odd. Hence $F(-s)=1-\sum_{j\ge0}s^{d_j}$, whose sum on the right after the minus sign increases strictly from zero to infinity as $s$ runs from zero to one. Each finite partial sum tends to its number of terms at the upper endpoint, giving divergence there. Thus exactly one $s\in(0,1)$ gives an exterior negative point $-1/s$. This completes the proof.

\begin{thebibliography}{9}
\bibitem{r42reiwa:SMJ} B. A. Rather, Zeros of the independent domination polynomial of the cozero divisor graphs of ring $\mathbb Z_n$, \emph{Siberian Math. J.} (2026), DOI: 10.1134/S0037446626060017. Problem stated in the abstract.
\bibitem{r42reiwa:DML} B. A. Rather, Independent domination polynomial for the cozero divisor graph of the ring of integers modulo $n$, \emph{Discrete Math. Lett.} \textbf{13} (2024), 36--43. DOI: 10.47443/dml.2023.215.
\bibitem{r42reiwa:J} R. Jentzsch, Untersuchungen zur Theorie der Folgen analytischer Funktionen, \emph{Acta Math.} \textbf{41} (1916), 219--251. DOI: 10.1007/BF02422945.
\end{thebibliography}

\par\endgroup

% Source: /workspace/scratch/81914cec2c78/work/archive/Tournament_Round_42_Final/Tournament_Round_42/keio/paper.tex
% Edited source: /workspace/scratch/81914cec2c78/work/audit_42_50/cleaned/r42_keio.tex
\begingroup
\newcommand{\cy}{\sum_{\rm cyc}}
\Needspace{9\baselineskip}
\section{\texorpdfstring{An interior-point quadratic inequality of Jian Liu}{An interior-point quadratic inequality of Jian Liu}}\label{section:r42keio}
\noindent{\small\sffamily\color{muted}Round 42 \enspace / \enspace Keio}\par\medskip


\begin{sectionabstract}
We prove the subtriangle-circumradius inequality proposed as Conjecture 7 in Liu (2002) by interpolating the weighted Erd\H os--Mordell inequality with a half-angle estimate. 
\end{sectionabstract}
Let $P$ be interior to a nondegenerate triangle $A_1A_2A_3$. Write $R_i=PA_i$, let $r_i$ be the perpendicular distance from $P$ to the line $A_jA_k$, and denote by $S_i$ the circumradius of $PA_jA_k$. All indices $(i,j,k)$ below are cyclic and $x_1,x_2,x_3$ are arbitrary real numbers.

\textbf{Theorem 1 (Liu's Conjecture 7, equation (41) in \cite{r42keio:source}).}
\[
 \cy\frac{x_i^2}{\sqrt{r_j+r_k}}\ \ge\ \cy\frac{x_jx_k}{\sqrt{S_i}}.
\]

\emph{Proof.} The area of $PA_jA_k$ equals both $|A_jA_k|r_i/2$ and $|A_jA_k|R_jR_k/(4S_i)$, so $S_i=R_jR_k/(2r_i)$. Put $x_i=\sqrt{R_i}u_i$. The desired claim is equivalent to
\begin{equation}
 \cy\frac{R_i}{\sqrt{r_j+r_k}}u_i^2\ge\cy\sqrt{2r_i}\,u_ju_k.
 \label{r42keio:eq:main}
\end{equation}

We use two endpoints. The first is the classical weighted Erd\H os--Mordell inequality (for example \cite[Thm.~1.2]{r42keio:jiang}):
\begin{equation}
 \cy R_i u_i^2\ge 2\cy r_i u_ju_k.\label{r42keio:eq:a}
\end{equation}
For the second, let $A_i$ also denote the angle at the $i$th vertex and write $v_i=\sin(A_i/2)$. The ray $A_iP$ splits that angle into $\alpha_i$ and $\beta_i$. Thus $r_j+r_k=R_i(\sin\alpha_i+\sin\beta_i)\le2R_i v_i$. Concavity of sine yields $\sum v_i\le3/2$. The function $w\mapsto w/(1+w)$ is increasing and concave on $[0,\infty)$; Jensen therefore gives $\sum v_i/(1+v_i)\le1$. By weighted Cauchy--Schwarz,
\[
 (\textstyle\sum u_i)^2\le\left(\sum v_i/(1+v_i)\right)\sum(1+1/v_i)u_i^2
 \le\sum (1+1/v_i)u_i^2.
\]
Subtracting $\sum u_i^2$ and using the preceding angle estimate proves the second endpoint:
\begin{equation}
 \cy\frac{R_i}{r_j+r_k}u_i^2\ge\cy u_ju_k.\label{r42keio:eq:b}
\end{equation}

For clarity we give the interpolation step, a special case of Liu's composite theorem \cite[Thm.~3]{r42keio:composite}. If an endpoint is $\sum d_i u_i^2\ge\sum e_i u_ju_k$ with positive coefficients, write $c_i=e_i/\sqrt{d_jd_k}$. Then $\sum c_i t_jt_k\le\sum t_i^2$ for every $t$. For the normalized endpoints $c_i,c'_i$, take $t_i\ge0$ and apply Cauchy--Schwarz over edges:
\[
 \cy\sqrt{c_i c'_i}\,t_jt_k
 \le\sqrt{\left(\cy c_i t_jt_k\right)\left(\cy c'_i t_jt_k\right)}
 \le\sum t_i^2.
\]
Taking absolute values extends this to all real $t_i$. Undoing normalization gives diagonal coefficients $\sqrt{d_i d'_i}$ and edge coefficients $\sqrt{e_i e'_i}$. Applying this to \eqref{r42keio:eq:a} and \eqref{r42keio:eq:b} gives \eqref{r42keio:eq:main}. At the center of an equilateral triangle all $S_i=r_j+r_k$ and the nonzero constant weight vector gives equality.\hfill$\square$

\begin{thebibliography}{9}
\bibitem{r42keio:source} J. Liu, Two results and several conjectures of a kind of geometry inequality, \emph{J. East China Jiaotong Univ.} 19(3) (2002), 89--94. \href{https://journal.ecjtu.edu.cn/hdjtdx/article/pdf/20020327}{Publisher PDF}, p.~94.
\bibitem{r42keio:jiang} W.-D. Jiang, A new weighted Erd\H os--Mordell type inequality, \emph{Creative Math. Inform.} 17 (2008), 152--155.
\bibitem{r42keio:composite} J. Liu, The composite theorem of ternary quadratic inequalities and its applications, \emph{RGMIA Research Report Collection} 11(4) (2008), article 13. \href{https://rgmia.org/papers/v11n4/CTTQIApp.pdf}{Author's report}.
\end{thebibliography}

\par\endgroup

% Source: /workspace/scratch/81914cec2c78/work/archive/Tournament_Round_43_Final/Tournament_Round_43/showa/paper.tex
% Edited source: /workspace/scratch/81914cec2c78/work/audit_42_50/cleaned/r43_showa.tex
\begingroup

\Needspace{9\baselineskip}
\section{\texorpdfstring{A mixed cycle bound for two polygonal conjectures of Liu}{A mixed cycle bound for two polygonal conjectures of Liu}}\label{section:r43showa}
\noindent{\small\sffamily\color{muted}Round 43 \enspace / \enspace Showa}\par\medskip


\begin{sectionabstract}
We prove the mixed square-root inequalities of Conjectures 3.1 and 3.5 in Jian Liu's 2011 RGMIA report. A weighted cycle Rayleigh bound and a weighted vertex-degree bound admit reciprocal rescalings that agree on the diagonal. Their geometric interpolation yields the required mixed term on the original polygon, as well as on its possibly nonsimple loop of doubled circumcenters.
\end{sectionabstract}
\subsection{Statements and scope}
Let $A_1,\ldots,A_n$ form an oriented counterclockwise simple polygon around an interior point $P$, $n\ge3$. Set $R_i=|PA_i|$, $r_i$ the distance from $P$ to the line through $A_i,A_{i+1}$, $\rho_i$ the circumradius of $PA_iA_{i+1}$, and $c=\cos(\pi/n)$. All indices are cyclic. Liu's numbered conjectures [1, (3.2), (3.15)] are
\begin{align}
 \sum_i\sqrt{(r_i+r_{i+1})R_{i+1}}&\ge\sqrt{2/c}\sum_i r_i,\label{r43showa:T1}\\
 \sum_i\sqrt{(R_i+R_{i+1})\rho_i}&\ge c^{-1/2}\sum_i R_i.\label{r43showa:T2}
\end{align}
We prove both, and the replacement of $r_i$ in \eqref{r43showa:T1} by the internal angle-bisector length $w_i$ suggested in [1, Remark 3.2]. For \eqref{r43showa:T2} we require the circumtriangles to be noncollinear so their scalar circumradii exist, an automatic condition for convex polygons and interior points. The source writes circumcenters after (3.15), while the triangle formula (3.14) and scalar units identify circumradii. The numbered Conjecture 3.5 specifies interior points. We do not assert the author's separate subsequent suggestion about every point in the plane.
\subsection{The mixed bound}
\begin{theorem}
Let $L_i>0$, $|\theta_i|<\pi$, $\sum_i\theta_i=2\pi$ and
$0\le h_i\le\sqrt{L_iL_{i+1}}\cos(\theta_i/2)$. Then
\[
 \sum_i\sqrt{L_i(h_{i-1}+h_i)}\ge\sqrt{2/c}\sum_i h_i.
\]
\end{theorem}
\begin{proof}
The weighted cyclic cosine bound is
\[
 \sum_j s_js_{j+1}\cos(\theta_j/2)\le c\sum_j s_j^2.
\]
Indeed form $z_j=s_j\exp(\tfrac i2\sum_{l<j}\theta_l)$, with $z_{n+1}=-z_1$. The left side is the real part of $\sum\bar z_jz_{j+1}$. The eigenvalues of the cyclic shift with negative wrap are the $n$th roots of $-1$; by Fourier diagonalization its Hermitian part has largest eigenvalue $c$.

For arbitrary $d_i>0$ apply this bound to $s_i=d_i\sqrt{L_i}$, obtaining $\sum_i h_i d_id_{i+1}\le c\sum_i L_i d_i^2$. In the opposite scaling, arithmetic-geometric means give
\[
 \sum_i\frac{h_i}{d_i d_{i+1}}\le
 \frac12\sum_i\frac{\lambda_i}{d_i^2},
 \qquad\lambda_i=h_{i-1}+h_i.
\]
By Cauchy--Schwarz the product of the two left sides is at least $(\sum_i h_i)^2$. Put $d_i^2=\sqrt{(\lambda_i+\varepsilon)/L_i}$. The product of the right sides tends, as $\varepsilon\downarrow0$, to $(c/2)(\sum_i\sqrt{L_i\lambda_i})^2$, including zero degrees. This gives the theorem.
\end{proof}
\subsection{Geometric realizations}
Translate $P$ to the origin. For nonzero vectors of lengths $a,b$ with signed principal direction change $\theta$, let $d$ be the distance between their endpoints. Since $d^2=(a-b)^2+4ab\sin^2(\theta/2)$, their line has distance $ab|\sin\theta|/d\le\sqrt{ab}\cos(\theta/2)$ from the origin. The collinear same-ray case is immediate. For coincident endpoints any line through that point has distance at most $a=b$. The same bound holds for the angle bisector length $2ab\cos(\theta/2)/(a+b)$ by AM--GM. If interpreted instead as distance to the closed segment, the bound persists: when the foot falls outside the segment, the closer endpoint length is, say, $a\le b\cos\theta\le b\cos^2(\theta/2)$.

The signed principal changes from $A_i$ to $A_{i+1}$ sum to $2\pi$ by winding number one; the edges avoid the origin. Applying the mixed theorem with $L_i=R_i$ and $h_i=r_i$ (or $w_i$) proves \eqref{r43showa:T1} and its stated variant. For the bisector in collinear incident triples use the continuous limiting value if desired.

For \eqref{r43showa:T2}, put $e_i=A_i/R_i$ and let $Q_i$ be twice the circumcenter vector of $0,A_i,A_{i+1}$. Equal distances give
\[
 |Q_i|=2\rho_i,\qquad e_i\cdot Q_{i-1}=e_i\cdot Q_i=R_i>0 .
\]
Thus the edge joining $Q_{i-1},Q_i$ is supported on a line of height $R_i$ (with that same line used for coincident endpoints). We verify the necessary signed angle sum rather than supposing the dual polygon is convex. Insert $Q_i$'s direction between $e_i$ and $e_{i+1}$. Both resulting principal angular increments have magnitude less than $\pi/2$ by the projection identities. Their sum is therefore exactly the principal change from $e_i$ to $e_{i+1}$. Regrouping these increments in pairs from $Q_{i-1}$ via $e_i$ to $Q_i$ likewise gives principal changes of magnitude less than $\pi$, with total $2\pi$. We can now use the theorem with vertex lengths $2\rho_i$ and adjacent edge heights $R_i,R_{i+1}$. Division by $\sqrt2$ gives \eqref{r43showa:T2}. A regular polygon centered at $P$ has $r_i=w_i=cR$, $2\rho_i=R/c$ and $R_i=R$, yielding equality in all three statements.
\paragraph{Provenance.}
Liu establishes the triangle cases and his Lemma 2.1 uses a quadratic-form composition principle. The ordinary weighted cycle/cosine estimate and reciprocal-weight Cauchy--Schwarz are tools, not priority claims. Our preceding round's circumradius power argument supplied the signed doubled-center transfer. Its power inequalities alone do not yield the mixed sums proved here. We make no assertion of external peer review or absolute historical priority.
\paragraph{References.}
[1] J. Liu, \emph{A new sharpening of the Erd\H{o}s--Mordell inequality and related inequalities}, RGMIA Research Report Collection 14 (2011), art.~85, esp. Sec.~3.\par
[2] H.-C. Lenhard, \emph{Verallgemeinerung und Versch\"arfung der Erd\H{o}s--Mordellschen Ungleichung f\"ur Polygone}, Arch. Math. 12 (1961), 311--314.

\par\endgroup

% Source: /workspace/scratch/81914cec2c78/work/archive/Tournament_Round_43_Final/Tournament_Round_43/heisei/paper.tex
% Edited source: /workspace/scratch/81914cec2c78/work/audit_42_50/cleaned/r43_heisei.tex
\begingroup
\newcommand{\Var}{\operatorname{Var}}
\Needspace{9\baselineskip}
\section{\texorpdfstring{Parameter log-convexity of the generalized modulus via a quantile contraction}{Parameter log-convexity of the generalized modulus via a quantile contraction}}\label{section:r43heisei}
\noindent{\small\sffamily\color{muted}Round 43 \enspace / \enspace Heisei}\par\medskip


\begin{sectionabstract}
A question of Baricz asks whether the zero-balanced generalized modulus is strictly log-convex in its denominator parameter. We prove this for all $0<a<c$ and $0<r<1$. A change of variable turns the second logarithmic derivative into a difference of variances. A comparison of translated densities makes this difference nonnegative and leaves a strictly positive variance term.
\end{sectionabstract}

\subsection{Question and statement}
Let $B(a,b)=\Gamma(a)\Gamma(b)/\Gamma(a+b)$, $F_z={}_2F_1(a,b;a+b;z)$, and $b=c-a$. The generalized modulus is
\[
 \mu_{a,c}(r)=\frac{B(a,b)}{2}\frac{F_{1-r^2}}{F_{r^2}}.
\]
Baricz \cite[p.~3227, Open Problems (a)]{r43heisei:baricz} asks about its logarithmic curvature as $c$ varies with $a,r$ fixed. This is different from the same paper's Corollary~2.4 about varying $a$ with $c$ fixed.
\begin{theorem}
For every $a>0$, $0<r<1$ and $c>a$, one has $\partial_c^2\log\mu_{a,c}(r)>0$. In fact $b\mapsto\log F_z$ is concave on $(0,\infty)$ for each $a>0$, $0<z<1$.
\end{theorem}

\subsection{The transformed measures}
For $0\le z<1$ write
\[
 h_z(s)=e^{-bs}(1-e^{-s})^{a-1}(1-ze^{-s})^{-a},\qquad
 I_z=\int_0^\infty h_z(s)\,ds .
\]
The substitution $t=e^{-s}$ in Euler's beta integral for $F_z$ yields $I_z=B(a,b)F_z$, $I_0=B(a,b)$. This integral identity also follows by expanding $(1-zt)^{-a}$ and integrating each nonnegative term against the beta weight. Let $X_z$ have density $f_z=h_z/I_z$, distribution function $G_z$ and quantile $Q_z=G_z^{-1}$ on $(0,1)$. All densities are positive and smooth on the interior of their common support $(0,\infty)$. At the origin $h_z$ is asymptotic to $(1-z)^{-a}s^{a-1}$, and at infinity it is asymptotic to $e^{-bs}$. Differentiation under the integral twice in $b$ is justified locally by these bounds, with $b$ bounded below by a positive number. Hence
\begin{equation}\label{r43heisei:var}
 \partial_b^2\log I_z=\Var(X_z).
\end{equation}
We compare these variances by a stronger, pointwise result.

\begin{lemma}\label{r43heisei:disp}
For $0<z<1$, $Q_z'(u)\le Q_0'(u)$ for each $0<u<1$. In particular $\Var(X_z)\le\Var(X_0)$.
\end{lemma}
\begin{proof}
First $f_z(s)/f_0(s)$ is strictly decreasing since its only nonconstant factor is $(1-ze^{-s})^{-a}$. The two positive densities have equal total mass, so their difference changes sign once from positive to negative. Integrating on the left or right of their crossing yields $G_z(s)>G_0(s)$ for each $s>0$. Thus $Q_z(u)<Q_0(u)$.

Fix any $d>0$. The ratio $R(s)=f_z(s)/f_0(s+d)$ has logarithmic derivative
\begin{equation}\label{r43heisei:der}
 (a-1)\left(\frac{1}{e^s-1}-\frac{1}{e^{s+d}-1}\right)
       -\frac{az}{e^s-z}.
\end{equation}
If $a\le1$ this is negative. If $a>1$, put $y=e^s>1$ and $k=e^{-d}\in(0,1)$. The ratio of the positive term in \eqref{r43heisei:der} to the magnitude of the negative term is
\[
 \frac{(a-1)(1-k)}{az}\cdot\frac{y}{y-k}\cdot\frac{y-z}{y-1}.
\]
Both displayed $y$-dependent factors are positive and strictly decreasing for $y>1$. Consequently $R$ is monotone or increases and then decreases. In particular $\{s:R(s)>1\}$ is an interval (possibly empty or unbounded).

Now let $D(s)=G_z(s)-G_0(s+d)$. Then $D(0)=-G_0(d)<0$, $D(\infty)=0$, and $D'(s)>0$ precisely on the interval just described. We claim that at every interior point $D(s)=0$ one has $D'(s)\ge0$. Indeed, at a point with negative derivative before the positive-derivative interval, $D$ has not increased from its negative start. At a point with negative derivative after that interval, $D$ cannot increase thereafter and decreases strictly in a neighborhood, so $D(s)>D(\infty)=0$. If the interval is empty there is no zero. These observations prove the claim.

Apply the claim to $d=Q_0(u)-Q_z(u)>0$ at $s=Q_z(u)$. It gives $f_z(Q_z(u))\ge f_0(Q_0(u))$. The inverse function theorem yields the asserted derivative inequality. Integrating this inequality from $u$ to $v$ gives
\[
 0\le Q_z(v)-Q_z(u)\le Q_0(v)-Q_0(u)\qquad(0<u<v<1).
\]
For independent uniform $U,V$ on $(0,1)$, the expected squared difference of $Q_z(U),Q_z(V)$ is twice the variance of $X_z$. The displayed gap comparison gives the variance conclusion.
\end{proof}

\begin{proof}[Proof of the theorem]
Using \eqref{r43heisei:var}, $\partial_b^2\log F_z=\Var(X_z)-\Var(X_0)\le0$ by Lemma~\ref{r43heisei:disp}. Write $z=r^2$. Replacing $F$ by $I/I_0$ gives
\[
 \partial_c^2\log\mu_{a,c}(r)=\Var(X_{1-z})+\Var(X_0)-\Var(X_z)
     \ \ge\ \Var(X_{1-z})>0 .
\]
Here $\partial_c=\partial_b$ for fixed $a$; strict positivity holds since $X_{1-z}$ has a positive density, not a point mass.
\end{proof}

\begin{thebibliography}{9}
\bibitem{r43heisei:baricz} \'A. Baricz, Tur\'an type inequalities for hypergeometric functions, \emph{Proc. Amer. Math. Soc.} 136 (2008), 3223--3229. Open Problems (a), p.~3227.
\end{thebibliography}

\par\endgroup

% Source: /workspace/scratch/81914cec2c78/work/archive/Tournament_Round_43_Final/Tournament_Round_43/keio/paper.tex
% Edited source: /workspace/scratch/81914cec2c78/work/audit_42_50/cleaned/r43_keio.tex
\begingroup
\newcommand{\cy}{\sum_{\mathrm{cyc}}}
\Needspace{9\baselineskip}
\section{\texorpdfstring{A coefficient certificate for a quadratic inequality involving internal angle bisectors}{A coefficient certificate for a quadratic inequality involving internal angle bisectors}}\label{section:r43keio}
\noindent{\small\sffamily\color{muted}Round 43 \enspace / \enspace Keio}\par\medskip


\begin{sectionabstract}
We give a global algebraic certificate for Conjecture 4 of Liu (2002). An angle-splitting argument eliminates the interior point. A rational minorant for each bisector reduces the three-variable quadratic inequality to four coefficient sign checks in two nonnegative variables.
\end{sectionabstract}
\textbf{Statement.} In any triangle with angles $A_i$ and respective opposite sides $a_i$, let $w_i$ be the internal bisector length at vertex $i$. Let $P$ be any interior point and let $R_i$ and $r_i$ be its distances to vertex $i$ and the opposite sideline. For all real $x_1,x_2,x_3$,
\[
 \cy \left(\frac{w_j+w_k}{r_i}\right)^2 x_i^2\ge
 16\cy \left(\frac{w_i}{R_i}\right)^2 x_jx_k. \tag{1}
\]
Here $(i,j,k)$ runs cyclically. This is \cite[Conjecture 4, (38)]{r43keio:liu}, not either of the two subtriangle questions on the next page.

\textbf{Reduction.} We may replace all $x_i$ by their absolute values. Set $x_i=r_i y_i$. The ray from vertex $i$ to $P$ splits its angle into $\alpha+\beta=A_i$. Since $r_jr_k/R_i^2=\sin\alpha\sin\beta\le\sin^2(A_i/2)$, it suffices to replace this quotient by its upper bound on the right. Indeed equality in this angular reduction holds at the incenter. Let $\Delta$ be the triangle area. Since $w_i=2a_ja_k\cos(A_i/2)/(a_j+a_k)$, we have $w_i\sin(A_i/2)=2\Delta/(a_j+a_k)$.

Let $s$ be the semiperimeter and write $X_i=\sqrt{s-a_i}$, so $a_i=X_j^2+X_k^2$. The half-angle identity and arithmetic--geometric mean yield
\[
 \frac{w_i}{2\Delta}
 =\frac{\sqrt{a_ja_k}}{(a_j+a_k)X_jX_k}
 \ge\frac{2a_ja_k}{(a_j+a_k)^2X_jX_k}.
\]
Put $T_i=(a_j+a_k)^2$, $L_i=X_i a_ja_k$, $U=X_1X_2X_3$, and $D_i=L_j/T_j+L_k/T_k$. On dividing by $(2\Delta)^2$ and using the minorants, it suffices to show
\[
 \cy D_i^2 y_i^2\ge 4U^2\cy y_jy_k/T_i .
\]
Set $E_i=L_jT_k+L_kT_j$ and $H=2U^2T_1T_2T_3$. Then $D_i=E_i/(T_jT_k)$. On putting $z_i=D_i y_i$ the last inequality is the positive semidefiniteness of the symmetric matrix with unit diagonal and with $(j,k)$ entry $-H/(E_jE_k)$. All denominators are positive. Consequently it suffices to check the signs of the three polynomials $P_i$ and the polynomial $P_4$:
\begin{align*}
 P_i &= E_jE_k-H,\\
 P_4 &= (E_1E_2E_3)^2-H^2(E_1^2+E_2^2+E_3^2)-2H^3 .
\end{align*}
The first three signs imply nonnegative principal minors of order two; the final polynomial is the determinant times $(E_1E_2E_3)^2$.

\textbf{Finite certificate.} By symmetry, reorder $X_1\ge X_2\ge X_3>0$. Homogeneity allows us to take $X_3=1$, $X_2=1+u$, $X_1=1+u+v$, with $u,v\ge0$. In this chart each of $P_1,P_2,P_3,P_4$ has nonnegative integer coefficients. For complete reproducibility we give the expansion procedure below. The counts of positive coefficients are $190,184,184,1516$; all other coefficients vanish. Their constant coefficients are $8192,8192,8192,0$, respectively.
\begin{enumerate}
\item Implement the preceding definitions by integer operations. Write each $P_i=F_i-G_i$ with the positive and negative parts exactly as displayed (not splitting individual coefficients). Both parts have nonnegative coefficients; their total degrees are at most $54$ because $L_i,T_i,E_i,H$ have degree at most $5,4,9,18$, respectively.
\item Evaluate the eight components at $u=v=1$. The maximum of $F_i(1,1)+G_i(1,1)$ is $1478598290021923104757392000000$. Hence every signed coefficient is strictly inside $(-B/2,B/2)$ for $B=2^{103}$.
\item For each $i$, set $n=F_i(B,B^{55})-G_i(B,B^{55})$. For $e=0,\ldots,55^2-1$, select the unique signed residue $c\equiv n\pmod B$ in $[-B/2,B/2)$; replace $n$ by $(n-c)/B$. Count positive residues and check there are no negative ones and that the final $n=0$. The coefficient recovered at step $e$ is that of $u^{e\bmod55}v^{\lfloor e/55\rfloor}$ by the degree bound and unique residues. Thus the reported signs certify the entire chart, rather than checking representative triangles.
\end{enumerate}
This proves (1) for the prescribed domain. The exact program below performs every integer operation in the certificate, without external dependencies.\hfill$\square$

\paragraph{Exact computational certificate.}
The following program contains the finite integer or rational calculation used in the proof.
{\small
\begin{Verbatim}
"""Exact bivariate coefficient certificate, with signed Kronecker packing."""
def build(u,v):
    X=(1+u+v,1+u,1)
    a=[X[(i+1)%3]**2+X[(i+2)%3]**2 for i in range(3)]
    T=[(a[(i+1)%3]+a[(i+2)%3])**2 for i in range(3)]
    L=[X[i]*a[(i+1)%3]*a[(i+2)%3] for i in range(3)]
    E=[L[(i+1)%3]*T[(i+2)%3]+L[(i+2)%3]*T[(i+1)%3] for i in range(3)]
    H=2*(X[0]*X[1]*X[2])**2*T[0]*T[1]*T[2]
    p=[(E[(i+1)%3]*E[(i+2)%3],H) for i in range(3)]
    p.append(((E[0]*E[1]*E[2])**2,H**2*sum(t*t for t in E)+2*H**3))
    return p
# All positive and negative component polynomials have degree <=54.
# Sum of absolute values of coefficients bounded by plus(1,1)+minus(1,1).
stride=55
bound=max(x+y for x,y in build(1,1))
bits=bound.bit_length()+2
B=1<<bits
print("stride =",stride,"bits =",bits,"one-norm bound =",bound)
allok=True
for idx,(p,q) in enumerate(build(B,B**stride)):
 n=p-q; coeff={}
 for e in range(stride**2):
  c=n%(B); c=c-B if 2*c>B else c; n=(n-c)//B
  if c: coeff[e%stride,e//stride]=c
 assert n==0 and max(i+j for i,j in coeff)<=54 and max(map(abs,coeff.values()))<=bound
 print(idx,len(coeff),min(coeff.values()),coeff.get((0,0),0),[ij for ij,c in coeff.items() if c<0][:8],flush=True)
 if min(coeff.values())<=0:allok=False
assert allok
\end{Verbatim}

}

\begin{thebibliography}{1}
\bibitem{r43keio:liu} J. Liu, Two results and several conjectures of a kind of geometry inequality (in Chinese), \emph{Journal of East China Jiaotong University} 19(3) (2002), 89--94, p.~93. \href{https://journal.ecjtu.edu.cn/hdjtdx/article/pdf/20020327}{Publisher PDF}.
\end{thebibliography}

\par\endgroup

% Source: /workspace/scratch/81914cec2c78/work/archive/Tournament_Round_44_Final/Tournament_Round_44/taisho/paper.tex
% Edited source: /workspace/scratch/81914cec2c78/work/audit_42_50/cleaned/r44_taisho.tex
\begingroup

\Needspace{9\baselineskip}
\section{\texorpdfstring{A height-two obstruction to algebraic differential equations for poset generating series}{A height-two obstruction to algebraic differential equations for poset generating series}}\label{section:r44taisho}
\noindent{\small\sffamily\color{muted}Round 44 \enspace / \enspace Taisho}\par\medskip


\begin{sectionabstract}
We give a negative answer to the $(3+1)$-free-poset branch of
Alejandro Morales's question on algebraic differential equations.
Arbitrary bipartite incidence matrices give too many height-two
orders for the counting coefficients to satisfy Maillet's necessary
growth condition. This applies to both unlabelled ordinary and labelled
exponential series. The other branch, concerning $(2+2)$-free posets,
is not addressed.
\end{sectionabstract}

\subsection*{Question and precise distinction}
For positive integers $a,b$, the poset $(a+b)$ is the disjoint
union of chains on $a$ and $b$ elements. A poset is $(a+b)$-free if
it has no induced subposet isomorphic to $(a+b)$.
Morales \cite{r44taisho:M} asks whether the counting generating functions of
$(2+2)$-free and $(3+1)$-free posets satisfy algebraic differential
equations (ADEs), after noting that they are not D-finite. The
question is on his research website; the related paper \cite{r44taisho:GMR}
supplies enumeration and asymptotics for $(3+1)$-free posets.
An ADE for a formal series $F\in\mathbb C[[x]]$ means an identity
\[
 P(x,F,F',\ldots,F^{(k)})=0,\qquad
 0\ne P\in\mathbb C[x,y_0,\ldots,y_k].
\]
Non-D-finiteness by itself excludes only equations linear in the
series and its derivatives. The negative answer below excludes all
such polynomials.

\begin{theorem}
Let $\mathcal C$ be any isomorphism-closed class containing all
finite posets of height at most two (height is the number of elements
in a longest chain). Denote by $u_n$ and $\ell_n$ its unlabelled and
labelled counts on $n$ elements. Neither
$\sum_{n\ge0}u_nx^n$ nor $\sum_{n\ge0}\ell_n x^n/n!$
satisfies an algebraic differential equation over $\mathbb C(x)$.
\end{theorem}
\begin{proof}
Fix disjoint labelled sets $A$ and $B$ of sizes $\lfloor n/2\rfloor$
and $\lceil n/2\rceil$. For any subset $S$ of $A\times B$ declare
that the strict comparabilities are precisely $a<b$ for $(a,b)\in S$.
These relations require no further pairs for transitivity: no element
can be both the target and the source of a strict pair. They define
height-at-most-two posets, including the empty matrix case. Distinct
subsets give distinct labelled orders. Each unlabelled order has at
most $n!$ labellings, irrespective of possible ambiguous layer
assignments at isolated elements. Consequently
\[
 \ell_n\ge 2^{\lfloor n^2/4\rfloor},\qquad
 u_n\ge\ell_n/n!\ge 2^{\lfloor n^2/4\rfloor}/n!.
\tag{1}
\]
The Maillet theorem on formal solutions of ordinary algebraic
differential equations \cite{r44taisho:MM} says that the coefficients of any
such solution admit a finite Gevrey bound
$|[x^n]F|\le C R^n(n!)^s$ with $C,R>0$ and finite $s\ge0$.
This assertion concerns formal solutions and does not presuppose
convergence. Rational coefficients in the independent variable may
be cleared to polynomial coefficients. For each series in the theorem
the logarithm of the lower bound in (1) is
$(\log 2)n^2/4-O(n\log n)$. A finite Gevrey bound has logarithm
$O(n\log n)$. The contradiction proves the theorem.
\end{proof}

\subsection*{Consequence and limits}
Because $(3+1)$ has a chain of three elements, every height-at-most-two
poset avoids it. Thus neither of the standard $(3+1)$-free
enumerators is differentially algebraic. More generally the same
conclusion holds when every member of a forbidden family has height
at least three. For example it holds for all $(a+b)$ with
$\max(a,b)\ge3$ and for simultaneous $(a+b),(c+d)$ avoidance when
both maxima are at least three.

By contrast, $(2+2)$ itself has height two and occurs among our
bipartite constructions. The proof therefore makes no claim about
the $(2+2)$ branch, the full classification problem, or a new theorem
on general Gevrey bounds. The consequence is a short application
of classical results, not a new asymptotic enumeration of posets.
\begin{thebibliography}{9}
\bibitem{r44taisho:M} A. H. Morales, \emph{Conjectures}, item on
structure and enumeration of $(3+1)$-free posets, Question~1,
\texttt{sites.google.com/view/ahmorales/research/conjectures}.
\bibitem{r44taisho:GMR} M. Guay-Paquet, A. H. Morales and E. Rowland,
Structure and enumeration of $(3+1)$-free posets,
\emph{Ann. Comb.} \textbf{18} (2014), 645--674; arXiv:1303.3652.
\bibitem{r44taisho:MM} R. Gontsov and I. Goryuchkina,
\emph{The Maillet--Malgrange type theorem for generalized power series},
arXiv:1601.05778. See the classical theorem recalled in the introduction.
\end{thebibliography}

\par\endgroup

% Source: /workspace/scratch/81914cec2c78/work/archive/Tournament_Round_44_Final/Tournament_Round_44/showa/paper.tex
% Edited source: /workspace/scratch/81914cec2c78/work/audit_42_50/cleaned/r44_showa.tex
\begingroup

\Needspace{9\baselineskip}
\section{\texorpdfstring{Tangent distances and side distances in a cyclic polygon}{Tangent distances and side distances in a cyclic polygon}}\label{section:r44showa}
\noindent{\small\sffamily\color{muted}Round 44 \enspace / \enspace Showa}\par\medskip


\begin{sectionabstract}
We give a proof of the cyclic polygon conjecture proposed by Thanh Oai Dao in 2016. The difference of the distance sums is affine on the polygon. At a vertex, each tangent distance is twice a square, while each side distance is twice the product of two successive half-angle sines. The sharp constant is thus the largest eigenvalue of a Dirichlet path averaging operator. No restriction on the position of the circumcenter is needed.
\end{sectionabstract}
\begin{theorem}
Let $A_0,\ldots,A_{n-1}$ be distinct vertices in circular order on a circle, $n\ge3$, and $P$ in their closed convex polygon. Write $d_i$ for the perpendicular distance to the line $A_iA_{i+1}$ and $D_i$ for the perpendicular distance to the tangent line at $A_i$. Indices wrap modulo $n$. Then
\[
 \sum_{i=0}^{n-1} D_i\ \ge\ \sec(\pi/n)\sum_{i=0}^{n-1}d_i .
\]
For strictly interior $P$, equality holds exactly when the polygon is regular. In the regular case equality holds for every $P$ in the closed polygon.
\end{theorem}
This is the side-line formulation of the cyclic/tangential polygon conjecture [1]: the lines of the associated tangential polygon are the tangents at the $A_i$. Perpendicular distance to a line (not minimum distance to a finite segment) is the convention of its previously established triangle instance [2]. The line formulation remains meaningful even for parallel adjacent tangents or when tangent intersections do not enclose the circumdisk.

\paragraph{Proof.}
By translation and scaling use the unit circle. Denote the unit vector at angle $\gamma$ by $u_\gamma$. Each tangent distance on the closed disk is the affine expression $1-P\cdot A_i$. The closed polygon lies in one closed half-plane of each side line, so each $d_i$ also agrees throughout it with its inward signed affine distance. In particular the difference of the two sides of the proposed inequality is affine. It suffices to prove the inequality at every vertex.

Fix a test vertex and rotate and relabel to put it at $P=A_0=(1,0)$. Lift the successive vertex arguments to
\[
 0=\alpha_0<\alpha_1<\cdots<\alpha_{n-1}<\alpha_n=2\pi.
\]
Set $s_i=\sin(\alpha_i/2)$; in particular $s_0=s_n=0$. The tangent distances are $D_i=1-\cos\alpha_i=2s_i^2$. For the side of endpoints at angles $\alpha_i,\alpha_{i+1}$, put $m_i=(\alpha_i+\alpha_{i+1})/2$, $\delta_i=(\alpha_{i+1}-\alpha_i)/2$. The chord equation is $X\cdot u_{m_i}=\cos\delta_i$. Every remaining vertex lies on the complementary arc and has scalar product with $u_{m_i}$ at most $\cos\delta_i$, regardless of the sign of this cosine. It follows that
\[
 d_i=\cos\delta_i-\cos m_i
     =2\sin(\alpha_i/2)\sin(\alpha_{i+1}/2)=2s_i s_{i+1}.
\]
Thus, with $c=\cos(\pi/n)$, the vertex assertion becomes
\[
 \sum_{i=0}^{n-1}s_i s_{i+1}\le c\sum_{i=1}^{n-1}s_i^2.
\]
For completeness take the symmetric operator $(Tf)_i=(f_{i-1}+f_{i+1})/2$ on $i=1,\ldots,n-1$ with zero endpoints. Each vector $f_i^{(k)}=\sin(ik\pi/n)$ is an eigenvector with eigenvalue $\cos(k\pi/n)$, $1\le k<n$. These $n-1$ distinct eigenvalues give a basis; the largest is $c$ and its eigenspace is one-dimensional. The claimed bound is the Rayleigh quotient inequality. This completes the proof at all vertices and hence on their convex hull.

If a nonnegative affine function is zero at an interior point of a two-dimensional convex polygon, it is identically zero on that polygon: a nonconstant affine function decreases along some sufficiently short segment starting from that point. Thus an interior equality forces equality in the path bound at every test vertex. At each such vertex $s_i$ must be a positive scalar multiple of $\sin(i\pi/n)$; in particular $s_1=s_{n-1}$. If the half-arcs on either side of that vertex have lengths $u,v$, this says $\sin u=\sin v$. Since $u,v>0$ and $u+v<\pi$ (there is at least one other positive half-arc), it implies $u=v$. All neighboring arcs therefore have equal length. Conversely, regularity gives $s_i=\sin(i\pi/n)$ in coordinates based at each vertex, so the affine deficit vanishes identically. This proves both the equality statement and sharpness.

\paragraph{Status and prior work.}
The triangle theorem is prior work; the 2016 question expressly asked for the general polygon case. We do not claim novelty of the discrete sine diagonalization, the triangle theorem, or an external priority review. The proof addresses the original planar side-line inequality; no statement about distances to finite side segments is asserted.
\paragraph{References.}
[1] T. O. Dao, \emph{An inequality in cyclic polygon and tangential polygon}, MathOverflow, question 244984, July 24, 2016.\\
{}[2] T. O. Dao, T. D. Nguyen and N. M. Pham, A strengthened version of the Erd\H{o}s--Mordell inequality, \emph{Forum Geom.} 16 (2016), 317--321, MR3556993.\\
{}[3] D. S. Marinescu and M. Monea, About a strengthened version of the Erd\H{o}s--Mordell inequality, \emph{Forum Geom.} 17 (2017), 197--202.\\
{}[4] T. Q. Hung, \emph{A family of weighted Erd\H{o}s--Mordell inequality and applications}, arXiv:2105.07885 (2021), Theorem 2.

\par\endgroup

% Source: /workspace/scratch/81914cec2c78/work/archive/Tournament_Round_44_Final/Tournament_Round_44/reiwa/paper.tex
% Edited source: /workspace/scratch/81914cec2c78/work/audit_42_50/cleaned/r44_reiwa.tex
\begingroup

\Needspace{9\baselineskip}
\section{\texorpdfstring{A sign grid for the secant-zeta unimodularity conjecture}{A sign grid for the secant-zeta unimodularity conjecture}}\label{section:r44reiwa}
\noindent{\small\sffamily\color{muted}Round 44 \enspace / \enspace Reiwa}\par\medskip


\begin{sectionabstract}
We prove the secant-zeta unimodularity conjecture of Lal\'in, Rodrigue and Rogers in every even weight. A meromorphic coefficient identity, justified on circles uniformly separated from its pole lattices, gives an absolutely convergent residue series. The leading residue dominates the tail on an explicit half-angle grid. Alternating signs then account for every root of the polynomial.
\end{sectionabstract}
\subsection{The target and coefficient}
For even $k\ge2$ the conjecture of \cite[Conj.~2]{r44reiwa:LRR} says that all roots of
\[
 P_k(x)=\sum_{m=0}^k {k\choose m}2^m(2^m-2)B_mE_{k-m}(x-x^m)(1+x)^{k-m}
\]
are on the unit circle. Bernoulli and Euler numbers have the signs in \cite{r44reiwa:LRR}, in particular $E_2=-1$. The $m=0$ summand gives degree $k+1$. At $x=1$ the polynomial vanishes. For real $y$ put $z=iy$ and $x=(1-z)/(1+z)$. Up to nonzero factors for finite $y$, its value is that of
\[
 A_k(z)=\sum_{m=0}^k{k\choose m}(2^{m-1}-1)B_m E_{k-m}
 \left((1+z)^{m-1}-(1-z)^{m-1}\right).
\]
In fact $P_k(x)=4x(1+x)^{k-1} A_k(z)$. The Taylor expansions of cosecant and secant in \cite[proof of Thm.~3]{r44reiwa:LRR} show that $[t^k]H_y(t)=-2 i^k A_k(iy)/k!$, where
\[
 H_y(t)=t\sec(t)\bigl(\csc((1+iy)t)-\csc((1-iy)t)\bigr).
\]
At a real pole of $\sec(t)$, the two sine arguments add to an odd multiple of $\pi$, so their sines agree. For $y>0$ the apparent singularity is removable. At a nonzero pole $b=m\pi/(1+iy)$, the residue of $H_y(t)/t^{k+1}$ is
\[
 \frac{(-1)^m}{(1+iy)b^k}\sec(b);
\]
there is a subtracted conjugate contribution at $\bar b$.
\subsection{A uniform alternation}
Write $\theta=\arctan y$, $u=e^{-2i\theta}$, and define the normalized real residue series
\[
 G_k(\theta)=\operatorname{Im}\left[e^{i(k-1)\theta}
 \sum_{m=1}^\infty\frac{(-1)^m}{m^k}
 \sec\left(\frac{m\pi}{2}(1+u)\right)\right].
\]
The series converges uniformly even at the endpoints of $[0,\pi/2]$. Indeed the secant argument is $s+ib$ with $0\le s\le m\pi$, $b^2=s(m\pi-s)$, and $|b|\ge\min(s,m\pi-s)\ge |\sin s|$. Consequently $|\cos(s+ib)|^2=1-\sin^2s+\sinh^2 b\ge1$.

\begin{proposition}
For even $k\ge4$ and $j=0,\ldots,k/2$, let $\theta_j=(j+1/2)\pi/(k+1)$. Then $(-1)^jG_k(\theta_j)>0$.
\end{proposition}
\begin{proof}
Since $e^{i(k+1)\theta_j}=(-1)^j i$, the first summand contributes to $(-1)^j G_k(\theta_j)$ the real part of $u/\sin(\pi u/2)$. Put $a=\pi/2$. The identity $h=\sin(au)/(au)=\int_0^1\cos(tau)\,dt$ gives $\operatorname{Re}h\ge\int_0^1\cos(at)\,dt=1/a$ and $|h|\le\int_0^1\cosh(at)\,dt=\sinh(a)/a$. Thus
\[
 \operatorname{Re}\frac{u}{\sin(au)}=\frac{\operatorname{Re}h}{a|h|^2}
 \ge\frac{1}{\sinh^2(a)}>\frac19.
\]
The other terms have combined modulus at most $\sum_{m\ge2}m^{-4}\le 1/16+1/24=5/48<1/9$. This proves the sign assertion.
\end{proof}
\subsection{Root count and the coefficient identity}
The contour argument below shows that $[t^k]H_y(t)$ equals the negative sum of the residues of $H_y(t)/t^{k+1}$ over all its nonzero poles. It is therefore a fixed nonzero imaginary multiple of $(1+y^2)^{(k-1)/2}G_k(\arctan y)$. The proposition and continuity then give $k/2$ distinct roots with $y>0$ between adjacent grid points (the last point is $\pi/2$). Conjugation and $x=1$ give $k+1$ roots in total. For $k=2$ directly $P_2=(x-1)(x^2+2x/3+1)$.
    
Fix $y>0$ and write $A=1+iy$. The positive moduli of poles of the three reciprocal trigonometric factors lie in two arithmetic progressions: the positive half-odd multiples of $\pi$, and the positive multiples of $\pi/|A|$. (The two cosecants give the same progression.) There is a $y$-dependent bound $C$ for the number of their points in any interval of length three. Choose $0<\delta<\min(1,1/(4C))$. In each sufficiently large interval $[N,N+1]$ the union of the $\delta$-neighborhoods of these moduli has total length at most $2\delta C<1$. Choose $R_N$ from the complement. The reverse triangle inequality shows that every $t$ on $|t|=R_N$ stays at least $\delta$ from every pole. After scaling by $A$ or $\bar A$, the sine arguments stay at least $|A|\delta$ from their zero lattices. A sine or cosine is bounded away from zero uniformly at fixed positive distance from its zeros: use periodicity and compactness in a bounded horizontal strip, and exponential growth outside the strip. Thus $H_y(t)=O_y(R_N)$ on these circles. The contour integral of $H_y(t)/t^{k+1}$ is $O_y(R_N^{1-k})$, which tends to zero for all $k\ge2$. The poles in each progression are simple and the non-real progressions are disjoint because $y>0$. The resulting residue series converges absolutely by the uniform secant bound already proved. The residue theorem now verifies the coefficient identity and hence the root count for every even weight.
\begin{thebibliography}{9}
\bibitem{r44reiwa:LRR} M. Lal\'in, F. Rodrigue and M. Rogers, Secant zeta functions, \emph{J. Math. Anal. Appl.} \textbf{409} (2014), 197--204. DOI: 10.1016/j.jmaa.2013.06.071.
\end{thebibliography}

\par\endgroup

% Source: /workspace/scratch/81914cec2c78/work/archive/Tournament_Round_44_Final/Tournament_Round_44/keio/paper.tex
% Edited source: /workspace/scratch/81914cec2c78/work/audit_42_50/cleaned/r44_keio.tex
\begingroup

\Needspace{9\baselineskip}
\section{\texorpdfstring{A sharp lower bound for alternating Ramanujan sine-cube sums}{A sharp lower bound for alternating Ramanujan sine-cube sums}}\label{section:r44keio}
\noindent{\small\sffamily\color{muted}Round 44 \enspace / \enspace Keio}\par\medskip


\begin{sectionabstract}
We prove the conjecture of Alzer and Kwong for the finite alternating sine-cube sums at frequencies $3^k$. The proof centers the sum at the candidate extremal phase $\pi/6$. For a phase to its left, two terms at neighboring frequencies, or one term and the truncation remainder, already compensate for the loss in the initial summand.
\end{sectionabstract}

\textbf{Theorem.} For every integer $n\ge0$ and every $x\in[0,\pi]$,
\[
 V_n(x):=\sum_{k=0}^n\frac{(-1)^k}{3^k}\sin^3(3^k x)>-\frac38.
\]
The constant $-3/8$ is the infimum over these $n,x$. This is the conjecture in \cite[Section 4, (27)]{r44keio:ak}.

\textbf{Proof.} For $n=0$ the inequality follows from $\sin x\ge0$. Assume that $n\ge1$. Since $3^k$ is odd, $V_n(\pi-x)=V_n(x)$. We may take $0\le x\le\pi/2$. For $x\ge\pi/6$, the first summand is at least $1/8$, and all remaining summands together are at least $-\sum_{k=1}^n3^{-k}=-1/2+3^{-n}/2$. This gives the required strict bound on this part of the interval.

Let $x<\pi/6$ and $\delta=\pi/6-x$. For $k\ge1$, $3^k\pi/6=3^{k-1}\pi/2$ has sine $(-1)^{k+1}$ and cosine zero. Expanding about this angle yields
\[
 V_n(x)+\frac38=\sin^3x-\frac18+
 \underbrace{\frac{3^{-n}}2+\sum_{k=1}^n3^{-k}\bigl(1-\cos^3(3^k\delta)\bigr)}_{B_n(\delta)} .
\]
The derivative $3\sin^2 x\cos x$ is increasing on $[0,\pi/6]$ (its derivative is $3\sin x(2-3\sin^2 x)$). Hence
\[
 \frac18-\sin^3 x\le\frac{3\sqrt3}{8}\delta.
\]

If $3^n\delta<\pi/6$, the remainder term alone gives $B_n(\delta)\ge3^{-n}/2>3\delta/\pi>(3\sqrt3/8)\delta$. Otherwise let $k$ be the first index in $\{1,\ldots,n\}$ for which $t=3^k\delta\ge\pi/6$. Since $\delta\le\pi/6$, we have $\pi/6\le t\le\pi/2$, also when $k=1$. The term indexed by $k$ contributes at least $3^{-k}\sin^2 t$, since $0\le\cos t\le1$. If $k<n$, the following term contributes at least $3^{-k}/3$ because $3t\in[\pi/2,3\pi/2]$ and $\cos(3t)\le0$. If $k=n$, the truncation remainder $3^{-k}/2$ suffices instead. The unused terms are all nonnegative. By the chord inequality $\sin t\ge 2t/\pi$ and the arithmetic--geometric mean inequality,
\[
 B_n(\delta)\ge 3^{-k}\left((2t/\pi)^2+\tfrac13\right)
 \ge\frac{4\delta}{\pi\sqrt3}>\frac{3\sqrt3}{8}\delta ,
\]
where $\pi<32/9$. Thus $V_n(x)>-3/8$ throughout.

Finally the authors' sharpness sequence \cite[Remark 1]{r44keio:ak} follows by setting $x=\pi/6$: then $V_n(\pi/6)=-3/8+3^{-n}/2$. This tends to the asserted lower bound as $n\to\infty$.\hfill$\square$

\begin{thebibliography}{1}
\bibitem{r44keio:ak} H.~Alzer and M.~K.~Kwong, On a sine series of Ramanujan, \emph{Axioms} 15 (2026), article 859, Section 4. \href{https://doi.org/10.3390/axioms15110859}{doi:10.3390/axioms15110859}.
\end{thebibliography}

\par\endgroup

% Source: /workspace/scratch/81914cec2c78/work/archive/Tournament_Round_45_Final/Tournament_Round_45/taisho/paper.tex
% Edited source: /workspace/scratch/81914cec2c78/work/audit_42_50/cleaned/r45_taisho.tex
\begingroup

\Needspace{9\baselineskip}
\section{\texorpdfstring{Two forced zeros of the irreducible-shadow polynomial}{Two forced zeros of the irreducible-shadow polynomial}}\label{section:r45taisho}
\noindent{\small\sffamily\color{muted}Round 45 \enspace / \enspace Taisho}\par\medskip


\begin{sectionabstract}
Han, Li and Reidys conjectured that the polynomial counting
irreducible shadows of genus $g$ by their number of arcs is divisible
by $z^{2g}(1+z)^{2g}$. Their minimum arc count accounts for the
first factor. We prove the second by working at $z=-1$ after rescaling
the genus marker. This exposes a cancellation in the known functional
equation and closes an induction for every positive genus.
\end{sectionabstract}
In the notation of \cite{r45taisho:HLR}, let $I_g(y)=\sum_m i_g(m)y^m$ count
irreducible shadows (reduced arc diagrams with connected crossing
graph) of genus $g$ on one backbone, with $m$ arcs. The inputs from
the original enumeration are that the arc count lies between $2g$
and $6g-2$ and that the two generating functions below satisfy
a formal identity. The conjecture is stated following the computed
polynomials for $1\le g\le8$ in Section~4 of \cite{r45taisho:HLR}.
\begin{theorem}
For each integer $g\ge1$,
\[
 y^{2g}(1+y)^{2g}\mid I_g(y)\qquad\text{in }\mathbb Z[y].
\]
\end{theorem}
\subsection*{The formal inputs}
Write $C(x,t)=\sum_{j\ge0}C_j(x)t^j$ for the generating function of
all matchings by arc count and genus, with $C_0(x)$ the Catalan series.
Theorems 2 and 3 of \cite{r45taisho:HLR} give
\begin{align}
 C-1-xC^2&=\sum_{j\ge1}t^j I_j\left(\frac{xC^2}{1-xC^2}\right), \tag{1}\\
 C_j(x)&=Q_j(x)(1-4x)^{1/2-3j}\qquad(j\ge1), \tag{2}
\end{align}
where each $Q_j$ is a polynomial divisible by $x^{2j}$.
Choose the change of variables from their Corollary~1,
\[
 x=\theta(y)=\frac{y(y+1)}{(2y+1)^2}.
\]
At $y=0$ the square-root branch satisfies
$\sqrt{1-4\theta(y)}=(2y+1)^{-1}$. Hence
\begin{equation}
 C_0(\theta)=\frac{2y+1}{y+1},\qquad
 C_j(\theta)=Q_j(\theta)(2y+1)^{6j-1}\quad(j\ge1).
 \tag{3}
\end{equation}
In particular the coefficients of $t$ in (1) after this substitution
are rational functions of $y$. Identity (1) remains valid in
$\mathbb Q(y)[[t]]$: coefficientwise it holds as germs at zero and
therefore as a rational identity. This point is needed since $C_0$
takes a different local form at $y=-1$.
\subsection*{A genus-weighted local ring}
\begin{proof}
Set $h=y+1$, $R=\mathbb Q[[h]]$ and $K=\mathbb Q((h))$.
Expand the rational identity there, and set $t=T/h^2$ in $K[[T]]$.
By (3) the coefficients $C_j(\theta)$ for $j\ge1$ are divisible by
$h^{2j}$ in $R$: $\theta$ has a simple zero at $h=0$ and $2y+1$ is
a unit. Thus
\[
 \widetilde C=C(\theta,T/h^2)=C_0(\theta)+D,\qquad
 D=\sum_{j\ge1}h^{-2j}C_j(\theta)T^j\in T R[[T]].
\]
The genus-zero term has a simple pole at $h=0$. Define
$B=\theta\widetilde C^2$, $B_0=\theta C_0(\theta)^2=y/h$ and
\[
 E=B-B_0=2\theta C_0(\theta)D+\theta D^2\in TR[[T]].
\]
Indeed $\theta C_0(\theta)=y/(2y+1)$ is regular. The genus-zero
identity $C_0(\theta)-1=B_0$ shows that
$\widetilde C-1-B=D-E$ is regular, despite the separate poles.
Also $1-B_0=1/h$. Therefore $W=B/(1-B)$ has the properties
\[
 1+W=\frac{1}{1-B}=\frac{h}{1-hE}\in hR[[T]],\qquad
 W\in R[[T]],\qquad W|_{T=0}=y.
\tag{4}
\]
After the rescaling (1) now reads
\[
 D-E=\sum_{j\ge1} T^j h^{-2j} I_j(W).
\tag{5}
\]
We induct on $g$. For $1\le j<g$, assuming
$(1+y)^{2j}\mid I_j(y)$ yields
$h^{-2j}I_j(W)\in R[[T]]$ by (4). Taking $[T^g]$ in (5), the left
side and all lower-genus contributions belong to $R$; no $j>g$ occurs.
The only remaining term is $h^{-2g}I_g(y)$, by $W|_{T=0}=y$.
It too belongs to $R$. For $g=1$ the same argument has no
lower-genus terms. Thus the polynomial $I_g$ has a zero of order
at least $2g$ at $y=-1$ in every genus. The known lower bound on
arc count gives $y^{2g}\mid I_g(y)$. The factors are relatively prime
and monic, so their product divides in $\mathbb Z[y]$.
\end{proof}
\paragraph{Checks and scope.}
An accompanying exact computation tabulates matchings by the
Harer--Zagier recurrence and independently extracts $I_g$ from (1)
through genus four. Both zero orders are $2,4,6,8$ in these four cases.
The proof establishes divisibility to at least the stated order, not
exact multiplicity at $-1$ or positivity of the quotients.
\begin{thebibliography}{9}
\bibitem{r45taisho:HLR} H. S. W. Han, T. J. X. Li and C. M. Reidys,
\emph{Combinatorics of $\gamma$-structures}, J. Comput. Biol. (2014),
DOI: 10.1089/cmb.2013.0128; preprint arXiv:1112.4151.
\end{thebibliography}

\par\endgroup

% Source: /workspace/scratch/81914cec2c78/work/archive/Tournament_Round_45_Final/Tournament_Round_45/showa/paper.tex
% Edited source: /workspace/scratch/81914cec2c78/work/audit_42_50/cleaned/r45_showa.tex
\begingroup

\Needspace{9\baselineskip}
\section{\texorpdfstring{An adjacent-edge correction to the polygonal area inequality}{An adjacent-edge correction to the polygonal area inequality}}\label{section:r45showa}
\noindent{\small\sffamily\color{muted}Round 45 \enspace / \enspace Showa}\par\medskip


\begin{sectionabstract}
We prove Dao's 2022 conjecture on the area of an arbitrary polygon and the squared differences of its consecutive side lengths. Rotating neighboring edge vectors through one regular turning angle makes their vector differences vanish on the regular mode. Fourier comparison with the area deficit gives an explicit remainder, including the critical pentagon case. The triangle endpoint reduces to the classical Finsler--Hadwiger inequality.
\end{sectionabstract}
\begin{theorem}
If a simple planar $n$-gon has area $A$ and consecutive side lengths $l_0,\ldots,l_{n-1}$, with indices modulo $n$, then
\[
 4\tan(\pi/n)A\le \sum_j l_j^2-\frac{1}{n-2}\sum_j(l_j-l_{j+1})^2 .
\]
No convexity assumption is required.
\end{theorem}
This is the question posted by T. O. Dao in 2022 [1]. Its correction is over consecutive pairs, not all unordered pairs as in a separate 2018 question. The coefficient is proved as proposed; we do not assert it is the optimal one at every $n$.
\subsection*{A Fourier comparison on the closed edge sequence}
Orient the polygon positively. Let its complex vertices be $z_j$, its edges $e_j=z_{j+1}-z_j$, and set $\alpha=\pi/n$, $\omega=\exp(2i\alpha)$. The following argument in fact allows arbitrary vertex sequences when $A$ denotes signed area. Use Fourier coefficients $b_k$ specified by $z_j=\sum_{k=0}^{n-1}b_k\omega^{jk}$. Then $e_j$ has coefficients $u_k=(\omega^k-1)b_k$ and $u_0=0$. By orthogonality of modes, the shoelace formula gives
\[
 \begin{split}
 2A&=\Im\sum_j\overline{z_j}z_{j+1}
 =n\sum_{k=1}^{n-1}\sin(2k\alpha)|b_k|^2,\\
 4A&=n\sum_{k=1}^{n-1}\cot(k\alpha)|u_k|^2,\qquad
 \sum_j l_j^2=n\sum_{k=1}^{n-1}|u_k|^2.
 \end{split}
\]
Reverse triangle inequality compares the adjacent-length variation to the twisted vector variation
\[
 \sum_j(l_j-l_{j+1})^2\le B:=
 \sum_j|e_j-\omega^{-1}e_{j+1}|^2
 =4n\sum_{k=1}^{n-1}\sin^2((k-1)\alpha)|u_k|^2 .
\]
Set $\Delta=\sum l_j^2-4\tan\alpha\,A$. The $k=1$ terms of both $\Delta$ and $B$ vanish. For $2\le k\le n-1$, the ratio of the coefficients of $|u_k|^2$ in $\Delta$ and $B$ is exactly
\[
 \frac{1-\tan\alpha\cot(k\alpha)}{4\sin^2((k-1)\alpha)}
 =\frac{1}{4\cos\alpha\sin(k\alpha)\sin((k-1)\alpha)} .
\]
The denominator is strictly positive. For $n\ge6$ it is at most $4\cos\alpha<4\le n-2$. For $n=4$, both $k=2,3$ give denominator $2=n-2$. For $n=5$, both sines are at most $\cos(\pi/10)$; consequently the denominator is at most
\[
 4\cos(\pi/5)\cos^2(\pi/10)
 =2\cos(\pi/5)(1+\cos(\pi/5))=\frac{5+3\sqrt5}{4}<3.
\]
Thus $\Delta\ge B/(n-2)$ for every $n\ge4$, proving the assertion in that range. In fact, for $n=5$ the Fourier comparison gives the larger coefficient $4/(5+3\sqrt5)$; this is sharp for the vector comparison over unrestricted edge sequences by choosing the extremal mode, but is not claimed to be sharp for the scalar side-length inequality.

\subsection*{Triangle endpoint}
For $n=3$ the conjecture is the Finsler--Hadwiger inequality. For a self-contained check, write the three side lengths as $y+z,z+x,x+y$ for positive $x,y,z$. The sum of their squares minus the sum of their consecutive squared differences is $4(xy+yz+zx)$. Heron's identity gives $A^2=xyz(x+y+z)$. The required inequality follows on squaring from
\[
 (xy+yz+zx)^2-3xyz(x+y+z)
 =\tfrac12\big((xy-yz)^2+(yz-zx)^2+(zx-xy)^2\big)\ge0 .
\]

\paragraph{What is and is not new.}
The classical polygonal isoperimetric inequality $4n\tan(\pi/n)A\le(\sum l_i)^2$ itself gives $\Delta\ge\sum(l_i-\bar l)^2$. Since $\sum(l_i-l_{i+1})^2\le4\sum(l_i-\bar l)^2$, this already treats $n\ge6$. Thus the finite low-order comparisons, particularly the pentagon, matter to the exact all-order conjecture. Stability estimates by Indrei and Nurbekyan [2] and the discrete spectral/curvature refinements of Kwong [3] are important prior work; neither Fourier diagonalization nor the triangle case is claimed as a new method or discovery. The particular twisted edge comparison is derived here explicitly so that the conclusion does not depend on an unspecified stability constant.

\paragraph{References.}
[1] T. O. Dao, \emph{Stronger conjectured inequality for area of a polygon}, MathOverflow question 426277, July 9, 2022.\\
{}[2] E. Indrei and L. Nurbekyan, \emph{On the stability of the polygonal isoperimetric inequality}, Advances in Mathematics 276 (2015), 62--86.\\
{}[3] K.-K. Kwong, \emph{A unified approach to higher order discrete and smooth isoperimetric inequalities}, Mathematika (2024), DOI: 10.1112/mtk.12238.

\par\endgroup

% Source: /workspace/scratch/81914cec2c78/work/archive/Tournament_Round_45_Final/Tournament_Round_45/keio/paper.tex
% Edited source: /workspace/scratch/81914cec2c78/work/audit_42_50/cleaned/r45_keio.tex
\begingroup
\newcommand{\cy}{\sum_{\rm cyc}}
\Needspace{9\baselineskip}
\section{\texorpdfstring{A median form of Liu's quadratic interior-point inequality}{A median form of Liu's quadratic interior-point inequality}}\label{section:r45keio}
\noindent{\small\sffamily\color{muted}Round 45 \enspace / \enspace Keio}\par\medskip


\begin{sectionabstract}
For an acute triangle we establish the ternary quadratic inequality proposed as Conjecture 5 of Liu (2002). After an angular reduction the squared medians are rational, but their cross products are not. A secant majorant uniform over the acute domain converts the determinant test to a polynomial with nonnegative coefficients in an ordered side chart.
\end{sectionabstract}
\textbf{Theorem.} Let $a_i$ be the sides of an acute triangle, $A_i$ its angles, and $m_i$ the medians. For an interior point $P$ let $R_i$ denote its distance from vertex $i$, and $r_i$ its distance from the opposite sideline. For all real $x_1,x_2,x_3$,
\[
 \cy\left(\frac{m_j+m_k}{r_i}\right)^2 x_i^2
 \ge 16\cy\left(\frac{m_i}{R_i}\right)^2 x_jx_k.
\]
This is equation (39), Conjecture 5 of \cite{r45keio:liu}. Indices $(i,j,k)$ are cyclic.

\textbf{Reduction to sides.} It suffices to treat nonnegative $x_i$. Substitute $x_i=r_i y_i$. If $A_iP$ splits the vertex angle into $\alpha$ and $\beta$, then
$r_jr_k/R_i^2=\sin\alpha\sin\beta\le\sin^2(A_i/2)=S_i$. Use $(m_j+m_k)^2\ge4m_jm_k$. We need only prove positive semidefiniteness of the unit-diagonal symmetric matrix with entries $-q_i$ in positions $(j,k)$, where
\[
 q_i=\frac{2m_iS_i}{\sqrt{m_jm_k}}.
\]
Indeed this is diagonal normalization of
$\sum m_jm_k y_i^2\ge4\sum m_i^2 S_i y_jy_k$. Its determinant is $1-\sum q_i^2-2q_1q_2q_3$. Since each $q_i$ is positive, nonnegative determinant implies $q_i<1$ and hence all the two-by-two principal minors are positive as well.

\textbf{Elimination of the remaining radicals.} Write $M_i=4m_i^2=2\sum a_j^2-3a_i^2$. Acuteness gives $0<a_i^2<(\sum a_j^2)/2$; consequently $\max M_i/\min M_i\le4$. Let $H_i=M_j+M_k$, $D_i=M_j-M_k$. The number $z=(D_i/H_i)^2$ belongs to $[0,9/25]$. On this interval the convex function $(1-z)^{-1/2}$ lies below the chord $1+25z/36$, so
\[
 q_i^2=\frac{4M_iS_i^2}{\sqrt{M_jM_k}}
 \le \frac{8M_iS_i^2}{H_i}\left(1+\frac{25D_i^2}{36H_i^2}\right),
 \qquad 2\prod q_i=16\prod S_i.
\]
Let $N_i=(a_i+a_j-a_k)(a_i+a_k-a_j)$ and $W_i=4a_ja_k$; then $S_i=N_i/W_i$. With $Z_i=W_i^2 H_i^3$, a sufficient condition for the determinant bound, after multiplying by $9\prod Z_i$, is $Q\ge0$, where
\begin{align*}
 Q={}&9\prod Z_i-144\prod N_i\prod W_i\prod H_i^3\\
 &-2\cy M_iN_i^2(36H_i^2+25D_i^2)Z_jZ_k .
\end{align*}

\textbf{Certificate for the entire acute domain.} By symmetry and homogeneity write the sides first as $a\ge b\ge c=1$. Then $a-b<1/(a+b)\le1/(2b)$. Set $b=1+u$ and $2b(a-b)=v/(1+v)$; thus $u,v\ge0$. Multiplying all sides by $2(1+u)(1+v)$ gives the polynomial substitution
\[
 c=2(1+u)(1+v),\qquad b=(1+u)c,\qquad a=b+v.
\]
It is enough to prove $Q\ge0$ on this (slightly enlarged) chart; no median-ratio assertion outside the acute domain is needed.

Expanding the displayed polynomial with these substitutions gives 1733 nonzero coefficients, all positive, the smallest being $4947802324992000000$. Here are exact reproduction details. The bidegree is bounded by $(60,30)$: each side has bidegree at most $(2,1)$ and $Q$ has total degree 30 in the sides. Put $B=2^{640}$, evaluate $q=Q(B,B^{61})$ and extract 1891 consecutive balanced base-$B$ digits, dividing $q$ by $B$ after removing each digit. These recover the bivariate coefficients in exponent order $i+61j$. There is no aliasing: the sums of absolute coefficients of the side polynomials are at most 17; those of $N_i,W_i,M_i$ are each below $2^{12}$, and those of $H_i,D_i$ below $2^{13}$. Submultiplicativity and the displayed formula for $Q$ bound its coefficient norm by $2^{210}$, comfortably less than $B/2$. The standard-library Python verifier below implements exactly these formulas and checks the signs and the final zero remainder. The calculation is a finite integer certificate for the universal polynomial inequality, not shape sampling.\hfill$\square$

\paragraph{Exact computational certificate.}
The following program contains the finite integer or rational calculation used in the proof.
{\small
\begin{Verbatim}
"""Exact proof by coefficient expansion of a rational determinant minorant."""
from math import prod
def certificate(u,v):
 c=2*(1+u)*(1+v);b=(1+u)*c;a=b+v
 A=[a,b,c]
 M=[2*sum(i*i for i in A)-3*i*i for i in A]
 N=[];W=[];H=[];D=[]
 for i in range(3):
  j,k=(i+1)%3,(i+2)%3
  N.append((A[i]+A[k]-A[j])*(A[i]+A[j]-A[k]))
  W.append(4*A[j]*A[k])
  H.append(M[j]+M[k]); D.append(M[j]-M[k])
 Z=[w*w*h**3 for w,h in zip(W,H)]
 Q=9*prod(Z)-144*prod(N)*prod(W)*prod(h**3 for h in H)
 for i in range(3):
  j,k=(i+1)%3,(i+2)%3
  Q-=2*M[i]*N[i]**2*(36*H[i]**2+25*D[i]**2)*Z[j]*Z[k]
 return Q
bits=640;B=1<<bits;stride=61
q=certificate(B,B**stride); C={}
for e in range(61*31):
 c=q%B
 if 2*c>B:c-=B
 q=(q-c)//B
 if c:C[e%stride,e//stride]=c
assert q==0
print('count',len(C),'negatives',sum(x<0 for x in C.values()),'minimum',min(C.values()),'maxbits',max(x.bit_length() for x in C.values()))
print('negative monomial indexes',[k for k,c in C.items() if c<0][:50])
assert all(x>0 for x in C.values())
\end{Verbatim}

}

\begin{thebibliography}{1}
\bibitem{r45keio:liu} J. Liu, Two results and several conjectures of a kind of geometry inequality, \emph{J. East China Jiaotong Univ.} 19(3) (2002), 89--94, p.~93. \href{https://journal.ecjtu.edu.cn/hdjtdx/article/pdf/20020327}{Publisher PDF}.
\end{thebibliography}

\par\endgroup
\clearpage
\part{Rounds 46--50}

% Source: /workspace/scratch/81914cec2c78/work/archive/Tournament_Round_46_Final/Tournament_Round_46/showa/paper.tex
% Edited source: /workspace/scratch/81914cec2c78/work/audit_42_50/cleaned/r46_showa.tex
\begingroup

\Needspace{9\baselineskip}
\section{\texorpdfstring{A centroid correction for Qi's log-concave convolution problem}{A centroid correction for Qi's log-concave convolution problem}}\label{section:r46showa}
\noindent{\small\sffamily\color{muted}Round 46 \enspace / \enspace Showa}\par\medskip


\begin{sectionabstract}
We answer Feng Qi's question about a stronger pointwise lower bound for convolutions of log-concave functions. Beyond the weighted average of the logarithms in the older bound, the correction requires only their values at the simplex centroid. A comparison of radial conditional means with an affine cone profile gives an explicit factor strictly larger than one for every positive centroid gap. Symmetric tent logarithms attain the improved bound in the two-factor case.
\end{sectionabstract}
\paragraph{Statement.}
Fix $n\ge2$, $d=n-1$, and $x>0$. Let $f_j$ be nonnegative locally square-integrable log-concave factors on the positive half-line, as in Qi [1]. Write $V=x^d/d!$ and
\[
 m=\frac{d}{x^d}\int_0^x(x-u)^{d-1}\sum_{j=1}^n\log f_j(u)\,du,\qquad
 M=\sum_{j=1}^n\log f_j(x/n),\quad g=M-m.
\]
In the nontrivial case assume $m$ finite. Then all logarithms are finite and continuous on $(0,x)$: a log-concave factor with finite logarithmic average is positive almost everywhere there, and its positivity set is an interval.
\begin{theorem}
One has $g\ge0$ and
\[
 (f_1*\cdots*f_n)(x)\ \ge\
 V e^M\,d\int_0^1 r^{d-1}\exp\!\left(-\frac{d+1}{d}gr\right)dr
 \ \ge\ V e^m .
\]
The second inequality is strict for $g>0$.
\end{theorem}
The last expression is exactly the previous convolution lower bound in [1, Lemma 3.4, (3.3)]. This is a pointwise strengthening on the full finite-log domain, not just a comparison of norms. For positive finite log-concave factors the mean may still be $-\infty$ due to endpoint behavior. In that case set the improved lower bound to zero, recovering the older bound. The same zero convention works if zero values are admitted under extended-log concavity and $m=-\infty$; no subtraction of infinities is needed. A finite concave logarithm on $(0,x)$ admits a global affine upper support there, so a logarithmic divergence can only be downward. The convolution at $x=0$ is zero under the usual half-line integral convention for $n\ge2$.

\paragraph{Proof.}
Let $U$ be uniform on $K=\{u_j\ge0:\sum_j u_j=x\}$ in its affine span. In coordinates $u_1,\ldots,u_{n-1}$ its volume is $V$. With $F(u)=\sum_j\log f_j(u_j)$ the iterated half-line convolution is $V\,\mathbb E e^{F(U)}$. Each coordinate has density $d(x-u)^{d-1}/x^d$ on $(0,x)$, by simplex slicing. Thus $\mathbb EF(U)=m$, and the barycenter $b$ has coordinates $x/n$, with $F(b)=M$. Concavity gives $M\ge m$.

Project $U$ radially from $b$ onto the boundary, calling the resulting boundary random variable $Y$, and write $U=b+R(Y-b)$. For every measurable angular set of rays, intersecting its cone with $b+r(K-b)$ scales its volume by $r^d$. Consequently $R$ is independent of $Y$ and has density $d r^{d-1}$, $0<r<1$. From $\mathbb EU=b$ and $\mathbb ER=d/(d+1)$ we obtain $\mathbb E Y=b$. Define
\[
 H(r)=\mathbb E\,F(b+r(Y-b)),\qquad 0\le r<1.
\]
On each fixed inner homothet the function is bounded and continuous. Hence $H$ is finite there, concave in $r$, and at most $M$ by Jensen. Since $H(0)=M$, it is nonincreasing. Its mean under the radial density is $m$. Let $L(r)=M-(d+1)gr/d$. This is nonincreasing too and has the same radial mean. The concavity of $H-L$ and equality at $r=0$ imply $(H(r)-L(r))/r$ is nonincreasing on $(0,1)$. Unless $H=L$ identically, equality of the weighted means therefore forces a crossing: $H\ge L$ before a point $r_0$, and $H\le L$ after, with common value $h_0$ there. For $r<r_0$ both values are at least $h_0$, so an exponential secant slope is at least $e^{h_0}$. For $r>r_0$ the slope is at most $e^{h_0}$ but $H-L\le0$. In both cases
\[
 e^{H(r)}-e^{L(r)}\ge e^{h_0}(H(r)-L(r)).
\]
Integration and conditional Jensen now yield $\mathbb E e^{F(U)}\ge \mathbb E e^{H(R)}\ge\mathbb E e^{L(R)}$. Finally $\mathbb E e^{L(R)}\ge e^m$ by Jensen, strictly if $L$ is nonconstant.

\paragraph{Strict gain and exact example.}
For every $n$, take $f_j(t)=e^{-t^2}$. The simplex second moment is $\mathbb E U_j^2=2x^2/[n(n+1)]$, giving $g=x^2(n-1)/[n(n+1)]>0$. For $n=2$ and fixed $x$, take both factors $f_j(t)=\exp(-a|t-x/2|)$, $a>0$. Then $F$ on the segment is $-axR$; the angular Jensen and the affine-profile comparison are equalities. The old normalized bound $e^{-ax/2}$ is strictly improved to the exact value $(1-e^{-ax})/(ax)$. A uniform factor strictly greater than one on all data would be impossible, since equal constants already give equality in the old formula. No claim of optimality over every class of separable factors in higher dimensions is needed or made.
\paragraph{Reference and scope.} [1] F. Qi, \emph{Completely monotonic degree of a function involving trigamma and tetragamma functions}, AIMS Mathematics 5 (2020), 3391--3407, DOI: 10.3934/math.2020219. The present result addresses its first open problem, Remark 7.3, not the completely monotonic degree question in Remark 7.4. An earlier version is arXiv:1301.0154. Jensen's inequality and cone-coordinate volume scaling are standard inputs; the proposed correction uses log-concavity specifically in the one-crossing comparison. No formal verification or external priority guarantee is asserted.

\par\endgroup

% Source: /workspace/scratch/81914cec2c78/work/archive/Tournament_Round_46_Final/Tournament_Round_46/keio/paper.tex
% Edited source: /workspace/scratch/81914cec2c78/work/audit_42_50/cleaned/r46_keio.tex
\begingroup
\newcommand{\cy}{\sum_{\mathrm{cyc}}}
\Needspace{9\baselineskip}
\section{\texorpdfstring{A quadratic inequality involving medians and exradii at an interior point}{A quadratic inequality involving medians and exradii at an interior point}}\label{section:r46keio}
\noindent{\small\sffamily\color{muted}Round 46 \enspace / \enspace Keio}\par\medskip


\begin{sectionabstract}
We establish Conjecture 6 of Liu (2002) on all acute triangles. We replace the mixed term in the median-plus-exradius square by two rational estimates, selecting the stronger type of estimate according to shape. Two explicit coefficient certificates then guarantee positive semidefiniteness of the reduced quadratic form.
\end{sectionabstract}
\textbf{Theorem.} Denote the sides, opposite angles, medians and opposite exradii of an acute triangle by $a_i,A_i,m_i,\rho_i$, respectively. For an interior point $P$, let $R_i$ denote its distance to vertex $i$ and $r_i$ its distance to the opposite sideline. Then for every real $x_1,x_2,x_3$,
\[
 \cy\left(\frac{m_i+\rho_i}{r_i}\right)^2 x_i^2
 \ge16\cy\left(\frac{m_i}{R_i}\right)^2 x_jx_k.
\]
The assertion is Conjecture 6, (40), in \cite{r46keio:liu}.

\textbf{Reduction.} Use absolute weights, since the off-diagonal coefficients on the right are positive. Substituting $x_i=r_i y_i$, one has $r_jr_k/R_i^2\le \sin^2(A_i/2)$: the ratio is the product of the sines of the two parts of the vertex angle cut out by $P$. Equality in this reduction is possible at the incenter. It remains to prove that the matrix with diagonal $(m_i+\rho_i)^2$ and $(j,k)$ entry $-8m_i^2\sin^2(A_i/2)$ is positive semidefinite.

Write $a=a_1\ge b=a_2\ge c=a_3$. Put $S=a+b+c$, $s=S/2$, $L_i=S-2a_i$, $X_i=L_i/2$, $M_i=2\sum a_j^2-3a_i^2=4m_i^2$ and $K=\Delta^2=S\prod L_i/16$ for the area $\Delta$. Then $\rho_i=\Delta/X_i$ and $\sin^2(A_i/2)=L_jL_k/(4a_ja_k)$.
We lower-bound the diagonal by
\[
 m_i^2+K/X_i^2+KF_i/X_i,\qquad
 F_i\le 2m_i/\Delta.
\]
For general shapes use $F_i=4/a_i$, since the median is at least the altitude. On the thin chart $b\ge2c$ instead put $F_1=F_2=2a/(bc)$ and retain $F_3=4/c$. Indeed, for an acute triangle
$4m_i^2\ge4m_1^2=2(b^2+c^2)-a^2\ge a^2$, while $\Delta\le bc/2$.

Make a diagonal congruence by $\mathrm{diag}(a_iL_i)$ and multiply by $4bc$. The resulting minorant matrix has diagonal entries
\[
 d_i=bc\bigl[a_i^2 L_i^2 M_i+a_i^2 L_i S L_j L_k
        +a_i^2 L_i^2 S L_j L_k F_i/2\bigr]
\]
and off-diagonal entries $-e_i$ in positions $(j,k)$, where
$e_i=2bc M_iL_j^2L_k^2$. These expressions are polynomials after inserting the appropriate $F_i$. All entries $d_i,e_i$ are positive on the original domain. It suffices to prove
\[
 T=d_1d_2d_3-\sum d_i e_i^2-2e_1e_2e_3\ge0.
\]
Indeed $T/(d_1d_2d_3)=1-\sum h_i^2-2h_1h_2h_3$, where $h_i=e_i/\sqrt{d_jd_k}>0$; it follows that every $h_i<1$, and all remaining principal minors are positive.

\textbf{Two polynomial charts.} Normalize $c=1$. In the thick chart $1\le b<2$, set $p=1+2u$, $q=1+u$, $b=p/q$. In the thin chart $2\le b$, set $p=2+u$, $q=1$, $b=p/q$. These parametrizations use $u\ge0$. Since
$a-b<1/(a+b)\le1/(2b)$ by acuteness, write $2b(a-b)=v/(1+v)$ with $v\ge0$. Clearing the common scale gives the three sides
\[
 c=2pq(1+v),\qquad b=2p^2(1+v),\qquad a=b+q^2v.
\]
Now $T$ has only nonnegative coefficients in each chart. The counts of its nonzero coefficients are 1222 and 1000 in the thick and thin charts, respectively. Their respective minima are $8153726976$ and $2415919104$.

Here is an exact reproduction of the coefficient test. Each cleared matrix entry is homogeneous of degree eight in the sides, so $T$ has degree 24. Each side has degree at most two in $u$ and one in $v$; hence the bidegree of $T$ is at most $(48,24)$. Use $B=2^{540}$ and evaluate $n=T(B,B^{49})$ for each chart. For $0\le j\le24$ and $0\le i\le48$, in order of increasing $49j+i$, remove the balanced residue $t$ of $n$ modulo $B$, then replace $n$ by $(n-t)/B$. Verify all digits are nonnegative and that the final quotient vanishes.

To see that these are the actual coefficients with no carries, note that the coefficient $1$-norm of each side polynomial is at most 40, that of $S$ is at most 120, that of $L_i$ at most 200, and that of $M_i$ is at most $9\cdot40^2$. Using the displayed formulas and submultiplicativity, each $d_i$ and $e_i$ has norm below $2^{70}$ (cancel $F_i$'s denominator first), and $T$ has norm below $2^{213}$. Thus each signed coefficient is strictly less than $B/2$ in absolute value. The Python verifier below evaluates the formulas in integer arithmetic, with no external packages. This computation proves the signs on the entire charts and completes the proof.\hfill$\square$

\paragraph{Exact computational certificate.}
The following program contains the finite integer or rational calculation used in the proof.
{\small
\begin{Verbatim}
"""Exact coefficient audit for the two shape charts; uses no dependencies."""
from math import prod
def compute(u,v,thin):
 p,q=(2+u,1) if thin else (1+2*u,1+u)
 c=2*p*q*(1+v); b=2*p*p*(1+v); a=b+q*q*v
 A=[a,b,c]; S=sum(A); X=[S-2*h for h in A]
 M=[2*sum(h*h for h in A)-3*h*h for h in A]
 diag=[];off=[]
 for i in range(3):
  j,k=(i+1)%3,(i+2)%3
  U=A[i];x=X[i];y=X[j];z=X[k]
  base=U*U*x*x*M[i]+U*U*x*S*y*z
  if thin and i<2:
   d=b*c*base+U*U*x*x*S*y*z*a
  else:
   d=b*c*(base+2*U*x*x*S*y*z)
  diag.append(d)
  off.append(2*b*c*M[i]*y*y*z*z)
 return prod(diag)-sum(diag[i]*off[i]**2 for i in range(3))-2*prod(off)
def verify(thin):
 bits=540; B=1<<bits;stride=49
 n=compute(B,B**stride,thin); C=[]
 for e in range(stride*25):
  c=n%B
  if 2*c>B:c-=B
  n=(n-c)//B
  if c:C.append(c)
 assert n==0
 print('thin' if thin else 'thick',len(C),sum(t<0 for t in C),min(C),max(t.bit_length() for t in C),flush=True)
 assert min(C)>0
for thin in [False,True]:verify(thin)
\end{Verbatim}

}

\begin{thebibliography}{1}
\bibitem{r46keio:liu} J. Liu, Two results and several conjectures of a kind of geometry inequality, \emph{Journal of East China Jiaotong University} 19(3) (2002), 89--94, Conjecture 6, p.~94. \href{https://journal.ecjtu.edu.cn/hdjtdx/article/pdf/20020327}{Publisher PDF}.
\end{thebibliography}

\par\endgroup

% Source: /workspace/scratch/81914cec2c78/work/archive/Tournament_Round_47_Final/Tournament_Round_47/reiwa/paper.tex
% Edited source: /workspace/scratch/81914cec2c78/work/audit_42_50/cleaned/r47_reiwa.tex
\begingroup
\newcommand{\Q}{\mathbb Q}
\newcommand{\Z}{\mathbb Z}
\newcommand{\Ph}{\Phi}
\newcommand{\z}{\zeta}
\Needspace{9\baselineskip}
\section{\texorpdfstring{The complete cyclotomic part of the principal Fekete polynomial at an odd semiprime}{The complete cyclotomic part of the principal Fekete polynomial at an odd semiprime}}\label{section:r47reiwa}
\noindent{\small\sffamily\color{muted}Round 47 \enspace / \enspace Reiwa}\par\medskip


\begin{sectionabstract}
For distinct odd primes \(p>q\), we prove that the cyclotomic factors of the polynomial enumerating the totatives of \(pq\) have exactly the orders proposed by Chidambaram, Min\'a\v{c}, Nguyen and T\^an. Their proof of exhaustiveness assumes an equidistribution conjecture. Instead, we apply Okada's cotangent independence theorem at arbitrary composite orders and use the Galois automorphisms that fix the smaller cyclotomic fields. The only exceptional stabilizer, of order two at the prime three, is ruled out by the cotangent distribution formula.
\end{sectionabstract}

\begin{theorem}\label{r47reiwa:T}
Let \(p>q\) be odd primes, \(n=pq\), and let \(F_n(x)=\sum_{1\leq a\leq n,\, (a,n)=1}x^a\).
Then \(\Ph_d\mid F_n\) if and only if \(d>1\) and at least one of the following holds:
\[
 d\mid q-1,\qquad d\mid p-1\ \hbox{and }d\neq q,\qquad
 d\mid\gcd(pq+1,p+q).
\]
\end{theorem}
This is the unconditional completeness statement proposed in \cite[Prop.~3.1, Thm.~3.2]{r47reiwa:CMNT}. The sufficient conditions were already proved there. That paper also gives the multiplicities: since \(n\) is odd, all these factors are simple \cite[Thm.~2.29]{r47reiwa:CMNT}. The present result concerns the list of distinct factors, not the irreducibility of the quotient after removing them and the monomial \(x\).

\paragraph{Reduction to an odd basis.}
Write \(\z_d=e^{2\pi i/d}\). For \(d=1\), \(F_n(1)=\varphi(n)>0\). If \(d>1\) divides \(n\), every unit class modulo \(d\) has \(\varphi(n)/\varphi(d)\) lifts among the totatives of \(n\). The sum of the primitive \(d\)-th roots is \(\mu(d)\) (by inclusion-exclusion); it is nonzero as \(d\) is squarefree. Thus \(F_n(\z_d)\ne0\), in agreement with the list. If \(d=2\), totatives pair as \(a,n-a\) with opposite parity. Hence \(F_n(-1)=0\), as predicted by \(2\mid q-1\).

Assume henceforth that \(d>2\) and \(d\nmid n\). No divisor of \(n\) is then divisible by \(d\). For \(d\nmid a\) define the symbol
\[
 A_d(a)=\frac{1+\z_d^a}{1-\z_d^a}=i\cot(\pi a/d).
\]
It is odd in \(a\), and \(\sigma_u(A_d(a))=A_d(ua)\) for the Galois automorphism \(\sigma_u:\z_d\mapsto\z_d^u\), \(u\in U_d=(\Z/d\Z)^\times\). A nonunit symbol can vanish, e.g.\ at \(a=d/2\); no argument uses its independence. By inclusion-exclusion of multiples of \(p,q\), or by summing periodic coefficients,
\[
 \frac{F_n(x)}{1-x^n}=\sum_{m\mid n}\mu(m)\frac{x^m}{1-x^m}.
\]
At \(x=\z_d\) all denominators are nonzero. Since \(2\z_d^m/(1-\z_d^m)=A_d(m)-1\), the desired divisibility is equivalent to
\begin{equation}\label{r47reiwa:E}
 A_d(1)-A_d(p)-A_d(q)+A_d(pq)=0 .
\end{equation}

Okada's theorem \cite{r47reiwa:O}, at derivative order zero, says that the \(A_d(a)\) for one representative from each pair \(\{\pm a\}\) in \(U_d\) are linearly independent over \(\Q\). Thus for a combination of \emph{unit} symbols \(\sum b_a A_d(a)\), a relation holds exactly when the signed coefficient function \(a\mapsto b_a-b_{-a}\) vanishes identically. Similarly an automorphism fixes such a combination precisely when multiplication by its index preserves this function. All Galois actions are on \(A_d(a)\) in \(\Q(\z_d)\), not on the real cotangent alone.

\paragraph{When \(d\) is coprime to \(pq\).}
If \(p\equiv1\) or \(q\equiv1\bmod d\), \eqref{r47reiwa:E} holds identically. Otherwise, both negative terms have indices different from \(1\), so in the signed coefficient at \(1\) the initial \(+1\) can only be cancelled by \(pq\equiv-1\bmod d\). With this condition the two positive terms cancel by oddness, and independence makes \(A_d(p)+A_d(q)=0\) equivalent to \(q\equiv-p\bmod d\). These are exactly the three listed possibilities in this case.

\paragraph{Exactly one common prime.}
Let \(r\in\{p,q\}\) divide \(d\), let \(s\) be the other prime, and put \(e=d/r>1\).
As \(A_d(ra)=A_e(a)\), equation \eqref{r47reiwa:E} reads
\begin{equation}\label{r47reiwa:R}
 A_d(1)-A_d(s)=A_e(1)-A_e(s).
\end{equation}
Let \(H=\ker(U_d\to U_e)\). Every element of \(H\) fixes the right side. The reduction map on units is surjective, and \(|H|=\varphi(d)/\varphi(e)\), which is \(r\) if \(r\mid e\) and \(r-1\) otherwise.

If \(s\equiv1\bmod d\), both sides vanish. Suppose otherwise. The signed coefficient function of the left side has positive support \(\{1,-s\}\) when \(s\not\equiv-1\bmod d\), and \(\{1\}\) when \(s\equiv-1\bmod d\). Every element of \(H\) must send \(1\) into this support by independence. Since \(r\ge3\), \(|H|\ge2\), so the singleton case is impossible. The sole remaining possibility is \(r=3\), \(3\nmid e\), and \(H=\{1,-s\}\). Thus \(s\equiv-1\bmod e\), and \(s\equiv1\bmod3\) since the nonidentity element of \(H\) projects to \(-1\) modulo three. If \(e=2\) these congruences give \(s\equiv1\bmod6\), already excluded.

For \(e>2\), among the three residues \(a\bmod 3e\) lying above \(1\bmod e\), the two unit residues are exactly \(H\); the third is \(3b\), where \(3b\equiv1\bmod e\). The cotangent distribution formula gives
\[
 \sum_{j=0}^{2}A_{3e}(1+je)=3A_e(1).
\]
Indeed \(\sum_{j=0}^2\cot(t+j\pi/3)=3\cot(3t)\) follows by logarithmically differentiating the triple-angle product for sine. Consequently the left side of \eqref{r47reiwa:R} equals \(3A_e(1)-A_e(b)\), whereas the right side equals \(2A_e(1)\). This gives \(A_e(b)=A_e(1)\). The fractional linear expression defining \(A_e\) is injective in its argument \(\z_e^b\), so \(b\equiv1\bmod e\). But \(3b\equiv1\bmod e\) then gives \(e\mid2\), a contradiction. Hence \eqref{r47reiwa:R} holds exactly for \(d\mid s-1\). In this case \(d\ne r\) because \(d\nmid pq\); if \(r=p\), \(d\ge p>q-1\) makes this condition vacuous. This agrees with the theorem's list.

\paragraph{Both primes divide \(d\).}
Put \(H_r=\ker(U_d\to U_{d/r})\) for \(r=p,q\). By the Chinese remainder theorem choose \(h\in H_p\) whose \(p\)-primary component modulo \(p^v\), \(p^v\Vert d\), is different from both \(1,-1\). This is possible because if \(v=1\), \(p\ge5\) gives at least four units, and if \(v>1\), the class \(1+p^{v-1}\) works. Take \(k\ne1\) in \(H_q\), possible since \(|H_q|\ge q-1\). The \(p\)-primary component of \(k\) is \(1\), so \(hk\not\equiv-1\bmod d\).

Apply \((1-\sigma_h)(1-\sigma_k)\) to \eqref{r47reiwa:E}. A symbol whose index is divisible by \(r\) lies in \(\Q(\z_{d/r})\), hence is fixed by \(H_r\). The group is abelian, so all three nonunit terms disappear. What remains is a relation with unit indices:
\[
 A_d(1)-A_d(h)-A_d(k)+A_d(hk)=0.
\]
But \(h,k\ne1\) and \(hk\ne-1\), so the positive coefficient at \(1\) cannot be cancelled in the signed coefficient function. This contradiction eliminates every order divisible by both primes. No such order belongs to the theorem's list. The cases exhaust all \(d\). \hfill\(\square\)

\paragraph{Scope and verification.}
The general rational-necklace framework of \cite{r47reiwa:M} explains families of sufficient cancellation conditions. Our necessity argument here is specific to two odd prime factors. It makes no assertion of equidistribution for arbitrary squarefree \(n\). As a bounded independent cross-check, integer polynomial division matches the three conditions for all \(3\le q<p\le31\) prime and \(1\le d\le144\) (6,480 tests). These computations play no role in the proof.
\begin{thebibliography}{3}
\bibitem{r47reiwa:CMNT} S. Chidambaram, J. Min\'a\v{c}, T. T. Nguyen and N. D. T\^an, Fekete polynomials of principal Dirichlet characters, \emph{J. Experimental Math.} \textbf{1} (2025), 51--93; arXiv:2307.14896.
\bibitem{r47reiwa:M} N. C. Minh, N. V. H. Minh, D. Nguyen and T. T. Nguyen, Cyclotomic factors of rational necklace functions, arXiv:2606.02324 (2026).
\bibitem{r47reiwa:O} T. Okada, On an extension of a theorem of S. Chowla, \emph{Acta Arith.} \textbf{38} (1981), 341--345.
\end{thebibliography}

\par\endgroup

% Source: /workspace/scratch/81914cec2c78/work/archive/Tournament_Round_48_Final/Tournament_Round_48/showa/paper.tex
% Edited source: /workspace/scratch/81914cec2c78/work/audit_42_50/cleaned/r48_showa.tex
\begingroup

\Needspace{9\baselineskip}
\section{\texorpdfstring{A harmonic comparison for the Lee metric on a cyclic group}{A harmonic comparison for the Lee metric on a cyclic group}}\label{section:r48showa}
\noindent{\small\sffamily\color{muted}Round 48 \enspace / \enspace Showa}\par\medskip


Let $H_n=\sum_{j=1}^n 1/j$ and give the cyclic group of order $n$ its word metric with generators $\{\pm1\}$. Its Lee harmonic sum and the Lagarias function are
\[
 S_n=\sum_{j=0}^{n-1}\frac1{1+\min(j,n-j)},\qquad L(x)=x+e^x\log x .
\]
The following settles the harmonic comparison posed in [1, Eq.~(2)].
\begin{theorem}
For every integer $n\ge7$, one has $H_n+2e^{H_n}\log H_n<L(S_n)$.
\end{theorem}
\begin{proof}
Let $D_n=S_n-H_n$. Collecting terms by distance gives
\[
 S_{2m}=2H_m-1+\frac1{m+1},\qquad S_{2m+1}=2H_{m+1}-1 .
\]
Thus $S_{n+1}-S_n=1/(1+\lfloor(n+1)/2\rfloor)>1/(n+1)$ for $n\ge2$, so $D_n$ increases strictly on that range. For $n\ge9$ we have
\[
 D_n\ge D_9=\frac{1859}{2520}>\frac7{10},
 \qquad e^{7/10}>\sum_{j=0}^3\frac{(7/10)^j}{j!}=\frac{12013}{6000}>2 .
\]
In particular $S_n>H_n>1$ and $e^{D_n}\log S_n>2\log H_n$.

For $n=7,8$ monotonicity gives
\[
 D_n\ge\frac{241}{420}>\frac{57}{100},\quad
 S_n\ge\frac{19}{6},\quad H_n\le\frac{761}{280}<\sum_{j=0}^7\frac1{j!}<e.
\]
The cubic exponential series at $57/100$ gives $e^{D_n}>1.76$. Using the first three positive terms of the series for $2\operatorname{arctanh}t$, with $t=13/25$, gives
\[
 \log(19/6)>2(t+t^3/3+t^5/5)>1.14 .
\]
Consequently $e^{D_n}\log S_n>1.76\cdot1.14>2>2\log H_n$. In both ranges multiply this bound relative to $2\log H_n$ by $e^{H_n}$, then add $S_n>H_n$, to obtain the claim.
\end{proof}
\noindent\textit{Scope.} The usual RH-equivalent Lagarias bound uses $H_n$ on its right, not the larger $S_n$; it is not established here. Neither is an inequality for arbitrary finite groups.
\smallskip
\par\noindent[1] user6671, \emph{A question about Lagarias inequality}, MathOverflow, May 10, 2019, question 331228, Eq.~(2). https://mathoverflow.net/q/331228.
\par\noindent[2] J.~C. Lagarias, \emph{An elementary problem equivalent to the Riemann hypothesis}, Amer. Math. Monthly 109 (2002), 534--543. arXiv:math/0008177.

\par\endgroup

% Source: /workspace/scratch/81914cec2c78/work/archive/Tournament_Round_48_Final/Tournament_Round_48/heisei/paper.tex
% Edited source: /workspace/scratch/81914cec2c78/work/audit_42_50/cleaned/r48_heisei.tex
\begingroup

\Needspace{9\baselineskip}
\section{\texorpdfstring{A reflected-density proof of median skewness for generalized inverse Gaussian laws}{A reflected-density proof of median skewness for generalized inverse Gaussian laws}}\label{section:r48heisei}
\noindent{\small\sffamily\color{muted}Round 48 \enspace / \enspace Heisei}\par\medskip


\begin{sectionabstract}
We answer the question after Proposition 2 of Hunter's note on skewed distributions. The gamma and beta cases are consequences of a previous convex-transform ordering theorem. For the L\'evy case we prove the stronger statement that all inverse gamma laws, and indeed all interior generalized inverse Gaussian laws, have strictly greater right than left tails at equal distances from their median. The proof uses one sign change in an explicit power series.
\end{sectionabstract}

\subsection{Hunter's question and the known cases}
For a continuous distribution function $F$ with unique median $m$, Hunter's right van Zwet condition \cite{r48heisei:hunter} is
\begin{equation}
 F(m-x)+F(m+x)\le1\qquad(x>0).
 \label{r48heisei:eq:main}
\end{equation}
The question following Proposition 2 of \cite{r48heisei:hunter} concerns gamma laws, L\'evy laws, and $\operatorname{Beta}(a,b)$ for $b\ge a\ge1$.

Here is the existing implication for two of these families.
Write $X\le_c Y$ when the increasing quantile map from $X$ to $Y$ is convex. If $X$ has a continuous symmetric distribution, its paired quantiles at levels $q$ and $1-q$ have midpoint equal to its median. Convexity of the map then gives
\[
 Q_Y(q)+Q_Y(1-q)\ge2Q_Y(1/2),\quad 0<q<1/2.
\]
For laws with positive densities on their support intervals this is equivalent to \eqref{r48heisei:eq:main}: when $m-x$ is in the support take $q=F(m-x)$, then $m+x\le Q_Y(1-q)$; when $m-x$ is below the support the inequality is automatic.

Arab, Oliveira and Wiklund \cite{r48heisei:aow}, Theorem 3, place the symmetric $\operatorname{Beta}(a,a)$ below $\operatorname{Beta}(a,b)$ in convex-transform order whenever $b\ge a>0$. Their beta-to-gamma limiting consequence (Proposition 4 in arXiv:2005.08204v4) similarly places $\operatorname{Beta}(a,a)$ below every $\operatorname{Gamma}(a,\theta)$ with $a,\theta>0$. Thus the gamma and beta assertions already follow, even with $a<1$ for beta. We do not claim them as new.

\subsection{Reflection across the mode and the median}
\textbf{Theorem.} Suppose $f(t)=C t^{p-1}\exp(-a t-b/t)$ for $t>0$. Assume either $a,b>0$ and $p\in\mathbb R$, or $a=0$, $b>0$ and $p<0$.
With $m$ its median, the inequality in \eqref{r48heisei:eq:main} is strict at every $x>0$.

\emph{Proof.} Extend $f$ by zero to the nonpositive axis. At a center $u>0$ and a distance $0<x<u$, put $y=x/u$, $z=y^2$, $A=au$ and $B=b/u$. Expanding the log ratio gives
\begin{align*}
 H_u(z):=\frac{1}{2y}\log\frac{f(u+x)}{f(u-x)}
 &= (p-1)\frac{\operatorname{artanh}(y)}{y}-A+\frac{B}{1-z}\\
 &=c_0+\sum_{n\ge1}c_nz^n,\\[-2pt]
 c_0&=p-1-A+B,\qquad c_n=B+\frac{p-1}{2n+1}.
\end{align*}
This series converges for $0<z<1$. The quantity $c_0=u f'(u)/f(u)$ strictly decreases from positive to negative as $u$ goes from zero to infinity, also when $a=0$ since then $p<0$. Let $v$ be its unique zero, the mode.

At $u=v$, every $c_n$ for $n\ge1$ is positive. Indeed, if $p\ge1$ this is immediate, whereas if $p<1$ then $B=A+1-p\ge1-p$. It follows that $f(v+x)>f(v-x)$ for $0<x<v$, and also for $x\ge v$ by the support. Integration yields $F(v)<1/2$, so $m>v$.

At $u=m$ we now have $c_0<0$. For $p\ge1$ every other coefficient is positive. For $p<1$ the sequence $(c_n)_{n\ge1}$ is strictly increasing toward $B>0$. In either case let $N\ge1$ be the first index with $c_N>0$. The terms $c_nz^{n-N}$ with $n<N$ are nondecreasing on $(0,1)$, strictly so for $n=0$, and those with $n>N$ are increasing. Consequently $z^{-N}H_m(z)$ is strictly increasing. It starts negative, and $H_m(z)$ tends to $+\infty$ as $z\uparrow1$: the pole of $B/(1-z)$ dominates the logarithmic term. There is exactly one zero.

Thus $d(x)=f(m+x)-f(m-x)$ changes sign exactly once from negative to positive, remaining positive beyond $x=m$. Its total integral is zero by the definition of the median. Every proper initial integral is therefore strictly negative (before the crossing use the initial integral; afterwards subtract the positive terminal integral). Since
\[
 \int_0^x d(s)\,ds=F(m+x)+F(m-x)-1,
\]
this proves the assertion. \hfill$\square$

\paragraph{L\'evy law.} The inverse gamma law of shape $\alpha>0$ arises with $p=-\alpha$, $a=0$, and arbitrary $b>0$. A L\'evy density with scale $s$ and lower endpoint $\mu$ has, after translating by $\mu$, precisely this form with $\alpha=1/2$ and $b=s/2$. Condition \eqref{r48heisei:eq:main} is unchanged by translation and positive scaling. This completes the three-family question in the source; positive moment skewness alone would not suffice to answer it.

\begin{thebibliography}{9}
\bibitem{r48heisei:hunter} T. Hunter, A note on the median of skewed distributions, \emph{Statistical Papers} \textbf{67} (2026), art. 64. \href{https://doi.org/10.1007/s00362-026-01831-0}{doi:10.1007/s00362-026-01831-0}.
\bibitem{r48heisei:aow} I. Arab, P. E. Oliveira and T. Wiklund, Convex transform order of Beta distributions with some consequences, \emph{Statistica Neerlandica} \textbf{75} (2021), 238--256. \href{https://doi.org/10.1111/stan.12233}{doi:10.1111/stan.12233}.
\end{thebibliography}

\par\endgroup

% Source: /workspace/scratch/81914cec2c78/work/archive/Tournament_Round_48_Final/Tournament_Round_48/keio/paper.tex
% Edited source: /workspace/scratch/81914cec2c78/work/audit_42_50/cleaned/r48_keio.tex
\begingroup
\newcommand{\cy}{\sum_{\rm cyc}}
\Needspace{9\baselineskip}
\section{\texorpdfstring{The exact maximal altitude exponent in Liu's quadratic triangle inequality}{The exact maximal altitude exponent in Liu's quadratic triangle inequality}}\label{section:r48keio}
\noindent{\small\sffamily\color{muted}Round 48 \enspace / \enspace Keio}\par\medskip


\begin{sectionabstract}
The maximal exponent in Liu's altitude inequality (33) is $\log 2/\log(4/3)$. A vanishing normalized tangent length forces this value; a derivative estimate at that limiting shape, together with concavity and exact interval bounds away from the extremizers, proves its attainability.
\end{sectionabstract}
\textbf{Theorem.} For any triangle with altitudes $h_i$, any interior point $P$ at distances $R_i$ from the vertices and $r_i$ from the opposite sidelines, and arbitrary real weights $x_i$, consider
\begin{equation}
 \cy\left(\frac{h_j+h_k}{r_i}\right)^k x_i^2\ge
 4^k\cy\left(\frac{h_i}{R_i}\right)^k x_jx_k.\label{r48keio:alt}
\end{equation}
Among positive exponents the universal range is precisely $0<k\le K$, where
\[
 K=\frac{\log2}{\log(4/3)}=2.40942083965\ldots.
\]
This answers the maximum-exponent question for (33) in \cite[Remark 2]{r48keio:liu}; it goes beyond the affirmative bound $k\le12/5$ requested in Conjecture 1.

\textbf{Exact scalar criterion.} Replacing $x_i$ by $|x_i|$ suffices. Let $a_i$ denote the side lengths, $s$ the semiperimeter, $A_i$ the angles, and $t_i=(s-a_i)/s$. Then $t_i>0$, $\sum t_i=1$; all such triples are realized by $a_i=1-t_i$. Normalize the diagonal by $z_i=x_i((h_j+h_k)/r_i)^{k/2}$. The coefficient of $z_jz_k$ on the right is at most $W_i^k$, with
\[
 W_i=\frac{4h_i\sin(A_i/2)}{\sqrt{(h_i+h_j)(h_i+h_k)}}.
\]
Indeed the ratio $\sqrt{r_jr_k}/R_i$ is at most $\sin(A_i/2)$ (Liu's Lemma 3), with simultaneous equality at the incenter. Now $h_i=2\Delta/a_i$ and $\sin^2(A_i/2)=(s-a_j)(s-a_k)/(a_ja_k)$ show that $W_i^2=16t_jt_k/((1+t_j)(1+t_k))$. Thus $W_i^k=d_jd_k$ where $d_i=(4t_i/(1+t_i))^{k/2}$. The normalized inequality is equivalent to $(\sum d_i z_i)^2\le\sum(2+d_i^2)z_i^2$, which by sharp weighted Cauchy--Schwarz holds for every $z$ precisely when
\[
 \sum_i g_k(t_i)\le1,\qquad
 g_k(t)=\frac{(4t/(1+t))^k}{2+(4t/(1+t))^k}.
\]
This reduction was used in our earlier note for $k=12/5$. Letting $t\to(0,1/2,1/2)$ now gives the necessary condition $2g_k(1/2)\le1$, equivalently $k\le K$. For $k>K$ a nearby nondegenerate isosceles triangle with $P$ the incenter violates the inequality by continuity and the sharpness of Cauchy--Schwarz.

\textbf{Analytic control of the two equality shapes.} Set $g=g_K$. Straight integer exponent comparisons give $12/5<K<29/12<5/2$. We have
\[
 g'(t)=\frac{K g(t)(1-g(t))}{t(1+t)},\qquad
 \operatorname{sgn}g''(t)=\operatorname{sgn}\{K(1-2g(t))-(1+2t)\}.
\]
For $t\ge1/4$, $g(t)>1/5$ since $(4/5)^K>(4/5)^3>1/2$. Consequently $g''<0$ there. If all $t_i\ge1/4$, Jensen gives $\sum g(t_i)\le3g(1/3)=1$.

Otherwise let $x\le y\le z$ denote the three coordinates. Thus $x\le1/4$ and $y\le1/2$. Consider first the corner $0\le x\le1/128$, $y\ge1/4$. There $z\ge y\ge1/4$, so averaging the two larger coordinates gives
\[
 g(x)+g(y)+g(z)\le H(x):=g(x)+2g((1-x)/2).
\]
Since $g(x)\le (4x)^K/2$, for $0<x\le1/128$ one has $g'(x)\le (K/2)4^K x^{K-1}\le40x\le5/16$. Meanwhile, $w=(1-x)/2\in[1/3,1/2]$ yields $g(w)\in[1/3,1/2]$ and $w(1+w)\le3/4$. Thus $g'(w)\ge2(2/9)/(3/4)=16/27$. Therefore $H'(x)<0$ and $H(x)\le H(0)=2g(1/2)=1$ by the defining equality for $K$.

\textbf{Complement, exact finite check.} The remaining domain is contained in the union
\[
 [0,1/128]\times[0,1/4]\quad\text{and}\quad
 [1/128,1/4]\times[0,1/2].
\]
For reproducibility take $N=2^{20}$. At any integer grid argument $t$ with $0\le t\le N$, use $(p,d)=(12,5)$ if $3t\le N$, and $(p,d)=(29,12)$ otherwise. In both cases raising $4t/(N+t)$ to $p/d$ gives an upper bound for raising it to $K$. Find the least integer $w\ge0$ with $w^d(N+t)^p\ge N^d(4t)^p$, and set $U(t)=w/(2N+w)$. This bounds $g(t/N)$ from above.

On a box $[l,h]\times[a,b]$ of scaled coordinates $(Nx,Ny)$, monotonicity gives the rational upper bound $U(h)+U(b)+U(N-l-a)$ for $g(x)+g(y)+g(1-x-y)$. Starting from the two scaled rectangles above, accept if the upper bound is strictly below one, and otherwise bisect each side and test the four children. The integer-only (exact rational) checker below returns 530 accepted boxes, depth at most 7, and a greatest upper bound of $2373141144949586801/2373204015597911675<1$. Each successful check covers a closed rectangle, not merely its grid points. This proves the endpoint across the entire simplex.

Finally, for $0<k<K$, H\"older interpolation with $\theta=k/K$ gives, for $z_i\ge0$,
\[
 \cy W_i^k z_jz_k\le
 (\cy W_i^K z_jz_k)^\theta(\cy z_jz_k)^{1-\theta}\le\sum z_i^2 .
\]
This concludes the proof of the maximum for \eqref{r48keio:alt}.

\paragraph{Exact computational certificate.}
The following program contains the finite integer or rational calculation used in the proof.
{\small
\begin{Verbatim}
from fractions import Fraction as F
from functools import lru_cache
assert 4**12<2**5*3**12 and 4**29>2**12*3**29 # 12/5 < K < 29/12
N=1<<20
@lru_cache(None)
def upper(t):
 # denominator d of upper exponent when larger than 1, lower if smaller
 p,d=(12,5) if 3*t<=N else (29,12)
 rhs=N**d*(4*t)**p
 coef=(N+t)**p
 l,h=-1,8*N
 while h-l>1:
  w=(l+h)//2
  if w**d*coef>=rhs:h=w
  else:l=w
 return F(h,2*N+h)
stack=[(0,N//128,0,N//4,0),(N//128,N//4,0,N//2,0)]
leaves=0;depth=0;bound=F(0)
while stack:
 l,h,a,b,k=stack.pop()
 m=upper(h)+upper(b)+upper(N-l-a)
 if m<1:
  leaves+=1;depth=max(depth,k);bound=max(bound,m);continue
 assert min(h-l,b-a)>=2,(l,h,a,b,m)
 X=(l+h)//2;Y=(a+b)//2
 for xx in [(l,X),(X,h)]:
  for yy in [(a,Y),(Y,b)]:stack.append((*xx,*yy,k+1))
print('exact exponent bounds 12/5 < K < 29/12')
print('accepted boxes',leaves,'maximum depth',depth,'largest rational upper bound',bound)
\end{Verbatim}

}

\begin{thebibliography}{1}
\bibitem{r48keio:liu} J. Liu, Two results and several conjectures of a kind of geometry inequality, \emph{J. East China Jiaotong Univ.} 19(3) (2002), 89--94, p.~93, Remark 2. \href{https://journal.ecjtu.edu.cn/hdjtdx/article/pdf/20020327}{Publisher PDF}.
\end{thebibliography}

\par\endgroup

% Source: /workspace/scratch/81914cec2c78/work/archive/Tournament_Round_49_Final/Tournament_Round_49/meiji/paper.tex
% Edited source: /workspace/scratch/81914cec2c78/work/audit_42_50/cleaned/r49_meiji.tex
\begingroup

\Needspace{9\baselineskip}
\section{\texorpdfstring{Elementary symmetric dominance and Liu's pedal-triangle power conjecture}{Elementary symmetric dominance and Liu's pedal-triangle power conjecture}}\label{section:r49meiji}
\noindent{\small\sffamily\color{muted}Round 49 \enspace / \enspace Meiji}\par\medskip

\begin{sectionabstract}
We prove the concave-power case of a distance inequality proposed by Liu in December 2012. A linear inequality of Liu supplies the first elementary symmetric coefficient. A pedal side-product identity and Gerretsen's inequality control the remaining two. An integral transform then gives all exponents between zero and one simultaneously.
\end{sectionabstract}
\paragraph{Theorem.}
Let $P$ be a strict interior point of a nondegenerate Euclidean triangle $ABC$. Let $U$ and $v$ denote the circumradius and inradius of the pedal triangle of $P$. For every $0<k<1$,
\[
 PA^k+PB^k+PC^k\ge(2U)^k+2(4v)^k .
\]
Equality holds exactly when $ABC$ is equilateral and $P$ is its center. This is Conjecture~3.1 in \cite{r49meiji:liu}, not its separate reciprocal Conjecture~3.2. The range $k\ge1$ is already Theorem~1.1 of that paper.

\paragraph{Elementary symmetric comparison.}
Set $X=(PA,PB,PC)$, $Y=(2U,4v,4v)$, and write $s$ for the pedal semiperimeter. The pedal area is positive: it is one quarter of the original area times $1-OP^2/R^2$, with $O,R$ the original circumcenter and radius. Denote by $e_j$ the elementary symmetric polynomials. Liu's linear case \cite[Thm.~1.1, $k=1$]{r49meiji:liu} gives $e_1(X)\ge e_1(Y)$, with precisely the stated equality case.

The pedal sides have lengths $X_1\sin A,X_2\sin B,X_3\sin C$. For instance, the adjacent projections of $P$ and $A$ belong to the circle with diameter $PA$, giving its opposite chord length $PA\sin A$. The pedal side product equals four times the pedal circumradius times the pedal area, i.e., $4Uvs$. Since $\prod\sin A\le(\sqrt3/2)^3$ by concavity of $\log\sin$, we obtain
\[
 e_3(X)\ge L:=\frac{32Uvs}{3\sqrt3}.
\]
The elementary inequality $s\ge3\sqrt3 v$ follows from $s/v=\sum\cot(\alpha/2)$, summed over the pedal angles, and convexity of $\cot$ on $(0,\pi/2)$. Thus $e_3(X)\ge 32Uv^2=e_3(Y)$.

By AM--GM, $e_2(X)\ge 3e_3(X)^{2/3}\ge 3L^{2/3}$. We claim
\[
 3L^{2/3}\ge16v(U+v)=e_2(Y).
\]
After cubing and dividing by $1024v^2>0$, this is just $U^2s^2\ge4v(U+v)^3$. By the lower Gerretsen inequality on the pedal triangle, $s^2\ge16Uv-5v^2$. Writing $t=U/v\ge2$ by Euler, the conclusion follows from
\[
 (16t-5)t^2-4(t+1)^3=(t-2)(12t^2+7t+2)\ge0.
\]
For completeness, Gerretsen's input follows by putting the sides equal to $b+c,c+a,a+b$ with $a,b,c>0$. Multiplying $s^2-16Uv+5v^2$ by $s$ gives $(a+b+c)^3-4(a+b)(b+c)(c+a)+5abc$, the nonnegative Schur cubic $\sum_{\rm cyc}a(a-b)(a-c)$.

\paragraph{Transfer to powers.}
Comparing the three coefficients of the product gives
\[
 \prod_{i=1}^3(1+hX_i)\ge\prod_{i=1}^3(1+hY_i)\qquad(h>0).
\]
Take logarithms and integrate against $h^{-1-k}\,dh$. For $b>0$,
\[
 \int_0^\infty\!\log(1+hb)\,h^{-1-k}\,dh
 = b^k\int_0^\infty\!\log(1+t)\,t^{-1-k}\,dt.
\]
The last integral is finite and strictly positive, by $\log(1+t)=O(t)$ at zero and $O(\log t)$ at infinity. Division by this common factor proves the theorem. If $ABC,P$ do not give equality in the linear case then the product inequality is strict for every $h>0$, giving strictness after integration. In the equilateral centered case, the distances and the three components of $Y$ all coincide.
\begin{thebibliography}{1}
\bibitem{r49meiji:liu} J. Liu, A pedal triangle inequality with the exponents, \textit{International Journal of Open Problems in Computer Science and Mathematics} 5 (2012), no.~4, 16--24; Theorem 1.1, Eq.~(6), and Conjecture 3.1.
\end{thebibliography}

\par\endgroup

% Source: /workspace/scratch/81914cec2c78/work/archive/Tournament_Round_49_Final/Tournament_Round_49/taisho/paper.tex
% Edited source: /workspace/scratch/81914cec2c78/work/audit_42_50/cleaned/r49_taisho.tex
\begingroup
\newcommand{\pair}[2]{\langle #1,#2\rangle}
\Needspace{9\baselineskip}
\section{\texorpdfstring{The height-to-depth ratio for random leftist trees}{The height-to-depth ratio for random leftist trees}}\label{section:r49taisho}
\noindent{\small\sffamily\color{muted}Round 49 \enspace / \enspace Taisho}\par\medskip


\begin{sectionabstract}
Nogueira computed the asymptotic constants for the expected height and average node depth of a uniform leftist tree, and conjectured in 2001 that the former is twice the latter. We prove this identity. The argument uses a bounded left eigenvector of the critical rank transfer operator to compare an exact quadratic difference at two arguments with the quadratic term of the height iteration. The common quadratic form cancels. All limits and pairings take place in ordinary sequence spaces; no equality is inferred from decimal agreement.
\end{sectionabstract}

\subsection{Statement and source normalizations}
Use Nogueira's size convention \cite{r49taisho:N}: a leftist tree is an ordered binary tree such that at each node the minimum distance to an empty subtree is at least as large for the left subtree as for the right. Rank (the $d$-number) is zero for the empty tree and one plus the minimum child rank otherwise. For the uniform tree of $n$ nodes, write $\xi_D\sqrt{\pi n}$ and $\xi_H\sqrt{\pi n}$ for the leading terms of the expected mean node depth and the expected height. They are specified by Proposition~14 and Theorem~3 of \cite{r49taisho:N}. The conjecture is in the paragraph after (104), p.~275, not a claim that the corresponding finite-size expectations are exactly proportional.
\begin{theorem} $\xi_H=2\xi_D$. \end{theorem}
We use the published asymptotic formulas and their defining equations as inputs; we do not give new proofs of the separate height and depth limit theorems. We call the space of absolutely summable real sequences $\ell^1$, keeping the notation $\ell$ for the counting series.

Let $\ell_k(x)$ be the generating function for rank $k$ trees, $\ell_0=1$ and $\ell=\sum_{k\ge0}\ell_k$. Thus $\ell_1(x)=x\ell(x)$. Write $\rho$ for the common critical argument, set $y(x)=(\ell_k(x))_{k\ge0}$, $y=y(\rho)$, and put
\[
 B^2=\frac{2(\ell_1(\rho)+\rho)}
             {\rho^4\ell_1''(\rho\ell_1(\rho))},\qquad B>0.
 \tag{1}
\]
The square-root singularity, (21),(38) in \cite{r49taisho:N}, gives
$\ell(\rho)-\ell(x)\sim B\sqrt{1-x/\rho}$. To see the normalization explicitly, the inverse of $\ell_1$ has second derivative $-\eta_2$ at its critical point, with
$\eta_2=\rho^3\ell_1''(\rho\ell_1(\rho))/(\rho+\ell_1(\rho))$.
The coefficient of $\sqrt{1-x/\rho}$ in $\ell_1(\rho)-\ell_1(x)$ is $\sqrt{2\rho/\eta_2}$; division by $\rho$ gives $B$.

Denote by $\alpha_k$ the limits $\ell'_k(x)/\ell'(x)$ as $x\uparrow\rho$. Equations (40)--(46) of \cite{r49taisho:N} give $\alpha_0=0$, $\alpha_k\ge0$ and $\sum_k\alpha_k=1$. Let $g_k=g_k(\rho)$ be the weights of (30)--(32) with $\gamma_1=\gamma_2=1$, and put $\sigma=\sigma_2(\rho)$ as in (97),(98). The two \emph{known} formulas (79),(104) then read
\[
 \xi_D=\frac{B}{2g_0}\sum_{k\ge0}g_k\alpha_k,\qquad
 \xi_H=\frac1{B|\sigma|}.
 \tag{2}
\]
In particular $g_0>0$ by Proposition~13 and the depth constant is positive. The missing identity is between $\sigma$ and the weighted sum in (2).

\subsection{The critical left functional}
For $u,v\in\ell^1$, put $Q(u,v)_0=0$ and
\[
 Q(u,v)_k=u_{k-1}\sum_{j\ge k-1}v_j,\quad k\ge1.
\]
This bilinear map has norm at most one. With $e_j$ the $j$th unit vector, rank decomposition gives
\[
 y(x)=e_0+xQ(y(x),y(x)),\qquad
 J(x)u=x\{Q(y(x),u)+Q(u,y(x))\},\qquad J=J(\rho).
\]
The second formula is the linearization. The vector $g$ defines a bounded functional: by (9),(32) of \cite{r49taisho:N}, $g_j(x)=g_0(x\ell_j(x))$; the arguments for $j\ge1$ tend to zero at $\rho$, and $g_0$ is analytic near zero and has finite positive boundary value $g_0(\rho)$ by (74). Each $g_j(x)$ extends continuously from below to $\rho$.

For a fixed $j$ use $\mu=e_j$ and $\beta=(I-J(x))e_j$ in the source's linear system (24)--(30). The support of $\beta$ lies in $0,\ldots,j+1$, directly from the definition of $Q$. Its solution has total $\sum_k\mu_k=1$, so
\[
 \sum_k g_k(x)\beta_k=1-f_0(x).
\]
By (62) of \cite{r49taisho:N}, $1-f_0(x)=g_0(x)/\ell_1'(x)\to0$ at the critical argument. Taking limits in this finite sum, then using density of finitely supported sequences and boundedness of $J$ and $g$, gives
\[
 \pair{g}{(I-J)u}=0\quad (u\in\ell^1).
 \tag{3}
\]
Since $Jy=2\rho Q(y,y)=2(y-e_0)$, (3) also gives $\pair{g}{y-e_0}=g_0$.

\subsection{Projection of two quadratic remainders}
Write $d=y-y(x)$ and $\varepsilon=\sqrt{1-x/\rho}$. Then $d/\varepsilon\to B\alpha$ in $\ell^1$. Indeed, coordinatewise convergence follows by integrating the derivative ratio limit against $\ell'(x)$ and using (1). All coordinates are nonnegative, and their total converges to $B=B\sum_k\alpha_k$. For any cutoff use convergence on its finite head and then convergence of the tail masses; letting the cutoff increase proves norm convergence.

Subtract the two counting equations at $\rho$ and $x$ without dropping any terms:
\[
 (I-J)d=(\rho-x)Q(y(x),y(x))-\rho Q(d,d).
\]
Apply (3), divide by $\varepsilon^2$, and use norm continuity of $Q$. The result is the normalization
\[
 \rho B^2\pair{g}{Q(\alpha,\alpha)}
   =\rho\pair{g}{Q(y,y)}=g_0 .
 \tag{4}
\]

Now put $c_k=\sigma_{k,2}(\rho)$. In (96)--(98) of \cite{r49taisho:N} the first-order height coefficients are $\alpha_k$, while $\sum_kc_k=\sigma$. With these substitutions, the second-order recurrence (97) reads exactly
\[
 c=Jc-\rho Q(\alpha,\alpha)-\sigma\alpha .
 \tag{5}
\]
For clarity, the summability required to use (3) is not implicit. From the convergent scalar series in (98), the tails $\sigma-\sum_{j<k-1}c_j$ are uniformly bounded. Also $\sum_{j\ge k-1}y_j\to0$, and the tails of $\alpha$ are at most one. The coordinate recurrence (97) therefore implies, for $k$ sufficiently large and a fixed $C$,
\[
 |c_k|\le \tfrac12|c_{k-1}|+
 C(y_{k-1}+\alpha_{k-1}+\alpha_k).
\]
Summing this through any upper index and absorbing the half-sum gives a uniform bound on partial sums of $|c_k|$. Thus $c\in\ell^1$ and (5) is an equality in that space.

Apply (3) to (5). By (4),
\[
 -\sigma\pair{g}{\alpha}
    =\rho\pair{g}{Q(\alpha,\alpha)}
    =\frac{g_0}{B^2}.
\]
The signs in (2) show $\pair{g}{\alpha}>0$, hence $\sigma<0$. Substitution in (2) gives $\xi_H=B\pair{g}{\alpha}/g_0=2\xi_D$, as asserted. The factor two comes from the quadratic transfer, not a disagreement in height conventions.

\begin{thebibliography}{9}
\bibitem{r49taisho:N} P.~Nogueira, On the combinatorics of leftist trees, \emph{Discrete Applied Mathematics} {\bf109} (2001), 253--278.
\url{https://doi.org/10.1016/S0166-218X(00)00250-X}.
\end{thebibliography}

\par\endgroup

% Source: /workspace/scratch/81914cec2c78/work/archive/Tournament_Round_49_Final/Tournament_Round_49/reiwa/paper.tex
% Edited source: /workspace/scratch/81914cec2c78/work/audit_42_50/cleaned/r49_reiwa.tex
\begingroup
\newcommand{\B}{\mathcal B}
\Needspace{9\baselineskip}
\section{\texorpdfstring{Altizio's \(2\pi\) lower bound for a geometric sum on the Stern exclusion disk}{Altizio's 2pi lower bound for a geometric sum on the Stern exclusion disk}}\label{section:r49reiwa}
\noindent{\small\sffamily\color{muted}Round 49 \enspace / \enspace Reiwa}\par\medskip


\begin{sectionabstract}
We prove the conjectured lower bound \(\min(n,2\pi)\) for the modulus of a geometric sum on the closed disk of radius one centered at two. A phase estimate of Altizio covers a short initial arc. At the endpoint of that arc, the second binomial term in a radial lower bound pays for the distance remaining to \(2\pi\). Monotonicity then covers the rest of the boundary. Thirty explicit polynomial sign checks handle the small degrees.
\end{sectionabstract}
\begin{theorem}\label{r49reiwa:T}
For every integer \(n\ge1\) and \(z\in\B=\{z:|z-2|\le1\}\),
\[
 |S_n(z)|\ge\min(n,2\pi),\qquad S_n(z)=1+z+\cdots+z^{n-1}.
\]
\end{theorem}
This is Remark 3(2) of Altizio \cite{r49reiwa:A}. His Theorem 8 obtains the lower threshold \(11/2\) in place of \(2\pi\). Here \(S_n\) denotes only the geometric sum, not the recursively defined Stern polynomial. The distinct question in his Remark 4 asks for a bound by \(n|z|\) at particular degrees.

\paragraph{Reduction to the boundary.} The only point of \(\B\) in the closed unit disk is \(1\); \(S_n(1)=n\), and every zero of \(S_n\) is a nontrivial \(n\)th root of unity. The maximum modulus principle applied to \(1/S_n\) thus reduces the problem to \(\partial\B\). By conjugation, parameterize its upper half as \(w(t)=2-e^{-it/n}\), \(0\le t\le\pi n\). Write \(y=|w(t)-1|=2\sin(t/(2n))\); then \(|w(t)|^2=1+2y^2\).

\paragraph{Large indices.} Suppose \(n\ge32\) and put
\[
 M=2\pi\left(1-\frac{2\pi}{n}\right).
\]
For \(0\le t\le M\), Proposition 5 of \cite{r49reiwa:A} gives
\[
 |S_n(w(t))|\ge n\left(1-\frac{t}{2\pi}\right)\ge 2\pi .
\]
For the remaining arc, the reverse triangle inequality gives, for \(y>0\),
\[
 |S_n(w(t))|\ge B_n(y):=\frac{(1+2y^2)^{n/2}-1}{y}.
\]
This lower bound is increasing, since \(B_n(y)=ny\int_0^1(1+2uy^2)^{n/2-1}\,du\).
The parameter \(y\) is also increasing throughout the upper boundary, so it suffices to treat \(t=M\).

Let \(y_0=2\sin(M/(2n))\). Taylor's theorem applied to \((1+v)^{n/2}\) at \(v=0\), with nonnegative third derivative on \(v\ge0\), yields
\[
 B_n(y_0)\ \ge\ ny_0+\tfrac12n(n-2)y_0^3.
\]
Using \(\sin x\ge x-x^3/6\) and \(M<2\pi<7\), we obtain
\[
 ny_0\ge M-\frac{M^3}{24n^2},\qquad y_0\ge\frac{99}{100}\frac{M}{n}.
\]
Consequently,
\[
 B_n(y_0)\ge M+\frac{M^3}{n}
 \left[\frac{1-2/n}{2}\left(\frac{99}{100}\right)^3-\frac{1}{24n}\right]
 \ge M+\frac{9M^3}{20n}.
\]
For the last bound, substitute \(n=32\) into the increasing bracket; its value is \(0.4535\ldots>9/20\). Now \(2\pi/n<1/4\) gives \(M\ge3\pi/2\). Since \((9/20)(3\pi/2)^3>4\pi^2\) by \(\pi>3\), we get \(B_n(y_0)>M+4\pi^2/n=2\pi\). This covers all \(t\ge M\).

\paragraph{Finite sign certificates.} Put \(s=|z-1|^2\in[0,4]\) on \(\partial\B\), so \(z+\bar z=2+s\) and \(r=z\bar z=1+2s\). Define \(T_0=2\), \(T_1=2+s\), \(T_n=(2+s)T_{n-1}-rT_{n-2}\), hence \(T_n=z^n+\bar z^n\). The polynomial
\[
 P_n(s)=\frac{r^n+1-T_n(s)}s-\min(n,44/7)^2
\]
equals \(|S_n(z)|^2-\min(n,44/7)^2\), including at \(s=0\) by continuity. For \(n=1\) it vanishes identically. For \(2\le n\le31\), divide by \(s\) as long as the constant coefficient vanishes. The remaining polynomial has positive values at both endpoints and zero real roots in \([0,4]\), proving \(P_n(s)\ge0\) throughout. Here is the complete finite check, using exact rational polynomial arithmetic in Python/SymPy. The root count uses a real-root counting algorithm, not samples.
\begin{Verbatim}
from sympy import symbols, Poly, QQ, Rational
X = Poly(symbols('s'), symbols('s'), domain=QQ)
one = X**0
r, a, b = 1+2*X, 2*one, 2+X
for n in range(1,32):
    if n > 1: a,b = b,(2+X)*b-r*a
    c = min(Rational(n), Rational(44,7))
    P = (r**n+one-b).exquo(X)-c*c
    if n == 1: assert P.is_zero; continue
    while P.eval(0) == 0: P = P.exquo(X)
    assert P.eval(0)>0 and P.eval(4)>0 and P.count_roots(0,4)==0
\end{Verbatim}

Since \(44/7>2\pi\), this proves the required finite cases and completes Theorem~\ref{r49reiwa:T}. \hfill\(\square\)

\begin{thebibliography}{1}
\bibitem{r49reiwa:A} D. Altizio, \emph{Zeros of Stern polynomials in the complex plane}, arXiv:2511.03847 (2025), Section 4.3, Proposition~5, Theorem~8 and Remark~3(2).
\end{thebibliography}

\par\endgroup

% Source: /workspace/scratch/81914cec2c78/work/archive/Tournament_Round_50_Final/Tournament_Round_50/meiji/paper.tex
% Edited source: /workspace/scratch/81914cec2c78/work/audit_42_50/cleaned/r50_meiji.tex
\begingroup

\Needspace{9\baselineskip}
\section{\texorpdfstring{Reciprocal symmetric bounds for the distances to three vertices}{Reciprocal symmetric bounds for the distances to three vertices}}\label{section:r50meiji}
\noindent{\small\sffamily\color{muted}Round 50 \enspace / \enspace Meiji}\par\medskip

\begin{sectionabstract}
We answer Liu's 2012 Conjecture 3.2 in the affirmative. The key comparison of the second reciprocal elementary symmetric functions follows from the squared norm of a weighted sum of outward normals and an inequality for the sides and inradius of the pedal triangle. This is separate from a comparison of the elementary symmetric functions of the distances themselves.
\end{sectionabstract}
\paragraph{Theorem.}
Let $P$ be an interior point of a nondegenerate triangle $ABC$ and let $U,v$ be respectively the circumradius and inradius of its pedal triangle. For every $0<k<1$,
\[
 \frac1{PA^k}+\frac1{PB^k}+\frac1{PC^k}
 \le\frac1{(2U)^k}+\frac2{(4v)^k}.
\]
Equality holds exactly for an equilateral triangle with $P$ at its center. This is Conjecture~3.2, Eq.~(30), of \cite{r50meiji:liu}. The reciprocal first-power endpoint is already part of Liu's Theorem~1.1.

\paragraph{First and third coefficients.}
Write $X=(PA,PB,PC)$, $Y=(2U,4v,4v)$ and $e_j$ for the $j$th elementary symmetric function. Reciprocals of triples below are componentwise. Liu's Eq.~(7), including its equality criterion, yields $e_1(1/X)\le e_1(1/Y)$.

The pedal triangle has positive area, since $P$ is strictly inside the original circumdisk and its pedal area is one quarter of the original area times $1-OP^2/R^2$ ($O,R$ are the original circumcenter and radius). Write $a_i$ for its side $X_i\sin A_i$, where $A_i$ is the corresponding original angle. This side formula is the chord formula in the circle having $PA_i$ as diameter, with $A_i$ here interpreted as its vertex. If $s$ is the pedal semiperimeter, then $\prod a_i=4Uvs$. By log-concavity of sine and convexity of cotangent on $(0,\pi/2)$,
\[
 \prod\sin A_i\le 3\sqrt3/8,\qquad
 s/v=\sum\cot(\alpha_i/2)\ge3\sqrt3,
\]
where $\alpha_i$ are the pedal angles. Hence $\prod X_i\ge32Uv^2=\prod Y_i$, giving the desired third coefficient inequality after inversion.

\paragraph{The second reciprocal coefficient.}
Since $1/X_i=\sin A_i/a_i$,
\[
 e_2(1/X)=\frac{\sum_{\rm cyc}a_i\sin A_j\sin A_\ell}{4Uvs}.
\]
The identity $2\sin A_j\sin A_\ell=\cos(A_j-A_\ell)+\cos A_i$ bounds the numerator by $s+\frac12\sum_i a_i\cos A_i$. To control the latter sum, take the outward unit normals $n_i$ to the original triangle's sidelines. Their scalar products satisfy $n_j\cdot n_\ell=-\cos A_i$. Put $d_i=\sqrt{a_ja_\ell/a_i}$. Since $d_jd_\ell=a_i$, positivity of $\|\sum_i d_in_i\|^2$ gives
\[
 2\sum_i a_i\cos A_i\le\sum_{\rm cyc}\frac{a_ja_\ell}{a_i}.
\]
We claim the right side is at most $sU/v$. Indeed it equals $4Uvs\sum_i a_i^{-2}$, so the claim is equivalent to $4v^2\sum_i a_i^{-2}\le1$. With $x=s-a_1$, $y=s-a_2$, $z=s-a_3$ one has $s=x+y+z$, $a_1=y+z$ and $v^2=xyz/s$ by Heron's formula. Thus $4v^2/a_1^2\le x/s$, and the cyclic sum proves the claim. In total,
\[
 e_2(1/X)\le\frac{s+sU/(4v)}{4Uvs}
 =\frac{U+4v}{16Uv^2}=e_2(1/Y).
\]

\paragraph{Passage to concave powers.}
The three coefficient comparisons imply
$\prod_i(1+h/X_i)\le\prod_i(1+h/Y_i)$ for all $h>0$.
Taking logarithms and integrating against $h^{-1-k}\,dh$ proves the theorem: substitution $t=hb$ gives
\[
 \int_0^\infty\log(1+hb)h^{-1-k}\,dh
 = b^k\int_0^\infty\log(1+t)t^{-1-k}\,dt .
\]
The final integral is finite and positive precisely in the range used here. If the linear reciprocal endpoint is strict then the product comparison is strict for every $h>0$. Conversely, in the equilateral centered case $X=Y$, proving the equality statement.
\begin{thebibliography}{1}
\bibitem{r50meiji:liu} J. Liu, A pedal triangle inequality with the exponents, \textit{International Journal of Open Problems in Computer Science and Mathematics} 5 (2012), no.~4, 16--24, Theorem 1.1, Eq.~(7), and Conjecture 3.2.
\end{thebibliography}

\par\endgroup

% Source: /workspace/scratch/81914cec2c78/work/archive/Tournament_Round_50_Final/Tournament_Round_50/taisho/paper.tex
% Edited source: /workspace/scratch/81914cec2c78/work/audit_42_50/cleaned/r50_taisho.tex
\begingroup

\Needspace{9\baselineskip}
\section{\texorpdfstring{A genus-two reduction of Barry's three-parameter Somos--6 conjecture}{A genus-two reduction of Barry's three-parameter Somos--6 conjecture}}\label{section:r50taisho}
\noindent{\small\sffamily\color{muted}Round 50 \enspace / \enspace Taisho}\par\medskip


\begin{sectionabstract}
Barry's Conjecture 7 (2022) gives a pair of polynomial coefficients for
a width-six bilinear recurrence among the principal Hankel determinants of a
quadratic moment series. Completing the square in its sextic model gives a
linear remainder. This identifies the coefficients with van der Poorten's
two curve invariants; Hone's continued-fraction normalization identifies the
actual Hankel sequence, not just the curve. Specialization of regular moment
determinants handles zero terms without division.
\end{sectionabstract}

\subsection{The precise identity}
Write $g(x)=\sum_{j\ge0}m_jx^j$, $g(0)=1$. Clearing the denominator of
the generalized Jacobi fraction in \cite[Conj.~7, p.~7]{r50taisho:B} gives
\begin{equation}
 \bigl(1-2x-(r+1)x^2-sx^3\bigr)g(x)
       =1-x+t x^3(1-x)g(x)^2. \label{r50taisho:eq:moment}
\end{equation}
It defines the $m_j$ uniquely over $\mathbb Z[r,s,t]$. Define the size-zero
determinant $\Delta_0=1$ and, for $k\ge1$,
$\Delta_k=\det(m_{i+j})_{0\le i,j<k}$; Barry's term is $h_n=\Delta_{n+1}$.
\begin{theorem}
For all $n\ge6$ there is a polynomial identity
\[
 h_n h_{n-6}=\alpha h_{n-1}h_{n-5}+\gamma h_{n-3}^2 ,
\]
with the coefficients of Barry's Conjecture 7:
\begin{align*}
 \alpha&=t^2(r+2)^2,\\
 \gamma &= t^3\bigl[
 r^3t+r^2(s+7t)+2r\{s^2+2(t+1)s+t(t+8)\}\\[-1mm]
 &\hspace{26mm}+s^3+s^2(3t+4)+s(t+2)(3t+2)+t(t^2+4t+12)
 \bigr].
\end{align*}
\end{theorem}
In particular the cross-multiplied statement specializes to every parameter
triple, including ones yielding identically zero denominators in a
quotient-form recurrence. The latter form has meaning only at nonzero
denominators. The source definition uses $n>6$; the polynomial formula also
covers $n=6$.

\subsection{The curve and the marked complete quotient}
Set $X=1/x$, $G(X)=X^{-1}g(1/X)$, and
\[
 P=X^3-2X^2-(r+1)X-s,\qquad A=P-2t,\qquad
 v=t(r+2),\qquad w=-\frac{s+t}{r+2}.
\]
Work initially over the parameter function field. Equation \eqref{r50taisho:eq:moment}
becomes $P G=X(X-1)+t(X-1)G^2$. Its discriminant has a second expression:
\begin{equation}
 F=P^2-4tX(X-1)^2=A^2-4v(X-w).             \label{r50taisho:eq:sextic}
\end{equation}
Indeed $P-X(X-1)^2=-(r+2)X-s$. Take $Y^2=F$ with
$Y=X^3+O(X^2)$ at the selected infinity. The root of the moment quadratic
with leading term $X^{-1}$ is
\begin{equation}
 G=\frac{2X(X-1)}{Y+P},\qquad
 \frac{1}{G}=\frac{Z+t}{X(X-1)},\qquad Z=\frac{Y+A}{2}. \label{r50taisho:eq:branch}
\end{equation}
Thus $Z^2-AZ+v(X-w)=0$. The denominator $X(X-1)$ divides $F-P^2$;
it is not an arbitrary quadratic attached to a curve of the right shape.
This is exactly the positive-degree complete quotient whose reciprocal
must generate Barry's moments.

\subsection{Applying the continued-fraction identities}
We state precisely which published results enter. In van der Poorten's
notation \cite[Thms.~3.1--3.2 and remarks following Thm.~3.1]{r50taisho:V}, generic
linear-partial-quotient steps have the form
\[
 \frac{Z+d_j(X+e_j)}{u_j (X^2-v_jX+w_j)} .
\]
For a cubic $A$ and the linear remainder in \eqref{r50taisho:eq:sextic}, put
$T_{j-1}T_{j+1}=d_jT_j^2$. Then
\begin{equation}
 T_{j-3}T_{j+3}
    =v^2T_{j-2}T_{j+2}-v^3 A(w) T_j^2.   \label{r50taisho:eq:somos}
\end{equation}
The printed depressed-cubic coefficient $g+wf+w^3$ in that paper is
$A(w)$: translation of $X$ removes the quadratic term, shifting the
argument and the cubic together without changing its value there, or the
leading coefficient of $d_j(X+e_j)$. The formula therefore applies in our
original coordinate.

To identify $T_j$, use a generic starting complete quotient
$Y_1=(Y+P_1)/(2U)$, where $U$ is monic quadratic, $P_1=A+2L$ with
$\deg L\le1$, and $U\mid F-P_1^2$. In van der Poorten's normalization this
is $(Z+L)/U$; in Hone's normalization at the $j$th step the numerator is
$Y+P_j$ and $P_j=A+2d_jX+O(1)$. Thus their $d_j$ have the same sign and
scale. With $\sum_{k\ge0}s_kX^{-k-1}=1/Y_1$ and the size-$k$ determinants
$\Delta_k$ of the $s$'s, Hone \cite[Thm.~4.1, (4.3), (4.14)]{r50taisho:H} gives
\[
 d_j=\frac{\Delta_j\Delta_{j-2}}{\Delta_{j-1}^2}\qquad(j\ge2).
\]
This identity is also invariant under translation of $X$: the moment
matrix changes by triangular binomial transformations of determinant one.
Hence we may take $T_j=\Delta_{j-1}$ for $j\ge1$. A multiplicative gauge
$c b^j$ in the $T_j$ is immaterial to \eqref{r50taisho:eq:somos}. For $n\ge6$,
take $j=n-1$, so $T_{j+3}=\Delta_{n+1}=h_n$,
$T_{j-3}=\Delta_{n-5}=h_{n-6}$, and $T_j=\Delta_{n-3}=h_{n-3}$.
All required factor identities use $d_{j-2},\ldots,d_{j+2}$ at indices
at least three; no negative-index Hankel determinant is inferred.

\subsection{Passing to the special marked pair}
One should not apply a generic two-sided continued fraction directly to
\eqref{r50taisho:eq:branch}: its linear correction $L=t$ is constant.
For each fixed $n$, first work on the smooth genus-two curve for generic
$r,s,t$; it is indeed generically smooth (at $r=s=t=1$,
$\gcd(F,F')=1$). Vary an ordered pair of finite points $(a,y_a),(b,y_b)$
on this fixed curve with $a\ne b$. Set $U=(X-a)(X-b)$ and choose the
unique $L$ of degree at most one with
$2L(a)=y_a-A(a)$, $2L(b)=y_b-A(b)$. This gives precisely the required
divisibility and yields a regular family of reciprocals
\[
   G'=\frac{2U}{Y+A+2L}.
\]
All its Laurent coefficients at the chosen infinity are regular on the
indicated open set: the interpolation only divides by $a-b$, and the
leading coefficient of $Y+A+2L$ is the constant 2.

For clarity about genericity, the generic effective degree-two divisor on
a genus-two curve gives a generic Jacobian class after subtracting a fixed
degree-two divisor. Indeed Riemann--Roch gives a unique effective
representative away from the canonical class. Pairs with repeated or
infinite points, coincident $X$-coordinates, or exceptional leading
coefficient of $L$, impose proper algebraic conditions. For the latter
one may vary $b$: $Y-A$ is not constant on the curve. On the resulting
dense chart the complete quotient is normal (its polynomial part is
linear, and the norm division gives a quadratic next denominator when the
correction has degree one). The continued-fraction step acts by
translation by the infinity difference on Jacobian classes
\cite[remarks after Thm.~3.2]{r50taisho:V}. Thus excluding the finitely many
translates of the complement of this dense open chart needed for a fixed
$n$ leaves a dense open set of admissible pairs. This argument does not
assume the class of the infinity difference is nontorsion.

Consequently \eqref{r50taisho:eq:somos} holds on a dense set for the size-indexed
determinants of $G'$. The bilinear difference of those determinants is
regular at every pair with $a\ne b$, so it vanishes at the pair
\[
 (a,y_a)=(0,P(0)),\qquad (b,y_b)=(1,P(1)).
\]
Here $F=P^2$ at both abscissas and the interpolation is simply $L=t$.
Thus $G'=G$ as required; no division by a vanishing Hankel determinant
has been made at specialization.

Finally $\alpha=v^2$, $\gamma=-v^3A(w)$. Substitution and expansion yield
the two displayed polynomials; in particular all factors $r+2$ in the
denominators cancel against $v^3$. Since every $h_n$ was already a
polynomial over $\mathbb Z[r,s,t]$, equality in the generic function field
is equality of polynomials, including at $r=-2$, $t=0$, and at parameters
where $F$ is singular. This proves the theorem.

\paragraph{Scope and check.}
At $r=s=t=1$ the coefficients are $(\alpha,\gamma)=(9,77)$ and the
initial determinants are $1,3,-10,-77,-807,437,468332$.
At $(-3,0,-1)$ one gets $(1,3)$, matching Barry's Example 9.
The labels and displayed data of Example 8 do not agree: at its printed
$(-2,-2,1)$, equation \eqref{r50taisho:eq:moment} gives $m_2=1$, not $0$.
This discrepancy is not a disproof of the general assertion. The
underlying genus-two and Hankel theorems are previously published; the
contribution here is the explicit reduction, including the chosen
divisor, parameters and specialization.

\begin{thebibliography}{9}
\bibitem{r50taisho:B} P. Barry, \emph{Conjectures on Somos 4, 6 and 8 sequences using
Riordan arrays and the Catalan numbers}, arXiv:2211.12637 (2022),
\url{https://arxiv.org/abs/2211.12637}.
\bibitem{r50taisho:V} A. J. van der Poorten, \emph{Curves of genus 2, continued
fractions, and Somos sequences}, J. Integer Seq. \textbf{8} (2005),
Article 05.3.4, \url{https://cs.uwaterloo.ca/journals/JIS/VOL8/Poorten2/vdp89.pdf}.
\bibitem{r50taisho:H} A. N. W. Hone, \emph{Continued fractions and Hankel determinants
from hyperelliptic curves}, Commun. Pure Appl. Math., arXiv:1907.05204,
\url{https://arxiv.org/abs/1907.05204}.
\end{thebibliography}

\par\endgroup

% Source: /workspace/scratch/81914cec2c78/work/archive/Tournament_Round_50_Final/Tournament_Round_50/showa/paper.tex
% Edited source: /workspace/scratch/81914cec2c78/work/audit_42_50/cleaned/r50_showa.tex
\begingroup

\Needspace{9\baselineskip}
\section{\texorpdfstring{An all-orders subgroup-sum bound for the full generating set}{An all-orders subgroup-sum bound for the full generating set}}\label{section:r50showa}
\noindent{\small\sffamily\color{muted}Round 50 \enspace / \enspace Showa}\par\medskip


\begin{sectionabstract}
For the generating set consisting of an entire finite group, we establish both versions of the Lagarias-type subgroup-sum question posed in 2019. Above order $63$ a short tuple count suffices. Below that cutoff, a random growing subgroup bounds the number of subgroups of each order using only divisor arithmetic, so neither a group classification nor a database is needed.
\end{sectionabstract}
Let $G$ be a group of order $N$, generated by $S=G$. Its harmonic sum $h=\sum_{g\in G}(1+|g|_S)^{-1}$ equals $(N+1)/2$: the identity has length zero and every other element has length one. Define $L(h)=h+e^h\log h$.
\begin{theorem}
For every finite group $G$ and $S=G$,
\[
 \max\left\{\sum_{H\le G}|H|,\ \sum_{H\le G}[G:H]\right\}\le L(h),
\]
with equality precisely for the trivial group.
\end{theorem}
The first sum is the original choice in [1]; the second is its author's revision. The choice $S=G$ is one of the explicit questions in [1]. We make no claim for general generating sets. In particular the singleton generator of a cyclic group, which recovers the RH-equivalent inequality, is not covered in general.

\paragraph{Large orders.}
As in Lucia's asymptotic argument [1], growing generators one by one at least doubles subgroup size at each new generator; hence each subgroup is generated by at most $r=\lfloor\log_2 N\rfloor$ elements. Padding with identities, there are at most $N^r$ subgroups, so each weighted sum is at most $N^{1+\log_2N}$. For $x\ge64$ put $f(x)=(1+\log_2 x)\log x$. Then $f(64)=42\log 2<32$ and $f'(x)=(1+2\log_2 x)/x$ is decreasing there, since $2/\log2 < 1+2\log_2 x$. Thus $f'(x)\le13/64<1/2$. It follows that for $N\ge64$ the sums are less than $e^{N/2}<L(h)$, the last step using $\log h>1$.

\paragraph{Small orders: a universal certificate.}
Start a random chain at $K=\{1\}$; at each proper $K$ choose $g$ uniformly from $G\setminus K$ and replace $K$ by $\langle K,g\rangle$. Fix $H\le G$ of order $m$. For divisors $k$ of $m$ define
\[
 A_m(m)=1,\quad
 A_m(k)=\frac{N-k}{m-k}
 \max_{\substack{k\mid j\mid m\\j>k}} A_m(j)\qquad(k<m).
\]
If $K\le H$ has size $k$, the conditional probability of hitting $H$ is at least $1/A_m(k)$. This follows by downward induction on $k$: with probability $(m-k)/(N-k)$ the chosen element is in $H$, in which case the next subgroup remains in $H$ and has an allowed order $j$. Since a chain hits at most one group of order $m$, summing these hitting probabilities over the subgroups $H$ gives
\[
 \#\{H\le G:|H|=m\}\le b_m:=\lfloor A_m(1)\rfloor .
\]
In particular the maximum of the two sums in the theorem is at most
\[
 B_N=\max\left\{\sum_{m\mid N} m b_m,\quad
                         \sum_{m\mid N}(N/m)b_m\right\}.
\]
For $1\le N<64$ compute the following division and comparisons over the rationals. The lower bound uses $\log h\ge 2(h-1)/(h+1)$, obtained by differentiating or expanding $2\operatorname{arctanh}((h-1)/(h+1))$, and $e^h\ge\sum_{j=0}^{10} h^j/j!$.
{\small
\begin{Verbatim}
"""Exact finite certificate. Uses only Python's standard library."""
from fractions import Fraction as Q
from functools import lru_cache
from math import factorial
def certificate(n):
    div=[v for v in range(1,n+1) if n%v==0]
    @lru_cache(None)
    def A(m,k):
        if k==m: return Q(1)
        return Q(n-k,m-k)*max(A(m,j) for j in div if j>k and j%k==0 and m%j==0)
    counts={m:int(A(m,1)) for m in div}
    B=max(sum(counts[m]*w(m) for m in div) for w in (lambda m:m, lambda m:n//m))
    h=Q(n+1,2)
    low=h+2*(h-1)/(h+1)*sum(h**j/factorial(j) for j in range(11))
    return B,low
for n in range(1,64):
    B,low=certificate(n)
    assert B==low if n==1 else B<low
    print(f'n={n}: order/index upper bound {B}; comparison passed.')
\end{Verbatim}

}

All 63 comparisons pass; for example $B_1=1,B_2=3,B_4=11,B_8=51,B_{16}=307$. The displayed Python implementation performs this finite certificate using exact rational arithmetic. Since $Q\le L(h)$, the proof is complete.

\paragraph{Source.} [1] user6671, \emph{A group theoretic interpretation of Lagarias inequality}, MathOverflow 330077 (26 April 2019; edited 24 May 2019), including Lucia's large-order answer, 27 April 2019. https://mathoverflow.net/q/330077.

\par\endgroup

% Source: /workspace/scratch/81914cec2c78/work/archive/Tournament_Round_50_Final/Tournament_Round_50/reiwa/paper.tex
% Edited source: /workspace/scratch/81914cec2c78/work/audit_42_50/cleaned/r50_reiwa.tex
\begingroup
\newcommand{\B}{\mathcal B}
\Needspace{9\baselineskip}
\section{\texorpdfstring{The first resonance determines Altizio's weighted extremal limit}{The first resonance determines Altizio's weighted extremal limit}}\label{section:r50reiwa}
\noindent{\small\sffamily\color{muted}Round 50 \enspace / \enspace Reiwa}\par\medskip


\begin{sectionabstract}
For the geometric sum on the closed unit disk centered at two, we prove the numerical prediction of Altizio that the minimum of its modulus after division by \(z\) tends to \(2\pi\). The disk geometry confines bounded values to an \(O(1/n)\) neighborhood of the tangency point. Under rescaling, the only possible limits for near-minimizers are the first two nonzero zeros of the limiting exponential integral.
\end{sectionabstract}

Let \(\B=\{z\in\mathbb C:|z-2|\le1\}\), let \(S_n(z)=\sum_{j=0}^{n-1}z^j\), and put \(G_n(z)=S_n(z)/z\). Since \(0\notin\B\), the minimum in the following theorem exists. Here \(S_n\) denotes a geometric sum, not a Stern polynomial.
\begin{theorem}\label{r50reiwa:main}
The numbers \(m_n=\min_{z\in\B}|G_n(z)|\) satisfy \(m_n\to2\pi\). Moreover, for any choice of near-minimizers \(a_n\in\B\) with \(|G_n(a_n)|-m_n\to0\),
\[
 \operatorname{dist}\bigl(n(a_n-1),\{-2\pi i,2\pi i\}\bigr)\longrightarrow0.
\]
\end{theorem}
The limit is the prediction in Remark 4 of \cite{r50reiwa:A}. It is distinct from both the pointwise inequality of Remark 3(2) for \(S_n\) without a divisor \(z\), and the degree-classification question about \(|S_n(z)|\ge n|z|\) elsewhere in Remark 4. Neither of those results is used below.

\paragraph{A sequence on the boundary.}
Let \(h=2\pi/n\) and \(w_n=2-e^{-ih}\), so \(w_n\in\partial\B\). Near \(h=0\), the logarithm chosen near one satisfies
\[
 w_n=1+ih+\tfrac12 h^2+O(h^3),\qquad
 \log w_n=ih+h^2+O(h^3).
\]
Thus \(n\log w_n=2\pi i+4\pi^2/n+O(n^{-2})\), and
\[
  G_n(w_n)=\frac{w_n^n-1}{w_n(w_n-1)}
       =\frac{4\pi^2/n+O(n^{-2})}{2\pi i/n+O(n^{-2})}
       \longrightarrow-2\pi i.
\]
This gives \(\limsup m_n\le2\pi\). A limit for a trial sequence alone does not yet give the minimum.

\paragraph{Localization for the lower bound.}
For \(z\in\B\), let \(y=|z-1|\). Expanding \(|(z-1)-1|^2\le1\) shows \(2\Re(z-1)\ge y^2\). Hence \(|z|^2\ge1+2y^2\) and \(1\le |z|\le3\). For \(n\ge2\) and \(y>0\), Bernoulli's inequality and the reverse triangle inequality give
\begin{equation}\label{r50reiwa:bound}
 |G_n(z)|\ge\frac{|z|^n-1}{|z|y}
 \ge\frac{(1+2y^2)^{n/2}-1}{|z|y}
 \ge\frac{ny}{|z|}.
\end{equation}
At \(z=1\) the last inequality holds directly. Thus every sequence \(a_n\in\B\) with \(|G_n(a_n)|\) bounded has \(u_n=n(a_n-1)\) bounded. In particular this holds for near-minimizers, and \(a_n\to1\).

Pass to any convergent subsequence \(u_n\to u\). Then
\[
 \frac{G_n(a_n)}{n}
   =\frac1{a_n n}\sum_{j=0}^{n-1}\left(1+\frac{u_n}{n}\right)^j
    \longrightarrow \int_0^1 e^{ut}\,dt.
\]
Indeed, uniformly in \(0\le j\le n\), the logarithmic exponent differs from \(ju/n\) by \(O(1/n)+o(1)\), since \(u_n\) is bounded and tends to \(u\); the remaining sums are Riemann sums. Boundedness forces the integral to vanish. Its value is \(1\) at \(u=0\) and \((e^u-1)/u\) otherwise, so \(u=2\pi ik\) for a nonzero integer \(k\).

Use exact minimizers and take the preceding subsequence from one realizing the lower limit of \(m_n\). By \eqref{r50reiwa:bound}, \(\liminf m_n\ge |u|\ge2\pi\). Along with the upper construction, this proves the claimed limit. For arbitrary near-minimizers we now have \(|G_n(a_n)|\to2\pi\). Inequality \eqref{r50reiwa:bound} gives \(|u|\le2\pi\) for each subsequential limit. The only remaining values are \(u=\pm2\pi i\), and boundedness gives the stated distance conclusion. \hfill\(\square\)

\begin{thebibliography}{1}
\bibitem{r50reiwa:A} D. Altizio, \emph{Zeros of Stern polynomials in the complex plane}, arXiv:2511.03847 (2025), Section~4.3, Remark~4.
\end{thebibliography}

\par\endgroup

% Source: /workspace/scratch/81914cec2c78/work/archive/Tournament_Round_50_Final/Tournament_Round_50/keio/paper.tex
% Edited source: /workspace/scratch/81914cec2c78/work/audit_42_50/cleaned/r50_keio.tex
\begingroup

\Needspace{9\baselineskip}
\section{\texorpdfstring{A median-ratio lower bound for the sum of three vertex distances}{A median-ratio lower bound for the sum of three vertex distances}}\label{section:r50keio}
\noindent{\small\sffamily\color{muted}Round 50 \enspace / \enspace Keio}\par\medskip

\begin{sectionabstract}
We prove the inequality proposed as Conjecture 3 by Liu and Chu in 2003. An equilateral-direction projection eliminates the variable point. The area determinant of two medians bounds the correction for their unequal lengths. A rational parametrization of the half-angles yields a polynomial certificate with seven charts. This establishes the statement, without asserting bibliographic priority.
\end{sectionabstract}
\paragraph{Theorem.} For every nondegenerate triangle $ABC$ of inradius $r$, with medians $m_a,m_b,m_c$, and every point $P$ in its plane,
\[
 PA+PB+PC\ge \frac{3r}{2}
 \left(\sqrt{m_b/m_c}+\sqrt{m_c/m_b}\right)^2. \tag{1}
\]
This is (29) of [1], not the authors' separate median-to-altitude Conjecture 4.
\paragraph{Reduction.} Write $s$ for semiperimeter, $D$ for area, $d=a^2+b^2+c^2$, $J=\sqrt3D$ and $F=\sqrt{d/2+2J}$. In complex coordinates the two numbers $A+\omega B+\omega^2 C$ ($\omega$ a primitive cube root of unity) have squared moduli $d/2\pm2\sqrt3D$ by expansion. Subtracting $P$ from all three vertices does not change either number. The triangle inequality therefore gives $PA+PB+PC\ge F$ regardless of the location of $P$.
Put $v=(b^2-c^2)^2$ and $N=4a^2+b^2+c^2+6J$. The absolute determinant of the median vectors from $B,C$ is $3D/2$: with $A$ at the origin these vectors are $C/2-B,B/2-C$. In particular $m_bm_c\ge3D/2$. Using the median formula twice,
\[
 (m_b^2-m_c^2)^2=9v/16,\qquad
 (m_b+m_c)^2\ge (4a^2+b^2+c^2)/4+3D\ge N/4 .
\]
Since $sr=D$, multiplying the right side of (1) by $s$ gives
\[
 6D+\frac{3D}{2m_bm_c}\frac{(m_b^2-m_c^2)^2}{(m_b+m_c)^2}
 \ \le\ 2\sqrt3J+\frac{9v}{4N}. \tag{2}
\]
It suffices to prove that $sF$ exceeds or equals the last expression. Both sides are nonnegative. Squaring, clearing $(4N)^2$, and using $144\sqrt3<252$ shows it suffices that
\[
 H=8N^2[s^2(d+4J)-24J^2]-81v^2-252JNv\ge0. \tag{3}
\]
\paragraph{Certificate for (3).} Pick one angle at most $60^\circ$ and another at least $60^\circ$. For these angles and the remaining one let $p,q,t$ be $\sqrt3$ times their respective half-angle tangents. We have $p\le1\le q$, $pq<3$ and $t=(3-pq)/(p+q)$. The following chart covers all admissible pairs:
\[
 p=\frac1{1+x},\qquad q=\frac{1+3(1+x)y}{1+y},\quad x,y\ge0;
 \qquad x=1/p-1,\quad y=\frac{q-1}{3/p-q}.
\]
Call the displayed numerators and denominators $p_1,p_2,q_1,q_2$ and set $t_1=3p_2q_2-p_1q_1$, $t_2=p_1q_2+p_2q_1$. Write $T=p_1q_1t_1$ and
\[
 (U,V,W)=(p_2q_1t_1,p_1q_2t_1,p_1q_1t_2)=T(1/p,1/q,1/t).
\]
By homogeneity normalize $r=T/\sqrt3$; hence $U,V,W$ are the contact lengths $r\cot(A_i/2)$. Thus $s=U+V+W$, $J=sT$, and the squared sides are
$A_i=(s-U_i)^2$ for $(U_0,U_1,U_2)=(U,V,W)$. Choose $a^2=A_i$ successively for $i=0,1,2$; (3) is symmetric in the two remaining sides $b,c$. All expressions in (3) now have integer coefficients.

For any polynomial $h(x,y)$ use the operations
\[
 Xh=h(x+y,y),\qquad Yh=h(x,x+y).
\]
Their domains together cover the full original nonnegative quadrant (subtract the smaller coordinate from the larger). Below are the number of nonzero coefficients and the smallest such coefficient, directly verified by polynomial arithmetic; the entry $XY$ means $Y(XH)$. Missing monomials have coefficient zero.
\[
\begin{array}{c|c|r|r}
 i&\text{successive operations}&\#&\min\\\hline
 0&X&686&419904\\ 0&Y&558&520224768\\
 1&X&685&2092959\\ 1&Y&558&1950842880\\
 2&XX&686&2092959\\ 2&XY&858&520224768\\
 2&Y&558&780337152
\end{array}
\]
Positivity of all listed coefficients gives (3) in all three cases, proving (1). Equality is attained for the equilateral triangle at its center. The short checker below performs the substitutions and coefficient tests in $\mathbb Z[x,y]$; it does not sample triangles.
\paragraph{Exact computational certificate.}
The following program contains the finite integer or rational calculation used in the proof.
{\small
\begin{Verbatim}
from sympy.polys.rings import ring
from sympy import ZZ
K,x,y=ring('x,y',ZZ)
p1,p2=K.one,1+x
q1,q2=1+3*p2*y,1+y
t1,t2=3*p2*q2-p1*q1,p1*q2+p2*q1
U,V,W=p2*q1*t1,p1*q2*t1,p1*q1*t2
T=p1*q1*t1
S=U+V+W; J=S*T
A,B,C=[(S-u)**2 for u in (U,V,W)]
def report(i,label,f):
 vals=f.values()
 print(i,label,len(vals),min(vals))
 assert all(z>0 for z in vals)

for i,(a,b,c) in enumerate(((A,B,C),(B,A,C),(C,A,B))):
 # Here a,b,c stand for the squared side lengths.
 d=a+b+c
 N=4*a+b+c+6*J
 v=(b-c)**2
 H=8*N*N*(S*S*(d+4*J)-24*J*J)-81*v*v-252*J*N*v
 # Positive quadrant = union of x >= y and y >= x.
 X=H.compose(x,x+y)
 Y=H.compose(y,x+y)
 if i<2:
  report(i,"X",X)
 else:
  report(i,"XX",X.compose(x,x+y))
  report(i,"XY",X.compose(y,x+y))
 report(i,"Y",Y)
\end{Verbatim}

}

\paragraph{Reference.} [1] J. Liu and X.-G. Chu, A New Geometry Inequality with the Fermat Problem, \emph{Journal of East China Jiaotong University} 20(1) (2003), 89--93, p.~92, Conjecture 3 (29). \url{https://journal.ecjtu.edu.cn/hdjtdxen/article/pdf/20030127}. A limited search found related later Chinese articles whose precise scope was not verified; historical priority is left open.

\par\endgroup


\clearpage
\part{Additional complete proofs}

% Additional source: /workspace/scratch/81914cec2c78/work/additional_geometry/r34_full_proof.tex
\begingroup

\section{\texorpdfstring{A full certificate for Liu's two-product pedal-triangle inequality}{A full certificate for Liu's two-product pedal-triangle inequality}}\label{section:add34keio}
\noindent{\small\sffamily\color{muted}Additional completion \enspace / \enspace Round 34 \enspace / \enspace Keio}\par\medskip


\begin{sectionabstract}
We prove Liu's Conjecture 3.6 (2013) for every nondegenerate triangle
and every interior point. A reversible polynomial reduction is covered
by all six barycentric order chambers. Four chambers have the square
corrections already present in the archived partial argument. In each
remaining chamber, extracting an explicitly nonnegative boundary
polynomial leaves a polynomial with nonnegative coefficients after
simple square corrections. The complete exact verifier is included.
\end{sectionabstract}

\subsection{Statement and geometric reduction}

Let $P$ be an interior point of a nondegenerate triangle $ABC$.
Write $a,b,c$ for its side lengths opposite $A,B,C$, $S$ for its area,
$R$ for its circumradius, and $O$ for its circumcenter.
Let $r_1,r_2,r_3$ be the perpendicular distances from $P$ to the
lines $BC,CA,AB$, and let $R_1=PA$, $R_2=PB$, $R_3=PC$.
Let $S_p$ be the area of the pedal triangle formed by those three
perpendicular projections.

\begin{theorem}\label{add34keio:main}
For every such triangle and every interior point $P$,
\begin{equation}\label{add34keio:goal}
 8r_1r_2r_3+R_1R_2R_3\ \le\ 8R^3\frac{S_p}{S}.
\end{equation}
Equality holds if and only if $ABC$ is equilateral and $P$ is its center.
\end{theorem}

This is Conjecture 3.6, equation (3.12), of Liu \cite{add34keio:liu}, also
Conjecture 5, equation (3.8), of the author's earlier preprint
\cite{add34keio:preprint}. The archived Round 34 Keio argument covered four of
six order chambers. The proof here includes all six and all auxiliary
identities needed for the completion. It does not assert historical priority.

We use the pedal-area identity
\begin{equation}\label{add34keio:pedal}
 \frac{S_p}{S}=\frac{R^2-OP^2}{4R^2}.
\end{equation}
For completeness, after translating the circumcenter to $0$ and
scaling $R$ to $1$, use complex vertex coordinates $A,B,C$ of modulus
one and point coordinate $p$. The foot on $BC$ is
$d=(B+C+p-BC\overline p)/2$, with cyclic formulas for $e,f$.
Substitution into the oriented-area expression gives
\[
 \operatorname{Im}(\overline d e+\overline e f+\overline f d)
 =\tfrac14(1-|p|^2)
 \operatorname{Im}(\overline A B+\overline B C+\overline C A).
\]
The ratio is positive for an interior point of the triangle, which is
strictly inside its circumdisk. This proves \eqref{add34keio:pedal} for ordinary
positive areas and also shows that the pedal triangle is nondegenerate.

Write $(u:v:w)$ for the positive barycentric coordinates of $P$ and put
\[
 D=u+v+w,\qquad T=(abc)^2,\qquad q=16S^2,
 \qquad L=a^2vw+b^2wu+c^2uv.
\]
The weighted variance identity for the vectors $A-O,B-O,C-O$ gives
\[
 R^2-OP^2=\frac{L}{D^2}.
\]
Indeed the difference between the weighted mean of their squared norms
and the squared norm of their weighted mean equals the sum over
unordered vertex pairs of the product of the normalized weights times
the squared pairwise distance. The area ratios similarly give
\[
 R=\frac{abc}{\sqrt q},\qquad
 r_1=\frac{\sqrt q\,u}{2aD},\quad
 r_2=\frac{\sqrt q\,v}{2bD},\quad
 r_3=\frac{\sqrt q\,w}{2cD}.
\]
Define
\begin{align*}
 Q_1&=(v+w)(c^2v+b^2w)-a^2vw,\\
 Q_2&=(w+u)(a^2w+c^2u)-b^2wu,\\
 Q_3&=(u+v)(b^2u+a^2v)-c^2uv.
\end{align*}
Expansion of $D(P-A)=v(B-A)+w(C-A)$ and its cyclic versions gives
$R_i^2=Q_i/D^2$. Each $Q_i$ is positive since $P$ is an interior
point distinct from every vertex.

Multiply \eqref{add34keio:goal} by the positive quantity $\sqrt q\,abcD^3$.
Using \eqref{add34keio:pedal}, the result is exactly
\[
 H:=2TDL-q^2uvw\ \ge\ \sqrt{qTQ_1Q_2Q_3}.
\]
Consequently the original inequality is equivalent to the two signs
\begin{equation}\label{add34keio:signs}
 H\ge0,\qquad F:=H^2-qTQ_1Q_2Q_3\ge0.
\end{equation}
This retains the sign condition required before squaring.

\subsection{The six complete order chambers}

Relabel the vertices so that $a\ge b\ge c$; the assertion is
symmetric under simultaneous relabeling of vertices, side lengths and
barycentric coordinates.  Let $s=(a+b+c)/2$ and put
$x=s-a$, $y=s-b$, $z=s-c$.  These are positive and satisfy
$x\le y\le z$.  By homogeneity normalize $x=1$, so
\[
 x=1,\qquad y=1+X,\qquad z=1+X+Y,\qquad X,Y\ge0,
\]
\[
 a=y+z,\qquad b=z+x,\qquad c=x+y,\qquad
 q=16(x+y+z)xyz.
\]
Independently scale the barycentric triple so that its minimum is one.
Its three sorted values can then be written
\[
 \ell_0=1,\qquad\ell_1=1+U,\qquad\ell_2=1+U+V,
 \qquad U,V\ge0.
\]
For a permutation $ijk$ of $012$ set
$(u,v,w)=(\ell_i,\ell_j,\ell_k)$. These six closed chambers cover
all positive barycentric triples, with ties included. The expression
$F$ is now a polynomial in $\mathbb Z[X,Y,U,V]$ for each chamber.

We use five fixed squares:
\begin{align*}
 B_1&=18432(V-16Y)^2,& B_2&=73728(V-8X)^2,\\
 B_3&=36864(U-8X)^2,& B_4&=73728(U-8Y)^2,\\
 B_5&=36864(V-8Y)^2.
\end{align*}
All are nonnegative. Two further nonnegative polynomials will handle
the chambers not covered by the archived certificate.

\subsection{The two additional boundary polynomials}

For the side lengths already defined put
\[
 \kappa_A=b^2+c^2-a^2,\quad
 \kappa_B=c^2+a^2-b^2,\quad
 \kappa_C=a^2+b^2-c^2,
\]
\[
 M=8a^2b^2c^2-q\kappa_C.
\]
Heron's identity $q=4a^2b^2-\kappa_C^2$ gives the useful identity
\begin{equation}\label{add34keio:Mpositive}
 M=8a^2b^2(a-b)^2+(2ab-\kappa_C)^2(4ab+\kappa_C)>0.
\end{equation}
The strict sign follows from
$2ab-\kappa_C=c^2-(a-b)^2>0$ and
$4ab+\kappa_C=(a+b)^2-c^2+2ab>0$, both consequences of the
strict triangle inequalities.

\paragraph{Chamber $102$.}
Here $(u,v,w)=(1+U,1,1+U+V)$. Define
\begin{equation}\label{add34keio:extra102}
 K_{102}=Ta^2b^2 V^4
 \left[b^2\kappa_A^2(1+U)^2+a^2\kappa_B^2+M(1+U)\right].
\end{equation}
This is nonnegative by \eqref{add34keio:Mpositive}. It is precisely the
coefficient of $V^4$ in $F$, multiplied by $V^4$.
One way to derive it is to take the coefficient of $w^4$ before
substituting the chamber. Directly from the displayed definitions,
that coefficient is
\[
 Ta^2b^2\bigl[b^2\kappa_A^2u^2+a^2\kappa_B^2v^2+Muv\bigr].
\]
Since $w=1+U+V$ and $F$ has degree at most four in $w$,
substitution of $u=1+U$, $v=1$ gives \eqref{add34keio:extra102}.
The verifier also checks the entire polynomial identity.

\paragraph{Chamber $120$.}
Here $(u,v,w)=(1+U,1+U+V,1)$. Set $u_0=U$, $v_0=U+V$ and define
\begin{equation}\label{add34keio:extra120}
\begin{split}
 K_{120}={}&Tc^4u_0^2v_0^2
 \left[(u_0+v_0)
       \left(u_0b^2\kappa_B^2+v_0a^2\kappa_A^2\right)
       +qc^2u_0v_0\right].
\end{split}
\end{equation}
Every factor and summand is nonnegative. This is the homogeneous
degree-six part of $F$ in $U,V$. Indeed $F$ is homogeneous of degree
six in the unnormalized barycentric variables. Its leading part in
this chamber is obtained by putting $(u,v,w)=(u_0,v_0,0)$, which gives
\[
 Tc^4u_0^2v_0^2\left[4T(u_0+v_0)^2
 -q\bigl((u_0+v_0)(b^2u_0+a^2v_0)-c^2u_0v_0\bigr)\right].
\]
Using $4T-qb^2=b^2\kappa_B^2$ and
$4T-qa^2=a^2\kappa_A^2$ reduces this expression to
\eqref{add34keio:extra120}. No assertion about a boundary point is used without
its displayed nonnegative algebraic expression.

\subsection{The full certificate and the equality case}

In chambers $102$ and $120$ use the respective $K$ above; in the other
four chambers put $K=0$. Subtract the indicated squares and write
$\mathcal R=F-K-B$. Exact expansion gives the following table:
\begin{center}
\begin{tabular}{cclrrrr}
\toprule
chamber&$K$&$B$&$\#H$&$\min H$&$\#\mathcal R$&$\min\mathcal R$\\
\midrule
012&0&0&378&2&3560&4\\
021&0&$B_1$&378&2&3569&4\\
102&$K_{102}$&$B_3$&378&2&3140&16\\
120&$K_{120}$&$B_3+B_4+B_5$&357&2&2797&4\\
201&0&$B_1+B_2+B_3$&378&2&3567&4\\
210&0&$B_2+B_3+B_4+B_5$&357&2&3365&4\\
\bottomrule
\end{tabular}
\end{center}
The counts and minima refer to nonzero coefficients; missing monomials
have coefficient zero.  In every chamber all nonzero coefficients of
$H$ and $\mathcal R$ are positive. Also the constant coefficient of
$H$ is $2304$, and each of the four coefficients
$[X^2]\mathcal R,[Y^2]\mathcal R,[U^2]\mathcal R,[V^2]\mathcal R$
is at least $1069056$. Thus, on every chamber,
\[
 H>0,\qquad
 F=K+B+\mathcal R\ge
 1069056(X^2+Y^2+U^2+V^2)\ge0.
\]
The two conditions in \eqref{add34keio:signs} follow on the full domain.
Furthermore $F=0$ is possible only when $X=Y=U=V=0$, which means
an equilateral triangle and equal barycentric coordinates.
For that configuration $P$ is the center, $r_i=R/2$, $R_i=R$ and
$S_p/S=1/4$, so both sides of \eqref{add34keio:goal} are $2R^3$.
All other configurations give strict inequality by the reversible
reduction. This proves Theorem~\ref{add34keio:main}, including its exact equality
statement.

\subsection{Complete exact verifier}

The following standalone Python 3 program uses SymPy's polynomial ring
over $\mathbb Z$. It constructs every polynomial from its defining
formula, checks the two extracted identities and all six coefficient
certificates, and verifies the coefficients used for strictness.
There is no finite sampling of triangles or points.
\begin{Verbatim}
"""Exact complete certificate for Liu (2013), Conjecture 3.6.

All six barycentric-order chambers are included. Requires Python 3 and
SymPy. The two extra polynomials are nonnegative by the formulas in the
proof; the remaining checks use integer coefficients only.
"""
from itertools import permutations
from sympy import ZZ
from sympy.polys.rings import ring

K, X, Y, U, V = ring('X,Y,U,V', ZZ)
x, y, z = K.one, 1+X, 1+X+Y
a, b, c = y+z, z+x, x+y
q, T = 16*(x+y+z)*x*y*z, (a*b*c)**2
kA, kB, kC = b*b+c*c-a*a, c*c+a*a-b*b, a*a+b*b-c*c
M = 8*a*a*b*b*c*c-q*kC
assert M == 8*a*a*b*b*(a-b)**2+(2*a*b-kC)**2*(4*a*b+kC)

B1 = 18432*(V-16*Y)**2
B2 = 73728*(V-8*X)**2
B3 = 36864*(U-8*X)**2
B4 = 73728*(U-8*Y)**2
B5 = 36864*(V-8*Y)**2
squares = {
    (0, 1, 2): K.zero,
    (0, 2, 1): B1,
    (1, 0, 2): B3,
    (1, 2, 0): B3+B4+B5,
    (2, 0, 1): B1+B2+B3,
    (2, 1, 0): B2+B3+B4+B5,
}
# Counts and minimum nonzero coefficients of H and of the residual.
expected = {
    (0, 1, 2): (378, 2, 3560, 4),
    (0, 2, 1): (378, 2, 3569, 4),
    (1, 0, 2): (378, 2, 3140, 16),
    (1, 2, 0): (357, 2, 2797, 4),
    (2, 0, 1): (378, 2, 3567, 4),
    (2, 1, 0): (357, 2, 3365, 4),
}
levels = (K.one, 1+U, 1+U+V)
for order in permutations(range(3)):
    u, v, w = [levels[i] for i in order]
    D = u+v+w
    L = a*a*v*w+b*b*w*u+c*c*u*v
    Q1 = (v+w)*(c*c*v+b*b*w)-a*a*v*w
    Q2 = (w+u)*(a*a*w+c*c*u)-b*b*w*u
    Q3 = (u+v)*(b*b*u+a*a*v)-c*c*u*v
    H = 2*T*D*L-q*q*u*v*w
    F = H*H-q*T*Q1*Q2*Q3
    extra = K.zero
    if order == (1, 0, 2):
        extra = T*a*a*b*b*V**4*(
            b*b*kA*kA*(1+U)**2+a*a*kB*kB+M*(1+U))
        # This is precisely the entire coefficient of V^4, times V^4.
        assert extra == K.from_dict({e: d for e, d in F.items()
                                     if e[3] == 4})
    if order == (1, 2, 0):
        u0, v0 = U, U+V
        extra = T*c**4*u0*u0*v0*v0*(
            (u0+v0)*(u0*b*b*kB*kB+v0*a*a*kA*kA)
            + q*c*c*u0*v0)
        # The degree-six homogeneous part in the barycentric increments.
        assert extra == K.from_dict({e: d for e, d in F.items()
                                     if e[2]+e[3] == 6})
    residual = F-extra-squares[order]
    assert min(H.values()) > 0 and H[(0, 0, 0, 0)] == 2304
    assert min(residual.values()) > 0
    assert F.get((0, 0, 0, 0), 0) == 0
    for j in range(4):
        exponent = tuple(2 if k == j else 0 for k in range(4))
        assert residual.get(exponent, 0) >= 1069056
    observed = (len(H), min(H.values()), len(residual),
                min(residual.values()))
    assert observed == expected[order]
    print(''.join(map(str, order)), observed)
print('PASS: all six chambers; equality only X=Y=U=V=0.')
\end{Verbatim}


\begin{thebibliography}{9}
\bibitem{add34keio:liu} J. Liu, On inequality $R_p<R$ of the pedal triangle,
\emph{Mathematical Inequalities \& Applications} \textbf{16}(3) (2013),
701--715, Conjecture 3.6, equation (3.12), p.~714.
\url{https://doi.org/10.7153/mia-16-53}.
\bibitem{add34keio:preprint} J. Liu, On the inequality $R_p<R$ of the pedal triangle,
\emph{RGMIA Research Report Collection} \textbf{14} (2011), Article 46,
Conjecture 5, equation (3.8).
\url{https://rgmia.org/papers/v14/v14a46.pdf}.
\end{thebibliography}

\par\endgroup

% Additional source: /workspace/scratch/81914cec2c78/work/additional_combinatorics/r43_full_composition_non_algebraicity.tex
\begingroup
\newcommand{\Qbar}{\overline{\mathbb Q}}
\newcommand{\E}{\mathbb E}
\newcommand{\Prob}{\mathbb P}
\clearpage
\section{\texorpdfstring{Equal-length restricted compositions: the full rational-alphabet nonalgebraicity theorem}{Equal-length restricted compositions: the full rational-alphabet nonalgebraicity theorem}}\label{section:add43taisho}
\noindent{\small\sffamily\color{muted}Additional completion \enspace / \enspace Round 43 \enspace / \enspace Taisho}\par\medskip


\begin{sectionabstract}
For every rational part series with nonnegative integer coefficients and at
least two supported positive part sizes, the generating function for ordered
$m$-tuples of compositions of the same size and the same length is
transcendental for every integer $m\ge3$.  The proof includes an explicit
lattice local limit calculation, so it treats finite and infinite alphabets,
nontrivial difference spans, and common divisors of the allowed sizes.
This supplies the full quantifiers missing from the archived Round 43 note.
The arithmetic amplitude obstruction is established prior work of
Banderier--Drmota; no novelty or priority is claimed.
\end{sectionabstract}

\subsection{Statement and attribution}
Let
\[
 p(z)=\sum_{j\ge1}a_jz^j\in\mathbb Q(z),\qquad
 a_j\in\mathbb Z_{\ge0},
 \qquad \mathcal P=\{j\ge1:a_j>0\}.
\]
An allowed part of size $j$ has $a_j$ colors.  Write
\[
 c_{n,k}=[z^n]p(z)^k,\qquad
 D_n^{(m)}=\sum_{k\ge0}c_{n,k}^{,m},\qquad
 D^{(m)}(z)=\sum_{n\ge0}D_n^{(m)}z^n.                 \tag{1}
\]
Thus $c_{n,k}$ counts compositions of $n$ with $k$ colored parts, and
$D_n^{(m)}$ counts ordered $m$-tuples having a common length.
The term $n=k=0$ represents the empty composition.  Every sum in (1) is
finite, since parts are positive.

\begin{theorem}\label{add43taisho:r43:main}
If $|\mathcal P|\ge2$, then $D^{(m)}(z)$ is transcendental over
$\mathbb C(z)$ for every integer $m\ge3$.
\end{theorem}

The original model of Banderier--Hitczenko assumes two distinct supported
relatively prime part sizes except where stated otherwise
\cite[Section 2]{add43taisho:r43:BH}.  Theorem~\ref{add43taisho:r43:main} therefore covers their
Conjecture 3.1 on its full original domain, and also permits a common divisor
of the part sizes.  Integer weights and two supported sizes already imply
the supercritical condition used in that model.  Both even and odd $m$ are
included, for every fixed $m$.

Banderier--Hitczenko obtained the Gaussian matching-length asymptotics.
Banderier--Drmota subsequently explained the nonalgebraicity obstruction
from the arithmetic constant for even $m$, as well as the exponent
obstruction for odd $m$ \cite[Proposition 7]{add43taisho:r43:BD}.
We give the full argument, including the lattice factors, instead of
treating the archived two-part example as the universal result.

\subsection{Normalization and the lattice local limit estimate}

There is a unique real number $\rho>0$ inside the radius of convergence of
$p$ such that
\[
                         p(\rho)=1.                  \tag{2}
\]
Here are details ensuring that no singularity case is omitted.
The function $p$ is strictly increasing on its positive convergence
interval and vanishes at zero.  If its radius $R$ exceeds $1$, two
positive integer coefficients give $p(1)\ge2$; the same argument works
for a polynomial.  If $R\le1$ is finite, nonnegativity implies that $R$
is a singularity.  Indeed, if $p$ were analytic at $R$, monotone
convergence would give
\[
 p^{(s)}(R)=\sum_{j\ge s}a_j(j)_sR^{j-s}\quad(s\ge0).
\]
For sufficiently small $h>0$, its convergent Taylor series at $R$ and
nonnegativity would then give
\[
 \sum_j a_j(R+h)^j
   =\sum_{s\ge0}\frac{p^{(s)}(R)}{s!}h^s<\infty,
\]
contradicting the radius $R$.  A singularity of a rational function is a
pole, and thus $p(t)\to+\infty$ as $t\uparrow R$.  This proves (2).
As $p$ is rational over $\mathbb Q$, the number $\rho$ is algebraic.

Define a positive integer-valued random variable $X$ by
\[
 \Prob(X=j)=a_j\rho^j,
 \qquad
 \mu=\E X=\rho p'(\rho),\qquad
 \sigma^2=\rho p'(\rho)+\rho^2p''(\rho)-\mu^2.        \tag{3}
\]
This is a probability law by (2).  Its moment generating function is
$p(\rho e^t)$, finite for $t$ in a neighborhood of zero.  The quantities
$\mu$ and $\sigma^2$ are algebraic, $\mu>0$, and $\sigma^2>0$ because
there are two different supported sizes.  For independent copies
$X_1,X_2,\ldots$ and $S_k=X_1+\cdots+X_k$,
\[
 \Prob(S_k=n)=\rho^n c_{n,k},\qquad
 \rho^{mn}D_n^{(m)}=\sum_{k\ge0}\Prob(S_k=n)^m.       \tag{4}
\]

We prove the probability estimate that will be used in (4).  It is stated
slightly more generally to make the lattice issue explicit.

\begin{lemma}\label{add43taisho:r43:llt}
Let $X$ be positive integer-valued, have an exponential moment in a
neighborhood of zero, and have mean $\mu>0$ and variance $\sigma^2>0$.
Assume the greatest common divisor of its supported values is $1$.
Let
\[
 d=\gcd\{|i-j|:i,j\in\operatorname{supp}X\}>0,
\]
and let $a$ be any one supported value.  For each fixed integer $m\ge1$,
\[
 \sum_{k\ge0}\Prob(S_k=n)^m
 \sim
 \frac{d^{m-1}\mu^{(m-3)/2}}
      {\sqrt m\,(2\pi\sigma^2)^{(m-1)/2}}
       n^{-(m-1)/2}.                                \tag{5}
\]
\end{lemma}

\begin{proof}
All supported values are congruent to $a$ modulo $d$, and
$\gcd(a,d)=1$.  Hence $\Prob(S_k=n)=0$ unless
$n\equiv ka\pmod d$; for fixed $n$ the allowed $k$ form one residue
class modulo $d$.

Put $\phi(t)=\E e^{itX}$.  For an admissible pair $(k,n)$, Fourier
inversion on this lattice is
\[
 \Prob(S_k=n)=\frac{d}{2\pi}
   \int_{-\pi/d}^{\pi/d}\phi(t)^k e^{-int}\,dt.       \tag{6}
\]
To see why there are no other points of modulus one in this interval,
equality in $|\E e^{itX}|\le1$ requires all $e^{itj}$ on the support
to be equal.  Thus $t(i-j)\in2\pi\mathbb Z$ for all supported $i,j$,
or $td\in2\pi\mathbb Z$.  Only $t=0$ satisfies this in
$[-\pi/d,\pi/d]$.  Consequently, outside any sufficiently small
neighborhood of zero, $|\phi(t)|$ is uniformly less than $1$.

The exponential moment gives, on a branch of the logarithm near zero,
\[
 \log\bigl(e^{-i\mu t}\phi(t)\bigr)
       =-\frac{\sigma^2t^2}{2}+O(|t|^3).
\]
For fixed sufficiently small $\eta>0$, it follows that, for
$|t|\le\eta$,
\[
 \left|(e^{-i\mu t}\phi(t))^k
             -e^{-k\sigma^2t^2/2}\right|
       \le Ck|t|^3e^{-c kt^2}.                      \tag{7}
\]
For example, interpolate between the two exponents and use that each
interpolated real part is at most $-c kt^2$.  The integral of the
right side of (7) is $O(k^{-1})$.  The part $|t|>\eta$ is exponentially
small.  Extending the Gaussian integral to the real line in (6) now
gives the uniform estimate
\[
 \Prob(S_k=n)=
 \frac{d}{\sigma\sqrt{2\pi k}}
 \exp\!\left(-\frac{(n-k\mu)^2}{2k\sigma^2}\right)
       +O(k^{-1})                                  \tag{8}
\]
for every admissible $n$.  The error is uniform in $n$, since its only
additional Fourier factor has modulus one.

We next control all $k$ outside a central window.  Since the moment
generating function is finite near zero, there are constants $K,t_0>0$
such that
\[
 \log\E e^{t(X-\mu)}\le Kt^2\qquad(|t|\le t_0).
\]
Exponential Markov inequalities with the sign of $t$ chosen to match
$n-k\mu$, and with
$|t|=\min\{t_0,|n-k\mu|/(2Kk)\}$, imply
\[
 \Prob(S_k=n)\le
 \exp\!\left[-c\min\left\{
       \frac{(n-k\mu)^2}{k},\ |n-k\mu|\right\}\right]
                                                        \tag{9}
\]
for a fixed $c>0$.  Only $1\le k\le n$ can contribute when $n>0$.
Fix $0<\delta<1/6$.  If
$|k-n/\mu|>n^{1/2+\delta}$, (9) is at most
$\exp(-c_1n^{2\delta})$, uniformly over these $k$.  Thus their total
contribution to the left side of (5) is exponentially small in
$n^{2\delta}$.

Inside the window $|k-n/\mu|\le n^{1/2+\delta}$, one has $k\asymp n$.
The Gaussian term in (8) is $O(n^{-1/2})$, as is the probability.
For nonnegative $u,v$,
$|u^m-v^m|\le m|u-v|\max(u,v)^{m-1}$.
Hence replacing the probabilities in the $m$th-power sum by their
Gaussian terms changes the sum by
\[
 O\!\left(n^{1/2+\delta}n^{-1}n^{-(m-1)/2}\right)
   =O(n^{-m/2+\delta})
   =o(n^{-(m-1)/2}).                               \tag{10}
\]
Write $s=(k-n/\mu)/\sqrt n$.  In this window the $m$th power of the
Gaussian term is, uniformly,
\[
 \left(\frac{d\sqrt\mu}{\sigma\sqrt{2\pi n}}\right)^m
 \exp\!\left(-\frac{m\mu^3s^2}{2\sigma^2}\right)(1+o(1)).
                                                        \tag{11}
\]
Indeed $k^{-m/2}=(n/\mu)^{-m/2}(1+o(1))$, and the error in replacing
$k$ by $n/\mu$ in the exponent is
$O(|k-n/\mu|^3/n^2)=O(n^{-1/2+3\delta})=o(1)$.
The mesh of the permitted values of $s$ is $d/\sqrt n$.
Their translate can depend on $n$, which does not affect the Riemann
sum limit.  The Gaussian tails are integrable and the windows expand to
the whole real line, so summing (11) gives
\[
 \left(\frac{d\sqrt\mu}{\sigma\sqrt{2\pi n}}\right)^m
 \frac{\sqrt n}{d}
 \int_{\mathbb R}\exp\!\left(-\frac{m\mu^3s^2}{2\sigma^2}\right)ds
 \,(1+o(1)).
\]
The integral equals $\sqrt{2\pi\sigma^2/(m\mu^3)}$.
Together with (9)--(10), this proves (5).
\end{proof}

\subsection{The algebraic-function obstruction}

We use two standard arithmetic facts: $\pi$ is transcendental, and an
algebraic function over $\Qbar(t)$ has a convergent local Puiseux
expansion at an algebraic point with coefficients in $\Qbar$.  The
second fact is the Newton--Puiseux theorem.  The following argument
shows precisely how they apply, without assumptions about other
singularities on the circle of convergence.

\begin{lemma}\label{add43taisho:r43:obstruction}
Let $b_n\in\Qbar$ and suppose
\[
 b_n\sim A\pi^{-(m-1)/2}n^{-(m-1)/2},
 \qquad A>0,\quad A\in\Qbar,
 \qquad m\ge3\text{ an integer}.
                                                        \tag{12}
\]
Then $B(t)=\sum_{n\ge0}b_nt^n$ is transcendental over $\mathbb C(t)$.
\end{lemma}

\begin{proof}
First, if a power series with algebraic coefficients is algebraic over
$\mathbb C(t)$, it is algebraic over $\Qbar(t)$.  To prove this, take
the finite collection of monomials $t^iB(t)^j$ occurring in a nonzero
complex polynomial relation.  Their coefficient vectors have entries
in $\Qbar$.  The rank of this finite-column, possibly infinite-row
matrix is the same over $\Qbar$ as over $\mathbb C$: rank is witnessed
by a finite nonzero minor, and a maximal independent finite collection
of rows spans all rows over either field.  A nontrivial complex kernel
therefore gives a nontrivial kernel over $\Qbar$.
Also, derivatives of an algebraic function are algebraic in
characteristic zero: differentiating its irreducible polynomial
equation expresses the derivative as a quotient in its algebraic
function field.  Thus, if $B$ were algebraic, every iterate of
$\theta=t\,d/dt$ applied to it would be algebraic over $\Qbar(t)$.

Suppose first that $m=2r+1$ is odd, $r\ge1$, and set
$G=\theta^{r-1}B$.  Its coefficients satisfy
$[t^n]G\sim c/n$, where $c=A\pi^{-r}>0$.  Since
$\sum_{n\ge1}t^n/n=\log(1/(1-t))$, comparison of tails gives
\[
 G(t)\sim c\log\frac1{1-t}\qquad(t\uparrow1).       \tag{13}
\]
Explicitly, after a finite initial segment the difference between the
two coefficient sequences is bounded in absolute value by
$\varepsilon/n$.  The finite segment is bounded on $0<t<1$ and the
tail is at most $\varepsilon\log(1/(1-t))$; then let
$\varepsilon\downarrow0$.
A local Puiseux expansion is either bounded at $1$, or has a first
negative rational exponent and grows as a nonzero constant times a
power of $(1-t)^{-1}$.  Neither behavior is (13), a contradiction.

Now let $m=2r$ be even, $r\ge2$, and set
$G=\theta^{r-1}B$.  Then
$[t^n]G\sim c n^{-1/2}$ with $c=A\pi^{1/2-r}$.
The coefficients of $(1-t)^{-1/2}$ are
$\binom{2n}{n}4^{-n}\sim(\pi n)^{-1/2}$.  The same tail comparison,
now multiplied by $(1-t)^{1/2}$, yields
\[
 \lim_{t\uparrow1}(1-t)^{1/2}G(t)
        =c\sqrt\pi=A\pi^{1-r}.                    \tag{14}
\]
If $G$ were algebraic over $\Qbar(t)$, its Puiseux expansion at the
algebraic point $1$ would make the nonzero finite limit in (14) an
algebraic leading coefficient, corresponding to exponent $-1/2$.
But $A\pi^{1-r}$ is transcendental.  For if it were algebraic, the
positive integer power $\pi^{r-1}$ would be algebraic, forcing $\pi$
itself to be algebraic.  This contradiction proves the even case too.
\end{proof}

\subsection{Completion for every rational alphabet}

\begin{proof}[Proof of Theorem~\ref{add43taisho:r43:main}]
First suppose $\gcd\mathcal P=1$, and let $d$ be the difference span
in Lemma~\ref{add43taisho:r43:llt}.  Equations (3)--(5) give
\[
 \rho^{mn}D_n^{(m)}
 \sim A\pi^{-(m-1)/2}n^{-(m-1)/2},\qquad
 A=\frac{d^{m-1}\mu^{(m-3)/2}}
         {\sqrt m\,(2\sigma^2)^{(m-1)/2}}>0.        \tag{15}
\]
All factors in $A$ are algebraic, including the positive square roots.
The series $B(t)=D^{(m)}(\rho^m t)$ has algebraic coefficients.
Lemma~\ref{add43taisho:r43:obstruction} makes $B$, and hence $D^{(m)}$,
transcendental.

For a general support, put $g=\gcd\mathcal P$.  The coefficients
$D_n^{(m)}$ vanish unless $g\mid n$.  Apply Lemma~\ref{add43taisho:r43:llt} to
$Y=X/g$, whose supported values have greatest common divisor $1$.
Its mean is $\mu/g$, its variance is $\sigma^2/g^2$, it has an
exponential moment, and it is still nonconstant.  Equation (4), with
$n=gN$, gives
\[
 \rho^{mgN}D_{gN}^{(m)}
 \sim A_g\pi^{-(m-1)/2}N^{-(m-1)/2},               \tag{16}
\]
where $A_g>0$ is algebraic by the same explicit formula, using the
mean, variance, and difference span of $Y$.
Consequently $E(t)=\sum_{N\ge0}D_{gN}^{(m)}t^N$ is transcendental,
by Lemma~\ref{add43taisho:r43:obstruction} and the algebraic rescaling
$t\mapsto\rho^{mg}t$.
Finally $D^{(m)}(z)=E(z^g)$.  If $D^{(m)}$ were algebraic over
$\mathbb C(z)$, substitution $z=t^{1/g}$ would make $E(t)$ algebraic
over $\mathbb C(t^{1/g})$, hence over $\mathbb C(t)$ by transitivity
of algebraic extensions.  This contradicts the preceding conclusion.
\end{proof}

\paragraph{What the completion adds.}
The result is uniform over the original class of rational part series
and over every integer tuple size $m>2$.  The proof does not restrict
the alphabet to $\{1,2\}$, does not assume difference span $1$, and
does not use numerical evidence to establish nonalgebraicity.  Every
analytic estimate and every reduction needed for these quantifiers is
given above.  The established Puiseux-amplitude method and the prior
composition asymptotics are attributed to their sources.

\begin{thebibliography}{9}
\bibitem{add43taisho:r43:BH}
C.~Banderier and P.~Hitczenko,
\emph{Enumeration and asymptotics of restricted compositions having the
same number of parts}, Discrete Applied Mathematics \textbf{160}
(2012), 2542--2554.
\href{https://doi.org/10.1016/j.dam.2011.12.011}{doi:10.1016/j.dam.2011.12.011}.
Author manuscript: \url{https://arxiv.org/abs/1201.6116}.

\bibitem{add43taisho:r43:BD}
C.~Banderier and M.~Drmota,
\emph{Formulae and asymptotics for coefficients of algebraic functions},
Combinatorics, Probability and Computing \textbf{24} (2015), 1--53.
\href{https://doi.org/10.1017/S0963548314000728}{doi:10.1017/S0963548314000728}.
The author manuscript, revised 22 March 2014, contains the explicit
Franel/composition arithmetic-amplitude discussion in Proposition~7:
\url{https://lipn.fr/~banderier/Papers/Alg.pdf}.
\end{thebibliography}

\par\endgroup

% Additional source: /workspace/scratch/81914cec2c78/work/additional_root/r43_pathak_sign_density.tex
\begingroup
\newcommand{\PP}{\mathbb P}
\newcommand{\EE}{\mathbb E}
\newcommand{\FF}{\mathcal F}
\clearpage
\section{\texorpdfstring{A continuous limiting law and a computable sign density for the odd primitive Pathak invariant}{A continuous limiting law and a computable sign density for the odd primitive Pathak invariant}}\label{section:add43reiwa}
\noindent{\small\sffamily\color{muted}Additional completion \enspace / \enspace Round 43 \enspace / \enspace Reiwa}\par\medskip


\begin{sectionabstract}
For negative fundamental discriminants, the invariant in the archived
Round 43 Reiwa problem has a continuous limiting distribution. Its
negative values have a natural density strictly between zero and one.
We specify that density by finite rational sums with certified upper
and lower bounds converging to the same number. The missing
anti-concentration step is proved using a prime-flip inequality and
Sperner's theorem. The result concerns the odd primitive population;
it makes no assertion about the separate imprimitive-even population.
\end{sectionabstract}

\subsection{The precise population and result}
Let
\[
 \FF(X)=\{D<0: D\text{ is a fundamental discriminant},\ |D|\le X\},
 \qquad \chi_D(n)=\left(\frac{D}{n}\right).
\]
Thus the characters are real, primitive, and odd, ordered by their
conductors. Define the invariant appearing in the archived sign-density
question, attributed there to Pathak \cite{add43reiwa:pathak}, by
\begin{equation}\label{add43reiwa:Idef}
 I_D=\frac12+\sum_{r=0}^{\infty}(-1)^r
                   \bigl(L(2r+1,\chi_D)-1\bigr).
\end{equation}
The term with $r=0$ is interpreted by the convergent Dirichlet series
for $L(1,\chi_D)$. All remaining terms can be summed absolutely.
Indeed,
\[
 \sum_{r\ge1}\sum_{n\ge2}n^{-2r-1}
  =\sum_{n\ge2}\frac1{n(n^2-1)}=\frac14.
\]
Consequently
\begin{equation}\label{add43reiwa:Iidentity}
 I_D=L(1,\chi_D)-H_D,\qquad
 H_D=\sum_{n\ge1}\frac{\chi_D(n)}{n(n^2+1)},
 \qquad \frac14<H_D<\frac34.
\end{equation}
Equivalently, $I_D=\sum_{n\ge1}\chi_D(n)n/(n^2+1)$, with the same
interpretation of its conditional convergence.

Let $\{X_p:p\text{ prime}\}$ be independent random variables with
\begin{equation}\label{add43reiwa:primelaw}
 \PP(X_p=0)=\frac1{p+1},\qquad
 \PP(X_p=1)=\PP(X_p=-1)=\frac{p}{2(p+1)}.
\end{equation}
Extend them completely multiplicatively to $X_n$, setting $X_1=1$.
In particular, a prime with value zero makes every positive multiple
of that prime have value zero. Define
\begin{equation}\label{add43reiwa:randomI}
 L=\prod_p(1-X_p/p)^{-1},\qquad
 H=\sum_{n\ge1}\frac{X_n}{n(n^2+1)},\qquad I=L-H.
\end{equation}
Convergence and positivity of the product will be proved below.
This is the standard quadratic-character prime law used by
Granville--Soundararajan \cite{add43reiwa:GS}; the proof will include the
finite-character approximation needed here.

\begin{theorem}\label{add43reiwa:main}
The empirical distributions of $I_D$, for $D\in\FF(X)$, converge
weakly to the law of $I$. The latter law has no atoms. In particular,
\[
 \lim_{X\to\infty}
 \frac{\#\{D\in\FF(X):I_D<0\}}{\#\FF(X)}
       =\delta:=\PP(I<0),\qquad 0<\delta<1,
\]
and the proportion with $I_D=0$ tends to zero. The number $\delta$
is specified by the explicit rational bounds in
Theorem~\ref{add43reiwa:compute}; those bounds give an algorithm to approximate
it to any prescribed accuracy.
\end{theorem}

This completes the limiting-law, nonzero-proportion, and evaluation
problem for the odd primitive population stated in the Round 43 note.
The evaluation is a convergent exact specification with rational
error bounds, not an asserted elementary closed form or an uncertified
decimal. The separate even-imprimitive question is not included in
the theorem. No claim of historical priority is made.

\subsection{The local character law}
For any fixed finite set of primes, the vector of values $\chi_D(p)$,
with $D$ uniformly distributed over $\FF(X)$, tends to the independent
law \eqref{add43reiwa:primelaw}. Here are the elementary counting details.

Write negative fundamental discriminants as either $D=-m$, with $m$
squarefree and $m\equiv3\pmod4$, or $D=-4m$, with $m$ squarefree
and $m\equiv1$ or $2\pmod4$. The usual squarefree sieve follows from
\[
 \mathbf1_{m\ \mathrm{squarefree}}=\sum_{a^2\mid m}\mu(a).
\]
For any fixed modulus, count each allowed residue class by substituting
$m=a^2b$ and applying the Chinese remainder theorem. The sum over
$a\le\sqrt X$ has an $O(\sqrt X)$ counting error (with a constant
depending on the fixed modulus); its main term is the product of
local proportions, since $\sum_a a^{-2}$ converges absolutely.
This proves uniformity among residue classes satisfying the fixed
local squarefree restrictions, together with independence at distinct
primes. It also gives
\begin{equation}\label{add43reiwa:count}
                    \#\FF(X)\sim\frac{3}{\pi^2}X.
\end{equation}

For an odd prime $p$, there are $p^2-1$ allowed residue classes
modulo $p^2$. Exactly $p-1$ of them are divisible by $p$. Among the
remaining $p(p-1)$ classes, the quadratic residue symbol has equally
many values $+1$ and $-1$. Thus the three local proportions are
$1/(p+1),p/(2(p+1)),p/(2(p+1))$.

At the prime $2$, the odd discriminants have density $2/\pi^2$ and
the even discriminants have density $1/\pi^2$, measured per unit $X$.
Within the odd population, $D\equiv1$ and $5\pmod8$ occur equally
often, giving character values $+1$ and $-1$ respectively. The even
population gives zero. Division by \eqref{add43reiwa:count} therefore gives
probability $1/3$ for each value, exactly \eqref{add43reiwa:primelaw} at $p=2$.
The same sieve with the product modulus proves its independence
from any fixed collection of odd primes. Restricting the sign of $D$
only changes the relevant residue classes, not these proportions.

\subsection{Mean-square control at the conditional endpoint}
We first record a useful explicit summable bound. Let
\[
 T_N=\sum_{\substack{m,n>N\\mn\text{ a square}}}\frac1{mn},
 \qquad h_N=\sum_{j=1}^N\frac1j.
\]
Writing $m=su^2,n=sv^2$ with $s$ squarefree gives
\begin{equation}\label{add43reiwa:Tbound}
 T_N=\sum_{s\ \mathrm{squarefree}}\frac1{s^2}
       \left(\sum_{u>\sqrt{N/s}}\frac1{u^2}\right)^2
       \le\frac{4(h_N+1)}{N}.
\end{equation}
For $s\le N$, use $\sum_{u>t}u^{-2}\le2/t$ for $t\ge1$;
for $s>N$, use $\sum_{u\ge1}u^{-2}<2$ and
$\sum_{s>N}s^{-2}\le1/N$. These give the displayed bound directly.

\begin{lemma}\label{add43reiwa:arithmeticL2}
One has
\[
 \lim_{N\to\infty}\limsup_{X\to\infty}
 \frac1{\#\FF(X)}\sum_{D\in\FF(X)}
 \left|L(1,\chi_D)-\sum_{n\le N}\frac{\chi_D(n)}n\right|^2=0.
\]
\end{lemma}
\begin{proof}
Put $M=\lfloor X^{3/4}\rfloor$. The P\'olya--Vinogradov bound and
partial summation give, uniformly in the primitive characters of
conductor at most $X$,
\[
 \left|\sum_{n>M}\frac{\chi_D(n)}n\right|
       \ll\frac{\sqrt X\log(2X)}{M}=o(1).
\]
To bound the middle tail $N<n\le M$, enlarge the nonnegative sum of
squares over fundamental discriminants to the sum over all negative
integers $-X\le d<0$. On expansion its terms involve
\[
              \sum_{-X\le d<0}\left(\frac{d}{mn}\right).
\]
If $mn$ is not a square, the Kronecker symbol in its numerator is a
nonprincipal Dirichlet character with a modulus at most $8mn$.
P\'olya--Vinogradov, also for imprimitive characters, bounds this
sum by $O(\sqrt{mn}\log(2mn))$. For completeness, the imprimitive
case follows from the primitive one by inclusion-exclusion at the
additional prime divisors; the factor from those divisors is at most
$\prod_{p\mid q/f}2\le2\sqrt{q/f}$, with the harmless factor $2$
covering $p=2,3$. Thus the bound $O(\sqrt q\log(2q))$ suffices.
After weighting by $1/(mn)$ and dividing by $X$, the total nonsquare
contribution is
\[
 O\left(\frac{\log(2M)}X
          \left(\sum_{n\le M}n^{-1/2}\right)^2\right)
       =O\left(\frac{M\log(2M)}X\right)=o(1).
\]
For square $mn$, the trivial bound $X$ leaves at most $T_N$.
Equation \eqref{add43reiwa:Tbound}, the preceding uniform tail, and
\eqref{add43reiwa:count} prove the lemma.
\end{proof}

For the random model, independence and symmetry give
\[
 \EE(X_mX_n)=0\quad(mn\text{ nonsquare}),\qquad
 0\le\EE(X_mX_n)\le1\quad(mn\text{ square}).
\]
Therefore $\sum_{n\le N}X_n/n$ is a Cauchy sequence in $L^2$;
call its limit $L_*$. Moreover,
\begin{equation}\label{add43reiwa:randomL2}
 \EE\left|L_*-\sum_{n\le N}\frac{X_n}n\right|^2\le T_N.
\end{equation}

We identify this limit with the positive Euler product in
\eqref{add43reiwa:randomI}. Put $Y_p=-\log(1-X_p/p)$. The uniform expansion
$Y_p=X_p/p+O(p^{-2})$ implies
\[
 \sum_p|\EE Y_p|<\infty,\qquad
 \sum_p\operatorname{Var}(Y_p)<\infty.
\]
The independent centered sums converge almost surely and in $L^2$
by the usual convergence theorem for sums of independent variables
with summable variances. Hence $\sum_pY_p$ is finite almost surely,
and the product $L=\exp(\sum_pY_p)$ is strictly positive and finite.
For a fixed finite prime set the geometric expansion of its product
is absolutely convergent and equals the sum of $X_n/n$ over integers
having no other prime factors. As the prime set expands through all
$p\le y$, the difference from $L_*$ has mean square at most $T_y$:
all excluded integers exceed $y$, and the same nonnegative
square-product bound used in \eqref{add43reiwa:randomL2} applies. Finite
truncations justify that estimate before passage to the limit.
Thus the products converge to $L_*$ in $L^2$ as well as to $L$
almost surely, so $L_*=L$.

Define
\[
 I_N=\sum_{n\le N}\frac{nX_n}{n^2+1},\qquad
 I_{D,N}=\sum_{n\le N}\frac{n\chi_D(n)}{n^2+1}.
\]
The absolute $H$-tails are bounded by $\sum_{n>N}n^{-3}\le1/(2N^2)$.
It follows from \eqref{add43reiwa:randomL2} and $(a+b)^2\le2a^2+2b^2$ that
\begin{equation}\label{add43reiwa:Vbound}
 \EE|I-I_N|^2\le V_N,
 \qquad V_N=\frac{8(h_N+1)}N+\frac1{2N^4}.
\end{equation}
Lemma~\ref{add43reiwa:arithmeticL2} gives the analogous vanishing mean-square
error for $I_D-I_{D,N}$, with the ordered limits $X\to\infty$ and
then $N\to\infty$.

By the finite-character law, $I_{D,N}$ converges in distribution to
$I_N$ for each fixed $N$. For any bounded Lipschitz function $f$, the
difference between its averages on $I_D$ and $I_{D,N}$ is bounded
by the Lipschitz constant times the mean absolute truncation error;
Cauchy--Schwarz makes that error tend to zero in the ordered limits.
The same argument applies to $I$ and $I_N$ using \eqref{add43reiwa:Vbound}.
Taking $X\to\infty$ first and then $N\to\infty$ proves
\begin{equation}\label{add43reiwa:weak}
                         I_D\ \Longrightarrow\ I.
\end{equation}

\subsection{A prime-flip proof that the law has no atoms}
This is the step that allows one to pass from \eqref{add43reiwa:weak} to sign
events. Fix a prime $p$ and all variables other than $X_p$. Outside
a null set the Euler product converges and remains convergent after
any finite collection of coordinate changes. Use superscripts $+$
and $-$ for the two values obtained by setting $X_p=+1$ and $-1$.
The product formula gives exactly
\begin{equation}\label{add43reiwa:Lflip}
                 L^+-L^-=\frac{2L^-}{p-1}.
\end{equation}
Absolute convergence of $H$ gives
\begin{align}
 |H^+-H^-|
 &\le2\sum_{\substack{j\ge1\\j\text{ odd}}}
             \sum_{\substack{n\ge1\\p\nmid n}}
           \frac1{p^jn((p^jn)^2+1)} \notag\\
 &\le\frac{2\zeta(3)}{p^3(1-p^{-6})}
       <\frac4{p^3}.                              \label{add43reiwa:Hflip}
\end{align}
For the last inequality, $\zeta(3)<3/2$ and $p\ge2$ suffice.
Combining \eqref{add43reiwa:Lflip}--\eqref{add43reiwa:Hflip}, if $L^-\ge\varepsilon>0$
and $p^2>2/\varepsilon$, then
\begin{equation}\label{add43reiwa:strictflip}
                         I^+>I^-.
\end{equation}
Also $L$ increases under every such negative-to-positive prime flip.

Fix a real number $t$, a positive $\varepsilon$, and $N$ distinct
primes with $p^2>2/\varepsilon$. Condition on all variables outside
these primes and on which of these $N$ variables are zero. Suppose
$M$ of them are nonzero. Their remaining signs are independent fair
bits. The set of sign assignments satisfying
\[
                         I=t,\qquad L\ge\varepsilon
\]
is an antichain in the Boolean lattice of these $M$ signs. Indeed,
between two distinct comparable assignments, flip the differing
coordinates one at a time from negative to positive. Starting from
the smaller assignment, $L$ remains at least $\varepsilon$ and
\eqref{add43reiwa:strictflip} makes $I$ strictly increase. The two endpoints
therefore cannot both belong to the displayed set.

Sperner's bound makes its conditional probability at most
\[
                  2^{-M}\binom{M}{\lfloor M/2\rfloor}.
\]
One quick proof of the bound counts permutations of the $M$ coordinates:
a subset of size $j$ lies on $j!(M-j)!$ of the permutation chains,
and an antichain meets each chain at most once. Thus
$\sum_{A\text{ in antichain}}\binom{M}{|A|}^{-1}\le1$, which
implies the assertion. The binomial ratio tends to zero as
$M\to\infty$, for instance by the central-binomial asymptotic.

Each selected prime is nonzero with probability $p/(p+1)\ge2/3$,
independently. Thus $\EE M\ge2N/3$ and
$\operatorname{Var}(M)\le N/4$; Chebyshev's inequality gives
$M\to\infty$ in probability as $N\to\infty$. The displayed
conditional bound is at most one and tends to zero in probability,
so its expectation tends to zero. We conclude
\[
                   \PP(I=t,\ L\ge\varepsilon)=0.
\]
Finally $L>0$ almost surely, so taking the union over
$\varepsilon=1,1/2,1/3,\ldots$ proves $\PP(I=t)=0$ for every
$t\in\mathbb R$.

\subsection{Both signs have positive mass}
The inequalities in \eqref{add43reiwa:Iidentity} also hold for the random $H$:
\[
                         \frac14<H<\frac34.
\]
Thus $L<1/4$ implies $I<0$, and $L>3/4$ implies $I>0$.
Both events for $L$ have positive probability. To verify this without
an assumption about the Euler-product density, let
$Z_y=\sum_{p>y}Y_p$. Its mean and variance tend to zero, so
$\PP(|Z_y|<1)>1/2$ for all sufficiently large $y$.
Independence makes this event independent of all signs at $p\le y$.

The elementary Euler divergence $\sum_p1/p=\infty$ gives
\[
 \prod_{p\le y}(1+1/p)^{-1}\longrightarrow0,
 \qquad
 \prod_{p\le y}(1-1/p)^{-1}\longrightarrow\infty.
\]
For completeness, convergence of $\sum_p1/p$ would bound the latter
products, whereas their positive expansions contain every integer
$n\le y$ and hence dominate $\sum_{n\le y}1/n$, a contradiction;
the first conclusion then follows from
$\log(1+1/p)=1/p+O(p^{-2})$.

Choose a sufficiently large fixed $y$ so that the first finite
product is below $1/(4e)$. The event $X_p=-1$ for all $p\le y$
has positive probability, and its intersection with $|Z_y|<1$ gives
$L<1/4$. For the other sign, choose $y$ with the second product
above $3e/4$ and force $X_p=+1$ for $p\le y$. This gives $L>3/4$
on another positive-probability event. Consequently
\[
                            0<\PP(I<0)<1.
\]
Weak convergence \eqref{add43reiwa:weak} and the absence of an atom at zero
now prove the sign-density assertion in Theorem~\ref{add43reiwa:main}.
For the equality event, the closed-set inequality for weak convergence
gives $\limsup\PP_X(I_D=0)\le\PP(I=0)=0$.

\subsection{An exact, computable value for the density}
For $N\ge1$ and a rational number $a$, define the following finite
rational sum, whose index runs over the primes at most $N$:
\begin{equation}\label{add43reiwa:finiteformula}
 F_N(a)=
 \sum_{\boldsymbol\epsilon\in\{-1,0,1\}^{\{p:p\le N\}}}
 \mathbf1\left\{
   \sum_{n=1}^N\frac{n}{n^2+1}
       \prod_{p^j\parallel n}\epsilon_p^j<a\right\}
 \prod_{p\le N}w_p(\epsilon_p),
\end{equation}
where $w_p(0)=1/(p+1)$ and $w_p(\pm1)=p/(2(p+1))$.
The empty product at $n=1$ is one. Every comparison in this formula
is a comparison of rational numbers, and $F_N(a)=\PP(I_N<a)$.

\begin{theorem}\label{add43reiwa:compute}
For every integer $N\ge1$ and every positive rational $\eta$,
\begin{equation}\label{add43reiwa:brackets}
 \max\{0,F_N(-\eta)-V_N/\eta^2\}
       \ \le\ \delta\ \le\
 \min\{1,F_N(\eta)+V_N/\eta^2\},
\end{equation}
with the explicit rational $V_N$ in \eqref{add43reiwa:Vbound}. If
$\eta_N=1/\lceil N^{1/8}\rceil$, both bounds in
\eqref{add43reiwa:brackets} converge to $\delta$. In particular,
\begin{equation}\label{add43reiwa:deltalimit}
                         \delta=\lim_{N\to\infty}F_N(0).
\end{equation}
\end{theorem}
\begin{proof}
By \eqref{add43reiwa:Vbound} and Chebyshev,
$\PP(|I-I_N|\ge\eta)\le V_N/\eta^2$.
If $I_N<-\eta$ and $|I-I_N|<\eta$, then $I<0$;
if $I<0$ and $|I-I_N|<\eta$, then $I_N<\eta$.
These inclusions prove \eqref{add43reiwa:brackets}.
Now $\eta_N\to0$ and
$V_N/\eta_N^2=O((\log N)N^{-3/4})\to0$.
Since $I_N\to I$ in probability and $\PP(I=0)=0$, the quantities
$F_N(-\eta_N)$, $F_N(\eta_N)$, and $F_N(0)$ all tend to
$\PP(I<0)$. For example, a difference between one of these sign
events and $\{I<0\}$ implies either a truncation error larger than
a fixed $b>0$ or $|I|\le2b$, for all sufficiently large $N$;
first let $N\to\infty$ and then $b\downarrow0$.
This proves every asserted convergence. To compute to a prescribed
rational tolerance, enumerate \eqref{add43reiwa:finiteformula} for successive
$N$ and stop when the two rational endpoints in \eqref{add43reiwa:brackets}
are within that tolerance. Their proved convergence guarantees
termination; no rate of anti-concentration is assumed.
\end{proof}

The finite enumeration is an exact specification, not an efficiency
claim. No decimal estimate for $\delta$ is needed in any preceding
proof. The continuously distributed random invariant, the strict
positive proportions on both sides, and the rational converging
brackets together supply the missing global sign conclusion.

\begin{thebibliography}{9}
\bibitem{add43reiwa:pathak} S. Pathak,
\emph{On Berndt--Straub conjectures about zeros of generalized
Ramanujan polynomials}, arXiv:2610.18065v1 (2026).
\url{https://arxiv.org/abs/2610.18065}.
The population and invariant used here are those explicitly recorded
in the supplied Round 43 Reiwa manuscript and its source notes.
\bibitem{add43reiwa:GS} A. Granville and K. Soundararajan,
\emph{The distribution of values of $L(1,\chi_d)$},
Geometric and Functional Analysis \textbf{13} (2003), 992--1028.
\url{https://arxiv.org/abs/math/0206031}.
The quadratic-character local law and random Euler-product model are
standard inputs explained there; the truncation and anti-concentration
arguments needed for the present invariant are given in full above.
\end{thebibliography}

\par\endgroup

% Additional source: /workspace/scratch/81914cec2c78/work/additional_combinatorics/r47_all_ferrers_gaps_non_algebraicity.tex
\begingroup

\clearpage
\section{\texorpdfstring{Nonalgebraicity for every strict Ferrers pattern with adjacent positive parts}{Nonalgebraicity for every strict Ferrers pattern with adjacent positive parts}}\label{section:add47taisho}
\noindent{\small\sffamily\color{muted}Additional completion \enspace / \enspace Round 47 \enspace / \enspace Taisho}\par\medskip


\begin{sectionabstract}
For every strict partition containing two positive parts that differ by
one, its Ferrers-pattern avoidance generating function is transcendental.
An exact telescoping recurrence expresses the series as a rational
function plus a rational linear combination of prefix-pattern series
whose largest two parts are consecutive.  The prefix ending at the last
unit gap has strictly dominant logarithmic growth: an earlier prefix
has either a smaller pole exponent or a smaller logarithmic degree.
This prevents the cancellation left unresolved in the archived one-gap
argument and covers arbitrarily many unit gaps.  The proof explicitly
uses the established coefficient asymptotic of Bloom--McNew for
nonstaircase patterns with consecutive largest parts; the staircase
radial asymptotic is proved directly.  No priority claim is made.
\end{sectionabstract}

\subsection{The full target}
A partition contains a Ferrers pattern if deleting rows and columns of
its Ferrers diagram and then justifying the remaining boxes produces
that pattern.  Let a strict pattern have positive part sizes, listed in
increasing order,
\[
 0=v_0<v_1<\cdots<v_k,\qquad b_i=v_i-v_{i-1}\ge1.
\]
Let $F(b_1,\ldots,b_k;q)$ be the ordinary generating function, by area,
for partitions avoiding this pattern, including the empty partition.
We suppress $q$ in the argument list when convenient.  The first gap
$b_1$ is not a difference between two positive parts.

\begin{theorem}\label{add47taisho:r47:main}
For every finite list of positive integers $b_1,\ldots,b_k$, if
$b_i=1$ for at least one index $2\le i\le k$, then
$F(b_1,\ldots,b_k;q)$ is transcendental over $\mathbb C(q)$.
\end{theorem}

This is the full statement of Bloom--McNew Conjecture 1.2 for strict
patterns \cite{add47taisho:r47:BM}.  There is no restriction on the number, depth,
or separation of the unit gaps.  For comparison, their Corollary 4.17
settles the case $b_k=1$, and the archived Round 47 argument settled a
single unit gap at arbitrary depth.  The additional issue addressed
below is cancellation among several prefix series.

\subsection{Exact greedy elimination}
An occupied row length of a partition is a positive integer occurring
as one of its parts.  Containment of the pattern is equivalent to the
existence of occupied lengths
\[
 0=j_0<j_1<\cdots<j_k,
 \qquad j_i-j_{i-1}\ge b_i.                         \tag{1}
\]
For necessity, the retained columns between the two selected row ends
must number at least $b_i$.  For sufficiency, choose exactly $b_i$
columns in $(j_{i-1},j_i]$ for each $i$ and retain the chosen rows.
The greedy algorithm that chooses the least available $j_i$ satisfying
(1) is therefore an exact test: induction shows that each greedy
choice is no larger than the corresponding choice in any witness.

\begin{lemma}[Telescoping recurrence]\label{add47taisho:r47:rec}
For every possibly empty gap list $B$, and integers $b\ge1$, $c\ge2$,
\[
 F(B,b,c)=
 \frac{F(B,b+c-1)-q^{c-1}F(B,b)}{1-q^{c-1}}.        \tag{2}
\]
\end{lemma}

\begin{proof}
Partitions for which the greedy prefix $B$ is not matched avoid all
three patterns in (2); their common contribution satisfies the identity.
Otherwise fix all occupied sizes and multiplicities through the final
matched size $h$; for empty $B$, put $h=0$.  Let $K=h+b$.
The sizes $h+1,\ldots,K-1$ are unrestricted in all three cases, so
factor out their generating function as well as the fixed prefix weight.
The relative contribution for avoiding $(B,b)$ is $1$.
For avoiding $(B,b+c-1)$, it is
\[
 P(K),\qquad P(j)=\prod_{t=0}^{c-2}(1-q^{j+t})^{-1}.
\]
For avoiding $(B,b,c)$, either no size at least $K$ occurs, or the
first such size is some $j\ge K$.  In the latter case $j$ has positive
multiplicity, the sizes $j+1,\ldots,j+c-1$ have arbitrary
multiplicities, and no size at least $j+c$ occurs.  The relative series
is therefore
\[
 1+\sum_{j\ge K}
      \frac{q^j}{\prod_{t=0}^{c-1}(1-q^{j+t})}
 =1+\frac{P(K)-1}{1-q^{c-1}}.
\]
The equality follows from
\[
 P(j)-P(j+1)
   =\frac{q^j(1-q^{c-1})}
           {\prod_{t=0}^{c-1}(1-q^{j+t})}
\]
and $P(j)\to1$ coefficientwise as $j\to\infty$.
This is exactly (2) for the fixed prefix.  Summing over all prefixes
proves the identity of formal power series, and hence of analytic
functions for $0<q<1$.
\end{proof}

For $2\le i\le k$ with $b_i=1$, put
\[
 F_i(q)=F(b_1,\ldots,b_i;q),\qquad
 v_i=\sum_{h=1}^i b_h.
\]

\begin{proposition}\label{add47taisho:r47:decomp}
Let $I=\{i:2\le i\le k,\ b_i=1\}$, assumed nonempty, and
$j=\max I$.  There are rational functions $R(q)$ and $C_i(q)$ such that
\[
 F(b_1,\ldots,b_k;q)=R(q)+\sum_{i\in I}C_i(q)F_i(q),              \tag{3}
\]
where, as $q\uparrow1$,
\[
 C_i(q)=O((1-q)^{-(k-i)}).                         \tag{4}
\]
If $s=k-j$, $u_h=b_{j+h}-1\ge1$ and
$U_r=u_1+\cdots+u_r$, then the coefficient of the last prefix is
\[
 C_j(q)=(-1)^s\prod_{r=1}^s\frac{q^{U_r}}{1-q^{U_r}}.          \tag{5}
\]
For $s=0$, the empty product is $1$.
\end{proposition}

\begin{proof}
Apply (2) repeatedly whenever the final gap is at least $2$.  Stop a
branch either when its final gap is $1$ and its length is at least
$2$, or when its length is $1$.  In the latter case
\[
 F(b;q)=\prod_{a=1}^{b-1}(1-q^a)^{-1}                         \tag{6}
\]
is rational, since all occupied row lengths must be less than $b$.

Each application of (2) reduces the length of its gap list by one.
Merging its last two gaps replaces $b,c$ by $b+c-1\ge2$; hence it
never creates a unit gap.  All changes are at the right end, so every
branch stopping at a unit gap ends at one of the original, unchanged
prefixes $(b_1,\ldots,b_i)$ with $i\in I$.  This proves (3).
Along a path from length $k$ to prefix length $i$ there are exactly
$k-i$ recurrence steps.  Each multiplying rational factor is either
$(1-q^h)^{-1}$ or $-q^h(1-q^h)^{-1}$ with $h\ge1$, and thus has a
simple pole at $1$.  Summing the finitely many paths gives (4).

For (5), only paths retaining the last unit gap at $j$ matter.
Let $H_s(u_1,\ldots,u_s)$ be its coefficient after eliminating the
nonunit tail.  Clearly $H_0=1$.  For $s=1$, the merged branch absorbs
the unit gap and cannot produce this prefix, so
\[
 H_1(u)=-\frac{q^u}{1-q^u}.
\]
For $s\ge2$, recurrence (2) gives
\[
 H_s(u_1,\ldots,u_s)=
 \frac{H_{s-1}(u_1,\ldots,u_{s-2},u_{s-1}+u_s)
      -q^{u_s}H_{s-1}(u_1,\ldots,u_{s-1})}
      {1-q^{u_s}}.                                      \tag{7}
\]
Write $h(W)=-q^W/(1-q^W)$.  The elementary identity
\[
 \frac{h(W+v)-q^v h(W)}{1-q^v}=h(W)h(W+v)             \tag{8}
\]
proves by induction that $H_s=\prod_{r=1}^s h(U_r)$.
Thus (5) holds, and in particular
\[
 C_j(q)\sim
   \frac{(-1)^s}{U_1\cdots U_s}(1-q)^{-s}.            \tag{9}
\]
This last coefficient cannot vanish.
\end{proof}

\subsection{Radial growth of the prefix series}
For $i\in I$, let $\ell_i\ge1$ be the number of consecutive unit gaps
ending at $i$, counting only gaps with index at least $2$.
Equivalently, $F_i$ has $\ell_i+1$ consecutive largest positive
pattern parts, and this is the maximal such initial run in decreasing
order.  Put $L(q)=\log(1/(1-q))$.

\begin{lemma}\label{add47taisho:r47:prefix}
For every $i\in I$ there is a constant $\gamma_i>0$ such that
\[
 F_i(q)\sim\gamma_i(1-q)^{-(v_i-1)}L(q)^{\ell_i}
 \qquad(q\uparrow1).                               \tag{10}
\]
\end{lemma}

\begin{proof}
First suppose the prefix pattern is not the staircase
$(v_i,v_i-1,\ldots,1)$.  Write $K=v_i-1$ and $\ell=\ell_i$,
so that in decreasing order the pattern is
\[
 (K+1,K,\ldots,K-\ell+1,a_0,a_1,\ldots),\qquad
 K-\ell>a_0\ge0,
\]
padding its remaining parts by zeros.  The established
Bloom--McNew asymptotic \cite[Theorem 4.16]{add47taisho:r47:BM} is
\[
 [q^n]F_i(q)\sim
 \frac{n^{K-1}(\log n)^\ell}
 {\ell!\,(K-1)!\displaystyle\prod_{r=0}^{K-\ell-1}
                      (K-\ell-a_r-r)}.              \tag{11}
\]
Every denominator factor is positive.  This is the precise external
enumeration theorem used in the present proof.

For completeness, the Abelian passage needed here is elementary:
if $a_n\sim c n^{K-1}(\log n)^\ell$ with $K\ge1$, then, writing
$q=e^{-t}$ and summing a Riemann sum with $u=nt$,
\[
 \sum_{n\ge0}a_nq^n\sim
 c\,t^{-K}\bigl(\log(1/t)\bigr)^\ell
 \int_0^\infty u^{K-1}e^{-u}\,du
 =c(K-1)!\,t^{-K}\bigl(\log(1/t)\bigr)^\ell.        \tag{12}
\]
To justify the endpoints, the normalized integrands are bounded by a
constant times $u^{K-1}(1+|\log u|^\ell)e^{-u}$ away from a finite
initial set of indices; that bound is integrable.  The finite set is
negligible after multiplication by $t^K/(\log(1/t))^\ell$.
Since $t\sim1-q$, equations (11)--(12) prove (10), with
\[
 \gamma_i=
 \frac{1}{\ell!\displaystyle\prod_{r=0}^{K-\ell-1}
                       (K-\ell-a_r-r)}>0.            \tag{13}
\]

It remains to check the staircase, including its small cases.
Here $v_i=i$, all its gaps are $1$, and $\ell_i=i-1=K$.
By the greedy test, a partition avoids $(K+1,K,\ldots,1)$ precisely
when it has at most $K$ distinct occupied sizes.  Set
\[
 a_j(q)=\frac{q^j}{1-q^j},\qquad A(q)=\sum_{j\ge1}a_j(q).
\]
The generating function is
$\sum_{r=0}^K e_r(a_1(q),a_2(q),\ldots)$, where $e_r$ is the
$r$th elementary symmetric sum.  With $q=e^{-t}$,
\[
 A(e^{-t})=\sum_{j\ge1}\frac1{e^{tj}-1}
       =t^{-1}\log(1/t)+O(t^{-1}).                  \tag{14}
\]
Indeed, for $j\le1/t$ use $(e^{tj}-1)^{-1}=(tj)^{-1}+O(1)$,
and the remaining geometric tail is $O(t^{-1})$.
Also
\[
 \sum_{j\ge1}a_j(e^{-t})^2\le t^{-2}\sum_{j\ge1}j^{-2}
       =O(t^{-2}).                                  \tag{15}
\]
For any fixed $r\ge2$, the contribution of ordered $r$-tuples with a
repeated index is at most
$\binom r2 A^{r-2}\sum_j a_j^2$.  Equations (14)--(15) therefore give
$r!e_r\sim A^r$; for $r=1$ this equality is exact.
The term $r=K$ dominates the finite sum, so
\[
 F_i(q)\sim\frac1{K!}(1-q)^{-K}L(q)^K.
\]
This proves (10) in every staircase case as well.
\end{proof}

\subsection{The last unit gap cannot cancel}

\begin{proof}[Proof of Theorem~\ref{add47taisho:r47:main}]
Use Proposition~\ref{add47taisho:r47:decomp}, and put
\[
 e_i=v_i-1+k-i\qquad(i\in I).
\]
Combining its pole bound (4) with Lemma~\ref{add47taisho:r47:prefix} gives
\[
 C_i(q)F_i(q)
       =O\bigl((1-q)^{-e_i}L(q)^{\ell_i}\bigr).       \tag{16}
\]
For the last unit gap $j=\max I$, the exact coefficient (9) improves
this to
\[
 C_j(q)F_j(q)\sim
 \frac{(-1)^{k-j}\gamma_j}{U_1\cdots U_{k-j}}
       (1-q)^{-e_j}L(q)^{\ell_j},                   \tag{17}
\]
with the denominator interpreted as $1$ when $j=k$.

For every $i\in I$ with $i<j$,
\[
 e_j-e_i=(v_j-v_i)-(j-i)
             =\sum_{h=i+1}^{j}(b_h-1)\ge0.          \tag{18}
\]
If this is positive, (16) is smaller than (17) by at least one
positive power of $1-q$, which dominates any fixed logarithmic
power.  If it is zero, all $b_{i+1},\ldots,b_j$ equal $1$.
Then the consecutive run is extended at every step, and
\[
                  \ell_j=\ell_i+j-i>\ell_i.
\]
In this case the pole exponents agree but the logarithmic degree in
(16) is smaller.  Thus in either case
\[
 C_i(q)F_i(q)=o(C_j(q)F_j(q)).                       \tag{19}
\]
There are finitely many terms.  Subtracting the rational function
$R$ in (3) and using (17)--(19), we obtain
\[
 F(b_1,\ldots,b_k;q)-R(q)
 \sim C(1-q)^{-e_j}L(q)^{\ell_j},\qquad C\ne0,
 \quad\ell_j\ge1.                                  \tag{20}
\]

An algebraic function over $\mathbb C(q)$ has a local Puiseux
expansion at $q=1$.  A nonzero such function has a first term
$c(1-q)^\alpha$ with $c\ne0$ and $\alpha\in\mathbb Q$; a branch
unbounded at $1$ has $\alpha<0$.  No such leading term can have the
behavior (20): for every $\varepsilon>0$ the magnitude in (20) is
$o((1-q)^{-e_j-\varepsilon})$, while its ratio to
$(1-q)^{-e_j}$ tends to infinity.  Therefore $F-R$ is not algebraic.
Since $R$ is rational, $F$ is not algebraic either.
\end{proof}

\paragraph{Scope and dependency.}
The proof covers all strict nonsuper-strict patterns in Conjecture 1.2,
not only a one-gap subfamily.  Its only specialized input is the
already proved nonstaircase coefficient formula (11) of Bloom--McNew.
The exact recurrence, finite decomposition, pole bounds, nonzero
coefficient, staircase cases, Abelian step, and cancellation argument
have all been supplied here.  No finite computation is needed for the
theorem.

\begin{thebibliography}{9}
\bibitem{add47taisho:r47:BM}
J.~Bloom and N.~McNew,
\emph{Counting pattern-avoiding integer partitions},
arXiv:1908.03953 (2019, revised 2020).
\url{https://arxiv.org/abs/1908.03953}.
Primary author manuscript with Conjecture~1.2, Theorem~4.16, and
Corollary~4.17:
\url{https://www.nathanmcnew.com/partitions.pdf}.
\end{thebibliography}

\par\endgroup

% Additional source: /workspace/scratch/81914cec2c78/work/additional_analysis/gamma_median/paper.tex
\begingroup

\clearpage
\section{\texorpdfstring{Strict decrease of Lyon's scaled gamma-median offset}{Strict decrease of Lyon's scaled gamma-median offset}}\label{section:add47heisei}
\noindent{\small\sffamily\color{muted}Additional completion \enspace / \enspace Round 47 \enspace / \enspace Heisei}\par\medskip


\begin{sectionabstract}
Let $m(k)$ be the median of the gamma distribution with shape $k>0$
and unit scale. We prove that $2^{1/k}m(k)-k$ has strictly negative
derivative on the entire positive axis. Two analytic estimates cover
the small- and large-shape tails. On the intervening compact interval,
convexity of the median reduces the derivative inequality to 1,344
explicit scalar comparisons. We verify all of them by directed integer
arithmetic, with proved remainder bounds for every infinite series.
The complete verifier is included.
\end{sectionabstract}

\subsection{The statement and its published inputs}
The median $m(k)$ is the unique solution of
\[
 P(k,m(k))=\frac12,\qquad
 P(k,t)=\frac1{\Gamma(k)}\int_0^t u^{k-1}e^{-u}\,du.
\]
Lyon \cite[``Theorems to prove'']{add47heisei:lyon} distinguished the two-sided
endpoint bounds for the function below from its conjectured monotonicity.
The result proved here is the monotonicity assertion itself, for every
positive shape.
\begin{theorem}\label{add47heisei:thm:global}
For every $k>0$,
\[
 A(k)=2^{1/k}m(k)-k\qquad\text{satisfies}\qquad A'(k)<0.
\]
\end{theorem}
We use the established gamma-median facts
\begin{equation}\label{add47heisei:eq:known}
 k-\frac13<m(k)<k,\qquad 0<m'(k)<1,\qquad m''(k)>0.
\end{equation}
The first bounds and the derivative theorem are recorded and proved in
Berg--Pedersen \cite[Introduction and Theorem 3.3]{add47heisei:bp1}; strict convexity
is their subsequent Theorem 1.1 \cite{add47heisei:bp2}. These inputs hold on all of
$(0,\infty)$. The implicit-function theorem also gives the differentiability
used below, since the gamma density is positive at its median.

Put $c=\log 2$. Direct differentiation gives
\begin{equation}\label{add47heisei:eq:derivative}
 A'(k)=e^{c/k}
 \left(m'(k)-\frac{cm(k)}{k^2}-e^{-c/k}\right).
\end{equation}
We will prove that the bracket is negative on three intervals whose union
is $(0,\infty)$.

\subsection{The large-shape tail}
For $x\ge0$, Taylor's theorem gives
$1-e^{-x}\le x-x^2/2+x^3/6$. Hence, for $k\ge7$,
\[
 \frac{k^2}{c}(1-e^{-c/k})
 \le k-\frac c2+\frac{c^2}{6k}
 <k-\frac13<m(k).
\]
For the strict middle inequality it suffices to use
$69/100<c<7/10$, because
\[
 \frac{69}{200}-\frac{49}{4200}=\frac13.
\]
Together with $m'(k)<1$, this proves that the bracket in
\eqref{add47heisei:eq:derivative} is strictly negative for $k\ge7$.

\subsection{The small-shape tail}
Throughout this section $0<k\le1/8$. Define
\[
 s(k)=e^{c/k}m(k),\qquad
 G(k)=\Gamma(1+k)^{1/k}.
\]
If $U$ has density $ku^{k-1}$ on $(0,1)$, put
\[
 H(k,t)=\mathbb E e^{-tU},\qquad
 h(k,t)=-\frac1k\log H(k,t).
\]
Rescaling the median integral gives
\begin{equation}\label{add47heisei:eq:implicit}
 s(k)=G(k)\exp(h(k,m(k))).
\end{equation}
Jensen's inequality gives
$0\le h(k,t)\le t/(1+k)$.
Log-convexity of $\Gamma$, between $1$ and $3/2$, gives
$G(k)\le\Gamma(3/2)^2=\pi/4$. Thus
\[
 0<s(k)\le\frac\pi4e^{1/8}<\frac9{10}.
\]
Here $e^{1/8}<8/7$ and $\pi<22/7$ already suffice.
In particular,
\begin{equation}\label{add47heisei:eq:smallestimates}
 m(k)<\frac{9}{2560}<\frac1{250},\qquad
 \frac{m(k)}{k^2}<\frac9{40}.
\end{equation}
For the second inequality, $e^{-c/k}/k^2$ is increasing on $(0,1/8]$,
since its logarithmic derivative is $(c-2k)/k^2>0$.
At $k=1/8$ its value is $1/4$.

We next bound the partial derivatives of $h$. Exponential tilting by
$e^{-tU}$ decreases the mean of $U$, so
\begin{equation}\label{add47heisei:eq:ht}
 0\le h_t(k,t)=\frac1k\mathbb E_t U
 \le\frac1{1+k}\le1.
\end{equation}
The assertion about the mean also follows directly by taking independent
copies $U,V$ and observing
$\mathbb E[(U-V)(e^{-tU}-e^{-tV})]\le0$.
Set
\[
 q(k,t)=\frac{1-H(k,t)}{k}
       =\int_0^1(1-e^{-tu})u^{k-1}\,du.
\]
At fixed $t$, $q_k\le0$, and $0\le q\le t/(1+k)$. Since
\[
 h=q\int_0^1\frac{dv}{1-kqv},
\]
this expression is increasing in $q$. Differentiating first at fixed $q$
and then using $q_k\le0$ yields
\[
 h_k(k,t)\le\frac{q^2}{2(1-kq)^2}.
\]
At $t=m(k)$, \eqref{add47heisei:eq:smallestimates} therefore gives $h_k<1/1000$.
All these differentiations under the integral are justified locally for
$k>0$ by the integrability of $u^{k-1}|\log u|$.

Finally, with $B=(\log G)'$,
\[
 B(k)=\frac1{k^2}\int_0^k u\psi'(1+u)\,du
 \le\frac{\pi^2}{12}<\frac56.
\]
The inequality follows from
$\psi'(z)=\sum_{j\ge0}(z+j)^{-2}$ and $\psi'(1)=\pi^2/6$.
Differentiate \eqref{add47heisei:eq:implicit} and use
$m'=e^{-c/k}s'+cm/k^2$. One obtains
\[
 (1-s h_t e^{-c/k})s'
 =s\left(B+h_k+h_t\frac{cm}{k^2}\right).
\]
The coefficient on the left is greater than $1-9/2560>0$. The right side
is strictly less than
\[
 \frac9{10}\left(\frac56+\frac1{1000}
                         +\frac7{10}\frac9{40}\right)
 <1-\frac9{2560}.
\]
Consequently $s'<1$, or $A'<0$, on $(0,1/8]$.

\subsection{A compact reduction that controls every shape}
Let $[l,r]$ have length $1/d$, and let $k\in[l,r]$.
Convexity of $m$ implies
\[
 m'(k)\le m'(r)
 \le d\bigl(m(r+1/d)-m(r)\bigr).
\]
Monotonicity of $m$ and of $e^{-c/k}$ also gives
\begin{equation}\label{add47heisei:eq:cell}
 m'(k)-\frac{cm(k)}{k^2}-e^{-c/k}
 \le D(l,r,d):=
 d\bigl(m(r+1/d)-m(r)\bigr)
 -\frac{cm(l)}{r^2}-e^{-c/l}.
\end{equation}
Thus a bound for one scalar $D$ proves the desired sign throughout an
entire closed interval; it is not an inference from sampled derivative
values.

\begin{lemma}\label{add47heisei:lem:certificate}
For the uniform subdivisions in the following table, every cell has
$D$ strictly smaller than the stated rational upper bound.
\[
\begin{array}{c|r|r|c}
\text{block}&d&\text{number of cells}&\text{upper bound for every }D\\\hline
{[1/8,1/4]}&4096&512&-29/100000\\
{[1/4,1/2]}&512&128&-7/5000\\
{[1/2,1]}&256&128&-1/125\\
{[1,2]}&128&128&-1/500\\
{[2,4]}&128&256&-1/1250\\
{[4,7]}&64&192&-9/100000
\end{array}
\]
\end{lemma}
\begin{proof}
The complete integer verifier below proves these 1,344 inequalities.
We give the enclosure formulas and error bounds here so that its checks
have a specified mathematical meaning.

Every interval is represented as $[L/2^{160},U/2^{160}]$ with integer
endpoints. Addition and subtraction are exact on endpoints; multiplication,
division by a positive interval, and multiplication by a rational round
outwards using integer floor and ceiling. No numerical library supplies a
gamma value, logarithm, exponential, median, or sign.

\paragraph{Logarithms and $\pi$.}
For $1\le x\le2$ put $z=(x-1)/(x+1)$. The identity
\[
 \log x=2\sum_{j=0}^{N-1}\frac{z^{2j+1}}{2j+1}+R_N,
 \qquad
 0\le R_N\le\frac{2z^{2N+1}}{(2N+1)(1-z^2)}
\]
is the integrated geometric series. We use $N=96$ and first scale every
positive argument by a power of two into $[1,2]$. The same formula gives
$c=\log2$. Machin's identity
$\pi=16\arctan(1/5)-4\arctan(1/239)$, with 80 alternating-series terms
for each arctangent, gives an interval for $\pi$; the first omitted term
bounds each remainder. These bounds also verify $69/100<c<7/10$.

\paragraph{Gamma normalization.}
For every shape to be checked, set $z=k+33$ and use
\[
 \log\Gamma(k+1)=\log\Gamma(z)-\sum_{j=1}^{32}\log(k+j).
\]
Stirling's remainder theorem \cite[\S5.11(ii)]{add47heisei:dlmf} gives
\begin{align*}
 \log\Gamma(z)
 &=\left(z-\tfrac12\right)\log z-z+\tfrac12\log(2\pi)
   +\sum_{j=1}^{10}\frac{B_{2j}}{2j(2j-1)z^{2j-1}}+R,\\
 0&<R<\frac{B_{22}}{22\cdot21\,z^{21}}.
\end{align*}
This signed first-omitted-term bound is valid for every positive real
$z$; the rational Bernoulli numbers through $B_{22}$ are listed in the
verifier. Here $z\ge33+1/8$.

\paragraph{The incomplete-gamma series and median brackets.}
Expanding $e^{t-u}$ inside the gamma integral and integrating the
nonnegative terms gives
\[
 P(k,t)=\frac{e^{-t}t^k}{\Gamma(k+1)}
 \sum_{j=0}^{\infty}a_j,
 \qquad a_0=1,\quad a_j=\frac{t^j}{(k+1)\cdots(k+j)}.
\]
For $0<t<8$ the tail after $j=160$ is at most
\[
 \frac{a_{161}}{1-t/(k+162)}.
\]
This follows because each subsequent term ratio is no greater than
$t/(k+162)<1$. Taking the logarithm of the bounded sum gives an interval
for
\[
 \log P(k,t)=k\log t-t-\log\Gamma(k+1)
                  +\log\left(\sum_{j\ge0}a_j\right).
\]
Start with $0<m(k)<k$ and perform 50 bisections. At each midpoint the
program proves either $\log P(k,t)+c<0$ or $\log P(k,t)+c>0$; an
unresolved sign would stop the program. Strict increase of the CDF in
$t$ therefore proves every resulting median enclosure. The program
encloses 1,350 distinct medians, including the right-hand secant points,
and verifies that every enclosure has width less than $10^{-12}$.

\paragraph{Exponentials.}
To enclose $e^{-x}$ for $0\le x\le8$, expand $e^{x/16}$ through degree
80. The positive remainder is bounded by its degree-81 term divided by
$1-(x/16)/82$. Four interval squarings and one reciprocal then enclose
$e^{-x}$. This applies to $x=c/l$ in all six blocks.

\paragraph{Final integer inequalities.}
Use the resulting intervals for every quantity in \eqref{add47heisei:eq:cell}.
If $U/2^{160}$ is an upper endpoint for a cell and the advertised bound
is $a/b$, the check is the integer inequality $bU<a2^{160}$.
All 1,344 such checks pass. The following coarser displayed fractions
are also strict upper bounds for the block maxima, rounded upward to
multiples of $10^{-9}$:
\[
\begin{gathered}
 -\frac{147357}{500000000},\quad
 -\frac{1482823}{1000000000},\quad
 -\frac{8290913}{1000000000},\\[3pt]
 -\frac{2313009}{1000000000},\quad
 -\frac{812639}{1000000000},\quad
 -\frac{47257}{500000000}.
\end{gathered}
\]
\end{proof}

\begin{proof}[Proof of Theorem~\ref{add47heisei:thm:global}]
The analytic arguments prove the derivative sign on $(0,1/8]$ and
$[7,\infty)$. Lemma~\ref{add47heisei:lem:certificate}, combined with the convexity
bound \eqref{add47heisei:eq:cell}, proves it at every point of $[1/8,7]$.
Equation~\eqref{add47heisei:eq:derivative} completes the proof.
\end{proof}

\subsection{Exact verifier}
The following is complete Python 3 code. It uses only arbitrary-precision
integers and the standard-library rational type. Its finite loops and
assertions implement the bounds proved above; the final output is
\texttt{CERTIFIED: 1344 cells; 1350 median enclosures; integer arithmetic only.}

\begin{Verbatim}
"""Exact certificate for A'(k)<0 on [1/8,7]. Python 3, standard library.

All numerical enclosures use integers with common denominator 2**160.
No binary floating-point or special-function library is used.
The accompanying paper proves the series and all remainder bounds.
"""
from fractions import Fraction

S = 1 << 160


def ceildiv(a, b):
    assert b > 0
    return -((-a) // b)


class I:
    def __init__(self, lo, hi=None):
        self.lo = lo
        self.hi = lo if hi is None else hi
        assert self.lo <= self.hi

    def __add__(self, other):
        return I(self.lo + other.lo, self.hi + other.hi)

    def __neg__(self):
        return I(-self.hi, -self.lo)

    def __sub__(self, other):
        return self + (-other)

    def __mul__(self, other):
        p = [self.lo * other.lo, self.lo * other.hi,
             self.hi * other.lo, self.hi * other.hi]
        return I(min(p) // S, ceildiv(max(p), S))

    def scale(self, n, d=1):
        assert d > 0
        p = [n * self.lo, n * self.hi]
        return I(min(p) // d, ceildiv(max(p), d))

    def __truediv__(self, other):
        assert other.lo > 0
        reciprocal = I(S * S // other.hi,
                       ceildiv(S * S, other.lo))
        return self * reciprocal


def Q(n, d=1):
    assert d > 0
    return I(n * S // d, ceildiv(n * S, d))


ZERO, ONE = Q(0), Q(1)


def log_near_one(n, d):
    """Enclose log(n/d) for 1 <= n/d <= 2 by 96 atanh terms."""
    assert d <= n <= 2 * d
    z = Q(n - d, n + d)
    z2 = z * z
    power, value = z, ZERO
    for j in range(96):
        value = value + power.scale(2, 2 * j + 1)
        power = power * z2
    # Tail <= 2*z**193 / (193*(1-z**2)).
    tail = power.scale(2, 193) / (ONE - z2)
    return I(value.lo, value.hi + tail.hi)


LOG2 = log_near_one(2, 1)
assert LOG2.lo > Q(69, 100).hi
assert LOG2.hi < Q(7, 10).lo


def log_point(n, d=S):
    assert n > 0 and d > 0
    exponent = 0
    while n < d:
        n *= 2
        exponent -= 1
    while n > 2 * d:
        d *= 2
        exponent += 1
    return log_near_one(n, d) + LOG2.scale(exponent)


def log_interval(x):
    return I(log_point(x.lo).lo, log_point(x.hi).hi)


def atan_reciprocal(q):
    """80 terms, alternating-series remainder, at x=1/q."""
    total = ZERO
    for j in range(80):
        total = total + Q((-1) ** j, (2 * j + 1) * q ** (2 * j + 1))
    # First omitted term is positive and bounds the remainder.
    return I(total.lo, total.hi + Q(1, 161 * q ** 161).hi)


PI = atan_reciprocal(5).scale(16) - atan_reciprocal(239).scale(4)
LOG2PI = log_interval(PI.scale(2))
BERNOULLI = [(1, 6), (-1, 30), (1, 42), (-1, 30), (5, 66),
             (-691, 2730), (7, 6), (-3617, 510), (43867, 798),
             (-174611, 330), (854513, 138)]


def log_gamma_plus_one(k):
    """Stirling at z=k+33, then 32 exact gamma recurrences."""
    z = k + Q(33)
    value = (z - Q(1, 2)) * log_interval(z) - z + LOG2PI.scale(1, 2)
    inverse = ONE / z
    square = inverse * inverse
    power = inverse
    for j, (n, d) in enumerate(BERNOULLI[:10], start=1):
        value = value + power.scale(n, d * (2 * j) * (2 * j - 1))
        power = power * square
    # B_22 > 0: remainder lies between zero and its next term.
    n, d = BERNOULLI[10]
    remainder = power.scale(n, d * 22 * 21)
    value = I(value.lo, value.hi + remainder.hi)
    for j in range(1, 33):
        value = value - log_interval(k + Q(j))
    return value


def log_cdf(k, t, log_gamma):
    """log P(k,t), using the positive incomplete-gamma series."""
    assert 0 < t.lo <= t.hi < 8 * S
    term, total = ONE, ONE
    for j in range(1, 161):
        term = term * t / (k + Q(j))
        total = total + term
    next_term = term * t / (k + Q(161))
    ratio = t / (k + Q(162))
    assert ratio.hi < S
    tail = next_term / (ONE - ratio)
    total = I(total.lo, total.hi + tail.hi)
    return k * log_interval(t) - t - log_gamma + log_interval(total)


def median(n, d):
    """50 rigorously signed bisections, starting from 0<m(k)<k."""
    k = Q(n, d)
    assert k.lo == k.hi
    log_gamma = log_gamma_plus_one(k)
    lo, hi = 0, k.hi
    for _ in range(50):
        mid = (lo + hi) // 2
        sign = log_cdf(k, I(mid), log_gamma) + LOG2
        if sign.hi < 0:
            lo = mid
        else:
            assert sign.lo > 0, (n, d, mid, sign.lo, sign.hi)
            hi = mid
    assert hi - lo < S // (10 ** 12)
    return I(lo, hi)


def exp_negative(x):
    """Enclose exp(-x) for 0<=x<=8; positive Taylor, then four squares."""
    assert 0 <= x.lo <= x.hi <= 8 * S
    y = x.scale(1, 16)
    term, total = ONE, ONE
    for j in range(1, 81):
        term = (term * y).scale(1, j)
        total = total + term
    first = (term * y).scale(1, 81)
    tail = first / (ONE - y.scale(1, 82))
    total = I(total.lo, total.hi + tail.hi)
    for _ in range(4):
        total = total * total
    return ONE / total


BLOCKS = [
    (Fraction(1, 8), Fraction(1, 4), 4096, Fraction(-29, 100000)),
    (Fraction(1, 4), Fraction(1, 2), 512, Fraction(-14, 10000)),
    (Fraction(1, 2), Fraction(1), 256, Fraction(-8, 1000)),
    (Fraction(1), Fraction(2), 128, Fraction(-2, 1000)),
    (Fraction(2), Fraction(4), 128, Fraction(-8, 10000)),
    (Fraction(4), Fraction(7), 64, Fraction(-9, 100000)),
]


def main():
    cells = 0
    cache = {}
    for left, right, d, claimed in BLOCKS:
        start = left * d
        count = (right - left) * d
        assert start.denominator == count.denominator == 1
        start, count = int(start), int(count)
        medians = []
        for i in range(count + 2):
            k = Fraction(start + i, d)
            if k not in cache:
                cache[k] = median(k.numerator, k.denominator)
            medians.append(cache[k])
        upper = -S * 100
        for i in range(count):
            l, r = Q(start + i, d), Q(start + i + 1, d)
            bound = ((medians[i + 2] - medians[i + 1]).scale(d)
                     - LOG2 * medians[i] / (r * r)
                     - exp_negative(LOG2 / l))
            assert bound.hi < 0
            # Prove the advertised rational upper bound as an integer check.
            assert bound.hi * claimed.denominator < claimed.numerator * S
            upper = max(upper, bound.hi)
        cells += count
        rounded = Fraction(ceildiv(upper * 10 ** 9, S), 10 ** 9)
        print(f"[{left},{right}]: {count} cells; max D < {rounded}; "
              f"claimed bound {claimed} verified", flush=True)
    assert cells == 1344
    print(f"CERTIFIED: {cells} cells; {len(cache)} median enclosures; "
          "integer arithmetic only.")


if __name__ == "__main__":
    main()
\end{Verbatim}


\begin{thebibliography}{9}
\bibitem{add47heisei:lyon} R. F. Lyon,
\emph{Tight bounds for the median of a gamma distribution},
PLoS ONE \textbf{18}(9) (2023), e0288601.
\href{https://doi.org/10.1371/journal.pone.0288601}{doi:10.1371/journal.pone.0288601}.
\bibitem{add47heisei:bp1} C. Berg and H. L. Pedersen,
\emph{The Chen--Rubin conjecture in a continuous setting},
\href{https://arxiv.org/abs/math/0411282}{arXiv:math/0411282}.
\bibitem{add47heisei:bp2} C. Berg and H. L. Pedersen,
\emph{Convexity of the median in the gamma distribution},
\href{https://arxiv.org/abs/math/0609442}{arXiv:math/0609442}, Theorem 1.1.
\bibitem{add47heisei:dlmf} NIST Digital Library of Mathematical Functions,
\href{https://dlmf.nist.gov/5.11}{\S5.11(i)--(ii)},
Stirling expansion (5.11.1) and its real positive remainder bound.
\end{thebibliography}

\par\endgroup

% Additional source: /workspace/scratch/81914cec2c78/work/additional_analysis/chernoff/paper.tex
\begingroup

\clearpage
\section{\texorpdfstring{The origin is the unique minimum of Chernoff density curvature}{The origin is the unique minimum of Chernoff density curvature}}\label{section:add49heisei}
\noindent{\small\sffamily\color{muted}Additional completion \enspace / \enspace Round 49 \enspace / \enspace Heisei}\par\medskip


\begin{sectionabstract}
For the standard Chernoff density $f$, we prove that
$\kappa(x)=-(\log f)''(x)$ has its unique global minimum at zero.
This gives the sharp form of Balabdaoui--Wellner's Conjecture 3.1.
An Airy convolution identity controls the unbounded tails. The remaining
compact comparison is proved by exact integer interval arithmetic with
explicit bounds for the Airy series, infinite trapezoidal rule, omitted
frequencies, and Taylor remainders. The complete verifier is included.
\end{sectionabstract}

\subsection{The sharp statement and its normalization}
Let $f$ denote the density of
\[
 Z=\operatorname*{argmax}_{t\in\mathbb R}\{W(t)-t^2\},
\]
where $W$ is standard two-sided Brownian motion. Put
$\kappa=-(\log f)''$.
\begin{theorem}\label{add49heisei:thm:sharp}
For all real $x$,
\[
 \kappa(x)\ge\kappa(0),
\]
with equality only at $x=0$. Moreover,
\[
 \frac{3405}{1000}<\kappa(0)<\frac{3406}{1000},\qquad
 \kappa(x)>\kappa(0)+\frac{x^2}{2}\quad(0<|x|\le1).
\]
\end{theorem}
The assertion resolves the sharp origin-curvature formulation in
Balabdaoui--Wellner \cite[Conjecture 3.1]{add49heisei:bw}; it is stronger than the
existence of an unspecified positive lower curvature bound.

We use their Airy factorization and log-concavity theorem
\cite[(1.3)--(2.5), Proposition 2.2]{add49heisei:bw}. In probability normalization,
there is a positive smooth log-concave density $g$ such that
\begin{equation}\label{add49heisei:eq:normalization}
 f(x)=C g(x)g(-x),\qquad
 M(s):=\int_{\mathbb R}e^{sx}g(x)\,dx
 =\frac{\operatorname{Ai}(0)}{\operatorname{Ai}(2^{-1/3}s)}.
\end{equation}
The positive constant $C$ has no effect on logarithmic derivatives.
Write
\[
 F(s)=\frac1{M(s)},\qquad v(x)=-(\log g)''(x).
\]
Then $F''(s)=sF(s)/2$, $F(0)=1$, and
\begin{equation}\label{add49heisei:eq:curvature}
 \kappa(x)=v(x)+v(-x).
\end{equation}
If $-\lambda_j$ are the zeros of $F$, arranged with
$0<\lambda_1<\lambda_2<\cdots$, its canonical product has the form
\begin{equation}\label{add49heisei:eq:product}
 M(s)=e^{ds}\prod_{j\ge1}\frac{e^{s/\lambda_j}}{1+s/\lambda_j},
 \qquad \sum_j\lambda_j^{-2}<\infty.
\end{equation}
All the zeros are negative real. These are standard Airy facts; see
\cite[Proposition 2.1]{add49heisei:bw} and \cite[\S\S9.2, 9.9]{add49heisei:dlmf}.
The transform is valid for $\Re s>-\lambda_1$.

\subsection{A complete analytic tail bound}
We record the Airy convolution argument, with the normalization needed
here; the identity also appears in Zhang--Zhou \cite[Lemma 7]{add49heisei:zz}.
\begin{lemma}\label{add49heisei:lem:tail}
For every real $x$,
\[
 v(x)\ge4x,\qquad v(x)\ge0,\qquad \kappa(x)\ge4|x|.
\]
\end{lemma}
\begin{proof}
Put $\Psi=\log M$ and
$r(u)=2u\sum_{j\ge1}e^{-\lambda_j u}$ for $u>0$.
The product gives
\[
 \Psi''(s)=\sum_{j\ge1}(s+\lambda_j)^{-2},\qquad
 \int_0^\infty r(u)\,du=2\sum_j\lambda_j^{-2}<\infty.
\]
The differential equation for $F$ yields
$\Psi''=(\Psi')^2-s/2$, and therefore
\[
 2\Psi''(s)M(s)=M''(s)-\frac{s}{2}M(s).
\]
Taking bilateral Laplace transforms and using their uniqueness gives
\begin{equation}\label{add49heisei:eq:airyconvolution}
 \int_0^\infty r(u)g(x+u)\,du
 =x^2g(x)+\frac12g'(x).
\end{equation}
For completeness, the transform on the left is $2\Psi''M$ by Tonelli's
theorem. On the right, differentiation of $M$ gives $M''$, and integration
by parts gives $\int e^{sx}g'(x)\,dx=-sM(s)$.
The boundary terms vanish: $e^{sx}g(x)$ is integrable and log-concave for
$s$ in a neighborhood of zero, hence tends to zero at both ends and has
finite total variation. Both sides of \eqref{add49heisei:eq:airyconvolution} are
continuous, so the transform identity holds everywhere.

For each $u\ge0$, log-concavity of $g$ makes $g(x+u)/g(x)$ nonincreasing
in $x$. Divide \eqref{add49heisei:eq:airyconvolution} by $g(x)>0$. Its left side is
nonincreasing, and its right side is $x^2+\tfrac12(\log g)'(x)$.
Differentiation gives $0\ge2x-v(x)/2$. Also $v\ge0$ by log-concavity.
Now use \eqref{add49heisei:eq:curvature} and its evenness.
\end{proof}

\subsection{What the compact certificate proves}
\begin{lemma}\label{add49heisei:lem:compact}
The normalized factor in \eqref{add49heisei:eq:normalization} satisfies
\[
 \frac{3405}{1000}<2v(0)<\frac{3406}{1000},\qquad
 v''(x)>\frac12\quad\text{for every }-1\le x\le1.
\]
\end{lemma}
The proof is an integer certificate described in the next three sections.
It encloses entire closed intervals, with every analytic truncation error
included. No conclusion is inferred from an ordinary floating-point grid.

\begin{proof}[Proof of Theorem~\ref{add49heisei:thm:sharp}, using Lemma~\ref{add49heisei:lem:compact}]
For $0<x\le1$, twice integrating $v''>1/2$ gives
\begin{align*}
 \kappa(x)-\kappa(0)
 &=\int_0^x(x-t)\{v''(t)+v''(-t)\}\,dt
 >\frac{x^2}{2}.
\end{align*}
Evenness handles negative $x$. For $|x|\ge1$, Lemma~\ref{add49heisei:lem:tail} gives
$\kappa(x)\ge4|x|\ge4>\kappa(0)$. The value at zero and the numerical
enclosure follow from Lemma~\ref{add49heisei:lem:compact}.
\end{proof}

\begin{corollary}\label{add49heisei:cor:sharpgaussian}
Set $\sigma_0=\kappa(0)^{-1/2}$ and let $\phi$ be the standard normal
density. For $\sigma>0$, the function
\[
 \rho_\sigma(x)=\frac{\sigma f(x)}{\phi(x/\sigma)}
\]
is log-concave if and only if $\sigma\ge\sigma_0$.
For the density $f_a$ of the maximizer of $W(t)-at^2$, $a>0$,
$f_a/\phi$ is log-concave if and only if $a\ge\sigma_0^{3/2}$.
\end{corollary}
\begin{proof}
The first assertion follows from
$-(\log\rho_\sigma)''=\kappa-\sigma^{-2}$ and Theorem~\ref{add49heisei:thm:sharp}.
Brownian scaling gives $f_a(x)=a^{2/3}f(a^{2/3}x)$, so its minimum
negative log-curvature is $a^{4/3}\kappa(0)$, proving the second assertion.
\end{proof}

\subsection{Airy evaluation and two elementary zero bounds}
Let $U,V$ solve $Y''=sY/2$ with
\[
 U(0)=1,\ U'(0)=0,\qquad V(0)=0,\ V'(0)=1.
\]
Then $F=U-dV$ for $d=-F'(0)>0$. For either power series the nonzero
coefficients are generated by
\begin{equation}\label{add49heisei:eq:series}
 a_{j+3}=\frac{a_j}{2(j+3)(j+2)}.
\end{equation}
There is no quadratic seed. To evaluate at $s=n/16$ or $s=in/16$, the
verifier accumulates the appropriate real and imaginary parts, with
outward integer rounding. After an omitted term is bounded by
$2\cdot2^{-230}$ and its ratio to the preceding term is less than $1/2$,
all later ratios are smaller. The entire omitted absolute tail is thus
less than $4\cdot2^{-230}$.

The Airy function is positive and decreasing on the nonnegative real
axis, so $0<F(48)<1$. This also follows from its negative initial slope,
absence of nonnegative zeros, and positive decay at infinity
\cite[\S\S9.2,9.7,9.9]{add49heisei:dlmf}. Consequently
\begin{equation}\label{add49heisei:eq:slope}
 \frac{U(48)-1}{V(48)}<d<\frac{U(48)}{V(48)}.
\end{equation}
The series and direct endpoint division enclose both fractions. The
verifier proves $1/2<d<3/5$, $F(1)>1/4$, and $F(-1)>1/4$.

We need only very coarse bounds for the Airy zeros:
\begin{equation}\label{add49heisei:eq:zeros}
 \lambda_1>1,\qquad \lambda_{16}<32.
\end{equation}
Here are elementary proofs that do not use numerical zero tables.
Write $y(t)=F(-t)$, so $y''+(t/2)y=0$, $y(0)=1$, and $y'(0)=d$.
If its first positive zero were $t_0\le1$, then before that zero
$0<y(t)\le1+dt<8/5$. The integral equation would give
\[
 y(t)=1+dt-\frac12\int_0^t(t-u)u y(u)\,du
 \ge1-\frac{2t^3}{15}>0\quad(0\le t\le t_0),
\]
a contradiction. Thus $\lambda_1>1$.

The following elementary Sturm comparison suffices for the other bound:
if $q\ge\alpha^2$ on an interval longer than $\pi/\alpha$, a nonzero
solution of $y''+qy=0$ has a zero in that interval. Indeed, if it had one
sign, integrate against a positive sine arch of length $\pi/\alpha$
strictly inside the interval. Integration by parts gives positive
endpoint terms, whereas the differential equation gives
$\int(\alpha^2-q)y\sin\le0$, a contradiction.
Apply this once on $(2,8)$, where $q=t/2\ge1$, and once on each of
15 disjoint subintervals of $(8,32)$ of length $8/5>\pi/2$, where
$q\ge4$. These supply 16 distinct zeros below 32.

For $u>-\lambda_1$, the canonical product gives
\begin{equation}\label{add49heisei:eq:modulus}
 |M(u+it)|=M(u)\prod_{j\ge1}
 \left(1+\frac{t^2}{(\lambda_j+u)^2}\right)^{-1/2}.
\end{equation}
In particular $|M(it)|\le1$. Log-convexity of the moment generating
function and the enclosures for $F(\pm1)$ give $M(u)<4$ for $|u|\le1$.
Retaining eight factors and using \eqref{add49heisei:eq:zeros} therefore gives
\begin{equation}\label{add49heisei:eq:stripdecay}
 |M(u+it)|\le4(1+t^2/33^2)^{-4}\qquad(|u|\le1).
\end{equation}

\subsection{Fourier inversion with explicit infinite errors}
Put $D_j(x)=2\pi g^{(j)}(x)$. The product, retaining arbitrarily many
factors, ensures absolute convergence of every differentiated Fourier
integral. Thus
\begin{equation}\label{add49heisei:eq:fourier}
 D_j(x)=\int_{\mathbb R}(-it)^j e^{-itx}M(it)\,dt.
\end{equation}
We require $0\le j\le4$ and $|x|\le1$. Let $h=1/16$, $T=32$, and define
the finite sum
\[
 \widetilde D_j(x)=h\sum_{n=-512}^{512}
       (-inh)^j e^{-inhx}M(inh).
\]
We now prove, uniformly in these $j,x$,
\begin{equation}\label{add49heisei:eq:fouriererror}
 |D_j(x)-\widetilde D_j(x)|<10^{-12}.
\end{equation}

\paragraph{Infinite trapezoidal rule.}
The integrand $H(t)=(-it)^j e^{-itx}M(it)$ is holomorphic in
$|\Im t|\le1$, because $\lambda_1>1$.
On either boundary line, \eqref{add49heisei:eq:stripdecay} and $e<3$ give
\[
 |H(t\pm i)|\le12(1+|t|)^4(1+t^2/33^2)^{-4}.
\]
Integrating separately over $|t|\le33$ and $|t|>33$ bounds each boundary
integral by
\[
 12\left(66\,34^4+\frac23\,33\,34^4\right)
 =12\cdot88\cdot34^4<2\cdot10^9.
\]
Contour shifting bounds the Fourier transform of $H$ at frequency
$\omega$ by the corresponding boundary integral times $e^{-|\omega|}$.
Poisson summation, or equivalently the Fourier series of its absolutely
convergent periodization, then bounds the trapezoidal error by
\[
 E_{\rm strip}
 <\frac{4\cdot10^9}{e^{2\pi/h}-1}
 <\frac{4\cdot10^9}{2^{96}-1}.
\]
The contour shifts are justified by the uniform polynomial decay in
\eqref{add49heisei:eq:stripdecay}; higher-factor versions give any additional decay
needed. The last inequality uses $\pi>3$ and $e>2$.

\paragraph{Omitted frequencies.}
The exact Airy evaluation proves
\[
 |\Re M(32i)|<10^{-25},\qquad |\Im M(32i)|<10^{-25},
\]
and hence $|M(32i)|<2\cdot10^{-25}$.
For $t\ge32$, taking the ratio of the products at $t$ and at 32 and
retaining 16 factors gives
\[
 |M(it)|\le2\cdot10^{-25}\,2^8(32/t)^{16}.
\]
Indeed, each retained squared ratio is at most $2\cdot32^2/t^2$;
all other ratios are at most one.
Since $t^j\le t^4$ here and the power $t^{-12}$ decreases,
\[
 E_{\rm tail}
 \le\frac{1024\cdot32^5}{11\cdot10^{25}}.
\]
This is obtained by bounding
$h\sum_{n\ge513}(nh)^{-12}$ by $\int_{32}^{\infty}t^{-12}\,dt$ and
including both frequency signs. The verifier checks the exact rational
inequality
\[
 \frac{4\cdot10^9}{2^{96}-1}
 +\frac{1024\cdot32^5}{11\cdot10^{25}}<10^{-12},
\]
which proves \eqref{add49heisei:eq:fouriererror}.

\subsection{Interval coverage of the compact domain}
The finite sums are evaluated by Taylor polynomials. The next details
specify all their errors, including those for the derivative orders used
internally by the verifier.

For $0\le k\le10$, expand the exponential in
$\widetilde D_k(c)$ through degree 160, for a real center $|c|\le1$.
The real-variable integral remainder for $e^{-itc}$ has modulus at most
$|tc|^{161}/161!$, since all its derivatives on the real segment have
modulus $|t|^{161}$. Using $|M(it)|\le1$, there are $1025$ terms and
$h=1/16$, so the center error is at most
\begin{equation}\label{add49heisei:eq:centererror}
 E_k^{\rm center}=\frac{1025\cdot32^{k+161}}{16\cdot161!}.
\end{equation}
Each of these is less than $10^{-20}$, as checked by integers.
The coefficients themselves are enclosed using \eqref{add49heisei:eq:series} and
\eqref{add49heisei:eq:slope}, with conjugate Fourier terms combined into real parts.

For a closed cell $[l,r]\subseteq[-1,1]$, let
$c=(l+r)/2$ and $\rho=(r-l)/2$. On that cell, expand
\[
 \widetilde D_j(x)=\sum_{q=0}^{6}
 \frac{\widetilde D_{j+q}(c)}{q!}(x-c)^q+R_j(x),\qquad0\le j\le4.
\]
The elementary finite-sum derivative bound
\[
 |\widetilde D_{j+7}(x)|\le\frac{1025}{16}\,32^{j+7}
\]
gives the uniform local remainder
\begin{equation}\label{add49heisei:eq:localerror}
 |R_j(x)|\le\frac{1025\cdot32^{j+7}}{16\cdot7!}\rho^7.
\end{equation}
Outward interval Horner evaluation of this local polynomial, widened by
\eqref{add49heisei:eq:localerror} and \eqref{add49heisei:eq:fouriererror}, therefore encloses
$D_0,\ldots,D_4$ at every point in the closed cell.

For the resulting derivative intervals form
\[
 (b,c_2,d_3,e_4)=\frac{(D_1,D_2,D_3,D_4)}{D_0}.
\]
The denominator is verified positive in every evaluation. Direct
differentiation gives
\begin{equation}\label{add49heisei:eq:vsecond}
 v=b^2-c_2,\qquad
 v''=-e_4+4bd_3+3c_2^2-12b^2c_2+6b^4.
\end{equation}
The entire computation uses intervals with integer endpoints and common
denominator $2^{230}$; every multiplication and quotient rounds outwards.

\begin{proof}[Proof of Lemma~\ref{add49heisei:lem:compact}]
The complete program below first evaluates the degenerate cell $[0,0]$
and verifies $3405/1000<2v(0)<3406/1000$.
For the remaining assertion it starts with the 512 closed cells of width
$1/256$ covering $[-1,1]$. A cell is accepted only if the lower endpoint
computed from \eqref{add49heisei:eq:vsecond} exceeds $1/2$; otherwise it is bisected.
The resulting 10,687 accepted closed cells cover the entire original
interval. The greatest refinement depth is seven, and every lower
endpoint exceeds $250011/500000>1/2$. All assertions terminate
successfully. The evaluation formulas and errors above prove the
asserted bounds at every point of these cells.
\end{proof}

\subsection{Complete exact verifier}
This Python 3 program uses only the standard library. No auxiliary data,
floating-point values, special-function library, or external certificate
file is required.

\begin{Verbatim}
"""Exact integer certificate for sharp Chernoff curvature. Python 3 only.

The proof supplies the Airy product, strip quadrature estimate, and tails.
This program verifies kappa(0)<4 and v''(x)>1/2 for every -1<=x<=1.
No floating point, numerical quadrature library, or special function is used.
"""
from math import factorial
from fractions import Fraction

S = 1 << 230


def ceildiv(a, b):
    assert b > 0
    return -((-a) // b)


class I:
    def __init__(self, lo, hi=None):
        self.lo = lo
        self.hi = lo if hi is None else hi
        assert self.lo <= self.hi

    def __add__(self, other):
        return I(self.lo + other.lo, self.hi + other.hi)

    def __neg__(self):
        return I(-self.hi, -self.lo)

    def __sub__(self, other):
        return self + (-other)

    def scale(self, n, d=1):
        assert d > 0
        a, b = n * self.lo, n * self.hi
        return I(min(a, b) // d, ceildiv(max(a, b), d))

    def __mul__(self, other):
        products = [a * b for a in (self.lo, self.hi)
                    for b in (other.lo, other.hi)]
        return I(min(products) // S, ceildiv(max(products), S))

    def __truediv__(self, other):
        assert other.lo > 0
        pairs = [(a * S, b) for a in (self.lo, self.hi)
                 for b in (other.lo, other.hi)]
        # Direct endpoint quotients retain precision even for huge U and V.
        return I(min(a // b for a, b in pairs),
                 max(ceildiv(a, b) for a, b in pairs))

    def widen(self, radius):
        assert radius >= 0
        return I(self.lo - radius, self.hi + radius)


def Q(n, d=1):
    assert d > 0
    return I(n * S // d, ceildiv(n * S, d))


ZERO, ONE = Q(0), Q(1)


def basis(seed, n, imaginary):
    """U or V at n/16 or i*n/16, with coefficient and geometric tail bounds."""
    assert seed in (0, 1) and n >= 0
    term = ONE if seed == 0 else Q(n, 16)
    result = [ZERO, ZERO]
    degree = seed
    while True:
        phase = degree % 4 if imaginary else 0
        result[phase % 2] = result[phase % 2] + (
            term if phase < 2 else -term)
        denominator = 2 * (degree + 2) * (degree + 3) * 16 ** 3
        term = term.scale(n ** 3, denominator)
        degree += 3
        # Subsequent ratios decrease. The entire omitted absolute tail <4/S.
        if term.hi <= 2 and 2 * n ** 3 < denominator:
            break
        assert degree < 4000
    return [z.widen(4) for z in result]


U48 = basis(0, 48 * 16, False)[0]
V48 = basis(1, 48 * 16, False)[0]
# F(48)=U(48)-d*V(48) lies strictly between zero and one.
SLOPE = I(((U48 - ONE) / V48).lo, (U48 / V48).hi)
assert SLOPE.lo > Q(1, 2).hi and SLOPE.hi < Q(3, 5).lo


def real_one(sign):
    """F(sign), sign=+/-1. Eighty coefficient terms and absolute tail <4/S."""
    assert sign in (-1, 1)
    parts = []
    for seed in (0, 1):
        term, value, degree = ONE, ZERO, seed
        for _ in range(80):
            value = value + (term if sign == 1 or degree % 2 == 0 else -term)
            term = term.scale(1, 2 * (degree + 2) * (degree + 3))
            degree += 3
        assert term.hi <= 2
        parts.append(value.widen(4))
    return parts[0] - SLOPE * parts[1]


assert real_one(1).lo > S // 4 and real_one(-1).lo > S // 4


def transform(n):
    """Real and imaginary parts of M(i*n/16)=1/F(i*n/16)."""
    u, v = basis(0, n, True), basis(1, n, True)
    re, im = u[0] - SLOPE * v[0], u[1] - SLOPE * v[1]
    denominator = re * re + im * im
    return re / denominator, (-im) / denominator


M = [transform(n) for n in range(513)]
assert all(-(S // (10 ** 25)) <= z.lo <= z.hi <= S // (10 ** 25)
           for z in M[-1])

# Analytically justified errors for D_j, 0<=j<=4, on [-1,1].
STRIP_ERROR = Fraction(4 * 10 ** 9, 2 ** 96 - 1)
TAIL_ERROR = Fraction(1024 * 32 ** 5, 11 * 10 ** 25)
assert STRIP_ERROR + TAIL_ERROR < Fraction(1, 10 ** 12)
FOURIER_ERROR = Q(1, 10 ** 12).hi

# Coefficients for the Taylor polynomial of the FINITE Fourier sum.
# Derivatives up to order 10 allow well-conditioned local expansions.
DEGREE = 160
POLY = [[ZERO for _ in range(DEGREE + 1)] for _ in range(11)]
for n, (re, im) in enumerate(M):
    for k in range(11):
        factor = Q((1 if n == 0 else 2) * n ** k, 16 ** (k + 1))
        for j in range(DEGREE + 1):
            phase = (re, im, -re, -im)[(k + j) % 4]
            POLY[k][j] = POLY[k][j] + factor * phase
            factor = factor.scale(n, 16 * (j + 1))

CENTER_ERRORS = [Q(1025 * 32 ** (k + 161), 16 * factorial(161)).hi
                 for k in range(11)]
assert max(CENTER_ERRORS) < Q(1, 10 ** 20).lo


def centers(center):
    """Enclose all eleven derivatives of the finite Fourier sum at a point."""
    x = I(center)
    result = []
    for k, p in enumerate(POLY):
        value = p[-1]
        for coefficient in reversed(p[:-1]):
            value = value * x + coefficient
        result.append(value.widen(CENTER_ERRORS[k]))
    return result


def derivatives(lo, hi):
    """Enclose D_0,...,D_4 over an entire interval, by local Taylor degree 6."""
    assert -S <= lo <= hi <= S
    center = (lo + hi) // 2
    delta = I(lo - center, hi - center)
    radius = max(center - lo, hi - center)
    c = centers(center)
    result = []
    for k in range(5):
        value = c[k + 6].scale(1, factorial(6))
        for j in reversed(range(6)):
            value = value * delta + c[k + j].scale(1, factorial(j))
        # |finite D_(k+7)| <= (1025/16)*32**(k+7).
        remainder = ceildiv(1025 * 32 ** (k + 7) * radius ** 7,
                            16 * factorial(7) * S ** 6)
        result.append(value.widen(remainder + FOURIER_ERROR))
    return result


def curvature_values(ds):
    assert ds[0].lo > 0
    b, c, d, e = [z / ds[0] for z in ds[1:]]
    v = b * b - c
    vv = (-e + (b * d).scale(4) + (c * c).scale(3)
          - (b * b * c).scale(12) + (b * b * b * b).scale(6))
    return v, vv


def main():
    v0, _ = curvature_values(derivatives(0, 0))
    kappa0 = v0.scale(2)
    assert kappa0.lo > Q(3405, 1000).hi
    assert kappa0.hi < Q(3406, 1000).lo
    print('CERTIFIED: 3405/1000 < kappa(0) < 3406/1000 < 4.', flush=True)
    completed = []
    maximum_depth = 0

    def check(lo, hi, depth):
        nonlocal maximum_depth
        _, value = curvature_values(derivatives(lo, hi))
        if value.lo > S // 2:
            completed.append(value.lo)
            maximum_depth = max(maximum_depth, depth)
            return
        assert hi - lo > S // (1 << 22)
        mid = (lo + hi) // 2
        check(lo, mid, depth + 1)
        check(mid, hi, depth + 1)

    for i in range(-256, 256):
        check(i * S // 256, (i + 1) * S // 256, 0)
    lower = Fraction(min(completed) * 10 ** 6 // S, 10 ** 6)
    print(f'CERTIFIED: v-second(x)>1/2 for every -1<=x<=1; '
          f'{len(completed)} closed cells; max refinement depth {maximum_depth}.')
    print(f'All interval lower endpoints exceed {lower}. Integer arithmetic only.')


if __name__ == '__main__':
    main()
\end{Verbatim}


\begin{thebibliography}{9}
\bibitem{add49heisei:bw} F. Balabdaoui and J. A. Wellner,
\emph{Chernoff's density is log-concave}, Bernoulli \textbf{20} (2014),
231--244.
\href{https://doi.org/10.3150/12-BEJ483}{doi:10.3150/12-BEJ483}.
\href{https://pmc.ncbi.nlm.nih.gov/articles/PMC3993999/}{Open full text}.
\bibitem{add49heisei:zz} X. Zhang and Q. Zhou,
\emph{Chernoff's Density Is Strongly Log-Concave} (2026),
\href{https://arxiv.org/abs/2607.18619}{arXiv:2607.18619v1}, Lemma 7.
\href{https://pdfs.assets.alphaxiv.org/2607.18619v1.pdf}{Original-paper PDF mirror}.
\bibitem{add49heisei:dlmf} NIST Digital Library of Mathematical Functions,
\href{https://dlmf.nist.gov/9.2}{\S9.2},
\href{https://dlmf.nist.gov/9.7}{\S9.7}, and
\href{https://dlmf.nist.gov/9.9}{\S9.9}: Airy's equation, decay, and zeros.
\end{thebibliography}

\par\endgroup

% Additional source: /workspace/scratch/81914cec2c78/work/additional_geometry/r49_full_proof.tex
\begingroup

\clearpage
\section{\texorpdfstring{A full proof of the Liu--Chu Fermat-distance conjecture}{A full proof of the Liu--Chu Fermat-distance conjecture}}\label{section:add49keio}
\noindent{\small\sffamily\color{muted}Additional completion \enspace / \enspace Round 49 \enspace / \enspace Keio}\par\medskip


\begin{sectionabstract}
We prove Conjecture 4, equation (30), of Liu and Chu (2003), for every
nondegenerate triangle and every point of its plane.  The proof separates
the arbitrary point from the triangle by a complex-coordinate lower bound.
A rational half-angle parametrization reduces the remaining shape
inequality to two integer polynomials. Five explicit cones have
nonnegative coefficients and cover the complete parameter domain.
The exact coefficient verifier is included in full.
\end{sectionabstract}

\subsection{The original assertion and its scope}

Let $ABC$ be any nondegenerate triangle, with side lengths
$a=BC$, $b=CA$, $c=AB$, semiperimeter $s$, area $\Delta$, and
inradius $r$.  Let $m_b,m_c$ and $h_b,h_c$ be the median lengths
and altitude lengths from $B,C$, respectively.
Liu and Chu \cite[p.~92, Conjecture 4, equation (30)]{add49keio:lc}
proposed the following assertion with $P$ an arbitrary point of the
plane, without an interior-point restriction.

\begin{theorem}\label{add49keio:full}
For every nondegenerate triangle $ABC$ and every point $P$ in its plane,
\[
 PA+PB+PC\ \ge\ 2r\left(1+\frac{m_b}{h_b}+\frac{m_c}{h_c}\right).
\]
Equality holds if and only if the triangle is equilateral and $P$ is its
center.
\end{theorem}

The archived Round 49 argument established this only when $B$ and $C$
straddle $60^\circ$.  The present proof includes that argument's
point-independent reduction, proves all remaining shape regions, and
therefore does not depend on a separate earlier manuscript. This is a
completion relative to that archive; no claim of historical priority is made.

\subsection{Removing the arbitrary point}\label{add49keio:point}

Put
\[
 d=a^2+b^2+c^2,\qquad
 E=(a^2-b^2)^2+(a^2-c^2)^2,\qquad
 F=\sqrt{d/2+2\sqrt3\,\Delta}.
\]
We first have $PA+PB+PC\ge F$ for all plane points $P$.
Indeed, represent the plane by complex numbers and let $\omega$ be a
nonreal cube root of unity. Since $1+\omega+\omega^2=0$,
\[
 |A+\omega B+\omega^2C|
 =|(A-P)+\omega(B-P)+\omega^2(C-P)|
 \le PA+PB+PC.
\]
Direct expansion for the two choices of $\omega$ gives the two squared
moduli $d/2\pm2\sqrt3\,\Delta$.  Taking the larger proves the claim.
This argument applies even when a triangle angle is at least $120^\circ$.

The median formula and Heron's identity give
\[
 4b^2m_b^2=16\Delta^2+(a^2-c^2)^2,\qquad
 4c^2m_c^2=16\Delta^2+(a^2-b^2)^2.
\]
Consequently the two-term Cauchy--Schwarz inequality yields
\[
 bm_b+cm_c\le \sqrt{16\Delta^2+E/2}.
\]
Because $rs=\Delta$, $h_b=2\Delta/b$ and $h_c=2\Delta/c$,
the right side in Theorem~\ref{add49keio:full}, multiplied by $s$, is
$2\Delta+bm_b+cm_c$.  It is therefore sufficient to prove the
stronger shape inequality
\begin{equation}\label{add49keio:shape}
 sF\ \ge\ 2\Delta+\sqrt{16\Delta^2+E/2}.
\end{equation}

For later use write $J=\sqrt3\,\Delta$ and set
\begin{align}
 L&=3s^2(d+4J)-40J^2-3E,\label{add49keio:Ldef}\\
 Q&=L^2-32J^2(32J^2+3E).\label{add49keio:Qdef}
\end{align}
These definitions give
\[
 L=6\bigl(s^2F^2-20\Delta^2-E/2\bigr),\qquad
 Q=L^2-\bigl(24\Delta\sqrt{16\Delta^2+E/2}\bigr)^2.
\]
Since $s,F,\Delta$ and the radical are positive, condition
\eqref{add49keio:shape} is equivalent to $L\ge0$ and $Q\ge0$.
In detail, these two signs give
$L\ge24\Delta\sqrt{16\Delta^2+E/2}$, hence
\[
 s^2F^2\ge
 20\Delta^2+E/2+4\Delta\sqrt{16\Delta^2+E/2}
 =\bigl(2\Delta+\sqrt{16\Delta^2+E/2}\bigr)^2.
\]
Taking the positive square root proves \eqref{add49keio:shape}; reversing these
steps proves necessity.  If $Q>0$ and $L>0$, all these inequalities are
strict at the final shape comparison.

\subsection{A parameter domain covering every triangle}

Among the three angles there is one at most $60^\circ$ and another at
least $60^\circ$. Choose such a pair and label their parameters
\[
 p=\sqrt3\tan(\alpha/2),\qquad
 q=\sqrt3\tan(\beta/2),\qquad
 t=\sqrt3\tan(\gamma/2),
\]
where $\alpha+\beta+\gamma=\pi$ and $\alpha\le\pi/3\le\beta$.
The half-angle addition formula gives
\[
 0<p\le1\le q,\qquad pq<3,\qquad
 t=\frac{3-pq}{p+q}>0.
\]
The complete range of $p,q$ has the parametrization
\begin{equation}\label{add49keio:param}
 p=\frac1{1+x},\qquad
 q=\frac{1+3(1+x)y}{1+y},\qquad x,y\ge0.
\end{equation}
Indeed its inverse is
$x=1/p-1$ and $y=(q-1)/(3/p-q)$, whose denominator is strictly
positive. Conversely every finite $x,y\ge0$ gives the asserted ranges
and hence a nondegenerate triangle. Ties at $60^\circ$ are allowed.

Use the polynomial numerators and denominators
\[
 p_1=1,\quad p_2=1+x,\quad q_1=1+3(1+x)y,\quad q_2=1+y,
\]
\[
 t_1=3p_2q_2-p_1q_1=2+3x,\qquad
 t_2=p_1q_2+p_2q_1,
\]
so that $p=p_1/p_2$, $q=q_1/q_2$, $t=t_1/t_2$.
All these quantities are positive on the indicated domain. Set
\begin{equation}\label{add49keio:contact}
 T=p_1q_1t_1,\quad
 U=p_2q_1t_1,\quad V=p_1q_2t_1,\quad W=p_1q_1t_2,
 \quad S=U+V+W.
\end{equation}
By similarity normalize the inradius to $T/\sqrt3$.
The three contact lengths are then
$r\cot(\alpha/2)=T/p=U$, $T/q=V$, and $T/t=W$.
Thus $s=S$, $J=ST$, and the respective squared sides are
\[
 A_0=(S-U)^2,\quad A_1=(S-V)^2,\quad A_2=(S-W)^2.
\]
The identity $3UVW=ST^2$, also checked below, agrees with Heron's
formula for this normalization.

The distinguished angle $A$ in Theorem~\ref{add49keio:full} can occupy any
index $i\in\{0,1,2\}$. Set $a^2=A_i$ and let $b^2,c^2$ be the
other two values. Their order is irrelevant since $E$ is symmetric in
$b,c$. For each index, definitions \eqref{add49keio:Ldef}--\eqref{add49keio:Qdef} now
give polynomials $L_i,Q_i\in\mathbb Z[x,y]$. Thus all assignments
of the original vertex labels are included.

\subsection{Five exact cones and their certificates}

The following substitutions use independent $u,v\ge0$:
\[
\begin{array}{c|c|c}
 i&\text{region in }(x,y)&(x,y)\\\hline
 0&x\le3y&(3u,u+v)\\
 0&x\ge3y&(3u+v,u)\\
 1&x\le y&(u,u+v)\\
 1&x\ge y&(u+v,u)\\
 2&\text{all}&(u,v)
\end{array}
\]
Each row parametrizes exactly its stated closed cone: for example the
first row has $u=x/3$, $v=y-x/3$, and the second has $u=y$,
$v=x-3y$. The analogous inverses for $i=1$ are immediate.
The rows therefore cover the whole nonnegative quadrant for every
possible index $i$, including their boundaries.

In each row, exact expansion gives strictly positive coefficients for
every nonzero term of both $L_i$ and $Q_i$. The following table records
the number of nonzero terms and the minimum nonzero coefficient.
\begin{center}
\begin{tabular}{ccrrrr}
\toprule
$i$&region&$\#L_i$&$\min L_i$&$\#Q_i$&$\min Q_i$\\
\midrule
0&$x\le3y$&153&4608&558&116785152\\
0&$x\ge3y$&180&1620&667&1364688\\
1&$x\le y$&153&4608&558&47775744\\
1&$x\ge y$&184&243&685&59049\\
2&all&94&243&336&59049\\
\bottomrule
\end{tabular}
\end{center}
Here $\min$ refers to coefficients, not values of the polynomial.
In every row the constant coefficient of $L_i$ is $4608$, while the
constant coefficient of $Q_i$ is zero. The coefficients of $u^2,v^2$
in $Q_i$, in row order, are
\[
 (122093568,233570304),\quad (122093568,26542080),\quad
 (47775744,106168320),
\]
\[
 (47775744,58392576),\qquad (26542080,106168320).
\]
Thus $L_i>0$ on every cone, and $Q_i>0$ unless $u=v=0$.
This proves \eqref{add49keio:shape} for every triangle, strictly unless
$x=y=0$. By \eqref{add49keio:param} the exceptional shape has
$p=q=t=1$, so it is exactly the equilateral triangle.

\subsection{Equality and completion}

For a nonequilateral triangle the strict version of \eqref{add49keio:shape},
together with the two bounds in Section~\ref{add49keio:point}, gives a strict inequality
in Theorem~\ref{add49keio:full} for every plane point $P$.
For an equilateral triangle of side length $a$, the right side is
$6r=\sqrt3a$. Let $G$ be its center, let $R=a/\sqrt3$, and put
$e_A=(A-G)/R$, cyclically. These three unit vectors sum to zero.
For every $P$, scalar projection gives
\[
 PA+PB+PC\ge
 (A-P)\mathbin{\cdot}e_A+(B-P)\mathbin{\cdot}e_B
 +(C-P)\mathbin{\cdot}e_C=3R=\sqrt3a.
\]
Equality requires each of the three vectors to be a nonnegative
multiple of its corresponding $e$ vector. In particular $P$ lies on
all three lines $AG,BG,CG$, whose common point is $G$.
Conversely $P=G$ attains equality. This completes the full theorem.

\subsection{Complete exact verifier}

The listing uses only Python 3 and SymPy's polynomial ring over
$\mathbb Z$. It constructs the formulas above and checks every
coefficient of each complete expanded polynomial. It does not evaluate
a sample of triangles or omit an unbounded part of a parameter cone.
\begin{Verbatim}
"""Exact five-chart certificate for Liu--Chu (2003), Conjecture 4.

Requires Python 3 and SymPy.  Every assertion is an integer polynomial
identity or coefficient check, and applies on an entire nonnegative cone.
"""
from sympy import ZZ
from sympy.polys.rings import ring

K, u, v = ring('u,v', ZZ)
charts = [
    (0, 'x<=3y', 3*u, u+v),
    (0, 'x>=3y', 3*u+v, u),
    (1, 'x<=y', u, u+v),
    (1, 'x>=y', u+v, u),
    (2, 'all', u, v),
]
# (#L, min nonzero L, #Q, min nonzero Q, [u^2]Q, [v^2]Q)
expected = [
    (153, 4608, 558, 116785152, 122093568, 233570304),
    (180, 1620, 667, 1364688, 122093568, 26542080),
    (153, 4608, 558, 47775744, 47775744, 106168320),
    (184, 243, 685, 59049, 47775744, 58392576),
    (94, 243, 336, 59049, 26542080, 106168320),
]
for chart, wanted in zip(charts, expected):
    i, name, x, y = chart
    p1, p2 = K.one, 1+x
    q1, q2 = 1+3*p2*y, 1+y
    t1 = 3*p2*q2-p1*q1
    t2 = p1*q2+p2*q1
    U, V, W = p2*q1*t1, p1*q2*t1, p1*q1*t2
    T = p1*q1*t1
    S, J = U+V+W, (U+V+W)*T
    assert 3*U*V*W == S*T*T
    A = [(S-z)**2 for z in (U, V, W)]
    a = A[i]
    b, c = [A[j] for j in range(3) if j != i]
    E = (a-b)**2+(a-c)**2
    L = 3*S*S*(a+b+c+4*J)-40*J*J-3*E
    Q = L*L-32*J*J*(32*J*J+3*E)
    assert min(L.values()) > 0
    assert min(Q.values()) > 0
    assert L[(0, 0)] == 4608
    assert Q.get((0, 0), 0) == 0
    assert Q[(2, 0)] > 0 and Q[(0, 2)] > 0
    observed = (len(L), min(L.values()), len(Q), min(Q.values()),
                Q[(2, 0)], Q[(0, 2)])
    assert observed == wanted
    print(i, name, observed)
print('PASS: all triangle shapes; equality shape only x=y=0.')
\end{Verbatim}


\begin{thebibliography}{9}
\bibitem{add49keio:lc} J. Liu and X.-G. Chu,
\emph{A New Geometry Inequality with the Fermat Problem},
Journal of East China Jiaotong University \textbf{20}(1) (2003),
89--93, Conjecture 4, equation (30), p.~92.
\url{https://journal.ecjtu.edu.cn/hdjtdxen/article/pdf/20030127}.
\end{thebibliography}

\par\endgroup

\end{document}
